% Options for packages loaded elsewhere
\PassOptionsToPackage{unicode}{hyperref}
\PassOptionsToPackage{hyphens}{url}
\documentclass[
]{article}
\usepackage{xcolor}
\usepackage[margin=1in]{geometry}
\usepackage{amsmath,amssymb}
\setcounter{secnumdepth}{5}
\usepackage{iftex}
\ifPDFTeX
  \usepackage[T1]{fontenc}
  \usepackage[utf8]{inputenc}
  \usepackage{textcomp} % provide euro and other symbols
\else % if luatex or xetex
  \usepackage{unicode-math} % this also loads fontspec
  \defaultfontfeatures{Scale=MatchLowercase}
  \defaultfontfeatures[\rmfamily]{Ligatures=TeX,Scale=1}
\fi
\usepackage{lmodern}
\ifPDFTeX\else
  % xetex/luatex font selection
\fi
% Use upquote if available, for straight quotes in verbatim environments
\IfFileExists{upquote.sty}{\usepackage{upquote}}{}
\IfFileExists{microtype.sty}{% use microtype if available
  \usepackage[]{microtype}
  \UseMicrotypeSet[protrusion]{basicmath} % disable protrusion for tt fonts
}{}
\makeatletter
\@ifundefined{KOMAClassName}{% if non-KOMA class
  \IfFileExists{parskip.sty}{%
    \usepackage{parskip}
  }{% else
    \setlength{\parindent}{0pt}
    \setlength{\parskip}{6pt plus 2pt minus 1pt}}
}{% if KOMA class
  \KOMAoptions{parskip=half}}
\makeatother
\usepackage{graphicx}
\makeatletter
\newsavebox\pandoc@box
\newcommand*\pandocbounded[1]{% scales image to fit in text height/width
  \sbox\pandoc@box{#1}%
  \Gscale@div\@tempa{\textheight}{\dimexpr\ht\pandoc@box+\dp\pandoc@box\relax}%
  \Gscale@div\@tempb{\linewidth}{\wd\pandoc@box}%
  \ifdim\@tempb\p@<\@tempa\p@\let\@tempa\@tempb\fi% select the smaller of both
  \ifdim\@tempa\p@<\p@\scalebox{\@tempa}{\usebox\pandoc@box}%
  \else\usebox{\pandoc@box}%
  \fi%
}
% Set default figure placement to htbp
\def\fps@figure{htbp}
\makeatother
\setlength{\emergencystretch}{3em} % prevent overfull lines
\providecommand{\tightlist}{%
  \setlength{\itemsep}{0pt}\setlength{\parskip}{0pt}}
\usepackage[utf8]{inputenc}
\usepackage[bookmarks=false]{hyperref}
\usepackage{xcolor}
\usepackage{pifont}
\usepackage{longtable}
\usepackage{bookmark}
\IfFileExists{xurl.sty}{\usepackage{xurl}}{} % add URL line breaks if available
\urlstyle{same}
\hypersetup{
  pdftitle={Relatório de Cobertura CNCT--CBO},
  pdfauthor={Equipe BM -- FGV/DGPE},
  hidelinks,
  pdfcreator={LaTeX via pandoc}}

\title{Relatório de Cobertura CNCT--CBO}
\usepackage{etoolbox}
\makeatletter
\providecommand{\subtitle}[1]{% add subtitle to \maketitle
  \apptocmd{\@title}{\par {\large #1 \par}}{}{}
}
\makeatother
\subtitle{Listagem de CBOs que não aparecem na seleção feita pela IA}
\author{Equipe BM -- FGV/DGPE}
\date{17/07/2025}

\begin{document}
\maketitle

\textbf{CBOs correspondidos e seus perfis (presentes no CNCT, e presentes na predição da IA)}

\section*{Tabela de Correspondências Positivas}

\textbf{CBOs correspondidos e seus perfis (presentes no CNCT, e presentes na predição da IA)}

\begin{longtable}{p{2cm}p{6cm}p{5cm}p{2cm}}
\caption{Cursos com CBOs Correspondidos no IA} \\ 
\hline
\textbf{Código} & \textbf{Curso CNCT} & \textbf{CBOs correspondidos} & \textbf{Cobertura CNCT por IA} \\ 
\hline
\endfirsthead
\hline
\textbf{Código} & \textbf{Curso CNCT} & \textbf{CBOs correspondidos} & \textbf{Cobertura CNCT por IA} \\ 
\hline
\endhead
\hline
\endfoot
\hline
\endlastfoot
11004 & Técnico em Agropecuária & 321110 & 1.00 \\ 
01013 & Técnico em Hemoterapia & 324220 & 1.00 \\ 
01015 & Técnico em Imobilizações Ortopédicas & 322605 & 1.00 \\ 
11001 & Técnico em Agricultura & 321105 & 1.00 \\ 
01008 & Técnico em Equipamentos Biomédicos & 915305 & 1.00 \\ 
01020 & Técnico em Nutrição e Dietética & 325210 & 1.00 \\ 
11008 & Técnico em Florestas & 321205, 321210 & 1.00 \\ 
01021 & Técnico em Óptica & 322305 & 1.00 \\ 
10015 & Técnico em Planejamento e Controle da Produção & 391125 & 1.00 \\ 
01023 & Técnico em Órteses e Próteses & 322505 & 1.00 \\ 
01027 & Técnico em Radiologia & 324115 & 1.00 \\ 
01030 & Técnico em Saúde Bucal & 322405, 322425, 322430 & 1.00 \\ 
11011 & Técnico em Mineração & 316305, 316315, 316320, 316330 & 1.00 \\ 
10004 & Técnico em Biotecnologia & 325305, 325310 & 1.00 \\ 
11005 & Técnico em Apicultura & 613405 & 1.00 \\ 
01018 & Técnico em Meteorologia & 311510 & 1.00 \\ 
02003 & Técnico em Eletromecânica & 300305 & 1.00 \\ 
02019 & Técnico em Mecânica de Precisão & 314105, 915105 & 1.00 \\ 
11017 & Técnico em Zootecnia & 323105 & 1.00 \\ 
06012 & Técnico em Saneamento & 312210 & 1.00 \\ 
10014 & Técnico em Petroquímica & 311205, 316325 & 1.00 \\ 
01009 & Técnico em Estética & 322130 & 1.00 \\ 
01003 & Técnico em Citopatologia & 324215 & 1.00 \\ 
01001 & Técnico em Agente Comunitário de Saúde & 322255 & 1.00 \\ 
01025 & Técnico em Podologia & 322110 & 1.00 \\ 
01010 & Técnico em Farmácia & 325105, 325110, 325115 & 1.00 \\ 
02017 & Técnico em Manutenção de Sistemas Metroferroviários & 914305 & 1.00 \\ 
08007 & Técnico em Viticultura e Enologia & 325005 & 1.00 \\ 
05009 & Técnico em Telecomunicações & 313305, 313310, 313315, 313320, 731305 & 0.83 \\ 
08004 & Técnico em Confeitaria & 848310 & 1.00 \\ 
10007 & Técnico em Cerâmica & 311305, 752305 & 1.00 \\ 
01002 & Técnico em Análises Clínicas & 324205 & 1.00 \\ 
11002 & Técnico em Agroecologia & 321105 & 1.00 \\ 
01026 & Técnico em Prótese Dentária & 322410 & 1.00 \\ 
05005 & Técnico em Manutenção e Suporte em Informática & 313220, 317210 & 1.00 \\ 
10017 & Técnico em Processamento da Madeira & 321205 & 1.00 \\ 
01031 & Técnico em Terapias Holísticas & 322125 & 1.00 \\ 
02014 & Técnico em Manutenção de Máquinas Industriais & 314410, 911305 & 1.00 \\ 
13004 & Técnico em Guia de Turismo & 511405 & 1.00 \\ 
02020 & Técnico em Mecatrônica & 300105, 300110 & 1.00 \\ 
01022 & Técnico em Optometria & 322305 & 1.00 \\ 
01004 & Técnico em Controle Ambiental & 311505 & 1.00 \\ 
04013 & Técnico em Secretariado & 351505 & 1.00 \\ 
01016 & Técnico em Massoterapia & 322120 & 1.00 \\ 
11010 & Técnico em Geologia & 316105, 316110, 316115, 316120, 316320 & 1.00 \\ 
02001 & Técnico em Automação Industrial & 300105 & 0.50 \\ 
10009 & Técnico em Curtimento & 311115, 760205, 762205, 762210, 762215, 762220 & 1.00 \\ 
03003 & Técnico em Biblioteconomia & 371105, 371110 & 1.00 \\ 
10020 & Técnico em Têxtil & 311605, 311610, 311615, 311620, 311625, 311710, 311720, 318420, 352320, 760105, 760120, 760125, 761105, 761110, 761810, 761815 & 0.80 \\ 
10016 & Técnico em Plásticos & 311405, 311410 & 1.00 \\ 
01007 & Técnico em Enfermagem & 322205, 322210, 322215, 322220, 322225, 322230, 322240, 322245, 322250 & 0.90 \\ 
01032 & Técnico em Veterinária & 519305 & 1.00 \\ 
04014 & Técnico em Seguros & 351710, 351735, 351740 & 1.00 \\ 
02009 & Técnico em Instrumentação Industrial & 313410 & 0.33 \\ 
02012 & Técnico em Manutenção Aeronáutica em Grupo Motopropulsor & 314310, 914105 & 1.00 \\ 
05002 & Técnico em Desenvolvimento de Sistemas & 317110 & 1.00 \\ 
10019 & Técnico em Química & 301105, 301110, 301115, 311105, 311205 & 1.00 \\ 
02015 & Técnico em Manutenção de Máquinas Navais & 314315, 914205 & 1.00 \\ 
06002 & Técnico em Agrimensura & 312305, 312310, 312315, 312320, 318110 & 1.00 \\ 
11006 & Técnico em Aquicultura & 321305, 321310, 321315, 321320 & 1.00 \\ 
02005 & Técnico em Eletrotécnica & 313105, 313110, 313115, 313120, 313125, 313130 & 0.86 \\ 
10018 & Técnico em Processos Gráficos & 318405, 371305, 371310, 760605 & 1.00 \\ 
02004 & Técnico em Eletrônica & 313205, 313210, 313215, 313220 & 1.00 \\ 
02018 & Técnico em Mecânica & 314105, 314110, 915105 & 1.00 \\ 
01005 & Técnico em Cuidados de Idosos & 516210 & 1.00 \\ 
02006 & Técnico em Fabricação Mecânica & 314205, 314210, 314610 & 0.75 \\ 
01017 & Técnico em Meio Ambiente & 311505 & 1.00 \\ 
10008 & Técnico em Construção Naval & 318215 & 0.33 \\ 
11009 & Técnico em Fruticultura & 321105, 325205 & 1.00 \\ 
04005 & Técnico em Contabilidade & 351105, 351110, 351115 & 1.00 \\ 
10021 & Técnico em Vestuário & 319105, 319110 & 1.00 \\ 
04003 & Técnico em Comércio Exterior & 342105, 342210, 351310, 354305 & 1.00 \\ 
10013 & Técnico em Petróleo e Gás & 301115, 316310, 316325 & 1.00 \\ 
06006 & Técnico em Estradas & 312205, 312320 & 1.00 \\ 
12005 & Técnico em Segurança do Trabalho & 351605, 351610 & 1.00 \\ 
06007 & Técnico em Geodésia e Cartografia & 312305, 312310, 312315, 318110 & 1.00 \\ 
02022 & Técnico em Metrologia & 313405, 352305, 352310, 352315 & 1.00 \\ 
13001 & Técnico em Agenciamento de Viagem & 354810, 354815 & 1.00 \\ 
02023 & Técnico em Outros - Eixo Controle e Processos Industriais & 721105, 721110 & 1.00 \\ 
12004 & Técnico em Prevenção e Combate a Incêndios & 510305, 517105, 517110 & 1.00 \\ 
07015 & Técnico em Mecânica de Aeronaves  & 314310, 914105, 953105 & 1.00 \\ 
10005 & Técnico em Calçados & 318815, 319105, 319110 & 1.00 \\ 
04001 & Técnico em Administração & 351305 & 1.00 \\ 
06011 & Técnico em Portos & 342210, 342605, 342610, 414215 & 0.80 \\ 
13003 & Técnico em Gastronomia & 513205 & 1.00 \\ 
09035 & Técnico em Teatro & 374205, 374210, 374215 & 1.00 \\ 
09024 & Técnico em Museologia & 371205, 371210 & 1.00 \\ 
02024 & Técnico em Refrigeração e Climatização & 725705 & 1.00 \\ 
13006 & Técnico em Lazer & 371410 & 1.00 \\ 
02010 & Técnico em Manutenção Aeronáutica em Aviônicos & 314310, 914105 & 1.00 \\ 
09005 & Técnico em Cenografia & 374205, 374210, 374215 & 1.00 \\ 
01029 & Técnico em Registros e Informações em Saúde & 415305 & 1.00 \\ 
02011 & Técnico em Manutenção Aeronáutica em Célula & 314310, 914105, 914110 & 1.00 \\ 
06008 & Técnico em Geoprocessamento & 312305, 312310, 312315, 312320, 318110 & 1.00 \\ 
08002 & Técnico em Alimentos & 325205, 818110 & 1.00 \\ 
05003 & Técnico em Informática & 317105, 317110, 317205, 317210 & 0.67 \\ 
11013 & Técnico em Pesca & 341220, 341225 & 1.00 \\ 
10011 & Técnico em Móveis & 318425, 318805, 319205 & 1.00 \\ 
03012 & Técnico em Secretaria Escolar & 351505 & 1.00 \\ 
09009 & Técnico em Design de Calçados & 318815 & 1.00 \\ 
02027 & Técnico em Soldagem & 314620, 724305, 724310, 724315, 724320, 724325 & 1.00 \\ 
10006 & Técnico em Celulose e Papel & 311105, 311110, 311705 & 0.50 \\ 
05008 & Técnico em Redes de Computadores & 313220, 317210 & 1.00 \\ 
02021 & Técnico em Metalurgia & 314620, 314705, 314710, 314720, 314725, 314730 & 1.00 \\ 
09015 & Técnico em Design Gráfico & 318405, 318410, 318415, 318420, 318425, 318430, 371305, 371310 & 1.00 \\ 
04018 & Técnico em Vendas & 354135, 354140, 354705, 521115 & 1.00 \\ 
10001 & Técnico em Açúcar e Álcool & 311105, 311110 & 0.40 \\ 
09022 & Técnico em Modelagem do Vestuário & 318810 & 1.00 \\ 
02002 & Técnico em Eletroeletrônica & 300305, 731135, 731150 & 1.00 \\ 
02025 & Técnico em Sistemas a Gás & 724115, 724130 & 1.00 \\ 
02013 & Técnico em Manutenção Automotiva & 314305, 314405 & 1.00 \\ 
09018 & Técnico em Fabricação de Instrumentos Musicais & 742105, 742110, 742115, 742120, 742125, 742130, 742135, 742140 & 1.00 \\ 
04011 & Técnico em Qualidade & 391205 & 1.00 \\ 
06018 & Técnico em Transporte Rodoviário & 342305 & 0.50 \\ 
01019 & Técnico em Necropsia & 328105, 328110, 516505 & 1.00 \\ 
11016 & Técnico em Recursos Pesqueiros & 321305, 321315, 321320, 341225 & 0.67 \\ 
01033 & Técnico em Vigilância em Saúde & 352210 & 0.50 \\ 
03007 & Técnico em Laboratório de Ciências da Natureza & 818110 & 1.00 \\ 
04009 & Técnico em Marketing & 354140 & 0.50 \\ 
09003 & Técnico em Artesanato & 752105, 791105, 791110, 791115, 791120, 791125, 791130, 791135, 791140, 791145, 791150, 791155, 791160 & 1.00 \\ 
09023 & Técnico em Multimídia & 317105, 317120 & 0.67 \\ 
02016 & Técnico em Manutenção de Máquinas Pesadas & 913105, 913110, 913115, 913120, 914420 & 1.00 \\ 
09001 & Técnico em Artes Circenses & 376205, 376210, 376215, 376220, 376225, 376230, 376235, 376240, 376245, 376250, 376255, 376325 & 0.75 \\ 
02008 & Técnico em Fundição & 314715, 721325, 722215, 722220, 722325, 771110, 821220, 822115 & 1.00 \\ 
13008 & Técnico em Serviços de Restaurante e Bar & 510130, 510135, 513405, 513415 & 1.00 \\ 
09008 & Técnico em Dança & 376105, 376110 & 1.00 \\ 
06001 & Técnico em Aeroportuário & 342535, 342545 & 1.00 \\ 
13005 & Técnico em Hospedagem & 422120, 513110, 513115 & 1.00 \\ 
04017 & Técnico em Transações Imobiliárias & 354410, 354605 & 1.00 \\ 
09011 & Técnico em Design de Interiores & 318015, 375105, 375110, 375115, 375120 & 0.71 \\ 
09030 & Técnico em Produção de Áudio e Vídeo & 372115, 373105, 373220, 374105, 374110, 374135, 374210, 374405, 374415, 766120 & 0.83 \\ 
04002 & Técnico em Comércio & 354140, 521105, 521115 & 0.60 \\ 
02026 & Técnico em Sistemas de Energia Renovável & 313110, 861105 & 1.00 \\ 
04008 & Técnico em Logística & 391125 & 0.33 \\ 
09033 & Técnico em Rádio e Televisão & 373105, 373220, 374105, 374125, 374130, 374145 & 0.55 \\ 
11007 & Técnico em Cafeicultura & 321105, 325205 & 1.00 \\ 
08006 & Técnico em Panificação & 513505, 848305 & 1.00 \\ 
08003 & Técnico em Cervejaria & 841710 & 1.00 \\ 
04004 & Técnico em Condomínio & 510110 & 1.00 \\ 
06005 & Técnico em Edificações & 312105, 318105, 514310 & 0.43 \\ 
06015 & Técnico em Transporte Aquaviário & 342605, 342610 & 1.00 \\ 
13002 & Técnico em Eventos & 354820, 354825 & 1.00 \\ 
05007 & Técnico em Programação de Jogos Digitais & 317120 & 1.00 \\ 
09014 & Técnico em Design de Móveis & 318425, 318805, 771105, 771110, 771115 & 0.83 \\ 
06014 & Técnico em Trânsito & 342305, 342310 & 1.00 \\ 
04015 & Técnico em Serviços Jurídicos & 351405, 351430 & 1.00 \\ 
06004 & Técnico em Desenho de Construção Civil & 318005, 318010, 318015, 318105, 318110, 318115, 318120, 318505, 318510 & 1.00 \\ 
09007 & Técnico em Conservação e Restauro & 768710 & 1.00 \\ 
03001 & Técnico em Alimentação Escolar & 841408 & 1.00 \\ 
03009 & Técnico em Treinamento e Instrução de Cães-Guia & 623005 & 1.00 \\ 
04007 & Técnico em Finanças & 353205, 353210, 353215, 353220 & 1.00 \\ 
07010 & Técnico em Guarda e Segurança & 517310 & 1.00 \\ 
03006 & Técnico em Infraestrutura Escolar & 514325 & 1.00 \\ 
06016 & Técnico em Transporte de Cargas & 342305, 342310 & 1.00 \\ 
05004 & Técnico em Informática para Internet & 317110 & 1.00 \\ 
06017 & Técnico em Transporte Metroferroviário & 342110, 783110 & 0.50 \\ 
03005 & Técnico em Desenvolvimento Comunitário & 515305, 515310, 515325 & 1.00 \\ 
08001 & Técnico em Agroindústria & 818110 & 1.00 \\ 
04012 & Técnico em Recursos Humanos & 351315, 411010 & 1.00 \\ 
\end{longtable}

\newpage
\section{ 01001 -- Técnico em Agente Comunitário de Saúde }

\textbf{Perfil do Curso (CNCT)}

O Técnico em Agente Comunitário de Saúde será habilitado para:

Orientar e acompanhar, sob a supervisão de profissional de nível
superior, indivíduos, suas famílias e a população em seu território,
levando em conta a interação com o processo saúde-doença no território.
Identificar e atuar nos múltiplos determinantes e condicionantes do
processo saúde e doença, para a promoção da saúde e redução de riscos à
saúde individual e da coletividade. Realizar o mapeamento e o
cadastramento de dados sociais, demográficos e de saúde, para contribuir
com a produção de informações e a construção de revisão contínua do
plano de ação em saúde para os territórios. Desenvolver suas atividades
norteadas pelas diretrizes, pelos princípios e estrutura organizacional
do SUS, bem como a partir dos referenciais éticos e políticos da
Educação Popular em Saúde. Promover a comunicação entre equipe
multidisciplinar (Equipe de Saúde da Família), unidade de saúde,
autoridades e comunidade. Promover a mobilização comunitária, ações
educativas e incentivar as atividades comunitárias, promovendo a
integração entre a equipe de saúde e a comunidade. Promover ações nas
áreas de vigilância em saúde e ambiental. Acompanhar e orientar, por
meio de visita domiciliar estabelecidas no planejamento das equipes, as
pessoas em situação de vulnerabilidade social e portadoras de doenças
crônicas e agravos que necessitam de maior número de visitas,
estimulando o autocuidado e a prevenção da exposição a fatores de
riscos, realizando procedimentos específicos nos casos indicados pela
equipe ou encaminhando quando necessário para a unidade de saúde de
referência.

Para a atuação como Técnico em Agente Comunitário de Saúde, são
fundamentais:

Conhecimentos das políticas públicas de saúde e compreensão de sua
atuação profissional frente às diretrizes, princípios e estrutura
organizacional do Sistema Único de Saúde (SUS). Conhecimentos referentes
ao âmbito da promoção da saúde, prevenção de agravos frequentes na
atenção primária, dirigidas a indivíduos, famílias, comunidades e
população. Conhecimentos referentes à educação popular em saúde, à
promoção da saúde dos indivíduos nos diferentes ciclos de vida, suas
famílias e sua comunidade e dos atributos derivados da atenção primária
da saúde. Conhecimentos relativos ao âmbito da promoção, da prevenção e
do monitoramento das situações de risco no âmbito da Vigilância em
Saúde. Conhecimentos referentes ao trabalho em equipe e
interdisciplinar, a comunicação em saúde, ao registro e informação em
saúde. Conhecimentos e saberes relacionados aos princípios das técnicas
aplicadas na área, sempre pautados numa postura humana e ética.
Atualização e aperfeiçoamento profissional por meio da Educação
Continuada. Conhecimentos e saberes relacionados à Política de Inclusão
e ao Atendimento Educacional Especializado.

\textbf{CBOs correspondidos e seus perfis (presentes no CNCT, e presentes na predição da IA)}

\textbf{CBO: 322255  ( Técnico em agente comunitário de saúde )}

Orienta e acompanha, sob a supervisão de profissional de nível superior,
indivíduos, suas famílias e a população em seu território, levando em
conta a interação com o processo saúde-doença. Efetua procedimentos de
admissão, apresentando-se à pessoa, situando-a no ambiente, controlando
sinais vitais, efetuando a verificação antropométrica (peso corporal,
altura, circunferência abdominal, entre outras medidas corporais), para
identificá-la e monitorar sua saúde. Auxilia equipe técnica em
procedimentos específicos, efetuando teste de glicemia, aferindo pressão
e temperatura axilar, acompanhando cartão de vacina do cidadão e
paciente por meio de carteira de saúde, a fim de identificar sinais de
doenças, agravos ou outros eventos de importância para saúde pública e
coletiva. Orienta e apoia a correta administração de medicação
prescrita, verificando medicamentos recebidos, acompanhando a
administração de medicação ao paciente no domicílio, durante o
tratamento. Realiza visitas domiciliares e a instituições. Planeja
roteiro com a equipe multidisciplinar, acompanhando-a nas visitas.
Realiza entrevistas e avalia condições do ambiente, visando identificar
necessidades do paciente. Planeja administração de vacina nos pacientes
acamados e domiciliados. Realiza visitas domiciliares regulares e
periódicas para acompanhar pacientes e identificar casos suspeitos de
doenças e agravos à saúde. Encaminha pacientes para unidade de saúde de
referência ou serviços de saúde, acompanhando-os aos centros médicos,
monitorando-os durante o período de tratamento. Identifica e cadastra
situações que interferem no curso das doenças ou que tenham importância
epidemiológica, comunicando, quando necessário, à autoridade sanitária
responsável. Orienta o bem-estar físico e mental do idoso, a prevenção
de quedas e acidentes domésticos. Realiza ações de prevenção de agravos
e curativas. Orienta a família de crianças e adolescentes sobre seu
crescimento e desenvolvimento físico e mental e verificando a situação
vacinal. Orienta gestante no pré-natal e a puérpera em relação ao
aleitamento materno e ao cuidado do recém-nascido nos primeiros seis
meses de vida, inclusive em visitas domiciliares regulares e periódicas.
Identifica condições de saúde existentes, detectando sinais de problemas
de ordem social, como negligência, violência doméstica, automutilação,
manifestações de doenças mentais e outros. Promove a saúde mental,
verificando pacientes com dependência química de álcool, de tabaco e de
outras drogas e pessoa em sofrimento psíquico, orientando a família a
procurar tratamento do profissional médico, a fim de prevenir tentativas
de suicídio e situações de risco. Promove a saúde da família,
participando da definição de território com a equipe, realizando
mapeamento com base em dados cadastrados (sociais, demográficos e de
saúde), cadastrando famílias, dados sociais e de saúde dos adscritos,
identificando grupos, famílias e indivíduos expostos a riscos ou em
vulnerabilidade social. Coordena o cuidado dos pacientes, organizando
grupos de promoção à saúde, realizando busca ativa de focos de risco à
saúde e de pacientes em descontinuidade de tratamento, comunicando
doenças, agravos e situações de importância local. Comunica ao médico
efeitos adversos de medicamentos e intercorrências. Registra dados nos
sistemas de informação do Ministério da Saúde e elabora relatório sobre
paciente. Utiliza tecnologia da informação e da comunicação para
registrar e consultar dados relacionados à sua função. Participa de
campanhas de saúde pública, reuniões comunitárias de grupos e comitês
representativos e das atividades de planejamento. Esclarece dúvidas no
que se refere à saúde e prevenção de doenças, viabilizando o diálogo
entre a equipe multidisciplinar, autoridades e comunidade. Cumpre
diretrizes do Sistema Único de Saúde, normas técnicas e regulamentadoras
de biossegurança, saúde e segurança no trabalho, de vigilância sanitária
e saúde ambiental.

\begin{Form}
\textbf{Esta correspondência está adequada?} \\ 
\ChoiceMenu[radio,radiosymbol=\ding{52},bordercolor=0 0.5 0,name=cbovalid_ 322255 ]{}{CBO deve ficar=sim, \hspace{6em} CBO não fica aqui=nao}
\end{Form}

\newpage
\section{ 01002 -- Técnico em Análises Clínicas }

\textbf{Perfil do Curso (CNCT)}

O Técnico em Análises Clínicas será habilitado para:

Executar, sob a supervisão do profissional responsável de nível
superior, processos operacionais necessários ao diagnóstico laboratorial
que compreendem a fase pré-analítica e analítica nos setores da
parasitologia, microbiologia, imunologia, hematologia, bioquímica,
biologia molecular, hormônios, toxicologia e líquidos corporais. Operar
aparato tecnológico de laboratório de saúde e equipamentos analíticos e
de suporte às atividades laboratoriais. Participar de campanhas
educativas e incentivar as atividades comunitárias de atenção primária,
promovendo a integração entre a equipe de saúde e a comunidade.
Recepcionar e cadastrar clientes e exames; realizar processos de coleta,
recepção, preparação e análise das amostras, colaborando ainda na
investigação e implantação de novas tecnologias biomédicas. Trabalhar de
acordo com as normas de biossegurança e qualidade, e aplicar as técnicas
adequadas no descarte de resíduos de serviços de saúde, protegendo os
indivíduos e o meio ambiente.

Para a atuação como Técnico em Análises Clínicas, são fundamentais:

Conhecimentos das políticas públicas de saúde e compreensão de sua
atuação profissional frente às diretrizes, princípios e estrutura
organizacional do Sistema Único de Saúde (SUS). Conhecimentos e saberes
relacionados aos princípios das técnicas aplicadas na área, sempre
pautados numa postura humana, ética e bioética. Capacidade de raciocínio
lógico, coordenação motora fina, capacidade de concentração e boa
acuidade (percepção) visual. Resolução de situações-problema,
comunicação, trabalho em equipe e interdisciplinar, tecnologias da
informação e da comunicação, gestão de conflitos e ética profissional.
Organização e responsabilidade. Iniciativa social. Determinação e
criatividade, buscando promover a humanização da assistência.
Atualização e aperfeiçoamento profissional por meio da Educação
Continuada.

\textbf{CBOs correspondidos e seus perfis (presentes no CNCT, e presentes na predição da IA)}

\textbf{CBO: 324205  ( Técnico em patologia clínica )}

Coleta material biológico, verificando o pedido de exame, confirmando o
preparo do paciente, fazendo ou orientando a assepsia da região de
coleta e identificando o material coletado. Coloca conservantes e
acondiciona a amostra. Pode receber material biológico já coletado,
confrontando-o com o pedido, conferindo suas condições e, se não
conforme, solicitando nova coleta. Prepara amostra do material
biológico, selecionando técnica, preparando soluções e reagentes,
diluindo material, centrifugando, sequenciando e homogeneizando as
amostras. Analisa material biológico, realizando testes
imuno-hematológicos, identificando e isolando microrganismos, realizando
análise microscópica e macroscópica e outros tipos, de acordo com
protocolos técnicos. Compara resultados com os parâmetros de
normalidade, com resultados anteriores e com os dados clínicos do
paciente. Encaminha os resultados dos exames para laudo clínico. Opera
equipamentos analíticos e de suporte, verificando seu funcionamento e
sua adequação e controlando sua temperatura. Solicita manutenção
preventiva e, quando necessária, manutenção corretiva dos equipamentos.
Aplica técnicas de descarte de resíduos biológicos e químicos, desinfeta
instrumental e equipamentos e esteriliza instrumentos. Realiza
acondicionamento de material para descarte responsável. Registra os
procedimentos do exame e transcreve resultados em sistemas
informatizados ou analógicos.

\begin{Form}
\textbf{Esta correspondência está adequada?} \\ 
\ChoiceMenu[radio,radiosymbol=\ding{52},bordercolor=0 0.5 0,name=cbovalid_ 324205 ]{}{CBO deve ficar=sim, \hspace{6em} CBO não fica aqui=nao}
\end{Form}

\newpage
\section{ 01003 -- Técnico em Citopatologia }

\textbf{Perfil do Curso (CNCT)}

O Técnico em Citopatologia será habilitado para:

Colaborar na investigação e implantação de novas tecnologias. Executar,
sob a supervisão do profissional responsável de nível superior,
atividades padronizadas de laboratório referentes aos exames
microscópicos em sua área técnica. Operar aparato tecnológico de
laboratório de saúde e equipamentos analíticos e de suporte. Participar
de campanhas educativas e incentivar as atividades comunitárias de
atenção primária, promovendo a integração entre a equipe de saúde e a
comunidade. Promover a comunicação com a sua equipe e com os
responsáveis técnicos. Realizar a análise microscópica para rastrear
células neoplásicas na amostra, estabelecendo relação das alterações
citológicas com o histórico clínico do paciente, elaborando um laudo
técnico que orientará o diagnóstico pelo responsável técnico. Receber e
preparar amostras para análise citopatológica. Trabalhar de acordo com
as normas de biossegurança e qualidade, e aplicar as técnicas adequadas
no descarte de resíduos de serviços de saúde, protegendo os indivíduos e
o meio ambiente.

Para a atuação como Técnico em Citopatologia, são fundamentais:

Conhecimento das políticas públicas de saúde e compreensão de sua
atuação profissional frente às diretrizes, princípios e estrutura
organizacional do Sistema Único de Saúde (SUS). Conhecimentos e saberes
relacionados aos princípios das técnicas aplicadas na área, sempre
pautados numa postura humana, ética e bioética. Capacidade de raciocínio
lógico, coordenação motora fina, capacidade de concentração e boa
acuidade (percepção) visual. Resolução de situações-problema,
comunicação, trabalho em equipe e interdisciplinar, tecnologias da
informação e da comunicação, gestão de conflitos e ética profissional.
Organização e responsabilidade; iniciativa social; determinação e
criatividade, buscando promover a humanização da assistência.
Atualização e aperfeiçoamento profissional por meio da Educação
Continuada.

\textbf{CBOs correspondidos e seus perfis (presentes no CNCT, e presentes na predição da IA)}

\textbf{CBO: 324215  ( Citotécnico )}

Recebe e tria material biológico, confrontando-o com o pedido médico,
conferindo suas condições e seu armazenamento. O material não-conforme é
rejeitado, solicitando-se nova amostra. Prepara amostra do material
biológico, selecionando técnicas, preparando soluções e reagentes,
sequenciando amostras, diluindo material biológico, homogeneizando
amostras, confeccionado e corando lâminas, centrifugando e inativando
material. Realiza análises citomorfológicas de líquidos, fluidos
orgânicos, secreções, esfregaços, raspados e outros materiais, por meio
da microscopia de amostras. Compara o resultado de exame com parâmetros
de normalidade, com resultados de exames anteriores e com dados clínicos
do paciente. Encaminha os resultados para o responsável pela emissão de
laudo técnico, para apoiar o diagnóstico precoce e a identificação de
doenças. Registra os procedimentos do exame e transcreve resultados
observados. Opera equipamentos analíticos e de suporte, checando seu
funcionamento, ajustando, programando e controlando temperatura.
Solicita manutenção preditiva e corretiva, quando necessária. Limpa e
desinfeta equipamentos e bancada, bem como esteriliza instrumentos. Pode
supervisionar equipe de auxiliares.

\begin{Form}
\textbf{Esta correspondência está adequada?} \\ 
\ChoiceMenu[radio,radiosymbol=\ding{52},bordercolor=0 0.5 0,name=cbovalid_ 324215 ]{}{CBO deve ficar=sim, \hspace{6em} CBO não fica aqui=nao}
\end{Form}

\newpage
\section{ 01004 -- Técnico em Controle Ambiental }

\textbf{Perfil do Curso (CNCT)}

O Técnico em Controle Ambiental será habilitado para:

Propor medidas para a minimização dos impactos e recuperação de
ambientes já degradados. Realizar ações de saúde ambiental nos
territórios. Promover monitoramento e ações sustentáveis de manejo
ambiental (hídrico, edáfico e atmosférico). Controlar processos
produtivos. Identificar o potencial poluidor de processos produtivos.
Monitorar e gerenciar os dados de controle das estações de tratamento de
água, esgoto, efluentes industriais, resíduos sólidos e emissões
atmosféricas. Executar coleta, medições in situ e análises
físico-químicas e microbiológicas das matrizes ambientais, operações e
processos unitários de tratamento. Avaliar as intervenções antrópicas e
utilizar tecnologias de prevenção, correção e monitoramento ambiental.
Realizar levantamentos ambientais. Realizar processos de educação
ambiental nos territórios e unidades de controle da poluição e reuso.
Identificar tecnologias apropriadas para o processo de produção
racional, redução de energia, reuso de águas residuárias, biomassa e
co-geração. Operar sistemas de tratamento de poluentes e de resíduos
sólidos. Executar análises de controle de qualidade ambiental. Realizar
vistorias ambiental e sanitária. Identificar e intervir nos problemas
relacionados aos fatores de riscos ambientais do território com o
propósito de contribuir para a melhoria da qualidade de vida e de
trabalho. Integrar ações de saúde do trabalhador com saúde ambiental.

Para a atuação como Técnico em Controle Ambiental, são fundamentais:

Conhecimentos das políticas públicas de saúde e compreensão de sua
atuação profissional frente às diretrizes, princípios e estrutura
organizacional do Sistema Único de Saúde (SUS). Conhecimentos das
políticas públicas de Meio Ambiente e compreensão de sua atuação
profissional frente às diretrizes, princípios e estrutura organizacional
do Sistema Nacional de Meio Ambiente (SISNAMA). Conhecimentos e saberes
relacionados a processos de sustentabilidade, territorialização e
monitoramento ambiental. Organização e responsabilidade. Resolução de
situações-problema, gestão de conflitos, trabalho em equipe de forma
colaborativa, comunicação e ética profissional. Atualização e
aperfeiçoamento profissional por meio da Educação Continuada.

\textbf{CBOs correspondidos e seus perfis (presentes no CNCT, e presentes na predição da IA)}

\textbf{CBO: 311505  ( Técnico de controle de meio ambiente )}

Prepara ambiente para análises, selecionando reagentes, vidrarias e
equipamentos. Identifica as amostras de efluentes coletadas, registrando
os pontos de coleta. Realiza testes e análises físico-químicas e
microbiológicas de efluentes. Interpreta resultados analíticos e elabora
laudos, relatórios e planilhas dos resultados. Pode encaminhar amostras
para análises externas complementares. Coordena processos de medição
para controle ambiental e tratamento de efluentes, monitorando os
parâmetros do processo e emitindo relatórios, para identificar os
aspectos ambientais e os impactos associados. Monitora a qualidade e o
uso das águas, bem como as características das emissões atmosféricas,
dos efluentes e dos resíduos sólidos, conforme manuais e procedimentos
técnicos. Pode coletar, orientar ou coordenar a coleta de amostras de
gases, solo, água e outros materiais contaminados, inclusive por
radiação. Implementa ações preventivas e corretivas, participando de
auditorias e propondo medidas para a mitigação dos impactos e
recuperação de ambientes já degradados. Executa tarefas de caráter
técnico referentes às questões ambientais da unidade fabril,
controlando, organizando e monitorando os resíduos e os efluentes e
verificando o licenciamento ambiental. Pode orientar o aproveitamento
dos resíduos industriais, buscando valorização econômica, ecoeficiência
e sustentabilidade do processo. Participa da implantação de programas de
qualidade ambiental, analisando indicadores e aplicando ferramentas da
qualidade. Desenvolve programas de educação ambiental junto à comunidade
e às instituições. Realiza campanhas de monitoramento e educação
ambiental, identificando tecnologias apropriadas para o processo de
produção racional e cuidados com o meio ambiente. Pode participar de
equipes de desenvolvimento de materiais biodegradáveis. Monitora o
funcionamento de máquinas, equipamentos e instrumentos de medição e
controle, mantendo-os em condições de uso e realizando ajustes no
processo, quando necessários. Monitora a segurança no trabalho,
interpretando mapas de risco; controlando o uso de equipamentos de
proteção individual e coletiva; e cumprindo procedimentos de emergência,
quando necessário.

\begin{Form}
\textbf{Esta correspondência está adequada?} \\ 
\ChoiceMenu[radio,radiosymbol=\ding{52},bordercolor=0 0.5 0,name=cbovalid_ 311505 ]{}{CBO deve ficar=sim, \hspace{6em} CBO não fica aqui=nao}
\end{Form}

\newpage
\section{ 01005 -- Técnico em Cuidados de Idosos }

\textbf{Perfil do Curso (CNCT)}

O Técnico em Cuidados de Idosos será habilitado para:

Acompanhar idosos a serviços previdenciários, de assistência social e de
saúde. Administrar medicação oral e tópica conforme prescrição médica
Auxiliar a pessoa idosa nas atividades básicas e instrumentais da vida
diária. Cuidar de pessoas idosas, dependentes ou independentes, nos
aspectos físico, mental, social e cultural. Estimular atividades
ocupacionais e de lazer. Identificar situações de urgência e
necessidades de primeiros socorros. Promover o envelhecimento ativo e a
saúde funcional. Promover os direitos e a dignidade da pessoa idosa.
Realizar atividades de promoção da saúde e cuidados. Recomendar
adequações ambientais que previnam acidentes. Zelar pela autonomia do
idoso e melhoria da qualidade de vida.

Para a atuação como Técnico em Cuidados de Idosos, são fundamentais:

Conhecimento das políticas públicas de saúde e compreensão de sua
atuação profissional frente às diretrizes, princípios e estrutura
organizacional do Sistema Único de Saúde (SUS). Respeito aos direitos da
pessoa idosa. Conhecimentos e saberes relacionados aos princípios das
técnicas aplicadas na área, sempre pautados numa postura humana e na
ética do cuidado. Conhecimento sobre biossegurança e sobre estratégias
para o autocuidado físico e mental. Resolução de situações-problema,
gestão de conflitos, trabalho em equipe de forma colaborativa,
comunicação e ética profissional. Atualização e aperfeiçoamento
profissional por meio da Educação Continuada.

\textbf{CBOs correspondidos e seus perfis (presentes no CNCT, e presentes na predição da IA)}

\textbf{CBO: 516210  ( Cuidador de idosos )}

Planeja o trabalho, a partir de informações -- levantadas junto ao
idoso, aos seus familiares e/ou aos representantes da instituição onde
ele vive -, sobre o estado de saúde e sobre a história e o estilo de
vida do idoso. Cuida da pessoa idosa em suas atividades da vida diária,
estimulando o seu grau de autonomia e independência. Ajuda-o, quando
necessário, na organização da sua rotina, na sua higiene corporal e na
sua imagem pessoal. Cuida da aparência da pessoa idosa -- unhas, cabelos
e vestuário - para aumentar a sua autoestima. Cuida da alimentação e da
hidratação do idoso, observando a qualidade e a validade dos alimentos,
verificando e dosando as propriedades nutricionais dos alimentos -
energéticos, construtores e reguladores - e estimulando a ingestão de
líquidos. Pode participar na elaboração do cardápio, preparar e servir a
refeição. Pode ajudar o idoso a preparar sua refeição. Cuida da saúde do
idoso, acompanhando-o aos exames, consultas e tratamentos de saúde e
transmitindo aos profissionais de saúde as mudanças no comportamento,
humor ou aparecimento de alterações físicas. Monitora o estado de saúde
do idoso de acordo com orientações da equipe multiprofissional e
premissas do cuidado humanizado. Pode cuidar da medicação oral da pessoa
idosa, em dose e horário prescritos pelo médico. Cuida do idoso em
recuperação de saúde, com limitação de mobilidade e com deficiência,
auxiliando-o nas suas atividades diárias. Promove atividades de lazer e
de ocupação do tempo livre, de acordo com o interesse do idoso. Estimula
a realização de trabalhos manuais, passeios e caminhadas. Incentiva a
prática de exercícios físicos, sempre que autorizados pelos
profissionais de saúde. Facilita a participação do idoso em atividades
culturais -- teatro, espetáculos de dança, cinema, entre outras -, que
sejam de seu interesse, acompanhando-o. Providencia livros, para
leitura. Pode orientar o idoso na digitação de mensagens e no envio de
imagens e vídeos, em aplicativo ou em rede social. Promove o bem-estar
da pessoa idosa, facilitando e estimulando sua comunicação com
familiares e amigos, quando possível. Mantém o ambiente, onde o idoso
vive, limpo e organizando, mantendo os objetos de uso arrumados nos
locais habituais. Comunica-se com os familiares ou responsáveis, para
prestar informações ou transmitir relatórios sobre o estado de saúde da
pessoa idosa. Zela pela segurança do idoso, prevenindo acidentes como
escorregões, quedas e lesões com objetos cortantes. Recomenda, à família
ou à instituição, adaptações no ambiente, para atender às condições de
saúde, mobilidade e segurança do idoso. Presta, ao idoso, cuidados de
primeiros socorros em caso de acidentes ou em outras situações de
emergência, solicitando auxílio aos serviços ou profissionais
especializados.

\begin{Form}
\textbf{Esta correspondência está adequada?} \\ 
\ChoiceMenu[radio,radiosymbol=\ding{52},bordercolor=0 0.5 0,name=cbovalid_ 516210 ]{}{CBO deve ficar=sim, \hspace{6em} CBO não fica aqui=nao}
\end{Form}

\newpage
\section{ 01007 -- Técnico em Enfermagem }

\textbf{Perfil do Curso (CNCT)}

O Técnico em Enfermagem será habilitado para:

Realizar, sob a supervisão do Enfermeiro, cuidados integrais de
enfermagem a indivíduos, família e grupos sociais vulneráveis ou não.
Atuar na promoção, prevenção, recuperação e reabilitação dos processos
saúde-doença em todo o ciclo vital. Participar do planejamento e
execução das ações de saúde junto à equipe multidisciplinar,
considerando as normas de biossegurança, envolvendo curativos,
administração de medicamentos e vacinas, nebulizações, banho de leito,
cuidados pós-morte, mensuração antropométrica e verificação de sinais
vitais. Preparar o paciente para os procedimentos de saúde. Participar
de comissões de certificação de serviços de saúde, tais como núcleo de
segurança do paciente, serviço de controle de infecção hospitalar,
gestão da qualidade dos serviços prestados à população, gestão de
riscos, comissões de ética de enfermagem, transplantes, óbitos e outros.
Colaborar com o Enfermeiro em ações de comissões de certificação de
serviços de saúde, tais como núcleo de segurança do paciente, serviço de
controle de infecção hospitalar, gestão da qualidade dos serviços
prestados à população, gestão de riscos, comissões de ética de
enfermagem, transplantes, óbitos e outros.

Para a atuação como Técnico em Enfermagem, são fundamentais:

Conhecimentos das políticas públicas de saúde e compreensão de sua
atuação profissional frente às diretrizes, princípios e estrutura
organizacional do Sistema Único de Saúde (SUS). Conhecimentos e saberes
relacionados aos princípios das técnicas aplicadas na área, sempre
pautados numa postura humana e ética. Resolução de situações-problema,
comunicação, trabalho em equipe e interdisciplinar, tecnologias da
informação e da comunicação, gestão de conflitos e ética profissional.
Organização e responsabilidade. Iniciativa social. Determinação e
criatividade, buscando promover a humanização da assistência.
Atualização e aperfeiçoamento profissional por meio da educação
continuada.

\textbf{CBOs correspondidos e seus perfis (presentes no CNCT, e presentes na predição da IA)}

\textbf{CBO: 322205  ( Técnico de enfermagem )}

Registra sinais vitais, como temperatura, pressão arterial, pulso ou
frequência respiratória do paciente. Auxilia a equipe técnica em
procedimentos invasivos, reanimação, preparação do paciente para exames
e cirurgia. Administra medicação prescrita, conferindo previamente a
medicação e os registros e a identificação do paciente. Pode administrar
medicamentos prescritos de alta vigilância e hemocomponentes. Mede e
registra a ingestão de alimentos e líquidos e a produção urinária e
fecal, relatando alterações para equipe médica ou de enfermagem. Detecta
sintomas que possam exigir atenção médica como hematomas, feridas
abertas, sangue na urina. Reposiciona pacientes acamados. Apoia o
paciente nas atividades da vida diária como sair da cama, tomar banho,
vestir-se, usar banheiro, ficar em pé, caminhar ou se exercitar. Presta
assistência de enfermagem à gestante - no período gestacional, no parto
e durante puerpério - e ao recém-nascido. Presta assistência de
enfermagem no período perioperatório. Executa ações de prevenção,
promoção, proteção, reabilitação e recuperação da saúde em serviços de
saúde extra-hospitalar. Presta assistência de enfermagem em cuidados
críticos e em cuidados paliativos, em domicílios e entidades de longa
permanência. Atua, ainda, em instituições que prestam atendimento
pré-hospitalar e serviços de diagnósticos, de resgate, remoção e
transporte de clientes. Pode participar de comissões de certificação de
serviços de saúde, tais como núcleo de segurança do paciente, serviço de
infecção hospitalar, gestão da qualidade, gestão de riscos, comissões de
ética de enfermagem, transplantes, óbitos e outros. Pode atuar em
programas de qualidade e certificação hospitalar. Pode atuar em
programas de educação.

\begin{Form}
\textbf{Esta correspondência está adequada?} \\ 
\ChoiceMenu[radio,radiosymbol=\ding{52},bordercolor=0 0.5 0,name=cbovalid_ 322205 ]{}{CBO deve ficar=sim, \hspace{6em} CBO não fica aqui=nao}
\end{Form}

\textbf{CBO: 322210  ( Técnico de enfermagem de terapia intensiva )}

Efetua os procedimentos de admissão do paciente de risco, em ambientes
pós-operatório, de isolamento ou de longa permanência. Monitora a
evolução dos pacientes, efetuando registros de sinais vitais e de
exames. Presta assistência humanizada e integral ao paciente, realizando
procedimentos de acesso venoso, de introdução de cateter nasogástrico e
vesical, de instalação de alimentação induzida, entre outros. Administra
medicação, verificando a conformidade da identificação do paciente e da
prescrição médica. Instala hemoderivados e quimioterápicos, de acordo
com a dosagem prescrita. Pode realizar instrumentação cirúrgica em
unidades de menor porte. Auxilia a equipe técnica em procedimentos
invasivos, em reanimação e em preparo do paciente para intervenções
cirúrgicas. Cuida do corpo pós morte. Participa de reuniões de troca de
plantões. Participa de programas de qualidade da assistência em UTI e
programas de controle de infecção hospitalar.

\begin{Form}
\textbf{Esta correspondência está adequada?} \\ 
\ChoiceMenu[radio,radiosymbol=\ding{52},bordercolor=0 0.5 0,name=cbovalid_ 322210 ]{}{CBO deve ficar=sim, \hspace{6em} CBO não fica aqui=nao}
\end{Form}

\textbf{CBO: 322215  ( Técnico de enfermagem do trabalho )}

Participa da organização e funcionamento do serviço de segurança e saúde
do trabalhador. Atua em situações de urgência e emergência, integrando a
equipe multidisciplinar de saúde e do serviço especializado em
engenharia e medicina do trabalho (SESMT). Auxilia na realização de
exame pré-admissional, periódicos e demissionais. Atende trabalhadores
portadores de doenças ou lesões de pouca gravidade, sob supervisão.
Realiza ações de prevenção de doenças ocupacionais e promoção da saúde e
segurança no trabalho. Participa de programas de prevenção de acidentes
e de reabilitação profissional. Auxilia na observação sistemática do
estado de saúde dos trabalhadores, da ocorrência de doenças
profissionais, lesões traumáticas e doenças epidemiológicas. Preenche
relatório com os dados observados. Pode fazer visitas domiciliares e
hospitalares, para acompanhar condições do trabalhador, nos casos de
acidentes ou doenças profissionais. Participa de campanhas de educação
sanitária e de vacinação. Pode atuar em programas de qualidade e de
certificação de serviços de saúde do trabalhador.

\begin{Form}
\textbf{Esta correspondência está adequada?} \\ 
\ChoiceMenu[radio,radiosymbol=\ding{52},bordercolor=0 0.5 0,name=cbovalid_ 322215 ]{}{CBO deve ficar=sim, \hspace{6em} CBO não fica aqui=nao}
\end{Form}

\textbf{CBO: 322220  ( Técnico de enfermagem psiquiátrica )}

Cuida de indivíduos ou grupos com problemas emocionais ou mentais,
monitorando o bem-estar físico e emocional dos pacientes. Relata
comportamentos incomuns ou doenças físicas à equipe de saúde,
registrando as ocorrências. Ajuda pacientes a realizar tarefas, como
tomar banho ou manter limpas as camas, roupas ou áreas de estar. Age na
prevenção de tentativas de suicídio e situações de risco, limitando o
espaço de circulação do paciente e demarcando limites de comportamento.
Administra medicamentos orais ou injeções hipodérmicas, seguindo as
prescrições médicas e procedimentos hospitalares. Protege pacientes
durante crises, por meios verbais ou físicos, quando necessários.
Estimula o paciente na expressão de sentimentos e o conduz para
atividades sociais e terapêuticas. Observa e influencia o comportamento
dos pacientes, dialogando, interagindo, ensinando e aconselhando.
Incentiva o paciente a desenvolver habilidades de trabalho e a
participar de atividades sociais, recreativas ou outras atividades
terapêuticas que aprimorem as habilidades interpessoais e estimulem as
relações sociais. Colabora na equipe de enfermeiros, médicos, psicólogos
ou terapeutas de reabilitação no cuidado e tratamento dos pacientes,
visando desinstitucionalizar os pacientes internados para o retorno ao
convívio comunitário. Requisita medicamentos do dispensário e mantém
registros de acordo com os procedimentos especificados. Pode participar
da preparação de novos funcionários sobre procedimentos a serem seguidos
com pacientes psiquiátricos. Pode entrar em contato com os parentes dos
pacientes para organizar reuniões familiares. Pode supervisionar
auxiliares de enfermagem em psiquiatria.

\begin{Form}
\textbf{Esta correspondência está adequada?} \\ 
\ChoiceMenu[radio,radiosymbol=\ding{52},bordercolor=0 0.5 0,name=cbovalid_ 322220 ]{}{CBO deve ficar=sim, \hspace{6em} CBO não fica aqui=nao}
\end{Form}

\textbf{CBO: 322225  ( Instrumentador cirúrgico )}

Dispõe os instrumentos cirúrgicos, colocando pinças, afastadores e
demais instrumentos na ordem de utilização durante o ato operatório.
Verifica quantidade de material cirúrgico e compressas. Testa pinças
anatômicas, hemostáticas e outros instrumentos cirúrgicos eletrônicos ou
robóticos, verificando seu funcionamento e articulações, para evitar
falhas na utilização. Paramenta-se e ajuda a equipe cirúrgica a
paramentar-se com roupa cirúrgica, luvas e máscaras. Mantém campo
estéril adequado durante procedimentos cirúrgicos. Prepara os pacientes
para a cirurgia, incluindo o posicionamento dos pacientes na mesa de
operações e a cobertura com campos cirúrgicos estéreis para evitar a
exposição. Passa instrumentos e suprimentos a cirurgiões e assistentes
de cirurgiões, segura afastadores e corta suturas, bem como realiza
outras tarefas, cumprindo as instruções do cirurgião. Monitora e avalia
continuamente as condições da sala de operações, incluindo as
necessidades do paciente e da equipe cirúrgica. Opera, monta, ajusta ou
monitora esterilizadores, luzes, máquinas de sucção ou equipamento de
diagnóstico para garantir condições adequadas à cirurgia. Auxilia a
remoção do paciente para sala de recuperação pós-operatória ou quarto.
Prepara e encaminha amostras de tecido coletadas para análise
laboratorial. Lava e esteriliza o equipamento usando germicidas e
esterilizadores. Conta esponjas, compressas, agulhas e instrumentos
antes e depois da operação. Faz a assepsia da mesa de instrumental,
limpando e esterilizando o local e as peças. Auxilia na arrumação da
sala cirúrgica, retirando e separando o instrumental e material usado,
para deixar a sala organizada para a cirurgia posterior.

\begin{Form}
\textbf{Esta correspondência está adequada?} \\ 
\ChoiceMenu[radio,radiosymbol=\ding{52},bordercolor=0 0.5 0,name=cbovalid_ 322225 ]{}{CBO deve ficar=sim, \hspace{6em} CBO não fica aqui=nao}
\end{Form}

\textbf{CBO: 322230  ( Auxiliar de enfermagem )}

Efetua procedimentos de admissão do paciente, procedendo sua
identificação e o acomodando no local de atendimento. Presta atendimento
básico ao paciente, ajudando-o nas refeições, a higienizar-se, a
banhar-se a vestir-se e a caminhar. Arruma e troca roupas do leito e
muda o decúbito do paciente no leito para aumentar o conforto e evitar
escaras. Colhe sinais vitais do paciente e procede os devidos registros.
Administra medicamentos prescritos, identificando e conferindo a
medicação a ser administrada - leito, nome e registro do paciente. Pode
ajudar a equipe de enfermagem em atividades complexas - como instalação
de sonda nasogástrica e outros procedimentos invasivos - e em reanimação
do paciente. Pode auxiliar nos procedimentos pré-cirúrgicos como
higienização do paciente e tricotomia. Segue rigorosamente os protocolos
de biossegurança, higienizando-se e protegendo-se a cada atividade junto
ao paciente. Pode auxiliar na remoção de pacientes e nos cuidados do
corpo pós-morte.

\begin{Form}
\textbf{Esta correspondência está adequada?} \\ 
\ChoiceMenu[radio,radiosymbol=\ding{52},bordercolor=0 0.5 0,name=cbovalid_ 322230 ]{}{CBO deve ficar=sim, \hspace{6em} CBO não fica aqui=nao}
\end{Form}

\textbf{CBO: 322240  ( Auxiliar de saúde (navegação marítima) )}

Efetua procedimentos de admissão do paciente no ambiente da enfermaria
da embarcação, arrola seus pertences, providencia sua higiene e
fornece-lhe roupa. Controla sinais vitais do paciente e mensura seu peso
e sua altura. Presta assistência ao paciente, ajudando-o a alimentar-se
- ou instalando alimentação induzida -, trocando curativos, mudando
decúbito no leito, aplicando clister para lavagem intestinal e
realizando outros procedimentos, sob supervisão de profissional de saúde
de nível superior. Administra medicação prescrita, identificando e
conferindo a medicação a ser administrada - leito, nome e registro do
paciente. Auxilia a equipe de saúde em atividades específicas, tais como
tricotomia, reanimação de paciente, procedimentos invasivos e realização
de testes e exames. No caso de saúde mental, implementa atividades
terapêuticas prescritas e protege pacientes durante crises. Controla a
administração de vacinas segundo exigências internacionais. Remove o
corpo do paciente após morte. Organiza o ambiente de trabalho,
providenciando material de consumo, controlando e fiscalizando validade
de medicamentos de uso do paciente e do posto de enfermagem. Identifica
pertences de pacientes, coloca etiquetas em prescrições médicas. Arruma
camas e rouparia. Orienta paciente e seus familiares e presta
informações ao enfermeiro e ao médico. Registra intercorrências e
procedimentos realizados. Faz registros de saúde da tripulação em meio
digital. Encaminha a dispensa de trabalho de tripulante doente e/ou
acidentado. Atua de acordo com protocolos internacionais de
biossegurança e normas de saúde e segurança no trabalho, lavando as mãos
antes e após cada procedimento, usando equipamento de proteção
individual (EPI), paramentado-se, precavendo-se contra efeitos adversos
dos produtos e seguindo protocolo em caso de contaminação ou acidente.
Providencia limpeza concorrente e terminal. Desinfeta aparelhos e
materiais. Transporta roupas para expurgo e descarta material
contaminado.

\begin{Form}
\textbf{Esta correspondência está adequada?} \\ 
\ChoiceMenu[radio,radiosymbol=\ding{52},bordercolor=0 0.5 0,name=cbovalid_ 322240 ]{}{CBO deve ficar=sim, \hspace{6em} CBO não fica aqui=nao}
\end{Form}

\textbf{CBO: 322245  ( Técnico de enfermagem da estratégia de saúde da família )}

Participa do processo de territorialização e mapeamento da área de
atuação da sua equipe, identificando grupos, famílias e indivíduos
expostos a riscos, inclusive aqueles relativos ao trabalho. Participa da
atualização contínua dessas informações, para subsidiar a definição de
prioridades a serem incluídas no planejamento local. Participa das
atividades de assistência básica em sua área de atuação, realizando
procedimentos de enfermagem na unidade de saúde familiar e, quando
indicado ou necessário, no domicílio e/ou nos demais espaços
comunitários (escolas, associações, entre outros). Cuida da saúde da
população local - realizando ações curativas e de prevenção de agravos -
e notifica doenças e agravos de notificação compulsória, sob orientação
de profissional de nível superior da área de saúde. Realiza a escuta
qualificada das necessidades dos usuários em todas as ações,
proporcionando atendimento humanizado e viabilizando o estabelecimento
do vínculo. Administra medicação prescrita, calculando dosagem,
executando antissepsia e acompanhando paciente na ingestão. Controla a
administração de vacinas. Auxilia a equipe médica, ajudando em
reanimação do paciente e coletando material para exame. Participa da
assistência convalescente e reabilitadora de pessoas doentes ou
acidentadas. Realiza ações de educação em saúde a grupos específicos e a
famílias em situação de risco, conforme planejamento da equipe.
Participa em campanhas de saúde, incentivando a participação da
comunidade e identificando parceiros e recursos disponíveis. Participa
de programas para proteger a saúde das crianças, incluindo palestras
sobre saúde infantil, saúde escolar, instruções em grupo para os pais e
programas de imunização. Participa de discussões de casos e de troca
informações técnicas. Registra intercorrências, procedimentos realizados
e administração de medicação. Providencia o registro das atividades
realizadas nos sistemas nacionais de informação na Atenção Básica.
Participa do gerenciamento dos insumos necessários para o funcionamento
da unidade de saúde familiar.

\begin{Form}
\textbf{Esta correspondência está adequada?} \\ 
\ChoiceMenu[radio,radiosymbol=\ding{52},bordercolor=0 0.5 0,name=cbovalid_ 322245 ]{}{CBO deve ficar=sim, \hspace{6em} CBO não fica aqui=nao}
\end{Form}

\textbf{CBO: 322250  ( Auxiliar de enfermagem da estratégia de saúde da família )}

Participa do processo de territorialização e mapeamento da área de
atuação da sua equipe, auxiliando na identificação de grupos, famílias e
indivíduos expostos a riscos, inclusive aqueles relativos ao trabalho.
Participa da atualização contínua dessas informações. Participa das
atividades de assistência básica em sua área de atuação, realizando
procedimentos de auxiliar de enfermagem na unidade de saúde familiar e,
quando indicado ou necessário, no domicílio e/ou nos demais espaços
comunitários (escolas, associações, entre outros). Auxilia na prestação
de assistência ao paciente, providenciando sua higienização, mensurando
seu peso e sua altura, trocando seus curativos, monitorando sua
evolução, entre outros procedimentos. Acolhe as pessoas da comunidade
local, escutando suas queixas e as encaminhado para avaliação de
enfermeiro ou para consulta médica. Auxilia na instrução da família
sobre cuidados maternos e infantis, sobre prevenção de doenças, sobre
cuidado e reabilitação de paciente e sobre outros assuntos relacionados
à saúde e ao bem-estar individual e comunitário. Participa de campanhas
de educação sanitária e de vacinação. Auxilia na busca ativa e na
notificação de doenças e agravos. Auxilia na organização do ambiente de
trabalho, providenciando material de consumo, controlando e fiscalizando
validade de medicamentos.

\begin{Form}
\textbf{Esta correspondência está adequada?} \\ 
\ChoiceMenu[radio,radiosymbol=\ding{52},bordercolor=0 0.5 0,name=cbovalid_ 322250 ]{}{CBO deve ficar=sim, \hspace{6em} CBO não fica aqui=nao}
\end{Form}

\newpage
\section{ 01008 -- Técnico em Equipamentos Biomédicos }

\textbf{Perfil do Curso (CNCT)}

O Técnico em Equipamentos Biomédicos será habilitado para:

Administrar e comercializar equipamentos biomédicos. Analisar e executar
os testes de calibração e aferição dos equipamentos médico-hospitalares.
Analisar, tecnicamente, os certificados de calibração e aferição.
Auxiliar na definição de medidas de controle de segurança e qualidade no
ambiente hospitalar referente ao uso de equipamentos
médico-hospitalares- odontológicos. Auxiliar na elaboração da
especificação técnica para aquisição de novos equipamentos. Auxiliar na
elaboração do planejamento da gestão da manutenção de equipamentos
médico-hospitalares-odontológicos. Coordenar o armazenamento e uso
adequado de equipamentos. Executar as ações planejadas de manutenção
preventiva, preditiva e corretiva, instalação, montagem, medições e
testes de equipamentos médicos. Executar instalação e manutenção de
equipamentos médico-hospitalares- odontológicos. Participar e realizar
treinamento operacional e de controle de segurança e qualidade de
equipamentos médicos para equipe médico-assistencial. Planejar e
executar instalação, montagem, medições e testes de equipamentos
biomédicos. Prestar assistência técnica no estudo e desenvolvimento de
projetos e pesquisas de tecnologias de equipamentos
médico-hospitalares-odontológicos. Realizar coleta de campo de dados que
auxiliam o planejamento e a gestão da manutenção de equipamentos médicos
e da infraestrutura das instalações hospitalares. Realizar e registrar
os procedimentos de manutenção preventiva, preditiva e corretiva de
equipamentos e instrumentos médico-hospitalares-odontológicos. Registrar
os procedimentos das ações planejadas de manutenção preventiva,
preditiva e corretiva de equipamentos e instrumentos
médico-hospitalares-odontológicos.

Para a atuação como Técnico em Equipamentos Biomédicos, são
fundamentais:

Conhecimentos das políticas públicas de saúde e compreensão de sua
atuação profissional frente às diretrizes, princípios e estrutura
organizacional do Sistema Único de Saúde (SUS). Conhecimentos e saberes
relacionados aos princípios das técnicas aplicadas na área. Resolução de
situações-problema, comunicação, trabalho em equipe e interdisciplinar,
tecnologias da informação e da comunicação, gestão de conflitos e ética
profissional. Organização e responsabilidade. Atualização e
aperfeiçoamento profissional por meio da educação continuada.

\textbf{CBOs correspondidos e seus perfis (presentes no CNCT, e presentes na predição da IA)}

\textbf{CBO: 915305  ( Técnico em manutenção de equipamentos e instrumentos médico-hospitalares )}

Planeja as atividades, analisando documentos técnicos de produtos e de
operações; interpretando procedimentos e normas técnicas; e selecionando
ferramentas, instrumentos, equipamentos, além de peças, acessórios e
materiais, para execução do trabalho. Realiza testes e ensaios em
equipamentos e instrumentos médico-hospitalares e odontológicos,
utilizando ferramentas, equipamentos e técnicas específicas, com o
objetivo de analisar parâmetros, identificar e diagnosticar defeitos.
Realiza ensaios químicos. Afere equipamentos e instrumentos. Testa
dispositivos de segurança de equipamentos. Testa conexões físicas e
protocolos de comunicação entre instrumentos, equipamentos e sistemas.
Avalia consumo de energia de equipamentos e instrumentos. Verifica
imagens produzidas por equipamentos, comparando com padrões. Testa
mecanismos de movimentação de equipamentos -- inclusive de robôs -- e
instrumentos. Realiza manutenção de equipamentos e instrumentos
médico-hospitalares e odontológicos, aplicando conhecimentos e técnicas
específicas, a fim de restabelecer plenas condições operacionais de
equipamentos e instrumentos. Manuseia instrumentos de medição; realiza
medições de potencial hidrogeniônico (pH) e medições de grandezas
elétricas e térmicas, dentre outras grandezas físicas. Verifica
parâmetros de peças e parâmetros de funcionamento; desmonta equipamentos
e instrumentos; identifica e elimina falhas. Manufatura peças para
reposição. Repara e substitui peças e conjuntos. Monta equipamentos e
instrumentos e ajusta conjuntos eletroeletrônicos e mecânicos. Realiza
testes de calibração e aferição, a fim de garantir plenas condições
operacionais de instrumentos e equipamentos médico-hospitalares e
odontológicos. Verifica certificados de calibração e aferição. Instala
equipamentos médico-hospitalares e equipamentos odontológicos, aplicando
conhecimentos e técnicas exigidas, a fim de disponibilizar os
equipamentos para funcionamento. Analisa arranjo físico do ambiente para
instalação de equipamentos; verifica disponibilidade de utilidades --
tais como alimentação elétrica, de água e de gás -; avalia as condições
ambientais -- como temperatura, umidade, luminosidade e fluxo de ar -,
para instalação de instrumentos, equipamentos e sistemas. Transporta
equipamentos e instrumentos, observando os cuidados necessários. Prepara
equipamentos para uso, definindo parâmetros de configuração. Coloca os
equipamentos em funcionamento (``startup''), completando o processo de
instalação. Executa procedimentos de manutenção preventiva, de acordo
com requisitos e especificações de fabricantes, substituindo peças,
quando necessário. Presta assistência técnica de equipamentos e
instrumentos médico-hospitalares e odontológicos, atendendo clientes --
de acordo com programação e cronograma de atendimento --, identificando
necessidades e provendo suporte técnico. Demonstra funcionamento de
equipamentos e instrumentos para clientes. Intermedeia sugestões de
clientes para a empresa e propõe melhorias para produtos. Atualiza-se
tecnologicamente, participando de treinamentos, acompanhando publicações
técnicas e consultando fontes diversificadas -- inclusive via internet.
Pode treinar equipes técnicas e usuários e realizar palestras técnicas.
Elabora documentos técnicos, tais como relatórios e procedimentos de
processos de manutenção e de funcionamento de instrumentos e
equipamentos. Faz orçamentos. Preenche ordens de serviço e emite
pareceres técnicos. Mantém local de trabalho limpo e organizado. Limpa e
descontamina equipamentos e instrumentos de trabalho. Separa sobras de
material para reciclagem e descarte, seguindo normas ambientais. Realiza
o trabalho de acordo com métodos de segurança e de higiene,
interpretando classificação de risco de áreas; e preservando saúde e
integridade física por meio de Equipamentos de Proteção Individual
(EPI). Cumpre protocolos e procedimentos de segurança e de higiene. Atua
de acordo com normas e medidas de biossegurança.

\begin{Form}
\textbf{Esta correspondência está adequada?} \\ 
\ChoiceMenu[radio,radiosymbol=\ding{52},bordercolor=0 0.5 0,name=cbovalid_ 915305 ]{}{CBO deve ficar=sim, \hspace{6em} CBO não fica aqui=nao}
\end{Form}

\newpage
\section{ 01009 -- Técnico em Estética }

\textbf{Perfil do Curso (CNCT)}

O Técnico em Estética será habilitado para:

Executar procedimentos estéticos faciais, corporais e capilares,
utilizando produtos cosméticos, técnicas e equipamentos com registro na
Agência Nacional de Vigilância Sanitária (Anvisa). Avaliar as condições
da pele por meio da anamnese, selecionar e executar procedimentos
estéticos faciais, corporais e capilares, além de orientar sobre os
cuidados específicos pós procedimento. Utilizar técnicas manuais,
associadas ou não a equipamentos, tecnologias e produtos cosméticos.
Tratar da promoção, proteção, manutenção e recuperação estética da pele.
Avaliar e selecionar as técnicas e os cosméticos mais apropriados de
acordo com as características e necessidades do cliente/paciente.).
Adotar os procedimentos de limpeza, desinfeção e esterilização dos
equipamentos, instrumentos e de todos os utensílios utilizados na
estética. Observar as prescrições médicas ou de outros profissionais da
saúde apresentadas pelo cliente, ou solicitar, após exame da situação,
avaliação médica e/ou de outro profissional da saúde, conforme
necessidade. Executar suas atividades em conformidade com as normas
vigentes da vigilância sanitária. Interagir com fornecedores e outros
profissionais de saúde, sobretudo o Esteticista, compondo equipes
multidisciplinares.

Para a atuação como Técnico em Estética, são fundamentais:

Conhecimentos das políticas públicas de saúde e compreensão de sua
atuação profissional frente às diretrizes, princípios e estrutura
organizacional do Sistema Único de Saúde (SUS). Conhecimentos e saberes
relacionados aos princípios das técnicas aplicadas na área e de
biossegurança, sempre pautados numa postura humana e ética. Resolução de
situações-problema, comunicação, trabalho em equipe e interdisciplinar,
gestão de conflitos e ética profissional. Organização, responsabilidade,
iniciativa social, determinação e criatividade, buscando promover a
humanização da assistência. Atualização e aperfeiçoamento profissional
por meio da educação continuada.

\textbf{CBOs correspondidos e seus perfis (presentes no CNCT, e presentes na predição da IA)}

\textbf{CBO: 322130  ( Esteticista )}

Prepara o ambiente de trabalho de acordo com os procedimentos a serem
realizados. Organiza a agenda com base em informações sobre os clientes.
Pesquisa diferentes cosméticos, identificando matérias-primas e
princípios ativos com base nas características químicas e funcionais de
cada componente, para selecionar os produtos a serem utilizados.
Recepciona o cliente, fornecendo informações sobre os serviços
oferecidos e esclarecendo as dúvidas sobre suas atividades. Efetua ou
atualiza cadastro, coletando informações gerais e de preferências do
cliente. Avalia as condições de saúde e os hábitos de vida do cliente
relacionados à área de estética. Indaga sobre a existência de
enfermidades correlacionadas ao sistema circulatório sanguíneo;
questiona sobre patologias decorrentes dos sistemas endócrino e nervoso;
e levanta outras informações sobre patologias que interferem na ação de
estética. Observa as prescrições médicas ou de outros profissionais da
saúde apresentadas pelo cliente. Realiza o levantamento dos hábitos de
vida do cliente que podem interferir na condição da saúde da pele e do
couro cabeludo. Realiza análise das condições e da aparência da face do
cliente, inspecionando a pele e identificando hidratação, biotipo e
fototipo cutâneo da face. Efetua análise das condições e da aparência do
corpo e coleta medidas antropométricas. Utiliza instrumento para
registrar a avaliação da saúde e as condições e a aparência da face e do
corpo do cliente. Faz encaminhamento ao médico ou a outro profissional
da área de saúde quando observar sintomas - de enfermidades, infecções,
irritações cutâneas, entre outros -, que inspirem cuidados
especializados. Define o plano de tratamento estético de baixa
complexidade para o cliente, especificando procedimentos estéticos
faciais, corporais e capilares a serem realizados com o uso de técnicas
manuais, associadas ou não a equipamentos, tecnologias e produtos
cosméticos com registro na Agência Nacional de Vigilância Sanitária
(Anvisa). Executa procedimentos estéticos faciais, tais como os
realizados para renovação celular, redução da secreção sebácea e
prevenção do processo de formação de manchas na pele da face. Orienta e
fornece informações referentes à terapia capilar. Realiza procedimentos
de relaxamento e bem-estar corporal em estética, executando manobras de
massagem relaxante com utilização de cosméticos; fazendo procedimentos
de hidratação; e efetuando manobras de drenagem linfática, conforme
necessidade da pele do corpo. Seleciona e aplica outros procedimentos
corporais -- tais como para estimulação celular e redução de medidas
corpóreas - de acordo com as necessidades do cliente, com base em sua
avaliação. Orienta o cliente sobre os cuidados específicos pós
procedimentos. Acompanha e registra, em sistema informatizado, a
evolução dos resultados alcançados. Define preço dos serviços de
estética e providencia o recebimento do pagamento. Pode vender
cosméticos aos clientes. Pode atuar com profissionais de nível superior
da área de saúde, auxiliando cirurgiões plásticos e dermatologistas no
fornecimento aos pacientes de cuidados com a pele pré e pós procedimento
cirúrgico; prestando apoio ao dermatologista na aplicação de peelings
químicos; entre outras atividades. Controla o estoque de produtos e
materiais, providenciando a aquisição para reposição. Mantém a
organização, a limpeza, a assepsia e a higienização dos ambientes. Adota
os procedimentos de limpeza, desinfeção e esterilização dos
equipamentos, instrumentos e utensílios utilizados na estética.
Providencia a manutenção de equipamentos, quando necessário. Descarta
material e produtos com validade vencida, seguindo os princípios
ambientais e as normas de biossegurança. Trabalha com segurança, usando
equipamentos de proteção individual e adotando postura ergonômica
correta durante o trabalho.

\begin{Form}
\textbf{Esta correspondência está adequada?} \\ 
\ChoiceMenu[radio,radiosymbol=\ding{52},bordercolor=0 0.5 0,name=cbovalid_ 322130 ]{}{CBO deve ficar=sim, \hspace{6em} CBO não fica aqui=nao}
\end{Form}

\newpage
\section{ 01010 -- Técnico em Farmácia }

\textbf{Perfil do Curso (CNCT)}

O Técnico em Farmácia será habilitado para:

Atender, sob a supervisão do farmacêutico, as prescrições de
medicamentos e cosméticos, interpretando a prescrição, separando e
fornecendo o produto solicitado e encaminhando ao farmacêutico casos
específicos. Auxiliar em processos administrativos relacionados ao
âmbito farmacêutico. Auxiliar na produção e controle de logística de
produtos em indústrias farmacêuticas e afins. Executar, como auxiliar,
as rotinas de compra, armazenamento e entrega de produtos farmacêuticos
e correlatos. Identificar e classificar produtos e formas farmacêuticas.
Participar da rotina de testes em laboratórios de pesquisa vinculados a
universidades, faculdades, institutos de pesquisa e indústrias
farmacêuticas. Realizar o controle e manutenção do estoque de produtos e
matérias-primas farmacêuticas. Realizar operações farmacotécnicas,
manipulação de formas farmacêuticas (alopáticas, fitoterápicas,
homeopáticas, cosméticas e afins). Realizar testes de controle de
qualidade.

Para a atuação como Técnico em Farmácia, são fundamentais:

Conhecimento das políticas públicas de saúde: organização, princípios e
diretrizes do Sistema Único de Saúde (SUS). Conhecimentos e saberes
relacionados aos princípios das técnicas aplicadas na área, sempre
pautados numa postura humana e ética. Conhecimentos e saberes
relacionados a processos de produção, sustentabilidade e logística.
Resolução de situações-problema, comunicação, trabalho em equipe e
interdisciplinar, tecnologias da informação e da comunicação, gestão de
conflitos e ética profissional. Organização e responsabilidade.
Iniciativa social. Determinação e criatividade, buscando promover a
humanização da assistência. Atualização e aperfeiçoamento profissional
por meio da Educação Continuada.

\textbf{CBOs correspondidos e seus perfis (presentes no CNCT, e presentes na predição da IA)}

\textbf{CBO: 325105  ( Auxiliar técnico em laboratório de farmácia )}

Prepara as atividades, verificando prioridades, organizando as
matérias-primas, separando as embalagens por tipo e examinando o
funcionamento dos equipamentos. Manipula a formulação, sob supervisão do
farmacêutico, realizando as operações de homogeneização, peneiramento,
filtração, trituração, emulsificação e solubilização dos componentes da
ficha de manipulação. Faz o encapsulamento de formulações sólidas e o
envase de semissólidos e líquidos. Realiza testes de qualidade, testando
a solubilidade das matérias-primas e determinando pH, a densidade e o
ponto de fusão dos materiais. Durante a execução dos testes, controla a
temperatura e a umidade do ambiente. Realiza a manutenção de rotina do
laboratório, higienizando equipamentos e utensílios e encaminhando para
descarte o material contaminado. Solicita manutenção preventiva e
corretiva de equipamentos. Controla estoques, utilizando softwares para
conferir produtos, verificar quantidade disponível de embalagens e
conferir quantidade, validade, lote e tempo de armazenamento de
matérias-primas. Solicita reposição de itens do estoque, controlando o
recebimento. Faz o descarte de materiais e produtos vencidos, seguindo
procedimento padrão. Documenta as atividades e os procedimentos da
manipulação farmacêutica, a entrada e a saída de estoque, as manutenções
de rotina e a dispensação de medicamentos. Trabalha de acordo com as
boas práticas de manipulação e dispensação, utilizando equipamentos de
proteção individual (EPI), aplicando técnicas de segurança e higiene
pessoal e separando material para descarte, conforme estabelecido em
procedimentos operacionais padronizados.

\begin{Form}
\textbf{Esta correspondência está adequada?} \\ 
\ChoiceMenu[radio,radiosymbol=\ding{52},bordercolor=0 0.5 0,name=cbovalid_ 325105 ]{}{CBO deve ficar=sim, \hspace{6em} CBO não fica aqui=nao}
\end{Form}

\textbf{CBO: 325110  ( Técnico em laboratório de farmácia )}

Interpreta a ficha de manipulação de produtos farmacêuticos e calcula,
separa e pesa os componentes da fórmula. Manipula a formulação,
realizando as operações de homogeneização, peneiramento, filtração,
trituração, emulsificação e solubilização dos componentes da ficha de
manipulação. Faz o encapsulamento de formulações sólidas e o envase de
semissólidos e líquidos. Confere fórmulas, examinando peso e volume,
revisando embalagens utilizadas e vistoriando rotulagem das fórmulas.
Realiza testes de qualidade, examinando características organolépticas,
testando a solubilidade das matérias-primas e determinando pH, a
densidade e o ponto de fusão dos materiais. Durante a execução dos
testes, controla a temperatura e a umidade do ambiente. Realiza a
manutenção de rotina do laboratório, higienizando equipamentos e
utensílios e encaminhando para descarte o material contaminado. Solicita
manutenção preventiva e corretiva de equipamentos. Pode recomendar a
substituição de equipamentos por novos, dotados de recursos de automação
na realização de operações farmacotécnicas. Controla estoques,
utilizando software para conferir o armazenamento de produtos,
matérias-primas e embalagens. Fraciona matérias-primas. Solicita
reposição de itens do estoque, emitindo ordens de compras. Faz a
dispensação de materiais e produtos vencidos, seguindo procedimento
padrão. Documenta as atividades e os procedimentos da manipulação
farmacêutica, a emissão de ordens para reposição de itens do estoque, as
manutenções de rotina e a dispensação de medicamentos. Trabalha de
acordo com as boas práticas de manipulação e dispensação, utilizando
equipamentos de proteção individual (EPI), aplicando técnicas de
segurança e higiene pessoal e separando material para descarte, conforme
estabelecido em procedimentos operacionais padronizados.

\begin{Form}
\textbf{Esta correspondência está adequada?} \\ 
\ChoiceMenu[radio,radiosymbol=\ding{52},bordercolor=0 0.5 0,name=cbovalid_ 325110 ]{}{CBO deve ficar=sim, \hspace{6em} CBO não fica aqui=nao}
\end{Form}

\textbf{CBO: 325115  ( Técnico em farmácia )}

Atende clientes, interpretando prescrições médicas, identificando
produtos e formas farmacêuticas, calculando a quantidade necessária do
medicamento - com base na relação entre posologia, tempo de tratamento e
legislação vigente -, sob a supervisão de farmacêutico. Verifica o prazo
de validade e a integridade da embalagem do produto. Orienta o cliente
sobre o uso correto, as possíveis reações adversas e a forma de
conservação dos medicamentos. Em farmácias hospitalares, separa
medicamentos, faz a distribuição e realiza a unitarização de doses de
medicamentos. Em farmácias e drogarias, comercializa produtos de higiene
pessoal, perfumaria e cosméticos, orientando sobre o uso, as condições
de armazenamento e a forma de descarte. Acompanha lançamento de novos
produtos - em especial medicamentos, dermocosméticos e cosmecêuticos -,
para manter a oferta de produtos atualizada. Pode atuar em programas da
qualidade e em processos de acreditação no segmento farmacêutico. Pode
propor ações promocionais de produtos de higiene pessoal, perfumaria e
cosméticos, com base nas estratégias de vendas do estabelecimento,
perfil dos clientes, tipos de produtos e estoque existente. Pode
auxiliar na realização de ações de promoção da saúde no segmento
farmacêutico, elaborando recursos e materiais educativos, prestando
informações ao público e realizando registros de acompanhamento e de
avaliação das ações implementadas. Controla o estoque de medicamentos e
de produtos de higiene pessoal, perfumaria e cosméticos, encaminhando e
acompanhando pedidos de aquisição de produtos para o setor responsável,
conforme procedimentos operacionais padronizados e autorização do
farmacêutico. Em farmácias e drogarias, encaminha produtos vencidos aos
órgãos competentes. Documenta as atividades e os procedimentos da
farmácia, registrando entrada e saída de estoques; as condições de
armazenamento e os prazos de validade dos produtos; a dispensação de
medicamentos; a aplicação de injetáveis; e a compra e a venda de
medicamentos de controle especial. Realiza a manutenção de rotina do
local de trabalho, executando higienização, solicitando manutenção
preventiva e corretiva de equipamentos e encaminhando materiais
contaminados para o descarte adequado. Trabalha de acordo com as boas
práticas de manipulação, utilizando equipamentos de proteção individual
(EPI) e aplicando técnicas de segurança e higiene pessoal. Pode
identificar situação de emergência, verificando o estado da vítima, os
sinais e os sintomas apresentados. Aciona o serviço de emergência local,
repassando as informações sobre o estado da vítima.

\begin{Form}
\textbf{Esta correspondência está adequada?} \\ 
\ChoiceMenu[radio,radiosymbol=\ding{52},bordercolor=0 0.5 0,name=cbovalid_ 325115 ]{}{CBO deve ficar=sim, \hspace{6em} CBO não fica aqui=nao}
\end{Form}

\newpage
\section{ 01013 -- Técnico em Hemoterapia }

\textbf{Perfil do Curso (CNCT)}

O Técnico em Hemoterapia será habilitado para:

Realizar processos de recepção, captação e pré-triagem clínica de
doadores de sangue. Coletar, receber, preparar e processar amostras
biológicas sanguíneas, provas sorológicas e imuno-hematológicas.
Desenvolver procedimentos técnicos assistenciais em serviços e unidades
de hemoterapia. Realizar procedimentos hemoterápicos. Colaborar, como
auxiliar, em pesquisas envolvendo cultura celular. Realizar produção
industrial de hemoderivados e kits diagnósticos. Controlar a qualidade
de reagentes, produtos, insumos e equipamentos, sob supervisão de um
profissional de nível superior.

Para a atuação como Técnico em Hemoterapia, são fundamentais:

Conhecimentos das políticas públicas de saúde e compreensão de sua
atuação profissional frente às diretrizes, princípios e estrutura
organizacional do Sistema Único de Saúde (SUS). Conhecimentos e saberes
relacionados aos princípios das técnicas aplicadas na área, sempre
pautados numa postura humana e ética com doadores, pacientes, familiares
e equipes de trabalho. Comunicação pró-ativa, articulando diferentes
processos de comunicação e de informação com foco na doação de sangue e
de medula óssea. Resolução de situações-problema, trabalho em equipe e
interdisciplinar, tecnologias da informação e da comunicação, gestão de
conflitos e ética profissional. Organização e responsabilidade.
Determinação e criatividade, buscando promover a humanização da
assistência. Atualização e aperfeiçoamento profissional por meio da
Educação Continuada.

\textbf{CBOs correspondidos e seus perfis (presentes no CNCT, e presentes na predição da IA)}

\textbf{CBO: 324220  ( Técnico em hemoterapia )}

Realiza a pré-triagem de doadores e prepara o doador, fazendo ou
orientando a assepsia da região de coleta. Coleta material biológico
como sangue e hemocomponentes por meio de coleta total do sangue ou por
aférese. Identifica o material do doador, coloca conservantes,
acondiciona e armazena a coleta. Pode receber sangue, confrontando o
material com o pedido médico, conferindo as condições de análise do
material, armazenado o material em conformidade e rejeitando material
não conforme. Prepara amostra do material biológico -- sangue,
hemocomponentes e hemoderivados --, selecionando a técnica, preparando
soluções e reagentes, sequenciando, diluindo e homogeneizando as
amostras. Analisa o sangue, hemocomponentes e os hemoderivados a partir
de protocolos técnicos e procedimentos específicos. Encaminha os
resultados para o responsável pela emissão de laudo clínico. Registra
transcrição de resultado em sistemas informatizados ou analógicos.
Prepara sangue, hemocomponentes e hemoderivados para transfusão. Aplica
técnicas de descarte de resíduos biológicos e químicos. Acondiciona
material para descarte Opera equipamentos e instrumentos analíticos e de
suporte, observando seu funcionamento e solicitando manutenção preditiva
e corretiva, quando necessária. Pode auxiliar no desenvolvimento de
novos produtos derivados do sangue.

\begin{Form}
\textbf{Esta correspondência está adequada?} \\ 
\ChoiceMenu[radio,radiosymbol=\ding{52},bordercolor=0 0.5 0,name=cbovalid_ 324220 ]{}{CBO deve ficar=sim, \hspace{6em} CBO não fica aqui=nao}
\end{Form}

\newpage
\section{ 01015 -- Técnico em Imobilizações Ortopédicas }

\textbf{Perfil do Curso (CNCT)}

O Técnico em Imobilizações Ortopédicas será habilitado para:

Confeccionar, aplicar e retirar, sob a supervisão de profissionais de
nível superior, aparelhos gessados como por exemplo: talas gessadas
(goteira, calhas) e enfeixamento com uso de material convencional e
sintético. Preparar e executar trações cutâneas, auxiliando o médico
ortopedista na instalação de trações esqueléticas e nas manobras de
redução manual de fratura e luxações. Preparar sala, fora do centro
cirúrgico, para procedimento simples, tais como: pequenas suturas de
redução manual, punções e infiltrações. Reconhecer as prescrições e
aplicar a técnica de confecção e modelagem das diversas imobilizações
ortopédicas. Reconhecer estruturas anatômicas de relevância para a área.
Analisar e avaliar as condições e tipos de fraturas, luxações para
melhor escolha da imobilização.

Para a atuação como Técnico em Imobilizações Ortopédicas, são
fundamentais:

Conhecimento das políticas públicas de saúde e compreensão de sua
atuação profissional frente às diretrizes, princípios e estrutura
organizacional do Sistema Único de Saúde (SUS). Conhecimento da técnica
de retirada das imobilizações ortopédicas. Conhecimento do material
utilizado para realizar as imobilizações e seus efeitos e reações, e
utilizar o material certificando-se da segurança. Conhecimento do
processo de reabilitação de lesões ortopédicas dos pacientes.
Conhecimentos e saberes relacionados aos princípios das técnicas
aplicadas na área, sempre pautados numa postura humana e na ética do
cuidado. Organização e responsabilidade. Iniciativa social. Determinação
e criatividade, buscando promover a humanização da assistência.
Resolução de situações-problemas, gestão de conflitos, trabalho em
equipe de forma colaborativa, comunicação e ética profissional.
Atualização e aperfeiçoamento profissional por meio da Educação
Continuada.

\textbf{CBOs correspondidos e seus perfis (presentes no CNCT, e presentes na predição da IA)}

\textbf{CBO: 322605  ( Técnico de imobilização ortopédica )}

Prepara o ambiente, organizando a sala de imobilizações e providenciando
a sua limpeza. Examina o funcionamento dos equipamentos e confere a
disponibilidade de material em estoque para realização dos
procedimentos. Verifica o material, acondicionando a quantidade
necessária. Prepara o paciente, recepcionando-o e autorizando (ou não) a
entrada de acompanhante. Verifica alergias do paciente aos materiais.
Certifica-se, com o paciente, sobre o local a ser imobilizado. Explica
ao paciente o procedimento que irá ser realizado, sanando todas as
dúvidas antes da realização. Prepara o procedimento, analisando o tipo
de imobilização com base na prescrição médica; examinando as condições
da área a ser imobilizada; liberando a área a ser imobilizada de anéis e
outros ornamentos; efetuando a assepsia do local a ser imobilizado; e
posicionando o paciente. Protege o paciente com avental, lençol e biombo
ou cortina. Informa ao médico sua análise da fratura ou luxação e relata
as condições da área a ser imobilizada. Confirma a escolha do tipo de
imobilização com o médico. Higieniza as mãos e coloca as luvas de
procedimento. Confecciona aparelhos de imobilização -- tais como tala
metálica, aparelhos gessados circulares, goteiras gessadas,
enfaixamentos, colar cervical, esparadrapagem e trações cutâneas -, com
materiais sintéticos. Remove resíduos de gesso do paciente. Orienta o
paciente sobre uso e conservação da imobilização. Encaminha o paciente
ao médico para avaliação, após o período previsto de imobilização.
Recebe autorização do médico para a retirada da imobilização ortopédica.
Explica ao paciente o procedimento a ser realizado. Pode bivalvar o
aparelho gessado; remover tala e/ou goteira gessada; cortar aparelho
gessado com cizalha; retirar aparelho gessado com serra elétrica
vibratória; remover enfaixamentos; retirar aparelho gessado com bisturi
ortopédico; remover aparelho sintético; ou remover talas metálicas.
Explica os cuidados que o paciente deve ter. Auxilia o médico
ortopedista nas reduções e trações esqueléticas e em imobilizações
realizadas no centro cirúrgico. Prepara material e instrumental para
procedimentos médicos. Confirma a integridade das imobilizações dos
pacientes internados. Segue orientação do médico para reforçar aparelho
gessado; colocar salto ortopédico; fender e frisar o aparelho gessado;
abrir janela no aparelho gessado e preparar modelagem de coto. Registra
informações técnicas e elabora relatório de plantão. Transcreve
resultados de procedimentos de imobilização ortopédica em sistemas
informatizados. Dialoga tecnicamente com os profissionais das várias
áreas de saúde. Providencia serviços de manutenção de equipamentos,
quando necessário. Providencia a limpeza e a higienização do local de
trabalho. Mantém o ambiente arejado. Realiza o acondicionamento dos
materiais perfurocortantes para descarte. Trabalha com segurança,
utilizando Equipamentos de Proteção Individual (EPI) -- tais como luvas,
máscara, avental, óculos e protetor auricular -; prevenindo acidentes;
mantendo postura ergonomicamente correta; precavendo-se contra efeitos
adversos dos produtos; e submetendo-se a exames médicos periódicos.

\begin{Form}
\textbf{Esta correspondência está adequada?} \\ 
\ChoiceMenu[radio,radiosymbol=\ding{52},bordercolor=0 0.5 0,name=cbovalid_ 322605 ]{}{CBO deve ficar=sim, \hspace{6em} CBO não fica aqui=nao}
\end{Form}

\newpage
\section{ 01016 -- Técnico em Massoterapia }

\textbf{Perfil do Curso (CNCT)}

O Técnico em Massoterapia será habilitado para:

Realizar práticas massoterapêuticas visando à promoção e manutenção da
saúde, com foco no equilíbrio físico e emocional do ser humano. Avaliar,
selecionar e aplicar a técnica adequada às necessidades do cliente
baseando-se nos conceitos anatômicos, fisiológicos, biomecânicos e
fisiopatológicos. Elaborar e executar planos de trabalho sob uma
perspectiva integral e com base na utilização de técnicas manuais,
observando as indicações e contraindicações específicas para o
atendimento, bem como as normas de biossegurança e ergonomia. Associar
sua prática profissional a determinadas terapias complementares e
integrativas não invasivas.

Para a atuação como Técnico em Massoterapia, são fundamentais:

Conhecimentos das políticas públicas de saúde e compreensão de sua
atuação profissional frente às diretrizes, princípios e estrutura
organizacional do Sistema Único de Saúde (SUS). Conhecimentos e saberes
relacionados aos princípios das técnicas aplicadas na área, sempre
pautados numa postura humana e ética. Respeitar as contraindicações das
técnicas em face das condições do cliente. Resolução de
situações-problema, comunicação, trabalho em equipe e interdisciplinar,
gestão de conflitos e ética profissional. Organização, responsabilidade,
iniciativa social, determinação e criatividade, buscando promover a
humanização da assistência. Atualização e aperfeiçoamento profissional
por meio da educação continuada.

\textbf{CBOs correspondidos e seus perfis (presentes no CNCT, e presentes na predição da IA)}

\textbf{CBO: 322120  ( Massoterapeuta )}

Prepara o ambiente de trabalho, higienizando-o para receber o cliente.
Organiza os prontuários e a agenda de trabalho com base em informações
sobre os clientes e atendimentos a serem realizados. Interpreta
relatórios e/ou prescrições médicas que orientam o seu trabalho.
Recepciona o cliente, utilizando técnicas de atendimento adequadas aos
diferentes perfis. Fornece informações sobre os serviços
massoterapêuticos oferecidos. Avalia as condições de saúde e verifica os
hábitos de vida do cliente, questionando-o sobre suas queixas,
expectativas e eventuais disfunções e patologias. Inspeciona e palpa a
pele e os grupos musculoesqueléticos e articulares. Identifica aspectos
biomecânicos do cliente, por meio da observação postural. Avalia a
condição dos tecidos moles do cliente, a qualidade e a função das
articulações, a força muscular e a amplitude de movimento. Realiza
avaliação energética do cliente. Identifica indicações e eventuais
contraindicações para o atendimento massoterapêutico, com base nas
informações levantadas. Define o plano de tratamento para o cliente,
especificando quais tipos de massagem devem ser usados. Pode consultar
outros profissionais de saúde - como fisioterapeutas, médicos e
psicólogos -, para tratamento multidisciplinar. Pode encaminhar o
cliente a outros profissionais, quando necessário. Preenche em sistema
informatizado a ficha de do cliente, detalhando a avaliação de saúde e
especificando o plano de tratamento. Prepara e mistura os óleos e aplica
as misturas na pele do cliente. Realiza procedimentos massoterapêuticos
orientais e ocidentais de relaxamento e estimulação muscular, de
remodelação corporal, de estímulo ao sistema linfático, de prevenção de
lesões e recuperação muscular no esporte. Aplica procedimentos
massoterapêuticos em microssistemas, tais como execução da reflexologia
podal, da reflexologia manual e da auriculoterapia. Pode aplicar
práticas complementares e integrativas não invasivas -- tais como
cromoterapia e aromaterapia - associadas às massagens. Presta
informações ao cliente sobre sua condição, orientando-o sobre medidas
preventivas e recomendando exercícios. Conclui o atendimento, definindo
preço do serviço de massoterapia e providenciando o recebimento do
pagamento. Elabora ações de divulgação e venda de serviços de
massoterapia, de acordo com o público-alvo. Controla o estoque de
produtos, materiais e utensílios, armazenando-os, conforme as indicações
do fabricante e a legislação vigente. Conserva o ambiente de trabalho
limpo, higienizado e organizado, providenciando a assepsia de
mobiliários e a esterilização do instrumental, de acordo com as normas
de biossegurança. Providencia serviço de manutenção de equipamentos,
quando necessário. Descarta material e produtos com validade vencida,
seguindo os princípios ambientais e as normas de biossegurança.
Acondiciona materiais perfurocortantes para descarte. Trabalha com
segurança, adotando postura ergonômica correta e utilizando equipamentos
de proteção individual.

\begin{Form}
\textbf{Esta correspondência está adequada?} \\ 
\ChoiceMenu[radio,radiosymbol=\ding{52},bordercolor=0 0.5 0,name=cbovalid_ 322120 ]{}{CBO deve ficar=sim, \hspace{6em} CBO não fica aqui=nao}
\end{Form}

\newpage
\section{ 01017 -- Técnico em Meio Ambiente }

\textbf{Perfil do Curso (CNCT)}

O Técnico em Meio Ambiente será habilitado para:

Coletar, armazenar e interpretar informações, dados e documentações
ambientais. Auxiliar na elaboração, na análise de projetos, nos
relatórios e estudos ambientais. Propor medidas para a minimização dos
impactos e recuperação de ambientes já degradados. Executar sistemas de
gestão ambiental. Organizar programas de educação ambiental com base no
monitoramento, na correção e prevenção das atividades antrópicas, na
conservação dos recursos naturais através de análises prevencionistas.
Organizar redução, reuso e reciclagem de resíduos e/ou recursos
utilizados em processos. Identificar os padrões de produção e consumo de
energia. Realizar levantamentos ambientais. Operar sistemas de
tratamento de poluentes e resíduos sólidos. Relacionar os sistemas
econômicos e suas interações com o meio ambiente. Realizar e coordenar o
sistema de coleta seletiva. Executar plano de ação e manejo de recursos
naturais. Elaborar relatório periódico das atividades e modificações dos
aspectos e impactos ambientais de um processo, indicando as
consequências de modificações. Realizar ações de saúde ambiental nos
territórios. Desenvolver tecnologias sociais ambientais. Promover ações
de manejo ambiental. Avaliar e monitorar sistema de tratamento e
abastecimento de água, bem como de esgotamento sanitário. Monitorar os
indicadores de qualidade do ar atmosférico. Executar ações de controle e
manejo da poluição. Realizar vistoria ambiental e sanitária. Realizar
monitoramento ambiental. Elaborar diagnóstico das condições
socioambientais, econômicas e culturais. Identificar e intervir nos
problemas de saúde relacionados aos fatores de riscos ambientais do
território com o propósito de contribuir para a melhoria da qualidade de
vida da população. Conhecer e utilizar sistemas de informação
geográficas para uso em atividades de geoprocessamento no trabalho
ambiental. Integrar ações da saúde do trabalhador com saúde ambiental.
Conhecer e integrar o sistema de saneamento ambiental bem como sua
relação com a saúde pública. Auditar sistemas de gestão ambiental. Atuar
nas áreas de educação, proteção e recuperação ambientais.

Para a atuação como Técnico em Meio Ambiente, são fundamentais:

Conhecimentos das políticas públicas de Meio Ambiente e compreensão de
sua atuação profissional frente às diretrizes, princípios e estrutura
organizacional do Sistema Nacional de Meio Ambiente (SISNAMA).
Conhecimentos das políticas públicas de saúde e compreensão de sua
atuação profissional frente às diretrizes, princípios e estrutura
organizacional do Sistema Único de Saúde (SUS). Conhecimentos e saberes
relacionados a processos de sustentabilidade, territorialização e
monitoramento ambiental. Organização, responsabilidade, resolução de
situações-problema, gestão de conflitos, trabalho em equipe de forma
colaborativa, comunicação e ética profissional. Atualização e
aperfeiçoamento profissional por meio da educação continuada. Visão
abrangente e integrada dos tópicos ambientais (água, ar, solo, fauna e
flora) e sua dinâmica. Orientação e controle de processos voltados às
áreas de conservação, pesquisa, proteção e defesa ambiental. Atuar em
equipes de gerenciamento ambiental de órgãos públicos e privados.

\textbf{CBOs correspondidos e seus perfis (presentes no CNCT, e presentes na predição da IA)}

\textbf{CBO: 311505  ( Técnico de controle de meio ambiente )}

Prepara ambiente para análises, selecionando reagentes, vidrarias e
equipamentos. Identifica as amostras de efluentes coletadas, registrando
os pontos de coleta. Realiza testes e análises físico-químicas e
microbiológicas de efluentes. Interpreta resultados analíticos e elabora
laudos, relatórios e planilhas dos resultados. Pode encaminhar amostras
para análises externas complementares. Coordena processos de medição
para controle ambiental e tratamento de efluentes, monitorando os
parâmetros do processo e emitindo relatórios, para identificar os
aspectos ambientais e os impactos associados. Monitora a qualidade e o
uso das águas, bem como as características das emissões atmosféricas,
dos efluentes e dos resíduos sólidos, conforme manuais e procedimentos
técnicos. Pode coletar, orientar ou coordenar a coleta de amostras de
gases, solo, água e outros materiais contaminados, inclusive por
radiação. Implementa ações preventivas e corretivas, participando de
auditorias e propondo medidas para a mitigação dos impactos e
recuperação de ambientes já degradados. Executa tarefas de caráter
técnico referentes às questões ambientais da unidade fabril,
controlando, organizando e monitorando os resíduos e os efluentes e
verificando o licenciamento ambiental. Pode orientar o aproveitamento
dos resíduos industriais, buscando valorização econômica, ecoeficiência
e sustentabilidade do processo. Participa da implantação de programas de
qualidade ambiental, analisando indicadores e aplicando ferramentas da
qualidade. Desenvolve programas de educação ambiental junto à comunidade
e às instituições. Realiza campanhas de monitoramento e educação
ambiental, identificando tecnologias apropriadas para o processo de
produção racional e cuidados com o meio ambiente. Pode participar de
equipes de desenvolvimento de materiais biodegradáveis. Monitora o
funcionamento de máquinas, equipamentos e instrumentos de medição e
controle, mantendo-os em condições de uso e realizando ajustes no
processo, quando necessários. Monitora a segurança no trabalho,
interpretando mapas de risco; controlando o uso de equipamentos de
proteção individual e coletiva; e cumprindo procedimentos de emergência,
quando necessário.

\begin{Form}
\textbf{Esta correspondência está adequada?} \\ 
\ChoiceMenu[radio,radiosymbol=\ding{52},bordercolor=0 0.5 0,name=cbovalid_ 311505 ]{}{CBO deve ficar=sim, \hspace{6em} CBO não fica aqui=nao}
\end{Form}

\newpage
\section{ 01018 -- Técnico em Meteorologia }

\textbf{Perfil do Curso (CNCT)}

O Técnico em Meteorologia será habilitado para:

Dar suporte direto ao Meteorologista, atuando desde a coleta de dados no
campo, na codificação, na decodificação e no registro dos elementos de
observação meteorológica, até a disponibilização desses dados com
qualidade e clareza. Realizar a leitura, a codificação, a decodificação
e o registro dos elementos de observação meteorológica. Analisar e
interpretar dados meteorológicos, produtos de satélites e de resultados
de modelos meteorológicos, organizando-os em bancos de dados. Proceder
com a instalação, a operação, a aferição e a manutenção de estações
meteorológicas e desenvolver sistemas computacionais para tratamento e
divulgação de informações meteorológicas. Atuar no apoio de atividades
relacionadas às áreas de agricultura, de energia, do meio ambiente, dos
recursos hídricos, da saúde, da defesa, do transporte, da construção
civil, dentre muitas outras.

Para a atuação como Técnico em Meteorologia, são fundamentais:

Conhecimentos e saberes relacionados aos princípios das técnicas
aplicadas na área, sempre pautados numa postura humana e ética.
Resolução de situações-problema, comunicação, trabalho em equipe e
interdisciplinar, tecnologias da informação e da comunicação, gestão de
conflitos e ética profissional. Organização, responsabilidade e
iniciativa social. Atualização e aperfeiçoamento profissional por meio
da educação continuada.

\textbf{CBOs correspondidos e seus perfis (presentes no CNCT, e presentes na predição da IA)}

\textbf{CBO: 311510  ( Técnico de meteorologia )}

Instala, opera, afere e mantém estações meteorológicas de superfície e
de altitude, sinóticas, climatológicas, agrometeorológicas, aeronáuticas
e especiais, tais como marítimas e fluviométricas. Realiza leitura,
interpreta e gera informações meteorológicas e climatológicas, por meio
de equipamentos específicos e de estações-radar meteorológicas, de
recepção de imagens de satélites e de radiodifusão. Controla a qualidade
das observações e dos dados meteorológicos. Participa da organização de
banco de dados meteorológicos, para previsões de curto e longo prazos.
Realiza estudos, participa de pesquisas e emite relatórios de impacto
ambiental, de diagnóstico da poluição do ar e de prevenção e dispersão
dos poluentes atmosféricos. Participa da preparação de relatórios ou de
mapas climáticos para análise, distribuição ou uso em transmissões
climáticas, utilizando softwares específicos. Emprega técnicas de
sensoriamento remoto a fim de gerar informações de interesses
meteorológicos para finalidades agrícolas, industriais, de turismo e de
lazer. Na área militar e de defesa, observa e descreve as condições
meteorológicas para o planejamento das missões terrestres, aéreas e
aquaviárias, indispensáveis à segurança de navegabilidade. Realiza
leitura, codificação, decodificação e registro dos elementos de
observação meteorológica necessários ao planejamento e à segurança da
navegação aérea e aquaviária. Aplica sistemas e métodos computacionais
para tratamento e divulgação de informações meteorológicas, realizando
prognósticos e emitindo boletins meteorológicos. Monitora o
funcionamento de máquinas, equipamentos e instrumentos utilizados,
solicitando manutenção periódica. Pode calibrar equipamentos e
instrumentos. Zela pela segurança no trabalho, interpretando mapa de
riscos; usando equipamentos de proteção individual e coletiva; e
cumprindo procedimentos de emergência, quando necessário.

\begin{Form}
\textbf{Esta correspondência está adequada?} \\ 
\ChoiceMenu[radio,radiosymbol=\ding{52},bordercolor=0 0.5 0,name=cbovalid_ 311510 ]{}{CBO deve ficar=sim, \hspace{6em} CBO não fica aqui=nao}
\end{Form}

\newpage
\section{ 01019 -- Técnico em Necropsia }

\textbf{Perfil do Curso (CNCT)}

O Técnico em Necropsia será habilitado para:

Auxiliar o perito médico legista e o médico patologista na realização de
exames necroscópicos. Executar técnicas de tanatopraxia visando à
conservação e reconstituição de cadáveres para traslados e cerimônias
póstumas. Preparar, por meio da dissecação e manipulação de soluções
conservantes, corpos e peças anatômicas para ensino e pesquisa. Prestar
serviço funeral aos responsáveis legais dos falecidos.

Para a atuação como Técnico em Necropsia, são fundamentais:

Conhecimento das políticas públicas de saúde: organização, princípios e
diretrizes do Sistema Único de Saúde (SUS). Conhecimentos e saberes
relacionados às diversas áreas de atuação que envolvem questões
pós-morte e suas consequências, sempre pautado numa postura humana e
ética. Resolução de situações-problemas, trabalho em equipe, tecnologias
da informação e da comunicação, gestão de conflitos e ética
profissional. Organização e responsabilidade; iniciativa social;
entusiasmo; empatia e respeito. Atualização e aperfeiçoamento
profissional por meio da Educação Continuada.

\textbf{CBOs correspondidos e seus perfis (presentes no CNCT, e presentes na predição da IA)}

\textbf{CBO: 328105  ( Embalsamador )}

Prepara as atividades de embalsamento, providenciando a limpeza e a
desinfecção das áreas onde são feitas a etapa preparatória e a etapa de
embalsamento dos corpos. Confere a disponibilidade dos materiais
específicos a serem utilizados, requisitando a compra dos faltantes.
Seleciona equipamentos, instrumentos e utensílios. Faz a identificação
do cadáver a ser embalsamado. Avalia as condições do cadáver,
encaminhando casos de morte violenta para o Instituto Médico Legal
(IML). Verifica se os restos mortais serão transportados a locais
distantes, outros estados ou exterior, para programar os procedimentos
especiais necessários. Registra as informações sobre o corpo a ser
embalsamado, tais como identificação, peso e altura. Faz a limpeza e a
higienização do corpo, lavando-o com uso de produtos químicos
específicos - como desinfetantes e germicidas -- para eliminar
bactérias, que podem causar a aceleração da decomposição. Seca o corpo
com secadores de ar quente. Observa o enrijecimento muscular do cadáver,
utilizando técnicas -- como massagem - para relaxar os músculos. Faz o
cálculo da quantidade de fluido embalsamador a ser utilizada, levando em
conta o peso do corpo. Executa a drenagem do sangue, fazendo uma incisão
e inserindo um cano ligado a uma máquina, que bombeia o sangue para fora
do corpo. Injeta fluido embalsamador - solução química normalmente
composta de água e formaldeído -, que preenche os espaços, para
prolongar a conservação do corpo e evitar que os órgãos desidratem e se
decomponham. Efetua nova incisão e insere instrumento cirúrgico acoplado
a um aspirador, para retirar gases do corpo. Realiza o fechamento das
incisões. Faz nova lavagem e desinfecção do corpo. Pode preparar cadáver
necropsiado para sua reconstituição, retirando vísceras e suturando o
corpo. Reconstitui o cadáver, retirando deformidades das mãos e da face
para restabelecer a forma original; recompondo o globo ocular; cerrando
bocas e olhos; e tamponando orifícios, para restituir-lhes o aspecto
normal. Restaura cadáver desfigurado e mutilado. Veste o corpo. Pode
fazer a necromaquiagem, utilizando creme hidratante e aplicando
maquiagem para recuperar aparência mais natural e esconder manchas na
pele. Prepara urna e nela insere o corpo. Pode embalar urna para viagem,
conforme exigência de companhias aéreas. Presta orientação para
estagiários e outras pessoas interessadas na área de embalsamamento de
corpos, ministrando aulas práticas em minicursos e proferindo palestras.
Pode desenvolver projetos específicos de embalsamamento de corpos para
museus. Mantém acervo de peças anatômicas conservadas para fins de
estudo, etiquetando peças e administrando banco de dados com registros
dos exemplares. Coleta material para pesquisa, desenvolvendo ferramentas
e acessórios específicos para essa atividade. Troca informações
técnicas. Conserva ferramentas, instrumentos e utensílios limpos,
esterilizados e acondicionados. Quando necessário, requisita serviços de
manutenção de equipamentos. Aplica procedimentos para cumprimento de
normas ambientais, adotando orientações de órgãos competentes para o
tratamento e a destinação de materiais infectados, de resíduos químicos,
de materiais biológicos e de outros resíduos descartados. Atua seguindo
as normas de segurança e de biossegurança, usando equipamentos de
proteção individual; enviando vestimenta de trabalho para a
descontaminação; tomando banho no término do expediente; lavando as mãos
após o manuseio de material; e tomando vacinas. Segue orientação para o
uso de soluções químicas, respeitando as normas de utilização.

\begin{Form}
\textbf{Esta correspondência está adequada?} \\ 
\ChoiceMenu[radio,radiosymbol=\ding{52},bordercolor=0 0.5 0,name=cbovalid_ 328105 ]{}{CBO deve ficar=sim, \hspace{6em} CBO não fica aqui=nao}
\end{Form}

\textbf{CBO: 328110  ( Taxidermista )}

Planeja as atividades de taxidermia, verificando se o corpo inteiro do
animal morto passará pelo processo ou apenas parte, como cabeça ou
cabeça e ombros. Providencia a limpeza e a higienização do ambiente de
trabalho. Seleciona equipamentos, instrumentos e utensílios, de acordo
com a espécie animal. Pode escolher ``software'' para auxiliar no estudo
da recomposição da estrutura do animal a ser taxidermizado e para a
criação de dioramas. Prepara soluções químicas. Cria ficha individual do
espécime animal, registrando peso; comprimento; circunferência do
abdome; e demais medidas. Avalia as condições do cadáver animal,
examinando a estrutura externa, verificando a presença de parasitas e
executando outras ações estabelecidas em protocolo de taxidermia da
espécie animal a ser conservada. Tira fotos do animal. Faz máscara
mortuária do animal morto, que fornece a imagem tridimensional das suas
feições e os detalhes de seu rosto. Coloca o corpo do animal na postura
correta -- com uso de arames, fios e apoios -, para que sua estrutura
corporal seja copiada, primeiramente numa forma de gesso e, depois, em
molde de resina. Usa o molde para criar um manequim - feito de gesso,
fibra, poliuretano ou outro material -, que reproduza o animal com a
maior fidelidade possível. Limpa e prepara o animal morto para separação
da pele, escamas, penas e outros elementos externos a serem usados.
Descarta órgãos e carcaça. Seleciona o processo de curtimento a ser
realizado, considerando que cada espécie animal possui pele
diferenciada, exigindo fórmulas próprias. Realiza o processo de
curtimento da pele. Executa banhos ácidos, que dissolvem resquícios de
sujeira e gordura, e banhos com produto químico preservativo. Usa
técnicas -- como engraxamento e amaciamento - para a pele ganhar
flexibilidade. Inicia a montagem, fixando as próteses sintéticas - de
olhos, nariz, orelhas e outras partes -- no manequim. Pode usar as fotos
tiradas, para reproduzir as formas e as características do animal.
Aplica cola especial no manequim, para adesão da pele. Modela a pele em
torno do manequim e faz o fechamento do revestimento com costura ou
colagem de adesivo. Pode fazer cobertura das partes sem pele e sem pelo
com tintas próprias para taxidermia, aplicadas com um aerógrafo. Faz a
secagem da peça concluída. Confecciona dioramas, pesquisando
características do animal e de seu habitat; montando projeto de
ambientação; preparando material e produzindo as réplicas de elementos
da natureza, para apresentar o espécime em cenário parecido com o local
onde ele vivia. Pode aplicar outras técnicas de taxidermia, como a
técnica de diafanização, que trata esqueleto e outras partes internas,
para mostrá-las visíveis em animal translúcido. Presta orientação a
estagiários e outras pessoas, mostrando museu de animais taxidermizados
a visitantes; ministrando aulas práticas; e proferindo palestras. Pode
desenvolver projetos para museus. Gerencia atividades comerciais,
comprando material e administrando fluxo de caixa. Adquire licenciamento
em órgãos competentes. Despacha animais taxidermizados para outros
estados ou países. Faz divulgação da taxidermia. Mantém acervo
científico, acompanhando pesquisadores em trabalho de campo e coletando
material. Troca informações técnicas. Administra banco de dados com
informação dos exemplares. Realiza manutenção de acervo. Pode
supervisionar trabalho da equipe, avaliando desempenho e desenvolvendo
treinamentos. Conserva ferramentas, instrumentos e utensílios limpos,
esterilizados e acondicionados. Requisita serviço de manutenção de
equipamentos. Realiza tratamento e destinação de resíduos de acordo com
normas ambientais. Atua de acordo com normas de segurança e de
biossegurança, usando equipamentos de proteção individual;
descontaminando vestimenta de trabalho; tomando banho após expediente;
lavando as mãos após manusear material; e tomando vacinas. Segue
orientação para o uso de soluções químicas, respeitando as normas de
utilização.

\begin{Form}
\textbf{Esta correspondência está adequada?} \\ 
\ChoiceMenu[radio,radiosymbol=\ding{52},bordercolor=0 0.5 0,name=cbovalid_ 328110 ]{}{CBO deve ficar=sim, \hspace{6em} CBO não fica aqui=nao}
\end{Form}

\textbf{CBO: 516505  ( Agente funerário )}

Presta primeiro atendimento à família do falecido, acolhendo-a e
identificando a circunstância e local do óbito. Verifica quem será o
contratante dos serviços funerários. Orienta a família quanto aos
procedimentos legais e solicita-lhe documentação do falecido e
declaração de óbito. Contata órgão competente para a liberação do corpo
e comunica à família o horário previsto. Consulta a família sobre sua
religião, sobre o velório e sobre a realização de sepultamento ou
cremação. Pode redigir nota para rádio e ou jornal, comunicando o
falecimento. Informa à família as opções de serviços funerários.
Solicita roupa para o falecido. Mede peso e altura do falecido. Mostra
as opções de urna mortuária, de acordo com as medidas, para escolha. Em
casos especiais -- como risco de contágio -, segue os procedimentos
legais de urna lacrada. Providencia o local do velório, a cerimônia
religiosa e o local para sepultamento ou cremação. Fornece as medidas da
urna para o cemitério. Encaminha o corpo para a preparação. Providencia
a preparação do corpo, avaliando seu estado e verificando a causa
mortis. Acompanha a higienização do corpo; a colocação de algodão para
tamponar o nariz e a boca; e a aplicação de material conservante, por
formolização, embalsamamento ou tanatopraxia. Veste o corpo. Providencia
a execução de restauração facial e de necromaquiagem. Ornamenta a urna.
Encaminha o corpo para o cerimonial. Pode trasladar o corpo, para fora
do município. Organiza o cerimonial, de acordo com a religião. Ornamenta
o local do velório e mantém a sua organização. Checa as alterações do
corpo. Pode fornecer serviços de copa. Pode alertar a família sobre a
necessidade de antecipação do sepultamento. Transporta a urna em veículo
funerário, conduzindo o cortejo fúnebre ao cemitério, onde o corpo será
sepultado. Pode transportar a urna ao crematório, para processo de
cremação do corpo. Prepara a documentação necessária. Confere e
complementa a declaração de óbito. Emite documentos para cartório e guia
de sepultamento para cemitério. Expede ata de embalsamamento,
formolização ou tanatopraxia. Registra histórico do falecido (Serviço de
Verificação de Óbito). Obtém documentação para traslado ou cremação.
Trabalha com segurança, utilizando equipamentos de proteção individual
-- tais como luvas, máscaras e aventais -- e tomando vacinas. Mantém
assepsia na sala de preparação e esteriliza os materiais utilizados,
considerando protocolos específicos de saúde - relacionados a patógenos
e perigos de contaminação - e procedimentos de higiene e biossegurança.
Acondiciona material para descarte, seguindo normas ambientais.

\begin{Form}
\textbf{Esta correspondência está adequada?} \\ 
\ChoiceMenu[radio,radiosymbol=\ding{52},bordercolor=0 0.5 0,name=cbovalid_ 516505 ]{}{CBO deve ficar=sim, \hspace{6em} CBO não fica aqui=nao}
\end{Form}

\newpage
\section{ 01020 -- Técnico em Nutrição e Dietética }

\textbf{Perfil do Curso (CNCT)}

O Técnico em Nutrição e Dietética será habilitado para:

Desenvolver, sob a supervisão de profissionais de nível superior,
atividades relacionadas à educação alimentar e nutricional de indivíduos
e comunidades, para prevenção e controle de carências nutricionais, de
doenças crônicas não transmissíveis e de doenças veiculadas por
alimentos. Realizar coleta de dados de interesse ao Serviço de Nutrição
e Dietética, bem como dados antropométricos para subsidiar a avaliação
nutricional a ser realizada pelo nutricionista. Realizar estudos das
necessidades nutricionais de indivíduos e coletividades, em todas as
fases do ciclo vital. Monitorar dietas de rotina sobre prescrição
dietética. Acompanhar e orientar a execução das atividades que compõe
toda a escala de produção de refeições, bem como as atividades de
controle de qualidade higiênico-sanitárias, atendendo às normas de
Segurança Alimentar e Nutricional. Auxiliar no planejamento e na
execução dos procedimentos de rotina, bem como orientar e monitorar as
atividades realizadas pela equipe de funcionários. Aplicar normas de
segurança do trabalho na produção de refeições e no comércio de
alimentos. Contribuir na elaboração de cardápios e na elaboração de
relatórios técnicos de não conformidades. Trabalhar também na pesquisa e
no desenvolvimento de novos produtos alimentícios. Atuar em diferentes
segmentos, sob orientação e supervisão do nutricionista.

Para a atuação como Técnico em Nutrição e Dietética, são fundamentais:

Conhecimentos das políticas públicas de saúde e compreensão de sua
atuação profissional frente às diretrizes, princípios e estrutura
organizacional do Sistema Único de Saúde (SUS). Conhecimentos e saberes
relacionados aos princípios das técnicas aplicadas na área, sempre
pautados numa postura humana e ética. Resolução de situações-problema,
comunicação, trabalho em equipe e interdisciplinar, tecnologias da
informação e da comunicação, gestão de conflitos e ética profissional.
Organização e responsabilidade. Iniciativa social. Determinação e
criatividade, buscando promover a humanização da assistência.
Atualização e aperfeiçoamento profissional por meio da Educação
Continuada.

\textbf{CBOs correspondidos e seus perfis (presentes no CNCT, e presentes na predição da IA)}

\textbf{CBO: 325210  ( Técnico em nutrição e dietética )}

Planeja as atividades, identificando público-alvo, estabelecendo o
cronograma de atividades, verificando a disponibilidade dos materiais,
insumos e gêneros alimentícios e selecionando os procedimentos para cada
atividade. Elabora receituário padrão, sob supervisão de profissional de
nível superior da área, coletando dados antropométricos para subsidiar a
avaliação nutricional, realizando inquérito alimentar e participando do
cálculo do valor calórico do cardápio de indivíduos sadios. Orienta
sobre o preparo de refeições, a composição dos alimentos e a ação dos
nutrientes. Supervisiona os processos de produção de alimentos e
refeições, recebendo insumos, produtos e gêneros alimentícios;
acompanhando o pré-preparo e o preparo de alimentos; monitorando
processos -- tais como de trituração, pasteurização, mistura,
fermentação e cocção -; controlando o tempo de produção; e coordenando
atividades de porcionamento, armazenamento, transporte e distribuição de
alimentos e de refeições. Controla a qualidade nas etapas de produção de
alimentos e refeições, preenchendo checklist - verificação de rotinas -,
identificando e corrigindo pontos críticos de controle. Coleta amostras,
acompanhando testes de desempenho e realizando as análises sensoriais
das matérias-primas e dos produtos. Verifica condições de armazenamento,
controlando data de vencimento dos produtos e acompanhando o controle
integrado de pragas e vetores. Avalia fornecedores, acompanhando testes
de desempenho de matérias-primas e insumos e realizando análise
comparativa com produtos concorrentes. Vistoria as instalações,
assegurando condições higiênico-sanitárias. Participa, sob supervisão de
profissional de nível superior da área, nas pesquisas e nos projetos
para melhoria, adequação e desenvolvimento de produtos, acompanhando as
necessidades do mercado, definindo estratégias e participando na
elaboração e na testagem de formulação dos produtos. Avalia a
aceitabilidade do produto e assessora a implementação das mudanças
aprovadas. Participa na elaboração da informação nutricional, que consta
na rotulagem do produto. Participa na divulgação dos aspectos
nutricionais do produto, orientando clientes internos e externos. Pode
promover a venda de insumos e produtos, levantando e interpretando
necessidades de clientes potenciais, realizando visitas técnicas para
apresentar novos produtos e prestando assistência técnica aos clientes.
Realiza atividades de educação nutricional, sob supervisão de
profissional de nível superior da área, identificando comportamentos
alimentares presentes em uma comunidade, por meio de estudos, análises e
pesquisas. Auxilia no planejamento de programas e elabora material
educativo relacionado à saúde, à alimentação e à nutrição, de acordo com
as diretrizes da Política Nacional de Alimentação e Nutrição. Participa
na execução de programas de educação alimentar e nutricional. Avalia e
registra os resultados dos programas educacionais, elaborando relatório.
Coordena e supervisiona equipes de trabalho atuantes no processo de
produção de refeições coletivas e dietas, participando na seleção de
pessoal, avaliando os resultados de desempenho e realizando treinamento
de rotinas operacionais. Controla as condições do ambiente de trabalho,
monitorando a temperatura. Verifica a segurança no local de trabalho,
conferindo o uso de equipamentos de proteção individual. Controla
desperdícios e destina adequadamente os resíduos gerados, fazendo
reciclagem e realizando descarte de acordo com as normas ambientais.

\begin{Form}
\textbf{Esta correspondência está adequada?} \\ 
\ChoiceMenu[radio,radiosymbol=\ding{52},bordercolor=0 0.5 0,name=cbovalid_ 325210 ]{}{CBO deve ficar=sim, \hspace{6em} CBO não fica aqui=nao}
\end{Form}

\newpage
\section{ 01021 -- Técnico em Óptica }

\textbf{Perfil do Curso (CNCT)}

O Técnico em Óptica será habilitado para:

Ser responsável técnico pelos estabelecimentos ópticos, centro de
adaptações de lentes de contato e laboratórios de surfaçagem e montagem.
Atuar como consultor óptico e como representante comercial de
equipamentos, de armações, de lentes oftálmicas e de contato. Emitir
laudos e pareceres técnicos relacionados aos produtos ópticos. Executar
fabricação de lentes em geral, montagem de óculos e adaptação de lentes
de contato. Implementar ações de gestão administrativa e estratégica no
segmento óptico. Indicar e comercializar produtos ópticos (lentes
oftálmicas, lentes de contato e armações para óculos) de acordo com a
dioptria (grau), com as medidas individuais, com a anatomia facial, com
a estética e os costumes comportamentais e laborais do usuário.
Interpretar, avaliar e aviar prescrições ópticas e laudos optométricos.
Realizar testes de acuidade visual de óptica oftálmica e contatologia.

Para a atuação como Técnico em Óptica, são fundamentais:

Conhecimento das políticas públicas de saúde: organização, princípios e
diretrizes do Sistema Único de Saúde (SUS). Conhecimentos e saberes
relacionados à profissão, sempre pautado numa postura humana e ética.
Resolução de situações-problemas, trabalho em equipe, tecnologias da
informação e da comunicação, gestão de conflitos e ética profissional.
Organização e responsabilidade; iniciativa social; entusiasmo; empatia e
respeito. Atualização e aperfeiçoamento profissional por meio da
Educação Continuada. Conhecimentos e saberes relacionados aos princípios
das técnicas aplicadas na área, voltadas à surfaçagem de lentes,
montagem de óculos, adaptação de lentes de contato.

\textbf{CBOs correspondidos e seus perfis (presentes no CNCT, e presentes na predição da IA)}

\textbf{CBO: 322305  ( Técnico em óptica e optometria )}

Planeja a fabricação das lentes de visão simples, especiais, bifocais e
multifocais - por meio dos processos de surfaçagem -, com base nas
especificações de ordem de serviço. Seleciona insumos, blocos,
equipamentos e ferramentas. Efetua cálculos de curvatura e espessura de
lentes. Executa a confecção de lentes, utilizando equipamentos e
softwares específicos. Dá acabamento às lentes, fazendo polimento e
adicionando tratamentos às lentes, como endurecimento, antirreflexo e
filtros. Faz a conferência das lentes, utilizando instrumentos
específicos, tais como lensômetro e espessímetro. Realiza processo de
montagem de lentes em armação com aro, armação de semiaro e sem aro,
conferindo a dioptria da lente, utilizando o lensômetro e marcando o
centro óptico e o eixo, de acordo com as especificações da ordem de
serviço. Executa a montagem dos óculos. Confere os óculos, de acordo com
as especificações da ordem de serviço. Pode confeccionar óculos de
segurança. Avalia a viabilidade (ou não) do uso de lentes de contato,
classificando as ametropias - conforme a prescrição do especialista - e
analisando as informações de anamnese, dos testes de filme lacrimal e
das tomadas de medidas. Adapta lentes de contato, definindo o material e
o desenho das lentes de contato; fazendo cálculos dos parâmetros;
selecionando as lentes de contato de teste; realizando assepsia das
lentes; testando o uso das lentes de contato; avaliando os resultados
das lentes de contato de teste; e executando os ajustes necessários.
Fornece instruções de uso de lentes de contato aos usuários. Aplica
prótese ocular, analisando a cavidade óptica para moldar e determinar as
características -- tamanho, cor, entre outras - da prótese. Confecciona
e ajusta a prótese ocular. Avalia e, quando necessário, readapta a
prótese. Realiza exames optométricos, posicionando e orientando o
cliente para realização de testes, de acordo com os procedimentos de
padrões técnicos. Mede acuidade visual, analisa estruturas externas e
internas dos olhos, mede pressão intraocular (tonometria), identifica
deficiências e anomalias relacionadas às alterações da função visual e
mede refração ocular (refratometria e retinoscopia). Faz diagnóstico
primário de enfermidades que acometem a saúde visual, encaminhando casos
patológicos a médicos especializados. Define compensações e auxílios
ópticos. Recepciona e atende cliente para o consumo de produtos e
serviços ópticos. Interpreta prescrição óptica do especialista,
identificando a ametropia do cliente. Sugere a armação mais adequada,
considerando estilo, atividade laboral e tipo físico do usuário. Indica
a lente mais adequada, de acordo com a armação escolhida e a prescrição
óptica do especialista. Ajusta óculos no rosto de cliente. Conserta
auxílios ópticos. Mantém registros de cliente. Gerencia estabelecimento
de artigos de óptica, administrando compras e vendas; controlando
estoque de mercadorias e materiais; providenciando manutenção de
equipamentos e instalações; administrando finanças; e controlando
qualidade de produtos e serviços. Administra pessoal, participando de
seleção, atribuindo tarefas e acompanhando registros de faltas e
atrasos. Supervisiona trabalho da equipe, avaliando desempenho e
promovendo programas de treinamento. Elabora pareceres e laudos técnicos
de produtos ópticos, conferindo óculos e lentes de contato; indicando
causas para não conformidade; e apresentando possíveis soluções. Redige
documentos de responsabilidade técnica. Promove ações de educação em
saúde visual, planejando ações de acordo com as necessidades do
público-alvo, ministrando palestras e cursos para a promoção de saúde e
prevenção de doenças visuais. Orienta o descarte de resíduos -- em
especial os gerados durante o processo de fabricação de lentes
oftálmicas -, de acordo com normas ambientais. Trabalha com segurança,
utilizando Equipamento de Proteção Individual (EPI), analisando mapa de
risco e prevenindo acidentes.

\begin{Form}
\textbf{Esta correspondência está adequada?} \\ 
\ChoiceMenu[radio,radiosymbol=\ding{52},bordercolor=0 0.5 0,name=cbovalid_ 322305 ]{}{CBO deve ficar=sim, \hspace{6em} CBO não fica aqui=nao}
\end{Form}

\newpage
\section{ 01022 -- Técnico em Optometria }

\textbf{Perfil do Curso (CNCT)}

O Técnico em Optometria será habilitado para:

Realizar exames optométricos. Adaptar auxílios ópticos para baixa visão.
Adaptar lentes de contato. Emitir laudos optométricos. Fomentar ações de
cuidado em saúde visual. Identificar e analisar os fatores determinantes
da saúde visual, para desenvolver, promover e executar ações que
permitam seu controle e acompanhamento adequado. Compor equipes
interdisciplinares para a detecção e tratamento das alterações visuais e
oculares.

Para a atuação como Técnico em Optometria, são fundamentais:

Conhecimentos das políticas públicas de saúde e compreensão de sua
atuação profissional frente às diretrizes, princípios e estrutura
organizacional do Sistema Único de Saúde (SUS). Conhecimentos e saberes
relacionados aos princípios das técnicas aplicadas na área, sempre
pautados numa postura humana e ética. Capacidade de raciocínio lógico,
coordenação motora fina, capacidade de concentração. Resolução de
situações-problema, boa comunicação, trabalho em equipe e
interdisciplinar, tecnologias da informação e da comunicação, gestão de
conflitos e ética profissional. Organização e responsabilidade.
Iniciativa social. Determinação e criatividade, buscando promover a
humanização da assistência. Atualização e aperfeiçoamento profissional
por meio da Educação Continuada. Conhecimentos e saberes relacionados
aos princípios das técnicas aplicadas na área, voltadas à surfaçagem de
lentes, montagem de óculos, adaptação de lentes de contato

\textbf{CBOs correspondidos e seus perfis (presentes no CNCT, e presentes na predição da IA)}

\textbf{CBO: 322305  ( Técnico em óptica e optometria )}

Planeja a fabricação das lentes de visão simples, especiais, bifocais e
multifocais - por meio dos processos de surfaçagem -, com base nas
especificações de ordem de serviço. Seleciona insumos, blocos,
equipamentos e ferramentas. Efetua cálculos de curvatura e espessura de
lentes. Executa a confecção de lentes, utilizando equipamentos e
softwares específicos. Dá acabamento às lentes, fazendo polimento e
adicionando tratamentos às lentes, como endurecimento, antirreflexo e
filtros. Faz a conferência das lentes, utilizando instrumentos
específicos, tais como lensômetro e espessímetro. Realiza processo de
montagem de lentes em armação com aro, armação de semiaro e sem aro,
conferindo a dioptria da lente, utilizando o lensômetro e marcando o
centro óptico e o eixo, de acordo com as especificações da ordem de
serviço. Executa a montagem dos óculos. Confere os óculos, de acordo com
as especificações da ordem de serviço. Pode confeccionar óculos de
segurança. Avalia a viabilidade (ou não) do uso de lentes de contato,
classificando as ametropias - conforme a prescrição do especialista - e
analisando as informações de anamnese, dos testes de filme lacrimal e
das tomadas de medidas. Adapta lentes de contato, definindo o material e
o desenho das lentes de contato; fazendo cálculos dos parâmetros;
selecionando as lentes de contato de teste; realizando assepsia das
lentes; testando o uso das lentes de contato; avaliando os resultados
das lentes de contato de teste; e executando os ajustes necessários.
Fornece instruções de uso de lentes de contato aos usuários. Aplica
prótese ocular, analisando a cavidade óptica para moldar e determinar as
características -- tamanho, cor, entre outras - da prótese. Confecciona
e ajusta a prótese ocular. Avalia e, quando necessário, readapta a
prótese. Realiza exames optométricos, posicionando e orientando o
cliente para realização de testes, de acordo com os procedimentos de
padrões técnicos. Mede acuidade visual, analisa estruturas externas e
internas dos olhos, mede pressão intraocular (tonometria), identifica
deficiências e anomalias relacionadas às alterações da função visual e
mede refração ocular (refratometria e retinoscopia). Faz diagnóstico
primário de enfermidades que acometem a saúde visual, encaminhando casos
patológicos a médicos especializados. Define compensações e auxílios
ópticos. Recepciona e atende cliente para o consumo de produtos e
serviços ópticos. Interpreta prescrição óptica do especialista,
identificando a ametropia do cliente. Sugere a armação mais adequada,
considerando estilo, atividade laboral e tipo físico do usuário. Indica
a lente mais adequada, de acordo com a armação escolhida e a prescrição
óptica do especialista. Ajusta óculos no rosto de cliente. Conserta
auxílios ópticos. Mantém registros de cliente. Gerencia estabelecimento
de artigos de óptica, administrando compras e vendas; controlando
estoque de mercadorias e materiais; providenciando manutenção de
equipamentos e instalações; administrando finanças; e controlando
qualidade de produtos e serviços. Administra pessoal, participando de
seleção, atribuindo tarefas e acompanhando registros de faltas e
atrasos. Supervisiona trabalho da equipe, avaliando desempenho e
promovendo programas de treinamento. Elabora pareceres e laudos técnicos
de produtos ópticos, conferindo óculos e lentes de contato; indicando
causas para não conformidade; e apresentando possíveis soluções. Redige
documentos de responsabilidade técnica. Promove ações de educação em
saúde visual, planejando ações de acordo com as necessidades do
público-alvo, ministrando palestras e cursos para a promoção de saúde e
prevenção de doenças visuais. Orienta o descarte de resíduos -- em
especial os gerados durante o processo de fabricação de lentes
oftálmicas -, de acordo com normas ambientais. Trabalha com segurança,
utilizando Equipamento de Proteção Individual (EPI), analisando mapa de
risco e prevenindo acidentes.

\begin{Form}
\textbf{Esta correspondência está adequada?} \\ 
\ChoiceMenu[radio,radiosymbol=\ding{52},bordercolor=0 0.5 0,name=cbovalid_ 322305 ]{}{CBO deve ficar=sim, \hspace{6em} CBO não fica aqui=nao}
\end{Form}

\newpage
\section{ 01023 -- Técnico em Órteses e Próteses }

\textbf{Perfil do Curso (CNCT)}

O Técnico em Órteses e Próteses será habilitado para:

Realizar, sob a supervisão de profissionais de nível superior, exames e
avaliação física para fins de medidas para órteses, próteses e meios
auxiliares de locomoção humanas. Realizar medidas para confeccionar,
modelar, ajustar e consertar órteses e próteses humanas, de acordo com
as necessidades físicas, psicológicas, econômicas e sociais do cliente,
como também com as possibilidades oferecidas por inovações tecnológicas.
Participar do projeto, da confecção, do ajuste e da avaliação de
órteses, próteses e meios auxiliares de locomoção. Aplicar tecnologias
para a melhoria da qualidade de vida do paciente. Acompanhar os
resultados do trabalho executado nos pacientes, atendendo a eventuais
necessidades de ajustes, de orientações ou adaptação, por solicitação
médica e fisioterapêutica e/ou de outros profissionais de áreas afins.
Avaliar e utilizar materiais e componentes relativos à produção de
órteses, de próteses e meios auxiliares de locomoção. Organizar o
processo de trabalho da oficina de produção, manutenção e adaptação de
órteses, de próteses e meios auxiliares de locomoção com base no perfil
epidemiológico do segmento da população com deficiência e nas demandas
locorregionais dos serviços da rede de atenção à saúde do SUS. Integrar
e assistir à equipe multidisciplinar nos processos de reabilitação e
readaptação, na avaliação e no acompanhamento das pessoas com
deficiência, mediante projetos terapêuticos singulares na perspectiva
sistêmica e integral.

Para a atuação como Técnico em Órteses e Próteses, são fundamentais:

Conhecimento das políticas públicas de saúde e compreensão de sua
atuação profissional frente às diretrizes, princípios e estrutura
organizacional do Sistema Único de Saúde (SUS). Conhecimentos e saberes
relacionados aos princípios políticos, sociais e éticos do SUS e bases
instrumentais para a atenção à saúde da pessoa com deficiência.
Conhecimentos e saberes relacionados ao Processo de trabalho em saúde e
especificidade do trabalho na área de atenção à saúde da pessoa com
deficiência. Conhecimento sobre os fundamentos e bases para a avaliação
e o cuidado à pessoa com deficiência física. Organização, gerência,
controle de qualidade e biossegurança em oficina de produção, manutenção
e adaptação de próteses e órteses e meios auxiliares de locomoção.
Resolução de situações-problema, comunicação, trabalho em equipe e
interdisciplinar, tecnologias da informação e da comunicação, gestão de
conflitos e ética profissional. Atualização e aperfeiçoamento
profissional por meio da educação continuada.

\textbf{CBOs correspondidos e seus perfis (presentes no CNCT, e presentes na predição da IA)}

\textbf{CBO: 322505  ( Técnico de ortopedia )}

Planeja as atividades, participando da elaboração de projeto de
confecção de prótese - peça ou aparelho de substituição dos membros ou
órgãos do corpo - e órtese, peça ou aparelho de correção ou
complementação de membros ou órgãos do corpo. Organiza o cronograma de
trabalho. Avalia diferentes recursos materiais, para escolher e elaborar
as especificações para aquisição de materiais e componentes. Seleciona
equipamentos, ferramentas e instrumentos. Realiza, sob supervisão de
profissional de nível superior, exames e avaliação física a fim de tomar
medidas para órteses e próteses. Efetua a anamnese. Pode encaminhar
solicitação e pedidos de exames complementares. Determina o design da
órtese e prótese, considerando as necessidades do paciente, as
prescrições médicas e as inovações em materiais e tecnologias
assistivas. Apresenta, ao paciente, órteses ou próteses adequadas ao seu
caso, considerando assim a sua escolha na definição do tipo a ser
confeccionado. Confecciona órtese ou prótese, interpretando as medidas
tomadas e obtendo molde negativo e molde positivo. Confecciona a
interface entre o corpo do paciente e a órtese ou prótese. Monta,
realiza prova e ajusta órtese e prótese, conforme protocolos e
procedimentos padronizados. Faz acabamento e entrega órtese e prótese
aos pacientes. Pode usar prototipagem e impressão 3D na confecção de
próteses e órteses. Realiza acompanhamento dos pacientes, fornecendo
assistência técnica durante a fase de adaptação. Atende a eventuais
necessidades de ajustes ou adaptação de órtese ou prótese, por
solicitação médica ou fisioterapêutica. Revisa periodicamente, conserta
e reforma órtese e prótese. Pode apresentar novas tecnologias em órtese
e prótese ao paciente, sugerindo substituição, diante das melhorias
introduzidas nas novas opções. Pode, sob a supervisão de profissionais
de nível superior, realizar tomada de medidas e executar confecção,
ajuste e avaliação de meios auxiliares de locomoção, tais como bengalas,
muletas canadenses, muletas axilares, andadores e cadeiras de rodas.
Orienta os pacientes sobre uso, conservação e higiene de órtese, prótese
e meios auxiliares de locomoção. Gerencia ateliê de órteses, próteses e
meios auxiliares de locomoção, participando na elaboração do projeto
físico (leiaute) do local; coordenando atividades administrativas e
financeiras; controlando estoque; e providenciando a manutenção de
máquinas, ferramentas e instalações. Subsidia tecnicamente a elaboração
de orçamentos. Zela pela qualidade dos produtos e do atendimento
prestado. Promove a segurança e a salubridade no ateliê. Orienta a
aplicação de novas tecnologias para a melhoria da qualidade de vida do
paciente. Administra pessoal, participando de seleção, atribuindo
tarefas e acompanhando registros de faltas e atrasos. Supervisiona
trabalho da equipe, avaliando desempenho e promovendo programas de
treinamento. Elabora pareceres e relatórios técnicos de qualidade e
segurança. Pode participar em projetos de equipe multidisciplinar para
promover processos de reabilitação e de readaptação, contribuindo com os
serviços de atenção à saúde e com os processos de inclusão social e
laboral da pessoa com deficiência. Arquiva e armazena, em sistemas
analógicos ou digitais, diagnósticos e pareceres sobre saúde de
pacientes. Compila informações e documentação técnica e mantém registros
de produção. Orienta o descarte de resíduos -- em especial os gerados
durante o processo de confecção de órteses e próteses -, de acordo com
normas ambientais. Trabalha com segurança, utilizando Equipamentos de
Proteção Individual (EPI), analisando mapa de risco e prevenindo
acidentes.

\begin{Form}
\textbf{Esta correspondência está adequada?} \\ 
\ChoiceMenu[radio,radiosymbol=\ding{52},bordercolor=0 0.5 0,name=cbovalid_ 322505 ]{}{CBO deve ficar=sim, \hspace{6em} CBO não fica aqui=nao}
\end{Form}

\newpage
\section{ 01025 -- Técnico em Podologia }

\textbf{Perfil do Curso (CNCT)}

O Técnico em Podologia será habilitado para:

Realizar ações de promoção da saúde e de prevenção das podopatias.
Avaliar condições da pele e anexos dos pés. Identificar lesões
elementares externas dos pés e realizar procedimentos podológicos em
diferentes tipos de afecções, utilizando técnicas de desbastamento e de
correção das unhas; aplicação de medicamentos tópicos prescritos em
receitas médicas e aplicação de curativos, em conformidade com as normas
e legislações vigentes. Selecionar e executar procedimentos de
higienização, proteção, tratamento e manutenção. Utilizar técnicas
manuais e equipamentos para podologia. Utilizar técnicas de acordo com
as características anatômicas, fisiológicas e fisiopatológicas dos pés.
Selecionar e aplicar procedimentos de profilaxia do ambiente e dos
instrumentais.

Para a atuação como Técnico em Podologia, são fundamentais:

Conhecimentos das políticas públicas de saúde e compreensão de sua
atuação profissional frente às diretrizes, princípios e estrutura
organizacional do Sistema Único de Saúde (SUS). Conhecimentos e saberes
relacionados aos princípios das técnicas aplicadas na área, sempre
pautados numa postura humana e ética Organização, responsabilidade,
iniciativa social. Atualização e aperfeiçoamento profissional por meio
da educação continuada.

\textbf{CBOs correspondidos e seus perfis (presentes no CNCT, e presentes na predição da IA)}

\textbf{CBO: 322110  ( Podólogo )}

Planeja as atividades, organizando o espaço de podologia conforme a
oferta de serviços e selecionando técnicas e recursos de trabalho para
realização dos procedimentos. Agenda atendimentos. Recebe o cliente,
apresentando os serviços e esclarecendo as dúvidas sobre suas
atividades. Efetua ou atualiza cadastro, coletando informações gerais e
de preferências do cliente. Realiza avaliação das condições de saúde do
cliente, identificando alterações patológicas - com base em informações
prestadas sobre as doenças preexistentes e hábitos de vida - e
levantando a ocorrência de patologias que interferem na ação em
podologia. Registra as informações em instrumento de avaliação, mapeando
as condições de saúde do cliente. Realiza avaliação podológica do
cliente, executando exame físico e inspeção visual, identificando as
podopatias e as lesões elementares externas dos pés. Realiza teste de
sensibilidade vibratório e pressórico e avalia os resultados, a partir
de técnicas específicas. Avalia as impressões plantares e marcha por
meio da observação do cliente. Pode encaminhar o cliente a médico,
quando forem observados sintomas indicativos de distúrbios sistêmicos -
como artrite ou diabetes - ou casos que inspirem maiores cuidados e
exames laboratoriais complementares. Elabora o plano de atendimento, a
partir da avaliação podológica. Insere tal plano em sistema
informatizado de registros de procedimentos. Realiza procedimentos da
profilaxia podal, de acordo com as normas de biossegurança, executando
higienização e antissepsia dos pés, a emoliência da pele e a limpeza
periungueal. Pode realizar a hidratação e o lixamento dos pés, além do
corte e o lixamento das unhas, conforme plano de atendimento. Desenvolve
tratamento de podopatias e deformidades podais por meio do uso de
instrumental perfurocortante, medicamentos de uso tópico, órteses e
contenção articular. Aplica procedimentos podológicos de correção de
unhas e efetua procedimentos podológicos nas onicopatias inflamatórias
ou infecciosas; nas hiperqueratoses dos pés; e nas dermatoses
infecciosas dos pés. Promove procedimentos de relaxamento e de
bem-estar, a partir da indicação podológica e necessidade do cliente.
Realiza ações de prevenção de podopatias, orientando o cliente sobre o
cuidado com os pés durante as atividades cotidianas. Conclui o
atendimento, definindo preço dos serviços com base no mercado e nos
custos de cada procedimento de podologia. Providencia o recebimento do
pagamento. Controla o estoque de materiais e produtos, armazenando-os
conforme suas funções, indicações do fabricante e legislação vigente.
Conserva o ambiente de trabalho limpo, higienizado e organizado,
providenciando a limpeza e a assepsia de mobiliários, utensílios,
equipamentos e instalações, de acordo com as normas de biossegurança.
Autoclava o instrumental e paramenta-se segundo os protocolos de
biossegurança. Providencia serviço de manutenção de equipamentos, quando
necessário. Descarta resíduos e materiais, seguindo os princípios
ambientais e as normas de biossegurança. Acondiciona materiais
perfurocortantes para descarte e lixo contaminado para incineração.
Trabalha com segurança, adotando postura ergonômica correta e utilizando
equipamentos de proteção individual, tais como luvas, máscaras e toucas.

\begin{Form}
\textbf{Esta correspondência está adequada?} \\ 
\ChoiceMenu[radio,radiosymbol=\ding{52},bordercolor=0 0.5 0,name=cbovalid_ 322110 ]{}{CBO deve ficar=sim, \hspace{6em} CBO não fica aqui=nao}
\end{Form}

\newpage
\section{ 01026 -- Técnico em Prótese Dentária }

\textbf{Perfil do Curso (CNCT)}

O Técnico em Prótese Dentária será habilitado para:

Confeccionar e reparar próteses dentárias, aparelhos ortodônticos,
aparelhos ortopédicos e dispositivos protéticos bucais. Gerenciar
laboratórios de prótese dentária. Prestar suporte técnico ao
cirurgião-dentista na fase laboratorial do processo de reabilitação
oral.

Para a atuação como Técnico em Prótese Dentária, são fundamentais:

Conhecimento das políticas públicas de saúde e compreensão de sua
atuação profissional frente às diretrizes, princípios e estrutura
organizacional do Sistema Único de Saúde (SUS). Conhecimentos e saberes
relacionados aos princípios das técnicas aplicadas na área, sempre
pautados numa postura humana e na ética do cuidado. Organização,
responsabilidade, iniciativa social, determinação e criatividade,
buscando promover a humanização da assistência. Resolução de
situações-problemas, gestão de conflitos, trabalho em equipe e
interdisciplinar de forma colaborativa, comunicação e ética
profissional. Atualização e aperfeiçoamento profissional por meio da
educação continuada.

\textbf{CBOs correspondidos e seus perfis (presentes no CNCT, e presentes na predição da IA)}

\textbf{CBO: 322410  ( Protético dentário )}

Planeja o trabalho, sugerindo tipos de prótese ao cliente, preparando
equipamentos e instrumental para uso e interpretando informações
técnicas. Confecciona próteses dentárias, preparando e avaliando o
modelo de trabalho e o preparo dos dentes, delineando a prótese,
confeccionando moldeiras e base de prova e confeccionando estruturas
metálicas, estruturas cerâmicas e estruturas de resina. Confecciona
barra metálica para prótese overdenture. Confecciona estrutura metálica
para prótese parcial removível. Confecciona provisório indireto para
prótese fixa, guia radiográfico e cirúrgico em resina acrílica. Realiza
planos de orientação, montagem balanceada, ceroplastia gengival,
acrilização e finalização de próteses. Restaura, faz acabamento,
polimento e limpeza de próteses dentárias, aparelhos e dispositivos,
seguindo protocolos e procedimentos padronizados. Pode confeccionar
próteses a partir de escaneamento intraoral e impressão 3D. Pode operar
equipamentos e instrumental automatizados e/ou informatizados para
confecção e restauro de próteses dentárias e aparelhos ortodônticos
removíveis. Armazena, manuseia e descarta produtos e resíduos materiais,
aplicando normas de biossegurança. Realiza assepsia, higienização,
desinfecção e esterilização de instrumentos e equipamentos. Empreende
procedimentos de controle e registro de insumos e equipamentos. Executa
acondicionamento e descarte responsável de materiais perfurocortantes e
biológicos. Atua no registro de procedimentos e na transcrição de
resultado em sistemas informatizados ou analógicos. Pode supervisionar e
orientar auxiliares de prótese dentária. Pode gerenciar laboratório de
prótese dentária, estruturando processos, fluxos e protocolos de
trabalho, elaborando a estratégia de comercialização dos produtos e
serviços com base na segmentação do público-alvo e administrando o
funcionamento.

\begin{Form}
\textbf{Esta correspondência está adequada?} \\ 
\ChoiceMenu[radio,radiosymbol=\ding{52},bordercolor=0 0.5 0,name=cbovalid_ 322410 ]{}{CBO deve ficar=sim, \hspace{6em} CBO não fica aqui=nao}
\end{Form}

\newpage
\section{ 01027 -- Técnico em Radiologia }

\textbf{Perfil do Curso (CNCT)}

O Técnico em Radiologia será habilitado para:

Aplicar, sob a supervisão de profissionais de nível superior, técnicas
de proteção radiológica e de biossegurança. Realizar exames de
radiodiagnóstico, considerando todo o processo de execução das técnicas
para aquisição de imagens radiológicas, que compreende: Acolher e
recepcionar o paciente. Proceder a revisão da anamnese. Orientar e
preparar o paciente para o exame. Posicionar o paciente e o equipamento.
Realizar a exposição. Processar e avaliar o padrão técnico da imagem.
Supervisionar as aplicações de técnicas em radiologia, em seus
respectivos setores. Utilizar radiação e outras formas de energia na
realização de procedimentos para obtenção de imagens diagnósticas, tais
como: radiologia convencional e digital, mamografia, densitometria,
hemodinâmica, tomografia computadorizada, ressonância magnética,
radiologia forense, radiologia veterinária, dentre outras.

Para a atuação como Técnico em Radiologia, são fundamentais:

Conhecimentos e saberes relacionados à anatomia e fisiologia humana.
Resolução de situações-problemas, trabalho em equipe e interdisciplinar,
tecnologias da informação e da comunicação, gestão de conflitos e ética
profissional. Organização e responsabilidade; iniciativa social;
entusiasmo; empatia e respeito. Atualização e aperfeiçoamento
profissional por meio da Educação Continuada.

\textbf{CBOs correspondidos e seus perfis (presentes no CNCT, e presentes na predição da IA)}

\textbf{CBO: 324115  ( Técnico em radiologia e imagenologia )}

Recebe pedido de exame e prontuário. Verifica as condições e as
necessidades físicas e psicológicas do paciente. Explica os
procedimentos do exame ao paciente, retirando suas próteses e o
posicionando em equipamento, conforme instrução e receituário médico.
Organiza a área de trabalho, equipamentos e acessórios de proteção
radiológica. Monta carrinho de medicamentos de emergência. Opera ou
auxilia na operação de equipamentos de radiologia e imagenologia. Zela
pela qualidade do atendimento prestado, pela qualidade técnica dos
exames, pela segurança do paciente, do técnico e do ambiente de
trabalho. Empreende procedimentos de controle e registro de insumos e
equipamentos. Realiza acondicionamento de materiais perfurocortantes
para descarte responsável. Atua no registro de procedimentos, de
identificação do paciente e transcrição de resultados do exame em
sistemas informatizados ou analógicos.

\begin{Form}
\textbf{Esta correspondência está adequada?} \\ 
\ChoiceMenu[radio,radiosymbol=\ding{52},bordercolor=0 0.5 0,name=cbovalid_ 324115 ]{}{CBO deve ficar=sim, \hspace{6em} CBO não fica aqui=nao}
\end{Form}

\newpage
\section{ 01029 -- Técnico em Registros e Informações em Saúde }

\textbf{Perfil do Curso (CNCT)}

O Técnico em Registros e Informações em Saúde será habilitado para:

Apoiar a administração e a coordenação de serviços de documentação,
serviços de registros e estatísticas de saúde (em diferentes formatos:
papel, eletrônico ou digital). Dar suporte ao atendimento do paciente.
Organizar os registros administrativos e epidemiológicos em sistemas de
informação ou em sistemas de Prontuários Eletrônico do Paciente (PEP).
Apoiar a elaboração de planejamento, de controle e de avaliação dos
serviços de saúde. Guardar, catalogar e manter registros. Produzir e
analisar indicadores demográficos, gerenciais e epidemiológicos.
Implementar e operar sistemas de informações em saúde e Sistemas de
Prontuário Eletrônico nas respectivas unidades. Identificar e extrair
dados relevantes, descartar dados, organizar e gerar dados que permitam
a elaboração de compreensões e de significados capazes de orientar a
tomada de decisão de técnicos e gestores com o máximo de eficácia.
Identificar e selecionar as fontes de informação. Identificar problemas
e acionar suportes de informática. Criar, testar e validar instrumentos
de coleta de dados. Fazer avaliação da coleta de dados: analisar,
depurar, recuperar, introduzir, preservar e validar dados. Identificar
instrumentos de tratamento de dados. Apoiar a supervisão do desempenho
dos sistemas. Disponibilizar dados, fazer apresentações e emitir
relatórios. Identificar pontos críticos e inconsistências nos dados.
Retroalimentar fonte de dados. Promover a articulação entre as áreas de
Tecnologias de Informação e Comunicação (TIC) e de saúde e a
transformação digital no contexto e processos do estabelecimento de
saúde que atua. Conhecer os princípios e regras estabelecidas para o
tratamento de dados pessoais de saúde. Atuar de forma ética para
garantir privacidade e confidencialidade dos dados pessoais de saúde.
Codificar os dados coletados utilizando sistemas de classificação que
dizem respeito a sinais e sintomas, diagnósticos, procedimentos,
medicamentos e outros, permitindo a comparabilidade entre diferentes
provedores, serviços e sistemas de saúde. Apoiar estratégias de educação
na área de registros e informações em saúde, junto aos profissionais de
saúde. Alimentar, sempre que necessário, os sistemas de informações em
saúde, de acordo com sua área de atuação, garantindo a fidedignidade,
disponibilidade e segurança dos dados;

Para a atuação como Técnico em Registros e Informações em Saúde, são
fundamentais:

Conhecimento das políticas públicas de saúde e compreensão de sua
atuação profissional frente às diretrizes, princípios e estrutura
organizacional do Sistema Único de Saúde (SUS). Conhecimentos e saberes
relacionados aos princípios das técnicas aplicadas na área, sempre
pautados numa postura humana e na ética do cuidado. Zelo pela qualidade
dos dados e das informações de saúde. Conhecimento e respeito aos
aspectos de privacidade e confidencialidade dos dados de saúde.
Interesse pela modernização e inovação. Reconhecimento dos benefícios e
uso crítico das tecnologias de informação e comunicação em saúde.
Organização, responsabilidade, flexibilidade e persistência.
Disponibilidade e iniciativa, buscando promover a humanização da
assistência e o foco no cidadão/paciente. Resolução de
situações-problema, gestão de conflitos, trabalho em equipe de forma
colaborativa, comunicação e ética profissional. Atualização e
aperfeiçoamento profissional por meio da educação continuada.
Compreensão do papel estratégico das informações, enquanto geradoras de
conhecimento para o processo de gestão do cuidado e do sistema de saúde.
Conhecimento sobre a Política Nacional de Informação, sobre Informática
em Saúde no Brasil, Sistemas Nacionais de Informação em Saúde e sobre
legislação correlata. Compreensão do seu processo de trabalho para que
possa atuar de forma crítica e integrada aos demais trabalhadores da
equipe responsável pelos cuidados de saúde. Conhecimento das normas de
saúde e segurança do trabalho, contribuindo para as ações relacionadas à
saúde do trabalhador de Registros e Informações em Saúde, de modo a
prevenir, diminuir ou eliminar os riscos à sua saúde.

\textbf{CBOs correspondidos e seus perfis (presentes no CNCT, e presentes na predição da IA)}

\textbf{CBO: 415305  ( Registrador de câncer )}

Planeja a coleta e o tratamento de dados para atualização do
RCBP-Registro de Câncer de Base Populacional de uma localidade definida.
Recebe lista de hospitais, clínicas e laboratórios, considerados fontes
notificadoras dos dados. Realiza coleta passiva de dados, recebendo-os
de hospitais, clínicas e laboratórios pela internet ou por outro meio de
comunicação. Faz a revisão, a padronização e a correção dos dados, para
introduzi-los em um banco temporário de registros de câncer. Realiza
coleta ativa de dados, planejando visitas às fontes notificadoras.
Elabora cronograma e agenda visitas. Visita hospitais, clínicas e
laboratórios, selecionando prontuários de paciente para inclusão no
banco de dados. Preenche dados identificadores em ficha de registro,
tais como dados pessoais, nome do estabelecimento de saúde, médicos
responsáveis pelo atendimento e fontes pagadoras do tratamento.
Acrescenta histórico do tratamento na ficha de registro, consultando no
prontuário procedência do paciente; instituição de origem que
diagnosticou o tumor; tratamento prévio em outra instituição; dados de
estadiamento do tumor (fase de desenvolvimento, extensão e gravidade da
doença). Checa a completude dos dados. Pode dirimir dúvidas junto à
coordenação do levantamento quanto aos dados contidos no prontuário.
Consulta lista de pacientes já registrados no banco de dados, separando
prontuários para preenchimento da ficha de seguimento. Coleta dados para
acompanhamento (``follow-up'') da doença, registrando a data da última
consulta do paciente na instituição e a ocorrência de novos tumores
primários. Localiza e contata pacientes com perda de seguimento,
registrando a atualização de seus dados. Organiza base de dados,
identificando e excluindo casos em duplicidade (mesma pessoa e mesmo
diagnóstico); corrigindo inconsistências identificadas nas fichas pelo
sistema; e agrupando casos em que o paciente tem vários tipos de tumor
primário. Complementa base de dados com data de óbito, obtida em fontes
oficiais de informações de mortalidade. Realiza codificação clínica dos
dados - da coleta passiva e da coleta ativa -, com o auxílio de
classificações internacionais de doenças. Consulta os códigos de
CID-Classificação Internacional de Doenças e Problemas Relacionados à
Saúde para dados topográficos (localização anatômica do tumor),
morfológicos (tipo de célula que forma o tumor) e para registrar o
código da causa da morte. Identifica os dados de estadiamento, a partir
da TNM--Classificação de Tumores Malignos (Classification of Malignant
Tumours) ou da Classificação da Figo-Federação Internacional de
Ginecologia e Obstetrícia. Consolida relatórios a serem submetidos à
validação da coordenação. Auxilia na preparação das informações para
publicações ou divulgação de dados. Pode alterar cadastro de fontes
notificadoras, informando à coordenação novas informações levantadas
durante a coleta ativa de dados. Organiza e atualiza arquivos -- físicos
e eletrônicos -- de documentos e informações para consulta, tais como
cadastro de fontes notificadoras e classificações internacionais de
doenças. Conserva o local de trabalho limpo e organizado. Zela pela
preservação, funcionamento e integridade dos recursos de trabalho.
Destina adequadamente os resíduos gerados, fazendo reciclagem e
realizando descarte de acordo com as normas ambientais. Trabalha com
segurança, utilizando equipamentos de proteção individual sempre que
necessário, especialmente durante os levantamentos de dados em
hospitais, clínicas e laboratórios.

\begin{Form}
\textbf{Esta correspondência está adequada?} \\ 
\ChoiceMenu[radio,radiosymbol=\ding{52},bordercolor=0 0.5 0,name=cbovalid_ 415305 ]{}{CBO deve ficar=sim, \hspace{6em} CBO não fica aqui=nao}
\end{Form}

\newpage
\section{ 01030 -- Técnico em Saúde Bucal }

\textbf{Perfil do Curso (CNCT)}

O Técnico em Saúde Bucal atua sob a supervisão do cirurgião-dentista e
será habilitado para:

Auxiliar na promoção, prevenção e controle de doenças bucais. Auxiliar
atividades clínicas voltadas ao restabelecimento da saúde, estética e
função mastigatória do indivíduo Participar de programas educativos
voltados à saúde bucal. Contribuir na realização de estudos
epidemiológicos em saúde bucal. Instrumentar o cirurgião-dentista.
Realizar tomadas radiográficas e fotográficas de uso odontológico.
Realizar escaneamento intraoral. Controlar estoques. Supervisionar a
manutenção dos equipamentos. Organizar o ambiente de trabalho
odontológico Exercer suas competências em âmbito hospitalar

Para a atuação como Técnico em Saúde Bucal, são fundamentais:

Conhecimento sobre as políticas públicas de saúde e compreensão de sua
atuação profissional frente às diretrizes, princípios e estrutura
organizacional do Sistema Único de Saúde (SUS). Conhecimentos e saberes
relacionados aos princípios das técnicas aplicadas na área, sempre
pautados numa postura humana e na ética do cuidado. Organização e
responsabilidade; iniciativa social; determinação e criatividade,
buscando promover a humanização da assistência. Resolução de
situações-problema, gestão de conflitos, trabalho em equipe e
interdisciplinar de forma colaborativa, comunicação e ética
profissional.

\textbf{CBOs correspondidos e seus perfis (presentes no CNCT, e presentes na predição da IA)}

\textbf{CBO: 322405  ( Técnico em saúde bucal )}

Planeja o trabalho, agendando consultas e realizando anamnese prévia do
paciente. Prepara equipamentos e instrumental, conforme programa de
atendimento da clínica. Recebe, prepara e acompanha pacientes na
realização de exames de diagnóstico ou de tratamento odontológico, sob
supervisão direta do cirurgião dentista. Auxilia na realização do
trabalho clínico odontológico, sob orientação do cirurgião-dentista.
Insere e distribui, no preparo cavitário, materiais odontológicos para
restaurações dentárias diretas. Seleciona, prova e prepara moldeiras.
Realiza profilaxia, aplicação tópica de flúor, remoção de biofilme e de
sutura. Realiza tomadas radiográficas intrabucais, processa filmes
radiográficos e realiza tomadas fotográficas de uso odontológico. Em
cirurgia, realiza o isolamento parcial ou total do campo operatório,
empreendendo a limpeza e antissepsia pré e pós operatória. Participa das
ações educativas de promoção da saúde e na prevenção das doenças bucais.
Participa da realização de levantamentos e estudos epidemiológicos. Pode
orientar e supervisionar equipe de auxiliares em saúde bucal. Higieniza,
realiza assepsia, desinfecção e esterilização de instrumentos e
equipamentos odontológicos e do ambiente de trabalho. Atua no registro
de procedimentos, transcrição de resultado em sistemas informatizados ou
analógicos. Empreende procedimentos de controle e registro de insumos e
equipamentos. Realiza acondicionamento e descarte responsável de
materiais perfurocortantes e biológicos. Armazena, manuseia e descarta
produtos e resíduos odontológicos, aplicando normas complementares de
biossegurança em todas as etapas laborais.

\begin{Form}
\textbf{Esta correspondência está adequada?} \\ 
\ChoiceMenu[radio,radiosymbol=\ding{52},bordercolor=0 0.5 0,name=cbovalid_ 322405 ]{}{CBO deve ficar=sim, \hspace{6em} CBO não fica aqui=nao}
\end{Form}

\textbf{CBO: 322425  ( Técnico em saúde bucal da estratégia de saúde da família )}

Participa do processo de territorialização e mapeamento da área de
atuação da equipe, identificando grupos, famílias e indivíduos expostos
a riscos e vulnerabilidades. Colabora no cadastramento e no levantamento
de dados de saúde das famílias e dos indivíduos. Participa do
planejamento das atividades e das definições das prioridades locais.
Organiza grupos para promoção da saúde bucal, de acordo com as
informações cadastrais obtidas. Atende usuários nas unidades básicas de
saúde, nos domicílios ou espaços comunitários e nas unidades
odontológicas móveis. Realiza o acolhimento do paciente nos serviços de
saúde bucal. Faz remoção do biofilme, de acordo com a indicação técnica
definida pelo cirurgião-dentista. Insere e distribui, no preparo
cavitário, materiais odontológicos na restauração dentária direta,
usando materiais e instrumentos indicados pelo cirurgião-dentista.
Realiza fotografias e tomadas de uso odontológico. Processa filme
radiográfico. Auxilia e instrumenta o cirurgião-dentista nas
intervenções clínicas. Seleciona moldeiras e prepara modelos em gesso.
Realiza a remoção de sutura, conforme indicação do cirurgião dentista.
Apoia as atividades dos auxiliares de saúde bucal e dos agentes
comunitários de saúde nas ações de prevenção e promoção da saúde bucal.
Participa do treinamento e capacitação de auxiliar em saúde bucal e de
agentes multiplicadores das ações de promoção à saúde. Pode
supervisionar equipe de auxiliares de saúde bucal. Participa das ações
educativas, atuando na promoção da saúde e na prevenção das doenças
bucais. Participa da realização de levantamentos e estudos
epidemiológicos. Pode ministrar palestras em campanhas epidemiológicas e
de saúde bucal. Executa a organização, limpeza, assepsia, desinfecção e
esterilização do instrumental, dos equipamentos odontológicos e do
ambiente de trabalho. Procede à limpeza e à antissepsia do campo
operatório, antes e após atos cirúrgicos. Armazena, manuseia e descarta
produtos e resíduos odontológicos, aplicando medidas de biossegurança.
Coordena a manutenção e a conservação dos equipamentos odontológicos.
Empreende procedimentos de controle e registro de insumos e
equipamentos. Registra as ações realizadas em sistema de informação.

\begin{Form}
\textbf{Esta correspondência está adequada?} \\ 
\ChoiceMenu[radio,radiosymbol=\ding{52},bordercolor=0 0.5 0,name=cbovalid_ 322425 ]{}{CBO deve ficar=sim, \hspace{6em} CBO não fica aqui=nao}
\end{Form}

\textbf{CBO: 322430  ( Auxiliar em saúde bucal da estratégia de saúde da família )}

Participa do processo de territorialização e mapeamento da área de
atuação da equipe, identificando grupos, famílias e indivíduos expostos
a riscos e vulnerabilidades. Colabora no cadastramento e no levantamento
de dados de saúde das famílias e dos indivíduos. Realiza o acolhimento
do paciente nos serviços de saúde bucal. Auxilia e instrumenta o
cirurgião-dentista nas intervenções clínicas. Organiza e executa
atividades de higiene bucal, processa filme radiográfico, auxilia na
seleção de moldeiras e na preparação de modelos em gesso. Pode
participar de ações educativas e de campanhas para promoção da saúde
bucal. Pode participar da realização de levantamentos e estudos
epidemiológicos. Executa organização, limpeza, assepsia, desinfecção e
esterilização do instrumental, dos equipamentos odontológicos e do
ambiente de trabalho. Auxilia na manutenção e na conservação dos
equipamentos odontológicos. Armazena, transporta, manuseia e descarta
produtos e resíduos odontológicos, aplicando medidas de biossegurança.

\begin{Form}
\textbf{Esta correspondência está adequada?} \\ 
\ChoiceMenu[radio,radiosymbol=\ding{52},bordercolor=0 0.5 0,name=cbovalid_ 322430 ]{}{CBO deve ficar=sim, \hspace{6em} CBO não fica aqui=nao}
\end{Form}

\newpage
\section{ 01031 -- Técnico em Terapias Holísticas }

\textbf{Perfil do Curso (CNCT)}

O Técnico em Terapias Holísticas será habilitado para:

Executar práticas corporais de descontração e relaxamento, advindas de
tradições da medicina aiurvédica indiana, como o Yoga e da Medicina
Tradicional Chinesa (MTC), como o Lian Gong e o Tai Chi Chuan. Aplicar
técnicas de massagens relaxantes e estimulantes, do tipo terapêutica e
estética. Conduzir práticas meditativas/contemplativas e danças
circulares. Utilizar técnicas oriundas da aromaterapia, cromoterapia,
musicoterapia, auriculoterapia, moxabustão e ventosaterapia, entre
outras. Participar de equipes multiprofissionais em serviços de saúde,
clínicas estéticas.

Para a atuação como Técnico em Terapias Holísticas, são fundamentais:

Conhecimentos das políticas públicas de saúde e compreensão de sua
atuação profissional frente às diretrizes, princípios e estrutura
organizacional do Sistema Único de Saúde (SUS). Conhecimentos e saberes
relacionados aos princípios das técnicas aplicadas na área, sempre
pautados numa postura humana e ética. Resolução de situações-problema,
comunicação, trabalho em equipe e interdisciplinar, tecnologias da
informação e da comunicação, gestão de conflitos e ética profissional.
Organização e responsabilidade. Iniciativa social. Determinação e
criatividade, buscando promover a humanização da assistência.
Atualização e aperfeiçoamento profissional por meio da Educação
Continuada.

\textbf{CBOs correspondidos e seus perfis (presentes no CNCT, e presentes na predição da IA)}

\textbf{CBO: 322125  ( Terapeuta holístico )}

Planeja as atividades, organizando o espaço conforme a oferta de
serviços e selecionando os recursos de trabalho. Agenda atendimentos.
Recebe o cliente, apresentando os serviços e esclarecendo as dúvidas
sobre suas atividades. Efetua ou atualiza cadastro, coletando
informações gerais e de preferências do cliente. Avalia as condições de
saúde e verifica os hábitos de vida do cliente, questionando-o sobre
suas queixas, expectativas e eventuais disfunções e patologias. Analisa
aspectos biomecânicos; verifica sistema muscular - força, temperatura e
tônus -; observa estado emocional; e outros sinais e sintomas. Define o
plano de tratamento para o cliente, estabelecendo os procedimentos
terapêuticos e/ou estéticos não invasivos. Pode consultar outros
profissionais de saúde - como fisioterapeutas, médicos e psicólogos -,
para tratamento multidisciplinar. Preenche em sistema informatizado a
ficha de do cliente, detalhando a avaliação das condições de saúde e
especificando o plano de tratamento. Executa práticas corporais de
descontração e relaxamento, advindas de tradições da medicina ayurvédica
indiana -- tal como o Yoga - e da medicina tradicional chinesa, como o
Lian Gong e o Tai Chi Chuan. Aplica técnicas de massagens relaxantes e
estimulantes, dos tipos terapêutica e estética. Conduz práticas
meditativas e contemplativas, além de danças circulares. Utiliza
técnicas não invasivas, tais como as oriundas da aromaterapia -
selecionando os óleos essenciais adequados a cada necessidade -; da
cromoterapia - fazendo a seleção de cores adequadas -; da moxaterapia -
localizando áreas, meridianos e pontos energéticos -; da ventosaterapia
- localizando meridianos, pontos energéticos e grupos musculares -;
entre outras. Considera as indicações, os benefícios e as possíveis
contraindicações de cada técnica. Pode recomendar ao cliente a
realização de exercícios e o uso de essências florais e de produtos
fitoterápicos. Controla o estoque de produtos e materiais,
providenciando a aquisição para reposição. Mantém a organização, a
limpeza, a assepsia e a higienização dos ambientes. Conserva
equipamentos, instrumentos e acessórios de trabalho limpos,
higienizados, acondicionados e em plenas condições de uso e
funcionamento. Requisita manutenção de equipamentos, quando necessário.
Descarta material e produtos com validade vencida, seguindo os
princípios ambientais e as normas de biossegurança. Trabalha com
segurança, usando equipamentos de proteção individual e adotando postura
ergonômica correta durante o trabalho.

\begin{Form}
\textbf{Esta correspondência está adequada?} \\ 
\ChoiceMenu[radio,radiosymbol=\ding{52},bordercolor=0 0.5 0,name=cbovalid_ 322125 ]{}{CBO deve ficar=sim, \hspace{6em} CBO não fica aqui=nao}
\end{Form}

\newpage
\section{ 01032 -- Técnico em Veterinária }

\textbf{Perfil do Curso (CNCT)}

O Técnico em Veterinária atua sob a supervisão do Médico-Veterinário e
será habilitado para:

Exercer atividade de apoio, de assistência e de acompanhamento do
trabalho do médico-veterinário. Preencher o cadastro do animal e
conferir seus dados. Registrar procedimentos especiais, tais como dieta
especial, jejum pré-cirúrgico, e outros previamente estabelecidos pelo
médico-veterinário. Pesar o animal, verificar sinais vitais dos animais,
observar e relatar as condições físicas, atitudes e comportamentos.
Orientar sobre cuidados gerais de higiene, conforme a prescrição e
orientação do médico-veterinário. Ministrar medicamentos, vacinas e
vermífugos e fazer curativos prescritos pelo Médico-Veterinário
responsável. Auxiliar nos primeiros socorros e nas manobras de parto e
cuidados neonatais. Preparar animais para procedimentos cirúrgicos,
realizando tricotomia, higiene do paciente e antissepsia da pele.
Auxiliar nos procedimentos de acesso intravenoso e no procedimento de
intubação do animal. Preparar animais e materiais para procedimentos
médico-veterinários. Realizar a contenção física do animal. Auxiliar na
coleta de material para exames clínicos. Identificar e embalar cadáver,
após constatação do óbito do animal pelo Médico-Veterinário. Realizar
cuidados gerais de limpeza, manutenção e esterilização de materiais e
equipamentos. Controlar estoques, solicitar material e repor
medicamentos e materiais.

Para a atuação como Técnico em Veterinária, são fundamentais:

Conhecimentos das políticas públicas de saúde e compreensão de sua
atuação profissional frente às diretrizes, princípios e estrutura
organizacional do Sistema Único de Saúde (SUS). Aptidão para atuar na
área de saúde animal, conhecimentos e habilidades técnicas relacionados
a procedimentos de apoio, de assistência e de acompanhamento do trabalho
do médico-veterinário. Conhecimento sobre as atividades de competência
privativa do médico-veterinário, atuação de forma colaborativa nos
trabalhos em equipe, visão sistêmica das atividades e processos,
capacidade de comunicação, autonomia, proatividade, respeito as
diversidades nos grupos de trabalho, resiliência frente aos problemas,
organização, responsabilidade, ética e consciente sobre a sua atuação
profissional e o bem-estar dos animais. Atualização e aperfeiçoamento
profissional por meio da Educação Continuada.

\textbf{CBOs correspondidos e seus perfis (presentes no CNCT, e presentes na predição da IA)}

\textbf{CBO: 519305  ( Auxiliar de veterinário )}

Atende cliente, entrevistando-o sobre as condições de saúde e o
comportamento de seu animal e sobre o motivo da consulta, para
preenchimento de ficha preliminar. Informa o cliente sobre normas e
regulamentos do estabelecimento. Transporta ou conduz o animal para
atendimento. Prepara o consultório, os instrumentos e os materiais
necessários para o atendimento do animal pelo médico veterinário. Pesa,
mede, controla sinais vitais - como temperatura e pressão arterial -- e
imobiliza o animal, preparando-o para a consulta. Auxilia o médico
veterinário durante a consulta. Presta apoio na coleta de material --
como sangue e urina -- e encaminha as amostras, identificadas pelo
médico veterinário, para exames laboratoriais. Prepara e organiza o
ambiente, seleciona a caixa cirúrgica e prepara o material para
procedimentos cirúrgicos. Realiza tricotomia, raspando os pelos da área
do animal que será submetida à cirurgia. Posiciona o animal na mesa de
cirurgia. Auxilia o médico veterinário durante procedimentos cirúrgicos,
passando-lhe os instrumentos e os materiais solicitados. Controla sinais
vitais do animal, tais como temperatura e pressão. Pode auxiliar no
procedimento de intubação. Presta cuidados pós-operatórios de rotina e
monitora os animais que se recuperam da cirurgia, informando ao médico
veterinário sobre quaisquer alterações ou sintomas incomuns. Pode usar
sistemas informatizados e eletrônicos de monitoração do comportamento e
da movimentação animal, tais como: câmeras, ``babás-eletrônicas'' e
sensores de controle de atividades do animal. Executa manobras de
contenção física do animal, segundo métodos ética e tecnicamente
adequados para sua espécie, seu porte e sua condição física. Alimenta os
animais, levando em consideração os casos especiais -- tais como dieta e
jejum pré-cirúrgico -- e os horários, estabelecidos pelo médico
veterinário. Administra medicamentos por via oral ou tópica, vacinas,
fármacos de desparasitação e vermifugação e outros prescritos por médico
veterinário, fazendo constar assinatura, data e hora no prontuário. Pode
auxiliar nos primeiros socorros, sob a orientação do médico veterinário.
Realiza procedimentos de assepsia e antissepsia. Faz curativos, quando
prescritos pelo médico veterinário. Pode banhar os animais, verificando
a existência de parasitas - como carrapatos e pulgas - e as condições
físicas do animal. Limpa dentes, ouvido e a região dos olhos. Realiza
atividades físicas com o animal, sob a orientação do médico veterinário.
Relata ao médico veterinário as condições físicas, as atitudes e o
comportamento do animal. Preenche e organiza fichas de atendimento e
documentação relativa a consultas, internações, estadas, exames,
cirurgias, medicamentos utilizados e serviços prestados, recolhendo as
devidas assinaturas do médico veterinário. Preenche o cadastro do
animal, em sistema informatizado. Auxilia no controle de estoques e na
reposição de medicamentos e de materiais hospitalares e clínicos.
Higieniza e desinfeta ambientes e equipamentos. Lava, higieniza,
desinfeta e esteriliza os instrumentos. Higieniza o local de estada dos
animais. Colabora na organização do ambiente de trabalho. Faz a
destinação correta de lixo hospitalar e clínico, separando e embalando
resíduos físicos, biológicos e químicos, bem como cadáveres dos animais
- após constatação do óbito pelo médico veterinário -, conforme normas
de biossegurança e de proteção pessoal.

\begin{Form}
\textbf{Esta correspondência está adequada?} \\ 
\ChoiceMenu[radio,radiosymbol=\ding{52},bordercolor=0 0.5 0,name=cbovalid_ 519305 ]{}{CBO deve ficar=sim, \hspace{6em} CBO não fica aqui=nao}
\end{Form}

\newpage
\section{ 01033 -- Técnico em Vigilância em Saúde }

\textbf{Perfil do Curso (CNCT)}

O Técnico em Vigilância em Saúde será habilitado para:

Atuar, sob a supervisão de profissionais de nível superior, de forma a
aplicar a normatização de produtos, processos, ambientes e serviços de
interesses a saúde. Controlar o fluxo de pessoas, animais, plantas e
produtos em portos, aeroportos e fronteiras. Desenvolver ações de
controle e monitoramento de doenças, endemias e de vetores. Desenvolver
ações de inspeção e fiscalização sanitárias. Elaborar e implementar,
junto com a população do território, ações educativas no âmbito das
vigilâncias: ambiental, sanitária, epidemiológica e saúde do trabalhador
para promoção da saúde. Investigar, monitorar e avaliar riscos e os
determinantes dos agravos e danos à saúde da população, relacionados ao
trabalho, assim como ao meio ambiente. Planejar e articular ações
intersetoriais para promoção da saúde prevenção, controle e
monitoramento de vetores e doenças endêmicas. Planejar, executar e
avaliar o processo de vigilância sanitária, epidemiológica, ambiental,
saúde do trabalhador e vigilância da situação de saúde. Realizar análise
territorial das condições de vida e de saúde da população, como também
identificar e intervir em situações de risco e de vulnerabilidade de
grupos populacionais e ambientes. Desenvolver ações de mobilização
social junto à comunidade para promoção e proteção da saúde. Realizar
ações de prevenção e controle de doenças e agravos à saúde em interação
com o Agente Comunitário de Saúde e a equipe de Atenção Básica.

Para a atuação como Técnico em Vigilância em Saúde, são fundamentais:

Conhecimento das políticas públicas de saúde: organização, princípios e
diretrizes do Sistema Único de Saúde (SUS). Conhecimentos e saberes
relacionados a processos de sustentabilidade, territorialização,
organização do SUS e vigilâncias. Resolução de situações-problemas,
trabalho em equipe, tecnologias da informação e da comunicação, gestão
de conflitos e ética profissional. Iniciativa e capacidade de
planejamento Organização, responsabilidade, iniciativa social,
entusiasmo, empatia e respeito. Atualização e aperfeiçoamento
profissional por meio da educação continuada.

\textbf{CBOs correspondidos e seus perfis (presentes no CNCT, e presentes na predição da IA)}

\textbf{CBO: 352210  ( Agente de saúde pública )}

Zela pela saúde pública, executando atividades para levantar e punir
atividades e projetos que infringem normas sanitárias. Fiscaliza locais
e atividades, solicitando documentação ao fiscalizado, levantando
informações técnicas referentes a questões sanitárias, coletando
produtos irregulares e enviando material para análise aos órgãos
competentes. Participa de operações especiais (blitz). Vistoria locais e
atividades - tais como indústrias, restaurantes, bares, escolas,
creches, asilos e hospitais -, verificando cumprimento das exigências
legais e técnicas, examinando condições sanitárias das instalações e dos
equipamentos e avaliando as condições de trabalho. Analisa tecnicamente
projetos, examinando o processo de licenciamento, verificando exigências
técnicas e elaborando parecer, laudo ou relatório técnico. Investiga
denúncias de problemas sanitários, levantando informações junto à
comunidade, coletando material para análise e apreendendo produtos
irregulares. Se constatada a veracidade da denúncia, enquadra legalmente
o caso. Reprime atividades ilegais, apreendendo equipamentos e
materiais, interditando estabelecimentos e atividades e inutilizando
produtos irregulares. Adverte, notifica, intima ou multa infrator,
podendo encaminhá-lo às autoridades competentes. Pode solicitar apoio à
polícia. Ministra palestras, promove educação sanitária e dá orientações
técnicas sobre saúde e serviços de profilaxia. Pode gerenciar
atividades, coordenando e supervisionando equipes, planejando operações
e administrando recursos humanos, financeiros e materiais. Emite
notificações, autorizações, intimações, licenças e ofícios. Registra
denúncias, preenche relatórios administrativos e formaliza proposta de
embargo, interdição e multa.

\begin{Form}
\textbf{Esta correspondência está adequada?} \\ 
\ChoiceMenu[radio,radiosymbol=\ding{52},bordercolor=0 0.5 0,name=cbovalid_ 352210 ]{}{CBO deve ficar=sim, \hspace{6em} CBO não fica aqui=nao}
\end{Form}

\newpage
\section{ 02001 -- Técnico em Automação Industrial }

\textbf{Perfil do Curso (CNCT)}

O Técnico em Automação Industrial será habilitado para:

Desenvolver e integrar soluções para sistemas de automação visando à
medição e ao controle de variáveis em processos industriais,
considerando as normas, os padrões e os requisitos técnicos de
qualidade, saúde e segurança e de meio ambiente. Empregar programas de
computação e redes industriais no controle de processos industriais.
Planejar, controlar e executar a instalação e a manutenção de
equipamentos automatizados e/ou sistemas robotizados para controle de
processos industriais. Realizar medições, testes e calibrações em
equipamentos eletroeletrônicos empregados em controle de processos
industriais. Instalar, configurar e operar tecnologias de manufatura
aditiva, sistemas ciberfísicos e processos de produção com internet das
coisas. Reconhecer tecnologias inovadoras presentes no segmento visando
a atender às transformações digitais na sociedade. Realizar
especificação, projeto, instalação, medição, teste, diagnóstico e
calibração de equipamentos e sistemas automatizados. Executar
procedimentos de controle de qualidade, operação e gestão de sistemas
automatizados e controle de processos.

Para atuação como Técnico em Automação Industrial, são fundamentais:

Conhecimentos e saberes relacionados aos processos de planejamento e
implementação de processos automatizados de modo a assegurar a saúde e a
segurança dos trabalhadores e dos usuários. Conhecimentos e saberes
relacionados à sustentabilidade do processo produtivo, às técnicas e aos
processos de produção, às normas técnicas, à liderança de equipes, à
solução de problemas técnicos e trabalhistas e à gestão de conflitos.

\textbf{CBOs correspondidos e seus perfis (presentes no CNCT, e presentes na predição da IA)}

\textbf{CBO: 300105  ( Técnico em mecatrônica - automação da manufatura )}

Projeta sistemas de automação da manufatura, analisando o processo de
manufatura e o produto manufaturado envolvidos no projeto. Avalia as
condições do local de trabalho para instalação de máquinas e
equipamentos. Define fluxo do processo, identifica alternativas
tecnológicas e esboça proposta para automação do processo. Faz
especificação de materiais e componentes empregados em projeto.
Considera aspectos de ergonomia, de segurança do trabalho, de economia
de energia e de preservação do meio ambiente na elaboração de projeto.
Elabora relatório de custo-benefício para análise técnico-financeira do
projeto. Elabora cronograma de implantação de projeto. Pode realizar
projetos de integração de sistemas de automação da manufatura,
estabelecendo interfaces entre equipamentos eletroeletrônicos e
mecânicos. Participa de projetos de implantação e utilização de
tecnologias típicas das fábricas digitais nos processos de manufatura.
Instala sistemas de automação da manufatura, interpretando documentação
do projeto e organizando materiais e recursos para o processo de
instalação. Monta componentes eletroeletrônicos e componentes mecânicos
em sistemas de automação da manufatura. Testa a operação do sistema, sem
processamento de matéria-prima. Acompanha teste de produção do sistema,
em processo. Faz correções e ajustes no sistema, conforme resultados de
testes. Pode realizar a integração de diversificados tipos de robôs nos
sistemas de automação da manufatura. Pode aperfeiçoar sistemas de
automação da manufatura já instalados, substituindo equipamentos,
atualizando programas e implementando novas tecnologias. Programa
equipamentos e sistemas de automação da manufatura, escrevendo programas
em linguagens diversificadas de programação ou utilizando recursos
avançados, tais como sistemas supervisórios, CAD-Desenho Assistido por
Computador (Computer Aided Design), CAM--Manufatura Auxiliada por
Computador (Computer Aided Manufacturing), dentre outros. Utiliza
sistemas analíticos e ambientes de desenvolvimento para criação de
aplicações. Desenvolve a programação de controladores lógicos
programáveis (CLP) e de máquinas-ferramentas com comando numérico
computadorizado (CNC). Programa parâmetros para acionamentos de potência
e de redes industriais, utilizando ferramentas de programação
específicas, tendo em vista a integração de equipamentos. Realiza a
manutenção de sistemas de automação da manufatura, avaliando gráficos de
tendências e relatórios de manutenção e planejando o processo e as
etapas de trabalho de manutenção. Planeja manutenção preventiva e
preditiva. Analisa falhas em sistemas. Executa procedimentos de
manutenção. Avalia a eficácia das soluções de manutenção implementadas.
Colabora no processo de aquisição de componentes, equipamentos e
sistemas, verificando características técnicas dos bens; selecionando
fornecedores; e acompanhando testes de funcionamento de máquinas e
equipamentos para elaboração e emissão de parecer técnico. Avalia
disponibilidade de peças de reposição. Coordena equipes de trabalho,
identificando competências técnicas e pessoais, atribuindo
responsabilidades e estabelecendo metas individuais. Forma equipe
multidisciplinar para análise de máquinas e equipamentos. Promove
integração entre setores da empresa e os integrantes de equipes
envolvidas em projetos. Reúne-se com a equipe de trabalho. Monitora a
execução de tarefas. Dá suporte técnico aos integrantes da equipe.
Treina usuários em operação e manutenção de sistemas de automação da
manufatura. Elabora documentação técnica de projetos de sistemas de
automação da manufatura. Documenta alterações de projeto ocorridas
durante a instalação de sistemas. Elabora relatório de aceitação de
equipamentos. Documenta o plano de ação de manutenção preventiva e
preditiva, as ações corretivas e as melhorias implementadas nos
sistemas. Mantém o local de trabalho limpo, seguro e organizado.
Conserva ferramentas, instrumentos e equipamentos. Utiliza equipamentos
de proteção.

\begin{Form}
\textbf{Esta correspondência está adequada?} \\ 
\ChoiceMenu[radio,radiosymbol=\ding{52},bordercolor=0 0.5 0,name=cbovalid_ 300105 ]{}{CBO deve ficar=sim, \hspace{6em} CBO não fica aqui=nao}
\end{Form}

\newpage
\section{ 02002 -- Técnico em Eletroeletrônica }

\textbf{Perfil do Curso (CNCT)}

O Técnico em Eletroeletrônica será habilitado para:

Planejar, controlar e executar a instalação e a manutenção de
equipamentos e instalações eletroeletrônicas industriais, considerando
as normas, os padrões e os requisitos técnicos de qualidade, saúde e
segurança e de meio ambiente. Projetar e instalar sistemas de
acionamentos, controles eletroeletrônicos e sistemas automáticos em
instalações industriais. Aplicar medidas para o uso eficiente da energia
elétrica e de fontes de energias alternativas. Realizar medições, testes
e calibrações de equipamentos eletroeletrônicos e inspecionar
componentes, produtos, serviços e atividades de profissionais da área de
eletroeletrônica. Reconhecer tecnologias inovadoras presentes no
segmento visando ao atendimento das transformações digitais
implementadas na sociedade.

Para atuação como Técnico em Eletroeletrônica, são fundamentais:

Conhecimentos e saberes relacionados aos processos de planejamento e
implementação de sistemas eletroeletrônicos de modo a assegurar a saúde
e a segurança dos trabalhadores e dos usuários. Conhecimentos e saberes
relacionados à sustentabilidade do processo produtivo, às técnicas e aos
processos de produção, às normas técnicas, à liderança de equipes, à
solução de problemas técnicos e trabalhistas e à gestão de conflitos.

\textbf{CBOs correspondidos e seus perfis (presentes no CNCT, e presentes na predição da IA)}

\textbf{CBO: 300305  ( Técnico em eletromecânica )}

Participa da elaboração de projetos eletromecânicos de máquinas,
equipamentos e instalações, auxiliando na especificação de componentes
eletromecânicos, avaliando características do local de implantação e
consultando normas de ergonomia e segurança do trabalho, conforme as
características do projeto. Elabora desenhos técnicos relacionados ao
projeto, bem como esquemas para orientar as instalações de todos os
elementos do projeto, de acordo com normas técnicas. Define listas de
materiais necessários. Colabora na elaboração do orçamento. Elabora
relatórios de custo-benefício, para análise técnico-financeira do
projeto. Planeja a execução de projetos eletromecânicos de máquinas,
equipamentos e instalações, interpretando o projeto para elaborar o
plano de trabalho da fabricação de componentes mecânicos; definir etapas
de produção a serem executadas na própria empresa e por terceiros;
relacionar pessoas e equipamentos necessários; e elaborar o cronograma
de atividades. Identifica necessidade de dispositivos e ferramentas para
melhoria dos recursos produtivos. Usina peças, interpretando desenhos
técnicos, preparando local de trabalho e máquinas e verificando as
dimensões da matéria-prima. Opera máquinas para usinagem de peças
mecânicas, inspecionando as dimensões e verificando o acabamento, para
assegurar a conformidades das peças usinadas ao projeto. Realiza
soldagem para unir peças mecânicas. Monta máquinas, equipamentos e
instalações, interpretando esquemas de montagem, preparando o local e
selecionando componentes eletromecânicos, ferramentas e instrumentos
para a montagem. Executa os trabalhos de montagem dos componentes,
observando sequência, encaixe, posicionamento e relação entre partes,
para compor sistemas de máquinas, equipamentos e instalações
eletromecânicas. Avalia as etapas da montagem para garantir conformidade
ao projeto. Faz entrega técnica de máquinas, equipamentos e instalações,
colaborando no planejamento, conferindo as condições físicas do local de
instalação do projeto e elaborando relatório da entrega. Testa o
funcionamento de máquinas, equipamentos e de instalações, no local de
entrega. Orienta usuários quanto à utilização, conservação e normas de
segurança de máquinas e equipamentos. Pode realizar o comissionamento de
máquinas, equipamentos e instalações eletromecânicas. Inspeciona
máquinas, equipamentos e instalações eletromecânicas, para avaliar
condições operacionais e gerais, com vistas à definição de parâmetros
para planos de manutenção. Realiza manutenção eletromecânica de
máquinas, equipamentos e instalações, participando na elaboração de
planos de manutenção, estabelecendo condições de segurança para execução
dos trabalhos, interpretando instruções e selecionando ferramentas e
instrumentos. Interage com sistemas digitais próprios das máquinas,
equipamentos e instalações, e com redes industriais, que auxiliam nos
serviços de manutenção eletromecânica preditiva, preventiva e corretiva.
Executa consertos em máquinas, equipamentos e instalações
eletromecânicas, identificando causas de defeitos, listando peças
danificadas, especificando componentes para reposição e definindo a
alternativa a ser seguida para sanar o problema. Elabora o orçamento do
serviço de manutenção. Substitui componentes danificados, testando
funcionamento após o conserto e atualizando registros de manutenção.
Executa serviços de modernização - ``retrofitting'' - de máquinas,
equipamentos e instalações eletromecânicas. Pode coordenar equipes de
trabalho, identificando competências técnicas e pessoais da equipe,
definindo metas e delegando responsabilidades sobre as etapas de
trabalho. Presta orientações técnicas e relacionadas aos aspectos de
segurança. Verifica o cumprimento de prazos definidos em cronogramas.
Mantém o local das atividades limpo e organizado. Faz o descarte de
resíduos seguindo normas ambientais.

\begin{Form}
\textbf{Esta correspondência está adequada?} \\ 
\ChoiceMenu[radio,radiosymbol=\ding{52},bordercolor=0 0.5 0,name=cbovalid_ 300305 ]{}{CBO deve ficar=sim, \hspace{6em} CBO não fica aqui=nao}
\end{Form}

\textbf{CBO: 731135  ( Montador de equipamentos elétricos )}

Seleciona e utiliza materiais, máquinas, equipamentos, instrumentos e
ferramentas, considerando o produto a ser montado. Monta equipamentos
elétricos, incluindo a montagem das placas de circuito impresso. Executa
processos de inserção, fixação, conexão e interligação de dispositivos e
acessórios. Aplica soldagem para interligação dos componentes.
Confecciona e instala bobinas elétricas, interpretando esquemas
elétricos. Compõe o gabinete e o painel do equipamento com a disposição,
conexão e interligação das placas. Fixa e interliga dispositivos
elétricos e eletrônicos. Finaliza a montagem, realizando ajustes
eletroeletrônicos e configurações. Testa equipamentos elétricos,
realizando teste das placas de circuito impresso, testes funcionais
durante o processo de montagem, teste final de impacto e testes
elétricos do equipamento. Confere valores e polaridade dos componentes e
verifica o acoplamento e a interligação de peças e bobinas. Inspeciona
equipamentos elétricos - incluindo a inspeção das placas de circuito
impresso -, por meio de verificação visual geral e de exame da
disposição dos componentes e dispositivos. Faz checagem visual do
equipamento e correção de defeitos de montagem, para atestar o produto.
Verifica a qualidade do processo de montagem de equipamentos elétricos,
no que se refere ao pleno funcionamento, às especificações técnicas e às
propriedades físicas e mecânicas. Prepara o produto para expedição,
fechando o equipamento, fazendo limpeza externa e organizando
acessórios, manual e lista de assistência técnica na embalagem. Prepara
e organiza o ambiente de trabalho, mantendo-o limpo e em condições de
uso e abastecendo-o com peças, componentes, materiais, equipamentos,
instrumentos e ferramentas. Segue rigorosamente normas regulamentadoras
e procedimentos de saúde e segurança no trabalho, utilizando
equipamentos de proteção individual e coletiva. Mantém tais equipamentos
organizados, acondicionados e em plenas condições de uso e
funcionamento. Preenche e organiza documentação técnica relacionada aos
processos, tais como relatório de controle de processo, ficha técnica,
ficha de rastreabilidade, requisição de materiais, entre outros.

\begin{Form}
\textbf{Esta correspondência está adequada?} \\ 
\ChoiceMenu[radio,radiosymbol=\ding{52},bordercolor=0 0.5 0,name=cbovalid_ 731135 ]{}{CBO deve ficar=sim, \hspace{6em} CBO não fica aqui=nao}
\end{Form}

\textbf{CBO: 731150  ( Montador de equipamentos eletrônicos )}

Seleciona e utiliza materiais, máquinas, equipamentos, instrumentos e
ferramentas, considerando o produto a ser montado. Monta equipamentos
eletrônicos, incluindo a montagem das placas de circuito impresso.
Executa processos de inserção, fixação, conexão e interligação de
dispositivos elétricos e eletrônicos e acessórios. Aplica soldagem para
interligação dos componentes. Confecciona e instala bobinas elétricas,
interpretando esquemas elétricos. Compõe gabinete e painel de controle
de equipamentos eletrônicos, dispondo, conectando e interligando
dispositivos elétricos e eletrônicos. Finaliza a montagem realizando
ajustes eletrônicos e configurações. Testa equipamentos eletrônicos,
realizando teste das placas de circuito impresso, testes funcionais
durante o processo de montagem, teste final de impacto e testes
elétricos. Confere valores e polaridade dos componentes e verifica o
acoplamento e a interligação de peças e bobinas. Inspeciona equipamentos
eletrônicos - incluindo a inspeção das placas de circuito impresso -,
por meio de verificação visual geral e exame da disposição dos
componentes e dispositivos. Faz checagem visual do equipamento e
correção de defeitos de montagem, atestando o produto. Verifica a
qualidade do processo de montagem de equipamentos eletrônicos, no que se
refere ao pleno funcionamento, às especificações técnicas e às
propriedades físicas e mecânicas. Prepara o produto para expedição,
fechando o equipamento, fixando acessórios, fazendo limpeza externa e
organizando manual e lista de assistência técnica na embalagem. Prepara
e organiza o ambiente de trabalho, mantendo-o limpo e em condições de
uso, abastecendo-o com peças, componentes, materiais, equipamentos,
instrumentos e ferramentas. Segue normas e procedimentos de saúde e
segurança no trabalho, utilizando equipamentos de proteção individual e
coletiva. Mantém tais equipamentos organizados, acondicionados e em
plenas condições de uso e funcionamento. Preenche e organiza
documentação técnica relacionada ao processo e ao produto, tais como
relatório de controle de processo, ficha técnica, ficha de
rastreabilidade, requisição de materiais, entre outros.

\begin{Form}
\textbf{Esta correspondência está adequada?} \\ 
\ChoiceMenu[radio,radiosymbol=\ding{52},bordercolor=0 0.5 0,name=cbovalid_ 731150 ]{}{CBO deve ficar=sim, \hspace{6em} CBO não fica aqui=nao}
\end{Form}

\newpage
\section{ 02003 -- Técnico em Eletromecânica }

\textbf{Perfil do Curso (CNCT)}

O Técnico em Eletromecânica será habilitado para:

Planejar, controlar e executar a instalação, a manutenção e a entrega
técnica de máquinas e equipamentos eletromecânicos industriais,
considerando as normas, os padrões e os requisitos técnicos de
qualidade, saúde e segurança e de meio ambiente. Elaborar projetos de
produtos relacionados a máquinas e equipamentos eletromecânicos
especificando materiais para construção mecânica e elétrica por meio de
técnicas de usinagem e soldagem. Realizar inspeção visual, dimensional e
testes em sistemas, instrumentos, equipamentos eletromecânicos,
pneumáticos e hidráulicos de máquinas. Reconhecer tecnologias inovadoras
presentes no segmento visando a atender às transformações digitais na
sociedade.

Para atuação como Técnico em Eletromecânica, são fundamentais:

Conhecimentos e saberes relacionados aos processos de planejamento,
produção e manutenção de equipamentos eletromecânicos de modo a
assegurar a saúde e a segurança dos trabalhadores e dos usuários.
Conhecimentos e saberes relacionados à sustentabilidade do processo
produtivo, às técnicas e aos processos de produção, às normas técnicas,
à liderança de equipes, à solução de problemas técnicos e trabalhistas e
à gestão de conflitos.

\textbf{CBOs correspondidos e seus perfis (presentes no CNCT, e presentes na predição da IA)}

\textbf{CBO: 300305  ( Técnico em eletromecânica )}

Participa da elaboração de projetos eletromecânicos de máquinas,
equipamentos e instalações, auxiliando na especificação de componentes
eletromecânicos, avaliando características do local de implantação e
consultando normas de ergonomia e segurança do trabalho, conforme as
características do projeto. Elabora desenhos técnicos relacionados ao
projeto, bem como esquemas para orientar as instalações de todos os
elementos do projeto, de acordo com normas técnicas. Define listas de
materiais necessários. Colabora na elaboração do orçamento. Elabora
relatórios de custo-benefício, para análise técnico-financeira do
projeto. Planeja a execução de projetos eletromecânicos de máquinas,
equipamentos e instalações, interpretando o projeto para elaborar o
plano de trabalho da fabricação de componentes mecânicos; definir etapas
de produção a serem executadas na própria empresa e por terceiros;
relacionar pessoas e equipamentos necessários; e elaborar o cronograma
de atividades. Identifica necessidade de dispositivos e ferramentas para
melhoria dos recursos produtivos. Usina peças, interpretando desenhos
técnicos, preparando local de trabalho e máquinas e verificando as
dimensões da matéria-prima. Opera máquinas para usinagem de peças
mecânicas, inspecionando as dimensões e verificando o acabamento, para
assegurar a conformidades das peças usinadas ao projeto. Realiza
soldagem para unir peças mecânicas. Monta máquinas, equipamentos e
instalações, interpretando esquemas de montagem, preparando o local e
selecionando componentes eletromecânicos, ferramentas e instrumentos
para a montagem. Executa os trabalhos de montagem dos componentes,
observando sequência, encaixe, posicionamento e relação entre partes,
para compor sistemas de máquinas, equipamentos e instalações
eletromecânicas. Avalia as etapas da montagem para garantir conformidade
ao projeto. Faz entrega técnica de máquinas, equipamentos e instalações,
colaborando no planejamento, conferindo as condições físicas do local de
instalação do projeto e elaborando relatório da entrega. Testa o
funcionamento de máquinas, equipamentos e de instalações, no local de
entrega. Orienta usuários quanto à utilização, conservação e normas de
segurança de máquinas e equipamentos. Pode realizar o comissionamento de
máquinas, equipamentos e instalações eletromecânicas. Inspeciona
máquinas, equipamentos e instalações eletromecânicas, para avaliar
condições operacionais e gerais, com vistas à definição de parâmetros
para planos de manutenção. Realiza manutenção eletromecânica de
máquinas, equipamentos e instalações, participando na elaboração de
planos de manutenção, estabelecendo condições de segurança para execução
dos trabalhos, interpretando instruções e selecionando ferramentas e
instrumentos. Interage com sistemas digitais próprios das máquinas,
equipamentos e instalações, e com redes industriais, que auxiliam nos
serviços de manutenção eletromecânica preditiva, preventiva e corretiva.
Executa consertos em máquinas, equipamentos e instalações
eletromecânicas, identificando causas de defeitos, listando peças
danificadas, especificando componentes para reposição e definindo a
alternativa a ser seguida para sanar o problema. Elabora o orçamento do
serviço de manutenção. Substitui componentes danificados, testando
funcionamento após o conserto e atualizando registros de manutenção.
Executa serviços de modernização - ``retrofitting'' - de máquinas,
equipamentos e instalações eletromecânicas. Pode coordenar equipes de
trabalho, identificando competências técnicas e pessoais da equipe,
definindo metas e delegando responsabilidades sobre as etapas de
trabalho. Presta orientações técnicas e relacionadas aos aspectos de
segurança. Verifica o cumprimento de prazos definidos em cronogramas.
Mantém o local das atividades limpo e organizado. Faz o descarte de
resíduos seguindo normas ambientais.

\begin{Form}
\textbf{Esta correspondência está adequada?} \\ 
\ChoiceMenu[radio,radiosymbol=\ding{52},bordercolor=0 0.5 0,name=cbovalid_ 300305 ]{}{CBO deve ficar=sim, \hspace{6em} CBO não fica aqui=nao}
\end{Form}

\newpage
\section{ 02004 -- Técnico em Eletrônica }

\textbf{Perfil do Curso (CNCT)}

O Técnico em Eletrônica será habilitado para:

Planejar, controlar e executar projetos eletrônicos com dispositivos e
tecnologias relacionadas às áreas de eletrônica analógica, digital, de
potência e microcontrolados. Executar e supervisionar a instalação e a
manutenção de equipamentos e sistemas eletrônicos e robotizados,
inclusive de telemetria e telecomunicações, considerando as normas, os
padrões e os requisitos técnicos de qualidade, saúde e segurança e de
meio ambiente. Realizar medições, testes, calibrações e comissionamento
de equipamentos eletrônicos. Reconhecer tecnologias inovadoras presentes
no segmento visando a atender às transformações digitais na sociedade.

Para atuação como Técnico em Eletrônica, são fundamentais:

Conhecimentos e saberes relacionados aos processos de planejamento e
implementação de sistema eletrônicos de modo a assegurar a saúde e a
segurança dos trabalhadores e dos usuários. Conhecimentos e saberes
relacionados à sustentabilidade do processo produtivo, às técnicas e
processos de produção, às normas técnicas, à liderança de equipes, à
solução de problemas técnicos e trabalhistas e à gestão de conflitos.

\textbf{CBOs correspondidos e seus perfis (presentes no CNCT, e presentes na predição da IA)}

\textbf{CBO: 313205  ( Técnico de manutenção eletrônica )}

Instala equipamentos e aparelhos eletrônicos, observando as condições do
ambiente, realizando ajustes e testes. Conserta e faz manutenção
corretiva em aparelhos e equipamentos eletrônicos, interpretando
esquemas elétricos, substituindo componentes defeituosos, podendo
utilizar equipamentos de diagnóstico eletrônico. Faz manutenções
preventivas e preditivas em equipamentos eletrônicos, de acordo com
planos de manutenção, podendo fazer uso da internet para suporte. Cria e
implementa dispositivos de automação. Integra dispositivos e sistemas
eletrônicos por meio de redes digitais de comunicação de dados, locais e
a distância. Simula o processo produtiva e libera a linha para produção
em massa. Desenvolve, monta e testa dispositivos de circuitos
eletrônicos, especificando componentes, calculando custos, incluindo a
possibilidade de uso de microcontroladores, microprocessadores e outros
dispositivos eletrônicos digitais programáveis. Descreve procedimentos
de trabalho, preenche laudos técnicos e formulário de reposição de peças
rejeitadas e emite relatórios técnicos. Registra ocorrências e elabora
gráficos de resultados. Orienta, treina e avalia o desempenho
operacional dos operadores.

\begin{Form}
\textbf{Esta correspondência está adequada?} \\ 
\ChoiceMenu[radio,radiosymbol=\ding{52},bordercolor=0 0.5 0,name=cbovalid_ 313205 ]{}{CBO deve ficar=sim, \hspace{6em} CBO não fica aqui=nao}
\end{Form}

\textbf{CBO: 313210  ( Técnico de manutenção eletrônica (circuitos de máquinas com comando numérico) )}

Conserta e faz manutenção corretiva em circuitos elétricos e eletrônicos
de máquinas com comando numérico, interpretando esquemas elétricos e
substituindo componentes defeituosos. Pode utilizar equipamentos de
diagnóstico eletrônico. Instala máquinas com comando numérico,
observando as condições do ambiente e realizando ajustes e testes. Faz
manutenções preventivas e preditivas em circuitos elétricos e
eletrônicos de máquinas com comando numérico, de acordo com planos de
manutenção, podendo fazer uso da internet para suporte. Propõe
alterações no processo produtivo, implementando dispositivos de
automação e integração de dispositivos e sistemas eletrônicos, por meio
de redes digitais de comunicação de dados, locais e a distância. Pode
realizar interligação de máquinas. Desenvolve, monta e testa
dispositivos de circuitos eletrônicos para máquinas com comando
numérico, especificando componentes e calculando custos. Pode usar
microcontroladores, microprocessadores e outros dispositivos eletrônicos
digitais programáveis. Descreve procedimentos de trabalho, preenche
laudos técnicos e formulário de reposição de peças rejeitadas e emite
relatórios técnicos. Registra ocorrências e elabora gráficos de
resultados. Orienta, treina e avalia o desempenho operacional dos
operadores.

\begin{Form}
\textbf{Esta correspondência está adequada?} \\ 
\ChoiceMenu[radio,radiosymbol=\ding{52},bordercolor=0 0.5 0,name=cbovalid_ 313210 ]{}{CBO deve ficar=sim, \hspace{6em} CBO não fica aqui=nao}
\end{Form}

\textbf{CBO: 313215  ( Técnico eletrônico )}

Conserta e faz manutenção corretiva em aparelhos e equipamentos
eletrônicos, interpretando esquemas elétricos e substituindo componentes
defeituosos. Pode utilizar equipamentos de diagnóstico eletrônico.
Instala equipamentos, aparelhos e dispositivos eletrônicos, observando
as condições do ambiente, realizando ajustes e testes. Faz manutenções
preventivas e preditivas em equipamentos eletrônicos, de acordo com
planos de manutenção, podendo fazer uso da internet para suporte.
Realiza medições, testes e calibrações de equipamentos. Desenvolve,
monta e testa dispositivos de circuitos eletrônicos, especificando
componentes, calculando custos, incluindo a possibilidade de uso de
microcontroladores, microprocessadores e outros dispositivos eletrônicos
digitais programáveis. Sugere alterações no processo de produção,
implementando dispositivos de automação e integração de dispositivos e
sistemas eletrônicos por meio de redes digitais de comunicação de dados,
locais e a distância. Participa de reuniões técnicas com pessoal interno
e externo. Redige procedimentos de trabalho e elabora gráficos de
resultados. Pode realizar treinamento, orientação e avaliação de
desempenho de operadores.

\begin{Form}
\textbf{Esta correspondência está adequada?} \\ 
\ChoiceMenu[radio,radiosymbol=\ding{52},bordercolor=0 0.5 0,name=cbovalid_ 313215 ]{}{CBO deve ficar=sim, \hspace{6em} CBO não fica aqui=nao}
\end{Form}

\textbf{CBO: 313220  ( Técnico em manutenção de equipamentos de informática )}

Conserta e faz manutenção corretiva em equipamentos de informática -
como computadores e seus periféricos, dispositivos de acesso a redes e
fontes chaveadas -, interpretando esquemas elétricos e substituindo
componentes defeituosos. Pode utilizar equipamentos de diagnóstico
eletrônico e fazer uso da internet para acesso remoto. Instala
equipamentos de informática, observando as condições do ambiente,
realizando ajustes e testes. Faz manutenções preventivas e preditivas em
equipamentos de informática, de acordo com planos de manutenção,
presencialmente ou online. Sugere alterações nos processos produtivos,
para melhoria do sistema, integrando dispositivos e sistemas eletrônicos
por meio de redes digitais de comunicação de dados, locais e a
distância. Desenvolve, monta e testa dispositivos de circuitos
eletrônicos para equipamentos de informática, especificando e calculando
custos de componentes. Pode usar microcontroladores, microprocessadores
e outros dispositivos eletrônicos digitais programáveis. Descreve
procedimentos de trabalho, preenche laudos técnicos e formulário de
reposição de peças rejeitadas e emite relatórios técnicos. Elabora
gráficos de resultados. Pode treinar, orientar e avaliar o desempenho de
operadores.

\begin{Form}
\textbf{Esta correspondência está adequada?} \\ 
\ChoiceMenu[radio,radiosymbol=\ding{52},bordercolor=0 0.5 0,name=cbovalid_ 313220 ]{}{CBO deve ficar=sim, \hspace{6em} CBO não fica aqui=nao}
\end{Form}

\newpage
\section{ 02005 -- Técnico em Eletrotécnica }

\textbf{Perfil do Curso (CNCT)}

O Técnico em Eletrotécnica será habilitado para:

Planejar, controlar e executar a instalação e a manutenção de sistemas e
instalações elétricas industriais, prediais e residenciais, considerando
as normas, os padrões e os requisitos técnicos de qualidade, saúde e
segurança e de meio ambiente. Elaborar e desenvolver projetos de
instalações elétricas industriais, prediais e residenciais, sistemas de
acionamentos elétricos e de automação industrial e de infraestrutura
para sistemas de telecomunicações em edificações. Aplicar medidas para o
uso eficiente da energia elétrica e de fontes energéticas alternativas.
Elaborar e desenvolver programação e parametrização de sistemas de
acionamentos eletrônicos industriais. Planejar e executar instalação e
manutenção de sistemas de aterramento e de descargas atmosféricas em
edificações residenciais, comerciais e industriais. Reconhecer
tecnologias inovadoras presentes no segmento visando a atender às
transformações digitais na sociedade.

Para atuação como Técnico em Eletrotécnica, são fundamentais:

Conhecimentos e saberes relacionados aos processos de planejamento e
implementação de sistemas elétricos de modo a assegurar a saúde e a
segurança dos trabalhadores e dos usuários. Conhecimentos e saberes
relacionados à sustentabilidade do processo produtivo, às técnicas e aos
processos de produção, às normas técnicas, à liderança de equipes, à
solução de problemas técnicos e trabalhistas e à gestão de conflitos.

\textbf{CBOs correspondidos e seus perfis (presentes no CNCT, e presentes na predição da IA)}

\textbf{CBO: 313105  ( Eletrotécnico )}

Elabora estudos e projetos referentes a sistemas elétricos de potência,
determinando escopo, aplicando normas e avaliando viabilidade econômica.
Dimensiona circuitos eletroeletrônicos e componentes. Participa do
desenvolvimento de produtos e elabora documentações técnicas, de acordo
com requisitos do projeto. Realiza projetos elétricos de potência -
conforme especificações e cronograma -, executando montagem e
comissionamento e colocando em operação. Participa do desenvolvimento de
processos em sistemas elétricos de potência, estabelecendo
procedimentos, normas e padrões. Determina fluxograma e define máquinas
e equipamentos. Realiza medições e ensaios e propõe melhorias. Aplica
tecnologias adequadas ao projeto e ao processo, considerando a crescente
utilização de fontes alternativas de energia, os programas de eficiência
energética e as tecnologias de automação aplicadas a sistemas elétricos.
Opera sistemas elétricos de geração, transmissão e distribuição de
energia - de acordo com normas, instruções e procedimentos -,
supervisionando e elaborando os programas de manobras. Pode utilizar
sistemas de automação. Executa manutenção preventiva, corretiva e
preditiva em sistemas elétricos de potência, de acordo com o plano de
manutenção, diagnosticando o desempenho dos equipamentos. Atua na
supervisão e no treinamento de equipe de trabalho, avaliando desempenho
e propondo programas de aperfeiçoamento e atualização profissional. Pode
trabalhar na área comercial, identificando necessidades do cliente,
pesquisando novos mercados, realizando vendas e prestando suporte
técnico.

\begin{Form}
\textbf{Esta correspondência está adequada?} \\ 
\ChoiceMenu[radio,radiosymbol=\ding{52},bordercolor=0 0.5 0,name=cbovalid_ 313105 ]{}{CBO deve ficar=sim, \hspace{6em} CBO não fica aqui=nao}
\end{Form}

\textbf{CBO: 313110  ( Eletrotécnico (produção de energia) )}

Elabora estudos e projetos em sistemas de produção de energia,
determinando escopo, aplicando normas e avaliando viabilidade econômica.
Dimensiona circuitos eletroeletrônicos e componentes. Participa do
desenvolvimento de produtos e elabora documentações técnicas, de acordo
com os requisitos do projeto. Realiza projetos em sistemas de produção
de energia - conforme especificações e cronograma -, executando montagem
e comissionamento e colocando em operação. Participa do desenvolvimento
de processos em sistemas de produção de energia, estabelecendo
procedimentos, normas e padrões. Determina fluxograma e define máquinas
e equipamentos. Realiza medições e ensaios e propõe melhorias. Aplica
tecnologias adequadas ao projeto e ao processo, considerando a crescente
utilização de fontes alternativas de energia, os programas de eficiência
energética e as tecnologias de automação aplicadas a sistemas elétricos.
Opera sistemas elétricos de geração, transmissão e distribuição de
energia - de acordo com normas, instruções e procedimentos -,
supervisionando e elaborando os programas de manobras. Pode utilizar
sistemas de automação. Executa manutenção preventiva, corretiva e
preditiva em sistemas elétricos, de acordo com o plano de manutenção,
diagnosticando o desempenho dos equipamentos. Atua na supervisão e no
treinamento de equipe de trabalho, avaliando desempenho e propondo
programas de aperfeiçoamento e atualização profissional.

\begin{Form}
\textbf{Esta correspondência está adequada?} \\ 
\ChoiceMenu[radio,radiosymbol=\ding{52},bordercolor=0 0.5 0,name=cbovalid_ 313110 ]{}{CBO deve ficar=sim, \hspace{6em} CBO não fica aqui=nao}
\end{Form}

\textbf{CBO: 313115  ( Eletrotécnico na fabricação, montagem e instalação de máquinas e equipamentos )}

Realiza projetos de máquinas e equipamentos elétricos - conforme
especificações e cronograma -, executando montagem, instalação e
comissionamento e colocando em operação. Participa do desenvolvimento de
processos de fabricação, montagem e instalação de máquinas e
equipamentos elétricos, estabelecendo procedimentos, normas e padrões.
Determina fluxograma e define máquinas e equipamentos. Realiza medições
e ensaios e propõe melhorias. Aplica tecnologias adequadas ao projeto e
ao processo, considerando a crescente utilização de fontes alternativas
de energia, os programas de eficiência energética e as tecnologias de
automação aplicadas a sistemas elétricos. Executa manutenção preventiva,
corretiva e preditiva em máquinas e equipamentos elétricos, de acordo
com o plano de manutenção. Atua na supervisão e no treinamento de equipe
de trabalho, avaliando desempenho e propondo programas de
aperfeiçoamento e atualização profissional.

\begin{Form}
\textbf{Esta correspondência está adequada?} \\ 
\ChoiceMenu[radio,radiosymbol=\ding{52},bordercolor=0 0.5 0,name=cbovalid_ 313115 ]{}{CBO deve ficar=sim, \hspace{6em} CBO não fica aqui=nao}
\end{Form}

\textbf{CBO: 313120  ( Técnico de manutenção elétrica )}

Executa manutenção elétrica preventiva e corretiva em sistemas elétricos
de potência, de acordo com o plano de manutenção. Realiza projetos em
sistemas elétricos de potência - conforme especificações e cronograma -,
executando montagem e comissionamento e colocando em operação. Participa
do desenvolvimento de processos em sistemas elétricos de potência,
estabelecendo procedimentos, normas e padrões. Determina fluxograma e
define máquinas e equipamentos. Realiza medições e ensaios e propõe
melhorias. Aplica tecnologias adequadas ao projeto e ao processo,
considerando a crescente utilização de fontes alternativas de energia,
os programas de eficiência energética e as tecnologias de automação
aplicadas a sistemas elétricos. Atua na supervisão e no treinamento de
equipe de trabalho, avaliando desempenho e propondo programas de
aperfeiçoamento e atualização profissional.

\begin{Form}
\textbf{Esta correspondência está adequada?} \\ 
\ChoiceMenu[radio,radiosymbol=\ding{52},bordercolor=0 0.5 0,name=cbovalid_ 313120 ]{}{CBO deve ficar=sim, \hspace{6em} CBO não fica aqui=nao}
\end{Form}

\textbf{CBO: 313125  ( Técnico de manutenção elétrica de máquina )}

Executa manutenção elétrica preventiva e corretiva em máquinas, de
acordo com o plano de manutenção. Realiza projetos em máquinas -
conforme especificações e cronograma -, executando montagem e
comissionamento e colocando em operação. Participa do desenvolvimento de
processos em sistemas elétricos de potência, estabelecendo
procedimentos, normas e padrões. Determina fluxograma e define máquinas
e equipamentos. Realiza medições e ensaios e propõe melhorias. Aplica
tecnologias adequadas ao projeto e ao processo, considerando a crescente
utilização de fontes alternativas de energia, os programas de eficiência
energética e as tecnologias de automação aplicadas a sistemas elétricos.
Atua na supervisão e no treinamento de equipe de trabalho, avaliando
desempenho e propondo programas de aperfeiçoamento e atualização
profissional.

\begin{Form}
\textbf{Esta correspondência está adequada?} \\ 
\ChoiceMenu[radio,radiosymbol=\ding{52},bordercolor=0 0.5 0,name=cbovalid_ 313125 ]{}{CBO deve ficar=sim, \hspace{6em} CBO não fica aqui=nao}
\end{Form}

\textbf{CBO: 313130  ( Técnico eletricista )}

Elabora estudos e projetos em sistemas elétricos de potência,
determinando escopo, aplicando normas e avaliando viabilidade econômica.
Dimensiona circuitos eletroeletrônicos e componentes. Participa do
desenvolvimento de produtos e elabora documentações técnicas, de acordo
com os requisitos do projeto. Realiza projetos em sistemas elétricos de
potência - conforme especificações e cronograma -, executando montagem e
comissionamento e colocando em operação. Participa do desenvolvimento de
processos em sistemas elétricos de potência, estabelecendo
procedimentos, normas e padrões. Determina fluxograma e define máquinas
e equipamentos. Realiza medições e ensaios e propõe melhorias. Aplica
tecnologias adequadas ao projeto e ao processo, considerando a crescente
utilização de fontes alternativas de energia, os programas de eficiência
energética e as tecnologias de automação aplicadas a sistemas elétricos.
Opera sistemas elétricos de potência - de acordo com normas, instruções
e procedimentos -, realizando a manobra dos equipamentos. Executa
manutenção preventiva e corretiva em sistemas elétricos de potência, de
acordo com o plano de manutenção. Atua na supervisão e no treinamento de
equipe de trabalho, avaliando desempenho e propondo programas de
aperfeiçoamento e atualização profissional. Pode trabalhar na área
comercial, pesquisando mercados, desenvolvendo clientes e elaborando
orçamentos.

\begin{Form}
\textbf{Esta correspondência está adequada?} \\ 
\ChoiceMenu[radio,radiosymbol=\ding{52},bordercolor=0 0.5 0,name=cbovalid_ 313130 ]{}{CBO deve ficar=sim, \hspace{6em} CBO não fica aqui=nao}
\end{Form}

\newpage
\section{ 02006 -- Técnico em Fabricação Mecânica }

\textbf{Perfil do Curso (CNCT)}

O Técnico em Fabricação Mecânica será habilitado para:

Desenvolver projetos, planejar, supervisionar e controlar atividades de
fundição, em usinagem convencional e computadorizada, em caldeiraria, em
soldagem e processos de conformação mecânica. Interpretar desenho
técnico. Selecionar, desenvolver e especi?car ferramental para os
processos produtivos. Executar ensaios mecânicos. Especi?car materiais e
insumos aplicados aos processos de fabricação mecânica. Controlar
estoques de produtos acabados.

Para atuação como Técnico em Fabricação Mecânica, são fundamentais:

Conhecimentos e saberes relacionados aos processos de planejamento e
operação das atribuições da área de modo a assegurar a saúde e a
segurança dos trabalhadores e dos futuros usuários e operadores de
empresas em processos de transformação em fabricação mecânica.
Conhecimentos e saberes relacionados à sustentabilidade do processo
produtivo, às normas e relatórios técnicos, à legislação da área, às
novas tecnologias relacionadas à indústria 4.0, à liderança de equipes,
à solução de problemas técnicos e à gestão de conflitos.

\textbf{CBOs correspondidos e seus perfis (presentes no CNCT, e presentes na predição da IA)}

\textbf{CBO: 314205  ( Técnico mecânico na fabricação de ferramentas )}

Pesquisa e estima as necessidades do mercado, avaliando a demanda
potencial e a viabilidade de execução do produto, tendo em vista a
seleção das características de ferramentas e dispositivos a serem
desenvolvidos. Identifica a relação entre custos e benefícios técnicos
para definir o produto a ser desenvolvido, de acordo com a visão de
marketing da empresa. Desenvolve ferramentas e dispositivos para
fabricação mecânica, planejando as fases do desenvolvimento; realizando
leitura e interpretação do desenho de projeto; e pesquisando novas
tecnologias para melhoria do produto. Define normas e métodos a serem
seguidos. Detalha especificações da matéria-prima. Preenche planilhas de
custo. Elabora protótipo e coordena as fases de sua execução. Acompanha
testes de protótipos. Planeja o processo produtivo para produção de
ferramentas e dispositivos. Desenvolve embalagens para o produto.
Elabora manuais e catálogos técnicos do produto. Pode propor
atualizações no desenho de projeto. Pode desenvolver ferramentas e
dispositivos com sensores, interfaces digitais e outros recursos
similares. Pode propor projetos considerando aspectos ergonômicos de
ferramentas e dispositivos. Desenvolve o processo de produção de
ferramentas e dispositivos para fabricação mecânica, descrevendo as
fases do processo de fabricação. Define máquinas, equipamentos e
ferramentas para o processo produtivo; leiaute e fluxo do processo; e
mão de obra interna e externa. Indica os parâmetros de controle
dimensional, de tratamento térmico e de fluido de corte. Estabelece
procedimentos para manuseio e transporte de ferramentas. Prepara
instruções de rotina de trabalho para operadores. Elabora programas para
máquinas-ferramentas - inclusive para comandos numéricos
computadorizados (CNC) -, utilizando recursos de CAD-Desenho Assistido
por Computador (Computer Aided Design) e CAM--Manufatura Auxiliada por
Computador (Computer Aided Manufacturing). Avalia a capabilidade e a
eficácia do processo de produção. Aplica ferramentas da qualidade total
ao longo das etapas de desenvolvimento do processo produtivo. Pode
desenvolver processos de fabricação por tecnologia aditiva. Providencia
recursos técnicos para o processo de produção, identificando recursos;
definindo materiais a serem utilizados; e acompanhando a fabricação de
ferramentas e dispositivos. Proporciona suporte técnico na elaboração de
projetos, tendo em vista a compatibilização do projeto com o processo
produtivo. Define processos de execução de ferramentas e dispositivos.
Acompanha o teste (``tryout'') e define a aprovação de ferramentas e
dispositivos. Propõe ações de melhoria contínua, para o processo e para
o produto, detectando falhas de processo, implantando soluções para
correção dessas falhas e implementando melhorias no processo. Pesquisa
novas tecnologias para aprimoramento de processos. Identifica novos
materiais, tendo em vista a melhoria de parâmetros de qualidade das
ferramentas. Realiza estudos para redução de custos industriais. Presta
suporte técnico para manutenção e assistência técnica do produto,
elaborando lista de verificação (``checklist'') da manutenção de
ferramentas e dispositivos; e identificando pontos críticos para
reposição de peças. Pode propor alterações no projeto de ferramentas e
dispositivos em função de indicadores dos processos de manutenção e
assistência técnica. Avalia, junto a cliente, o produto fornecido,
quando solicitado. Treina técnicos, internos e externos, no uso do
produto. Pode supervisionar equipes de trabalho, distribuindo tarefas,
avaliando desempenho e desenvolvendo treinamentos. Mantém o local de
trabalho limpo, seguro e organizado. Conserva ferramentas, instrumentos
e equipamentos de trabalho. Trabalha com segurança, usando equipamentos
de proteção; participando de ações para prevenção de acidentes;
identificando situações de risco e propondo soluções para eliminá-las.
Realiza armazenamento e descarte de material de acordo com preceitos de
proteção do meio ambiente.

\begin{Form}
\textbf{Esta correspondência está adequada?} \\ 
\ChoiceMenu[radio,radiosymbol=\ding{52},bordercolor=0 0.5 0,name=cbovalid_ 314205 ]{}{CBO deve ficar=sim, \hspace{6em} CBO não fica aqui=nao}
\end{Form}

\textbf{CBO: 314210  ( Técnico mecânico na manutenção de ferramentas )}

Planeja manutenção preventiva e corretiva de ferramentas e dispositivos
para fabricação mecânica, utilizando dados de indicadores de manutenção;
elaborando cronograma de manutenção; e definindo normas e métodos a
serem seguidos no processo de manutenção. Controla uso e desgaste de
ferramentas e dispositivos, utilizando redes, sistemas de sensoriamento
e programas de computador, tendo em vista a eficiência do processo de
manutenção. Realiza manutenção preventiva e corretiva de ferramentas e
dispositivos para fabricação mecânica, utilizando lista de verificação
(``checklist'') de manutenção; e identificando pontos críticos para
reposição de peças. Define máquinas, equipamentos, ferramentas e
instrumentos a serem utilizados no processo de manutenção. Estabelece e
executa procedimentos de manuseio e transporte de equipamentos,
ferramentas e instrumentos, especialmente em atividades externas de
manutenção. Realiza controle dimensional. Analisa parâmetros de fluidos
de corte empregados nas operações com ferramentas e dispositivos. Faz
especificação de peças, materiais e outros insumos para manutenção.
Realiza procedimentos de detecção e de análise de falhas em ferramentas
e dispositivos. Elabora e implementa soluções para as falhas detectadas
e analisadas. Pode substituir partes ou peças em ferramentas e
dispositivos. Pode realizar procedimentos de manutenção em ferramentas e
dispositivos com sensores, interfaces digitais e outros recursos
similares. Presta assistência técnica a ferramentas e dispositivos para
fabricação mecânica, definindo normas e métodos a serem seguidos no
processo de assistência técnica; definindo especificações de peças,
materiais e outros insumos para o processo; e elaborando o plano de
visita técnica. Presta suporte técnico a clientes e visitantes. Avalia,
junto a cliente, o produto fornecido. Testa e aprova ferramentas e
dispositivos em campo. Pode propor alterações em ferramentas e
dispositivos, em função de resultados do processo de assistência
técnica. Propõe ações de melhoria contínua para o produto, com base nos
indicadores e em outras informações dos processos de manutenção e de
assistência técnica, aplicando ferramentas de qualidade total.
Identifica novas tecnologias para melhoria do produto e para
aprimoramento dos processos de manutenção e assistência técnica.
Identifica novos materiais, tendo em vista a melhoria de parâmetros de
qualidade das ferramentas. Treina técnicos internos e externos na
utilização de ferramentas e dispositivos. Pode supervisionar equipes de
trabalho, distribuindo tarefas, avaliando desempenho e desenvolvendo
treinamentos. Pode prestar suporte técnico na elaboração de projetos de
ferramentas e dispositivos, quando solicitado. Pode propor alterações no
projeto de ferramentas e dispositivos. Participa em estudos para redução
de custos da manutenção e emite pareceres técnicos sobre controle dos
custos. Mantém o local de trabalho limpo, seguro e organizado. Conserva
ferramentas, instrumentos e equipamentos de trabalho. Trabalha em
conformidade com as normas de segurança, utilizando Equipamentos de
Proteção Individual (EPI) e Equipamentos de Proteção Coletiva (EPC);
participando de ações para prevenção de acidentes; identificando
situações de risco e propondo soluções para eliminá-las. Orienta
clientes sobre riscos na utilização de ferramentas e dispositivos.
Trabalha integradamente com o departamento de segurança da empresa.
Realiza armazenamento e descarte de material de acordo com preceitos de
proteção do meio ambiente.

\begin{Form}
\textbf{Esta correspondência está adequada?} \\ 
\ChoiceMenu[radio,radiosymbol=\ding{52},bordercolor=0 0.5 0,name=cbovalid_ 314210 ]{}{CBO deve ficar=sim, \hspace{6em} CBO não fica aqui=nao}
\end{Form}

\textbf{CBO: 314610  ( Técnico em caldeiraria )}

Planeja processos de caldeiraria, interpretando projetos e consultando
normas técnicas aplicáveis. Estuda plantas para determinar localizações,
relações ou dimensões das peças e equipamentos. Define os procedimentos
de execução e inspeção. Elabora o plano de soldagem. Estima o material
de consumo e de aplicação. Seleciona os procedimentos para cada
atividade. Verifica as condições de segurança do ambiente. Organiza
equipamentos, ferramentas e materiais no local de trabalho. Define a
sequência de execução dos serviços de caldeiraria. Elabora a programação
da produção. Desenvolve projetos de caldeiraria pesada, caldeiraria leve
e média e chaparia, utilizando software de auxílio ao projeto e à
análise, com recursos de simulação. Opera processos de caldeiraria,
determinando equipamentos e métodos de trabalho, aplicando conhecimentos
de metalurgia, geometria e técnicas de caldeiraria. Monitora processos
de caldeiraria, utilizando procedimentos e sistemas específicos, tendo
em vista a garantia da qualidade exigida para as peças e os
equipamentos. Utiliza sistemas de gerenciamento de informações de
caldeiraria, tendo em vista a melhoria da produtividade e da qualidade
dos processos. Desenvolve modelos para projetos de caldeiraria, usando
cálculos matemáticos baseados em informações de desenhos de projeto. Ao
receber notificações de mau funcionamento de equipamentos ou de
materiais defeituosos, define as medidas necessárias para solução dos
problemas operacionais dos processos de caldeiraria. Coordena as
atividades de execução dos serviços de caldeiraria de acordo com o
cronograma e a programação da produção. Verifica a montagem, tendo em
vista o cumprimento dos requisitos do projeto. Controla o consumo do
material. Realiza o controle dimensional. Libera o serviço executado
para inspeção do controle de qualidade. Coordena a realização de ensaios
mecânicos em peças e equipamentos de caldeiraria, aplicando
procedimentos específicos, para garantir o nível de qualidade exigido
para o produto e o atendimento aos requisitos do projeto. Inspeciona os
processos, examinando matéria-prima para os serviços de caldeiraria;
observando consumíveis de soldagem; e verificando peças e equipamentos
de caldeiraria montados ou componentes individuais - como tubos,
conexões, válvulas, controles ou mecanismos auxiliares -, para localizar
falhas de processo. Examina caldeiras, vasos de pressão, tanques, cubas
ou outros artefatos de caldeiraria para localizar defeitos - como
vazamentos, pontos fracos ou seções com defeito -, para operações de
reparo. Implementa soluções alternativas para corrigir falhas em
processos de caldeiraria. Supervisiona processos de caldeiraria -- no
desenvolvimento de projetos, na fabricação, na montagem, na instalação e
na manutenção de artefatos de caldeiraria --, tendo em vista a eficácia
e a eficiência dos serviços. Aplica as normas técnicas na execução dos
serviços de caldeiraria. Pode supervisionar equipes de trabalho,
distribuindo tarefas, avaliando desempenho e desenvolvendo treinamentos.
Incorpora tecnologias inovadoras aos processos de caldeiraria, definindo
aplicações para equipamentos de trabalho automatizados ou incluindo
recursos robóticos nos processos de calandragem e soldagem. Define
técnicas, métodos, máquinas e equipamentos para operar com novos
materiais, tais como novas ligas metálicas, novos tipos de chapas de
aço, novos sensores e outros recursos de automação industrial. Segue
normas regulamentadoras e procedimentos de segurança e saúde no trabalho
em processos de caldeiraria, utilizando e orientando pessoal para
utilizar equipamentos de proteção individual e coletiva. Promove o
descarte e a destinação de resíduos de acordo com procedimentos e normas
de preservação ambiental.

\begin{Form}
\textbf{Esta correspondência está adequada?} \\ 
\ChoiceMenu[radio,radiosymbol=\ding{52},bordercolor=0 0.5 0,name=cbovalid_ 314610 ]{}{CBO deve ficar=sim, \hspace{6em} CBO não fica aqui=nao}
\end{Form}

\newpage
\section{ 02008 -- Técnico em Fundição }

\textbf{Perfil do Curso (CNCT)}

O Técnico em Fundição será habilitado para:

Elaborar, desenvolver e executar projetos de moldes para fundição de
metais respeitando procedimentos e normas técnicas, de qualidade, de
saúde e segurança e de meio ambiente. Realizar a gestão das etapas de
fusão e beneficiamento de materiais ferrosos e não ferrosos para o
processo de fundição. Elaborar ensaios e análises químicas dos metais
empregados em processos de fundição. Zelar pelo atendimento de questões
ambientais específicas relacionadas ao processamento de fundição de
metais. Controlar defeitos e garantir a qualidade dos produtos
fabricados. Simular processos de fundição pela utilização de softwares
específicos. Reconhecer tecnologias inovadoras presentes no segmento
visando a atender às transformações digitais na sociedade.

Para atuação como Técnico em Fundição, são fundamentais:

Conhecimentos e saberes relacionados aos processos de planejamento e
controle nos processos de fundição de modo a assegurar a saúde e a
segurança de trabalhadores e futuros usuários. Conhecimentos e saberes
relacionados à sustentabilidade do processo produtivo, às técnicas e aos
processos de produção, às normas técnicas, à liderança de equipes, à
solução de problemas técnicos e trabalhistas e à gestão de conflitos.

\textbf{CBOs correspondidos e seus perfis (presentes no CNCT, e presentes na predição da IA)}

\textbf{CBO: 314715  ( Técnico de fundição em siderurgia )}

Planeja as atividades relacionadas aos processos siderúrgicos de
fundição, descrevendo-as, definindo responsabilidades para cada etapa
dos processos e estabelecendo prazos de entrega de produtos e serviços.
Identifica necessidade de melhoria de produtos -- no âmbito de processos
siderúrgicos de fundição - e participa da elaboração de projetos de
fabricação de produtos siderúrgicos. Mantém contato com clientes e
fornecedores para troca de informações. Controla processos siderúrgicos
de fundição, realizando cálculos de área, volume e peso de produtos e
equipamentos siderúrgicos. Calcula o balanço de massa e as quantidades
de fundentes para adição aos processos. Calcula o balanço térmico dos
processos e calcula as quantidades de produtos para adição, a fim de
realizar a correção química dos processos. Pode utilizar sistemas
simuladores de processos siderúrgicos. Verifica desgaste de máquinas,
equipamentos e outros recursos de trabalho utilizados na fundição
siderúrgica. Coordena paradas de equipamentos. Define itens de controle
da qualidade dos processos siderúrgicos de fundição e monitora
resultados dos indicadores desses itens de controle. Inspeciona produtos
visualmente. Realiza ensaios físicos e químicos e interpreta seus
resultados. Interpreta resultados de ensaios mecânicos e metalográficos.
Avalia defeitos de produtos e analisa suas causas, a fim de elaborar
planos de ações preventivas e corretivas para os processos. Analisa
aproveitamento de produtos não conformes. Presta assistência técnica
para clientes, analisando suas reclamações e lhes sugerindo a melhor
forma de utilização de produtos. Pode adaptar processos às
especificações de clientes, a partir da identificação de suas
necessidades. Elabora planos experimentais de novos produtos e de
implantação de novas tecnologias, fazendo uso de materiais clássicos e
de novos materiais na produção siderúrgica; introduzindo equipamentos de
maior eficiência energética; utilizando fontes alternativas de energia
nos processos; e aplicando tecnologias digitais e de automatização,
tendo em vista a eficiência, a rentabilidade, a segurança e a redução de
desperdícios. Divulga os planos experimentais para os profissionais
envolvidos nos processos. Programa e monitora a execução das etapas de
planos experimentais. Avalia e valida resultados dos planos
experimentais desenvolvidos. Identifica atividades técnicas e
operacionais, elabora manuais de procedimentos, normaliza e revisa
padrões e procedimentos dos processos siderúrgicos de fundição.
Interpreta desenhos técnicos. Elabora relatórios técnicos para
documentação dos processos siderúrgicos de fundição, tais como
relatórios técnicos de produção; relatórios de planos experimentais;
relatórios de visitas e de assistência técnica; relatórios de controle
de qualidade; e relatórios de reclamação de clientes. Elabora informes
técnicos dos processos. Promove os meios para desenvolvimento
profissional da equipe de trabalho, com focalização nos processos
siderúrgicos de fundição. Identifica necessidades de desenvolvimento de
pessoal, para elaborar plano de desenvolvimento. Ministra treinamento
e/ou solicita contratação de terceiros para realização dos programas.
Avalia a eficácia dos treinamentos ministrados. Providencia serviços --
inclusive terceirizados -- de manutenção de máquinas, equipamentos e
outros recursos utilizados em processos de fundição. Atua de acordo com
as diretrizes adotadas pela indústria siderúrgica para o setor de
fundição, a fim de atender exigências de órgãos fiscalizadores e
requisitos de clientes e de mercado, no que se refere às condições do
local de trabalho, ao meio ambiente e ao desenvolvimento sustentável.
Aplica normas e procedimentos de segurança, utilizando equipamentos de
proteção individual e coletiva e desenvolvendo ações de prevenção de
acidentes de trabalho. Promove o descarte e a destinação de resíduos de
acordo com diretrizes normatizadas e procedimentos.

\begin{Form}
\textbf{Esta correspondência está adequada?} \\ 
\ChoiceMenu[radio,radiosymbol=\ding{52},bordercolor=0 0.5 0,name=cbovalid_ 314715 ]{}{CBO deve ficar=sim, \hspace{6em} CBO não fica aqui=nao}
\end{Form}

\textbf{CBO: 721325  ( Polidor de metais )}

Planeja o trabalho de polimento, identificando o material a ser polido e
o grau de polimento, de acordo com o padrão de qualidade exigido. Estuda
desenhos e especificações das peças metálicas a serem polidas, para
determinar os procedimentos de polimento, bem como planejar
configurações da máquina e sequências operacionais. Estuda novos
métodos, equipamentos e materiais empregados nas atividades de
polimentos de metais, avaliando a viabilidade de aplicação e o alcance
de melhores desempenhos. Seleciona equipamento ou máquina e ferramenta
para polir, baseando-se no desenho da peça, especificações ou outras
instruções, bem como nos conhecimentos dos requisitos de máquina e de
produção, para assegurar o bom resultado do trabalho. Define parâmetros
-- tempo, quantidade, tipo de material e granulometria -- para aplicação
de abrasivo na superfície metálica a ser polida. Estabelece velocidades
da politriz, em rotações por minuto, de acordo com tabelas de
especificações. Define a medida mínima do instrumento polidor -- rebolo,
ferramenta giratória em forma de escova circular ou similares -- e
estima a produção de polimento. Executa polimento por processo manual,
montando a roda de polir -- em forma de escova ou disco -- na politriz,
ajustando o perfil da roda -- pano ou sisal -- com a ferramenta de
corte. Unta a ferramenta escolhida, utilizando cera ou outra substância,
para facilitar o trabalho. Coloca a máquina em funcionamento, pondo em
contato a ferramenta abrasiva com os diferentes lados da peça metálica a
ser polida. Pressiona o instrumento polidor conforme o estado da
superfície e tipo de acabamento, observando o intervalo de tempo
necessário e estabelecendo a sequência de movimentos, para eliminar as
imperfeições da superfície e obter o polimento desejado. Identifica fim
da vida útil da ferramenta de polir. Substitui disco ou escova e mantém
o equipamento de polir em ordem, utilizando ferramentas próprias, para
melhorar ou modificar o tipo de polimento. Executa polimento por
processo automático, seguindo instruções de trabalho de polimento e
programando a máquina para polir. Monta a roda de polir. Coloca chips e
esferas em equipamentos vibratórios de polir. Coloca a peça a ser polida
na politriz. Retira peça polida da máquina. Controla desgaste de chips,
esferas e rodas de polir. Repõe chips e esferas. Substitui rodas de
polir, quando necessário. Pode limpar os objetos metálicos em banho
químico, antes de poli-los. Pode realizar polimento a laser. Controla a
qualidade, verificando as dimensões das peças ou lotes no início do
processo e definindo quantidade - totalidade ou amostras das peças - a
inspecionar. Identifica os pontos a inspecionar e realiza a inspeção do
produto, visualmente ou fazendo uso de equipamentos. Examina e mede as
peças para garantir que as superfícies e as dimensões atendam às
especificações ao final do processo. Identifica material a ser refugado
ou retrabalhado. Executa manutenção básica - como limpeza e lubrificação
de peças - das máquinas e ferramentas de polir. Remove e substitui peças
gastas ou quebradas, usando ferramentas manuais. Trabalha com segurança,
seguindo normas de segurança e de medicina do trabalho. Ajusta as
condições ergonômicas. Utiliza equipamentos de proteção individual.
Participa do levantamento de riscos ambientais e de ações preventivas
contra incêndios e acidentes. Presta primeiros socorros. Conserva o
ambiente de trabalho limpo e organizado. Destina adequadamente os
resíduos gerados, reciclando materiais e realizando descarte de acordo
com as normas ambientais.

\begin{Form}
\textbf{Esta correspondência está adequada?} \\ 
\ChoiceMenu[radio,radiosymbol=\ding{52},bordercolor=0 0.5 0,name=cbovalid_ 721325 ]{}{CBO deve ficar=sim, \hspace{6em} CBO não fica aqui=nao}
\end{Form}

\textbf{CBO: 722215  ( Operador de acabamento de peças fundidas )}

Prepara o trabalho, inspecionando a peça fundida, no início do processo
de acabamento. Verifica a existência, o formato e a resistência de
rebarbas e examina as irregularidades da superfície de peças fundidas.
Retira resíduos de areia e excesso de metal -- rebarbas -- das peças
fundidas, operando equipamentos de ar comprimido ou marteletes
pneumáticos. Homogeneíza as peças, utilizando máquinas de corte,
esmerilhadora, esmeril fixo e marteletes. Remove irregularidades da
superfície das peças fundidas, separando os canais e os alimentadores
resultantes da fundição das peças, com o uso de discos de corte, esmeril
e marteletes. Remove defeitos da fundição na superfície de peças
fundidas, recuperando-as por meio de ensaios de ultrassom, partículas
magnéticas ou líquidos penetrantes. Pode preparar equipamento de
jateamento robotizado. Faz a limpeza da peça, operando equipamento de
jateamento com granalhas, retirando carepas, areia e oxidação
superficial. Pinta e armazena as peças acabadas nos locais
identificados. Conserva o local das atividades limpo e organizado.
Mantém materiais e apetrechos organizados e acondicionados. Observa as
condições de funcionamento de ferramentas, máquinas e equipamentos
utilizados, executando pequenos reparos elétricos, mecânicos,
refratários e hidráulicos, para obter o bom andamento do serviço.
Solicita serviço de manutenção de ferramentas, máquinas e equipamentos,
quando necessário. Destina adequadamente os resíduos gerados, realizando
a reciclagem de materiais e o descarte, de acordo com as normas
ambientais. Trabalha com segurança, usando equipamentos de proteção
individual e atuando na prevenção de acidentes.

\begin{Form}
\textbf{Esta correspondência está adequada?} \\ 
\ChoiceMenu[radio,radiosymbol=\ding{52},bordercolor=0 0.5 0,name=cbovalid_ 722215 ]{}{CBO deve ficar=sim, \hspace{6em} CBO não fica aqui=nao}
\end{Form}

\textbf{CBO: 722220  ( Operador de máquina centrifugadora de fundição )}

Prepara a operação da máquina, fixando o molde metálico e aprontando-o
para o recebimento do metal líquido. Inspeciona a máquina quanto a
riscos, eficiência operacional, mau funcionamento, desgaste ou
vazamentos. Limpa o pó ou as impurezas metálicas existentes no molde,
utilizando escovas apropriadas, para assegurar a perfeita moldação da
peça. Fixa uma tampa na extremidade do molde, empregando ferramentas
específicas, para evitar derramamento do metal. Faz o aquecimento prévio
do molde, para deixá-lo na temperatura adequada ao recebimento do metal.
Opera a máquina, colocando-a em movimento, atuando sobre os comandos,
para possibilitar a rotação do molde. Aplica um agente isolante no molde
em rotação, pulverizando-lhe a superfície interior, para evitar que o
metal adira às paredes no molde. Vaza o metal líquido no molde. Faz o
molde girar durante o tempo necessário para possibilitar a solidificação
do metal. Retira a tampa do molde metálico e a peça fundida, utilizando
tenazes ou pistão hidráulico, para possibilitar nova operação da
máquina. Retira e substitui o molde metálico no equipamento, conforme a
programação da produção. Pode coletar amostras do metal líquido para
análise, durante o processo. Inspeciona visualmente a qualidade das
peças prontas. Participa das ações de melhoramento contínuo dos
processos de fundição, podendo operar equipamentos automatizados de
fundição por centrifugação. Pode operar equipamentos de fundição por
centrifugação de materiais de baixo ponto de fusão - tais como ligas de
estanho ou chumbo -, em moldes de borracha, para produzir peças
metálicas de geometrias complexas. Mantém registros de leituras de
instrumentos e resultados de turnos de produção para entrada em banco de
dados de computadores. Executa manutenção básica, como limpeza e
lubrificação de peças de máquinas e equipamentos. Remove e substitui
peças gastas ou quebradas, usando ferramentas manuais. Atua de modo a
colaborar para a redução de riscos e impactos socioambientais. Trabalha
com segurança, usando equipamentos de proteção individual e atuando na
prevenção de acidentes.

\begin{Form}
\textbf{Esta correspondência está adequada?} \\ 
\ChoiceMenu[radio,radiosymbol=\ding{52},bordercolor=0 0.5 0,name=cbovalid_ 722220 ]{}{CBO deve ficar=sim, \hspace{6em} CBO não fica aqui=nao}
\end{Form}

\textbf{CBO: 722325  ( Operador de equipamentos de preparação de areia )}

Prepara as atividades, interpretando as ordens de produção. Confere as
condições de segurança e de funcionamento dos equipamentos. Examina o
nível do silo, verifica a medida da areia no misturador e inspeciona a
qualidade da areia. Identifica os tipos de areias de moldação e macharia
-- naturais, sintéticos e semissintéticos --, utilizados na confecção de
moldes e machos, de acordo com as ordens de produção. Seleciona e
providencia os materiais -- areias e resinas aglomerantes --, para
permitir a mistura adequada na confecção de moldes e machos. Opera
equipamentos de preparação de areia, acionando dispositivos de comando
mecânico ou eletrônico no painel da máquina. Processa a mistura da
areia, adicionando-lhe os aglomerantes necessários, dentro das
proporções e especificações estabelecidas, a fim de prepará-la para a
moldação ou macharia. Transporta os materiais componentes da areia de
fundição, utilizando carrinhos manuais ou operando equipamentos
mecanizados de transporte de materiais -- como correias transportadoras
--, para fazer a distribuição ou abastecer os silos. Efetua a
recuperação de areias após a utilização de moldes e machos nos processos
de fundição, processando a eliminação de impurezas, para utilizá-las
novamente. Retira do molde o peso de lastro ou grampo e o excesso de
metal sólido. Extrai a peça, separando os canais para a área de
acabamento e limpeza. Retira o excesso de areia da peça. Regenera os
diferentes tipos de areia de moldação ou macharia. Preenche formulário
de controle da produção de areia preparada. Faz a requisição dos
materiais necessários, em função das quantidades estocadas, para manter
abastecidos os depósitos. Limpa os equipamentos de preparação de areia.
Solicita manutenção dos misturadores, correias transportadoras, silos,
peneiras e outros recursos de trabalho utilizados. Desenvolve as
atividades observando as exigências ambientais e as medidas para
prevenção de acidentes. Usa equipamentos de proteção individual.

\begin{Form}
\textbf{Esta correspondência está adequada?} \\ 
\ChoiceMenu[radio,radiosymbol=\ding{52},bordercolor=0 0.5 0,name=cbovalid_ 722325 ]{}{CBO deve ficar=sim, \hspace{6em} CBO não fica aqui=nao}
\end{Form}

\textbf{CBO: 771110  ( Modelador de madeira )}

Planeja o trabalho, verificando ordens de serviço e de produção e
interpretando projetos, desenhos e especificações técnicas para a
confecção e restauração de matrizes de madeira e derivados. Seleciona
matéria-prima, insumos e componentes, instrumentos de medição e
controle, ferramentas manuais, ferramentas elétricas portáteis, máquinas
e equipamentos, para a execução do trabalho. Define roteiro para as
atividades de confecção e restauração. Pode elaborar orçamento. Elabora
croqui ou esboço e realiza cálculos para dimensionamento do produto,
conforme documentos de solicitação. Pode utilizar software de apoio para
desenhar o projeto completo em computador, especificando os materiais e
as operações para a confecção do produto. Confecciona gabaritos ou
moldes para execução das peças de madeira e derivados, de acordo com
especificações. Prepara máquinas e ferramentas. Afia ferramentas e
ajusta parâmetros operacionais, tais como de posicionamento e geometria
de corte, velocidade e pressão da ferramenta, entre outros. Confecciona
matrizes de madeira e derivados, aplicando processos de usinagem,
programando, operando e regulando as máquinas para obter o produto
conforme o projeto. Realiza cálculos e traça os contornos da peça,
utilizando instrumentos de desenho e traçagem e ajustando o
posicionamento e a geometria do corte. Usina as partes, definindo a
sequência de operações a serem aplicadas para cortar, serrar, furar,
aplainar, tornear, lixar, fresar e laminar. Ajusta os elementos e aplica
materiais e dispositivos de fixação. Pode programar e utilizar
máquinas-ferramenta com comando numérico computadorizado (CNC) e
programas para usinagem das peças integrados a ferramentas de software
de CAD/CAM. Monta matrizes de madeira e derivados utilizando elementos
de fixação e colocando acessórios e ferragens para reajuste. Coloca
arremates finais. Verifica as peças e os produtos confeccionados, para
garantir a qualidade das matrizes de madeira e derivados, avaliando a
matéria-prima, a resistência física, aferindo as formas e as dimensões e
as condições do acabamento. Compara as características do produto final
com os requisitos do cliente ou do projeto, providenciando ajustes,
quando necessários. Desmonta o produto confeccionado, no todo ou em
parte, para posterior entrega. Prepara a expedição de matrizes de
madeira e derivados conforme demanda do cliente, embalando separadamente
os acessórios e os elementos principais; definindo as ferramentas a
serem utilizadas para fazer o transporte e a entrega do produto; e
apondo os arremates finais. Restaura matrizes de madeira e derivados,
avaliando produtos para determinar a viabilidade de restauração,
confeccionando peças para reposição, substituindo peças danificadas,
reapertando elementos de fixação e preparando o produto para acabamento.
Organiza o local de trabalho, para o desenvolvimento das atividades.
Mantém os materiais, instrumentos e acessórios de trabalho organizados e
acondicionados. Mantém as máquinas, os equipamentos e as ferramentas em
plenas condições de uso e funcionamento. Realiza a limpeza e a
conservação dos recursos materiais. Controla desperdícios de material.
Separa, identifica e classifica resíduos -- sólidos e líquidos -, para
descarte, reúso e reciclagem. Preenche documentação técnica relacionada
aos processos, tais como ordens de serviço, fichas de controle de
processo, ficha técnica do produto, requisição de materiais, entre
outros. Interpreta normas, manuais técnicos e catálogos. Utiliza
equipamentos de proteção individual e coletiva.

\begin{Form}
\textbf{Esta correspondência está adequada?} \\ 
\ChoiceMenu[radio,radiosymbol=\ding{52},bordercolor=0 0.5 0,name=cbovalid_ 771110 ]{}{CBO deve ficar=sim, \hspace{6em} CBO não fica aqui=nao}
\end{Form}

\textbf{CBO: 821220  ( Forneiro e operador (refino de metais não ferrosos) )}

Prepara máquinas, equipamentos e materiais, conferindo a disponibilidade
de matérias-primas e insumos, pesando minérios e fundentes,
inspecionando visualmente máquinas e equipamentos e comunicando-se com
as áreas envolvidas. Solicita manutenção para máquinas e equipamentos,
quando necessário. Opera e controla o funcionamento de fornos para
refino de metais não ferrosos, acionando comandos apropriados,
executando o carregamento do forno, orientando a entrada dos materiais e
acendendo a carga, para processar as operações de fusão. Opera
equipamentos auxiliares. Controla a temperatura do forno, regulando a
combustão, a fim de assegurar a conformidade do processo com as
especificações técnicas. Acompanha o processo de fusão, observando a
coloração e a consistência do metal ou procedendo à leitura dos
indicadores de temperatura, para determinar o momento oportuno de início
da corrida. Pode realizar a operação e o controle com o auxílio de
painéis ou telas de sistemas de supervisão de processos, observando seus
gráficos e instrumentos -- físicos ou virtuais --, e efetuando as
intervenções necessárias para manter o fluxo de operação dentro dos
parâmetros especificados para o processo. Vaza o metal fundido,
providenciando a abertura do forno e desobstruindo sua entrada, para
possibilitar o vazamento. Retira amostras de metal e escória, usando
ferramentas e técnicas específicas, para análise e controle da qualidade
do produto. Interrompe o processo, para eliminar bloqueios ou
atolamentos. Executa a granulagem da escória. Controla as
características físico-químicas da matéria-prima e do produto, medindo a
temperatura do metal líquido e monitorando visualmente o metal e a
escória. Controla parâmetros operacionais -- tais como pressão, vazão e
temperatura --, para mantê-los nos níveis especificados para o processo.
Pode coletar amostras do metal fundido para testes de laboratório, a fim
de garantir a conformidade com as especificações. Na fase de tratamentos
secundários, realiza a desgaseificação para tratar metais não ferrosos.
Participa das ações de melhoria nas atividades com fornos de refino de
metais não ferrosos, fazendo uso de tecnologias de automação e controle
- incluindo os sistemas de supervisão -, para aprimoramento dos
parâmetros tecnológicos do processo e elevação do nível de qualidade dos
vários tipos de metais não ferrosos produzidos. Colabora no
desenvolvimento de novos processos de fusão e refino de metais não
ferrosos. Mantém o local de trabalho limpo e organizado. Zela pela
conservação e pela manutenção de máquinas e equipamentos. Cumpre normas
de segurança pessoal e ambiental, usando Equipamentos de Proteção
Individual (EPI) e Equipamentos de Proteção Coletiva (EPC), e coletando
lixos e resíduos industriais. Respeita sinalização de segurança e
comunica acidentes e incidentes. Opera sistema de despoeiramento.

\begin{Form}
\textbf{Esta correspondência está adequada?} \\ 
\ChoiceMenu[radio,radiosymbol=\ding{52},bordercolor=0 0.5 0,name=cbovalid_ 821220 ]{}{CBO deve ficar=sim, \hspace{6em} CBO não fica aqui=nao}
\end{Form}

\textbf{CBO: 822115  ( Forneiro de fundição (forno de redução) )}

Prepara o forno elétrico de redução para fundição de metais e ligas
metálicas, estabelecendo a sequência de operações e definindo a
temperatura para fusão. Prepara ferramentas e instrumentos e verifica as
condições de funcionamento de equipamentos e acessórios. Aquece os
recipientes de transporte de metal e calafeta as calhas de vazamento.
Aplica desmoldantes em calhas e bacias de vazamento. Prepara a carga
para carregar o forno, inspecionando a sucata e as matérias-primas;
selecionando a carga; e controlando a quantidade de carga. Carrega o
forno, removendo a abóbada do forno; colocando os materiais a fundir no
forno; e adicionando fundentes. Utiliza equipamentos de movimentação de
carga para carregar e distribuir a carga no interior do forno. Realiza
fundição de metais e ligas metálicas em forno elétrico de redução,
monitorando o funcionamento do forno; acionando o sistema de aquecimento
e de resfriamento; e verificando as condições físico-químicas da carga.
Retira escória do forno e limpa sistema de exaustão de fornos, calhas de
vazamento e escórias das bicas. Abre orifícios de escoamento de metal
líquido. Transfere cargas entre fornos. Ajusta a composição química de
ligas metálicas, adicionando ligas ao metal líquido; homogeneizando o
metal; e completando a carga de metal no forno. Desgaseifica e
escorifica o metal fundido. Controla a temperatura dos metais. Coleta
amostras de material para análises laboratoriais, avalia os resultados
de análises de composição química do metal e, se necessário, corrige o
teor de elementos químicos. Realiza o escoamento do metal fundido,
escumando o metal, desobstruindo calhas de escoamento, basculando o
forno e extraindo metal do forno. Armazena o material extraído para
transporte. Registra ocorrências técnicas e operacionais - em meio
físico ou digital -, detalhando consumo de insumos; produção diária; e
paradas de equipamentos e de máquinas. Registra dados do processo, como
temperatura e dimensões do metal e do produto. Registra os resultados de
testes e ensaios. Requisita materiais de consumo. Elabora relatórios
técnicos. Prepara o forno de redução para manutenção, desligando-o;
resfriando-o; esgotando lastros do forno; e retirando a camada
refratária e os acessórios elétricos e mecânicos. Analisa o histórico de
funcionamento do forno. Repara o revestimento refratário ou realiza o
revestimento com material refratário novo. Substitui abóbadas de fornos.
Resfria e limpa os recipientes de transporte do metal fundido. Mantém o
ambiente de trabalho limpo e organizado. Presta apoio à política de
preservação ambiental, reaproveitando materiais e descartando resíduos
conforme procedimentos adotados pela empresa. Trabalha com segurança,
acionando sistemas de segurança de equipamentos e utilizando
equipamentos de proteção individual e coletiva. Participa de ações de
prevenção de acidentes propostas pela empresa. Identifica condições e
atos inseguros, relatando-os aos membros da CIPA-Comissão Interna de
Prevenção de Acidentes. Manuseia equipamentos de combate a incêndio.
Identifica necessidade de substituição de equipamentos de proteção e
verifica a validade e as condições de equipamentos de combate a
incêndio. Providencia primeiros socorros.

\begin{Form}
\textbf{Esta correspondência está adequada?} \\ 
\ChoiceMenu[radio,radiosymbol=\ding{52},bordercolor=0 0.5 0,name=cbovalid_ 822115 ]{}{CBO deve ficar=sim, \hspace{6em} CBO não fica aqui=nao}
\end{Form}

\newpage
\section{ 02009 -- Técnico em Instrumentação Industrial }

\textbf{Perfil do Curso (CNCT)}

O Técnico em Instrumentação Industrial será habilitado para:

Implantar, configurar e manter sistemas de instrumentação e controle de
processos industriais. Elaborar, desenvolver e executar projetos de
instalações de instrumentos de medição de variáveis industriais, de
acordo com normas técnicas, ambientais, de saúde e segurança no trabalho
e de qualidade. Realizar programação, parametrização, medições e testes
de equipamentos por meio de instrumentação industrial. Reconhecer
tecnologias inovadoras presentes no segmento visando a atender às
transformações digitais na sociedade.

Para atuação como Técnico em Instrumentação Industrial, são
fundamentais:

Conhecimentos e saberes relacionados aos processos de planejamento e
implementação de controle de processos industriais por meio de
instrumentação de modo a assegurar a saúde e a segurança dos
trabalhadores. Conhecimentos e saberes relacionados à sustentabilidade
do processo produtivo, às técnicas e processos de produção, às normas
técnicas, à liderança de equipes, à solução de problemas técnicos e
trabalhistas e à gestão de conflitos

\textbf{CBOs correspondidos e seus perfis (presentes no CNCT, e presentes na predição da IA)}

\textbf{CBO: 313410  ( Técnico em instrumentação )}

Participa no desenvolvimento de projetos de sistemas de medição e
controle, identificando as variáveis envolvidas no processo,
especificando instrumentos e determinando posições de medição e
controle. Examina a possibilidade de uso de instrumentação inteligente
e/ou interligada por meio de redes industriais digitais e de Internet
das Coisas Industrial (IIoT), propondo tipos de sistema adequados às
aplicações. Implanta sistemas de instrumentação e controle em processos
industriais, analisando tecnicamente a aquisição de produtos e serviços
de medição e controle, estabelecendo objetivos, indicando fornecedores
potenciais, testando e comparando características técnicas de produtos e
serviços de fornecedores diversificados. Emite parecer técnico.
Determina valores de grandezas para o controle de processos, lendo e
interpretando desenhos técnicos, selecionando métodos e princípios de
medição e realizando ensaios - físicos, elétricos, mecânicos,
eletrônicos, dentre outros -, correspondentes às grandezas a serem
medidas. Opera equipamentos, sistemas e instrumentos de medição e
controle. Garante a calibração dos equipamentos, instrumentos e sistemas
utilizados no processo, selecionando prestadores de serviços de
calibração. Monitora a validade de calibração. Pode efetuar calibração,
em laboratórios credenciados, seguindo metodologias específicas. Pode
emitir laudos e certificados de calibração de equipamentos e
instrumentos de medição e controle. Participa da gestão da documentação
técnica, elaborando, atualizando e analisando criticamente procedimentos
e instruções de trabalho, cadastrando instrumentos de medição e controle
e elaborando fichas e formulários de registros. Participa da gestão dos
procedimentos do sistema de confiabilidade metrológica, colaborando na
elaboração de sistema de codificação, estabelecendo frequência,
controlando prazos e validando resultados da calibração. Participa de
auditorias internas e externas. Participa do planejamento de manutenção
preventiva, preditiva e corretiva de instrumentos de medição e controle.
Realiza a manutenção de instrumentos de medição e controle,
desinstalando os instrumentos para realização dos serviços de
manutenção, identificando disfunções e suas causas, reparando ou
substituindo componentes danificados, reinstalando e testando os
instrumentos.

\begin{Form}
\textbf{Esta correspondência está adequada?} \\ 
\ChoiceMenu[radio,radiosymbol=\ding{52},bordercolor=0 0.5 0,name=cbovalid_ 313410 ]{}{CBO deve ficar=sim, \hspace{6em} CBO não fica aqui=nao}
\end{Form}

\newpage
\section{ 02010 -- Técnico em Manutenção Aeronáutica em Aviônicos }

\textbf{Perfil do Curso (CNCT)}

O Técnico em Manutenção de Aeronática em Aviônicos será habilitado para:

Programar, controlar e executar manutenção preventiva e corretiva dos
sistemas elétricos e eletrônicos de navegação, comunicação,
monitoramento e controle de aeronaves atendendo às normas e aos padrões
técnicos de qualidade, saúde e segurança e de meio ambiente. Aplicar
procedimentos de manuais de fabricantes, publicações técnicas e normas
nacionais e internacionais do setor aeronáutico. Diagnosticar as
condições dos instrumentos que compõem uma aeronave e fazer testes de
comissionamento e de performance em equipamentos de aeronaves. Indicar
os processos de manutenção a serem executados na revisão de aeronaves,
bem como orientar o balizamento de aeronaves.

Para atuação como Técnico em Manutenção de Aeronática em Aviônicos, são
fundamentais:

Conhecimentos e saberes relacionados aos processos de planejamento e
manutenção de aeronaves de modo a assegurar a saúde e a segurança dos
trabalhadores e dos usuários. Conhecimentos e saberes relacionados à
sustentabilidade do processo produtivo, às técnicas e aos processos de
manutenção na aviação, às normas técnicas, à constante atenção às
atualizações presentes no setor de aviação, à liderança de equipes, à
solução de problemas técnicos e trabalhistas e à gestão de conflitos.

\textbf{CBOs correspondidos e seus perfis (presentes no CNCT, e presentes na predição da IA)}

\textbf{CBO: 314310  ( Técnico mecânico (aeronaves) )}

Prepara ensaios de fabricação de equipamentos, estruturas e sistemas
mecânicos de aviões e helicópteros, verificando cronograma do processo
de fabricação e interpretando desenhos de projeto, para elaborar croquis
e fabricar ou adaptar peças, ferramentas e dispositivos a serem
utilizados nos ensaios. Realiza ensaios de fabricação de equipamentos,
estruturas e sistemas mecânicos de aviões e helicópteros, abrangendo
fuselagem, montantes, naceles, capotas de motor, carenagens, rotores e
outras superfícies aerodinâmicas de célula, trens de aterrissagem,
sistemas de freio e demais equipamentos mecânicos da aeronave, incluindo
seus acessórios e controles. Identifica necessidade de ajustes
dimensionais, adapta peças e acessórios e torna modificações adequadas
às características da aeronave. Participa dos testes finais de ensaios e
montagem. Prepara a montagem de componentes em aviões e helicópteros,
interpretando desenhos técnicos, verificando a disponibilidade de itens
no estoque para as montagens requeridas e analisando condições do local
de instalação e dos componentes, para realizar o preparo do local.
Executa, seguindo as instruções de trabalho, a montagem de componentes
em bancada e os instala em sistemas da aeronave de acordo com o projeto,
utilizando ferramentas, instrumentos e equipamentos. Elimina
interferências prejudiciais à segurança, assegurando ajustes entre peças
e destas com a estrutura. Realiza testes funcionais de aviões e
helicópteros, cumprindo procedimentos, seguindo parâmetros estabelecidos
pelo fabricante e utilizando equipamentos específicos e calibrados.
Executa testes iniciais para confirmação dos problemas apresentados,
testa sistemas operacionais mecânicos, de estrutura da aeronave e de
aviônica aplicada ao controle de equipamentos e sistemas mecânicos.
Simula processos, realiza análise dimensional por etapas de montagem e
participa da execução de provas de inclinação da aeronave. Libera
processos para itens de segurança. Prepara a execução de manutenção de
estruturas, equipamentos e sistemas mecânicos de aviões e helicópteros,
identificando sua forma de utilização, seus defeitos e suas falhas.
Executa os reparos, utilizando ferramentas e equipamentos apropriados,
substituindo e regulando componentes e utilizando técnicas de proteção
química e de acabamento de superfície. Nivela a linha de eixo do motor
da aeronave. Presta assessoria técnica interna e externa, treinando
equipes de trabalho, acompanhando o processo de montagem e manutenção,
provendo suporte técnico a equipes de trabalho e clientes, realizando
atendimento de emergência e fazendo visita técnica. Analisa resultados
da assessoria prestada, com a finalidade de assegurar a satisfação dos
clientes. Participa do programa de qualidade da empresa, seguindo as
normas internas e as definidas por organismos externos competentes,
utilizando ferramentas da qualidade, seguindo procedimentos de serviço e
identificando pontos críticos para propor melhorias, tendo em vista o
alcance de metas de qualidade. Fornece informações técnicas para
execução do controle estatístico do processo. Elabora documentação
técnica. Faz orçamento de serviços, com estimativa de tempo e de custo.
Emite parecer técnico. Elabora plano de ação referente às
não-conformidades detectadas em produtos e processos. Elabora relatório
técnico e de campo, identificando os componentes instalados e
registrando a não-conformidade relacionada ao componente substituído ou
reparado. Documenta serviços concluídos, fornecendo registros para
liberação das fases de processos e do programa de qualidade aos
clientes. Emite documento para liberação final do produto. Conserva o
local de trabalho limpo e organizado. Descarta resíduos seguindo normas
ambientais. Trabalha com segurança, utilizando Equipamentos de Proteção
Individual (EPI) e Equipamentos de Proteção Coletiva (EPC). Participa de
programas de prevenção de acidentes, colaborando nas ações e propondo
melhorias nas condições de segurança.

\begin{Form}
\textbf{Esta correspondência está adequada?} \\ 
\ChoiceMenu[radio,radiosymbol=\ding{52},bordercolor=0 0.5 0,name=cbovalid_ 314310 ]{}{CBO deve ficar=sim, \hspace{6em} CBO não fica aqui=nao}
\end{Form}

\textbf{CBO: 914105  ( Mecânico de manutenção de aeronaves, em geral )}

Planeja o trabalho, verificando ordens de serviço. Interpreta manuais,
desenhos técnicos mecânicos e elétricos de motopropulsores, estruturas,
sistemas diversificados da aeronave e componentes aviônicos, de acordo
com o tipo -- aviões e helicópteros, de grande, médio e pequeno porte --
de aeronave. Analisa ``diretivas de aeronavegabilidade'' -- notificações
ou avisos de que existem deficiências de segurança conhecidas,
associadas às aeronaves, que podem estar relacionadas à manutenção --,
para cumprir diretrizes de segurança aeronáutica. Utiliza sistemas de
informação de apoio à manutenção. Identifica novos materiais utilizados
na construção de peças de motores e de estruturas. Repara motores
convencionais e motores a reação (a jato), removendo-os, corrigindo
defeitos, instalando acessórios, executando ensaios não destrutivos e
realizando revisão geral. Balanceia e instala motores. Coleta amostras
de fluidos para análise laboratorial. Pode reparar aeronaves com
motorização elétrica. Repara estruturas de aeronaves, inspecionando-as a
fim de detectar defeitos e verificando fixação de rebites. Repara
chapeamento, aplica materiais compósitos e trata estruturas contra
corrosão. Alinha superfícies em relação à fuselagem. Repara sistema de
combustível, examinando tipo de combustível, verificando nível e pressão
do sistema e realizando destanqueamento (retirada de combustível dos
tanques). Remove componentes; repara tanques; descontaminha o sistema;
substitui componentes; e instala tanques. Realiza revisão geral de
componentes e abastece os tanques. Realiza manutenção do sistema de trem
de pouso, removendo rodas, pneus, trens de pouso, amortecedores,
atuadores e sistema de freios; revisando sistema de emergência e sistema
de freio automático; montando e ajustando componentes. Lubrifica
componentes mecânicos e realiza testes operacionais em componentes do
trem de pouso. Executa manutenção de sistemas elétricos e eletrônicos,
verificando acionamento elétrico dos diversos sistemas; reparando e
instalando componentes. Faz manutenção do sistema de comandos de voo,
inspecionando funcionamento de componentes, removendo comandos e
balanceando superfícies de comando. Regula e instala comandos de voo.
Realiza manutenção do sistema interior de aeronaves, substituindo
componentes de poltronas e revisando estado geral de poltronas e de
cintos de segurança. Repara componentes do sistema da copa (``galley'').
Conserta sistemas de som, vídeo, iluminação e comunicação. Substitui
pisos e trilhos e repara componentes das toaletes. Executa manutenção em
sistemas diversificados, cumprindo rotinas de inspeção de manutenção e
executando serviços de substituição de peças ou equipamentos. Substitui
peças dos sistemas de oxigênio; ar-condicionado; pressurização; água
potável; e de detritos. Substitui equipamentos de emergência. Substitui
peças do sistema antichuva e antigelo, bem como do sistema de detecção e
extinção de fogo. Pode realizar manutenção de sistemas hidráulicos.
Repara sistemas de hélices e de rotores de helicópteros, removendo-os e
substituindo engrenagens, selos mecânicos e rolamentos. Balanceia e
monta os sistemas. Ajusta conjunto de sincronismo de rotores. Faz testes
de componentes da aeronave, de acordo com procedimentos específicos, a
fim de comprovar seu funcionamento. Testa funcionamento de motores; de
sistemas de alimentação e geração de energia elétrica; e de sistema
antiderrapagem. Realiza testes de freio do rotor de helicópteros e
testes de funcionamento de rotores e hélices de helicópteros. Preenche
relatórios de avarias e fichas de inspeção. Mantém o local de trabalho
limpo e organizado. Descarta resíduos seguindo normas ambientais.
Trabalha com segurança, aplicando conceitos de ergonomia e utilizando
Equipamentos de Proteção Individual (EPI). Sinaliza áreas de manutenção
e delimita áreas de risco. Opera equipamentos de extinção de incêndio.
Orienta balizamento de aeronaves. Realiza aterramento elétrico da
aeronave em manutenção.

\begin{Form}
\textbf{Esta correspondência está adequada?} \\ 
\ChoiceMenu[radio,radiosymbol=\ding{52},bordercolor=0 0.5 0,name=cbovalid_ 914105 ]{}{CBO deve ficar=sim, \hspace{6em} CBO não fica aqui=nao}
\end{Form}

\newpage
\section{ 02011 -- Técnico em Manutenção Aeronáutica em Célula }

\textbf{Perfil do Curso (CNCT)}

O Técnico em Manutenção Aeronática em Célula será habilitado para:

Programar, controlar e executar manutenção preventiva e corretiva de
aeronaves. Aplicar procedimentos de manuais de fabricantes, publicações
técnicas e normas nacionais e internacionais do setor aeronáutico.
Identificar a sequência adequada de atividades na desmontagem e montagem
de aeronaves atendendo às normas e aos padrões técnicos de qualidade,
saúde e segurança e de meio ambiente. Diagnosticar as condições de
operação das diferentes partes da aeronave. Realizar inspeção visual e
fazer testes de comissionamento e de performance em equipamentos que
mantêm a célula das aeronaves em condições de disponibilidade para o
voo. Coordenar tarefas de limpeza, lubrificação, pequenos reparos,
desmontagem, montagem, substituição, testagem e regulagem de peças,
equipamentos e sistemas. Reparar estruturas de aeronaves atendendo às
normas e aos padrões técnicos de qualidade, saúde e segurança e de meio
ambiente. Realizar manutenção em sistemas de trem de pouso, hidráulicos
e pneumáticos, comando de voo e interiores de aeronaves.

Para atuação como Técnico em Manutenção Aeronática em Célula, são
fundamentais:

Conhecimentos e saberes relacionados aos processos de planejamento e
manutenção de aeronaves de modo a assegurar a saúde e a segurança dos
trabalhadores e dos usuários. Conhecimentos e saberes relacionados à
sustentabilidade do processo produtivo. Conhecimentos e saberes
relacionados às técnicas e processos de manutenção na aviação, às normas
técnicas, à constante atenção às atualizações presentes no setor de
aviação, à liderança de equipes, à solução de problemas técnicos e
trabalhistas e à gestão de conflitos.

\textbf{CBOs correspondidos e seus perfis (presentes no CNCT, e presentes na predição da IA)}

\textbf{CBO: 314310  ( Técnico mecânico (aeronaves) )}

Prepara ensaios de fabricação de equipamentos, estruturas e sistemas
mecânicos de aviões e helicópteros, verificando cronograma do processo
de fabricação e interpretando desenhos de projeto, para elaborar croquis
e fabricar ou adaptar peças, ferramentas e dispositivos a serem
utilizados nos ensaios. Realiza ensaios de fabricação de equipamentos,
estruturas e sistemas mecânicos de aviões e helicópteros, abrangendo
fuselagem, montantes, naceles, capotas de motor, carenagens, rotores e
outras superfícies aerodinâmicas de célula, trens de aterrissagem,
sistemas de freio e demais equipamentos mecânicos da aeronave, incluindo
seus acessórios e controles. Identifica necessidade de ajustes
dimensionais, adapta peças e acessórios e torna modificações adequadas
às características da aeronave. Participa dos testes finais de ensaios e
montagem. Prepara a montagem de componentes em aviões e helicópteros,
interpretando desenhos técnicos, verificando a disponibilidade de itens
no estoque para as montagens requeridas e analisando condições do local
de instalação e dos componentes, para realizar o preparo do local.
Executa, seguindo as instruções de trabalho, a montagem de componentes
em bancada e os instala em sistemas da aeronave de acordo com o projeto,
utilizando ferramentas, instrumentos e equipamentos. Elimina
interferências prejudiciais à segurança, assegurando ajustes entre peças
e destas com a estrutura. Realiza testes funcionais de aviões e
helicópteros, cumprindo procedimentos, seguindo parâmetros estabelecidos
pelo fabricante e utilizando equipamentos específicos e calibrados.
Executa testes iniciais para confirmação dos problemas apresentados,
testa sistemas operacionais mecânicos, de estrutura da aeronave e de
aviônica aplicada ao controle de equipamentos e sistemas mecânicos.
Simula processos, realiza análise dimensional por etapas de montagem e
participa da execução de provas de inclinação da aeronave. Libera
processos para itens de segurança. Prepara a execução de manutenção de
estruturas, equipamentos e sistemas mecânicos de aviões e helicópteros,
identificando sua forma de utilização, seus defeitos e suas falhas.
Executa os reparos, utilizando ferramentas e equipamentos apropriados,
substituindo e regulando componentes e utilizando técnicas de proteção
química e de acabamento de superfície. Nivela a linha de eixo do motor
da aeronave. Presta assessoria técnica interna e externa, treinando
equipes de trabalho, acompanhando o processo de montagem e manutenção,
provendo suporte técnico a equipes de trabalho e clientes, realizando
atendimento de emergência e fazendo visita técnica. Analisa resultados
da assessoria prestada, com a finalidade de assegurar a satisfação dos
clientes. Participa do programa de qualidade da empresa, seguindo as
normas internas e as definidas por organismos externos competentes,
utilizando ferramentas da qualidade, seguindo procedimentos de serviço e
identificando pontos críticos para propor melhorias, tendo em vista o
alcance de metas de qualidade. Fornece informações técnicas para
execução do controle estatístico do processo. Elabora documentação
técnica. Faz orçamento de serviços, com estimativa de tempo e de custo.
Emite parecer técnico. Elabora plano de ação referente às
não-conformidades detectadas em produtos e processos. Elabora relatório
técnico e de campo, identificando os componentes instalados e
registrando a não-conformidade relacionada ao componente substituído ou
reparado. Documenta serviços concluídos, fornecendo registros para
liberação das fases de processos e do programa de qualidade aos
clientes. Emite documento para liberação final do produto. Conserva o
local de trabalho limpo e organizado. Descarta resíduos seguindo normas
ambientais. Trabalha com segurança, utilizando Equipamentos de Proteção
Individual (EPI) e Equipamentos de Proteção Coletiva (EPC). Participa de
programas de prevenção de acidentes, colaborando nas ações e propondo
melhorias nas condições de segurança.

\begin{Form}
\textbf{Esta correspondência está adequada?} \\ 
\ChoiceMenu[radio,radiosymbol=\ding{52},bordercolor=0 0.5 0,name=cbovalid_ 314310 ]{}{CBO deve ficar=sim, \hspace{6em} CBO não fica aqui=nao}
\end{Form}

\textbf{CBO: 914105  ( Mecânico de manutenção de aeronaves, em geral )}

Planeja o trabalho, verificando ordens de serviço. Interpreta manuais,
desenhos técnicos mecânicos e elétricos de motopropulsores, estruturas,
sistemas diversificados da aeronave e componentes aviônicos, de acordo
com o tipo -- aviões e helicópteros, de grande, médio e pequeno porte --
de aeronave. Analisa ``diretivas de aeronavegabilidade'' -- notificações
ou avisos de que existem deficiências de segurança conhecidas,
associadas às aeronaves, que podem estar relacionadas à manutenção --,
para cumprir diretrizes de segurança aeronáutica. Utiliza sistemas de
informação de apoio à manutenção. Identifica novos materiais utilizados
na construção de peças de motores e de estruturas. Repara motores
convencionais e motores a reação (a jato), removendo-os, corrigindo
defeitos, instalando acessórios, executando ensaios não destrutivos e
realizando revisão geral. Balanceia e instala motores. Coleta amostras
de fluidos para análise laboratorial. Pode reparar aeronaves com
motorização elétrica. Repara estruturas de aeronaves, inspecionando-as a
fim de detectar defeitos e verificando fixação de rebites. Repara
chapeamento, aplica materiais compósitos e trata estruturas contra
corrosão. Alinha superfícies em relação à fuselagem. Repara sistema de
combustível, examinando tipo de combustível, verificando nível e pressão
do sistema e realizando destanqueamento (retirada de combustível dos
tanques). Remove componentes; repara tanques; descontaminha o sistema;
substitui componentes; e instala tanques. Realiza revisão geral de
componentes e abastece os tanques. Realiza manutenção do sistema de trem
de pouso, removendo rodas, pneus, trens de pouso, amortecedores,
atuadores e sistema de freios; revisando sistema de emergência e sistema
de freio automático; montando e ajustando componentes. Lubrifica
componentes mecânicos e realiza testes operacionais em componentes do
trem de pouso. Executa manutenção de sistemas elétricos e eletrônicos,
verificando acionamento elétrico dos diversos sistemas; reparando e
instalando componentes. Faz manutenção do sistema de comandos de voo,
inspecionando funcionamento de componentes, removendo comandos e
balanceando superfícies de comando. Regula e instala comandos de voo.
Realiza manutenção do sistema interior de aeronaves, substituindo
componentes de poltronas e revisando estado geral de poltronas e de
cintos de segurança. Repara componentes do sistema da copa (``galley'').
Conserta sistemas de som, vídeo, iluminação e comunicação. Substitui
pisos e trilhos e repara componentes das toaletes. Executa manutenção em
sistemas diversificados, cumprindo rotinas de inspeção de manutenção e
executando serviços de substituição de peças ou equipamentos. Substitui
peças dos sistemas de oxigênio; ar-condicionado; pressurização; água
potável; e de detritos. Substitui equipamentos de emergência. Substitui
peças do sistema antichuva e antigelo, bem como do sistema de detecção e
extinção de fogo. Pode realizar manutenção de sistemas hidráulicos.
Repara sistemas de hélices e de rotores de helicópteros, removendo-os e
substituindo engrenagens, selos mecânicos e rolamentos. Balanceia e
monta os sistemas. Ajusta conjunto de sincronismo de rotores. Faz testes
de componentes da aeronave, de acordo com procedimentos específicos, a
fim de comprovar seu funcionamento. Testa funcionamento de motores; de
sistemas de alimentação e geração de energia elétrica; e de sistema
antiderrapagem. Realiza testes de freio do rotor de helicópteros e
testes de funcionamento de rotores e hélices de helicópteros. Preenche
relatórios de avarias e fichas de inspeção. Mantém o local de trabalho
limpo e organizado. Descarta resíduos seguindo normas ambientais.
Trabalha com segurança, aplicando conceitos de ergonomia e utilizando
Equipamentos de Proteção Individual (EPI). Sinaliza áreas de manutenção
e delimita áreas de risco. Opera equipamentos de extinção de incêndio.
Orienta balizamento de aeronaves. Realiza aterramento elétrico da
aeronave em manutenção.

\begin{Form}
\textbf{Esta correspondência está adequada?} \\ 
\ChoiceMenu[radio,radiosymbol=\ding{52},bordercolor=0 0.5 0,name=cbovalid_ 914105 ]{}{CBO deve ficar=sim, \hspace{6em} CBO não fica aqui=nao}
\end{Form}

\textbf{CBO: 914110  ( Mecânico de manutenção de sistema hidráulico de aeronaves (serviços de pista e hangar) )}

Planeja o trabalho, verificando ordens de serviço. Interpreta manuais e
desenhos técnicos mecânicos e elétricos de sistemas hidráulicos da
aeronave e de componentes elétricos e eletrônicos diretamente associados
aos sistemas hidráulicos, de acordo com o tipo -- aviões e helicópteros,
de grande, médio e pequeno porte -- de aeronave. Analisa ``diretivas de
aeronavegabilidade'' -- notificações ou avisos de que existem
deficiências de segurança conhecidas, associadas às aeronaves, que podem
estar relacionadas à manutenção --, a fim de cumprir diretrizes de
segurança aeronáutica. Utiliza sistemas de informação de apoio à
manutenção. Realiza manutenção de sistemas hidráulicos, na pista ou no
hangar, executando testes e examinando uniões, juntas e outros
componentes do sistema -- utilizando instrumentos ou por observação
direta --, para identificação de defeitos. Remove o componente
defeituoso, isolando o circuito hidráulico, utilizando ferramentas
comuns e específicas, para realizar sua manutenção na oficina do hangar.
Substitui mangueiras e tubulações. Elimina vazamentos, variação de
pressão e nível de óleo, entre outros defeitos. Instala componentes --
revisados ou novos --, empregando ferramentas apropriadas e fluidos
hidráulicos, para restabelecer o pleno funcionamento do sistema
hidráulico. Monitora níveis de fluidos hidráulicos, de lubrificantes e
de pressão do sistema. Procede à renovação do fluido hidráulico do
sistema e elimina o ar existente na canalização, utilizando chaves e
equipamentos especiais, para devolver as condições exigidas ao sistema.
Substitui componentes que tenham atingido o tempo-limite, usando
ferramentas convenientes e fluido hidráulico, para garantir a segurança
de operação das aeronaves. Revisa funcionamento de bombas, atuadores,
reservatórios, válvulas, filtros e outros componentes. Ajusta
componentes mecânicos dos sistemas hidráulicos. Realiza testes de
componentes de sistemas hidráulicos da aeronave, de acordo com
procedimentos específicos, a fim de comprovar que o seu funcionamento
atende às especificações técnicas. Testa funcionamento de bombas,
atuadores, válvulas e outros componentes. Comprova o funcionamento do
sistema revisado, observando indicações de instrumentos de painéis, para
certificar-se da exatidão dos serviços concluídos. Pode realizar
manutenção de sistemas elétricos e eletrônicos diretamente associados
aos sistemas hidráulicos da aeronave, verificando acionamento
eletro-hidráulico; reparando e instalando componentes elétricos e
eletrônicos; e interpretando leitura de instrumentos de medidas
elétricas e eletrônicas. Comunica à supervisão eventuais anormalidades
encontradas em outras partes da aeronave - oralmente ou por escrito - a
fim de colaborar com outros serviços de manutenção aeronáutica. Elabora
documentos, preenchendo relatórios de avarias e fichas de inspeção.
Mantém o local de trabalho limpo e organizado. Descarta resíduos
seguindo normas ambientais. Trabalha com segurança, utilizando
Equipamentos de Proteção Individual (EPI), aplicando conceitos de
ergonomia nas atividades de manutenção, sinalizando áreas de manutenção
e delimitando áreas de risco para preservar a integridade física das
pessoas. Opera equipamentos de extinção de incêndio. Orienta balizamento
de aeronaves. Realiza aterramento elétrico de aeronave em manutenção.

\begin{Form}
\textbf{Esta correspondência está adequada?} \\ 
\ChoiceMenu[radio,radiosymbol=\ding{52},bordercolor=0 0.5 0,name=cbovalid_ 914110 ]{}{CBO deve ficar=sim, \hspace{6em} CBO não fica aqui=nao}
\end{Form}

\newpage
\section{ 02012 -- Técnico em Manutenção Aeronáutica em Grupo Motopropulsor }

\textbf{Perfil do Curso (CNCT)}

O Técnico em Manutenção de Aeronáutica em Grupo Motopropulsor será
habilitado para:

Programar, controlar e executar manutenção preventiva e corretiva em
motopropulsor de aeronaves. Aplicar procedimentos de manuais de
fabricantes, publicações técnicas e normas nacionais e internacionais do
setor aeronáutico. Identificar a sequência adequada de atividades na
desmontagem e montagem de motopropulsor de aeronaves atendendo às normas
e aos padrões técnicos de qualidade, saúde e segurança e de meio
ambiente. Diagnosticar as condições de operação das diferentes partes do
motopropulsor de aeronaves. Realizar inspeção visual e fazer testes de
comissionamento e de performance em motopropulsor que mantém as
aeronaves em condições de disponibilidade para o voo. Coordenar tarefas
de limpeza, lubrificação, pequenos reparos, desmontagem, montagem,
substituição, testagem e regulagem de peças, equipamentos e sistemas de
motopropulsores. Reparar estruturas de aeronaves atendendo às normas e
aos padrões técnicos de qualidade, saúde e segurança e de meio ambiente.
Realizar manutenção em sistemas de trem de pouso, hidráulicos e
pneumáticos, comando de voo e interiores de aeronaves. Realizar inspeção
visual do grupo motopropulsor visando ao reconhecimento de anomalias.
Reparar motores convencionais e à reação de acordo com a ordem de
serviço.

Para atuação como Técnico em Manutenção de Aeronáutica em Grupo
Motopropulsor, são fundamentais:

Conhecimentos e saberes relacionados aos processos de planejamento e
manutenção de motopropulsores de aeronaves de modo a assegurar a saúde e
a segurança dos trabalhadores e dos usuários. Conhecimentos e saberes
relacionados à sustentabilidade do processo produtivo, às técnicas e aos
processos de manutenção na aviação, às normas técnicas, à constante
atenção às atualizações presentes no setor de aviação, à liderança de
equipes, à solução de problemas técnicos e trabalhistas e à gestão de
conflitos.

\textbf{CBOs correspondidos e seus perfis (presentes no CNCT, e presentes na predição da IA)}

\textbf{CBO: 314310  ( Técnico mecânico (aeronaves) )}

Prepara ensaios de fabricação de equipamentos, estruturas e sistemas
mecânicos de aviões e helicópteros, verificando cronograma do processo
de fabricação e interpretando desenhos de projeto, para elaborar croquis
e fabricar ou adaptar peças, ferramentas e dispositivos a serem
utilizados nos ensaios. Realiza ensaios de fabricação de equipamentos,
estruturas e sistemas mecânicos de aviões e helicópteros, abrangendo
fuselagem, montantes, naceles, capotas de motor, carenagens, rotores e
outras superfícies aerodinâmicas de célula, trens de aterrissagem,
sistemas de freio e demais equipamentos mecânicos da aeronave, incluindo
seus acessórios e controles. Identifica necessidade de ajustes
dimensionais, adapta peças e acessórios e torna modificações adequadas
às características da aeronave. Participa dos testes finais de ensaios e
montagem. Prepara a montagem de componentes em aviões e helicópteros,
interpretando desenhos técnicos, verificando a disponibilidade de itens
no estoque para as montagens requeridas e analisando condições do local
de instalação e dos componentes, para realizar o preparo do local.
Executa, seguindo as instruções de trabalho, a montagem de componentes
em bancada e os instala em sistemas da aeronave de acordo com o projeto,
utilizando ferramentas, instrumentos e equipamentos. Elimina
interferências prejudiciais à segurança, assegurando ajustes entre peças
e destas com a estrutura. Realiza testes funcionais de aviões e
helicópteros, cumprindo procedimentos, seguindo parâmetros estabelecidos
pelo fabricante e utilizando equipamentos específicos e calibrados.
Executa testes iniciais para confirmação dos problemas apresentados,
testa sistemas operacionais mecânicos, de estrutura da aeronave e de
aviônica aplicada ao controle de equipamentos e sistemas mecânicos.
Simula processos, realiza análise dimensional por etapas de montagem e
participa da execução de provas de inclinação da aeronave. Libera
processos para itens de segurança. Prepara a execução de manutenção de
estruturas, equipamentos e sistemas mecânicos de aviões e helicópteros,
identificando sua forma de utilização, seus defeitos e suas falhas.
Executa os reparos, utilizando ferramentas e equipamentos apropriados,
substituindo e regulando componentes e utilizando técnicas de proteção
química e de acabamento de superfície. Nivela a linha de eixo do motor
da aeronave. Presta assessoria técnica interna e externa, treinando
equipes de trabalho, acompanhando o processo de montagem e manutenção,
provendo suporte técnico a equipes de trabalho e clientes, realizando
atendimento de emergência e fazendo visita técnica. Analisa resultados
da assessoria prestada, com a finalidade de assegurar a satisfação dos
clientes. Participa do programa de qualidade da empresa, seguindo as
normas internas e as definidas por organismos externos competentes,
utilizando ferramentas da qualidade, seguindo procedimentos de serviço e
identificando pontos críticos para propor melhorias, tendo em vista o
alcance de metas de qualidade. Fornece informações técnicas para
execução do controle estatístico do processo. Elabora documentação
técnica. Faz orçamento de serviços, com estimativa de tempo e de custo.
Emite parecer técnico. Elabora plano de ação referente às
não-conformidades detectadas em produtos e processos. Elabora relatório
técnico e de campo, identificando os componentes instalados e
registrando a não-conformidade relacionada ao componente substituído ou
reparado. Documenta serviços concluídos, fornecendo registros para
liberação das fases de processos e do programa de qualidade aos
clientes. Emite documento para liberação final do produto. Conserva o
local de trabalho limpo e organizado. Descarta resíduos seguindo normas
ambientais. Trabalha com segurança, utilizando Equipamentos de Proteção
Individual (EPI) e Equipamentos de Proteção Coletiva (EPC). Participa de
programas de prevenção de acidentes, colaborando nas ações e propondo
melhorias nas condições de segurança.

\begin{Form}
\textbf{Esta correspondência está adequada?} \\ 
\ChoiceMenu[radio,radiosymbol=\ding{52},bordercolor=0 0.5 0,name=cbovalid_ 314310 ]{}{CBO deve ficar=sim, \hspace{6em} CBO não fica aqui=nao}
\end{Form}

\textbf{CBO: 914105  ( Mecânico de manutenção de aeronaves, em geral )}

Planeja o trabalho, verificando ordens de serviço. Interpreta manuais,
desenhos técnicos mecânicos e elétricos de motopropulsores, estruturas,
sistemas diversificados da aeronave e componentes aviônicos, de acordo
com o tipo -- aviões e helicópteros, de grande, médio e pequeno porte --
de aeronave. Analisa ``diretivas de aeronavegabilidade'' -- notificações
ou avisos de que existem deficiências de segurança conhecidas,
associadas às aeronaves, que podem estar relacionadas à manutenção --,
para cumprir diretrizes de segurança aeronáutica. Utiliza sistemas de
informação de apoio à manutenção. Identifica novos materiais utilizados
na construção de peças de motores e de estruturas. Repara motores
convencionais e motores a reação (a jato), removendo-os, corrigindo
defeitos, instalando acessórios, executando ensaios não destrutivos e
realizando revisão geral. Balanceia e instala motores. Coleta amostras
de fluidos para análise laboratorial. Pode reparar aeronaves com
motorização elétrica. Repara estruturas de aeronaves, inspecionando-as a
fim de detectar defeitos e verificando fixação de rebites. Repara
chapeamento, aplica materiais compósitos e trata estruturas contra
corrosão. Alinha superfícies em relação à fuselagem. Repara sistema de
combustível, examinando tipo de combustível, verificando nível e pressão
do sistema e realizando destanqueamento (retirada de combustível dos
tanques). Remove componentes; repara tanques; descontaminha o sistema;
substitui componentes; e instala tanques. Realiza revisão geral de
componentes e abastece os tanques. Realiza manutenção do sistema de trem
de pouso, removendo rodas, pneus, trens de pouso, amortecedores,
atuadores e sistema de freios; revisando sistema de emergência e sistema
de freio automático; montando e ajustando componentes. Lubrifica
componentes mecânicos e realiza testes operacionais em componentes do
trem de pouso. Executa manutenção de sistemas elétricos e eletrônicos,
verificando acionamento elétrico dos diversos sistemas; reparando e
instalando componentes. Faz manutenção do sistema de comandos de voo,
inspecionando funcionamento de componentes, removendo comandos e
balanceando superfícies de comando. Regula e instala comandos de voo.
Realiza manutenção do sistema interior de aeronaves, substituindo
componentes de poltronas e revisando estado geral de poltronas e de
cintos de segurança. Repara componentes do sistema da copa (``galley'').
Conserta sistemas de som, vídeo, iluminação e comunicação. Substitui
pisos e trilhos e repara componentes das toaletes. Executa manutenção em
sistemas diversificados, cumprindo rotinas de inspeção de manutenção e
executando serviços de substituição de peças ou equipamentos. Substitui
peças dos sistemas de oxigênio; ar-condicionado; pressurização; água
potável; e de detritos. Substitui equipamentos de emergência. Substitui
peças do sistema antichuva e antigelo, bem como do sistema de detecção e
extinção de fogo. Pode realizar manutenção de sistemas hidráulicos.
Repara sistemas de hélices e de rotores de helicópteros, removendo-os e
substituindo engrenagens, selos mecânicos e rolamentos. Balanceia e
monta os sistemas. Ajusta conjunto de sincronismo de rotores. Faz testes
de componentes da aeronave, de acordo com procedimentos específicos, a
fim de comprovar seu funcionamento. Testa funcionamento de motores; de
sistemas de alimentação e geração de energia elétrica; e de sistema
antiderrapagem. Realiza testes de freio do rotor de helicópteros e
testes de funcionamento de rotores e hélices de helicópteros. Preenche
relatórios de avarias e fichas de inspeção. Mantém o local de trabalho
limpo e organizado. Descarta resíduos seguindo normas ambientais.
Trabalha com segurança, aplicando conceitos de ergonomia e utilizando
Equipamentos de Proteção Individual (EPI). Sinaliza áreas de manutenção
e delimita áreas de risco. Opera equipamentos de extinção de incêndio.
Orienta balizamento de aeronaves. Realiza aterramento elétrico da
aeronave em manutenção.

\begin{Form}
\textbf{Esta correspondência está adequada?} \\ 
\ChoiceMenu[radio,radiosymbol=\ding{52},bordercolor=0 0.5 0,name=cbovalid_ 914105 ]{}{CBO deve ficar=sim, \hspace{6em} CBO não fica aqui=nao}
\end{Form}

\newpage
\section{ 02013 -- Técnico em Manutenção Automotiva }

\textbf{Perfil do Curso (CNCT)}

O Técnico em Manutenção Automotiva será habilitado para:

Programar, controlar e executar planos de manutenção preventiva em
veículos automotores seguindo as normas técnicas dos respectivos
fabricantes. Executar manutenção preventiva e corretiva de acordo com
diagnósticos em sistemas elétricos e mecânicos em veículos automotores
com ciclo otto e/ou diesel por meio de ferramentas e instrumentos de
medição, atendendo às normas e aos padrões técnicos de qualidade, saúde
e segurança e de meio ambiente. Controlar a emissão de gases poluentes e
reparar defeitos eletrônicos como uso de dispositivos de teste e/ou
scanners. Identificar a conformidade de documentações legais que
permitam que o veículo esteja apto a ser utilizado em vias públicas.
Reconhecer tecnologias inovadoras presentes no segmento, tais como
veículos elétricos e híbridos. Reconhecer tecnologias inovadoras
presentes no segmento visando a atender às transformações digitais na
sociedade.

Para atuação como Técnico em Manutenção Automotiva, são fundamentais:

Conhecimentos e saberes relacionados aos processos de manutenção de
veículos automotivos de modo a assegurar a saúde e a segurança dos
trabalhadores e dos usuários. Conhecimentos e saberes relacionados à
sustentabilidade dos processos, às normas técnicas, à liderança de
equipes, à solução de problemas técnicos e trabalhistas e à gestão de
conflitos.

\textbf{CBOs correspondidos e seus perfis (presentes no CNCT, e presentes na predição da IA)}

\textbf{CBO: 314305  ( Técnico em automobilística )}

Prepara ensaios de fabricação de veículos automotores terrestres,
verificando cronograma do processo de fabricação e interpretando
desenhos de projeto, para elaborar croquis e fabricar ou adaptar peças,
ferramentas e dispositivos a serem utilizados nos ensaios. Realiza
ensaios de fabricação de veículos automotores terrestres, identificando
necessidade de ajustes dimensionais. Adapta peças e acessórios e torna
modificações adequadas às características do veículo. Participa de
testes finais de ensaios de montagem. Prepara a montagem de componentes
em veículos automotores, interpretando desenhos técnicos, verificando a
disponibilidade de itens no estoque para as montagens requeridas e
analisando as condições do local de instalação e dos componentes, para
realizar o preparo do local e possíveis ajustes na estrutura em função
de alinhamento. Executa, seguindo instruções de trabalho, a montagem de
componentes em bancada e os instala em sistemas veiculares de acordo com
o projeto, utilizando ferramentas, instrumentos e equipamentos de
montagem e instalação. Elimina interferências prejudiciais à segurança
do veículo, assegurando ajustes entre peças e destas com a estrutura.
Pode montar componentes em veículos nas linhas de montagem flexíveis e
com utilização de robôs. Realiza testes funcionais dos veículos
automotores, cumprindo procedimentos, seguindo parâmetros estabelecidos
pelo fabricante e utilizando equipamentos específicos e calibrados.
Executa testes iniciais para confirmação dos problemas apresentados e
testa sistemas operacionais elétricos, eletrônicos, mecânicos e de
estrutura veicular. Simula processos, realiza análise dimensional por
etapas de montagem e libera processos para itens de segurança. Prepara a
manutenção de estruturas e sistemas mecânicos e eletroeletrônicos de
veículos automotores terrestres, identificando a forma de utilização do
veículo, seus defeitos e suas falhas. Executa reparos estruturais e
relacionados aos componentes mecânicos e eletroeletrônicos - incluindo
componentes de sistemas de eletrônica embarcada e típicos de carros
elétricos -, usando ferramentas e equipamentos apropriados, substituindo
e regulando componentes e utilizando técnicas de proteção química e de
acabamento de superfície. Presta assessoria técnica interna e externa,
treinando equipes de trabalho, acompanhando o processo de montagem e
manutenção e realizando atendimento de emergência. Provê suporte técnico
a equipes de trabalho e a clientes, realizando visita técnica. Analisa
resultados da assessoria prestada, com a finalidade de assegurar a
satisfação dos clientes. Participa do programa de qualidade da empresa,
seguindo normas internas e definidas por organismos externos
competentes, utilizando ferramentas da qualidade, seguindo procedimentos
de serviço e identificando pontos críticos para propor melhorias, tendo
em vista o alcance de metas de qualidade gerais da empresa e das
operações ao longo do processo. Fornece informações técnicas para
execução do controle estatístico do processo e participa de auditorias
internas e externas. Elabora documentação técnica. Faz orçamento de
serviços, com estimativa de tempo e de custo. Elabora plano de ação
referente às não-conformidades detectadas em produtos e processos.
Elabora parecer técnico e relatório técnico e de campo, registrando
não-conformidade relacionada ao componente substituído ou reparado.
Registra a identificação dos componentes instalados. Documenta serviços
concluídos, fornecendo registros para liberação das fases de processos e
do programa de qualidade aos clientes. Emite documento para liberação
final do produto. Conserva o local de trabalho limpo e organizado.
Descarta resíduos seguindo normas ambientais. Trabalha com segurança,
utilizando Equipamentos de Proteção Individual (EPI) e Equipamentos de
Proteção Coletiva (EPC). Participa de programas de prevenção de
acidentes, colaborando nas ações e propondo melhorias nas condições de
segurança.

\begin{Form}
\textbf{Esta correspondência está adequada?} \\ 
\ChoiceMenu[radio,radiosymbol=\ding{52},bordercolor=0 0.5 0,name=cbovalid_ 314305 ]{}{CBO deve ficar=sim, \hspace{6em} CBO não fica aqui=nao}
\end{Form}

\textbf{CBO: 314405  ( Técnico de manutenção de sistemas e instrumentos )}

Planeja a manutenção mecânica de sistemas e instrumentos industriais,
dimensionando equipe, elaborando cronograma e selecionando ferramentas e
materiais. Estima despesas com a atividade e analisa relação
custo-benefício. Seleciona técnicas de manutenção e elabora indicadores
de performance. Estabelece metodologia de acompanhamento do plano de
manutenção, cria relatório e programa reuniões de acompanhamento.
Executa manutenção mecânica preditiva, preventiva e corretiva de
sistemas e instrumentos industriais. Analisa indicadores de manutenção -
tais como tempo entre falhas, frequência de manutenção e temperatura dos
sistemas e instrumentos, dentre outros -, para avaliar a capabilidade do
processo. Controla os processos de manutenção mecânica de sistemas e
instrumentos industriais, aplicando técnicas para a solução de problemas
em manutenção, analisando as especificações técnicas e verificando o
cumprimento dos procedimentos. Monitora o funcionamento dos sistemas e
instrumentos industriais. Elabora desenho para peças de reposição e
dispositivos. Interage com sistemas de monitoramento em tempo real -
interpretando dados, indicadores e outras informações -, para subsidiar
o controle digital da manutenção de sistemas e instrumentos industriais.
Participa na manutenção aditiva - para eliminar possibilidades de falhas
não perceptíveis de segurança e proteção - e na manutenção proativa -
para eliminar causas primárias de falhas -, melhorando o desempenho e a
vida útil dos sistemas e instrumentos. Realiza testes, ensaios mecânicos
e ensaios especiais - tais como de líquidos penetrantes, raio x,
fractolográficos e macros metalográficos (visuais) -, em sistemas e
instrumentos industriais. Propõe melhorias nas condições de utilização
de sistemas e instrumentos industriais, realizando estudos para redução
de tempos de parada por manutenção, adaptando peças e componentes,
recomendando novos métodos e melhores condições ergonômicas e de
iluminação, dentre outros aspectos intervenientes nas condições de
trabalho. Presta assessoria técnica de manutenção, realizando
demonstração e instalação técnica e negociando serviços junto ao
cliente. Participa na elaboração de procedimentos para aplicação do
sistema de qualidade. Pode realizar calibração mecânica de instrumentos
e equipamentos industriais em laboratórios credenciados, seguindo
metodologias específicas, aplicando métodos de medição conforme normas,
especificando equipamentos para realização da calibração e definindo
critérios de aprovação da calibração. Coleta e analisa dados para
emissão de certificado de calibração. Elabora procedimentos técnicos e
administrativos, interpretando normas técnicas; emitindo instruções de
testes e ensaios; e descrevendo procedimentos para manutenção e para
funcionamento de equipamentos, instrumentos e sistemas. Prepara
fluxograma do processo de manutenção. Elabora especificação de critérios
para compra de materiais. Acompanha auditorias de processos técnicos e
administrativos. Ministra programas de treinamento para pessoal de
manutenção e para clientes internos e externos, elaborando o material
didático. Conserva o local de trabalho limpo e organizado. Descarta
resíduos seguindo normas ambientais. Aplica técnicas de segurança,
implementando medidas para evitar acidente de trabalho, orientando na
utilização de normas e participando de auditorias e inspeções de
segurança.

\begin{Form}
\textbf{Esta correspondência está adequada?} \\ 
\ChoiceMenu[radio,radiosymbol=\ding{52},bordercolor=0 0.5 0,name=cbovalid_ 314405 ]{}{CBO deve ficar=sim, \hspace{6em} CBO não fica aqui=nao}
\end{Form}

\newpage
\section{ 02014 -- Técnico em Manutenção de Máquinas Industriais }

\textbf{Perfil do Curso (CNCT)}

O Técnico em Manutenção de Máquinas Industriais será habilitado para:

Planejar, controlar e executar atividades relativas à manutenção
mecânica nos níveis preventivos, preditivos e corretivos. Realizar a
confecção de peças e componentes mecânicos para manutenção de máquinas e
equipamentos. Desenvolver projetos dedicados à manutenção mecânica.
Atuar na manutenção dos sistemas automatizados de máquinas e
equipamentos industriais, atendendo às normas e aos padrões técnicos de
qualidade, saúde e segurança e de meio ambiente. Reconhecer tecnologias
inovadoras presentes no segmento visando a atender às transformações
digitais na sociedade.

Para atuação como Técnico em Manutenção de Máquinas Industriais, são
fundamentais:

Conhecimentos e saberes relacionados aos processos de manutenção de
equipamentos mecânicos e construção de modo a assegurar a saúde e a
segurança dos trabalhadores e usuários. Conhecimentos e saberes
relacionados à sustentabilidade do processo produtivo, às normas
técnicas, à liderança de equipes, à solução de problemas técnicos e
trabalhistas, à tomada de decisões e à gestão de conflitos.

\textbf{CBOs correspondidos e seus perfis (presentes no CNCT, e presentes na predição da IA)}

\textbf{CBO: 314410  ( Técnico em manutenção de máquinas )}

Planeja a manutenção mecânica de máquinas industriais, dimensionando
equipe, elaborando cronograma e selecionando ferramentas e materiais.
Estima despesas com a atividade, analisando a relação custo-benefício.
Descreve planos de manutenção, selecionando técnicas e definindo
indicadores de desempenho. Estabelece metodologia de acompanhamento do
plano de manutenção, criando relatório de acompanhamento e programando
reuniões. Executa a manutenção mecânica preditiva, preventiva e
corretiva de máquinas industriais. Analisa indicadores de manutenção -
tais como tempo entre falhas, frequência de manutenção e temperatura das
máquinas, dentre outros -, para avaliar a capabilidade do processo.
Controla os processos de manutenção mecânica de máquinas industriais,
aplicando técnicas para a solução de problemas em manutenção, analisando
as especificações técnicas e verificando o cumprimento dos
procedimentos. Monitora o funcionamento de máquinas industriais. Elabora
desenho para peças de reposição e dispositivos. Interage com sistemas de
monitoramento em tempo real, interpretando dados, indicadores e outras
informações, para subsidiar o controle digital da manutenção de
máquinas. Participa na manutenção aditiva - para eliminar possibilidades
de falhas não perceptíveis de segurança e proteção - e na manutenção
proativa - para eliminar causas primárias de falhas -, melhorando o
desempenho e a vida útil das máquinas. Realiza testes, ensaios mecânicos
e ensaios especiais - tais como de líquidos penetrantes, raio x,
fractolográficos e macros metalográficos (visuais) -, em máquinas
industriais. Propõe melhorias nas condições de utilização de máquinas
industriais, realizando estudos para redução de tempos de parada por
manutenção, adaptando peças e componentes, recomendando novos métodos e
melhores condições ergonômicas e de iluminação, dentre outros aspectos
intervenientes nas condições de trabalho. Presta assessoria técnica de
manutenção, realizando demonstração e instalação técnica e negociando
serviços junto ao cliente. Pode executar reforma (retrofitting) em
máquinas industriais. Participa na elaboração de procedimentos para
aplicação do sistema de qualidade. Elabora procedimentos técnicos e
administrativos, interpretando normas técnicas; emitindo instruções de
testes e ensaios; e descrevendo procedimentos para manutenção e para
funcionamento de máquinas. Prepara fluxograma do processo de manutenção.
Elabora especificação de critérios para compra de materiais. Acompanha
auditorias de processos técnicos e administrativos. Ministra programas
de treinamento para pessoal de manutenção e para clientes internos e
externos, elaborando o material didático. Conserva o local de trabalho
limpo e organizado. Descarta resíduos seguindo normas ambientais. Aplica
técnicas de segurança, implementando medidas para evitar acidente de
trabalho, orientando na utilização de normas e participando de
auditorias e inspeções de segurança.

\begin{Form}
\textbf{Esta correspondência está adequada?} \\ 
\ChoiceMenu[radio,radiosymbol=\ding{52},bordercolor=0 0.5 0,name=cbovalid_ 314410 ]{}{CBO deve ficar=sim, \hspace{6em} CBO não fica aqui=nao}
\end{Form}

\textbf{CBO: 911305  ( Mecânico de manutenção de máquinas, em geral )}

Planeja atividades de manutenção mecânica, analisando históricos de
manutenção das máquinas e equipamentos, verificando dados de sistemas de
controle de redes industriais - locais e a distância - e considerando
relato do operador sobre aspectos operacionais como vibração, ruído e
temperatura dos elementos e do conjunto. Verifica disponibilidade de
material e ferramental, estima tempo do serviço, define cronograma e
estabelece a quantidade de mão de obra e a composição da equipe. Executa
serviços de manutenção mecânica de equipamentos, interpretando manuais e
desenhos técnicos. Realiza ensaios físicos e mecânicos nas máquinas e em
seus componentes, para detectar possíveis falhas. Encaminha amostras
para ensaios laboratoriais e analisa os resultados. Desmonta e remonta
conjuntos mecânicos, medindo seus parâmetros dimensionais. Substitui ou
recondiciona peças. Calibra, alinha e nivela equipamentos, máquinas e
conjuntos mecânicos. Testa o funcionamento das máquinas e dos
equipamentos. Por meio de inspeção em períodos regulares, avalia
condições de funcionamento e desempenho de componentes, equipamentos e
máquinas industriais. Lubrifica máquinas, componentes e ferramentas,
desobstruindo os condutos para completar nível de lubrificantes, de
acordo com especificações técnicas. Verifica alinhamento e nivelamento
das máquinas. Instrui operadores sobre causas de incidência de quebras e
defeitos, como medida preventiva. Pode executar serviços de modernização
- ``retrofitting'' - de máquinas e equipamentos. Documenta informações
técnicas, elaborando relatórios de ocorrências, emitindo pareceres
técnicos, registrando modificações e melhorias nos equipamentos e
máquinas bem como resultados de testes realizados em peças e materiais.
Mantém registros do tempo despendido por serviço de manutenção. Mantém o
local de trabalho limpo, seguro, organizado e adequadamente iluminado.
Requisita materiais e componentes e especifica materiais e peças para
compra. Em atendimento às normas ambientais, adapta componentes e
acessórios, seleciona peças e materiais para descarte e reciclagem,
reutiliza matéria-prima de peças substituídas e prioriza o uso de
produtos químicos não poluentes. Pode propor e desenvolver, após
aprovação, dispositivos para redução de riscos de acidentes.

\begin{Form}
\textbf{Esta correspondência está adequada?} \\ 
\ChoiceMenu[radio,radiosymbol=\ding{52},bordercolor=0 0.5 0,name=cbovalid_ 911305 ]{}{CBO deve ficar=sim, \hspace{6em} CBO não fica aqui=nao}
\end{Form}

\newpage
\section{ 02015 -- Técnico em Manutenção de Máquinas Navais }

\textbf{Perfil do Curso (CNCT)}

O Técnico em Manutenção de Máquinas Navais será habilitado para:

Planejar, controlar e executar tarefas de manutenção e instalação de
máquinas navais, equipamentos eletro-hidráulicos e de refrigeração,
motores de combustão interna, turbinas a gás e caldeiras de navios.
Elaborar documentação técnica de atividades operacionais. Interpretar
informações de sensores de medidas físicas, térmicas e mecânicas.
Recuperar componentes de motores e de equipamentos navais e comissionar
motores e equipamentos. Utilizar requisitos de sistemas de qualidade e
preservação ambiental.

Para atuação como Técnico em Manutenção de Máquinas Navais, são
fundamentais:

Conhecimentos e saberes relacionados aos processos de planejamento e
manutenção de máquinas navais de modo a assegurar a saúde e a segurança
dos usuários. Conhecimentos e saberes relacionados à sustentabilidade do
processo de recuperação de peças, às técnicas, às normas técnicas, à
liderança de equipes, à solução de problemas técnicos e trabalhistas e à
gestão de conflitos.

\textbf{CBOs correspondidos e seus perfis (presentes no CNCT, e presentes na predição da IA)}

\textbf{CBO: 314315  ( Técnico mecânico (embarcações) )}

Prepara ensaios de fabricação de estruturas, máquinas, equipamentos e
sistemas mecânicos de embarcações, verificando cronograma do processo de
fabricação e interpretando desenhos de projeto, para elaborar croquis e
fabricar ou adaptar peças, ferramentas e dispositivos a serem utilizados
nos ensaios. Realiza ensaios de fabricação de estruturas, máquinas,
equipamentos e sistemas mecânicos de embarcações, identificando
necessidades de ajustes dimensionais. Adapta peças e acessórios e torna
modificações adequadas às características da embarcação. Participa de
testes finais de ensaios e montagem. Prepara a montagem de componentes
em embarcações, interpretando desenhos técnicos, verificando a
disponibilidade de itens no estoque para as montagens requeridas e
analisando condições do local de instalação e dos componentes, para
realizar o preparo do local. Executa, seguindo instruções de trabalho, a
montagem de componentes em bancada e os instala em sistemas e estruturas
da embarcação de acordo com o projeto, utilizando ferramentas,
instrumentos e equipamentos de montagem e instalação. Elimina
interferências prejudiciais à segurança, assegurando ajustes entre peças
e destas com a estrutura. Realiza testes funcionais de embarcações,
cumprindo procedimentos, seguindo parâmetros estabelecidos pelo
fabricante e utilizando equipamentos específicos e calibrados. Executa
testes iniciais para confirmação dos problemas apresentados e realiza
testes de sistemas operacionais mecânicos, de estrutura da embarcação e
de eletrônica embarcada e automação naval, quando associados aos
controles de sistemas mecânicos. Simula processos e realiza análise
dimensional por etapas de montagem. Libera processos para itens de
segurança. Prepara a execução de manutenção de estruturas, máquinas,
equipamentos e sistemas mecânicos de embarcações, identificando sua
forma de utilização, seus defeitos e suas falhas. Executa reparos,
utilizando ferramentas e equipamentos apropriados para substituir e
regular componentes, além de aplicar técnicas de tratamento de
superfícies. Participa do programa de qualidade da empresa, seguindo as
normas internas e as definidas por organismos externos competentes,
utilizando ferramentas da qualidade, seguindo procedimentos de serviço e
identificando pontos críticos para propor melhorias, tendo em vista o
alcance das metas de qualidade. Fornece informações técnicas para
execução do controle estatístico do processo. Presta assessoria técnica
interna e externa, atuando no estaleiro, porto ou offshore; treinando
equipes de trabalho; acompanhando o processo de montagem e manutenção;
realizando atendimento de emergência; provendo suporte técnico a equipes
de trabalho e a clientes; e realizando visita técnica. Analisa
resultados da assessoria prestada, com a finalidade de assegurar a
satisfação dos clientes. Elabora documentação técnica. Faz orçamento de
serviços, com estimativa de tempo e de custo. Elabora parecer técnico.
Emite plano de ação referente às não-conformidades detectadas em
produtos e processos. Emite relatório técnico e de campo, identificando
os componentes instalados e registrando a não-conformidade relacionada
ao componente substituído ou reparado. Documenta serviços concluídos,
fornecendo registros para liberação das fases de processos e do programa
de qualidade aos clientes. Emite documento para liberação final do
produto. Conserva o local de trabalho limpo e organizado. Descarta
resíduos seguindo normas ambientais. Trabalha com segurança, utilizando
Equipamentos de Proteção Individual (EPI) e Equipamentos de Proteção
Coletiva (EPC). Participa de programas de prevenção de acidentes,
colaborando nas ações e propondo melhorias nas condições de segurança.

\begin{Form}
\textbf{Esta correspondência está adequada?} \\ 
\ChoiceMenu[radio,radiosymbol=\ding{52},bordercolor=0 0.5 0,name=cbovalid_ 314315 ]{}{CBO deve ficar=sim, \hspace{6em} CBO não fica aqui=nao}
\end{Form}

\textbf{CBO: 914205  ( Mecânico de manutenção de motores e equipamentos navais )}

Planeja o trabalho de manutenção de motores e equipamentos mecânicos
navais, identificando tipos de motores e de equipamentos e analisando
suas características de funcionamento; examinando manuais de fabricação,
de montagem, de manutenção e catálogos de componentes; e interpretando
desenhos técnicos mecânicos. Seleciona ferramentas para manutenção e
identifica necessidades de fabricação de ferramentas especiais. Estima
recursos humanos para manutenção. Verifica disponibilidade de peças e de
componentes sobressalentes. Estima tempo de execução da manutenção.
Define tipo e periodicidade da manutenção. Utiliza sistemas de
informação de apoio à manutenção. Realiza manutenção em motores -- a
gasolina, a álcool ou a diesel, dentre outros tipos de motores a
combustão -- e equipamentos mecânicos navais -- tais como bombas, eixos,
redutores, sistemas e máquinas auxiliares --, examinando motores e
equipamentos. Desmonta equipamentos e realiza inspeção dimensional em
componentes. Identifica e diagnostica causas de avarias. Substitui,
alinha e regula componentes. Ajusta folgas entre componentes. Monta e
alinha equipamentos. Recupera componentes de motores e equipamentos
mecânicos navais em manutenção, realizando operações de ajustagem em
componentes; identificando necessidade de tratamento de superfície de
peças; medindo peças e componentes; recondicionando peças e componentes;
e identificando necessidade de balanceamento dinâmico das peças
recuperadas. Executa ajustes circunstanciais de emergência. Identifica
necessidade de fabricação de peças sobressalentes, elabora croqui das
peças e faz especificação de materiais para sua produção. Realiza testes
em motores e equipamentos mecânicos navais, de acordo com procedimentos
específicos, verificando parâmetros de tolerância; e analisando
intensidade de vibração de equipamentos rotativos. Faz ajustes de
motores e equipamentos em função de resultados de testes. Verifica
alinhamento e deflexão em equipamentos. Testa potência e temperatura de
motores e equipamentos; testa pressão de óleos lubrificantes e
combustíveis; e realiza testes hidrostáticos. Testa componentes de
eletrônica embarcada e de automação naval -- diretamente relacionados
aos motores e equipamentos mecânicos navais --, aplicando procedimentos
específicos e podendo removê-los, substituí-los ou instalá-los, a fim de
restabelecer o pleno funcionamento das partes em manutenção. Mantém-se
atualizado acerca das inovações tecnológicas da área de mecânica naval,
consultando fontes diversificadas de informação, interpretando e
analisando informações técnicas de motores e equipamentos mecânicos,
para realizar os serviços de manutenção em elevados níveis de qualidade
e eficiência. Elabora documentação técnica, emitindo pareceres técnicos;
registrando históricos de manutenção; especificando materiais e peças
para requisição; e registrando resultados de medições e testes. Elabora
relatórios e orçamentos de serviços e de peças. Conserva o local de
trabalho limpo e organizado. Mantém equipamentos, ferramentas e
instrumentos organizados, acondicionados e em plenas condições de uso e
funcionamento. Preserva o meio ambiente, identificando materiais
poluentes e fazendo seu descarte de acordo com normas ambientais.
Trabalha com segurança, definindo equipamentos de proteção individual de
acordo com as atividades e identificando localização de equipamentos de
proteção coletiva. Identifica e isola áreas de risco. Identifica
condições inseguras e previne atos inseguros. Solicita visto de
liberação de equipamentos e de locais de trabalho. Sugere treinamentos
de segurança no trabalho.

\begin{Form}
\textbf{Esta correspondência está adequada?} \\ 
\ChoiceMenu[radio,radiosymbol=\ding{52},bordercolor=0 0.5 0,name=cbovalid_ 914205 ]{}{CBO deve ficar=sim, \hspace{6em} CBO não fica aqui=nao}
\end{Form}

\newpage
\section{ 02016 -- Técnico em Manutenção de Máquinas Pesadas }

\textbf{Perfil do Curso (CNCT)}

O Técnico em Manutenção de Máquinas Pesadas será habilitado para:

Planejar, controlar e executar atividades relativas à manutenção de
máquinas pesadas automotoras sobre pneus e esteiras. Realizar atividades
de inspeção, atualização tecnológica, elaboração de planos de manutenção
e projetos em máquinas pesadas. Planejar, controlar e executar
procedimentos de desmontagem, montagem, medição, lubrificação e ensaios
dos sistemas mecânicos e automatizados em máquinas pesadas, atendendo às
normas e aos padrões técnicos de qualidade, saúde e segurança e de meio
ambiente. Executar a instalação de acessórios e equipamentos em máquinas
pesadas.

Para atuação como Técnico em Manutenção de Máquinas Pesadas, são
fundamentais:

Conhecimentos e saberes relacionados aos processos de manutenção de
equipamentos mecânicos pesados e construção de modo a assegurar a saúde
e a segurança dos trabalhadores e usuários. Conhecimentos e saberes
relacionados à sustentabilidade do processo produtivo, às normas
técnicas, à liderança de equipes, à solução de problemas técnicos e
trabalhistas, à tomada de decisões e à gestão de conflitos.

\textbf{CBOs correspondidos e seus perfis (presentes no CNCT, e presentes na predição da IA)}

\textbf{CBO: 913105  ( Mecânico de manutenção de aparelhos de levantamento )}

Planeja atividades de manutenção, interpretando desenhos, projetos e
catálogos e consultando manuais técnicos. Seleciona ferramental e
equipamentos auxiliares. Define peças para reposição. Seleciona mão de
obra de acordo com a atividade. Estima tempo de realização da
manutenção. Realiza manutenção em aparelhos de levantamento e
movimentação de cargas, executando a desmontagem do equipamento para
obter pleno acesso às partes e procedendo à montagem, após as ações de
manutenção. Instala, substitui ou troca peças e acessórios do
equipamento, de acordo com procedimentos de manutenção. Monta
rolamentos. Efetua regulagem de motores, conjuntos e sistemas de freios.
Lubrifica e limpa o equipamento. Alinha conjuntos de transmissão. Ajusta
peças. Modifica parâmetros de desempenho do equipamento. Recupera
redutores. Repara motores e cilindros hidráulicos. Sana vazamentos
hidráulicos e pneumáticos. Troca embuchamento, correias, cabos de aço,
roldanas e revestimentos. Substitui filtros, fluidos e conectores
eletrônicos. Prepara peças para montagem de aparelhos de levantamento e
movimentação de cargas, adaptando peças e conjuntos e conferindo medidas
de peças. Realiza operações de corte, soldagem e usinagem, utilizando
ferramentas, máquinas e instrumentos específicos, na preparação das
peças para montagem. Analisa a qualidade das peças e executa sua
limpeza. Inspeciona funcionamento de aparelhos de levantamento e
movimentação de cargas, interpretando ordens de serviço e especificações
técnicas, analisando informações do operador e coletando dados e
informações de fontes diversificadas para determinar os equipamentos que
requerem reparo ou manutenção e identificar suas falhas e defeitos.
Vistoria sistema de arrefecimento e verifica restrição de sistemas de
ar. Analisa a estrutura do equipamento e vistoria sistemas hidráulico,
pneumático, elétrico e eletrônico. Analisa desgastes de chapas, de peças
de sistemas rodantes e de correias e correntes. Verifica conservação de
cabos de aço. Confere sistema de lubrificação e regulagem de motores.
Testa funcionamento de aparelhos de levantamento e movimentação de
cargas, utilizando instrumentos e métodos específicos, a fim de avaliar
o desempenho do equipamento e de suas partes. Configura e opera os
equipamentos, atuando nos dispositivos de comando e controle para
regular as operações durante os testes. Testa rendimento de potência de
motores. Confere pressão de sistemas hidráulicos e pneumáticos. Mede
rotação de motores. Testa sistemas vibratórios do equipamento. Confere
engrenamento de sistemas de transmissão. Testa tensionamento de
correias. Confere folgas de eixos e mancais. Analisa tipos e índices de
contaminação de fluidos. Coleta amostras para análise laboratorial.
Confere calibração de balanças. Testa capacidade de içamento de cargas e
capacidade de produção do equipamento. Monitora condições de
funcionamento do equipamento. Testa sistemas de embreagens e unidades
compressoras. Mede emissão de poluentes. Testa sistemas eletrônicos.
Registra informações técnicas, elaborando relatórios, atualizando
históricos de máquinas e registrando resultados de testes. Registra
informações de reparo e manutenção. Preenche requisições, emite ordens
de serviço e notas de serviços prestados. Registra horas trabalhadas de
peças sobressalentes. Elabora orçamentos. Busca manter-se atualizado,
tendo em vista a diversificação dos equipamentos, as novas tecnologias
específicas para levantamento e movimentação de cargas e os
desdobramentos dos movimentos tecnológicos que impactam a indústria.
Trabalha com segurança, desligando equipamentos para manutenção,
bloqueando chaves de partida de equipamentos e despressurizando sistemas
hidráulicos e pneumáticos. Utiliza Equipamentos de Proteção Individual
(EPI). Identifica áreas de risco e isola a área de trabalho. Limpa e
organiza o local de trabalho, acondicionando peças e ferramentas.
Seleciona materiais para reciclagem e acondiciona resíduos para
descarte.

\begin{Form}
\textbf{Esta correspondência está adequada?} \\ 
\ChoiceMenu[radio,radiosymbol=\ding{52},bordercolor=0 0.5 0,name=cbovalid_ 913105 ]{}{CBO deve ficar=sim, \hspace{6em} CBO não fica aqui=nao}
\end{Form}

\textbf{CBO: 913110  ( Mecânico de manutenção de equipamento de mineração )}

Planeja atividades de manutenção, interpretando desenhos, projetos e
catálogos e consultando manuais técnicos. Seleciona ferramental e
equipamentos auxiliares. Define peças para reposição. Seleciona mão de
obra de acordo com a atividade. Estima tempo de realização da
manutenção. Realiza manutenção em equipamentos de mineração, executando
a desmontagem do equipamento para obter pleno acesso às partes e
procedendo à montagem, após as ações de manutenção. Instala, substitui
ou troca peças e acessórios do equipamento, de acordo com procedimentos
de manutenção. Monta rolamentos. Efetua regulagem de motores, conjuntos
e sistemas de freios. Lubrifica e limpa o equipamento. Alinha conjuntos
de transmissão. Ajusta peças. Modifica parâmetros de desempenho do
equipamento. Recupera redutores. Repara motores e cilindros hidráulicos.
Sana vazamentos hidráulicos e pneumáticos. Troca embuchamento, correias,
cabos de aço, roldanas e revestimentos. Substitui filtros, fluidos e
conectores eletrônicos. Prepara peças para montagem de equipamentos de
mineração, adaptando peças e conjuntos e conferindo medidas de peças.
Realiza operações de corte, soldagem e usinagem, utilizando ferramentas,
máquinas e instrumentos específicos, na preparação das peças para
montagem. Analisa a qualidade das peças e executa sua limpeza.
Inspeciona funcionamento de equipamentos de mineração, interpretando
ordens de serviço e especificações, analisando informações do operador e
coletando dados e informações de fontes diversificadas, para determinar
necessidades de reparo ou manutenção e identificar falhas e defeitos.
Vistoria sistema de arrefecimento e verifica restrição de sistemas de
ar. Analisa a estrutura do equipamento e vistoria sistemas hidráulico,
pneumático, elétrico e eletrônico. Analisa desgastes de chapas, de peças
de sistemas rodantes e de correias e correntes. Verifica conservação de
cabos de aço. Confere sistema de lubrificação e regulagem de motores.
Verifica estado de conservação de moitões. Testa funcionamento de
equipamentos de mineração, utilizando instrumentos e métodos
específicos, a fim de avaliar o desempenho do equipamento e de suas
partes. Configura e opera os equipamentos, atuando nos dispositivos de
comando e controle para regular as operações durante os testes. Testa
rendimento de potência de motores. Confere pressão de sistemas
hidráulicos e pneumáticos. Mede rotação de motores. Testa sistemas
vibratórios do equipamento. Confere engrenamento de sistemas de
transmissão. Testa tensionamento de correias. Confere folgas de eixos e
mancais. Testa cabeçotes de perfuratrizes. Analisa tipos e índices de
contaminação de fluidos. Coleta amostras para análise laboratorial.
Confere calibração de balanças. Testa capacidade de içamento de cargas e
capacidade de produção do equipamento. Monitora condições de
funcionamento do equipamento. Testa sistemas de embreagens e unidades
compressoras. Mede emissão de poluentes. Testa sistemas eletrônicos.
Registra informações técnicas, elaborando relatórios, atualizando
históricos de máquinas e registrando resultados de testes. Registra
informações de reparo e manutenção do equipamento. Preenche requisições
de insumos, emite ordens de serviço e notas de serviços prestados.
Registra horas trabalhadas de peças sobressalentes, para controle.
Elabora orçamentos de serviços. Busca manter-se atualizado, tendo em
vista a diversificação dos equipamentos, as novas tecnologias
específicas para mineração e os desdobramentos dos movimentos
tecnológicos que impactam a indústria, em geral, e a mineração. Trabalha
com segurança, desligando equipamentos para manutenção, bloqueando
chaves de partida de equipamentos e despressurizando sistemas
hidráulicos e pneumáticos. Utiliza Equipamentos de Proteção Individual
(EPI). Identifica áreas de risco e isola a área de trabalho. Limpa e
organiza o local de trabalho, acondicionando peças e ferramentas.
Seleciona materiais para reciclagem e acondiciona resíduos para
descarte.

\begin{Form}
\textbf{Esta correspondência está adequada?} \\ 
\ChoiceMenu[radio,radiosymbol=\ding{52},bordercolor=0 0.5 0,name=cbovalid_ 913110 ]{}{CBO deve ficar=sim, \hspace{6em} CBO não fica aqui=nao}
\end{Form}

\textbf{CBO: 913115  ( Mecânico de manutenção de máquinas agrícolas )}

Planeja atividades de manutenção, interpretando desenhos, projetos e
catálogos e consultando manuais técnicos. Seleciona ferramental e
equipamentos auxiliares. Define peças para reposição. Seleciona mão de
obra de acordo com a atividade. Estima tempo de realização da
manutenção. Realiza a manutenção em máquinas agrícolas, executando a
desmontagem do equipamento para obter pleno acesso às partes e
procedendo à montagem, após as ações de manutenção. Instala, substitui
ou troca peças e acessórios do equipamento, de acordo com procedimentos
de manutenção. Monta rolamentos. Efetua regulagem de motores, conjuntos
e sistemas de freios. Lubrifica e limpa o equipamento. Alinha conjuntos
de transmissão. Ajusta peças. Modifica parâmetros de desempenho do
equipamento. Recupera redutores. Repara motores e cilindros hidráulicos.
Sana vazamentos hidráulicos e pneumáticos. Troca embuchamento, correias,
cabos de aço, roldanas e revestimentos. Substitui filtros, fluidos e
conectores eletrônicos. Prepara as peças para montagem de máquinas
agrícolas, adaptando peças e conjuntos e conferindo medidas de peças.
Realiza operações de corte, soldagem e usinagem, utilizando ferramentas,
máquinas e instrumentos específicos, na preparação das peças para
montagem. Analisa a qualidade das peças e executa sua limpeza.
Inspeciona o funcionamento de máquinas agrícolas, interpretando ordens
de serviço e especificações, analisando informações do operador e
coletando dados e informações de fontes diversificadas, para determinar
necessidade de reparo ou manutenção e identificar falhas e defeitos.
Vistoria transportadores. Mede malha de telas. Verifica estado de
conservação de moitões. Vistoria sistema de arrefecimento e verifica
restrição de sistemas de ar. Analisa a estrutura do equipamento e
vistoria sistemas hidráulico, pneumático, elétrico e eletrônico. Analisa
desgastes de chapas, de peças de sistemas rodantes e de correias e
correntes. Verifica conservação de cabos de aço. Confere sistema de
lubrificação e regulagem de motores. Testa o funcionamento de máquinas
agrícolas, utilizando instrumentos e métodos específicos, a fim de
avaliar o desempenho do equipamento e de suas partes. Configura e opera
os equipamentos, atuando nos dispositivos de comando e controle para
regular as operações durante os testes. Testa rendimento de potência de
motores. Confere pressão de sistemas hidráulicos e pneumáticos. Mede
rotação de motores. Testa sistemas vibratórios do equipamento. Confere
engrenamento de sistemas de transmissão. Testa tensionamento de
correias. Confere folgas de eixos e mancais. Analisa tipos e índices de
contaminação de fluidos. Coleta amostras para análise laboratorial.
Confere calibração de balanças. Testa capacidade de produção do
equipamento. Monitora condições de funcionamento do equipamento. Testa
sistemas de embreagens e unidades compressoras. Mede emissão de
poluentes. Testa eletroímãs e sistemas eletrônicos. Registra informações
técnicas, elaborando relatórios, atualizando históricos de máquinas e
registrando resultados de testes. Registra informações de reparo e
manutenção do equipamento. Preenche requisições de insumos, emite ordens
de serviço e notas de serviços prestados. Registra horas trabalhadas de
peças sobressalentes, para controle. Elabora orçamentos de serviços.
Busca manter-se atualizado, tendo em vista a diversificação dos
equipamentos, as novas tecnologias específicas para máquinas agrícolas e
os desdobramentos dos movimentos tecnológicos que impactam a indústria,
em geral, e as operações agrícolas. Trabalha com segurança, desligando
equipamentos para manutenção, bloqueando chaves de partida de
equipamentos e despressurizando sistemas hidráulicos e pneumáticos.
Utiliza Equipamentos de Proteção Individual (EPI). Identifica áreas de
risco e isola a área de trabalho. Limpa e organiza o local de trabalho,
acondicionando peças e ferramentas. Seleciona materiais para reciclagem
e acondiciona resíduos para descarte.

\begin{Form}
\textbf{Esta correspondência está adequada?} \\ 
\ChoiceMenu[radio,radiosymbol=\ding{52},bordercolor=0 0.5 0,name=cbovalid_ 913115 ]{}{CBO deve ficar=sim, \hspace{6em} CBO não fica aqui=nao}
\end{Form}

\textbf{CBO: 913120  ( Mecânico de manutenção de máquinas de construção e terraplenagem )}

Planeja atividades de manutenção, interpretando desenhos, projetos e
catálogos e consultando manuais técnicos. Seleciona ferramental e
equipamentos auxiliares. Define peças para reposição. Seleciona mão de
obra de acordo com a atividade. Estima tempo de realização da
manutenção. Realiza manutenção em máquinas de construção e
terraplenagem, executando a desmontagem do equipamento para obter pleno
acesso às partes e procedendo à montagem, após as ações de manutenção.
Instala, substitui ou troca peças e acessórios do equipamento, conforme
procedimentos de manutenção. Monta rolamentos. Efetua regulagem de
motores, conjuntos e sistemas de freios. Lubrifica e limpa o
equipamento. Alinha conjuntos de transmissão. Ajusta peças. Modifica
parâmetros de desempenho do equipamento. Recupera redutores. Repara
motores e cilindros hidráulicos. Sana vazamentos hidráulicos e
pneumáticos. Troca embuchamento, correias, cabos de aço, roldanas e
revestimentos. Substitui filtros, fluidos e conectores eletrônicos.
Prepara peças para montagem de máquinas de construção e terraplenagem,
adaptando peças e conjuntos e conferindo medidas de peças. Realiza
operações de corte, soldagem e usinagem, utilizando ferramentas,
máquinas e instrumentos específicos, na preparação das peças para
montagem. Analisa a qualidade das peças e executa sua limpeza.
Inspeciona o funcionamento de máquinas de construção e terraplenagem,
interpretando ordens de serviço e especificações técnicas, analisando
informações do operador e coletando dados e informações de fontes
diversificadas, para determinar os equipamentos que requerem reparo ou
manutenção e identificar suas falhas e defeitos. Vistoria sistema de
arrefecimento e verifica restrição de sistemas de ar. Mede malha de
telas. Analisa a estrutura do equipamento e vistoria sistemas
hidráulico, pneumático, elétrico e eletrônico. Vistoria transportadores.
Analisa desgastes de chapas, de peças de sistemas rodantes e de correias
e correntes. Verifica conservação de cabos de aço. Confere sistema de
lubrificação e regulagem de motores. Testa funcionamento de máquinas de
construção e terraplenagem, utilizando instrumentos e métodos
específicos, para avaliar o desempenho do equipamento e de suas partes.
Configura e opera equipamentos, atuando nos dispositivos de comando e
controle para regular as operações durante os testes. Testa rendimento
de potência de motores. Confere pressão de sistemas hidráulicos e
pneumáticos. Testa cabeçotes de perfuratrizes. Mede rotação de motores.
Testa sistemas vibratórios do equipamento. Confere engrenamento de
sistemas de transmissão. Testa tensionamento de correias. Confere folgas
de eixos e mancais. Analisa tipos e índices de contaminação de fluidos.
Coleta amostras para análise laboratorial. Testa capacidade de içamento
de cargas e capacidade de produção do equipamento. Monitora condições de
funcionamento do equipamento. Testa sistemas de embreagens e unidades
compressoras. Mede emissão de poluentes. Testa eletroímãs e sistemas
eletrônicos. Registra informações técnicas, elaborando relatórios,
atualizando históricos de máquinas e registrando resultados de testes.
Registra informações de reparo e manutenção do equipamento. Preenche
requisições de insumos, emite ordens de serviço e notas de serviços
prestados. Elabora orçamentos de serviços. Busca manter-se atualizado,
tendo em vista a diversificação dos equipamentos, as novas tecnologias
específicas para construção e terraplenagem e os desdobramentos dos
movimentos tecnológicos que impactam a indústria, em geral, e a
construção civil. Trabalha com segurança, desligando equipamentos para
manutenção, bloqueando chaves de partida de equipamentos e
despressurizando sistemas hidráulicos e pneumáticos. Utiliza
Equipamentos de Proteção Individual (EPI). Identifica áreas de risco e
isola a área de trabalho. Limpa e organiza o local de trabalho,
acondicionando peças e ferramentas. Seleciona materiais para reciclagem
e acondiciona resíduos para descarte.

\begin{Form}
\textbf{Esta correspondência está adequada?} \\ 
\ChoiceMenu[radio,radiosymbol=\ding{52},bordercolor=0 0.5 0,name=cbovalid_ 913120 ]{}{CBO deve ficar=sim, \hspace{6em} CBO não fica aqui=nao}
\end{Form}

\textbf{CBO: 914420  ( Mecânico de manutenção de tratores )}

Elabora plano de manutenção mecânica de motores, sistemas e partes de
tratores de rodas e esteiras, interpretando desenhos e normas técnicas,
diagnosticando falhas e identificando o trabalho a ser realizado, para
estimar tempo de execução, orçar o serviço e requisitar material.
Confirma o plano de manutenção com o cliente e com o fornecedor de
material. Seleciona equipamentos, ferramentas e instrumentos, de acordo
com o trabalho a ser executado. Realiza manutenção mecânica de motores,
removendo e desmontando o motor. Monta e reinstala o motor, executando o
ajuste de válvulas. Remove o sistema de transmissão, efetua ajuste de
montagem, limpa filtros e reinstala o sistema. Ajusta componentes
mecânicos, elétricos, hidráulicos e pneumáticos. Regula e faz a sangria
do sistema de freios. Drena filtros da linha de alimentação. Regula o
sistema de ignição e injeção. Ajusta cubos de rodas e esteiras. Regula
altura e lubrifica articulações da suspensão. Alinha e substitui braços
do sistema de direção. Substitui peças dos diversos sistemas de tratores
de rodas e esteiras, identificando as partes com defeitos de fabricação
ou desgastes pelo uso, para realizar troca de agregados de suspensão e
outros, tais como compressor, alternador e bomba d'água. Troca fluidos
hidráulicos, sistema de embreagem, válvula injetora, componentes
eletroeletrônicos, amortecedores, barras estabilizadoras, hastes de
reação, molas, válvulas pneumáticas, buchas, pivôs, rolamentos de rodas,
componentes da unidade hidráulica, radiadores, sensores térmicos,
mangueiras, correias, polias, vedantes de óleo e água e bolsas de ar.
Troca componentes de esteiras. Repara componentes, realizando a
manutenção corretiva por meio de ações de reparo em sistema elétrico e
sistemas de arrefecimento, carga e partida, sobrealimentação e
escapamento. Repara válvulas pneumáticas, bomba de combustível, válvula
injetora diesel, bomba injetora, bomba hidráulica da direção, caixa de
direção, chassis, unidade hidráulica e bomba d'água. Realiza manutenção
da bomba hidráulica e do compressor de ar. Pode executar soldagem
elétrica e soldagem oxiacetilênica, como parte das atividades. Testa
desempenho de componentes e sistemas dos tratores de rodas e esteiras,
realizando testes de motor no veículo, testando sistema de transmissão e
válvulas injetoras e efetuando testes hidráulicos e pneumáticos. Testa
circuitos eletroeletrônicos, pressão de alimentação e vazão. Verifica o
funcionamento do sistema de arrefecimento. Examina o funcionamento de
alavanca, cabo e demais componentes do sistema de freios. Verifica as
condições da tubulação e do tanque de combustível. Realiza teste de
emissão de poluentes. Avalia a satisfação do cliente com os serviços
executados. Organiza e providencia a limpeza do local de trabalho.
Mantém equipamentos, ferramentas, instrumentos e acessórios de trabalho
organizados, acondicionados e em plenas condições de uso e
funcionamento. Zela pela segurança no trabalho, identificando áreas de
risco, vestindo equipamentos de proteção individual e consultando
recomendações de segurança contidas nos manuais. Controla desperdícios
de material. Separa, identifica e classifica resíduos -- sólidos e
líquidos -, para descarte, reúso e reciclagem. Realiza o descarte de
acordo com as normas ambientais.

\begin{Form}
\textbf{Esta correspondência está adequada?} \\ 
\ChoiceMenu[radio,radiosymbol=\ding{52},bordercolor=0 0.5 0,name=cbovalid_ 914420 ]{}{CBO deve ficar=sim, \hspace{6em} CBO não fica aqui=nao}
\end{Form}

\newpage
\section{ 02017 -- Técnico em Manutenção de Sistemas Metroferroviários }

\textbf{Perfil do Curso (CNCT)}

O Técnico em Manutenção de Sistemas Metroferroviários será habilitado
para:

Planejar, controlar e executar atividades relativas à manutenção
mecânica nos níveis preventivos, preditivos e corretivos em componentes
de vagões, locomotivas e máquinas metroferroviárias. Planejar, controlar
e executar a fabricação e a montagem de peças em componentes de vagões,
locomotivas e máquinas metroferroviárias. Realizar reformas, testes de
comissionamento e de performance em equipamentos metroferroviários.
Especificar equipamentos e insumos para processos de manutenção
metroferroviária atendendo às normas e aos padrões técnicos de
qualidade, saúde e segurança e de meio ambiente. Realizar medições,
testes, calibrações e comissionamento de equipamentos eletrônicos
empregados em locomotivas e máquinas metroferroviárias. Instalar e
configurar equipamentos e/ou instrumentos destinados à automação de
locomotivas e máquinas metroferroviárias. Atuar na manutenção dos
sistemas automatizados em componentes de vagões, locomotivas e máquinas
metroferroviárias, atendendo às normas e aos padrões técnicos de
qualidade, saúde e segurança e de meio ambiente. Reconhecer as
tecnologias empregadas nos sistemas de sinalização em malhas
metroferroviárias. Reconhecer tecnologias inovadoras presentes no
segmento visando a atender às transformações digitais na sociedade.

Para atuação como Técnico em Manutenção de Sistemas Metroferroviários,
são fundamentais:

Conhecimentos e saberes relacionados aos processos de manutenção de
sistemas metroferroviários e construção de modo a assegurar a saúde e a
segurança dos trabalhadores e usuários. Conhecimentos e saberes
relacionados à sustentabilidade do processo produtivo, às normas
técnicas, à liderança de equipes, à solução de problemas técnicos e
trabalhistas, à tomada de decisões e à gestão de conflitos.

\textbf{CBOs correspondidos e seus perfis (presentes no CNCT, e presentes na predição da IA)}

\textbf{CBO: 914305  ( Mecânico de manutenção de veículos ferroviários )}

Planeja atividades de manutenção, interpretando procedimentos, manuais
técnicos, desenhos e esquemas, de acordo com os diversos tipos de
veículos ferroviários. Programa atividades de manutenção, determinando
prioridades de execução, selecionando materiais e ferramentas e
estimando tempo de execução. Faz uso de sistemas de informação de apoio
à manutenção. Utiliza técnicas de diagnóstico, identificando problemas
nos veículos ferroviários e suas causas, para definir ações de
manutenção. Realiza manutenção em veículos ferroviários, lubrificando
componentes e conjuntos. Aplica e controla torque em elementos de
fixação, ajusta folgas entre componentes, nivela caixa do veículo e
alinha conjunto rotativo motriz. Repara componentes e conjuntos e
substitui peças e componentes, realizando operações de corte de peças
metálicas e não metálicas e soldando componentes metálicos, entre outras
operações. Repara acabamentos do interior de veículos. Realiza medições
e testes em peças, componentes e veículos ferroviários, utilizando
instrumentos específicos, comparando resultados com dados técnicos e
interpretando resultados, a fim de garantir atendimento de
especificações. Mede peças e componentes para realização de testes de
fadiga. Mede taxa de compressão e emissão de poluentes. Testa
equipamentos, sistemas de freio manual e sistemas hidráulicos e
pneumáticos. Realiza testes de potência do sistema motriz e testes de
estanqueidade. Testa rendimento dos compressores, avalia especificações
de vibração e ruído, monitora taxa de desgaste de componentes e verifica
trincas estruturais em componentes. Pode testar componentes de
eletrônica embarcada e de automação -- diretamente relacionados aos
elementos mecânicos de máquinas dos veículos ferroviários --, podendo
removê-los, substituí-los ou instalá-los, a fim de restabelecer o pleno
funcionamento das partes em manutenção. Manobra equipamentos
ferroviários, avaliando condições de operação de equipamentos e
acionando sistemas do veículo. Conduz veículos em pátios de manobra da
oficina, cumprindo requisitos de segurança, para avaliar condições de
tráfego de vias. Reforma veículos ferroviários, desmontando conjuntos e
equipamentos, avaliando condição de equipamentos e componentes e
solicitando materiais e peças. Define substituição de componentes.
Propõe modificações em peças e equipamentos e confecciona peças.
Substitui partes, analisando compatibilidade entre componentes e
equipamentos. Recupera estruturas. Monta equipamentos e conjuntos.
Inspeciona veículos ferroviários, avaliando informações do histórico de
equipamentos, verificando condições gerais, examinando parâmetros
específicos de funcionamento e interpretando dados e indicadores de
análises. Inspeciona equipamentos visualmente. Retém veículos sem
condições de serviço e libera veículos em condições de tráfego. Registra
os resultados das análises de inspeção. Mantém-se atualizado acerca das
inovações tecnológicas da área de mecânica ferroviária, consultando
fontes diversificadas de informação, interpretando e analisando
informações técnicas de veículos ferroviários, para realizar os serviços
de manutenção em elevados níveis de qualidade, eficiência e segurança.
Pode supervisionar equipe de trabalhadores de manutenção, avaliando
desempenho e promovendo programas de treinamento. Administra pessoal,
participando de seleção de membros da equipe, atribuindo tarefas e
acompanhando registros de faltas e atrasos. Conserva o local de trabalho
limpo e organizado. Mantém equipamentos, ferramentas e instrumentos
organizados, acondicionados e em plenas condições de uso e
funcionamento. Atua na preservação do meio ambiente, selecionando
resíduos para reciclagem e realizando descarte dos demais, de acordo com
normas ambientais. Trabalha com segurança, utilizando Equipamentos de
Proteção Individual (EPI). Orienta e monitora o uso de EPI pela equipe.
Sinaliza e isola áreas de trabalho. Identifica condições e atos
inseguros, atuando para eliminá-los.

\begin{Form}
\textbf{Esta correspondência está adequada?} \\ 
\ChoiceMenu[radio,radiosymbol=\ding{52},bordercolor=0 0.5 0,name=cbovalid_ 914305 ]{}{CBO deve ficar=sim, \hspace{6em} CBO não fica aqui=nao}
\end{Form}

\newpage
\section{ 02018 -- Técnico em Mecânica }

\textbf{Perfil do Curso (CNCT)}

O Técnico em Mecânica será habilitado para:

Programar, controlar e executar processos de fabricação mecânica para
máquinas e equipamentos mecânicos atendendo às normas e aos padrões
técnicos de qualidade, saúde e segurança e de meio ambiente. Planejar,
aplicar e controlar procedimentos de instalação, de manutenção e
inspeção mecânica de máquinas e equipamentos. Elaborar projetos de
produtos relacionados a máquinas e equipamentos mecânicos especificando
materiais para construção mecânica por meio de técnicas de usinagem,
soldagem e conformação mecânica. Realizar inspeção visual, dimensional e
testes em sistemas, instrumentos e equipamentos mecânicos, pneumáticos,
hidráulicos e eletromecânicos de máquinas. Reconhecer tecnologias
inovadoras presentes no segmento visando a atender às transformações
digitais na sociedade.

Para atuação como Técnico em Mecânica, são fundamentais:

Conhecimentos e saberes relacionados aos processos de planejamento,
produção e manutenção de equipamentos mecânicos de modo a assegurar a
saúde e a segurança dos trabalhadores e dos usuários. Conhecimentos e
saberes relacionados à sustentabilidade do processo produtivo, às
técnicas e aos processos de produção, às normas técnicas, à liderança de
equipes, à solução de problemas técnicos e trabalhistas e à gestão de
conflitos.

\textbf{CBOs correspondidos e seus perfis (presentes no CNCT, e presentes na predição da IA)}

\textbf{CBO: 314105  ( Técnico em mecânica de precisão )}

Elabora projetos relacionados a máquinas, equipamentos e sistemas
mecânicos de alta precisão, fazendo desenhos técnicos, especificando
componentes, definindo leiaute e preparando a documentação técnica.
Participa da execução técnica do projeto. Monta máquinas, equipamentos e
sistemas mecânicos de alta precisão, de acordo com manuais e desenhos
técnicos. Realiza ajustes dimensionais, de posição, detecção e possível
correção de falhas de projeto, testando o funcionamento. Instala
máquinas, equipamentos e sistemas mecânicos de alta precisão, observando
as condições do ambiente, realizando ajustes e testes. Realiza a
manutenção de máquinas, equipamentos e sistemas mecânicos de alta
precisão, por meio de detecção de falhas e substituição de peças e
componentes. Desenvolve processos de fabricação e montagem relacionados
à mecânica de precisão, estabelecendo sequência de operações, tempos e
métodos, de acordo com recursos disponíveis. Realiza teste de ajuste
final e ações de controle e aprimoramento da qualidade e produtividade.
Pode usar novas técnicas de fabricação a partir das tecnologias do pó,
da nanotecnologia e de métodos de microfabricação. Pode realizar
treinamento operacional dos processos de fabricação e montagem
desenvolvidos. Redige relatórios técnicos. Pode trabalhar tecnicamente
na compra de produtos, avaliando as condições técnicas de contratos e a
especificação de serviços e avaliando desempenho de fornecedores. Pode
trabalhar tecnicamente na venda de produtos, analisando necessidades de
clientes, fazendo demonstração de produto e prestando assistência
técnica.

\begin{Form}
\textbf{Esta correspondência está adequada?} \\ 
\ChoiceMenu[radio,radiosymbol=\ding{52},bordercolor=0 0.5 0,name=cbovalid_ 314105 ]{}{CBO deve ficar=sim, \hspace{6em} CBO não fica aqui=nao}
\end{Form}

\textbf{CBO: 314110  ( Técnico mecânico )}

Elabora projetos relacionados a máquinas, equipamentos e sistemas
mecânicos, elaborando desenhos técnicos, especificando componentes e
definindo de leiaute. Pode incluir automação. Participa da execução
técnica do projeto. Monta máquinas, equipamentos e sistemas mecânicos,
de acordo com manuais e desenhos técnicos, realizando ajustes
dimensionais, de posição, detecção e possível correção de falhas de
projeto e testes de funcionamento. Instala máquinas, equipamentos e
sistemas mecânicos observando as condições do ambiente, realizando
ajustes, testes e treinamento operacional. Realiza a manutenção de
máquinas, equipamentos e sistemas mecânicos, por meio de detecção de
falhas e de substituição de peças e componentes. Desenvolve processos
mecânicos de fabricação e montagem, estabelecendo sequência de
operações, tempos e métodos, de acordo com recursos disponíveis. Realiza
teste de ajuste final e ações de controle e aprimoramento da qualidade e
produtividade. Pode incorporar novas técnicas de fabricação a partir das
tecnologias do pó. Prepara documentação técnica, observando
procedimentos e métodos de registro. Pode trabalhar tecnicamente na
venda e em pós-venda de produtos.

\begin{Form}
\textbf{Esta correspondência está adequada?} \\ 
\ChoiceMenu[radio,radiosymbol=\ding{52},bordercolor=0 0.5 0,name=cbovalid_ 314110 ]{}{CBO deve ficar=sim, \hspace{6em} CBO não fica aqui=nao}
\end{Form}

\textbf{CBO: 915105  ( Técnico em manutenção de instrumentos de medição e precisão )}

Planeja as atividades de manutenção de instrumentos, equipamentos e
sistemas de medição e de precisão, orientando-se por instruções de
trabalho; interpretando normas e desenhos técnicos; e selecionando
ferramentas. Utiliza sistemas de informação de apoio à manutenção.
Realiza a desmontagem de instrumentos, equipamentos e sistemas de
medição e de precisão, limpando-os, identificando defeitos e falhas e
analisando suas causas, a fim de definir a estratégia de reparação.
Descontamina manômetros. Seleciona peças e componentes para reposição.
Desenvolve gabaritos. Recupera peças e fabrica ferramentas e peças
específicas para reparação. Pode construir peças para manutenção por
meio de métodos de fabricação aditiva. Repara instrumentos, equipamentos
e sistemas de medição e de precisão, substituindo peças e componentes
mecânicos e eletrônicos, inclusive interfaces com redes. Realiza
lubrificação. Monta instrumentos, equipamentos e sistemas, efetuando
acabamento em seus gabinetes, carcaças e flanges. Realiza testes de
funcionamento em instrumentos, equipamentos e sistemas de medição e de
precisão, aplicando procedimentos de testes, a fim de validar o
atendimento de normas e especificações técnicas. Testa repetibilidade de
medições e testa histerese (fenômeno próprio de sistemas físicos cujas
propriedades dependem de sua história precedente). Realiza testes em
protótipos de instrumentos, equipamentos e sistemas. Verifica desvios de
indicação de medição em relação a padrões. Realiza ajustes e calibração
em instrumentos, equipamentos e sistemas de medição e de precisão, de
acordo com especificações técnicas e procedimentos de laboratórios
acreditados. Seleciona padrões para calibração e verifica classes de
tolerância. Executa calibração, comparando valores-padrão com indicações
dos instrumentos, equipamentos e sistemas; interpretando unidades de
medida; convertendo unidades de referência em unidades de medição
correlacionadas; realizando os ajustes necessários; e executando a
lacração após conclusão do processo de calibração. Embala e transporta
instrumentos, equipamentos e sistemas de medição e de precisão,
observando os procedimentos e cuidados necessários, em função dos
requisitos dos serviços de manutenção. Pode instalar instrumentos,
equipamentos e sistemas de medição e de precisão. Orienta clientes
quanto às condições de instalação e uso de instrumentos, equipamentos e
sistemas. Executa adequações de instrumentos, equipamentos e sistemas de
medição e de precisão em função de projetos, aplicando conhecimentos de
tecnologia mecânica e eletrônica -- analógica e digital --, montando
protótipos, adaptando protótipos a normas técnicas e implementando
correções em projetos. Pode treinar novos funcionários nas atividades de
manutenção. Elabora documentação técnica, emitindo laudos técnicos;
registrando alterações efetivadas em projetos; emitindo certificados de
funcionamento e de calibração de instrumentos, equipamentos e sistemas
de medição e de precisão. Elabora especificações técnicas. Emite
relatórios de itens fora das especificações. Elabora instruções de
reparos. Registra nível de temperatura e teor de umidade de ambientes,
especialmente do ambiente de calibração. Cumpre procedimentos
específicos para trabalhos em ``sala limpa''. Atua de acordo com
diretrizes de preservação ambiental, selecionando materiais para
reciclagem e realizando descarte de resíduos de acordo com normas
ambientais. Mantém equipamentos, ferramentas e instrumentos organizados,
acondicionados e em plenas condições de uso e funcionamento. Trabalha
com segurança, utilizando Equipamentos de Proteção Individual (EPI) e
organizando o local de trabalho. Identifica e isola área de riscos de
acidentes. Interage com setores de segurança do trabalho de empresas
clientes, a fim de cumprir normas e procedimentos específicos de
segurança no trabalho em atividades externas.

\begin{Form}
\textbf{Esta correspondência está adequada?} \\ 
\ChoiceMenu[radio,radiosymbol=\ding{52},bordercolor=0 0.5 0,name=cbovalid_ 915105 ]{}{CBO deve ficar=sim, \hspace{6em} CBO não fica aqui=nao}
\end{Form}

\newpage
\section{ 02019 -- Técnico em Mecânica de Precisão }

\textbf{Perfil do Curso (CNCT)}

O Técnico em Mecânica de Precisão será habilitado para:

Programar, controlar e executar as atividades de desmontagem e montagem
de sistemas mecânicos de precisão. Planejar os processos de manutenção
de máquinas e equipamentos mecânicos de precisão respeitando
procedimentos e normas técnicas, de qualidade, de saúde e segurança e de
meio ambiente. Diagnosticar as condições dos elementos de máquinas que
compõem sistemas mecânicos de precisão. Indicar processos de fabricação
mecânica com tolerâncias dimensionais adequadas aos projetos de máquinas
e equipamentos mecânicos de precisão. Realizar inspeção visual,
dimensional e testes em sistemas mecânicos de precisão, instrumentos e
equipamentos mecânicos, pneumáticos, hidráulicos e eletromecânicos de
máquinas. Reconhecer tecnologias inovadoras presentes no segmento
visando a atender às transformações digitais na sociedade.

Para atuação como Técnico em Mecânica de Precisão, são fundamentais:

Conhecimentos e saberes relacionados aos processos de planejamento,
produção e manutenção de dispositivos de precisão de modo a assegurar a
saúde e a segurança dos trabalhadores e dos usuários. Conhecimentos e
saberes relacionados à sustentabilidade do processo produtivo, às
técnicas e aos processos de produção, às normas técnicas, à liderança de
equipes, à solução de problemas técnicos e trabalhistas e à gestão de
conflitos.

\textbf{CBOs correspondidos e seus perfis (presentes no CNCT, e presentes na predição da IA)}

\textbf{CBO: 314105  ( Técnico em mecânica de precisão )}

Elabora projetos relacionados a máquinas, equipamentos e sistemas
mecânicos de alta precisão, fazendo desenhos técnicos, especificando
componentes, definindo leiaute e preparando a documentação técnica.
Participa da execução técnica do projeto. Monta máquinas, equipamentos e
sistemas mecânicos de alta precisão, de acordo com manuais e desenhos
técnicos. Realiza ajustes dimensionais, de posição, detecção e possível
correção de falhas de projeto, testando o funcionamento. Instala
máquinas, equipamentos e sistemas mecânicos de alta precisão, observando
as condições do ambiente, realizando ajustes e testes. Realiza a
manutenção de máquinas, equipamentos e sistemas mecânicos de alta
precisão, por meio de detecção de falhas e substituição de peças e
componentes. Desenvolve processos de fabricação e montagem relacionados
à mecânica de precisão, estabelecendo sequência de operações, tempos e
métodos, de acordo com recursos disponíveis. Realiza teste de ajuste
final e ações de controle e aprimoramento da qualidade e produtividade.
Pode usar novas técnicas de fabricação a partir das tecnologias do pó,
da nanotecnologia e de métodos de microfabricação. Pode realizar
treinamento operacional dos processos de fabricação e montagem
desenvolvidos. Redige relatórios técnicos. Pode trabalhar tecnicamente
na compra de produtos, avaliando as condições técnicas de contratos e a
especificação de serviços e avaliando desempenho de fornecedores. Pode
trabalhar tecnicamente na venda de produtos, analisando necessidades de
clientes, fazendo demonstração de produto e prestando assistência
técnica.

\begin{Form}
\textbf{Esta correspondência está adequada?} \\ 
\ChoiceMenu[radio,radiosymbol=\ding{52},bordercolor=0 0.5 0,name=cbovalid_ 314105 ]{}{CBO deve ficar=sim, \hspace{6em} CBO não fica aqui=nao}
\end{Form}

\textbf{CBO: 915105  ( Técnico em manutenção de instrumentos de medição e precisão )}

Planeja as atividades de manutenção de instrumentos, equipamentos e
sistemas de medição e de precisão, orientando-se por instruções de
trabalho; interpretando normas e desenhos técnicos; e selecionando
ferramentas. Utiliza sistemas de informação de apoio à manutenção.
Realiza a desmontagem de instrumentos, equipamentos e sistemas de
medição e de precisão, limpando-os, identificando defeitos e falhas e
analisando suas causas, a fim de definir a estratégia de reparação.
Descontamina manômetros. Seleciona peças e componentes para reposição.
Desenvolve gabaritos. Recupera peças e fabrica ferramentas e peças
específicas para reparação. Pode construir peças para manutenção por
meio de métodos de fabricação aditiva. Repara instrumentos, equipamentos
e sistemas de medição e de precisão, substituindo peças e componentes
mecânicos e eletrônicos, inclusive interfaces com redes. Realiza
lubrificação. Monta instrumentos, equipamentos e sistemas, efetuando
acabamento em seus gabinetes, carcaças e flanges. Realiza testes de
funcionamento em instrumentos, equipamentos e sistemas de medição e de
precisão, aplicando procedimentos de testes, a fim de validar o
atendimento de normas e especificações técnicas. Testa repetibilidade de
medições e testa histerese (fenômeno próprio de sistemas físicos cujas
propriedades dependem de sua história precedente). Realiza testes em
protótipos de instrumentos, equipamentos e sistemas. Verifica desvios de
indicação de medição em relação a padrões. Realiza ajustes e calibração
em instrumentos, equipamentos e sistemas de medição e de precisão, de
acordo com especificações técnicas e procedimentos de laboratórios
acreditados. Seleciona padrões para calibração e verifica classes de
tolerância. Executa calibração, comparando valores-padrão com indicações
dos instrumentos, equipamentos e sistemas; interpretando unidades de
medida; convertendo unidades de referência em unidades de medição
correlacionadas; realizando os ajustes necessários; e executando a
lacração após conclusão do processo de calibração. Embala e transporta
instrumentos, equipamentos e sistemas de medição e de precisão,
observando os procedimentos e cuidados necessários, em função dos
requisitos dos serviços de manutenção. Pode instalar instrumentos,
equipamentos e sistemas de medição e de precisão. Orienta clientes
quanto às condições de instalação e uso de instrumentos, equipamentos e
sistemas. Executa adequações de instrumentos, equipamentos e sistemas de
medição e de precisão em função de projetos, aplicando conhecimentos de
tecnologia mecânica e eletrônica -- analógica e digital --, montando
protótipos, adaptando protótipos a normas técnicas e implementando
correções em projetos. Pode treinar novos funcionários nas atividades de
manutenção. Elabora documentação técnica, emitindo laudos técnicos;
registrando alterações efetivadas em projetos; emitindo certificados de
funcionamento e de calibração de instrumentos, equipamentos e sistemas
de medição e de precisão. Elabora especificações técnicas. Emite
relatórios de itens fora das especificações. Elabora instruções de
reparos. Registra nível de temperatura e teor de umidade de ambientes,
especialmente do ambiente de calibração. Cumpre procedimentos
específicos para trabalhos em ``sala limpa''. Atua de acordo com
diretrizes de preservação ambiental, selecionando materiais para
reciclagem e realizando descarte de resíduos de acordo com normas
ambientais. Mantém equipamentos, ferramentas e instrumentos organizados,
acondicionados e em plenas condições de uso e funcionamento. Trabalha
com segurança, utilizando Equipamentos de Proteção Individual (EPI) e
organizando o local de trabalho. Identifica e isola área de riscos de
acidentes. Interage com setores de segurança do trabalho de empresas
clientes, a fim de cumprir normas e procedimentos específicos de
segurança no trabalho em atividades externas.

\begin{Form}
\textbf{Esta correspondência está adequada?} \\ 
\ChoiceMenu[radio,radiosymbol=\ding{52},bordercolor=0 0.5 0,name=cbovalid_ 915105 ]{}{CBO deve ficar=sim, \hspace{6em} CBO não fica aqui=nao}
\end{Form}

\newpage
\section{ 02020 -- Técnico em Mecatrônica }

\textbf{Perfil do Curso (CNCT)}

O Técnico em Mecatrônica será habilitado para:

Projetar, instalar e operar equipamentos automatizados e/ou robotizados
empregados em processos de manufatura considerando as normas, os padrões
e os requisitos técnicos de qualidade, saúde e segurança e de meio
ambiente. Realizar programação, parametrização, medições e testes de
equipamentos automatizados em processos de manufatura. Realizar
integração de equipamentos mecânicos e eletrônicos utilizados em
processos de manufatura. Reconhecer tecnologias inovadoras presentes no
segmento visando a atender às transformações digitais na sociedade.

Para atuação como Técnico em Mecatrônica, são fundamentais:

Conhecimentos e saberes relacionados ao planejamento e implementação de
processos automatizados de manufatura de modo a assegurar a saúde e a
segurança dos trabalhadores e dos usuários. Conhecimento e saberes
relacionados à sustentabilidade do processo produtivo, às técnicas e
processos de produção, às normas técnicas, à liderança de equipes, à
solução de problemas técnicos e trabalhistas e à gestão de conflitos.

\textbf{CBOs correspondidos e seus perfis (presentes no CNCT, e presentes na predição da IA)}

\textbf{CBO: 300105  ( Técnico em mecatrônica - automação da manufatura )}

Projeta sistemas de automação da manufatura, analisando o processo de
manufatura e o produto manufaturado envolvidos no projeto. Avalia as
condições do local de trabalho para instalação de máquinas e
equipamentos. Define fluxo do processo, identifica alternativas
tecnológicas e esboça proposta para automação do processo. Faz
especificação de materiais e componentes empregados em projeto.
Considera aspectos de ergonomia, de segurança do trabalho, de economia
de energia e de preservação do meio ambiente na elaboração de projeto.
Elabora relatório de custo-benefício para análise técnico-financeira do
projeto. Elabora cronograma de implantação de projeto. Pode realizar
projetos de integração de sistemas de automação da manufatura,
estabelecendo interfaces entre equipamentos eletroeletrônicos e
mecânicos. Participa de projetos de implantação e utilização de
tecnologias típicas das fábricas digitais nos processos de manufatura.
Instala sistemas de automação da manufatura, interpretando documentação
do projeto e organizando materiais e recursos para o processo de
instalação. Monta componentes eletroeletrônicos e componentes mecânicos
em sistemas de automação da manufatura. Testa a operação do sistema, sem
processamento de matéria-prima. Acompanha teste de produção do sistema,
em processo. Faz correções e ajustes no sistema, conforme resultados de
testes. Pode realizar a integração de diversificados tipos de robôs nos
sistemas de automação da manufatura. Pode aperfeiçoar sistemas de
automação da manufatura já instalados, substituindo equipamentos,
atualizando programas e implementando novas tecnologias. Programa
equipamentos e sistemas de automação da manufatura, escrevendo programas
em linguagens diversificadas de programação ou utilizando recursos
avançados, tais como sistemas supervisórios, CAD-Desenho Assistido por
Computador (Computer Aided Design), CAM--Manufatura Auxiliada por
Computador (Computer Aided Manufacturing), dentre outros. Utiliza
sistemas analíticos e ambientes de desenvolvimento para criação de
aplicações. Desenvolve a programação de controladores lógicos
programáveis (CLP) e de máquinas-ferramentas com comando numérico
computadorizado (CNC). Programa parâmetros para acionamentos de potência
e de redes industriais, utilizando ferramentas de programação
específicas, tendo em vista a integração de equipamentos. Realiza a
manutenção de sistemas de automação da manufatura, avaliando gráficos de
tendências e relatórios de manutenção e planejando o processo e as
etapas de trabalho de manutenção. Planeja manutenção preventiva e
preditiva. Analisa falhas em sistemas. Executa procedimentos de
manutenção. Avalia a eficácia das soluções de manutenção implementadas.
Colabora no processo de aquisição de componentes, equipamentos e
sistemas, verificando características técnicas dos bens; selecionando
fornecedores; e acompanhando testes de funcionamento de máquinas e
equipamentos para elaboração e emissão de parecer técnico. Avalia
disponibilidade de peças de reposição. Coordena equipes de trabalho,
identificando competências técnicas e pessoais, atribuindo
responsabilidades e estabelecendo metas individuais. Forma equipe
multidisciplinar para análise de máquinas e equipamentos. Promove
integração entre setores da empresa e os integrantes de equipes
envolvidas em projetos. Reúne-se com a equipe de trabalho. Monitora a
execução de tarefas. Dá suporte técnico aos integrantes da equipe.
Treina usuários em operação e manutenção de sistemas de automação da
manufatura. Elabora documentação técnica de projetos de sistemas de
automação da manufatura. Documenta alterações de projeto ocorridas
durante a instalação de sistemas. Elabora relatório de aceitação de
equipamentos. Documenta o plano de ação de manutenção preventiva e
preditiva, as ações corretivas e as melhorias implementadas nos
sistemas. Mantém o local de trabalho limpo, seguro e organizado.
Conserva ferramentas, instrumentos e equipamentos. Utiliza equipamentos
de proteção.

\begin{Form}
\textbf{Esta correspondência está adequada?} \\ 
\ChoiceMenu[radio,radiosymbol=\ding{52},bordercolor=0 0.5 0,name=cbovalid_ 300105 ]{}{CBO deve ficar=sim, \hspace{6em} CBO não fica aqui=nao}
\end{Form}

\textbf{CBO: 300110  ( Técnico em mecatrônica - robótica )}

Projeta sistemas robóticos, avaliando condições do local de trabalho
para instalação de máquinas e equipamentos; identificando alternativas
tecnológicas; e especificando materiais e componentes empregados em
projetos. Considera aspectos de ergonomia, de segurança do trabalho, de
segurança robótica e de preservação do meio ambiente na elaboração de
projeto. Elabora relatório de custo-benefício para análise
técnico-financeira de projeto. Elabora cronograma de implantação de
projeto. Pode realizar projetos de integração de sistemas robóticos com
outros sistemas de automação, estabelecendo interfaces entre
equipamentos eletroeletrônicos e mecânicos. Pode aperfeiçoar sistemas
robóticos já instalados, substituindo equipamentos, atualizando
programas e implementando novas tecnologias. Pode realizar projetos de
sistemas robóticos utilizando robôs colaborativos. Participa de projetos
de implantação e utilização de tecnologias robóticas em diversificados
tipos de indústria, tais como automobilísticas, aeroespaciais, de óleo e
gás, dentre outras. Instala sistemas robóticos, interpretando
documentação do projeto e organizando materiais e recursos para o
processo de instalação. Monta componentes eletroeletrônicos e
componentes mecânicos em sistemas robóticos. Testa a operação do
sistema, sem processamento de matéria-prima. Acompanha teste de produção
do sistema, em processo. Faz correções e ajustes no sistema, conforme
resultados de testes. Pode realizar a integração de robôs em sistemas de
automação. Programa robôs e sistemas robóticos, escrevendo programas em
linguagens diversificadas de programação ou utilizando recursos
avançados, tais como sistemas supervisórios, CAD-Desenho Assistido por
Computador (Computer Aided Design), CAM--Manufatura Auxiliada por
Computador (Computer Aided Manufacturing), dentre outros. Utiliza
sistemas analíticos e ambientes de desenvolvimento para criação de
aplicações. Desenvolve a programação de robôs por demonstração. Na
integração entre robôs e outras máquinas e equipamentos, desenvolve a
programação de controladores lógicos programáveis (CLP); escreve ou
altera programas de máquinas-ferramentas com comando numérico
computadorizado (CNC); e programa parâmetros para acionamentos de
potência e de redes industriais, utilizando ferramentas de programação
específicas. Programa interfaces eletrônicas, para permitir a conexão
digital de robôs com servidores externos, tendo em vista a integração de
sistemas robóticos a grandes redes, como IIoT-Internet Industrial das
Coisas (Industrial Internet of Things) ou equivalente. Realiza a
manutenção de sistemas robóticos, avaliando gráficos de tendências e
relatórios de manutenção e planejando o processo e as etapas de trabalho
de manutenção. Planeja manutenção preventiva e preditiva. Analisa falhas
em sistemas. Executa procedimentos de manutenção. Avalia a eficácia das
soluções de manutenção implementadas. Colabora no processo de aquisição
de componentes, equipamentos e sistemas, verificando características
técnicas dos itens a serem adquiridos; selecionando fornecedores; e
acompanhando testes de funcionamento de máquinas e equipamentos para
elaboração e emissão de parecer técnico. Avalia disponibilidade de peças
de reposição. Coordena equipes de trabalho, formando equipe
multidisciplinar para análise de máquinas e equipamentos; reunindo-se
com equipes de trabalho; monitorando a execução de tarefas; e
propiciando suporte técnico aos integrantes de equipe. Treina usuários
em operação e manutenção de sistemas robóticos. Elabora documentação
técnica de projetos de sistemas robóticos. Documenta alterações de
projeto ocorridas durante a instalação de sistemas. Elabora relatório de
aceitação de equipamentos. Documenta o plano de ação de manutenção
preventiva e preditiva, as ações corretivas e as melhorias implementadas
nos sistemas. Mantém o local de trabalho limpo, seguro e organizado.
Conserva ferramentas, instrumentos e equipamentos. Utiliza equipamentos
de proteção.

\begin{Form}
\textbf{Esta correspondência está adequada?} \\ 
\ChoiceMenu[radio,radiosymbol=\ding{52},bordercolor=0 0.5 0,name=cbovalid_ 300110 ]{}{CBO deve ficar=sim, \hspace{6em} CBO não fica aqui=nao}
\end{Form}

\newpage
\section{ 02021 -- Técnico em Metalurgia }

\textbf{Perfil do Curso (CNCT)}

O Técnico em Metalurgia será habilitado para:

Realizar a gestão das etapas de obtenção e transformação de materiais
ferrosos e não ferrosos. Elaborar ensaios e análises químicas dos metais
e suas ligas, respeitando procedimentos e normas técnicas de qualidade,
de saúde e segurança e de meio ambiente. Controlar a execução dos
processos metalúrgicos de transformação térmica e mecânica dos
materiais. Interpretar e desenvolver projetos por meio de técnicas de
usinagem e soldagem. Reconhecer tecnologias inovadoras presentes no
segmento visando a atender às transformações digitais na sociedade.
Reconhecer os processos de manufatura aditiva empregados na metalurgia.

Para atuação como Técnico em Metalurgia, são fundamentais:

Conhecimentos e saberes relacionados aos processos de produção de metais
de modo a assegurar a saúde e a segurança dos trabalhadores.
Conhecimentos e saberes relacionados à sustentabilidade do processo
produtivo. Conhecimentos e saberes relacionados às técnicas e aos
processos de produção na metalurgia, às normas técnicas, à liderança de
equipes, à solução de problemas técnicos e trabalhistas e à gestão de
conflitos.

\textbf{CBOs correspondidos e seus perfis (presentes no CNCT, e presentes na predição da IA)}

\textbf{CBO: 314620  ( Técnico em soldagem )}

Planeja processos de soldagem -- em estruturas, chapas, perfis,
tubulações e outros artefatos industriais metalmecânicos --,
interpretando projetos e consultando normas técnicas. Define a aplicação
de processos diversificados de soldagem -- tais como eletrodo revestido;
soldagem por gás inerte de tungstênio-TIG (Tungsten Inert Gas); soldagem
com utilização de gás inerte para proteção da poça de fusão-MIG (Metal
Inert Gas); soldagem com utilização de gás ativo para proteção da poça
de fusão-MAG (Metal Active Gas); oxigás; arame tubular; arco submerso; e
plasma --, conforme requisitos de projetos e de posições dos produtos
soldados. Elabora o plano de qualificação dos procedimentos de soldagem.
Analisa desenhos de engenharia, projetos, especificações, esboços,
ordens de serviço e planilhas de dados de segurança de material para
planejar o leiaute e a montagem dos cenários das operações de soldagem.
Elabora o plano de soldagem. Define os procedimentos de execução e
inspeção. Estima o material de consumo e de aplicação. Seleciona os
procedimentos para cada atividade. Verifica as condições de segurança do
ambiente. Organiza equipamentos, ferramentas e materiais no local de
trabalho. Define a sequência de execução dos serviços de soldagem.
Elabora a programação da produção. Opera processos de soldagem,
determinando equipamentos e métodos de soldagem necessários e aplicando
conhecimentos de metalurgia, geometria, técnicas e posições de soldagem.
Monitora processos de soldagem, utilizando procedimentos e sistemas
específicos, tendo em vista a garantida da qualidade exigida para os
produtos soldados. Utiliza sistemas de gerenciamento de informações de
soldagem, para a melhoria da produtividade e da qualidade dos processos.
Desenvolve modelos para projetos de soldagem, usando cálculos
matemáticos baseados em informações dos desenhos de projeto. Ao receber
notificações de mau funcionamento de equipamentos ou de materiais
defeituosos, define as medidas necessárias para solução dos problemas
operacionais. Qualifica procedimentos de soldagem, identificando chapas
ou tubos de testes e acompanhando o processo de qualificação. Coordena a
realização de ensaios destrutivos e não destrutivos em produtos
soldados. Inspeciona matéria-prima e consumíveis de soldagem. Verifica
as condições de secagem, estocagem e manuseio de consumíveis de
soldagem. Acompanha e inspeciona a soldagem no decorrer dos processos.
Supervisiona processos de soldagem, tendo em vista a eficácia e a
eficiência dos serviços. Aplica normas técnicas na execução dos serviços
de soldagem. Coordena as equipes na execução dos serviços de soldagem,
de acordo com o cronograma e a programação da produção. Verifica a
montagem, para o cumprimento dos requisitos do projeto no que se refere
à soldagem. Controla o consumo do material. Controla condições de
secagem, estocagem e manuseio dos consumíveis de soldagem. Realiza o
controle dimensional. Verifica a conclusão da soldagem; as
identificações da soldagem; e os aspectos visual e dimensional da
soldagem. Coordena as operações de tratamento térmico. Libera o serviço
executado para inspeção do controle de qualidade. Pode supervisionar
equipes de trabalho, distribuindo tarefas, avaliando desempenho e
desenvolvendo treinamentos. Incorpora tecnologias inovadoras aos
processos de soldagem, tais como uso de soldagem a laser; e aplicações
para máquinas inteligentes, com princípios de automação e robótica.
Define técnicas, métodos, máquinas e equipamentos para operar com novos
materiais, como aços de alta resistência combinados com revestimentos
galvanizados; alumínio; aços inoxidáveis; e outros materiais não
ferrosos. Segue normas regulamentadoras e procedimentos de segurança e
saúde no trabalho em processos de soldagem, orientando e fiscalizando a
utilização de equipamentos de proteção individual e coletiva. Promove o
descarte e a destinação de resíduos de acordo com procedimentos e normas
de preservação ambiental.

\begin{Form}
\textbf{Esta correspondência está adequada?} \\ 
\ChoiceMenu[radio,radiosymbol=\ding{52},bordercolor=0 0.5 0,name=cbovalid_ 314620 ]{}{CBO deve ficar=sim, \hspace{6em} CBO não fica aqui=nao}
\end{Form}

\textbf{CBO: 314705  ( Técnico de acabamento em siderurgia )}

Planeja as atividades relacionadas aos processos siderúrgicos de
acabamento, descrevendo-as, definindo responsabilidades para cada etapa
dos processos e estabelecendo prazos de entrega de produtos e serviços.
Identifica necessidade de melhoria de produtos -- no âmbito dos
processos siderúrgicos de acabamento - e participa da elaboração de
projetos de fabricação de produtos siderúrgicos. Mantém contato com
clientes e fornecedores para troca de informações. Controla processos
siderúrgicos de acabamento, realizando cálculos de deformações mecânicas
e cálculos de área, volume e peso de produtos e equipamentos
siderúrgicos. Calcula o balanço de massa e as quantidades de fundentes
para adição aos processos. Pode utilizar sistemas simuladores de
processos siderúrgicos. Verifica desgaste de máquinas, equipamentos e
outros recursos de trabalho utilizados no acabamento siderúrgico.
Coordena paradas de equipamentos. Pode utilizar recursos de logística
integrada e software para monitoramento dos processos. Define itens de
controle da qualidade dos processos siderúrgicos de acabamento e
monitora resultados dos indicadores desses itens de controle. Inspeciona
produtos visualmente. Realiza ensaios físicos e químicos e interpreta
seus resultados. Interpreta resultados de ensaios mecânicos e
metalográficos. Avalia defeitos de produtos e analisa suas causas, a fim
de elaborar planos de ações preventivas e corretivas para os processos.
Analisa aproveitamento de produtos não conformes. Presta assistência
técnica para clientes, analisando suas reclamações e lhes sugerindo a
melhor forma de utilização de produtos. Pode adaptar processos às
necessidades específicas -- nível de qualidade, composição do produto,
entre outras -- de cada cliente, tendo em vista as tendências de
customização no setor siderúrgico. Elabora planos experimentais de novos
produtos e de implantação de novas tecnologias, fazendo uso de materiais
clássicos e de novos materiais na produção siderúrgica; introduzindo
equipamentos de maior eficiência energética; e utilizando fontes
alternativas de energia nos processos. Aplica tecnologias digitais e de
automatização nos processos, tendo em vista a eficiência, a
rentabilidade, a segurança e a redução de desperdícios. Divulga os
planos experimentais para os profissionais envolvidos nos processos.
Programa e monitora a execução das etapas de planos experimentais.
Avalia e valida resultados dos planos experimentais desenvolvidos.
Identifica atividades técnicas e operacionais, elabora manuais de
procedimentos, normaliza e revisa padrões e procedimentos dos processos
siderúrgicos de acabamento. Interpreta desenhos técnicos. Elabora
relatórios técnicos para documentação dos processos siderúrgicos de
acabamento, tais como relatórios técnicos de produção; relatórios de
planos experimentais; relatórios de visitas e de assistência técnica;
relatórios de controle de qualidade; e relatórios de reclamação de
clientes. Elabora informes técnicos dos processos. Promove os meios para
o desenvolvimento profissional da equipe de trabalho -- com focalização
nos processos siderúrgicos de acabamento --, identificando necessidades
de desenvolvimento de pessoal, para elaborar o plano de desenvolvimento.
Ministra treinamento e/ou solicita contratação de terceiros para
realização dos programas. Avalia a eficácia dos treinamentos
ministrados. Providencia serviços -- inclusive terceirizados -- de
manutenção de máquinas, equipamentos e outros recursos utilizados nos
processos siderúrgicos de acabamento. Mantém o ambiente de trabalho
limpo e organizado. Aplica normas e procedimentos de segurança,
utilizando equipamentos de proteção individual e coletiva e
desenvolvendo ações de prevenção de acidentes de trabalho. Promove o
descarte e a destinação de resíduos de acordo com diretrizes
normatizadas e procedimentos, a fim de manter os processos alinhados aos
conceitos e às práticas do desenvolvimento sustentável.

\begin{Form}
\textbf{Esta correspondência está adequada?} \\ 
\ChoiceMenu[radio,radiosymbol=\ding{52},bordercolor=0 0.5 0,name=cbovalid_ 314705 ]{}{CBO deve ficar=sim, \hspace{6em} CBO não fica aqui=nao}
\end{Form}

\textbf{CBO: 314710  ( Técnico de aciaria em siderurgia )}

Planeja as atividades relacionadas aos processos siderúrgicos de
aciaria, descrevendo-as, definindo responsabilidades para cada etapa dos
processos e estabelecendo prazos de entrega de produtos e serviços.
Identifica necessidade de melhoria de produtos -- no âmbito da aciaria -
e participa da elaboração de projetos de fabricação de produtos
siderúrgicos. Mantém contato com clientes e fornecedores para troca de
informações. Controla processos siderúrgicos de aciaria, realizando
cálculos de volume e peso de produtos e equipamentos siderúrgicos.
Calcula o balanço térmico dos processos e calcula as quantidades de
produtos para adição, a fim de realizar a correção química dos
processos. Verifica desgaste de máquinas, equipamentos e outros recursos
de trabalho utilizados na aciaria. Coordena paradas de equipamentos.
Pode utilizar recursos de logística integrada e software para
monitoramento dos processos. Define itens de controle da qualidade dos
processos de aciaria e monitora resultados dos indicadores desses itens
de controle. Inspeciona produtos visualmente. Realiza ensaios físicos e
químicos e interpreta seus resultados. Interpreta resultados de ensaios
mecânicos e metalográficos. Avalia defeitos de produtos e analisa suas
causas, a fim de elaborar planos de ações preventivas e corretivas para
os processos. Analisa aproveitamento de produtos não conformes. Presta
assistência técnica para clientes, analisando suas reclamações e lhes
sugerindo a melhor forma de utilização de produtos. Pode adaptar
processos às especificações de clientes, a partir da identificação de
suas necessidades. Elabora planos experimentais de novos produtos e de
implantação de novas tecnologias, introduzindo equipamentos de maior
eficiência energética; utilizando fontes alternativas de energia nos
processos; e aplicando tecnologias digitais e de automatização em todas
as operações de aciaria. Divulga os planos experimentais para os
profissionais envolvidos nos processos. Programa e monitora a execução
das etapas de planos experimentais. Avalia e valida resultados dos
planos experimentais desenvolvidos. Acompanha as inovações em aciaria, a
fim de implantar modernizações em conversores a oxigênio; automação em
sistemas de carregamento de gusa e de vazamento de aço; projetos de
melhoria da qualidade final do aço; e inovações nas áreas de corrida.
Pode desenvolver processos de aciaria com uso de fornos elétricos, para
produção de aço a partir de diferentes cargas de matéria-prima, tais
como ferro-gusa, ferro-esponja, finos de ferro-esponja, entre outras.
Identifica atividades técnicas e operacionais, elabora manuais de
procedimentos, normaliza e revisa padrões e procedimentos dos processos
de aciaria. Interpreta desenhos técnicos. Elabora relatórios técnicos
para documentação dos processos siderúrgicos de aciaria, tais como
relatórios técnicos de produção; relatórios de planos experimentais;
relatórios de visitas e de assistência técnica; relatórios de controle
de qualidade; e relatórios de reclamação de clientes. Elabora informes
técnicos dos processos. Promove os meios para desenvolvimento
profissional da equipe de trabalho, com focalização nos processos de
aciaria. Identifica necessidades de desenvolvimento de pessoal, para
elaborar plano de desenvolvimento. Ministra treinamento e/ou solicita
contratação de terceiros para realizar os programas. Avalia a eficácia
dos treinamentos ministrados. Providencia serviços -- inclusive
terceirizados -- de manutenção de máquinas, equipamentos e outros
recursos utilizados em processos de aciaria. Mantém o ambiente de
trabalho limpo e organizado. Aplica normas e procedimentos de segurança,
utilizando equipamentos de proteção individual e coletiva e
desenvolvendo ações de prevenção de acidentes de trabalho. Promove o
descarte e a destinação de resíduos de acordo com diretrizes
normatizadas e procedimentos, para manter os processos alinhados aos
conceitos e às práticas do desenvolvimento sustentável.

\begin{Form}
\textbf{Esta correspondência está adequada?} \\ 
\ChoiceMenu[radio,radiosymbol=\ding{52},bordercolor=0 0.5 0,name=cbovalid_ 314710 ]{}{CBO deve ficar=sim, \hspace{6em} CBO não fica aqui=nao}
\end{Form}

\textbf{CBO: 314720  ( Técnico de laminação em siderurgia )}

Planeja as atividades relacionadas aos processos siderúrgicos de
laminação, descrevendo-as, definindo responsabilidades para cada etapa
dos processos e estabelecendo prazos de entrega de produtos e serviços.
Identifica necessidade de melhoria de produtos -- no âmbito dos
processos siderúrgicos de laminação - e participa da elaboração de
projetos de fabricação de produtos siderúrgicos. Mantém contato com
clientes e fornecedores para troca de informações. Controla processos
siderúrgicos de laminação, realizando cálculos de deformações mecânicas
e cálculos de área, volume e peso de produtos e equipamentos
siderúrgicos. Pode utilizar sistemas simuladores de processos
siderúrgicos. Verifica desgaste de máquinas, equipamentos e outros
recursos de trabalho utilizados na laminação siderúrgica. Coordena
paradas de equipamentos. Pode utilizar recursos de logística integrada e
software para monitoramento dos processos. Define itens de controle da
qualidade dos processos siderúrgicos de laminação e monitora resultados
dos indicadores desses itens de controle. Inspeciona produtos
visualmente. Realiza ensaios físicos e químicos e interpreta seus
resultados. Interpreta resultados de ensaios mecânicos e metalográficos.
Avalia defeitos de produtos e analisa suas causas, a fim de elaborar
planos de ações preventivas e corretivas para os processos. Analisa
aproveitamento de produtos não conformes. Presta assistência técnica
para clientes, analisando suas reclamações e lhes sugerindo a melhor
forma de utilização de produtos. Pode adaptar processos às necessidades
específicas - nível de qualidade, composição do produto, entre outras -
de cada cliente, tendo em vista as tendências de customização no setor
siderúrgico. Elabora planos experimentais de novos produtos e de
implantação de novas tecnologias, introduzindo inovações nos processos e
aplicando tecnologias digitais e de automatização, a fim de obter
controle mais acurado e processos mais flexíveis. Divulga os planos
experimentais para os profissionais envolvidos nos processos. Programa e
monitora a execução das etapas dos planos experimentais. Avalia e valida
resultados dos planos experimentais desenvolvidos. Identifica atividades
técnicas e operacionais, elabora manuais de procedimentos, normaliza e
revisa padrões e procedimentos dos processos siderúrgicos de laminação.
Interpreta desenhos técnicos. Elabora relatórios técnicos para
documentação dos processos siderúrgicos de laminação, tais como
relatórios técnicos de produção; relatórios de planos experimentais;
relatórios de visitas e de assistência técnica; relatórios de controle
de qualidade; e relatórios de reclamação de clientes. Elabora informes
técnicos dos processos. Promove os meios para o desenvolvimento
profissional da equipe de trabalho -- com focalização nos processos
siderúrgicos de laminação --, identificando necessidades de
desenvolvimento de pessoal, para elaborar o plano de desenvolvimento.
Ministra treinamento e/ou solicita contratação de terceiros para
realização dos programas. Avalia a eficácia dos treinamentos
ministrados. Providencia serviços -- inclusive terceirizados -- de
manutenção de máquinas, equipamentos e outros recursos utilizados nos
processos siderúrgicos de laminação. Mantém o ambiente de trabalho limpo
e organizado. Aplica normas e procedimentos de segurança em processos
siderúrgicos de laminação, utilizando equipamentos de proteção
individual e coletiva e desenvolvendo ações de prevenção de acidentes de
trabalho. Promove o descarte e a destinação de resíduos de acordo com
diretrizes normatizadas e procedimentos, a fim de manter os processos
alinhados aos conceitos e às práticas do desenvolvimento sustentável.

\begin{Form}
\textbf{Esta correspondência está adequada?} \\ 
\ChoiceMenu[radio,radiosymbol=\ding{52},bordercolor=0 0.5 0,name=cbovalid_ 314720 ]{}{CBO deve ficar=sim, \hspace{6em} CBO não fica aqui=nao}
\end{Form}

\textbf{CBO: 314725  ( Técnico de redução na siderurgia (primeira fusão) )}

Planeja as atividades relacionadas aos processos siderúrgicos de
redução, descrevendo-as, definindo responsabilidades para cada etapa dos
processos e estabelecendo prazos de entrega de produtos e serviços.
Identifica necessidade de melhoria de produtos -- no âmbito de processos
siderúrgicos de redução - e participa da elaboração de projetos de
fabricação de produtos siderúrgicos. Mantém contato com clientes e
fornecedores para troca de informações. Controla processos siderúrgicos
de redução, realizando cálculos de área, volume e peso de produtos e
equipamentos siderúrgicos. Calcula o balanço de massa e as quantidades
de fundentes para adição aos processos. Calcula o balanço térmico dos
processos e calcula as quantidades de produtos para adição, a fim de
realizar a correção química dos processos. Pode utilizar sistemas
simuladores de processos siderúrgicos. Verifica desgaste de máquinas,
equipamentos e outros recursos de trabalho utilizados nos processos de
redução siderúrgica. Coordena paradas de equipamentos. Define itens de
controle da qualidade dos processos siderúrgicos de redução e monitora
resultados dos indicadores desses itens de controle. Inspeciona produtos
visualmente. Realiza ensaios físicos e químicos e interpreta seus
resultados. Interpreta resultados de ensaios mecânicos e metalográficos.
Avalia defeitos de produtos e analisa suas causas, a fim de elaborar
planos de ações preventivas e corretivas para os processos. Analisa
aproveitamento de produtos não conformes. Presta assistência técnica
para clientes, analisando suas reclamações e lhes sugerindo a melhor
forma de utilização de produtos. Pode adaptar processos às
especificações de clientes, a partir da identificação de suas
necessidades. Elabora planos experimentais de novos produtos e de
implantação de novas tecnologias, introduzindo equipamentos de melhor
desempenho técnico-ambiental e aplicando tecnologias de instrumentação,
automação e controle de processos, tendo em vista o incremento da
produtividade. Divulga os planos experimentais para os profissionais
envolvidos nos processos. Programa e monitora a execução das etapas de
referidos planos. Avalia e valida resultados dos planos experimentais
desenvolvidos. Identifica atividades técnicas e operacionais, elabora
manuais de procedimentos, normaliza e revisa padrões e procedimentos dos
processos siderúrgicos de redução. Interpreta desenhos técnicos. Elabora
relatórios técnicos para documentação dos processos siderúrgicos de
redução, tais como relatórios técnicos de produção; relatórios de planos
experimentais; relatórios de visitas e de assistência técnica;
relatórios de controle de qualidade; e relatórios de reclamação de
clientes. Elabora informes técnicos dos processos. Promove os meios para
desenvolvimento profissional da equipe de trabalho, com focalização nos
processos siderúrgicos de redução. Identifica necessidades de
desenvolvimento de pessoal, para elaborar o plano de desenvolvimento.
Ministra treinamento e/ou solicita contratação de terceiros para
realização dos programas. Avalia a eficácia dos treinamentos
ministrados. Providencia serviços -- inclusive terceirizados -- de
manutenção de máquinas, equipamentos e outros recursos utilizados em
processos de fundição. Atua de acordo com as diretrizes adotadas pela
indústria siderúrgica para o setor de redução, a fim de atender
exigências de órgãos fiscalizadores e requisitos de clientes e de
mercado, no que se refere às condições do local de trabalho, ao meio
ambiente e ao desenvolvimento sustentável. Aplica normas e procedimentos
de segurança no trabalho em processos siderúrgicos de redução,
utilizando equipamentos de proteção individual e coletiva e
desenvolvendo ações de prevenção de acidentes de trabalho. Promove o
descarte e a destinação de resíduos de acordo com diretrizes
normatizadas e procedimentos.

\begin{Form}
\textbf{Esta correspondência está adequada?} \\ 
\ChoiceMenu[radio,radiosymbol=\ding{52},bordercolor=0 0.5 0,name=cbovalid_ 314725 ]{}{CBO deve ficar=sim, \hspace{6em} CBO não fica aqui=nao}
\end{Form}

\textbf{CBO: 314730  ( Técnico de refratário em siderurgia )}

Planeja as atividades que integram os serviços relacionados aos
materiais refratários, descrevendo-as, definindo responsabilidades para
cada etapa das atividades e estabelecendo prazos de entrega de serviços.
Identifica necessidade de melhoria de produtos -- no âmbito dos serviços
relacionados aos materiais refratários - e participa da elaboração de
projetos de fabricação de produtos siderúrgicos. Mantém contato com
clientes e fornecedores para troca de informações. Em função dos
materiais refratários utilizados nos processos siderúrgicos, realiza
cálculos de área, volume e peso de produtos e equipamentos siderúrgicos;
calcula o balanço de massa e as quantidades de fundentes para adição aos
processos; e calcula as quantidades de produtos para adição, a fim de
realizar a correção química dos processos. Pode utilizar sistemas
simuladores de processos siderúrgicos. Verifica desgaste de máquinas,
equipamentos e outros recursos de trabalho utilizados nos serviços.
Coordena paradas de equipamentos. Define itens de controle da qualidade
dos serviços relacionados aos materiais refratários e monitora
resultados dos indicadores desses itens de controle. Inspeciona
visualmente os materiais refratários, tendo em vista a eficácia dos
serviços. Realiza ensaios físicos e químicos e interpreta seus
resultados. Interpreta resultados de ensaios mecânicos e metalográficos.
Avalia defeitos de produtos e analisa suas causas, a fim de elaborar
planos de ações preventivas e corretivas para os processos, com
focalização nos materiais refratários. Presta assistência técnica para
clientes -- internos e externos --, adaptando serviços às especificações
de clientes, a partir da identificação de suas necessidades. Elabora
planos experimentais de novos produtos e de implantação de novas
tecnologias, fazendo uso de materiais refratários tecnologicamente
avançados; utilizando novos tipos de produtos refratários; selecionando
e desenvolvendo fornecedores de materiais refratários. Divulga os planos
experimentais para os profissionais envolvidos nos processos. Programa e
monitora a execução das etapas dos planos experimentais. Avalia e valida
resultados dos planos experimentais desenvolvidos. Realiza experiências
e testes com materiais refratários. Desenvolve pesquisas e projetos de
aplicação de revestimentos refratários para altos-fornos e outros
equipamentos siderúrgicos. Planeja e desenvolve planos de manutenção --
preditiva, preventiva e corretiva -- de materiais refratários, para
todos os processos e etapas da produção siderúrgica. Identifica
atividades técnicas e operacionais, elabora manuais de procedimentos,
normaliza e revisa padrões e procedimentos dos serviços relacionados aos
materiais refratários utilizados na produção siderúrgica. Interpreta
desenhos técnicos. Elabora relatórios técnicos para documentação dos
serviços relacionados aos materiais refratários, tais como relatórios de
planos experimentais; relatórios de visitas e de assistência técnica;
relatórios de manutenção; e relatórios de pesquisas e projetos. Elabora
informes técnicos dos serviços. Promove os meios para o desenvolvimento
profissional da equipe de trabalho, com focalização nos serviços
relacionados aos materiais refratários. Identifica necessidades de
desenvolvimento de pessoal, para elaborar o plano de desenvolvimento.
Ministra treinamento e/ou solicita contratação de terceiros para
realização dos programas. Avalia a eficácia dos treinamentos
ministrados. Aplica normas e procedimentos de segurança no trabalho,
utilizando equipamentos de proteção individual e coletiva e
desenvolvendo ações de prevenção de acidentes de trabalho. Promove o
descarte e a destinação de resíduos de acordo com diretrizes
normatizadas e procedimentos, a fim de manter os processos alinhados aos
conceitos e práticas do desenvolvimento sustentável.

\begin{Form}
\textbf{Esta correspondência está adequada?} \\ 
\ChoiceMenu[radio,radiosymbol=\ding{52},bordercolor=0 0.5 0,name=cbovalid_ 314730 ]{}{CBO deve ficar=sim, \hspace{6em} CBO não fica aqui=nao}
\end{Form}

\newpage
\section{ 02022 -- Técnico em Metrologia }

\textbf{Perfil do Curso (CNCT)}

O Técnico em Metrologia será habilitado para:

Planejar, controlar e executar manutenção de sistemas, equipamentos,
métodos e padrões de medição. Controlar e assegurar a qualidade e o
correto funcionamento dos instrumentos de medição ou medidas
materializadas. Empregar técnicas e conceitos metrológicos na indústria
e no setor de serviços. Supervisionar e realizar perícia metrológica no
campo da metrologia legal. Realizar ensaios em instrumentos de medição
usados nas indústrias e a calibração de padrões nas áreas de acústica,
vibrações, mecânica, elétrica, telecomunicações, térmica, química,
materiais, óptica, vazão, temperatura, pressão e nível atendendo às
normas e aos padrões técnicos de qualidade, saúde e segurança e de meio
ambiente. Reconhecer tecnologias inovadoras presentes no segmento
visando a atender às transformações digitais na sociedade.

Para atuação como Técnico em Metrologia, são fundamentais:

Conhecimentos e saberes relacionados aos processos de planejamento e
manutenção em instrumentos de medição de modo a assegurar a saúde e a
segurança dos trabalhadores e dos consumidores. Conhecimentos e saberes
relacionados à sustentabilidade do processo produtivo, às técnicas e aos
processos de calibração, às normas técnicas, à liderança de equipes, à
solução de problemas técnicos e trabalhistas e à gestão de conflitos.

\textbf{CBOs correspondidos e seus perfis (presentes no CNCT, e presentes na predição da IA)}

\textbf{CBO: 313405  ( Técnico em calibração )}

Prepara o serviço de calibração, identificando especificações técnicas,
monitorando as condições ambientais para a calibração, inspecionando
visualmente os itens a serem calibrados e determinando o procedimento de
calibração. Realiza aferição e calibração de padrões, equipamentos,
sistemas e instrumentos de medição e controle. Calcula e interpreta os
resultados das medições realizadas no processo, comparando com padrões,
para realizar os ajustes necessários nos equipamentos, sistemas e
instrumentos. Realiza medições de grandezas, lendo e interpretando
desenhos técnicos, selecionando métodos e princípios de medição. Executa
ensaios físicos, elétricos, mecânicos e eletrônicos, dentre outros,
correspondentes às grandezas a serem medidas. Opera padrões,
equipamentos, sistemas e instrumentos de medição e controle. Emite
laudos e certificados da calibração de equipamentos e instrumentos de
medição e controle. Arquiva documentação técnica e gerencial. Participa
da gestão dos procedimentos do sistema de confiabilidade metrológica de
laboratórios de calibração credenciados, elaborando sistema de
codificação, estabelecendo frequência, controlando prazos e validando
resultados da calibração. Participa de auditorias internas e externas.
Participa da gestão da documentação técnica de laboratórios
credenciados, elaborando, atualizando e analisando criticamente
procedimentos e instruções de trabalho, cadastrando instrumentos de
medição e controle, elaborando fichas e formulários de registros.
Contribui no desenvolvimento de projetos de sistemas de medição e
controle, identificando as variáveis envolvidas no processo e avaliando
o desempenho de instrumentos e sistemas em desenvolvimento, realizando a
calibração de protótipos e novos produtos. Realiza manutenção de
instrumentos de medição e controle, em bancada, para fins de calibração,
identificando disfunções e suas causas, relacionando custos e benefícios
da manutenção, reparando ou substituindo componentes danificados.
Realiza a calibração dos instrumentos após os trabalhos de manutenção.

\begin{Form}
\textbf{Esta correspondência está adequada?} \\ 
\ChoiceMenu[radio,radiosymbol=\ding{52},bordercolor=0 0.5 0,name=cbovalid_ 313405 ]{}{CBO deve ficar=sim, \hspace{6em} CBO não fica aqui=nao}
\end{Form}

\textbf{CBO: 352305  ( Metrologista )}

Realiza manutenção de instrumentos, equipamentos, sistemas e padrões
metrológicos, em bancada ou no local das instalações, para fins de
calibração e garantia da qualidade metrológica. Repara ou substitui
componentes danificados. Identifica disfunções e suas causas. Realiza a
calibração dos instrumentos usados nos trabalhos de manutenção.
Relaciona custos e benefícios da manutenção. Realiza calibração de
instrumentos, equipamentos, sistemas e padrões metrológicos - incluindo
sensores -, em laboratórios credenciados, seguindo metodologias
específicas, identificando especificações técnicas, monitorando as
condições ambientais para a calibração, inspecionando visualmente os
itens a serem calibrados e determinando o procedimento de calibração.
Realiza os ajustes necessários nos instrumentos, equipamentos e
sistemas. Participa da gestão da documentação técnica e da gestão dos
procedimentos do sistema de confiabilidade metrológica de laboratórios
de calibração credenciados. Realiza ensaios, testes e análises, seguindo
a legislação pertinente, para aferir o desempenho de produtos e serviços
em face de padrões metrológicos. Efetua testes prévios e seletivos de
produtos pré-medidos, coletando amostras nos pontos de venda e
realizando análises laboratoriais, para certificar de que a quantidade
anotada no rótulo corresponde ao conteúdo da embalagem. Realiza testes
físicos, óticos, acústicos, elétricos, mecânicos e químicos,
relacionados à qualidade de produtos e serviços. Arqueia e afere tanques
de fluido fixo ou móvel, utilizando instrumentos de medição e padrões de
volume, para atestar a conformidade do reservatório com as
especificações. Colabora na fiscalização metrológica de instrumentos,
medidas materializadas, produtos e serviços, examinando e verificando
instrumentos e medidas materializadas e vistoriando oficinas de
manutenção. Pode colaborar em perícias metrológicas, examinando e
aferindo a totalidade das condições de instrumentos de medição e medida
materializada, determinando suas qualidades metrológicas. Supervisiona
atividades metrológicas, inclusive as atividades específicas de
laboratórios credenciados. Supervisiona e coordena equipe, programando e
acompanhando as suas atividades e avaliando os trabalhos executados.
Disponibiliza meios para a execução das atividades e controla
instrumentos, equipamentos, sistemas e materiais utilizados nas
atividades. Participa de grupos de estudos, especialmente sobre
adequação dos sistemas de medição aos conceitos da manufatura avançada.
Capacita recursos humanos na área de metrologia, elaborando material
didático, preparando aulas, ministrando treinamentos internos e
avaliando o desempenho dos treinandos.

\begin{Form}
\textbf{Esta correspondência está adequada?} \\ 
\ChoiceMenu[radio,radiosymbol=\ding{52},bordercolor=0 0.5 0,name=cbovalid_ 352305 ]{}{CBO deve ficar=sim, \hspace{6em} CBO não fica aqui=nao}
\end{Form}

\textbf{CBO: 352310  ( Agente fiscal de qualidade )}

Fiscaliza instrumentos metrológicos, medidas materializadas, produtos,
marcas de conformidade e serviços, examinando aspectos formais com
vistas à existência de todas as informações obrigatórias e verificando a
marca de conformidade. Realiza fiscalizações em operações especiais
(blitz). Coleta produtos, notifica o fiscalizado e verifica o
cumprimento da notificação. Acompanha liberações. Autua o responsável
por infração. Interdita produtos e instrumentos e os apreende, cautelar
e definitivamente. Inutiliza produtos irregulares. Coloca marca de
verificação em instrumentos e medidas materializadas. Realiza ensaios,
testes e análises, seguindo a legislação pertinente a cada produto.
Efetua testes prévios e seletivos de produtos pré-medidos, coletando
amostras nos pontos de venda e realizando análises laboratoriais, para
certificar que a quantidade anotada no rótulo corresponde ao conteúdo da
embalagem. Realiza testes físicos, óticos, acústicos, elétricos,
mecânicos e químicos relacionados à qualidade de produtos e serviços.
Inspeciona segurança no transporte de cargas perigosas. Supervisiona as
atividades de fiscalização da qualidade, organizando operações especiais
e encaminhando produtos para ensaios de laboratório. Supervisiona e
coordena equipe, programando e acompanhando as suas atividades e
avaliando os trabalhos executados. Disponibiliza meios para a execução
das atividades e controla equipamentos, instrumentos e materiais
utilizados na fiscalização. Participa de grupos de estudos. Capacita
recursos humanos na área da fiscalização da qualidade, elaborando
material didático, preparando aulas, ministrando treinamentos internos e
avaliando desempenho dos treinandos. Ministra palestras e colabora na
disseminação de conceitos metrológicos e da qualidade. Orienta
tecnicamente o público -- consumidor e fiscalizado --, atendendo
denúncias, analisando reclamações, informando sobre procedimentos
técnicos e administrativos e prestando esclarecimentos sobre legislação.
Participa da elaboração de material de orientação. Registra o processo
de fiscalização, emitindo notificações, autos, pareceres técnicos,
termos de reprovação e certificados de aprovação em inspeção. Preenche
fichas cadastrais e relatórios diários.

\begin{Form}
\textbf{Esta correspondência está adequada?} \\ 
\ChoiceMenu[radio,radiosymbol=\ding{52},bordercolor=0 0.5 0,name=cbovalid_ 352310 ]{}{CBO deve ficar=sim, \hspace{6em} CBO não fica aqui=nao}
\end{Form}

\textbf{CBO: 352315  ( Agente fiscal metrológico )}

Realiza fiscalização metrológica de instrumentos, medidas
materializadas, produtos e serviços, examinando os aspectos formais,
conferindo a existência de todas as informações obrigatórias, analisando
aspectos quantitativos e verificando a fidelidade de informações.
Examina instrumentos e medidas materializadas e vistoria oficinas de
manutenção. Realiza fiscalizações em operações especiais (blitz). Coleta
produtos, notifica o fiscalizado e verifica o cumprimento da
notificação. Acompanha liberações. Autua o responsável por infração.
Interdita produtos e instrumentos e os apreende, cautelar e
definitivamente. Inutiliza produtos irregulares. Realiza testes e
análises, rastreando padrões metrológicos e realizando conservação de
padrões e instrumentos metrológicos. Testa produtos e instrumentos in
loco, realizando ensaios e testes de acordo com a legislação, elaborando
laudos técnicos metrológicos e laudos periciais de instrumentos e
produtos fraudados. Efetua testes prévios e seletivos de produtos
pré-medidos, coletando amostras nos pontos de venda, realizando análises
laboratoriais, para certificar que a quantidade anotada no rótulo
corresponde ao conteúdo da embalagem. Verifica instrumentos e medidas
materializadas, realizando verificações inicial, periódica e eventual,
lacrando o instrumento e colocando marca de verificação. Cobra taxas de
serviços metrológicos. Emite documento para recebimento do valor de
pagamento pelos serviços prestados. Supervisiona as atividades de
fiscalização metrológica, organizando operações especiais. Supervisiona
e coordena equipe, programando e acompanhando as suas atividades e
avaliando os trabalhos executados. Disponibiliza meios para a execução
das atividades e controla equipamentos, instrumentos e materiais
utilizados na fiscalização. Participa de grupos de estudos. Capacita
recursos humanos na área de fiscalização metrológica, elaborando
material didático, preparando aulas, ministrando treinamentos internos e
avaliando desempenho dos treinandos. Ministra palestras e colabora na
disseminação de conceitos de fiscalização metrológica. Orienta
tecnicamente o público -- consumidor e fiscalizado -- informando sobre
procedimentos técnicos e administrativos e prestando esclarecimentos
sobre legislação. Participa da elaboração de material de orientação.
Atende denúncias e reclamações. Registra o processo de fiscalização e
verificação, emitindo certificados de verificação, laudos, termos,
intimações, notificações, autos e pareceres técnicos. Preenche relatório
de verificação, fichas cadastrais e relatórios diários.

\begin{Form}
\textbf{Esta correspondência está adequada?} \\ 
\ChoiceMenu[radio,radiosymbol=\ding{52},bordercolor=0 0.5 0,name=cbovalid_ 352315 ]{}{CBO deve ficar=sim, \hspace{6em} CBO não fica aqui=nao}
\end{Form}

\newpage
\section{ 02023 -- Técnico em Outros - Eixo Controle e Processos Industriais }

\textbf{Perfil do Curso (CNCT)}

O Técnico em Ferramentaria será habilitado para:

Elaborar, desenvolver e executar projetos de estampos para a fabricação
de peças por processos de conformação de metais, moldes para a
fabricação de produtos poliméricos, dispositivos e gabaritos. Programar,
controlar e executar processos de manutenção de moldes de injeção e/ou
sopro e de estampos de acordo com normas técnicas, ambientais, de saúde
e segurança no trabalho e de qualidade. Zelar pelo atendimento de
questões ambientais específicas relacionadas ao processamento de
materiais poliméricos e chapas metálicas. Controlar defeitos e garantir
a qualidade dos produtos fabricados a partir de matérias-primas
poliméricas e chapas metálicas. Simular processos de injeção de
polímeros e conformação de chapas pela utilização de softwares
específicos. Reconhecer tecnologias inovadoras presentes no segmento
visando a atender às transformações digitais na sociedade.

Para atuação como Técnico em Ferramentaria, são fundamentais:

Conhecimentos e saberes relacionados aos processos de planejamento,
confecção e manutenção de moldes e estampos de modo a assegurar a saúde
e a segurança dos trabalhadores e futuros usuários. Conhecimentos e
saberes relacionados à sustentabilidade do processo produtivo, às
técnicas e aos processos de produção, às normas técnicas, à liderança de
equipes, à solução de problemas técnicos e trabalhistas e à gestão de
conflitos.

\textbf{CBOs correspondidos e seus perfis (presentes no CNCT, e presentes na predição da IA)}

\textbf{CBO: 721105  ( Ferramenteiro )}

Planeja o processo de construção, utilizando normas técnicas e
especificações e analisando desenhos do projeto. Segue os desenhos e
desenvolve protótipos. Visualiza e calcula dimensões, tamanhos, formas e
tolerâncias de montagens, com base nas especificações. Define a
sequência de operações, seleciona máquinas e ferramentas e identifica o
material a ser utilizado. Seleciona os metais e ligas a serem usados com
base em propriedades como dureza ou tolerância ao calor. Verifica o
emprego de itens padronizados, estima tempos das etapas e segue o
cronograma do processo. Projeta gabaritos, acessórios e modelos para uso
como auxiliar de trabalho na fabricação de peças ou produtos. Durante o
planejamento, pode identificar incorreções no projeto e realizar
modificações, para correção. Pode auxiliar no desenvolvimento e no
projeto de novas ferramentas e matrizes, usando software de design e
manufatura auxiliados por computador e simuladores. Constrói ferramentas
e dispositivos, organizando a distribuição do material, consultando
catálogos técnicos, utilizando parâmetros de corte e definindo a
sequência de montagem e construção. Elabora croqui de peças e modifica
projetos. Pode construir ferramentas para usinagem EDM -- Usinagem por
descarga elétrica. Afia ferramentas de corte. Configura e opera
máquinas-ferramentas convencionais ou a comando numérico computadorizado
(CNC) de múltiplos eixos -- como tornos, fresadoras, retificadoras ou
centros de usinagem -- para cortar, furar, retificar ou modelar peças de
acordo com as dimensões e acabamentos prescritos. Pode configurar e
operar máquinas de impressão 3D para construção de peças por manufatura
aditiva. Configura e opera prensas para perfurar e fazer orifícios nas
peças para montagem. Corta, modela e apara material metálico para obter
dimensões ou formas especificadas, usando serras e tesouras elétricas e
ferramentas manuais. Executa o tratamento térmico de materiais. Conforma
a quente, utilizando equipamentos oxiacetilênicos. Solda peças e
conjuntos. Instala circuitos hidráulicos e pneumáticos. Dá acabamento na
superfície de machos e matrizes e faz o ajuste final (try-out).
Desenvolve ferramentas especiais ou manuais e de corte para máquinas,
constrói dispositivos para usinagem e de controle do produto. Constrói
estampos de corte, dobra, repuxo e corte fino, limando punções e
matrizes, perfilando punções e trefilando matrizes. Determina folgas
entre punções e matrizes. Executa eletroerosão a fio e por penetração.
Constrói moldes de sopro, de injeção e de eletroerosão, confeccionando
eletrodos para eletroerosão, usinando canais de injeção e furando o
conjunto extrator. Usina circuitos de refrigeração. Ajusta câmara a
quente, gavetas, postiços e núcleos rotativos. Faz o ajuste de
fechamento do molde. Constrói modelos e moldes metálicos para fundição,
projetando caixa de machos, coquilhas e modelos. Pode copiar modelos,
utilizando escâneres industriais tridimensionais e outras máquinas
especiais similares. Conforma modelos e moldes a quente. Usina canais de
alimentação e respiro. Confecciona guias. Ajusta fechamento de
coquilhas. Executa manutenção preventiva e corretiva, identificando
falhas no ferramental, desmontando/ montando e avaliando os desgastes
nos estampos, moldes e dispositivos. Inspeciona matrizes acabadas quanto
à uniformidade das superfícies, conformidade do contorno e defeitos.
Administra tempos de execução das etapas da manutenção. Faz controle
dimensional, utilizando equipamentos de medição como pinças, blocos de
medida, paquímetros e micrômetros. Traça material a ser usinado. Mede
dureza de componentes submetidos a tratamento térmico. Controla peças
usinadas e produtos acabados, verificando dimensões, alinhamentos e
folgas quanto à conformidade com as especificações e fazendo ajustes,
quando necessário. Pode configurar e operar máquinas de medição
tridimensional por coordenadas. Conserva o local de trabalho limpo e
organizado. Atua na prevenção de acidentes.

\begin{Form}
\textbf{Esta correspondência está adequada?} \\ 
\ChoiceMenu[radio,radiosymbol=\ding{52},bordercolor=0 0.5 0,name=cbovalid_ 721105 ]{}{CBO deve ficar=sim, \hspace{6em} CBO não fica aqui=nao}
\end{Form}

\textbf{CBO: 721110  ( Ferramenteiro de mandris, calibradores e outros dispositivos )}

Planeja o processo de construção, utilizando normas técnicas e
especificações e analisando desenhos do projeto. Segue desenhos e
desenvolve protótipos. Visualiza e calcula dimensões, tamanhos, formas e
tolerâncias de montagens, com base nas especificações. Define a
sequência de operações, seleciona máquinas e ferramentas e identifica o
material a ser utilizado na construção. Seleciona os metais e ligas a
serem usados com base em propriedades como dureza ou tolerância ao
calor. Verifica o emprego de itens padronizados, estima tempos das
etapas e segue o cronograma do processo de construção. Durante o
planejamento, pode identificar incorreções no projeto e realizar
modificações, para correção. Pode auxiliar no desenvolvimento e no
projeto de novas ferramentas e dispositivos, usando software de design e
manufatura auxiliados por computador e simuladores. Constrói ferramentas
e dispositivos mecânicos de usinagem, organizando a distribuição do
material, consultando catálogos técnicos, utilizando parâmetros de corte
e definindo a sequência de montagem e construção. Elabora croqui de
peças e modifica projetos. Afia ferramentas de corte. Configura e opera
máquinas-ferramentas convencionais ou a comando numérico computadorizado
(CNC) de múltiplos eixos -- como tornos, fresadoras, retificadoras ou
centros de usinagem -- para cortar, furar, retificar ou modelar peças de
acordo com as dimensões e acabamentos prescritos. Pode configurar e
operar máquinas de impressão 3D para construção de peças por manufatura
aditiva. Configura e opera prensas para perfurar e fazer orifícios nas
peças para montagem. Corta, modela e apara material metálico para obter
dimensões ou formas especificadas, usando serras e tesouras elétricas e
ferramentas manuais. Executa o tratamento térmico de materiais. Conforma
a quente utilizando equipamentos oxiacetilênicos. Solda peças e
conjuntos. Instala circuitos hidráulicos e pneumáticos. Executa
eletroerosão a fio e por penetração. Confecciona eletrodos para
eletroerosão. Executa a disposição, o ajuste e o assentamento de
mandris, calibradores, suportes e outros dispositivos mecânicos de
usinagem, utilizando instrumentos de medição, ferramentas manuais,
máquinas-ferramentas, formas e outros meios. Mede, marca e escreve sobre
material metálico para fazer a disposição. Levanta, posiciona e prende
os dispositivos em máquinas - tornos, fresadoras, retificadoras,
furadeiras ou prensas -, placas de superfícies ou mesas de trabalho,
utilizando guinchos, blocos em V ou placas auxiliares. Fura peças e
superfícies para fixação. Determina folgas, ajusta gavetas e postiços.
Executa polimento de superfícies planas e com contornos de peças ou
ferramentas, usando raspadores, pedras abrasivas, limas, lixas e esmeris
movidos a energia elétrica. Faz ajustes finais e testes de usinagem para
certificar-se da precisão das operações realizadas e fazer modificações,
se necessário. Executa manutenção preventiva e corretiva, identificando
falhas no ferramental, recuperando peças e utilizando o plano de
produção para modificar ferramentas. Administra tempos de execução das
etapas da manutenção. Realiza testes para garantir que as peças atendam
às especificações, fazendo os ajustes necessários. Faz controle
dimensional, utilizando equipamentos de medição como pinças, blocos de
medida, paquímetros e micrômetros. Traça material a ser usinado. Mede
dureza de componentes submetidos a tratamento térmico. Controla peças
usinadas e produtos acabados, verificando dimensões, alinhamentos e
folgas quanto à conformidade com as especificações e fazendo ajustes
quando necessário. Conserva o local de trabalho limpo e organizado. Atua
na prevenção de acidentes.

\begin{Form}
\textbf{Esta correspondência está adequada?} \\ 
\ChoiceMenu[radio,radiosymbol=\ding{52},bordercolor=0 0.5 0,name=cbovalid_ 721110 ]{}{CBO deve ficar=sim, \hspace{6em} CBO não fica aqui=nao}
\end{Form}

\newpage
\section{ 02024 -- Técnico em Refrigeração e Climatização }

\textbf{Perfil do Curso (CNCT)}

O Técnico em Refrigeração e Climatização será habilitado para:

Planejar, controlar e executar a instalação e a manutenção em
equipamentos de refrigeração e climatização residencial, comercial e
industrial, seguindo legislação vigente, normas técnicas, ambientais, de
saúde e segurança no trabalho e utilizando as boas práticas. Avaliar e
dimensionar máquinas e equipamentos para utilização em projetos de
instalação de refrigeração e climatização. Reconhecer tecnologias
inovadoras presentes no segmento visando à eficiência energética e ao
bem-estar do usuário.

Para atuação como Técnico em Refrigeração e Climatização, são
fundamentais:

Conhecimentos e saberes relacionados aos processos de instalação e
manutenção de aparelhos de refrigeração e climatização de modo a
assegurar a saúde e a segurança dos usuários. Conhecimentos e saberes
relacionados à sustentabilidade do processo, às técnicas e aos processos
de manuseio dos gases, às normas técnicas, à liderança de equipes, à
solução de problemas técnicos e trabalhistas e à gestão de conflitos.

\textbf{CBOs correspondidos e seus perfis (presentes no CNCT, e presentes na predição da IA)}

\textbf{CBO: 725705  ( Mecânico de refrigeração )}

Avalia o local da instalação, verificando dimensões e identificando
fontes de calor. Especifica a capacidade do equipamento de refrigeração
e climatização, em função da carga térmica calculada. Elabora documento
com dados do local. Especifica materiais e acessórios de refrigeração e
climatização. Analisa o projeto de instalação e pesquisa outras fontes
de informação -- tais como catálogos -, para definir tipo, modelo,
tensão, fonte de alimentação do equipamento e outras características do
equipamento. Requisita materiais e acessórios e os confere, quando do
recebimento. Instala equipamentos de refrigeração e climatização,
selecionando ferramentas e equipamentos e medindo o local de instalação
para posicionamento do equipamento. Interpreta normas e procedimentos de
instalação. Instala tubulações e drenos, interliga unidades evaporadoras
e condensadoras e efetua as instalações elétricas. Instala os
equipamentos de automação e controle do sistema. Instala ramais de
dutos, marcando posições ou locais de instalação. Confecciona as peças
componentes do ramal de dutos -- segmentos de dutos --, utilizando
ferramentas, instrumentos e equipamentos específicos. Monta e fixa as
peças, conforme as técnicas requeridas, para formar os ramais de dutos.
Acopla as juntas elásticas para conter vibração, fixa as grelhas de
insuflamento e retorno. Efetua o isolamento térmico dos ramais de dutos.
Calafeta as juntas de conexão e acopla os registros de regulagem de ar à
instalação. Monta tubulações de refrigeração e climatização,
estabelecendo o percurso da instalação e dimensionando os comprimentos
das tubulações, conforme o percurso. Nivela tubulações e equipamentos.
Fixa e solda tubos e conexões. Prepara a instalação para pressurização,
efetuando sua limpeza e tamponamento, para pressurizar a tubulação com
nitrogênio. Monitora a pressão manométrica da tubulação pressurizada,
identificando vazamentos para corrigi-los, quando encontrados. Efetua o
isolamento térmico da tubulação. Aplica vácuo em sistemas de
refrigeração e climatização, despressurizando o sistema e conectando-o a
uma bomba de vácuo. Instala vacuômetro e monitora a pressão do vácuo.
Desmonta os equipamentos de vácuo, após a operação. Carrega os sistemas
de refrigeração e climatização com fluido refrigerante, conectando
manômetros -- de alta e baixa pressão -- e cilindro de fluido
refrigerante. Expurga o ar da mangueira dos manômetros. Abre as válvulas
de serviço do equipamento e injeta o fluído refrigerante, controlando
sua pressão. Realiza testes nos sistemas de refrigeração e climatização,
verificando as condições de alimentação elétrica dos equipamentos --
motores elétricos e compressores -- e acionando-os. Verifica o sentido
de rotação dos motores. Controla tensão e corrente elétricas, pressão
manométrica e outras variáveis, utilizando instrumentos e aparelhos de
medição, para permitir uma avaliação do consumo energético e da eficácia
do sistema. Monitora o superaquecimento e o sub-resfriamento, bem como o
funcionamento dos dispositivos de proteção, automação e controle.
Desconecta garrafas de gás e manômetro. Preenche relatório de testes,
apresenta o equipamento instalado e orienta o usuário quanto ao
funcionamento do sistema. Realiza a manutenção preventiva e corretiva
dos equipamentos e sistemas de refrigeração e climatização, procedendo à
inspeção de manutenção, executando testes e outros procedimentos de
avaliação da integridade das instalações elétricas, das tubulações, dos
ramais de dutos e dos dispositivos de automação e controle, para avaliar
a eficácia do sistema e identificar defeitos e suas causas. Executa as
intervenções saneadoras, recuperando ou substituindo peças ou módulos
danificados, realizando ajustes e outros procedimentos, para
restabelecer a plenitude de funcionamento de equipamentos e sistemas.
Preenche relatórios de manutenção, reportando as circunstâncias dos
serviços prestados. Pode montar equipamentos com tecnologia que dispensa
gás refrigerante.

\begin{Form}
\textbf{Esta correspondência está adequada?} \\ 
\ChoiceMenu[radio,radiosymbol=\ding{52},bordercolor=0 0.5 0,name=cbovalid_ 725705 ]{}{CBO deve ficar=sim, \hspace{6em} CBO não fica aqui=nao}
\end{Form}

\newpage
\section{ 02025 -- Técnico em Sistemas a Gás }

\textbf{Perfil do Curso (CNCT)}

O Técnico em Sistemas a Gás será habilitado para:

Realizar instalação, operação, montagem e manutenção de equipamentos de
sistemas de gases combustíveis e utilidades industriais. Coordenar
processos de utilização de equipamentos, soldagem de tubulação de
polietileno, acessórios de sistemas de combustão a gás. Coordenar
manutenção, produção, transporte, distribuição e entrega de gás natural
e gás liquefeito de petróleo. Identificar problemas e buscar soluções de
geração, transmissão e distribuição de gás. Projetar instalações
prediais de gás e de conversão entre equipamentos, além de realizar a
manutenção destes sistemas atendendo às normas e aos padrões técnicos de
qualidade, saúde e segurança e de meio ambiente.

Para atuação como Técnico em Sistemas a Gás, são fundamentais:

Conhecimentos e saberes relacionados aos processos de planejamento e
operação de plantas de gás de modo a assegurar a saúde e a segurança dos
trabalhadores e dos usuários. Conhecimentos e saberes relacionados à
sustentabilidade do processo produtivo, às técnicas e aos processos de
produção e distribuição de gases, às normas técnicas, à liderança de
equipes, à solução de problemas técnicos e trabalhistas e à gestão de
conflitos.

\textbf{CBOs correspondidos e seus perfis (presentes no CNCT, e presentes na predição da IA)}

\textbf{CBO: 724115  ( Instalador de tubulações )}

Prepara e organiza o trabalho, analisando projeto e interpretando
plantas e desenhos de tubulação. Estabelece cronograma para execução do
serviço. Especifica e quantifica os materiais e seleciona máquinas,
equipamentos e ferramentas, para a realização das atividades. Pode
elaborar orçamento do serviço. Prepara o local para a instalação,
inspecionando e isolando a área. Assinala os pontos de colocação das
tubulações, junções e furos em paredes, muros e escavações do solo,
utilizando ferramentas de traçagem ou marcação. Abre paredes, lajes,
pisos ou valas para introduzir tubos metálicos e partes anexas. Fixa
suportes. Prepara as tubulações, cortando tubos, abrindo roscas,
alinhando tubos conforme ângulo especificado, encaixando conexões,
encurvando tubos, colando e pintando tubulações. Distribui, assenta e
veda tubulações, interligando redes a ramais (pontos de consumo).
Posiciona e fixa todos os elementos, baseando-se nas especificações do
projeto, para montar a linha de condução dos fluidos. Utiliza soldagem,
no processo de montagem da tubulação. Combina uso de sistemas de
tubagens por meio de pressão e por tubagens de escoamento, com ou sem
bombas de pressão. Identifica a pressão do fluido. Instala e regula
peças de utilização, encaixando as tubulações de fluidos, líquidos,
gases e outras substâncias. Instala sistemas de drenagem, distribuição e
irrigação. Instala tanques, torneiras e válvulas. Liga a tubulação aos
aparelhos e demais componentes da instalação - como registros, bombas e
aquecedor - utilizando ferramentas manuais e material de vedação, para
completar as instalações projetadas. Realiza as instalações utilizando
canos, tubos rígidos, flexíveis e conexões em materiais de polietileno
de alta densidade (PEAD), cloreto de polivinil clorado (CPVC) e
polipropileno (PPR). Pode confeccionar peças de tubulação industrial.
Testa a estanqueidade, examinando a existência de vazamentos. Monitora
teste no manômetro e na rede. Efetua as correções necessárias e refaz
testes. Procede a ensaios de funcionamento dos equipamentos e
dispositivos acessórios das redes e realiza ajustes, se necessários.
Fecha as escavações abertas no solo e os furos e os rasgos nas paredes.
Realiza manutenção de tubulações industriais, emendando tubos,
desobstruindo ralos e substituindo tubos, registros e torneiras. Pode
instalar sistemas de combate a incêndio, com extintores, mangueiras e
sprinklers. Preenche relatório dos serviços executados. Pode emitir
recibos e notas fiscais de serviços. Realiza manutenção de primeiro
nível de máquinas e equipamentos, limpando-os e os lubrificando.
Identifica falhas e defeitos, requisitando manutenção corretiva. Mantém
ferramentas e instrumentos de trabalho limpos, organizados,
acondicionados e em plenas condições de uso e funcionamento. Controla
estoque de materiais, podendo requisitar reposição. Controla
desperdícios de materiais, selecionando os reutilizáveis. Faz o descarte
de resíduos, de acordo com as normas ambientais. Zela pela segurança,
identificando e sinalizando áreas de risco e prevenindo acidente e
incêndio. Utiliza Equipamentos de Proteção Individual (EPI) e
Equipamentos de Proteção Coletiva (EPC). Pode prestar primeiros
socorros.

\begin{Form}
\textbf{Esta correspondência está adequada?} \\ 
\ChoiceMenu[radio,radiosymbol=\ding{52},bordercolor=0 0.5 0,name=cbovalid_ 724115 ]{}{CBO deve ficar=sim, \hspace{6em} CBO não fica aqui=nao}
\end{Form}

\textbf{CBO: 724130  ( Instalador de tubulações de gás combustível (produção e distribuição) )}

Prepara e organiza o trabalho, interpretando projeto de implantação de
sistema de geração e distribuição de gás combustível. Estabelece
cronograma para execução do serviço. Seleciona máquinas, equipamentos,
ferramentas, instrumentos e materiais, para a realização das atividades.
Pode elaborar orçamento do serviço. Pode usar BIM - Modelagem da
Informação da Construção (Building Information Modeling), para
conhecimento do projeto e para posterior registro das atividades
executadas. Examina e marca o local definido para instalação, conforme
projeto. Abre paredes, lajes, pisos ou valas e fixa suportes, para
encaixar as tubulações de gás combustível. Separa materiais conforme
medidas e tipos, para o processo de pré-montagem de tubulações. Corta e
encurva tubos. Abre roscas, encaixa conexões e ponteia tubulações. Pinta
as tubulações, identificando-as com cores, conforme finalidade. Usa
processos de corte e de união por soldagem, brasagem forte e brasagem
fraca. Distribui, assenta e veda tubulações, interligando redes a ramais
em pontos de centrais de suprimento de gás combustível. Combina uso de
sistemas de tubagens com ou sem bombas de pressão, identificando e
regulando pressão do transporte de gás combustível. Instala acessórios,
peças e equipamentos, tais como tanques, válvulas e atuadores. Realiza
testes operacionais de pressão de gás combustível, regulando pressão e
drenando tubulações. Realiza teste de estanqueidade, para assegurar que
não existe vazamento nas redes canalizadas. Monitora teste no manômetro
e na rede. Quando necessário, corrige falhas na vedação e refaz o teste.
Libera rede para uso. Instala sistemas de combate a incêndio, com
extintores, mangueiras e sprinklers. Executa manutenção em centrais de
suprimento e redes de transporte e distribuição de gás combustível,
identificando falhas e defeitos. Desativa o sistema de distribuição,
levanta causa do problema identificado, substitui componentes e
equipamentos defeituosos, estanca vazamentos, testa os reparos
realizados e reativa o sistema de distribuição. Preenche relatórios de
serviços. Pode emitir recibos e notas fiscais de serviços. Realiza
manutenção de primeiro nível de máquinas e equipamentos usados durante
os serviços, limpando-os e os lubrificando. Pode requisitar manutenção
corretiva. Mantém ferramentas e instrumentos de trabalho limpos,
organizados, acondicionados e em plenas condições de uso e
funcionamento. Controla estoque de materiais, podendo requisitar
reposição. Controla desperdícios de materiais, selecionando os
reutilizáveis. Faz o descarte de resíduos, de acordo com as normas
ambientais. Zela pela segurança, identificando e sinalizando áreas de
risco e prevenindo acidente e incêndio. Utiliza Equipamentos de Proteção
Individual (EPI) e Equipamentos de Proteção Coletiva (EPC). Pode prestar
primeiros socorros.

\begin{Form}
\textbf{Esta correspondência está adequada?} \\ 
\ChoiceMenu[radio,radiosymbol=\ding{52},bordercolor=0 0.5 0,name=cbovalid_ 724130 ]{}{CBO deve ficar=sim, \hspace{6em} CBO não fica aqui=nao}
\end{Form}

\newpage
\section{ 02026 -- Técnico em Sistemas de Energia Renovável }

\textbf{Perfil do Curso (CNCT)}

O Técnico em Sistemas de Energia Renovável será habilitado para:

Planejar, controlar e executar projetos de instalação, operação,
montagem e manutenção de sistemas de geração, transmissão e distribuição
de energia elétrica de fontes renováveis. Coordenar atividades de
utilização e conservação de energia e fontes alternativas (energia
eólica, solar e hidráulica). Seguir especificações técnicas e de
segurança na montagem de projetos de viabilidade de geração de energia
elétrica proveniente de fonte eólica, solar e hidráulica em substituição
às convencionais. Desenvolver novas formas produtivas para a geração de
energias renováveis e eficiência energética, bem como adotar medidas
para o uso eficiente de energia elétrica. Identificar e propor soluções
para problemas de gestão energética, para questões decorrentes da
geração, transmissão e distribuição de energia elétrica.

Para atuação como Técnico em Sistemas de Energia Renovável, são
fundamentais:

Conhecimentos e saberes relacionados aos processos de planejamento e
instalação de sistemas de energia renovável de modo a assegurar a saúde
e a segurança dos trabalhadores e dos usuários. Conhecimentos e saberes
relacionados à sustentabilidade do processo produtivo, às técnicas e aos
processos de produção limpa, às normas técnicas, à liderança de equipes,
à solução de problemas técnicos e trabalhistas e à gestão de conflitos.

\textbf{CBOs correspondidos e seus perfis (presentes no CNCT, e presentes na predição da IA)}

\textbf{CBO: 313110  ( Eletrotécnico (produção de energia) )}

Elabora estudos e projetos em sistemas de produção de energia,
determinando escopo, aplicando normas e avaliando viabilidade econômica.
Dimensiona circuitos eletroeletrônicos e componentes. Participa do
desenvolvimento de produtos e elabora documentações técnicas, de acordo
com os requisitos do projeto. Realiza projetos em sistemas de produção
de energia - conforme especificações e cronograma -, executando montagem
e comissionamento e colocando em operação. Participa do desenvolvimento
de processos em sistemas de produção de energia, estabelecendo
procedimentos, normas e padrões. Determina fluxograma e define máquinas
e equipamentos. Realiza medições e ensaios e propõe melhorias. Aplica
tecnologias adequadas ao projeto e ao processo, considerando a crescente
utilização de fontes alternativas de energia, os programas de eficiência
energética e as tecnologias de automação aplicadas a sistemas elétricos.
Opera sistemas elétricos de geração, transmissão e distribuição de
energia - de acordo com normas, instruções e procedimentos -,
supervisionando e elaborando os programas de manobras. Pode utilizar
sistemas de automação. Executa manutenção preventiva, corretiva e
preditiva em sistemas elétricos, de acordo com o plano de manutenção,
diagnosticando o desempenho dos equipamentos. Atua na supervisão e no
treinamento de equipe de trabalho, avaliando desempenho e propondo
programas de aperfeiçoamento e atualização profissional.

\begin{Form}
\textbf{Esta correspondência está adequada?} \\ 
\ChoiceMenu[radio,radiosymbol=\ding{52},bordercolor=0 0.5 0,name=cbovalid_ 313110 ]{}{CBO deve ficar=sim, \hspace{6em} CBO não fica aqui=nao}
\end{Form}

\textbf{CBO: 861105  ( Operador de central hidrelétrica )}

Planeja o trabalho, considerando ordens de serviço; interpretando
procedimentos, normas e manuais técnicos; e analisando gráficos,
esquemas elétricos e fluxogramas. Manobra equipamentos de geração de
energia elétrica, em centrais hidrelétricas de diversificados portes --
Usinas Hidrelétricas de Energia (UHE), com reservatório ou a fio de
água; Pequenas Centrais Hidrelétricas (PCH) ou Centrais Geradoras
Hidrelétricas (CGH) --, acionando seus mecanismos e sistemas, para
iniciar, manter ou interromper seu funcionamento. Analisa ordens de
manobra; interpreta leiaute da planta de geração; libera equipamentos de
geração para funcionamento; e aciona turbinas, geradores e equipamentos
auxiliares. Aciona comportas para vazão do excesso de água nas usinas
hidrelétricas. Sincroniza unidade geradora. Analisa plano de isolamento;
bloqueia funcionamento de equipamentos; e desacopla unidade geradora do
sistema. Executa manobras para interligação do sistema de geração ao
sistema de transmissão de energia elétrica, operando equipamentos de
telecomando e/ou dispositivos de automação, de acordo com procedimentos
específicos, a fim de alimentar -- com tensões e correntes elétricas das
saídas dos geradores -- o sistema de transmissão. Determina o número de
geradores em operação e balanceia cargas entre geradores. Isola e libera
linhas e equipamentos de transmissão; bloqueia ou desbloqueia
equipamentos de transmissão, de acordo com as circunstâncias. Controla a
produção de energia elétrica em centrais hidrelétricas, regulando o
funcionamento dos geradores; utilizando instrumentos, equipamentos e
sistemas de supervisão e controle, bem como elementos de sensoriamento e
de redes, para atender à demanda de corrente elétrica. Controla volume
de água do reservatório e vazão de água nas turbinas. Controla variação
de velocidade das turbinas. Ajusta parâmetros dos equipamentos de
produção de energia elétrica. Ajusta níveis de tensão dos geradores.
Monitora sistemas de geração de energia elétrica de centrais
hidrelétricas, utilizando aparelhos e métodos específicos, para avaliar
o funcionamento do conjunto de máquinas e equipamentos de transformação
de energia hidráulica em elétrica. Inspeciona funcionamento de turbinas,
de geradores e de equipamentos auxiliares. Monitora pressão,
temperatura, nível e vazão; e monitora tensão, corrente, frequência e
potência elétricas no sistema de geração. Mantém registros dos dados
monitorados. Executa atividades de apoio à manutenção, de acordo com
procedimentos de rotina, para manter máquinas e equipamentos de geração
de energia elétrica em condições de operação nas centrais hidrelétricas.
Identifica defeitos, registra dados sobre condições de máquinas e
equipamentos, solicitando manutenção, quando necessário. Acompanha
serviços de manutenção de máquinas e equipamentos. Pode realizar
pequenos reparos. Testa máquinas e equipamentos. Registra número de
horas de funcionamento de máquinas e equipamentos em operação. Limpa
máquinas e equipamentos. Interage com outros setores das centrais
hidrelétricas e com outras instituições relacionadas à geração de
energia elétrica, operando instrumentos de comunicação e participando de
reuniões e discussões técnicas. Elabora documentação técnica,
preenchendo formulários -- físicos ou eletrônicos -- e redigindo
relatórios técnicos. Mantém o ambiente de trabalho limpo e organizado.
Atua segundo procedimentos de proteção ao meio ambiente, participando de
ações para minimizar impactos ambientais. Trabalha segundo procedimentos
de segurança e de saúde ocupacional, utilizando Equipamentos de Proteção
Individual (EPI); e isolando áreas de riscos físicos, químicos,
biológicos e ergonômicos. Identifica atos e condições inseguras, testa
equipamentos de segurança e simula emergências. Participa de programas
de segurança no trabalho e de prevenção de doenças que possam se
originar no ambiente de trabalho.

\begin{Form}
\textbf{Esta correspondência está adequada?} \\ 
\ChoiceMenu[radio,radiosymbol=\ding{52},bordercolor=0 0.5 0,name=cbovalid_ 861105 ]{}{CBO deve ficar=sim, \hspace{6em} CBO não fica aqui=nao}
\end{Form}

\newpage
\section{ 02027 -- Técnico em Soldagem }

\textbf{Perfil do Curso (CNCT)}

O Técnico em Soldagem será habilitado para:

Planejar e coordenar a execução de atividades de soldagem em estruturas
metálicas e tubulações industriais, respeitando procedimentos e normas
técnicas, de qualidade, de saúde, de segurança e de meio ambiente.
Elaborar ensaios destrutivos e não destrutivos em produtos soldados.
Selecionar processos de soldagem e corte calibrando máquinas e indicando
consumíveis. Utilizar processos de soldagem e corte manual e
automatizado por meio de eletrodo revestido, TIG, MIG, MAG, oxigás, arco
submerso, brasagem e plasma. Reconhecer tecnologias inovadoras presentes
no segmento visando a atender às transformações digitais na sociedade.

Para atuação como Técnico em Soldagem, são fundamentais:

Conhecimentos e saberes relacionados aos processos de soldagem de modo a
assegurar a saúde e a segurança dos trabalhadores, bem como assegurar a
qualidade dos projetos soldados. Conhecimentos e saberes relacionados à
sustentabilidade do processo produtivo, às técnicas e aos processos de
soldagem, às normas técnicas, à liderança de equipes, à solução de
problemas técnicos e trabalhistas e à gestão de conflitos.

\textbf{CBOs correspondidos e seus perfis (presentes no CNCT, e presentes na predição da IA)}

\textbf{CBO: 314620  ( Técnico em soldagem )}

Planeja processos de soldagem -- em estruturas, chapas, perfis,
tubulações e outros artefatos industriais metalmecânicos --,
interpretando projetos e consultando normas técnicas. Define a aplicação
de processos diversificados de soldagem -- tais como eletrodo revestido;
soldagem por gás inerte de tungstênio-TIG (Tungsten Inert Gas); soldagem
com utilização de gás inerte para proteção da poça de fusão-MIG (Metal
Inert Gas); soldagem com utilização de gás ativo para proteção da poça
de fusão-MAG (Metal Active Gas); oxigás; arame tubular; arco submerso; e
plasma --, conforme requisitos de projetos e de posições dos produtos
soldados. Elabora o plano de qualificação dos procedimentos de soldagem.
Analisa desenhos de engenharia, projetos, especificações, esboços,
ordens de serviço e planilhas de dados de segurança de material para
planejar o leiaute e a montagem dos cenários das operações de soldagem.
Elabora o plano de soldagem. Define os procedimentos de execução e
inspeção. Estima o material de consumo e de aplicação. Seleciona os
procedimentos para cada atividade. Verifica as condições de segurança do
ambiente. Organiza equipamentos, ferramentas e materiais no local de
trabalho. Define a sequência de execução dos serviços de soldagem.
Elabora a programação da produção. Opera processos de soldagem,
determinando equipamentos e métodos de soldagem necessários e aplicando
conhecimentos de metalurgia, geometria, técnicas e posições de soldagem.
Monitora processos de soldagem, utilizando procedimentos e sistemas
específicos, tendo em vista a garantida da qualidade exigida para os
produtos soldados. Utiliza sistemas de gerenciamento de informações de
soldagem, para a melhoria da produtividade e da qualidade dos processos.
Desenvolve modelos para projetos de soldagem, usando cálculos
matemáticos baseados em informações dos desenhos de projeto. Ao receber
notificações de mau funcionamento de equipamentos ou de materiais
defeituosos, define as medidas necessárias para solução dos problemas
operacionais. Qualifica procedimentos de soldagem, identificando chapas
ou tubos de testes e acompanhando o processo de qualificação. Coordena a
realização de ensaios destrutivos e não destrutivos em produtos
soldados. Inspeciona matéria-prima e consumíveis de soldagem. Verifica
as condições de secagem, estocagem e manuseio de consumíveis de
soldagem. Acompanha e inspeciona a soldagem no decorrer dos processos.
Supervisiona processos de soldagem, tendo em vista a eficácia e a
eficiência dos serviços. Aplica normas técnicas na execução dos serviços
de soldagem. Coordena as equipes na execução dos serviços de soldagem,
de acordo com o cronograma e a programação da produção. Verifica a
montagem, para o cumprimento dos requisitos do projeto no que se refere
à soldagem. Controla o consumo do material. Controla condições de
secagem, estocagem e manuseio dos consumíveis de soldagem. Realiza o
controle dimensional. Verifica a conclusão da soldagem; as
identificações da soldagem; e os aspectos visual e dimensional da
soldagem. Coordena as operações de tratamento térmico. Libera o serviço
executado para inspeção do controle de qualidade. Pode supervisionar
equipes de trabalho, distribuindo tarefas, avaliando desempenho e
desenvolvendo treinamentos. Incorpora tecnologias inovadoras aos
processos de soldagem, tais como uso de soldagem a laser; e aplicações
para máquinas inteligentes, com princípios de automação e robótica.
Define técnicas, métodos, máquinas e equipamentos para operar com novos
materiais, como aços de alta resistência combinados com revestimentos
galvanizados; alumínio; aços inoxidáveis; e outros materiais não
ferrosos. Segue normas regulamentadoras e procedimentos de segurança e
saúde no trabalho em processos de soldagem, orientando e fiscalizando a
utilização de equipamentos de proteção individual e coletiva. Promove o
descarte e a destinação de resíduos de acordo com procedimentos e normas
de preservação ambiental.

\begin{Form}
\textbf{Esta correspondência está adequada?} \\ 
\ChoiceMenu[radio,radiosymbol=\ding{52},bordercolor=0 0.5 0,name=cbovalid_ 314620 ]{}{CBO deve ficar=sim, \hspace{6em} CBO não fica aqui=nao}
\end{Form}

\textbf{CBO: 724305  ( Brasador )}

Solda peças metálicas por brasagem de acordo com os métodos de
aquecimento a serem aplicados -- brasagem por chama, em forno, por
indução, por resistência, por imersão e por infravermelho --. Define o
metal de adição em função do metal de base, do método de aquecimento, do
desenho da junta e da proteção. Controla o enchimento das folgas entre
os materiais de base, aplicando fluxos, verificando o efeito capilar ou
capilaridade, para assegurar a qualidade do processo de brasagem.
Verifica as faixas de medidas das folgas, de acordo com tabelas,
considerando o metal de adição, o tipo de fluxo e o tipo de junta
utilizada. Controla a molhabilidade do metal de base tendo em vista o
preenchimento das folgas e a qualidade da brasagem. Controla a
temperatura de fusão (do metal de adição) e a temperatura de brasagem
(do metal de base). Executa a limpeza da superfície dos materiais de
base, aplicando solventes industriais para eliminação de óleos ou
graxas. Pode eliminar possível corrosão causada pela aplicação de
fluxos. Prepara peças metálicas para o processo de brasagem, aplicando
operações de chanfrar, goivar e escovar peças, identificando posições de
soldagem e inspecionando visualmente as condições para a soldagem.
Prepara equipamentos e consumíveis para soldagem de peças metálicas por
brasagem, definindo o tipo de gás, o equipamento de soldagem, o bico e a
regulagem do maçarico. Regula parâmetros de corte e de soldagem e dos
manômetros. Substitui acessórios de acordo com o processo. Controla o
aquecimento local da peça e o derretimento do material consumível no
processo, a velocidade e o resfriamento da soldagem e realiza o
acabamento da solda, aplicando removedores para retirada de sujidade.
Organiza o local de trabalho para soldagem de peças metálicas por
brasagem, em observância às Instruções de Execução e Inspeção de
Soldagem (IEIS), com consulta a desenhos e especificações técnicas.
Executa o processo de brasagem de acordo normas e procedimentos técnicos
de manuseio de equipamentos e consumíveis, providenciando ferramentas e
isolando a área com anteparos. Segue rigorosamente normas
regulamentadoras e procedimentos de segurança e saúde no trabalho em
processos de soldagem. Aplica procedimentos de qualidade e produtividade
na realização das atividades, considerando redução de tempo, maior
produtividade, cumprimento de prazos e eliminação de desperdícios.
Organiza o local de trabalho, selecionando materiais de trabalho,
delimitando e limpando a área de atuação, acondicionando equipamentos e
consumíveis, descartando resíduos de acordo com normas e procedimentos.

\begin{Form}
\textbf{Esta correspondência está adequada?} \\ 
\ChoiceMenu[radio,radiosymbol=\ding{52},bordercolor=0 0.5 0,name=cbovalid_ 724305 ]{}{CBO deve ficar=sim, \hspace{6em} CBO não fica aqui=nao}
\end{Form}

\textbf{CBO: 724310  ( Oxicortador a mão e a  máquina )}

Corta peças metálicas, com maçarico de corte utilizando chama de gás
combustível e oxigênio, para reutilização de material ou desmonte de
estruturas metálicas. Maneja maçaricos e máquinas, ajustando e acionando
os seus comandos, para cortar lâminas e outras peças metálicas nas
dimensões desejadas. Prepara peças para corte por chama oxigás,
identificando a posição de corte e considerando as operações de
chanfragem e goivagem. Define o tipo de corte -- leve, médio ou pesado
-- em função da peça a ser cortada e de acordo com a superfície que se
deseja obter -- corte reto, corte duplo e corte com chanfro --. Define o
tipo de maçarico -- injetores e misturadores --, o tipo de gás, o bico
de maçarico, os bicos de gás combustível e oxigênio e os tipos de
chamas, de acordo com as especificações para a espessura a ser cortada.
Ajusta as válvulas dos cilindros e a pressão de trabalho e substitui
acessórios de acordo o processo e a técnica aplicados. Regula parâmetros
de corte, de manômetros, do maçarico e da máquina, controlando
velocidade e direção do corte. Pode utilizar guias para executar cortes
retos. Controla o tempo de fusão do material, sem ou com adição de
eletrodo, para possibilitar o corte das peças. Controla o aquecimento
local da peça, a velocidade e o resfriamento do corte e realiza o
acabamento. Realiza inspeção visual na peça cortada, identificando
defeitos por escória ou por trincas e eliminando-os por meio de técnicas
e métodos de reparo. Executa limpeza de peças, escovando-as e aplicando
removedores para retirada de óleos e graxas. Organiza o local de
trabalho, delimitando e limpando a área de atuação, seguindo
procedimentos de manuseio e acondicionamento de máquinas, equipamentos,
acessórios e consumíveis. Segue Instruções de Execução e Inspeção de
Soldagem (IEIS), com consulta a desenhos e especificações técnicas.
Segue rigorosamente normas regulamentadoras e procedimentos de segurança
e saúde no trabalho em processos de soldagem, cumprindo a utilização de
equipamentos de proteção individual e coletiva. Mantém tais equipamentos
organizados, acondicionados e em plenas condições de uso e
funcionamento. Aplica procedimentos de qualidade e produtividade na
realização das atividades. Cumpre normas de preservação ambiental,
descartando e destinando resíduos de acordo com normas e procedimentos.
Controla o funcionamento de máquinas e equipamentos e as condições de
plena utilização de consumíveis, acessórios e materiais, requisitando
manutenção e substituições.

\begin{Form}
\textbf{Esta correspondência está adequada?} \\ 
\ChoiceMenu[radio,radiosymbol=\ding{52},bordercolor=0 0.5 0,name=cbovalid_ 724310 ]{}{CBO deve ficar=sim, \hspace{6em} CBO não fica aqui=nao}
\end{Form}

\textbf{CBO: 724315  ( Soldador )}

Une peças de ligas metálicas, por pressão ou fusão, considerando
conjuntos, estruturas, tubos, chapas e perfis. Define e prepara o
processo de soldagem de acordo com a peça a ser executada, seguindo
desenhos e especificações técnicas. Prepara peças metálicas para
soldagem, realizando inspeção visual, preaquecendo a peça com maçaricos,
chanfrando, escovando e goivando as peças. Aplica removedores para
retirada de sujidade da superfície. Define o material de adição de
acordo com o metal de base e com o processo a ser aplicado, o tipo de
eletrodo, o bico do maçarico e tipo de gás, a regulagem de manômetros.
Controla o aquecimento local da peça, o derretimento do consumível no
processo de soldagem e o resfriamento da peça. Controla as posições de
soldagem e o ângulo do maçarico. Controla e ajusta parâmetros de
soldagem de acordo com as Instruções de Execução e Inspeção de Soldagem
(IEIS). Verifica o acabamento de superfície da solda na peça, a
qualidade do cordão de solda, no que se refere a tensões, trincas e
deformações, eliminando escórias e fumus. Pode executar operações
básicas de medição, limagem, serragem, furação, corte de chapa, entre
outras. Pode montar e pontear elementos de peças metálicas. Pode
executar soldagem de manutenção e recuperação de peças metálicas - no
que se refere a trincas ou fraturas, por exemplo -, considerando os
custos da matéria-prima, energia, insumos e mão de obra envolvida.
Organiza o local de trabalho, isolando e limpando a área de atuação e
seguindo procedimentos de manuseio e de armazenamento de máquinas,
equipamentos, acessórios e consumíveis. Segue Instruções de Execução e
Inspeção de Soldagem (IEIS), com consulta a desenhos, documentação e
especificações técnicas. Segue rigorosamente as normas regulamentadoras
e os procedimentos de segurança e saúde no trabalho em processos de
soldagem, cumprindo a utilização de equipamentos de proteção individual
e coletiva. Mantém tais equipamentos organizados, acondicionados e em
plenas condições de uso e funcionamento. Aplica procedimentos de
qualidade e de produtividade na realização das atividades. Cumpre normas
de preservação ambiental, descartando e destinando resíduos de acordo
com normas e procedimentos. Controla o funcionamento de máquinas e
equipamentos e as condições de plena utilização de consumíveis,
acessórios e materiais, requisitando manutenção e substituições.

\begin{Form}
\textbf{Esta correspondência está adequada?} \\ 
\ChoiceMenu[radio,radiosymbol=\ding{52},bordercolor=0 0.5 0,name=cbovalid_ 724315 ]{}{CBO deve ficar=sim, \hspace{6em} CBO não fica aqui=nao}
\end{Form}

\textbf{CBO: 724320  ( Soldador a  oxigás )}

Une peças metálicas, de metais ferrosos e não-ferrosos, para fabricar
peças ou reparar solda, fazendo uso do processo de soldagem a oxigás.
Define máquinas, equipamentos, materiais e consumíveis e orienta-se por
desenhos, especificações e documentação. Prepara peças para soldagem,
executando operações de chanfragem e goivagem, realizando inspeção
visual e aplicando removedores para retirada de sujidade. Pode executar
operações de dobramento e desempeno. Define o material de adição de
acordo com o metal de base, os tipos de gases, o tipo e o bico do
maçarico e o tipo da chama -- redutora ou neutra --. Define o
posicionamento de realização da soldagem tendo em vista o controle do
ponto de fusão, de acordo com a espessura do material a ser soldado.
Ajusta e controla a pressão do equipamento, a vazão e a velocidade de
propagação dos gases, o volume de material da poça de fusão, a
temperatura da chama, de acordo com especificações técnicas e parâmetros
de soldagem a gás. Controla o aquecimento local da peça e o derretimento
do material consumível no processo, no que se refere à velocidade e ao
resfriamento da solda e do acabamento da superfície soldada, regulando
manômetros e maçaricos. Controla a qualidade do cordão de solda no que
se refere à espessura e características mecânicas, selecionando varetas
(material de adição) com comprimento e espessura compatíveis com a
aplicação. Atua, em todas as etapas do processo de soldagem, desde a
preparação até o teste final, de acordo com a Instrução de Execução e
Inspeção de Soldagem -- IEIS, selecionando processos e todos os testes
aplicáveis para garantir a estanqueidade, a integridade, a resistência e
a estética da solda. Organiza o local de trabalho, isolando e limpando a
área de atuação, seguindo procedimentos de manuseio e armazenamento de
máquinas, equipamentos, acessórios e consumíveis. Segue rigorosamente
normas regulamentadoras e procedimentos de segurança e saúde no trabalho
em processos de soldagem, cumprindo a utilização de equipamentos de
proteção individual e coletiva. Mantém tais equipamentos organizados,
acondicionados e em plenas condições de uso e funcionamento. Aplica
procedimentos de qualidade e produtividade na realização das atividades.
Cumpre normas de preservação ambiental, descartando e destinando
resíduos de acordo com normas e procedimentos. Controla o funcionamento
de máquinas e equipamentos, a integridade de consumíveis, acessórios e
materiais, mantendo-os em plenas condições de uso. Requisita, quando
necessário, manutenção e substituições.

\begin{Form}
\textbf{Esta correspondência está adequada?} \\ 
\ChoiceMenu[radio,radiosymbol=\ding{52},bordercolor=0 0.5 0,name=cbovalid_ 724320 ]{}{CBO deve ficar=sim, \hspace{6em} CBO não fica aqui=nao}
\end{Form}

\textbf{CBO: 724325  ( Soldador elétrico )}

Executa soldagem elétrica para fabricação, reparo e recuperação de peças
metálicas, nas posições horizontal, verticais ascendente e descendente,
plana e sobrecabeça, aplicando técnicas de abertura e manutenção do arco
elétrico. Pode fazer uso de ferramentas automáticas. Define o tipo de
soldagem elétrica a ser aplicado, considerando a peça a ser fundida ou
revestida e de acordo com desenhos, documentação e especificações
técnicas. Define máquinas, equipamentos, materiais, consumíveis e
acessórios a serem utilizados na soldagem, de acordo com o processo a
ser aplicado, seguindo Instrução de Execução e Inspeção de Soldagem
(IEIS). Prepara as peças para soldagem, verificando visualmente as
condições da superfície e aplicando removedores para retirada de óleos e
graxas. Aquece previamente a peça com maçarico, chanfrando-a,
escovando-a e goivando-a. Seleciona o arco elétrico e o material de
adição de acordo com o metal de base. Define tipo, espessura e
comprimento do eletrodo, de acordo com o tipo de processo a ser
aplicado, considerando também a posição em que a solda será realizada e
as propriedades da solda requerida. Define o tipo de revestimento de
soldagem, considerando a proteção do metal, a adição de elementos de
liga, o direcionamento do arco elétrico, a posição de soldagem, as
propriedades mecânicas do metal de solda, entre outros fatores. Controla
os parâmetros variáveis de soldagem - tais como tensão e corrente, ponto
de fusão do eletrodo, tensão e comprimento do arco elétrico, velocidade
e temperatura da soldagem e vazão de gases -, eliminando o sopro
magnético e fazendo o acabamento da solda. Organiza o local de trabalho,
isolando e limpando a área de atuação, seguindo procedimentos de
manuseio e armazenamento de máquinas, equipamentos, acessórios e
consumíveis. Segue rigorosamente normas regulamentadoras e procedimentos
de segurança e saúde no trabalho em processos de soldagem, cumprindo a
utilização de equipamentos de proteção individual e coletiva. Mantém
tais equipamentos organizados, acondicionados e em plenas condições de
uso e funcionamento. Aplica procedimentos de qualidade e produtividade
na realização das atividades. Cumpre normas de preservação ambiental,
descartando e destinando resíduos de acordo com normas e procedimentos.
Controla o funcionamento de máquinas e equipamentos e as condições de
plena utilização de consumíveis, acessórios e materiais. Requisita
manutenção e substituições.

\begin{Form}
\textbf{Esta correspondência está adequada?} \\ 
\ChoiceMenu[radio,radiosymbol=\ding{52},bordercolor=0 0.5 0,name=cbovalid_ 724325 ]{}{CBO deve ficar=sim, \hspace{6em} CBO não fica aqui=nao}
\end{Form}

\newpage
\section{ 03001 -- Técnico em Alimentação Escolar }

\textbf{Perfil do Curso (CNCT)}

O Técnico em Alimentação Escolar será habilitado para:

Organizar e executar fluxos de aquisição e armazenamento de alimentos e
insumos. Organizar, controlar e executar os processos de higienização de
alimentos. Preparar, selecionar e conservar alimentos, conforme cardápio
escolar e orientações nutricionais. Calcular o quantitativo de alimentos
para merenda escolar, considerando as porções diárias e a aquisição
mensal junto à gestão da escola. Preparar variedades de receitas,
considerando as características regionais associadas ao cardápio
escolar. Organizar e controlar a cozinha escolar para o preparo e o
fornecimento da alimentação. Utilizar técnicas de higiene e segurança do
trabalho desde a aquisição dos alimentos ao descarte de resíduos.
Realizar o papel de educador alimentar sob supervisão de nutricionista.

Para atuação como Técnico em Alimentação Escolar, são fundamentais:

Conhecimentos e saberes relacionados à prática da alimentação escolar, a
valores nutricionais dos alimentos, a variações culinárias, a
especificidades regionais alimentícias e a porções alimentares.
Princípios e práticas da organização da cozinha escolar, da conservação,
do armazenamento e de manejo de alimentos, de descarte de resíduos, de
técnicas de segurança e higiene do trabalho. Responsabilidade com a
formação de hábitos saudáveis de alimentação e com o cumprimento das
legislações vigentes. Capacidade de se comunicar assertivamente, de
colaborar e mediar conflitos, de solucionar possíveis problemas durante
o processo de preparo e fornecimento da merenda escolar. Habilidade para
lidar com imprevistos, demonstrando estabilidade emocional e foco para
solução de problemas dentro dos processos que envolvem a alimentação
escolar nas instituições de ensino.

\textbf{CBOs correspondidos e seus perfis (presentes no CNCT, e presentes na predição da IA)}

\textbf{CBO: 841408  ( Cozinhador (conservação de alimentos) )}

Prepara as atividades, de acordo com ficha técnica, ordem de serviço e
programação da produção. Seleciona utensílios, instrumentos,
matérias-primas e insumos, de acordo com as especificações do processo e
dos produtos. Inspeciona as condições operacionais de máquinas e
equipamentos. Interpreta as receitas e os procedimentos de execução,
determinando as proporções requeridas. Efetua as medidas de massa,
volume e temperatura das matérias-primas e demais insumos, tais como
temperos, conservantes e espessantes. Prepara os alimentos, higienizando
as matérias-primas; aplicando técnicas de cortes de vegetais; realizando
a desossa, limpeza e porcionamento de produtos cárneos; entre outras
atividades. Opera equipamentos de cocção - tais como caldeiras, panelas
de pressão, cubas de vácuo ou de evaporação e aparelhos de frigir - e
controla a temperatura, o tempo e a pressão de cozimento. Examina o
produto cozido, verificando sua consistência e conferindo o atendimento
aos parâmetros estabelecidos na ficha técnica de produção e na ordem de
serviço. Opera equipamentos de embalagem e etiquetagem, acondicionando
os produtos para armazenamento e discriminando, na etiqueta, o tipo de
produto, a data de preparo, o prazo de validade, entre outras
informações. Opera câmara fria para armazenar e conservar as
matérias-primas, os insumos e os produtos. Emite relatório de produção,
para prestar informações sobre o processo realizado e orientar a sua
continuidade. Mantém as instalações limpas, organizadas, higienizadas e
desinfetadas. Limpa, organiza e higieniza equipamentos, instrumentos e
utensílios. Destina adequadamente os resíduos gerados, realizando
descarte de acordo com as normas ambientais. Controla a reposição de
matérias-primas e insumos, examinando a programação da produção, as
fichas técnicas e o estoque disponível. Recebe e armazena as
mercadorias, controlando seu estado de conservação. Verifica o
funcionamento das máquinas e dos equipamentos de cozinha. Desliga
máquinas e equipamentos ao ouvir alarmes de mau funcionamento,
requisitando a necessária manutenção ou substituição.

\begin{Form}
\textbf{Esta correspondência está adequada?} \\ 
\ChoiceMenu[radio,radiosymbol=\ding{52},bordercolor=0 0.5 0,name=cbovalid_ 841408 ]{}{CBO deve ficar=sim, \hspace{6em} CBO não fica aqui=nao}
\end{Form}

\newpage
\section{ 03003 -- Técnico em Biblioteconomia }

\textbf{Perfil do Curso (CNCT)}

O Técnico em Biblioteconomia será habilitado para:

Executar atividades técnico-administrativas e socioeducativas
relacionadas à rotina de bibliotecas e de centros de documentação e de
informação. Organizar e recuperar acervos físicos e virtuais. Atender ao
público e orientá-lo. Disseminar e organizar informações em ambientes
físicos e virtuais. Executar ações de conservação de documentos.
Organizar o espaço físico da biblioteca e/ou centro de informação e o
ambiente destinado ao usuário. Desenvolver projetos e ações
socioculturais. Colaborar na criação e aplicação de política de
desenvolvimento de coleção. Auxiliar no processamento técnico do acervo.
Manter e promover intercâmbio bibliográfico com outras unidades
congêneres. Auxiliar na divulgação dos produtos e serviços. Realizar
disseminação seletiva e ética da informação em meios físicos, virtuais e
digitais.

Para atuação como Técnico em Biblioteconomia, são fundamentais:

Conhecimentos e saberes relacionados aos processos de organização dos
espaços físicos de bibliotecas e centros de informação para
desenvolvimento de projetos e ações socioculturais. Conhecimentos e
saberes relacionados aos processos de colaboração, criação e aplicação
da política de desenvolvimento de coleção, bem como aos processos de
interpretação e aplicação de normas do exercício profissional.
Conhecimentos e saberes para identificar o perfil do usuário e
orientá-lo em relação à pesquisa em diversas fontes de informação e em
relação à utilização das ferramentas da tecnologia para construção de
banco de dados. Proatividade, liderança, organização, ética
profissional, confiança, empatia e tolerância. Comunicação assertiva com
diversos públicos e capacidade de promover a difusão de conhecimento e
melhoria do atendimento. Competência para manter e promover intercâmbio
bibliográfico com outras unidades.

\textbf{CBOs correspondidos e seus perfis (presentes no CNCT, e presentes na predição da IA)}

\textbf{CBO: 371105  ( Auxiliar de biblioteca )}

Atende o usuário nas formas presencial e a distância, informando a ele
sobre o regulamento, os recursos e a forma de funcionamento da
biblioteca ou unidade de documentação e de informação. Confecciona o
cartão de identificação do usuário. Empresta material do acervo ao
usuário, localizando o documento no acervo e controlando reserva,
empréstimo, renovação e devolução. Pode aplicar sanções ao usuário, pela
não devolução do material após o período de empréstimo. Orienta o
usuário na preservação do acervo. Participa do processo de disseminação
de informações, preparando painel para exposição das novas aquisições da
biblioteca, divulgando eventos culturais e ajudando na organização de
teleconferências. Auxilia na organização de atividades socioculturais,
tais como saraus culturais, exposições e feiras de livros. Pode operar
equipamentos audiovisuais, durante as atividades. Pode participar na
realização da biblioteca itinerante e ajudar nas campanhas de doação de
livros. Trata a documentação, auxiliando na seleção e no tombamento de
documentos, para incorporação no acervo. Confere a existência de
defeitos nos documentos. Magnetiza, etiqueta e carimba os documentos.
Digitaliza fotografias, jornais, artigos, monografia e outros
documentos, para inserção na base de dados da biblioteca. Participa da
manutenção do acervo, inventariando-o e o mantendo em ordem, de acordo
com o sistema de classificação adotado. Pode participar do remanejamento
do acervo. Realiza trabalhos de reparação de documentos, higienizando-os
e selecionando-os para encadernação. Prepara a encadernação e confere os
documentos encadernados. Verifica a inserção de documentos no acervo
digital. Executa atividades técnico-administrativas, expedindo malotes e
correios, controlando os estoques de material de consumo, auxiliando no
inventário de bens patrimoniais não bibliográficos e reproduzindo
documentos. Auxilia na elaboração de relatórios estatísticos e em
reuniões de planejamento e avaliação. Coleta dados estatísticos,
preenchendo planilhas. Participa na organização e na manutenção do
ambiente, arrumando a disposição do mobiliário e auxiliando no controle
do uso, disposição e manutenção dos equipamentos. Auxilia na aplicação
de novos conceitos de organização dos espaços da biblioteca, tornando o
ambiente mais interativo, informal e agradável, facilitando a pesquisa
física e, também, o uso de sistemas informatizados. Controla o fluxo de
usuários. Controla as condições de higiene e limpeza do ambiente.

\begin{Form}
\textbf{Esta correspondência está adequada?} \\ 
\ChoiceMenu[radio,radiosymbol=\ding{52},bordercolor=0 0.5 0,name=cbovalid_ 371105 ]{}{CBO deve ficar=sim, \hspace{6em} CBO não fica aqui=nao}
\end{Form}

\textbf{CBO: 371110  ( Técnico em biblioteconomia )}

Atende o usuário nas formas presencial e a distância, informando-o sobre
o funcionamento, o regulamento e os recursos da biblioteca ou unidade de
documentação e de informação. Cadastra o usuário, confeccionando seu
cartão de identificação. Atualiza periodicamente o cadastro de usuários.
Atende pedido do usuário, pesquisando o acervo para empréstimo de
material. Controla reserva, empréstimo, renovação e devolução de
material. Cobra devolução, após período de empréstimo, podendo aplicar
sanções ao usuário. Orienta o usuário na preservação do acervo. Auxilia
no trabalho escolar do usuário, orientando-o em pesquisa bibliográfica,
mostrando as normas de apresentação de trabalhos acadêmicos e ajudando
na editoração. Auxilia nas atividades de ensino a distância e na
organização de teleconferências. Monitora visitas à biblioteca.
Participa do estudo das demandas existentes e potenciais. Realiza
empréstimos entre bibliotecas, pesquisando bases de dados e localizando
material no acervo. Faz levantamentos bibliográficos, reservando
material. Organiza ações socioculturais e de extensão, participando da
elaboração de programas culturais em conjunto com a comunidade;
auxiliando na busca de parcerias para viabilizar a realização das
atividades; e fazendo contato com profissionais, lideranças e
instituições. Auxilia na organização, divulgação e realização de feiras
do livro, campanhas de doação e exposições. Controla a aquisição de
documentos, desdobrando e arquivando fichas catalográficas, auxiliando
na indexação e na elaboração de resumos. Participa do processo de
consistência e de alimentação da base de dados. Elabora clipping. Trata
documentação, participando na aquisição, seleção, catalogação,
classificação, cadastramento e tombamento de documentos, para
incorporação no acervo. Identifica publicações, participando da
organização da hemeroteca. Pode auxiliar na indexação de documentos no
acervo digital. Cadastra produção científica do corpo docente,
arquivando-a, física ou digitalmente. Realiza ou monitora a
digitalização de materiais. Realiza a manutenção do acervo,
inventariando-o e mantendo-o de acordo com o sistema de classificação
adotado. Participa do remanejamento do acervo. Controla permuta e acervo
de duplicatas de documentos. Substitui documentos. Executa trabalho de
reparação de documentos, selecionando-os para a encadernação e
conferindo os documentos encadernados. Realiza atividades
técnico-administrativas, participando na gestão administrativa da
biblioteca e em reuniões de planejamento e avaliação. Participa na
elaboração de projetos. Coleta dados, preenche planilhas e elabora
relatórios estatísticos. Colabora na elaboração do regimento interno da
biblioteca e na elaboração de manuais de procedimentos. Auxilia na
operação de sistemas de contratos eletrônicos. Elabora correspondências.
Organiza arquivos administrativos. Controla os estoques de material de
consumo, auxiliando na aquisição. Participa na aquisição de mobiliário e
equipamentos. Auxilia no inventário de bens patrimoniais não
bibliográficos. Pode utilizar sistema de gestão de biblioteca para
inserir e obter dados para o controle administrativo. Pode realizar
venda de publicações e materiais correlatos. Reproduz documentos.
Participa na manutenção do ambiente, organizando a disposição do
mobiliário e dos equipamentos. Elabora a sinalização do ambiente.
Auxilia no controle do uso e da manutenção dos equipamentos. Controla o
fluxo de usuários. Participa da aplicação de novos conceitos de
organização dos espaços da biblioteca, tornando o ambiente mais
interativo, informal e agradável, facilitando a pesquisa e o uso de
sistemas de informação. Controla as condições de higiene e limpeza do
ambiente. Auxilia no reaproveitamento de materiais e no descarte de
resíduos, de acordo com as normas ambientais.

\begin{Form}
\textbf{Esta correspondência está adequada?} \\ 
\ChoiceMenu[radio,radiosymbol=\ding{52},bordercolor=0 0.5 0,name=cbovalid_ 371110 ]{}{CBO deve ficar=sim, \hspace{6em} CBO não fica aqui=nao}
\end{Form}

\newpage
\section{ 03005 -- Técnico em Desenvolvimento Comunitário }

\textbf{Perfil do Curso (CNCT)}

O Técnico em Desenvolvimento Comunitário será habilitado para:

Organizar grupos de interesse em comunidades. Promover ações de
integração da comunidade e de aproximação positiva. Articular temáticas
de cultura, educação, esporte e lazer, meio ambiente, saúde, turismo,
trabalho e renda. Identificar potencialidades, necessidades, demandas
sociais, riscos e ameaças às condições de vida locais. Ler e interpretar
informações geradas a partir dos dados coletados no território e nos
aparelhos de serviços públicos e acadêmicos. Elaborar projetos e
programas sociais. Desenvolver ações temáticas para o desenvolvimento de
lideranças comunitárias. Executar campanhas socioeducativas.

Para atuação como Técnico em Desenvolvimento Comunitário, são
fundamentais:

Conhecimentos e saberes relativos à articulação de atores sociais, ao
fortalecimento da atuação em rede de organizações públicas e privadas.
Saberes relativos à valorização de memórias e identidades locais que
fortaleçam as expressões culturais comunitárias. Conhecimentos
relacionados ao associativismo e à democracia participativa, à
governança e à gestão territorial transparente e sustentável, às formas
de organização e de incidência política. Conhecimentos e habilidades
relacionados ao acesso e ao manejo de informações de interesse local, à
elaboração e à execução de projetos e programas comunitários. Capacidade
de lidar com imprevistos e construir soluções. Capacidade de administrar
conflitos e exercer a conciliação e a liderança. Proatividade,
criatividade, resiliência, flexibilidade e capacidade de persuasão.

\textbf{CBOs correspondidos e seus perfis (presentes no CNCT, e presentes na predição da IA)}

\textbf{CBO: 515305  ( Educador social )}

Entrevista pessoas individualmente, em famílias ou em grupos, avaliando
suas situações, capacidades e problemas, para definir quais serviços
podem atender às suas necessidades. Planeja ações em projetos ou
programas, definindo objetivos, metas, metodologias e estratégias de
trabalho, a partir do mapeamento das necessidades de grupos e pessoas ou
territórios atendidos. Atua em projetos ou programas, desenvolvendo
ações que estimulam a socialização e a construção de projetos de vida.
Desenvolve atividades recreativas, esportivas, culturais e de lazer,
para promoção de bem-estar e de melhor qualidade de vida. Encaminha
pessoas vulneráveis aos serviços da comunidade - como colocação de
emprego, aconselhamento de dívidas, assistência jurídica, moradia,
tratamento médico ou assistência financeira - fornecendo informações
sobre como, para onde ir e como se inscrever. Encaminha jovens a
programas de educação básica e de formação profissional, para
prepará-los para a vida e o trabalho. Aconselha indivíduos, grupos,
famílias ou comunidades sobre questões que incluem saúde mental,
pobreza, desemprego, abuso de substâncias, abuso físico, reabilitação,
ajuste social, assistência infantil ou assistência médica. Faz a
interlocução entre estudantes, lares, escolas, serviços familiares,
clínicas de orientação infantil, tribunais, serviços de proteção,
médicos e outros contatos, para ajudar crianças em situação de
vulnerabilidade pessoal e social. Orienta grupos e pessoas -
principalmente jovens e adolescentes - em situação de risco e/ou
vulnerabilidade social, sobre a possível violação de seus direitos e os
encaminha a entidades e a serviços de apoio. Quando necessário, denuncia
a situação e solicita resgates. Analisa práticas e avalia os resultados
das ações realizadas, elaborando relatório de atendimento e
acompanhamento. Mantém registros de histórico de casos. Estabelece
parcerias com entidades - públicas ou privadas -, para assegurar o
atendimento das pessoas, na perspectiva da proteção e defesa de seus
direitos.

\begin{Form}
\textbf{Esta correspondência está adequada?} \\ 
\ChoiceMenu[radio,radiosymbol=\ding{52},bordercolor=0 0.5 0,name=cbovalid_ 515305 ]{}{CBO deve ficar=sim, \hspace{6em} CBO não fica aqui=nao}
\end{Form}

\textbf{CBO: 515310  ( Agente de ação social )}

Auxilia no planejamento de projetos ou programas sociais, que visam
estimular a socialização e a construção de projetos de vida. Desenvolve
ações de natureza socioeducativa, atividades recreativas, esportivas,
culturais e de lazer para pessoas - em especial crianças, jovens e
idosos - em situação de risco e/ou vulnerabilidade social, violência e
exploração física e psicológica. Realiza atendimento interno - em
instituições - ou externo - em visitas domiciliares e espaços públicos
-, para identificação de possíveis situações de violações de direitos de
pessoas ou de grupos vulneráveis. Participa na verificação de denúncias
e de pedidos de ajuda de famílias em situação de risco, que poderão ser
encaminhadas a instituições e a serviços de apoio e proteção. Auxilia na
análise das práticas, na avaliação do resultado das ações e na
elaboração de relatórios sobre a efetividade do atendimento realizado.
Pode auxiliar na redefinição de ações e estratégias. Pode participar do
estabelecimento de parcerias com entidades públicas ou privadas.

\begin{Form}
\textbf{Esta correspondência está adequada?} \\ 
\ChoiceMenu[radio,radiosymbol=\ding{52},bordercolor=0 0.5 0,name=cbovalid_ 515310 ]{}{CBO deve ficar=sim, \hspace{6em} CBO não fica aqui=nao}
\end{Form}

\textbf{CBO: 515325  ( Socioeducador )}

Planeja e realiza ações socioeducativas para jovens infratores, que
possam contribuir para o desenvolvimento de suas competências para ser e
conviver. Acompanha a equipe técnica em visitas domiciliares, para
diálogo com educando e familiares e para identificação de necessidades.
Recepciona e acolhe o educando, criando vínculo que facilita o
reconhecimento das suas aptidões e habilidades, a sua conscientização
sobre regras e normas, a sua reflexão sobre riscos, o resgate de sua
autoestima e a consideração das suas alternativas de vida. Desenvolve
atividades recreativas, esportivas, culturais, de formação profissional,
de lazer, além de oficinas de atividades artísticas. Desenvolve com os
adolescentes atividades da vida diária -- alimentação, higiene pessoal e
do ambiente -, para criar hábitos saudáveis de convivência. Faz o
acompanhamento pedagógico dos adolescentes. Avalia a adesão a medidas
socioeducativas, contribuindo com a elaboração do diagnóstico
polidimensional e outros, como Plano Individual de Atendimento -- PIA.
Avalia processos de trabalho, considerando casos, ações, práticas e
estratégias. Elabora relatórios - com avaliação dos programas realizados
e com análise de estratégias utilizadas -, para subsidiar decisões sobre
adoção de ações socioeducativas cada vez mais eficazes. Quando atua em
instituições de privação de liberdade, realiza procedimentos de
segurança.

\begin{Form}
\textbf{Esta correspondência está adequada?} \\ 
\ChoiceMenu[radio,radiosymbol=\ding{52},bordercolor=0 0.5 0,name=cbovalid_ 515325 ]{}{CBO deve ficar=sim, \hspace{6em} CBO não fica aqui=nao}
\end{Form}

\newpage
\section{ 03006 -- Técnico em Infraestrutura Escolar }

\textbf{Perfil do Curso (CNCT)}

O Técnico em Infraestrutura Escolar será habilitado para:

Realizar a manutenção preventiva e corretiva de equipamentos e
instalações escolares. Organizar e conservar espaços físicos. Promover a
construção de hábitos de preservação e de manutenção do ambiente e do
patrimônio escolar. Identificar e buscar soluções para problemas de
infraestrutura. Organizar o espaço escolar.

Para atuação como Técnico em Infraestrutura Escolar, são fundamentais:

Conhecimentos e saberes relacionados aos processos de manutenção
preventiva e corretiva de equipamentos e instalações prediais escolares.
Capacidade de organização, conservação e preservação do ambiente e do
patrimônio escolar. Conhecimentos e habilidades para detectar problemas
e solucioná-los nas instalações prediais e de infraestrutura na
organização do espaço escolar. Habilidade e desenvoltura para o trabalho
em equipe e para a gestão de conflitos. Habilidade comunicativa.

\textbf{CBOs correspondidos e seus perfis (presentes no CNCT, e presentes na predição da IA)}

\textbf{CBO: 514325  ( Trabalhador da manutenção de edificações )}

Executa manutenção elétrica, preventiva e corretiva, para manter
equipamentos e instalações elétricas e de iluminação em perfeitas
condições de funcionamento. Identifica avarias e executa os reparos
necessários. Executa o serviço de manutenção hidráulica, verificando o
funcionamento, consertando e trocando instalações. Limpa e substitui
filtros, desentope ralos, pias e vasos sanitários, além de outros
reparos necessários. Executa serviços gerais de carpintaria, marcenaria
e alvenaria, reparando trincas e rachaduras, impermeabilizando
superfícies, recuperando pinturas, repondo cerâmica, consertando móveis,
substituindo portas, reparando divisórias, vedando fendas e emendas e
realizando ajustes diversos. Zela pela manutenção, organização e
conservação dos equipamentos e ferramentas. Controla o estoque de
produtos, materiais, peças, componentes, ferramentas e equipamentos,
verificando quantidades e data de validade. Providencia aquisição e
substituição, a fim de evitar interrupções das atividades de manutenção.
Orienta o trabalho de auxiliares de manutenção.

\begin{Form}
\textbf{Esta correspondência está adequada?} \\ 
\ChoiceMenu[radio,radiosymbol=\ding{52},bordercolor=0 0.5 0,name=cbovalid_ 514325 ]{}{CBO deve ficar=sim, \hspace{6em} CBO não fica aqui=nao}
\end{Form}

\newpage
\section{ 03007 -- Técnico em Laboratório de Ciências da Natureza }

\textbf{Perfil do Curso (CNCT)}

O Técnico em Laboratório de Ciências da Natureza será habilitado para:

Organizar laboratórios didáticos de ciências (física, química e
biologia) conforme normas de segurança, saúde ocupacional e preservação
ambiental. Proceder à montagem de experimentos reunindo equipamentos e
materiais de consumo para serem utilizados em aulas experimentais e
ensaios de pesquisa. Manusear equipamentos tecnológicos associados às
atividades desenvolvidas nos laboratórios de ciências da natureza.
Auxiliar nas atividades de práticas experimentais relacionadas aos
conhecimentos e saberes das ciências da natureza. Preparar soluções,
reagentes, peças, equipamentos e outros materiais utilizados em
experimentos. Proceder à limpeza e à conservação de instalações,
equipamentos e materiais dos laboratórios. Proceder ao controle de
estoque dos materiais de consumo dos laboratórios. Organizar e gerenciar
pequenos depósitos e/ou almoxarifados dos laboratórios. Utilizar
softwares interligados aos maquinários e equipamentos do laboratório.

Para atuação como Técnico em Laboratório de Ciências da Natureza, são
fundamentais:

Conhecimentos e saberes relacionados à física, química e biologia.
Habilidades de classificação de materiais e equipamentos, de acordo com
a finalidade de cada um, nas mais variadas atividades e experimentos.
Conhecimentos acerca dos equipamentos tecnológicos do laboratório e de
respectiva forma de manutenção básica. Conhecimento da norma de
segurança do trabalho e demais normas inerentes ao espaço de atuação.
Conhecimento dos diversos tipos de materiais e ferramentas utilizadas
para realizar o manuseio seguro e a limpeza com produtos adequados.
Noções de organização e catalogação de produtos e experimentos. Domínio
de equipamentos tecnológicos, ferramentais digitais e softwares
específicos para o desenvolvimento de atividades no laboratório.
Pensamento crítico e ação colaborativa. Competência comunicativa.\\
Proatividade, auto-organização e criatividade. Capacidade para mediar
conflitos e propor soluções para questões imprevistas.

\textbf{CBOs correspondidos e seus perfis (presentes no CNCT, e presentes na predição da IA)}

\textbf{CBO: 818110  ( Auxiliar de laboratório de análises físico-químicas )}

Auxilia na realização de análises físico-químicas no laboratório,
interpretando as ordens de serviço. Planeja as atividades, selecionando
método de análise, programando as etapas de trabalho e relacionando
materiais, equipamentos e instrumentos necessários. Verifica o
suprimento de materiais. Efetua cálculos, preenchendo fichas e
formulários. Prepara vidrarias para serem utilizadas em análises
físico-químicas, lavando-as, secando-as e esterilizando-as, a fim de
evitar contaminações. Identifica e armazena as vidrarias. Prepara as
soluções para as análises laboratoriais, selecionando vidrarias e
equipamentos e pesando ou medindo o volume dos reagentes. Mistura os
reagentes, para homogeneizar a solução preparada. Padroniza as soluções
e mede os parâmetros necessários para a sua caracterização. Armazena as
soluções, rotulando os frascos com as informações para a sua
identificação, seguindo as boas práticas de laboratório. Interpreta
manual de operações e prepara equipamentos de medição e ensaios,
identificando tensão elétrica. Seleciona e monta acessórios. Verifica a
aferição dos equipamentos. Coleta amostras para análise laboratorial,
interpretando as instruções do plano de amostragem. Etiqueta as amostras
e registra os seus dados, realizando os procedimentos para preservação
das características das amostras. Prepara as amostras para a análise,
organizando a bancada de trabalho, separando reagentes e soluções,
selecionando meios de cultura e outros insumos. Realiza análises,
interpretando dados. Emite laudos. Mantém arquivos físicos e/ou digitais
de fichas, formulários e laudos. Organiza literaturas técnicas, para
consulta. Conserva o ambiente limpo, higienizado e organizado.
Acondiciona e etiqueta materiais químicos e amostras para armazenamento.
Executa o descarte dos resíduos de acordo com as normas ambientais.
Limpa e organiza acessórios e instrumentos de trabalho. Solicita serviço
de manutenção de máquinas e equipamentos, quando necessário. Trabalha
com segurança, prevenindo acidentes e usando equipamentos de proteção
individual. Presta atenção aos rótulos de reagentes e obedece às
especificidades neles descritas.

\begin{Form}
\textbf{Esta correspondência está adequada?} \\ 
\ChoiceMenu[radio,radiosymbol=\ding{52},bordercolor=0 0.5 0,name=cbovalid_ 818110 ]{}{CBO deve ficar=sim, \hspace{6em} CBO não fica aqui=nao}
\end{Form}

\newpage
\section{ 03009 -- Técnico em Treinamento e Instrução de Cães-Guia }

\textbf{Perfil do Curso (CNCT)}

O Técnico em Treinamento e Instrução de Cães-Guia será habilitado para:

Treinar cães-guia para pessoas com deficiência visual ou cega.
Desenvolver e aplicar técnicas de adestramento que permitem ao cão
tornar-se apto à condução (mobilidade física) de pessoas com deficiência
visual ou cegas. Coordenar o processo de introdução do cão em família
socializadora. Selecionar matrizes de cães para servir de reprodutoras.
Selecionar filhotes para o ingresso em programa de cães-guia. Gerenciar
espaços definidos para reprodução e treinamento de cães para a atuação
como cães-guia. Formar duplas usuário x cão-guia.

Para a atuação como Técnico em Treinamento e Instrução de Cães-Guia, são
fundamentais:

Domínio e conhecimento da formação de treinador e de instrutor de
cães-guia e suas responsabilidades, bem como da gestão dos centros de
treinamento. Conhecimento acerca da legislação em vigor no país no que
se refere às pessoas com deficiência e correlata à formação em questão.
Domínio sobre o comportamento do cão-guia por meio do conhecimento da
fisiologia e comportamento específico dos cães na condução da pessoa com
deficiência visual. Saberes relacionados à seleção de cães que serão
treinados para a função de guia de pessoas com deficiência visual e para
as ?famílias socializadoras?, bem como saberes relacionados à condução,
com destreza, e à etapa de introdução e acompanhamento dos filhotes nas
referidas famílias. Saberes relacionados ao gerenciamento do processo
multidisciplinar de recrutamento e seleção de candidatos a usuários de
cães-guia. Conhecimentos e saberes relacionados ao processo de formação
das duplas entre a pessoa com deficiência visual e o cão-guia e de
acompanhamento até a aposentadoria do cão-guia. Competência
comunicativa, estabilidade emocional para atuação em situações
imprevistas, persistência, autoconfiança, autocontrole, foco e
constância para a realização do processo de treinamento e instrução de
cães-guia. Habilidade de estímulo para a percepção de circunstâncias
exteriores ao treinamento. Competência comunicativa e capacidade para
atuação com imprevistos. Empatia, sensibilidade e colaboração no
processo de introdução e de acompanhamento dos cães nas
comunidades/famílias.

\textbf{CBOs correspondidos e seus perfis (presentes no CNCT, e presentes na predição da IA)}

\textbf{CBO: 623005  ( Adestrador de animais )}

Programa o adestramento, analisando as informações disponíveis sobre o
animal, verificando o objetivo a ser atingido e estabelecendo as fases a
cumprir. Prepara o material de trabalho, examinando o estado de
conservação do instrumental necessário, para assegurar a execução das
atividades programadas. Estabelece contato com o animal, aclimatando-o
ao ambiente. Interage com animal, para familiarizá-lo com o contato e
com a voz humana. Amansa o animal. Avalia o animal para determinar seu
temperamento, sua capacidade e sua aptidão para o adestramento. Produz
técnicas de enriquecimento comportamental, recompensando comportamentos
esperados e restringindo comportamentos fora do padrão. Dá comandos aos
animais e estabelece níveis de autonomia e obediência, em processos
rotineiros e/ou automatizados. Condiciona o animal, induzindo, repetindo
e reforçando comportamentos. Condiciona a movimentação do animal em
itinerários e espaços determinados. Condiciona a conduta adequada do
animal em relação aos seres humanos, acompanhando e dando segurança à
pessoa, inclusive em vias e locais públicos. Condiciona a conduta do
animal em relação a outros animais, sociabilizando-o em grupos. Promove
comportamentos, ações e reações do animal em coletividade, objetivando a
interação saudável e amistosa em bandos e rebanhos. Cria situações e
objetos de ensino e aprendizagem, apresentando e representando contextos
de alimentação, busca e captura de objetos. Adestra animal para o
convívio doméstico. Desenvolve e aplica programas e rotinas de
exercícios e passeios com animais. Prepara animal para participar de
competições e de espetáculos de entretenimento. Pode realizar doma
racional de cavalos, técnica de adestramento que busca o entendimento
entre o domador e o animal, sem o uso da força e sem utilização de
técnicas violentas. Pode treinar animais para atuarem como farejadores.
Pode domar animais para lida de carga, tração e pastoreio. Pode
alimentar, lavar, higienizar e prestar outros cuidados ao animal. Pode
observar a condição física do animal para detectar doenças ou condições
insalubres que requeiram cuidados de veterinários. Pode usar circuito de
vigilância eletrônica, para observar comportamento e desempenho do
animal. Orienta o proprietário do animal. Presta informações ao
proprietário, tais como a forma de tratamento do animal e as
características da raça. Organiza e conserva limpo o local de
adestramento de animais. Acondiciona equipamentos, instrumentos e
acessórios. Preenche formulários e elabora relatórios, relatando
ocorrências e apresentando os resultados do adestramento. Pode prestar
primeiros socorros a animais.

\begin{Form}
\textbf{Esta correspondência está adequada?} \\ 
\ChoiceMenu[radio,radiosymbol=\ding{52},bordercolor=0 0.5 0,name=cbovalid_ 623005 ]{}{CBO deve ficar=sim, \hspace{6em} CBO não fica aqui=nao}
\end{Form}

\newpage
\section{ 03012 -- Técnico em Secretaria Escolar }

\textbf{Perfil do Curso (CNCT)}

O Técnico em Secretaria Escolar será habilitado para:

Assessorar os gestores escolares e suas equipes, aplicando as técnicas
secretariais em atividades de operacionalização
administrativas/financeiras e pedagógicas. Intermediar os
relacionamentos internos e externos, visando a metas e objetivos das
partes interessadas (Diretoria Escolar, Docentes, Discentes, Órgãos
Educacionais, Pais de Alunos, Fornecedores, Prestadores de Serviços).
Administrar, triar, manusear, armazenar e preservar informações gerais,
administrativas, financeiras e de legislação da instituição/escola, do
corpo docente e discente, tanto físicas quanto digitais, bem como
organizar os fluxos informacionais. Gerar e elaborar documentos
administrativos e pedagógicos da vida acadêmica dos alunos, como
matrícula e processos de transferências, registro e controle de
frequência, mapeamento do histórico escolar, tanto por meio físico
quanto digital, atendendo às orientações da direção, seguindo a
legislação em vigor e as exigências dos órgãos de regulação. Realizar,
de forma eficaz, a comunicação interna e externa. Organizar eventos
internos (reuniões, eventos educacionais, confraternizações). Organizar
e preparar viagens e/ou locomoção (na própria cidade/comunidade)
referentes às atividades e demandas da instituição de ensino. Atender
aos alunos, aos professores/educadores, à equipe administrativa, aos
dirigentes e líderes educacionais, aos pais, à comunidade e aos demais
profissionais envolvidos no processo educacional, além de manter uma boa
relação com eles.

Para atuação como Técnico em Secretaria Escolar, são fundamentais:

Conhecimentos e saberes relacionados às técnicas secretariais e da
administração; às noções de estatística e matemática financeira, de
planejamento estratégico; às técnicas de informática para utilizar
sistema operacional, pacote office, plataformas online, aplicativos,
equipamentos eletrônicos e de multimídias. Habilidades de liderança.
Eficácia para atingir resultados e capacidade para apoiar a diretoria na
concretização das atividades. Ética para o cumprimento das exigências
legais em território nacional. Competência comunicativa em idioma
nacional e estrangeiro - oral e escrita, para elaboração de textos
educacionais e corporativos. Noções sobre inteligência emocional para
administrar as emoções e alcançar os objetivos. Capacidade de
auto-organização, atuação com imprevistos e proposição de soluções.
Capacidade para o trabalho colaborativo e em equipe, para a comunicação
e a mediação de conflitos.

\textbf{CBOs correspondidos e seus perfis (presentes no CNCT, e presentes na predição da IA)}

\textbf{CBO: 351505  ( Técnico em secretariado )}

Realiza atividades secretariais de apoio ao gestor, auxiliando-o nas
suas atividades e no cumprimento de seus compromissos. Executa serviços
de apoio, reservando passagens aéreas e hotéis em viagens nacionais e
internacionais, prestando suporte administrativo e logístico em eventos
corporativos, elaborando pautas e atas de reuniões e realizando outras
atividades de acordo com a agenda e as prioridades definidas. Estabelece
os canais de comunicação do gestor com interlocutores. Recepciona e
atende clientes internos e externos, prestando-lhes informações,
encaminhando-os para o setor responsável pelo atendimento a suas
solicitações e verificando sua satisfação. Pode cadastrar e atualizar as
informações de clientes e fornecedores. Elabora cartas, memorandos,
ofícios, relatórios e outros documentos. Confere, revisa e corrige
textos, promovendo ajustes embasados em normas técnicas. Pode redigir
textos em idioma estrangeiro. Organiza, controla a circulação e arquiva
documentos. Pesquisa, localiza e recupera documentos, física ou
digitalmente. Pesquisa dados em fontes diversas, inclusive em páginas da
Internet. Organiza e sintetiza informações, podendo armazená-las em meio
físico ou eletrônico. Executa atividades financeiras de apoio aos
processos administrativos, orçando produtos e serviços, conferindo e
encaminhando pagamento, prestando contas e finalizando contratos e
aquisições de acordo com os processos organizacionais. Administra e
organiza recursos para execução dos trabalhos, requisitando material de
escritório, controlando estoque, providenciando manutenção de
equipamentos e aparelhos. Pode supervisionar equipe de trabalho,
distribuindo tarefas, acompanhando o andamento dos serviços e avaliando
o desempenho de cada profissional.

\begin{Form}
\textbf{Esta correspondência está adequada?} \\ 
\ChoiceMenu[radio,radiosymbol=\ding{52},bordercolor=0 0.5 0,name=cbovalid_ 351505 ]{}{CBO deve ficar=sim, \hspace{6em} CBO não fica aqui=nao}
\end{Form}

\newpage
\section{ 04001 -- Técnico em Administração }

\textbf{Perfil do Curso (CNCT)}

O Técnico em Administração será habilitado para:

Executar operações administrativas de planejamento, pesquisas, análise e
assessoria no que tange à gestão de pessoal, de materiais e produção, de
serviços, gestão financeira, orçamentária e mercadológica. Utilizar
sistemas de informação e aplicar conceitos e modelos de gestão em
funções administrativas, sejam operacionais, de coordenação, de chefia
intermediária ou de direção superior, sob orientação. Elaborar
orçamentos, fluxos de caixa e demais demonstrativos financeiros.
Elaborar e expedir relatórios e documentos diversos. Auxiliar na
elaboração de pareceres e laudos para tomada de decisões.

Para atuação como Técnico em Administração, são fundamentais:

Conhecimentos e saberes relacionados à área administrativa, com atuação
em conformidade com as legislações e diretrizes de conduta, como também
com as normas de saúde e segurança do trabalho, pautada em ações
empreendedoras e inovadoras, com foco em geração de novas oportunidades
de negócio e geração de renda. Exercício da profissão pautado no
comprometimento com necessidades, desejos e percepção da realidade
social de clientes, além de respeito à diversidade e à sustentabilidade.

\textbf{CBOs correspondidos e seus perfis (presentes no CNCT, e presentes na predição da IA)}

\textbf{CBO: 351305  ( Técnico em administração )}

Auxilia na elaboração, na implementação e no acompanhamento de
planejamento estratégico. Participa do planejamento e da execução de
melhorias nos processos organizacionais. Participa da execução dos
procedimentos de recrutamento, seleção e integração de pessoas. Executa
ações referentes às rotinas de admissão e de demissão. Acompanha
frequência, afastamentos, férias e processos de transferências de
pessoal. Executa ações referentes à avaliação de desempenho e ao
desenvolvimento de pessoas. Organiza propostas de treinamentos.
Participa da elaboração de planos de cargos e salários. Presta
informações acerca de direitos trabalhistas. Atua na organização e na
execução de ações relacionadas à qualidade de vida, saúde e segurança
nos ambientes de trabalho. Auxilia na elaboração da folha de pagamento,
lançando e conferindo informações na base de dados e gerando guias para
recolhimento de encargos sociais. Executa atividades nas áreas fiscal e
financeira, controlando fluxos de documentos fiscais, fiscalizando
recolhimento de encargos públicos de firmas terceirizadas e controlando
documentação de serviços terceirizados ao longo do contrato. Participa
da gestão de materiais e patrimônio, conferindo recebimento e
movimentação de materiais e controlando estoques. Participa do
planejamento e do desenvolvimento das ações de marketing. Atua na área
de compras, organizando processos de contratação de serviços,
participando de processos licitatórios e efetivando compras de produtos.
Assessora a área de vendas, recomendando margem de desconto, contatando
e orientando clientes e acompanhando venda e pós-venda. Elabora,
organiza e controla documentos da organização. Executa operações
administrativas relativas a protocolos e arquivos, confecção e expedição
de documentos. Acompanha validade de documentos legais. Auxilia na
execução de atividades relacionadas às operações logísticas. Opera
sistemas de informações gerenciais de pessoal, de finanças, de controle
patrimonial e de materiais. Pode apoiar a operacionalização de ações de
comércio exterior. Pode atuar na supervisão e no treinamento de equipe
de trabalho, avaliando desempenho e analisando necessidades de
treinamento.

\begin{Form}
\textbf{Esta correspondência está adequada?} \\ 
\ChoiceMenu[radio,radiosymbol=\ding{52},bordercolor=0 0.5 0,name=cbovalid_ 351305 ]{}{CBO deve ficar=sim, \hspace{6em} CBO não fica aqui=nao}
\end{Form}

\newpage
\section{ 04002 -- Técnico em Comércio }

\textbf{Perfil do Curso (CNCT)}

O Técnico em Comércio será habilitado para:

Aplicar métodos de comercialização de bens e serviços em loja física ou
virtual. Efetuar controle quantitativo e qualitativo de produtos, preços
e tributos. Coordenar e controlar a armazenagem em estabelecimento
comercial. Elaborar planilha de custos. Identificar demanda e comunicar
previsões a fornecedores. Ofertar serviços correlatos aos produtos
comercializados. Operacionalizar planos de marketing e de comunicação.
Executar atividades voltadas à logística, a recursos humanos e à
comercialização.

Para atuação como Técnico em Comércio, são fundamentais:

Conhecimentos e saberes relacionados ao funcionamento da área comercial
e de prestação de serviços, de modo a atuar em conformidade com as
legislações e diretrizes de conduta, como também com as normas de saúde
e segurança do trabalho. Atuação de forma proativa em atividades de
comercialização de produtos e serviços, com visão empreendedora,
comunicação clara e cordial, comprometimento com necessidades e desejos
de clientes e respeito a demais stakeholders.

\textbf{CBOs correspondidos e seus perfis (presentes no CNCT, e presentes na predição da IA)}

\textbf{CBO: 354140  ( Técnico em atendimento e vendas )}

Analisa o comportamento do mercado consumidor, participando do
desenvolvimento de estudos mercadológicos. Estuda o cliente potencial da
empresa, analisando seu perfil e suas expectativas, visando ao
estabelecimento de possíveis estratégias de fidelização. Avalia qual
estratégia de marketing é mais adequada para cada tipo de cliente.
Analisa vantagens e desvantagens no uso de comércio eletrônico
(e-commerce), dependendo do tipo de cliente. Pesquisa ações dos
concorrentes, para analisar o que é oferecido aos clientes. Providencia
o cadastramento de clientes potenciais e efetivos, atualizando os dados
sempre que necessário. Pode usar cadastro de clientes centralizado da
plataforma de Gestão de Relacionamento com o Cliente (Customer
Relationship Management - CRM). Contata clientes, usando estratégias
digitais (e-mail, redes sociais ou outras) e/ou realizando visitas.
Comunica-se com os clientes, identificando suas necessidades e
expectativas e analisando seu potencial de compras. Divulga produtos e
serviços aos clientes, apresentando características e benefícios do
produto ou serviço, esclarecendo dados técnicos, informando preços e
detalhando formas de pagamento. Pode fornecer amostras e convidar
clientes para visitar a empresa. Concretiza vendas presencialmente, por
telefone ou por meios digitais, de acordo com as características dos
clientes. Acompanha o cliente após a venda -- inclusive por meio de
visita -- para sanar suas dúvidas, providenciar apoio técnico e
verificar suas opiniões e suas sugestões. Analisa reclamações e
insatisfações, para propor melhorias no processo de vendas. Elabora
relatórios e estatísticas sobre atendimento ao cliente e sobre o
processo de vendas. Presta assistência à equipe de vendedores de
produtos e serviços, com relação aos procedimentos adequados a cada tipo
de cliente nas etapas de pré-venda, venda e pós-venda. Participa da
elaboração de planos e estratégias de vendas, sugerindo melhorias no
atendimento ao cliente e no processo de vendas.

\begin{Form}
\textbf{Esta correspondência está adequada?} \\ 
\ChoiceMenu[radio,radiosymbol=\ding{52},bordercolor=0 0.5 0,name=cbovalid_ 354140 ]{}{CBO deve ficar=sim, \hspace{6em} CBO não fica aqui=nao}
\end{Form}

\textbf{CBO: 521105  ( Vendedor em comércio atacadista )}

Realiza levantamento e análise de informações técnicas e pesquisa preços
do mercado, em relação às mercadorias em venda. Identifica as
necessidades dos potenciais compradores, no contexto do comércio
atacadista. Seleciona os recursos necessários às ações de vendas, de
acordo com as estratégias comerciais e de marketing da empresa.
Estabelece contato com clientes potenciais -- presencialmente ou on-line
-, apresentando a oferta de produtos, dentro da área de interesse e
necessidade de compra de cada um deles. Recebe e examina editais de
licitação de órgãos públicos, podendo apresentar, ao seu supervisor,
proposta de participação no processo licitatório. Recepciona e atende o
comprador, orientando-o na escolha da mercadoria. Apresenta as
características, os benefícios e o preço de cada opção de produto,
detalhando formas de pagamento. Consulta a disponibilidade de produtos
escolhidos, em estoque. Apresenta opções de data e de forma de entrega.
Prepara documentos para o fechamento da venda, conforme procedimentos
fiscais e da empresa. Realiza ações de pós-venda. Pode analisar
resultados de pesquisa de satisfação, para levantar subsídios para
melhoria do processo de venda. Pode promover ações de fidelização, de
acordo com o perfil de cada cliente. Encaminha reclamações, sugestões e
solicitações ao setor responsável. Elabora ou atualiza o cadastro com
dados e informações sobre o comprador, em sistema informatizado. Pode
concretizar vendas on-line por correspondência (e-mail), telefone, chat
ou aplicativos, tanto para vendas no país como para o exterior. Pode
contatar clientes para eventuais cobranças. Elabora relatórios de
visitas a clientes e de vendas presenciais e on-line.

\begin{Form}
\textbf{Esta correspondência está adequada?} \\ 
\ChoiceMenu[radio,radiosymbol=\ding{52},bordercolor=0 0.5 0,name=cbovalid_ 521105 ]{}{CBO deve ficar=sim, \hspace{6em} CBO não fica aqui=nao}
\end{Form}

\textbf{CBO: 521115  ( Promotor de vendas )}

Efetua pesquisas sobre o potencial de mercado e a atuação da
concorrência, para subsidiar reuniões de definição de estratégias de
comercialização de linhas de produto. Prospecta um novo ponto de venda
(PDV), em sua região de atuação. Providencia a limpeza, a montagem e a
organização do ponto de venda. Orienta o abastecimento de mercadorias e
produtos. Orienta a montagem de áreas de exposição, considerando o
leiaute do ponto de venda e o plano de merchandising. Pode abordar e
orientar clientes sobre mercadorias, produtos e serviços no ponto de
venda, seguindo técnicas de atendimento. Acompanha o desempenho e os
resultados do ponto de venda. Promove produtos, visando ao aumento de
vendas e à fixação da marca. Realiza visitas -- presenciais e virtuais -
de divulgação de linhas de produto e de apresentação de estratégias de
comercialização, junto ao comércio varejista e a clientes em
perspectiva. Presta assistência aos revendedores e às equipes de venda,
demonstrando os produtos e fornecendo informações sobre seus manuseios e
usos. Em caso de solicitação de informações específicas de
aplicabilidade de um produto, pode fazer encaminhamento das questões
para a área técnica da empresa produtora. Pode auxiliar o vendedor nas
vendas do produto promovido. Pode distribuir brindes e oferecer
degustação de produtos. Pode promover produtos em exposições, feiras e
eventos. Elabora e organiza documentação em meio impresso e digital,
para a promoção de vendas junto ao comércio varejista. Prepara
documentação subsidiária para orientação a vendedores dos pontos de
venda, em meio impresso e digital.

\begin{Form}
\textbf{Esta correspondência está adequada?} \\ 
\ChoiceMenu[radio,radiosymbol=\ding{52},bordercolor=0 0.5 0,name=cbovalid_ 521115 ]{}{CBO deve ficar=sim, \hspace{6em} CBO não fica aqui=nao}
\end{Form}

\newpage
\section{ 04003 -- Técnico em Comércio Exterior }

\textbf{Perfil do Curso (CNCT)}

O Técnico em Comércio Exterior será habilitado para:

Prestar apoio a análises de mercado por intermédio da aplicação de
regras e políticas cambiais específicas de países envolvidos nas
negociações. Executar e controlar atividades inerentes ao processo de
exportação e importação. Cumprir os trâmites aduaneiros em operações de
importação e exportação. Elaborar cálculos de custos, preços e tributos.
Utilizar canais informatizados de órgãos reguladores, como Receita
Federal, Siscomex e Inmetro. Executar procedimentos de transporte,
armazenamento e logística internacional.

Para atuação como Técnico em Comércio Exterior, são fundamentais:

Conhecimentos e saberes relacionados à área de exportação e importação,
com atuação em conformidade com as legislações e diretrizes de conduta,
como também com as normas de saúde e segurança do trabalho, pautada em
ações empreendedoras e inovadoras, com foco em geração de novas
oportunidades de negócio internacionais. Atuação com base na construção
de relacionamentos positivos, respeito à diversidade, atenção à
sustentabilidade, trabalho em equipe e resolução efetiva de conflitos.

\textbf{CBOs correspondidos e seus perfis (presentes no CNCT, e presentes na predição da IA)}

\textbf{CBO: 342105  ( Analista de transporte em comércio exterior )}

Controla, programa e coordena operações em diferentes modais de
transportes, incluindo embarque, transbordo e desembarque e carga em
regime de trânsito aduaneiro. Verifica as condições de segurança dos
meios de transportes e dos equipamentos utilizados, como também do
acondicionamento, embalagem e etiquetagem da carga. Realiza pesquisa de
mercado, consultando preços de serviços de transporte, analisando a
viabilidade de novos produtos e serviços, identificando clientes
potenciais, rotas de transporte e volume e demanda de carga, utilizando
softwares e sistemas de informação especializados. Define, contrata e
negocia parceria com fornecedores de serviços, atualizando o cadastro.
Atende clientes, organizados por demanda, região e tipo de produto,
elaborando proposta de serviços e verificando sua satisfação. Controla
recursos financeiros e insumos, calculando custos e receitas, conferindo
faturas de pagamentos e recebimentos de serviços, bem como elaborando
documentação para desembaraço de cargas. Auxilia na elaboração de
contratos, confere, organiza e arquiva documentos de transporte de
comércio exterior, alimentando sistemas informatizados internos e
externos. Remaneja equipamentos e veículos, requisitando manutenção,
quando necessária. Orienta funcionários e estagiários, avaliando seus
desempenhos e analisando suas necessidades de treinamento.

\begin{Form}
\textbf{Esta correspondência está adequada?} \\ 
\ChoiceMenu[radio,radiosymbol=\ding{52},bordercolor=0 0.5 0,name=cbovalid_ 342105 ]{}{CBO deve ficar=sim, \hspace{6em} CBO não fica aqui=nao}
\end{Form}

\textbf{CBO: 342210  ( Despachante aduaneiro )}

Executa atividades relacionadas ao despacho aduaneiro de mercadorias -
inclusive bagagem de viajante - na importação, na exportação e na
internação, realizadas por qualquer via de transporte. Elabora e/ou
supervisiona a elaboração de documentos relativos ao despacho aduaneiro.
Providencia credenciamento de cliente junto à Receita Federal. Analisa
cartas de crédito. Elabora fatura proforma e fatura comercial. Prepara
lista de embarque. Executa e/ou supervisiona a classificação de
mercadorias, analisando a documentação técnica das mercadorias;
verificando suas funções de uso e aplicação; e observando seu material
constitutivo. Analisa necessidade de anuência de órgãos pertinentes e,
quando necessário, consulta a Receita Federal sobre dúvidas de
classificação. Analisa enquadramento legal das alíquotas de impostos de
importação e exportação. Enquadra mercadorias na Tarifa Externa Comum
(TEC), na Nomenclatura da Associação Latino-americana de Integração
(Naladi) e na Tabela de Incidência do Imposto sobre Produtos
Industrializados (TIPI). Preenche Certificado de Origem; Conhecimento de
Transporte internacional; Documento de Arrecadação de Receitas Federais
(DARF); guia de Imposto sobre Operações relativas à Circulação de
Mercadorias e sobre Prestações de Serviços de Transporte Interestadual e
Intermunicipal e de Comunicação (ICMS); além de outros formulários
necessários às operações de comércio exterior. Realiza e/ou supervisiona
a atividade de entrada de dados no Sistema Integrado de Comércio
Exterior (Siscomex) e sistemas auxiliares. Coleta dados de operações e
registra licenças e declarações de importação e de despacho de
exportação. Efetua registros no módulo Registro de Operações Financeiras
do Sistema Registro Declaratório Eletrônico (RDE-ROF) e de atos
concessórios de regime aduaneiro especial (``Drawback''). Faz Declaração
de Trânsito Aduaneiro (DTA). Efetua acompanhamento da tramitação de
documentos relacionados ao despacho aduaneiro. Solicita deferimento de
licenças de importação e de declarações retificadoras. Subscreve
documentos, inclusive termos de responsabilidade. Faz recebimento de
intimações, de notificações, de autos de infração, de despachos, de
decisões e de outros atos e termos processuais relacionados com o
procedimento de despacho aduaneiro. Realiza ou supervisiona o
recebimento de mercadorias e bagagens. Paga impostos e taxas. Pode
apresentar desistência de vistoria aduaneira. Executa acompanhamento de
vistoria da mercadoria na conferência aduaneira, observando a
conferência física e a retirada de amostra de mercadoria para
assistência técnica e perícia. Avalia exigências da receita federal e
demais órgãos pertinentes, presta esclarecimentos e apresenta
documentação. Retira comprovantes de importação e exportação junto à
receita federal. Presta assessoria a importadores e exportadores,
identificando suas necessidades; informando sobre legislação pertinente;
propondo alternativas de logística de transportes; e avaliando
benefícios fiscais. Calcula custos de operação. Pleiteia benefícios
fiscais. Monitora processos administrativos e informa a cliente sobre
tramitação de processos. Realiza contratação de serviços de terceiros,
supervisionando a cotação de preços de armazenagem e preços de fretes
nacionais e internacionais; orçando seguros; e negociando preços e
prazos. Confere cobrança de terceiros por serviços prestados e
providencia o pagamento. Supervisiona equipe de trabalho, distribuindo
tarefas, avaliando desempenho e desenvolvendo treinamentos. Monitora a
limpeza e a organização do ambiente de trabalho. Controla a aplicação de
procedimentos da empresa para cumprimento de normas ambientais,
inspecionando o reaproveitamento de sobras de material e o descarte de
resíduos.

\begin{Form}
\textbf{Esta correspondência está adequada?} \\ 
\ChoiceMenu[radio,radiosymbol=\ding{52},bordercolor=0 0.5 0,name=cbovalid_ 342210 ]{}{CBO deve ficar=sim, \hspace{6em} CBO não fica aqui=nao}
\end{Form}

\textbf{CBO: 351310  ( Técnico em administração de comércio exterior )}

Auxilia no processo de negociações comerciais internacionais,
considerando as condicionantes legais, verificando a adequação de
produtos e serviços diante de aspectos culturais e ambientais,
apresentando possíveis parcerias e fornecedores e estimando custos das
importações e preço das exportações. Planeja as operações de transporte
nas transações de comércio exterior, analisando a logística de
transporte, verificando a segurança dos equipamentos e cargas nos modais
de transporte e identificando as melhores formas de acondicionamento,
embalagem e etiquetagem da carga. Seleciona modalidade de transporte
para cada trajeto. Programa, coordena e monitora as operações em
diferentes modais de transportes, incluindo embarque, transbordo e
desembarque. Planeja e seleciona o tipo de armazenamento, de acordo com
as características da carga e com a melhor relação custo-benefício.
Programa e monitora as operações do processo de armazenamento. Solicita
e acompanha as operações de câmbio realizadas por operadoras e bancos.
Providencia a documentação de câmbio, seguindo a legislação cambial
vigente. Elabora a documentação que instrui os processos de exportação.
Confere a documentação de importação, de acordo com a negociação, a
legislação e as normas técnicas nacionais e internacionais. Formula o
despacho aduaneiro de importação ou de exportação, de acordo com a
documentação vigente e as normas aduaneiras. Confere, organiza e arquiva
documentos físicos e digitais. Controla recursos financeiros e insumos,
calculando custos de exportação e/ou importação, verificando despesas e
receitas e conferindo faturas de pagamentos e recebimentos de serviços.
Faz uso de sistemas de informação especializados para a área de comércio
internacional. Pode atuar na supervisão e no treinamento de equipe de
trabalho, avaliando desempenho e analisando necessidades de treinamento.

\begin{Form}
\textbf{Esta correspondência está adequada?} \\ 
\ChoiceMenu[radio,radiosymbol=\ding{52},bordercolor=0 0.5 0,name=cbovalid_ 351310 ]{}{CBO deve ficar=sim, \hspace{6em} CBO não fica aqui=nao}
\end{Form}

\textbf{CBO: 354305  ( Analista de exportação e importação )}

Analisa mercado internacional de produtos e serviços, para identificar
oportunidades de negócio e fornecer subsídios para definição de
políticas de importação ou exportação. Examina aspectos cambiais e
creditícios. Analisa aspectos legais e normativos. Verifica as
tendências de consumo nos países, para adequar produtos e embalagens aos
diferentes mercados. Analisa logística de movimentação. Elabora o plano
de exportação, estabelecendo tipos e quantidades de mercadorias e
avaliando concorrentes. Participa da definição de preços de produtos e
serviços. Avalia a capacidade produtiva da mercadoria a exportar, para
projetar curvas de vendas e definir logística. Operacionaliza a
exportação, negociando condições de venda de produtos e serviços.
Elabora proposta comercial, incluindo especificações técnicas do produto
e tipos de embalagens e acondicionamento. Comunica ao produtor a
aceitação da proposta, acertando a compatibilização do cronograma de
produção com o contrato de venda. Contrata despachante, elabora
documentação para exportação e providencia pagamento de encargos.
Acompanha o andamento da produção de bens e serviços, para liberação da
mercadoria comprometida em tempo hábil de embarque. Contrata serviços de
terceiros para realização de transporte da mercadoria. Providencia meios
do recebimento do valor acordado, por intermédio de transações em moeda
estrangeira. Conclui o processo junto aos órgãos competentes,
acompanhando a finalização no Sistema Integrado de Comércio Exterior --
SISCOMEX. Guarda a documentação pelo prazo legal. Processa a importação
de mercadorias, estudando a viabilidade do processo e emitindo planilha
de custos. Verifica a conformidade da documentação e da nomenclatura das
mercadorias, em relação à legislação vigente. Faz o pedido das
mercadorias. No recebimento, instrui o desembaraço aduaneiro e verifica
a documentação. Pode negociar com fornecedores a troca de mercadorias
com defeito ou em desacordo com o pedido. Providencia meios de pagamento
por intermédio de transações em moeda estrangeira. Aciona a seguradora,
quando necessário. Realiza atividades de promoção de comércio exterior,
organizando missões empresariais, prestando assistência técnica e
operacional aos clientes em feiras e eventos e providenciando propaganda
de produtos e serviços. Atende consultas de clientes, orienta
exportadores e importadores e presta informações sobre o andamento de
processos de importação ou exportação. Desenvolve fornecedores. Pode
ministrar palestras técnicas da área.

\begin{Form}
\textbf{Esta correspondência está adequada?} \\ 
\ChoiceMenu[radio,radiosymbol=\ding{52},bordercolor=0 0.5 0,name=cbovalid_ 354305 ]{}{CBO deve ficar=sim, \hspace{6em} CBO não fica aqui=nao}
\end{Form}

\newpage
\section{ 04004 -- Técnico em Condomínio }

\textbf{Perfil do Curso (CNCT)}

O Técnico em Condomínio será habilitado para:

Executar atividades administrativas voltadas a recursos humanos,
financeiros e de gestão de suprimentos e materiais, de acordo com
convenção condominial. Supervisionar a conservação e manutenção predial.
Conduzir reuniões e assembleias de condomínios. Elaborar atas e
relatórios de prestação de contas. Promover a integração dos condôminos.

Para atuação como Técnico em Condomínio, são fundamentais:

Conhecimentos e saberes relacionados aos processos de gestão
condominial, de modo a atuar em conformidade com as legislações e
diretrizes de órgãos reguladores, como também com as normas de saúde e
segurança do trabalho. Atuação pautada na construção de relacionamentos
positivos, atenção à sustentabilidade, resolução efetiva de conflitos e
respeito à diversidade.

\textbf{CBOs correspondidos e seus perfis (presentes no CNCT, e presentes na predição da IA)}

\textbf{CBO: 510110  ( Administrador de edifícios )}

Planeja rotinas de trabalho em administração de edifícios e condomínios,
elaborando cronogramas de execução de tarefas, definindo as funções da
equipe de trabalho e programando o uso de equipamentos. Descreve
procedimentos para execução dos serviços, levando em consideração normas
do condomínio. Atende condôminos, verificando a pertinência de
reclamações, registrando solicitações e solucionando os problemas
apresentados. Atua para promover a integração dos condôminos. Estabelece
a grade de horário da portaria, estabelecendo os turnos de trabalho.
Orienta e avalia o comportamento dos porteiros com os moradores. Pode
substituir a presença de porteiros por serviços de portaria remota,
cumprindo decisão de assembleia condominial. Supervisiona a conservação
predial, verificando o estado das instalações e providenciando os
pequenos reparos necessários. Em caso de reforma, realiza a parte que
lhe compete, no cumprimento da norma da ABNT- Associação Brasileira de
Normas Técnicas, que estabelece como as obras no edifício -- inclusive
no interior de uma unidade residencial ou de um conjunto comercial -
devem ser acompanhadas, planejadas e executadas. Monitora o
funcionamento dos equipamentos, providenciando reparos e contratando
serviços de manutenção corretiva, quando necessário. Pode apresentar
proposta de atualização de equipamentos. Supervisiona a limpeza e a
organização das áreas comuns do edifício. Verifica os resultados dos
serviços de dedetização. Controla o consumo de produtos e materiais,
realizando a aquisição para reposição no estoque. Controla os recursos
financeiros, verificando as despesas e as receitas previstas e
realizadas. Coordena a equipe de trabalho, dimensionando pessoal para
cada setor, atribuindo tarefas, delegando responsabilidades e orientando
a execução dos serviços. Controla conflitos na equipe. Administra
recursos humanos, participando da seleção de pessoal, controlando
absenteísmo e estabelecendo escala de horários e folgas. Supervisiona a
equipe de trabalho, avaliando o desempenho e identificando o potencial
das pessoas. Controla tempo e qualidade na realização de tarefas. Pode
solicitar autoavaliação de cada membro da equipe. Realiza reuniões de
avaliação da equipe. Levanta as necessidades de treinamento,
providenciando as atividades de capacitação. Instrui a equipe sobre
regulamento do condomínio e organiza treinamento para execução de novos
serviços e para uso de novos equipamentos e produtos. Simula atividades
para aprendizagem. Elabora relatórios de operação e de avaliação,
preparando informações sobre quebra e reposição de materiais,
registrando ocorrências nos setores e apresentando o nível de satisfação
dos condôminos. Participa de reuniões e assembleias de condomínios. Pode
elaborar atas e relatórios de prestação de contas. Controla
desperdícios, com o reaproveitamento do material, sempre que possível.
Orienta o descarte de resíduos de acordo com as normas ambientais. Zela
pela segurança das pessoas, assegurando o cumprimento das normas;
identificando e eliminando as condições inseguras; orientando e
conscientizando a equipe para o uso de Equipamentos de Proteção
Individual (EPI) e Equipamentos de Proteção Coletiva (EPC).

\begin{Form}
\textbf{Esta correspondência está adequada?} \\ 
\ChoiceMenu[radio,radiosymbol=\ding{52},bordercolor=0 0.5 0,name=cbovalid_ 510110 ]{}{CBO deve ficar=sim, \hspace{6em} CBO não fica aqui=nao}
\end{Form}

\newpage
\section{ 04005 -- Técnico em Contabilidade }

\textbf{Perfil do Curso (CNCT)}

O Técnico em Contabilidade será habilitado para:

Executar processos administrativos e contábeis. Classificar documentos
contábeis, fiscais e não fiscais. Calcular tributos federais, estaduais
e municipais. Prestar atendimento à fiscalização e apresentar
documentos, livros e relatórios contábeis. Elaborar planos de
determinação das taxas de depreciação e exaustão dos bens materiais e de
amortização dos valores imateriais. Ordenar os fatos contábeis por
débito e crédito. Apurar haveres, direitos e obrigações legais.

Para atuação como Técnico em Contabilidade, são fundamentais:

Conhecimentos e saberes relacionados aos processos financeiros e
contábeis empresariais, de modo a atuar em conformidade com as
legislações e diretrizes de órgãos reguladores, como também com as
normas de saúde e segurança do trabalho, sempre sob a supervisão de um
contabilista. Atuação pautada em decisões responsáveis baseadas em
conceitos éticos construtivos e relacionamentos positivos, trabalho em
equipe e resolução efetiva de conflitos.

\textbf{CBOs correspondidos e seus perfis (presentes no CNCT, e presentes na predição da IA)}

\textbf{CBO: 351105  ( Técnico de contabilidade )}

Realiza processos administrativos e contábeis para constituição de uma
empresa, efetuando o enquadramento tributário. Faz registros da empresa
junto aos órgãos públicos competentes e providencia documentos para
envio à junta comercial e/ou aos cartórios. Identifica, organiza e
confere documentos fiscais e contábeis. Analisa a documentação,
verificando a sua validade e conformidade e separando-a de acordo com a
sua natureza. Classifica os documentos, em função do seu conteúdo, e
registra os dados referentes à sua movimentação. Efetua lançamentos
contábeis, elabora balancetes de verificação e concilia contas.
Escritura livros contábeis e fiscais, com o uso de programas e sistemas
informatizados. Realiza a Escrituração Contábil Digital (ECD) para
encaminhar informações contábeis e fiscais -- tais como Livro Diário,
Livro Razão, Livro Balancetes Diários, Balanços e fichas de lançamento
-- à Receita Federal do Brasil, pelo Sistema Público de Escrituração
Digital (Sped). Apura impostos, juros, taxas e tarifas a receber e a
pagar. Atende a obrigações fiscais acessórias. Presta informações para
auditorias interna e externa. Operacionaliza a contabilidade de custos,
relacionando custos operacionais e não operacionais. Apura estoques.
Contabiliza custos orçados em comparação aos realizados. Elabora
relatório de custo. Efetua a contabilidade gerencial, realizando
previsão orçamentária, compilando informações contábeis, analisando
evolução de contas, efetuando análises comparativas e preparando fluxo
de caixa. Executa planejamento tributário e elabora balanço social.
Realiza o controle patrimonial, verificando a entrada e a saída de
ativos imobilizados e de outros itens patrimoniais. Calcula depreciação
e reavaliação de bens. Calcula juros sobre patrimônio em formação.
Amortiza gastos e custos incorridos. Procede equivalência patrimonial.
Realiza inventário do patrimônio. Elabora balanço patrimonial. Presta
atendimento à fiscalização por intermédio da apresentação de documentos,
livros e relatórios contábeis. Pode entrar em contato com os órgãos de
fiscalização, para regularizar documentos da empresa ou instituição.
Pode prestar consultoria empresarial, examinando os documentos fiscais e
a documentação contábil, diagnosticando os problemas da empresa e
reestruturando o plano de contas. Pode colaborar na gestão de processos
relativos a recursos humanos, tais como verificação dos termos de
dissídios coletivos; formalização de contratos de trabalho;
processamento de folha de pagamento; cálculos de rescisão de contrato de
trabalho; e preenchimento de relatórios de informações sociais para
órgãos governamentais. Interage com as demais áreas da empresa ou
instituição que possuem interface com a área contábil, para alinhamento
de conceitos e ações. Organiza e arquiva -- em meio físico e eletrônico
-- informações e documentos administrativos e contábeis.

\begin{Form}
\textbf{Esta correspondência está adequada?} \\ 
\ChoiceMenu[radio,radiosymbol=\ding{52},bordercolor=0 0.5 0,name=cbovalid_ 351105 ]{}{CBO deve ficar=sim, \hspace{6em} CBO não fica aqui=nao}
\end{Form}

\textbf{CBO: 351110  ( Chefe de contabilidade (técnico) )}

Prepara-se para as atividades de supervisão, interpretando mudanças em
leis e normas relacionadas às áreas contábil, fiscal e trabalhista.
Supervisiona a elaboração de atos constitutivos de novas empresas. Atua
no encerramento de empresas. Coordena as atividades de contabilidade
geral, formando peças contábeis. Desenvolve plano de contas. Atende a
obrigações fiscais acessórias. Monitora a escrituração de livros
contábeis e fiscais, com o uso de programas e sistemas informatizados.
Supervisiona a elaboração da Escrituração Contábil Digital (ECD) e o
encaminhamento das informações contábeis e fiscais -- tais como Livro
Diário, Livro Razão, Livro Balancetes Diários, Balanços e fichas de
lançamento -- à Receita Federal do Brasil, pelo Sistema Público de
Escrituração Digital (Sped). Coordena a operacionalização da
contabilidade de custos, identifica custo gerencial e administrativo e
demonstra custo incorrido e/ou orçado. Coordena a execução da
contabilidade gerencial e analisa o comportamento das contas. Monitora a
execução do planejamento tributário e a elaboração do balanço social.
Acompanha resultados e efetua análises comparativas, fornecendo
subsídios aos administradores da empresa ou instituição. Coordena a
realização do controle patrimonial, analisando a entrada e a saída de
ativos imobilizados e o inventário do patrimônio. Acompanha a elaboração
do balanço patrimonial. Colabora com a gestão da área de recursos
humanos, analisando as mudanças da legislação trabalhista e verificando
os termos dos dissídios coletivos. Participa da implantação de plano de
cargos e salários. Atende à fiscalização, prestando informações e
esclarecimentos, apresentando documentos, livros e relatórios e, quando
necessário, providenciando defesa administrativa. Pode participar de
equipe de assessoria e consultoria contábil, examinando documentos
fiscais, analisando documentação trabalhista, societária e contábil e
diagnosticando problemas. Propõe reestruturação de planos de contas e
outras ações necessárias. Administra equipe de trabalho da área
contábil, participando da seleção de pessoal e controlando absenteísmo.
Observa o cumprimento de normas e procedimentos da empresa ou
instituição. Supervisiona o desempenho da equipe, monitorando e
avaliando o seu trabalho. Identifica necessidade de atualização ou
aperfeiçoamento de membros da equipe. Treina ou providencia programação
de treinamento de pessoal. Interage com os supervisores de outras áreas
da empresa ou instituição, para troca de informações e análise de
resultados. Elabora relatórios - com gráficos e análise de informações -
para gestores da empresa ou instituição. Pode utilizar Sistema Integrado
de Gestão Empresarial - ERP ou outro sistema de gestão. Examina a
organização dos arquivos -- em meio físico e eletrônico -- de
informações e documentos administrativos e contábeis. Verifica a limpeza
e a organização do local de trabalho da equipe de contabilidade.
Controla o reaproveitamento de materiais e o descarte de resíduos, de
acordo com normas ambientais.

\begin{Form}
\textbf{Esta correspondência está adequada?} \\ 
\ChoiceMenu[radio,radiosymbol=\ding{52},bordercolor=0 0.5 0,name=cbovalid_ 351110 ]{}{CBO deve ficar=sim, \hspace{6em} CBO não fica aqui=nao}
\end{Form}

\textbf{CBO: 351115  ( Consultor contábil (técnico) )}

Participa da realização de diagnóstico dos problemas da empresa ou
instituição, analisando a documentação societária, examinando documentos
fiscais, inventariando documentação trabalhista e estudando a
documentação contábil. Participa da elaboração de proposta de serviços
de consultoria, propondo medidas para solucionar os problemas nas áreas
contábil, fiscal e trabalhista. Em tal proposta, leva em consideração a
legislação vigente - das esferas municipal, estadual e federal -
relacionada às áreas fiscal, tributária e trabalhista. Presta
consultoria de contabilidade, indicando procedimentos padronizados e
legais para execução de atividades como registros, lançamentos,
classificação e conferência de documentos fiscais e contábeis. Orienta
sobre procedimentos de elaboração de balancetes, fluxo de caixa,
demonstrativos de resultado e balanço patrimonial. Mostra os
procedimentos para conciliação de contas, apuração de impostos, cálculo
de depreciação de bens e apuração de estoques. Orienta a análise da
evolução de contas e a elaboração de previsão orçamentária e de
planejamento tributário. Mostra o correto preenchimento de declarações
de impostos e outras obrigações contábeis legais. Orienta a
reestruturação do plano de contas e a definição de centros de custos.
Apresenta parâmetros legais e procedimentos adequados para efetivar a
gestão de processos relativos a recursos humanos. Expõe proposta de
utilização de aplicativos específicos para gestão contábil e financeira.
Orienta a utilização de sistemas para transmissão de dados e informações
aos diferentes órgãos, como a Receita Federal do Brasil. Apresenta
proposta de treinamento de equipe que, na empresa ou instituição,
operacionalizará as medidas propostas pela consultoria.

\begin{Form}
\textbf{Esta correspondência está adequada?} \\ 
\ChoiceMenu[radio,radiosymbol=\ding{52},bordercolor=0 0.5 0,name=cbovalid_ 351115 ]{}{CBO deve ficar=sim, \hspace{6em} CBO não fica aqui=nao}
\end{Form}

\newpage
\section{ 04007 -- Técnico em Finanças }

\textbf{Perfil do Curso (CNCT)}

O Técnico em Finanças será habilitado para:

Realizar operações relativas a serviços e produtos financeiros de
empréstimos, financiamentos, investimentos e bancários. Elaborar e
analisar fluxos de caixa e demais relatórios financeiros. Efetuar
lançamentos contábeis, ordens de pagamento e de contas a pagar e
receber. Analisar mercado de capitais, contratos bancários e legislações
contábil, tributária, trabalhista e do consumidor. Coletar informações e
elaborar orçamento empresarial. Elaborar relatórios de controle de
custos, de gastos e de despesas gerais. Atuar de acordo com princípios
da educação financeira organizacional e pessoal. Utilizar sistemas
informatizados na execução de atividades financeiras.

Para atuação como Técnico em Finanças, são fundamentais:

Conhecimentos e saberes relacionados à área financeira e bancária, para
atuação em conformidade com as legislações e preceitos da Governança
Corporativa e Compliance, como também com normas de saúde e segurança do
trabalho. Competências socioemocionais e atributos comportamentais
relacionados à comunicação clara e cordial, respeito à diversidade,
trabalho colaborativo, flexibilidade na resolução de conflitos e
protagonismo na análise e solução de problemas.

\textbf{CBOs correspondidos e seus perfis (presentes no CNCT, e presentes na predição da IA)}

\textbf{CBO: 353205  ( Técnico de operações e serviços bancários – câmbio )}

Atende cliente -- pessoa física ou jurídica - para compra e venda de
moeda estrangeira, venda de cartões internacionais e compra de cheques
internacionais. Realiza os serviços bancários relativos a câmbio
comercial e de turismo, conferindo documentação do cliente e da operação
que envolve moeda estrangeira. Enquadra a operação às normas internas e
externas. Calcula custo de operação de câmbio (VET - Valor efetivo
total). Realiza negociações de taxas. Emite contrato de câmbio, carta de
crédito e outros documentos. Registra as operações de câmbio, inclusive
para importação, exportação e remessa financeira, no Sistema Integrado
de Registro de Operações de Câmbio (Sistema Câmbio) - Sistema de
Informação do Banco Central (Sisbacen). Efetua transações de troca de
moedas entre bancos, como câmbio e cobertura ou segurança (hedge).
Realiza controle contábil e faz lançamentos de acertos. Controla
operações de financiamento, alimentando o sistema com informações e
verificando documentação pendente. Notifica clientes em débito e
encaminha processos para cobrança judicial. Executa baixa ou liquidação
de operações. Elabora relatórios sobre posição de carteiras. Participa
da gestão do sistema de normas internas e externas, fazendo sua
atualização e inspecionando sua aplicação. Instrui e orienta
funcionários sobre a aplicação das normas e regras, podendo contatar
órgão governamentais -- tais como Receita Federal e Banco Central do
Brasil -, para solicitação de esclarecimentos ou informações
complementares. Presta atendimento a clientes internos e externos --
inclusive instituições e clientes internacionais -, esclarecendo dúvidas
e dando orientações. Obtém autorização de cliente para operações.
Encaminha reclamações de clientes para o setor competente. Presta
informações a outras áreas da instituição financeira sobre operações
cambiais. Pode treinar funcionários. Recebe malotes, checando os
documentos recebidos. Pode organizar e arquivar -- em meio físico e
eletrônico -- documentos.

\begin{Form}
\textbf{Esta correspondência está adequada?} \\ 
\ChoiceMenu[radio,radiosymbol=\ding{52},bordercolor=0 0.5 0,name=cbovalid_ 353205 ]{}{CBO deve ficar=sim, \hspace{6em} CBO não fica aqui=nao}
\end{Form}

\textbf{CBO: 353210  ( Técnico de operações e serviços bancários - crédito imobiliário )}

Atende clientes interessados em financiamento para aquisição de imóveis
residenciais e comerciais. Executa os serviços bancários relativos ao
crédito imobiliário, conferindo documentação do cliente, do imóvel a ser
financiado e demais documentos que compõem essa operação. Enquadra a
operação às normas internas e externas. Calcula custo da operação de
financiamento (CET -- Custo efetivo total). Realiza negociação de taxas.
Controla a utilização do Fundo de Garantia por Tempo de Serviço (FGTS).
Calcula saldo devedor e juro. Emite carta de crédito, contrato e outros
documentos. Efetua a operação de financiamento imobiliário e solicita
liberação de recursos. Registra a operação no Sistema de Informações de
Créditos do Banco Central (SCR). Acompanha registro de alienação
fiduciária ou hipoteca de imóveis em cartório de registro. Realiza
controle contábil, concilia valores e procede lançamento de acerto,
quando necessário. Verifica a disponibilidade de recursos obrigatórios
junto a órgãos controladores. Controla a operação de financiamento,
alimentando o sistema com informações e verificando documentação
pendente. Acompanha inadimplência, notifica clientes em débito, repactua
dívidas e encaminha processos para cobrança judicial. Controla processos
enviados para cobrança judicial. Executa baixa ou liquidação de
operações. Elabora relatórios sobre posição das carteiras de
financiamento. Participa da gestão do sistema de normas internas e
externas, fazendo sua atualização e inspecionando sua aplicação. Instrui
e orienta funcionários sobre a aplicação das normas e regras, podendo
contatar órgãos governamentais -- tais como Receita Federal e Banco
Central do Brasil -, para solicitação de esclarecimentos ou informações
complementares. Presta atendimento a clientes internos e externos,
esclarecendo dúvidas e dando orientações. Calcula saldos devedores e
juros e pesquisa informações sobre operações e documentação, para
prestar esclarecimentos aos clientes. Obtém autorização de cliente para
operações. Encaminha reclamações de clientes para o setor competente.
Presta informações a outras áreas da instituição financeira sobre
operações de crédito imobiliário. Pode treinar funcionários.

\begin{Form}
\textbf{Esta correspondência está adequada?} \\ 
\ChoiceMenu[radio,radiosymbol=\ding{52},bordercolor=0 0.5 0,name=cbovalid_ 353210 ]{}{CBO deve ficar=sim, \hspace{6em} CBO não fica aqui=nao}
\end{Form}

\textbf{CBO: 353215  ( Técnico de operações e serviços bancários - crédito rural )}

Atende clientes interessados em recursos para custeio da produção,
comercialização de produtos agropecuários e investimentos na produção
rural. Executa serviços bancários relativos ao crédito rural, conferindo
documentação do cliente e da propriedade rural e demais documentos que
compõem a operação. Enquadra a operação às normas internas e externas.
Calcula custo de operações de financiamento (CET -- Custo efetivo
total). Realiza negociações de taxas. Emite contrato, certidão e outros
documentos. Efetua a operação de crédito rural e solicita liberação de
recursos. Registra a operação - inclusive o controle das coordenadas
geodésicas da propriedade - no Sistema de Operações do Crédito Rural e
do Proagro (Sicor) do Banco Central do Brasil. Efetua controle contábil,
concilia valores e procede lançamento de acerto, quando necessário.
Verifica disponibilidade de recursos obrigatórios junto a órgãos
controladores. Controla a operação de financiamento, alimentando o
sistema com informações e verificando documentação pendente. Encaminha
processos para cobrança judicial, controlando-os. Executa baixa ou
liquidação de operação. Elabora relatórios sobre a posição de carteira
de financiamento. Participa da gestão do sistema de normas internas e
externas, fazendo sua atualização e inspecionando sua aplicação. Instrui
e orienta funcionários sobre a aplicação das normas e regras, podendo
contatar órgãos governamentais -- tais como Receita Federal e Banco
Central do Brasil -, para solicitação de esclarecimentos ou informações
complementares. Presta atendimento a clientes internos e externos,
esclarecendo dúvidas e dando orientações. Calcula saldos devedores e
juros e pesquisa operações e documentação, para prestar informações aos
clientes. Encaminha reclamações de clientes para o setor competente.
Presta informações a outras áreas da instituição financeira sobre
operações de crédito rural. Pode treinar funcionários.

\begin{Form}
\textbf{Esta correspondência está adequada?} \\ 
\ChoiceMenu[radio,radiosymbol=\ding{52},bordercolor=0 0.5 0,name=cbovalid_ 353215 ]{}{CBO deve ficar=sim, \hspace{6em} CBO não fica aqui=nao}
\end{Form}

\textbf{CBO: 353220  ( Técnico de operações e serviços bancários - leasing )}

Executa serviços bancários relativos a operações de leasing financeiro e
leasing operacional. Confere documentação do cliente, documentação dos
bens móveis e imóveis arrendados e demais documentos que compõem a
operação de arrendamento mercantil. Enquadra a operação às normas
internas e externas. Calcula custo de operações de financiamento (CET --
Custo efetivo total). Realiza negociação de taxas. Emite contratos,
aditivos e outros documentos. Registra a operação no Sistema de
Informações de Créditos (SCR) do Banco Central do Brasil. Efetua
controle contábil, concilia valores e procede lançamento de acerto,
quando necessário. Controla operação de arrendamento mercantil,
verificando documentação pendente, recebendo e averiguando documentos e
alimentando sistema com informações. Acompanha inadimplência, notifica
clientes em débito, repactua dívidas e encaminha processos para cobrança
judicial. Controla processos enviados para cobrança judicial. Executa
baixa ou liquidação de operação. Elabora relatórios sobre posição de
carteira de arrendamento mercantil. Participa da gestão do sistema de
normas internas e externas, fazendo sua atualização e inspecionando sua
aplicação. Instrui e orienta funcionários sobre a aplicação das normas e
regras, podendo contatar órgãos governamentais -- tais como Receita
Federal e Banco Central do Brasil -, para solicitação de esclarecimentos
ou informações complementares. Presta atendimento a clientes internos e
externos, esclarecendo dúvidas, atualizando dados cadastrais e dando
orientações. Calcula saldos devedores e juros e pesquisa operações e
documentação, para prestar informações aos clientes. Obtém autorização
de cliente para operações. Encaminha reclamações de cliente para o setor
competente. Presta informações a outras áreas da instituição financeira
sobre operações de leasing. Pode treinar funcionários.

\begin{Form}
\textbf{Esta correspondência está adequada?} \\ 
\ChoiceMenu[radio,radiosymbol=\ding{52},bordercolor=0 0.5 0,name=cbovalid_ 353220 ]{}{CBO deve ficar=sim, \hspace{6em} CBO não fica aqui=nao}
\end{Form}

\newpage
\section{ 04008 -- Técnico em Logística }

\textbf{Perfil do Curso (CNCT)}

O Técnico em Logística será habilitado para:

Auxiliar no planejamento, operacionalização e controle da cadeia
produtiva e seu fluxo logístico. Executar procedimentos relacionados a
suprimentos, produção, recebimento, armazenagem e distribuição de
produtos, fazendo uso das tecnologias de informação e comunicação.
Identificar agentes da cadeia de suprimentos. Elaborar relatórios
operacionais para tomada de decisões.

Para atuação como Técnico em Logística, são fundamentais:

Conhecimentos e saberes relacionados à área operacional, de produção e
de prestação de serviços das organizações, de modo a atuar em
conformidade com as legislações e diretrizes de conduta, como também com
as normas de saúde e segurança do trabalho. Atuação de forma proativa na
resolução de situações-problema do mundo do trabalho, desenvolvendo
competências socioemocionais e atributos comportamentais relacionados à
sustentabilidade e ao trabalho colaborativo.

\textbf{CBOs correspondidos e seus perfis (presentes no CNCT, e presentes na predição da IA)}

\textbf{CBO: 391125  ( Técnico de planejamento de produção )}

Planeja a produção, a partir da análise da demanda e da capacidade de
produção. Realiza a programação da produção, analisando ordens de
produção, determinando prioridades, definindo cronograma e estabelecendo
a sequência de produção. Estabelece parâmetros para controle dos
processos de produção. Controla o desenvolvimento dos planos de
produção, com base em parâmetros de controle e monitoramento do fluxo
produtivo. Identifica desvios e define ações corretivas. Propõe
melhorias no processo de produção. Define estoque de segurança de
suprimentos necessários para atender às demandas de produção.
Providencia entrega e distribuição de suprimentos para agilizar o fluxo
de materiais e cumprir os cronogramas de produção. Pode revisar os
cronogramas de produção por falta de trabalhadores, atraso de entrega de
material ou outros problemas. Trata as informações do ciclo
planejamento-produção, elaborando gráficos e relatórios de controle.
Disponibiliza e atualiza informações. Interage com profissionais
envolvidos nas diversas etapas de produção, com clientes e com
fornecedores, para coordenar as atividades de produção, para resolver
reclamações ou para eliminar atrasos. Auxilia na definição das
instruções de trabalho.

\begin{Form}
\textbf{Esta correspondência está adequada?} \\ 
\ChoiceMenu[radio,radiosymbol=\ding{52},bordercolor=0 0.5 0,name=cbovalid_ 391125 ]{}{CBO deve ficar=sim, \hspace{6em} CBO não fica aqui=nao}
\end{Form}

\newpage
\section{ 04009 -- Técnico em Marketing }

\textbf{Perfil do Curso (CNCT)}

O Técnico em Marketing será habilitado para:

Projetar e implementar planos de marketing. Realizar análises de vendas,
preços e produtos. Desenvolver projetos de comunicação, ?delização de
clientes e relação com fornecedores ou outras entidades. Desenvolver,
implementar e gerenciar estratégias de marketing digital.
Operacionalizar apresentação dos serviços e produtos no ponto de venda.
Elaborar estudos de mercado.

Para atuação como Técnico em Marketing, são fundamentais:

Conhecimentos e saberes relacionados à área comercial e de negócios das
organizações, de modo a atuar em conformidade com as legislações e
diretrizes de conduta, como também com as normas de saúde e segurança do
trabalho, demonstrando visão empreendedora. Competências socioemocionais
e atributos comportamentais relacionados à comunicação clara e cordial,
respeito à diversidade, trabalho colaborativo e protagonismo na análise
e resolução de problemas voltados ao mundo do trabalho.

\textbf{CBOs correspondidos e seus perfis (presentes no CNCT, e presentes na predição da IA)}

\textbf{CBO: 354140  ( Técnico em atendimento e vendas )}

Analisa o comportamento do mercado consumidor, participando do
desenvolvimento de estudos mercadológicos. Estuda o cliente potencial da
empresa, analisando seu perfil e suas expectativas, visando ao
estabelecimento de possíveis estratégias de fidelização. Avalia qual
estratégia de marketing é mais adequada para cada tipo de cliente.
Analisa vantagens e desvantagens no uso de comércio eletrônico
(e-commerce), dependendo do tipo de cliente. Pesquisa ações dos
concorrentes, para analisar o que é oferecido aos clientes. Providencia
o cadastramento de clientes potenciais e efetivos, atualizando os dados
sempre que necessário. Pode usar cadastro de clientes centralizado da
plataforma de Gestão de Relacionamento com o Cliente (Customer
Relationship Management - CRM). Contata clientes, usando estratégias
digitais (e-mail, redes sociais ou outras) e/ou realizando visitas.
Comunica-se com os clientes, identificando suas necessidades e
expectativas e analisando seu potencial de compras. Divulga produtos e
serviços aos clientes, apresentando características e benefícios do
produto ou serviço, esclarecendo dados técnicos, informando preços e
detalhando formas de pagamento. Pode fornecer amostras e convidar
clientes para visitar a empresa. Concretiza vendas presencialmente, por
telefone ou por meios digitais, de acordo com as características dos
clientes. Acompanha o cliente após a venda -- inclusive por meio de
visita -- para sanar suas dúvidas, providenciar apoio técnico e
verificar suas opiniões e suas sugestões. Analisa reclamações e
insatisfações, para propor melhorias no processo de vendas. Elabora
relatórios e estatísticas sobre atendimento ao cliente e sobre o
processo de vendas. Presta assistência à equipe de vendedores de
produtos e serviços, com relação aos procedimentos adequados a cada tipo
de cliente nas etapas de pré-venda, venda e pós-venda. Participa da
elaboração de planos e estratégias de vendas, sugerindo melhorias no
atendimento ao cliente e no processo de vendas.

\begin{Form}
\textbf{Esta correspondência está adequada?} \\ 
\ChoiceMenu[radio,radiosymbol=\ding{52},bordercolor=0 0.5 0,name=cbovalid_ 354140 ]{}{CBO deve ficar=sim, \hspace{6em} CBO não fica aqui=nao}
\end{Form}

\newpage
\section{ 04011 -- Técnico em Qualidade }

\textbf{Perfil do Curso (CNCT)}

O Técnico em Qualidade será habilitado para:

Elaborar manuais, procedimentos, diagnósticos e relatórios de processos
de qualidade das organizações. Registrar o controle da qualidade.
Executar auditorias internas da qualidade. Acompanhar auditorias
externas. Divulgar procedimentos de qualidade. Propor ações de
informação e formação específica. Identificar inconformidades em
serviços, produtos e processos e suas possíveis causas. Propor ações
corretivas e preventivas. Interpretar conjunto de mecanismos e
procedimentos de integridade, controle e auditoria. Executar atividades
voltadas à prevenção, à detecção e resolução de desvios, a fraudes, a
irregularidades e atos ilícitos praticados nas organizações.

Para atuação como Técnico em Qualidade, são fundamentais:

Conhecimentos e saberes relacionados ao funcionamento da área de
qualidade e compliance, de modo a atuar em conformidade com as
legislações e diretrizes de conduta, como também com as normas de saúde
e segurança do trabalho. Atuação de forma proativa em atividades
técnico-administrativas de gestão e controle da qualidade, pautadas nos
preceitos da governança corporativa, demonstrando comprometimento com os
padrões e modelos institucionais.

\textbf{CBOs correspondidos e seus perfis (presentes no CNCT, e presentes na predição da IA)}

\textbf{CBO: 391205  ( Inspetor de qualidade )}

Planeja e desenvolve a inspeção de insumos e produtos - peças, conjuntos
e material em processo de produção -, definindo amostras e aplicando
procedimentos e métodos de inspeção de acordo com sistemas, normas e
instruções específicas. Define e opera equipamentos laboratoriais,
instrumentos e ferramentas para realização de ensaios. Pode fazer coleta
e envio de amostras para análises. Aplica ferramentas da qualidade e
metodologia estatística para avaliação da conformidade, de acordo com o
objeto da inspeção. Verifica conformidade de processos, controlando as
condições do ambiente em que se desenvolve o processo e observando
condições de uso de equipamentos, instrumentos e ferramentas. Analisa
relatórios de controle de processos para constatar o cumprimento de
normas e procedimentos e informar setores responsáveis sobre processos,
produtos e serviços não-conformes. Na inspeção de recebimento, na
inspeção realizada durante o processo de produção e na inspeção de
produtos acabados, examina documentos, verifica local de recebimento e
inspeção, avalia condições de descarregamento e verifica higiene e
limpeza do local de armazenamento de insumos e produtos. Pode executar
ações corretivas em função de reclamações do mercado. Libera produtos e
serviços inspecionados, verificando especificações e padrões e
conferindo condições de transporte e higiene.

\begin{Form}
\textbf{Esta correspondência está adequada?} \\ 
\ChoiceMenu[radio,radiosymbol=\ding{52},bordercolor=0 0.5 0,name=cbovalid_ 391205 ]{}{CBO deve ficar=sim, \hspace{6em} CBO não fica aqui=nao}
\end{Form}

\newpage
\section{ 04012 -- Técnico em Recursos Humanos }

\textbf{Perfil do Curso (CNCT)}

O Técnico em Recursos Humanos será habilitado para:

Organizar rotina diária dos processos de gestão de pessoas inerentes à
relação de emprego/trabalho existente entre empresa e empregado, bem
como documentos da área de recursos humanos. Processar cálculos de folha
de pagamento. Registrar informações governamentais, de fiscalizações, de
processos trabalhistas e de auditoria interna em recursos humanos.
Organizar e realizar ações de recrutamento e seleção. Realizar
atividades diárias para desenvolvimento de pessoas e retenção de
talentos. Organizar rotinas relativas às políticas de remuneração e
cargos. Realizar atividades relativas à concessão de benefícios.
Acompanhar e organizar processos administrativos de higiene e segurança
do trabalho. Organizar e realizar ações de inclusão de Pessoas com
Deficiência (PCDs) no ambiente de trabalho.

Para atuação como Técnico em Recursos Humanos, são fundamentais:

Conhecimentos e saberes relacionados à área de pessoal, para atuação em
conformidade com as legislações e diretrizes de conduta, como também com
as normas de saúde e segurança do trabalho. Competências socioemocionais
e atributos comportamentais relacionados à comunicação clara e cordial,
respeito à diversidade, atenção à sustentabilidade, trabalho
colaborativo, flexibilidade na resolução de conflitos e protagonismo na
análise e solução de problemas.

\textbf{CBOs correspondidos e seus perfis (presentes no CNCT, e presentes na predição da IA)}

\textbf{CBO: 351315  ( Agente de recrutamento e seleção )}

A partir de recebimento de solicitação de preenchimento de vaga,
estabelece etapas, elabora cronograma e organiza recursos necessários
para os processos de recrutamento e seleção. Elabora proposta de edital,
em parceria com o gestor da área que solicitou o preenchimento de vaga.
Após aprovação do edital, providencia sua divulgação nos canais de
comunicação. Verifica a possibilidade de participação de profissional do
quadro no processo seletivo, realizando triagem preliminar de currículos
e analisando históricos de desempenho, conforme requisitos e critérios
preestabelecidos no edital da vaga. Convoca candidatos e os orienta a
respeito das etapas do processo de recrutamento e seleção, de acordo com
procedimentos da área de recursos humanos. Verifica documentos e dados
de candidatos para processos de seleção. Explica as políticas de pessoal
e informa o salário e os benefícios oferecidos aos candidatos a emprego.
Executa ou acompanha cada fase do processo seletivo, podendo
providenciar impressão de materiais, fazer correção de provas, realizar
entrevistas, acompanhar dinâmicas de grupo ou aplicar outras estratégias
de seleção. Confere os documentos dos candidatos selecionados e os
encaminha para admissão. Apresenta o novo contratado às equipes de
trabalho, realizando sua ambientação no trabalho, de acordo com
procedimentos da área de recursos humanos. Reúne os dados dos processos
de recrutamento, seleção e ambientação e os inclui em sistemas de
informações gerenciais de pessoal. Organiza, arquiva e encaminha
documentos. Pode propor melhorias no processo de recrutamento, seleção e
ambientação, identificando oportunidades de inovação e pesquisando boas
práticas no mercado.

\begin{Form}
\textbf{Esta correspondência está adequada?} \\ 
\ChoiceMenu[radio,radiosymbol=\ding{52},bordercolor=0 0.5 0,name=cbovalid_ 351315 ]{}{CBO deve ficar=sim, \hspace{6em} CBO não fica aqui=nao}
\end{Form}

\textbf{CBO: 411010  ( Assistente administrativo )}

Executa atividades de apoio aos setores administrativos, providenciando
compra e distribuindo recursos materiais, providenciando reparos de
equipamentos, intermediando contatos, solicitando recursos para viagens,
realizando prestação de contas, controlando execução de serviços gerais
- limpeza, transporte e vigilância - e monitorando expedição e
recebimentos de malotes. Opera sistemas de correio eletrônico e coordena
o fluxo de informações, internamente ou com outras organizações. Mantém
agenda e calendário de eventos, preparando materiais como folhetos e
convites. Executa atividades de apoio ao setor de recursos humanos,
auxiliando no controle de pessoal - afastamentos, férias e horas extras
-, realizando procedimentos de recrutamento e seleção e prestando
suporte administrativo à área de treinamento. Executa atividades de
apoio aos setores contábil e financeiro, coletando informações
financeiras e índices econômicos, digitando notas de lançamentos
contábeis e conferindo notas fiscais, faturas de pagamentos e boletos.
Mantém sistemas de arquivos - físico ou eletrônico -, registrando e
atualizando informações. Controla documentação, triando, distribuindo,
classificando, conferindo dados e datas, registrando entrada e saída e
identificando irregularidades nos documentos. Presta apoio na elaboração
de documentos, digitando textos e planilhas, preenchendo formulários e
cadastros, redigindo atas, prestando apoio operacional para elaboração
de manuais técnicos e elaborando minutas de correspondência. Elabora
organogramas, fluxogramas e cronogramas. Insere e atualiza informações
nos bancos de dados. Atende os clientes e os fornecedores, identificando
seus perfis, fornecendo informações sobre produtos e serviços,
esclarecendo dúvidas e registrando reclamações.

\begin{Form}
\textbf{Esta correspondência está adequada?} \\ 
\ChoiceMenu[radio,radiosymbol=\ding{52},bordercolor=0 0.5 0,name=cbovalid_ 411010 ]{}{CBO deve ficar=sim, \hspace{6em} CBO não fica aqui=nao}
\end{Form}

\newpage
\section{ 04013 -- Técnico em Secretariado }

\textbf{Perfil do Curso (CNCT)}

O Técnico em Secretariado será habilitado para:

Executar atividades voltadas ao planejamento organizacional e
operacional. Prestar assessoramento a gestores(as) de organizações de
diferentes portes e segmentos econômicos. Utilizar técnicas secretariais
e ferramentas tecnológicas em atividades relativas ao fluxo processual
de gestão, organização e registro administrativo, de informação e de
relacionamento com clientes internos e externos. Apoiar atividades de
gestão financeira, orçamentos, pagamentos e prestação de contas.

Para atuação como Técnico em Secretariado, são fundamentais:

Conhecimentos e saberes relacionados ao funcionamento das organizações,
de modo a atuar em conformidade com as legislações e diretrizes de
conduta, como também com as normas de saúde e segurança do trabalho.
Atuação de forma proativa em atividades de mediação, de resolução de
conflitos, de situações-problema e trabalho em equipe, com comunicação
clara e cordial e respeito à diversidade.

\textbf{CBOs correspondidos e seus perfis (presentes no CNCT, e presentes na predição da IA)}

\textbf{CBO: 351505  ( Técnico em secretariado )}

Realiza atividades secretariais de apoio ao gestor, auxiliando-o nas
suas atividades e no cumprimento de seus compromissos. Executa serviços
de apoio, reservando passagens aéreas e hotéis em viagens nacionais e
internacionais, prestando suporte administrativo e logístico em eventos
corporativos, elaborando pautas e atas de reuniões e realizando outras
atividades de acordo com a agenda e as prioridades definidas. Estabelece
os canais de comunicação do gestor com interlocutores. Recepciona e
atende clientes internos e externos, prestando-lhes informações,
encaminhando-os para o setor responsável pelo atendimento a suas
solicitações e verificando sua satisfação. Pode cadastrar e atualizar as
informações de clientes e fornecedores. Elabora cartas, memorandos,
ofícios, relatórios e outros documentos. Confere, revisa e corrige
textos, promovendo ajustes embasados em normas técnicas. Pode redigir
textos em idioma estrangeiro. Organiza, controla a circulação e arquiva
documentos. Pesquisa, localiza e recupera documentos, física ou
digitalmente. Pesquisa dados em fontes diversas, inclusive em páginas da
Internet. Organiza e sintetiza informações, podendo armazená-las em meio
físico ou eletrônico. Executa atividades financeiras de apoio aos
processos administrativos, orçando produtos e serviços, conferindo e
encaminhando pagamento, prestando contas e finalizando contratos e
aquisições de acordo com os processos organizacionais. Administra e
organiza recursos para execução dos trabalhos, requisitando material de
escritório, controlando estoque, providenciando manutenção de
equipamentos e aparelhos. Pode supervisionar equipe de trabalho,
distribuindo tarefas, acompanhando o andamento dos serviços e avaliando
o desempenho de cada profissional.

\begin{Form}
\textbf{Esta correspondência está adequada?} \\ 
\ChoiceMenu[radio,radiosymbol=\ding{52},bordercolor=0 0.5 0,name=cbovalid_ 351505 ]{}{CBO deve ficar=sim, \hspace{6em} CBO não fica aqui=nao}
\end{Form}

\newpage
\section{ 04014 -- Técnico em Seguros }

\textbf{Perfil do Curso (CNCT)}

O Técnico em Seguros será habilitado para:

Analisar proposta de seguro. Avaliar risco e outras condições
comerciais. Elaborar proposta de apólices e contratos para as diversas
modalidades de seguro. Elaborar pareceres técnicos e relatórios.
Analisar e conferir propostas de contratos de seguros e prestar outros
serviços de apoio e suporte técnico. Controlar contas correntes
relativas a prêmios e sinistros. Organizar fatos contábeis.
Operacionalizar cálculos de prêmios e outros procedimentos.

Para atuação como Técnico em Seguros, são fundamentais:

Conhecimentos e saberes relacionados aos processos de formalização e
contratação de seguros, de modo a atuar em conformidade com as
legislações e diretrizes de órgãos reguladores, como também com as
normas de saúde e segurança do trabalho. Exercício profissional pautado
no respeito à diversidade, por intermédio de comunicação clara e
cordial, demonstrando atenção à sustentabilidade, trabalho colaborativo,
flexibilidade na resolução de conflitos e protagonismo na análise e
solução de problemas.

\textbf{CBOs correspondidos e seus perfis (presentes no CNCT, e presentes na predição da IA)}

\textbf{CBO: 351710  ( Analista de sinistros )}

Realiza abertura de processos de sinistros, confrontando a reivindicação
do segurado com os termos da sua apólice. Estuda evidências da
ocorrência do sinistro, avaliando causas e efeitos e analisando
informações sobre natureza e extensão das avarias. Confere a veracidade
dessas informações, confrontando-as com Boletim de Ocorrência (B.O.) e
laudos apresentados. Pode solicitar informações complementares, vistoria
ao bem sinistrado ou, mesmo, sindicância sobre o sinistro, para evitar
ilícitos. Examina as reivindicações investigadas pelos avaliadores de
seguros, averiguando ainda as reivindicações questionáveis, para avaliar
se o pagamento deve ser autorizado. Encerra o processo, providenciando
alienação do salvado de sinistro, calculando o valor de indenização e,
se for o caso, encaminhando o processo para pagamento ao segurado.
Presta informações e, sempre que possível, atende a reivindicações dos
segurados, agilizando o envio de documentações para reduzir o tempo para
conclusão do processo.

\begin{Form}
\textbf{Esta correspondência está adequada?} \\ 
\ChoiceMenu[radio,radiosymbol=\ding{52},bordercolor=0 0.5 0,name=cbovalid_ 351710 ]{}{CBO deve ficar=sim, \hspace{6em} CBO não fica aqui=nao}
\end{Form}

\textbf{CBO: 351735  ( Técnico de resseguros )}

Analisa termos, condições e preços propostos pela seguradora,
subsidiando a decisão dos subscritores sobre aceitação de um novo
contrato de resseguro. Pode negociar os termos com a companhia de
seguro. Prepara documentação sobre sinistros pendentes e pagos, calcula
prêmios por excesso de perdas, registra crédito e débitos, entre outros
processos de controle dos contratos existentes. Analisa os riscos de uma
carteira e o potencial passivo para o ressegurador de indenizar uma
seguradora. A partir das informações disponíveis, pode calcular a
probabilidade estatística de uma reivindicação ser feita, para que a
resseguradora analise tendências e posicione-se em relação à natureza do
risco. Pode participar da criação de programas de resseguros sob medida,
que protejam os interesses da resseguradora e também contribuam para que
os clientes atinjam seus objetivos. Apoia corretores de resseguros,
subscritores, advogados, gerentes de contas e outros membros da equipe.

\begin{Form}
\textbf{Esta correspondência está adequada?} \\ 
\ChoiceMenu[radio,radiosymbol=\ding{52},bordercolor=0 0.5 0,name=cbovalid_ 351735 ]{}{CBO deve ficar=sim, \hspace{6em} CBO não fica aqui=nao}
\end{Form}

\textbf{CBO: 351740  ( Técnico de seguros )}

Analisa proposta de seguro, comparando-a com as condições estabelecidas
pela companhia, estudando relatórios de inspeções, laudos médicos e
outras informações do cliente. Avalia risco e outras condições, para
subsidiar a decisão de aceitação ou não da solicitação. Elabora proposta
de apólices e contratos para as diversas modalidades de seguro, fixando
as tarifas, prêmios, valor de resgate e outras condições, para atender
às solicitações de empresas ou particulares. Executa atividades
administrativas e financeiras para apoiar às diferentes áreas da
companhia, tais como enviar avisos de renovação de seguros, juntar
documentação para licitações, controlar contas correntes relativas a
prêmios e sinistros e organizar fatos contábeis. Elabora parecer técnico
e relatórios, pesquisando legislação, analisando resultados de carteira,
avaliando a adesão dos clientes aos vários produtos e serviços
oferecidos. Apura rentabilidade e estuda outras informações
sistematizadas disponíveis em sistemas informatizados. Presta
atendimento a clientes, informando sobre a oferta de produtos,
respondendo ou providenciando resposta às reclamações e encaminhando
solicitações às diversas áreas da companhia de seguro.

\begin{Form}
\textbf{Esta correspondência está adequada?} \\ 
\ChoiceMenu[radio,radiosymbol=\ding{52},bordercolor=0 0.5 0,name=cbovalid_ 351740 ]{}{CBO deve ficar=sim, \hspace{6em} CBO não fica aqui=nao}
\end{Form}

\newpage
\section{ 04015 -- Técnico em Serviços Jurídicos }

\textbf{Perfil do Curso (CNCT)}

O Técnico em Serviços Jurídicos será habilitado para:

Executar atividades administrativas de planejamento, organização,
direção e controle em rotinas de escritórios de advocacia e demais
organizações que dispõem de departamento jurídico. Prestar suporte e
apoio técnico-administrativo a profissionais da área jurídica.
Acompanhar, gerenciar e arquivar documentos e processos de natureza
jurídica. Prestar atendimento receptivo ao público.

Para atuação como Técnico em Serviços Jurídicos, são fundamentais:

Conhecimentos e saberes relacionados ao direito, de modo a atuar em
conformidade com as legislações e diretrizes de órgãos reguladores, como
também com as normas de saúde e segurança do trabalho. Atuação de forma
proativa em atividades de mediação, de resolução de conflitos, de
situações-problema e trabalho em equipe, com comunicação clara e cordial
e respeito à diversidade.

\textbf{CBOs correspondidos e seus perfis (presentes no CNCT, e presentes na predição da IA)}

\textbf{CBO: 351405  ( Escrevente )}

Prepara e organiza o trabalho, estabelecendo a ordem de execução das
atividades -- com prioridade para a elaboração de documentos legais de
urgência - e controlando a chegada de novos documentos. Expede
documentos em cartórios, tais como cartas precatórias e rogatórias;
ofícios; certidões de nascimento, casamento e óbito; carta de
arrematação e adjudicação; certidões negativas e positivas de protestos;
guias de sepultamento; alvarás; cartas de sentença; entre outros.
Registra protestos de títulos, sustação de protestos, depósitos
judiciais, penhoras e outros documentos. Cumpre determinações legais e
judiciais, lavrando e executando atos, dando publicidade a atos,
remetendo autos aos tribunais e acompanhando fases de processos.
Coadjuva nas audiências e nas sessões de mediação e arbitragem,
organizando a sala, lendo sentenças e outras peças dos autos, lavrando
atas de audiência e colhendo assinaturas das partes e das testemunhas.
Pode elaborar relatórios estatísticos. Organiza arquivos, classificando
documentos e conservando cópias de segurança em diversas mídias.
Digitaliza documentos. Utiliza sistemas, plataformas informatizadas,
internet e redes internas de comunicação, para execução dos trabalhos.
Realiza pesquisa em plataformas, repositórios e bancos de dados de
processos judiciais. Presta atendimento ao público, informando andamento
de processos, autenticando documentos, reconhecendo firmas, entre outros
serviços. Conserva o local de trabalho limpo e organizado. Verifica a
disponibilidade de materiais de escritório. Mantém equipamentos,
instrumentos e acessórios de trabalho limpos, acondicionados e em plenas
condições de uso e funcionamento. Requisita manutenção de equipamentos,
quando necessário. Realiza a destinação de resíduos, fazendo
reaproveitamento e providenciando descarte de acordo com as normas
ambientais.

\begin{Form}
\textbf{Esta correspondência está adequada?} \\ 
\ChoiceMenu[radio,radiosymbol=\ding{52},bordercolor=0 0.5 0,name=cbovalid_ 351405 ]{}{CBO deve ficar=sim, \hspace{6em} CBO não fica aqui=nao}
\end{Form}

\textbf{CBO: 351430  ( Auxiliar de serviços jurídicos )}

Planeja as atividades, priorizando ações e providenciando recursos
materiais para execução das tarefas. Recebe documentos de processos
jurídicos, registrando a entrada e encaminhando ao profissional ou à
equipe responsável pelo processo. Controla prazos estabelecidos para
tramitação dos documentos. Expede - preferencialmente por meio
eletrônico - documentos, tais como ofícios, certidões, alvarás, entre
outros. Remete autos aos tribunais. Providencia a entrega de
correspondência legal para testemunhas, oficiais do tribunal e outros
destinatários. Registra a emissão de petições iniciais, atas de
audiências e outros documentos. Organiza arquivos físicos e digitais,
classificando, reproduzindo e digitalizando documentos e processos.
Acompanha a consulta dos arquivos, auxiliando os consulentes, quando
necessário. Presta apoio às pesquisas em plataformas, repositórios e
bancos de dados, para identificação de leis, decisões judiciais em
processos e outras informações. Organiza a realização de reuniões e de
audiências, encaminhando pauta aos participantes, testando
apresentações, reproduzindo documentos, organizando o local e tomando
outras providências solicitadas. Pode preparar relatórios estatísticos,
a partir de dados fornecidos pelas equipes técnicas. Presta atendimento
ao público, utilizando softwares para consultar e posteriormente
informar andamento de processos; identificando o profissional ou a
equipe técnica para encaminhar questões recebidas; e respondendo sobre
recebimento e expedição de documentos. Conserva o local de trabalho
limpo e organizado. Mantém instrumentos e equipamentos de trabalho
limpos e em plenas condições de uso e funcionamento. Monitora o
funcionamento do sistema de informática. Requisita a aquisição de
materiais para reposição, quando necessário. Solicita serviços de
manutenção para reparos em equipamentos. Realiza a destinação de
resíduos, providenciando descarte de acordo com as normas ambientais.

\begin{Form}
\textbf{Esta correspondência está adequada?} \\ 
\ChoiceMenu[radio,radiosymbol=\ding{52},bordercolor=0 0.5 0,name=cbovalid_ 351430 ]{}{CBO deve ficar=sim, \hspace{6em} CBO não fica aqui=nao}
\end{Form}

\newpage
\section{ 04017 -- Técnico em Transações Imobiliárias }

\textbf{Perfil do Curso (CNCT)}

O Técnico em Transações Imobiliárias será habilitado para:

Executar atividades de intermediação na compra, venda, permuta e locação
de imóveis, sejam terrenos ou edificações. Realizar captação, vistoria e
demonstração de imóveis. Prestar assessoria na identificação de
oportunidades de negócios, no processo de transferências, estruturações
e registros imobiliários. Orientar quanto ao financiamento imobiliário.
Avaliar imóveis para determinar valor de mercado.

Para atuação como Técnico em Transações Imobiliárias, são fundamentais:

Conhecimentos e saberes relacionados ao mercado imobiliário, de modo a
atuar em conformidade com a legislação profissional e do setor, com
diretrizes de conduta e com normas de saúde e segurança do trabalho.
Atuação de forma proativa, comunicando-se de forma clara e cordial,
demonstrando desinibição e comprometimento com necessidades, desejos e
percepção da realidade social de clientes, além de respeito à
diversidade e à sustentabilidade.

\textbf{CBOs correspondidos e seus perfis (presentes no CNCT, e presentes na predição da IA)}

\textbf{CBO: 354410  ( Avaliador de imóveis )}

Recebe solicitação de avaliação de bem imóvel, levantando a
identificação do solicitante. Verifica o objetivo da avaliação, podendo
destinar-se a provar o valor em processos judiciais e extrajudiciais -
tais como divórcios, inventários, indenização em casos de
desapropriação, entre outros -; estabelecer valor para aluguel e venda
do imóvel; ou outra finalidade similar. Seleciona método de avaliação de
imóveis -- como o comparativo (mais usado), o evolutivo ou o involutivo
-- e os procedimentos a serem utilizados no trabalho. Vistoria o imóvel
para fazer sua caracterização, registrando tipo -- como casa,
apartamento ou outro - e tamanho (metragem). Avalia as condições
estruturais internas de casa, apartamento ou imóvel similar, verificando
suas dependências e observando itens - como rede elétrica, rede
hidráulica, pintura, forro, telhado, vidros e acabamento -, além do
estado geral de conservação. Assinala ressalvas, como necessidade de
reforma. Avalia as condições externas do imóvel, verificando itens da
casa -- como jardins, instalação elétrica externa, entre outros -- ou do
apartamento, como estado de conservação de fachada do prédio, portaria,
elevador, área de lazer e vagas em estacionamento. Fotografa partes de
interesse do imóvel. Avalia a localização e a posição geográfica do
imóvel. Analisa o entorno, para verificar a proximidade de oferta de
serviços essenciais (como escolas e hospitais), estabelecimentos
comerciais (farmácia, supermercados, entre outros), acesso a meios de
transporte público, espaços com natureza e lazer. Coleta informações com
a vizinhança sobre segurança - verificando histórico de assaltos ou
outros episódios de violência - e sobre poluição sonora no local.
Considera tendências ou mudanças iminentes no local, que possam
influenciar a avaliação dos valores futuros do imóvel. Pode filmar áreas
do entorno. Analisa a documentação do imóvel, verificando se não há
fator negativo, como gravame de hipoteca ou gravame de área de
preservação. Recolhe cópias -- físicas ou digitais - de plantas baixas;
de comprovantes de pagamentos de impostos; de escritura; e de outros
documentos, para anexá-las ao laudo. Pesquisa o mercado imobiliário da
região, coletando dados sobre aluguéis e sobre preços alcançados em
vendas de imóveis semelhantes ou comparáveis ao bem em avaliação. Faz
tratamento e análise dos dados levantados, de acordo com a metodologia
adotada de avaliação. Elabora o laudo da avaliação, apresentando os
resultados da análise realizada e atribuindo valor ao imóvel, em
determinada data de referência e para uma finalidade predefinida.
Organiza e arquiva cópia de documentos - físicos e digitais --
relacionados ao processo de avaliação.

\begin{Form}
\textbf{Esta correspondência está adequada?} \\ 
\ChoiceMenu[radio,radiosymbol=\ding{52},bordercolor=0 0.5 0,name=cbovalid_ 354410 ]{}{CBO deve ficar=sim, \hspace{6em} CBO não fica aqui=nao}
\end{Form}

\textbf{CBO: 354605  ( Corretor de imóveis )}

Planeja as atividades de intermediação de compra, venda, permuta e
locação de imóveis -- edificações e terrenos urbanos ou rurais -,
identificando o perfil de potenciais clientes e verificando suas
preferências. Prepara pesquisa do mercado imobiliário, delimitando
região de interesse, distinguindo zoneamento e obtendo dados regionais.
Efetua pesquisa direcionada do mercado, visitando a região, consultando
fontes de informações e verificando a disponibilidade de imóveis. Visita
imóveis, verificando suas condições físicas, interpretando seus projetos
arquitetônicos e examinando sua documentação. Define valores de
comercialização. Combina formas de trabalho com os proprietários.
Elabora estratégias de comercialização, fazendo adequação de valores e
condições de pagamento ao mercado e preparando processo de divulgação.
Solicita documento com autorização do proprietário para realizar
divulgação em diversos meios, como ``site'' da internet, fôlder,
revista, entre outros. Cria material de divulgação adequado a cada meio
de comunicação. Pode fotografar diferentes dependências dos imóveis,
para uso na divulgação. Pode usar dados da pesquisa realizada na região,
para apresentar o entorno do imóvel. Entrevista potenciais clientes,
identificando seu perfil e necessidades, para intermediação de negócios
imobiliários. Seleciona imóveis de acordo com a necessidade de cada
cliente. Oferece alternativas de negócios e esclarece dúvidas dos
clientes. Informa seus honorários, incidentes nas transações
imobiliárias. Vistoria os imóveis, inventariando mobiliário,
equipamentos e acessórios. Faz a intermediação de negócios, mostrando os
imóveis a clientes e participando na negociação de valores entre as
partes. Esclarece dúvidas dos clientes, como as restrições do zoneamento
da região. Formaliza a proposta da transação, para conclusão de
negócios. Solicita e confere documentação das partes, examinando
documentos e, quando necessário, requerendo complementação. Apresenta
contratos e aceites das partes, para confirmação das transações. Presta
assessoramento aos clientes após transação, acompanhando a entrega do
imóvel. Solicita indicações de potenciais clientes, para oferecer novos
negócios. Opera sistemas informatizados, inserindo os resultados da
pesquisa do mercado imobiliário e cadastrando imóveis visitados e
potenciais clientes. Atualiza o cadastro de clientes com as informações
das transações imobiliárias realizadas. Faz a organização, a
classificação e o arquivo de documentos - físicos e digitais --
relacionados às transações imobiliárias.

\begin{Form}
\textbf{Esta correspondência está adequada?} \\ 
\ChoiceMenu[radio,radiosymbol=\ding{52},bordercolor=0 0.5 0,name=cbovalid_ 354605 ]{}{CBO deve ficar=sim, \hspace{6em} CBO não fica aqui=nao}
\end{Form}

\newpage
\section{ 04018 -- Técnico em Vendas }

\textbf{Perfil do Curso (CNCT)}

O Técnico em Vendas será habilitado para:

Identificar produtos e serviços da empresa e canais de venda adequados
às respectivas especificidades. Caracterizar os perfis de clientes.
Coletar informações sobre a concorrência e o mercado em geral. Planejar
e promover a venda de produtos e serviços. Organizar o ambiente de
venda. Fidelizar clientes promovendo serviços de apoio e atendimento
pós-venda. Organizar e gerenciar arquivos com informações de clientes.
Realizar prospecção de novos clientes.

Para atuação como Técnico em Vendas, são fundamentais:

Conhecimentos e saberes relacionados ao funcionamento da área comercial,
de modo a atuar em conformidade com as legislações e diretrizes de
conduta, como também com as normas de saúde e segurança do trabalho.
Atuação de forma proativa e com visão empreendedora em atividades de
venda de produtos e serviços, demonstrando desinibição e comprometimento
com necessidades e desejos de clientes, comunicando-se de forma clara e
cordial, com respeito à diversidade.

\textbf{CBOs correspondidos e seus perfis (presentes no CNCT, e presentes na predição da IA)}

\textbf{CBO: 354135  ( Técnico de vendas )}

Avalia o comportamento do mercado consumidor, participando do
desenvolvimento de estudos mercadológicos. Estuda os produtos e os
serviços da empresa, podendo contribuir no desenvolvimento de inovações
para atender às expectativas do mercado consumidor. Pesquisa as ações
dos concorrentes. Caracteriza os clientes potenciais, analisando suas
demandas e necessidades. Prepara os ambientes de venda que poderão ser
usados para dispor produtos ou serviços ao mercado consumidor. Auxilia
no desenvolvimento do plano de vendas, estabelecendo participação
esperada no mercado, avaliando estratégia de marketing e analisando
vantagens competitivas do comércio eletrônico (e-commerce). Organiza os
arquivos dos clientes já atendidos e realiza prospecção de novos
clientes. Apresenta e divulga produtos e serviços, presencialmente e por
meio digital. Realiza visitas e estabelece contatos em áreas internas de
empresas, utilizando banco de dados com informações de clientes e
seguindo princípios de roteirização logística. Relaciona-se com os
clientes, seguindo procedimentos do tipo de canal de vendas. Concretiza
a venda, apresentando proposta ao cliente, calculando custo do produto e
estabelecendo prazos de entrega dos produtos. Pode negociar contratos e
ajudar na obtenção de crédito. Organiza e elabora documentos para o
fechamento da venda, conforme procedimentos fiscais e da empresa.
Realiza atendimento pós-venda, coordenando assessoria técnica ao
cliente, esclarecendo dúvidas, verificando opiniões do cliente e
colhendo sugestões para aprimorar o produto e a política de vendas. Atua
no sentido de fidelizar o cliente. Elabora relatórios e estatísticas.

\begin{Form}
\textbf{Esta correspondência está adequada?} \\ 
\ChoiceMenu[radio,radiosymbol=\ding{52},bordercolor=0 0.5 0,name=cbovalid_ 354135 ]{}{CBO deve ficar=sim, \hspace{6em} CBO não fica aqui=nao}
\end{Form}

\textbf{CBO: 354140  ( Técnico em atendimento e vendas )}

Analisa o comportamento do mercado consumidor, participando do
desenvolvimento de estudos mercadológicos. Estuda o cliente potencial da
empresa, analisando seu perfil e suas expectativas, visando ao
estabelecimento de possíveis estratégias de fidelização. Avalia qual
estratégia de marketing é mais adequada para cada tipo de cliente.
Analisa vantagens e desvantagens no uso de comércio eletrônico
(e-commerce), dependendo do tipo de cliente. Pesquisa ações dos
concorrentes, para analisar o que é oferecido aos clientes. Providencia
o cadastramento de clientes potenciais e efetivos, atualizando os dados
sempre que necessário. Pode usar cadastro de clientes centralizado da
plataforma de Gestão de Relacionamento com o Cliente (Customer
Relationship Management - CRM). Contata clientes, usando estratégias
digitais (e-mail, redes sociais ou outras) e/ou realizando visitas.
Comunica-se com os clientes, identificando suas necessidades e
expectativas e analisando seu potencial de compras. Divulga produtos e
serviços aos clientes, apresentando características e benefícios do
produto ou serviço, esclarecendo dados técnicos, informando preços e
detalhando formas de pagamento. Pode fornecer amostras e convidar
clientes para visitar a empresa. Concretiza vendas presencialmente, por
telefone ou por meios digitais, de acordo com as características dos
clientes. Acompanha o cliente após a venda -- inclusive por meio de
visita -- para sanar suas dúvidas, providenciar apoio técnico e
verificar suas opiniões e suas sugestões. Analisa reclamações e
insatisfações, para propor melhorias no processo de vendas. Elabora
relatórios e estatísticas sobre atendimento ao cliente e sobre o
processo de vendas. Presta assistência à equipe de vendedores de
produtos e serviços, com relação aos procedimentos adequados a cada tipo
de cliente nas etapas de pré-venda, venda e pós-venda. Participa da
elaboração de planos e estratégias de vendas, sugerindo melhorias no
atendimento ao cliente e no processo de vendas.

\begin{Form}
\textbf{Esta correspondência está adequada?} \\ 
\ChoiceMenu[radio,radiosymbol=\ding{52},bordercolor=0 0.5 0,name=cbovalid_ 354140 ]{}{CBO deve ficar=sim, \hspace{6em} CBO não fica aqui=nao}
\end{Form}

\textbf{CBO: 354705  ( Representante comercial autônomo )}

Realiza planejamento das atividades de representação comercial autônoma
- de empresas nacionais ou estrangeiras -, fazendo pesquisa de mercado,
analisando potencial de clientes e elaborando plano estratégico.
Estabelece metas de venda de produtos e serviços. Realiza divulgação de
produtos e serviços, definindo estratégias de marketing e requisitando
material de divulgação. Envia mala direta ao conjunto de clientes e
divulga produtos por meio de multimídia. Informa aos clientes sobre
alterações nos produtos e serviços, ressaltando as melhorias. Prepara
cadastro de novos clientes e estabelece contato com eles por telefone,
celular, e-mail e outros meios de comunicação. Realiza vendas aos novos
clientes. Planeja itinerário de visitas, agendando horários. Explica os
objetivos da visita, demonstra benefícios e qualidades de produtos e
serviços e fornece amostras. Auxilia na escolha, prestando informações
técnicas dos produtos. Apresenta tabelas de preços, explica formas de
pagamento e efetiva vendas. Recebe clientes para vendas, identificando
suas necessidades e apresentando proposta ao cliente. Calcula custo do
produto e serviço. Estabelece prazos de entrega dos produtos. Finaliza
as vendas, emitindo pedidos de produtos e serviços. Esclarece dúvidas
dos clientes referentes aos pedidos e aos contratos de prestação de
serviços. Monitora o cumprimento do prazo de entrega dos produtos.
Acompanha clientes pós-venda, verificando suas opiniões. Filtra
informações para melhoria de serviços, dando retorno às sugestões dos
clientes e enviando-lhes brindes. Renova e altera contratos de prestação
de serviços. Interage com as diversas áreas da empresa representada.
Requisita mostruário de produtos e solicita troca de produtos
defeituosos. Obtém informações de pré-venda com telemarketing. Comunica
transações comerciais realizadas e os resultados de vendas. Sugere
políticas de vendas, propondo ações para melhoria dos resultados. Visita
área industrial da empresa, participando de reuniões sobre promoção de
produtos e serviços. Mantém-se atualizado, participando de cursos sobre
novas estratégias de marketing. Participa de eventos, tais como feiras e
palestras relacionadas à área. Expõe produtos em locais definidos pela
empresa representada.

\begin{Form}
\textbf{Esta correspondência está adequada?} \\ 
\ChoiceMenu[radio,radiosymbol=\ding{52},bordercolor=0 0.5 0,name=cbovalid_ 354705 ]{}{CBO deve ficar=sim, \hspace{6em} CBO não fica aqui=nao}
\end{Form}

\textbf{CBO: 521115  ( Promotor de vendas )}

Efetua pesquisas sobre o potencial de mercado e a atuação da
concorrência, para subsidiar reuniões de definição de estratégias de
comercialização de linhas de produto. Prospecta um novo ponto de venda
(PDV), em sua região de atuação. Providencia a limpeza, a montagem e a
organização do ponto de venda. Orienta o abastecimento de mercadorias e
produtos. Orienta a montagem de áreas de exposição, considerando o
leiaute do ponto de venda e o plano de merchandising. Pode abordar e
orientar clientes sobre mercadorias, produtos e serviços no ponto de
venda, seguindo técnicas de atendimento. Acompanha o desempenho e os
resultados do ponto de venda. Promove produtos, visando ao aumento de
vendas e à fixação da marca. Realiza visitas -- presenciais e virtuais -
de divulgação de linhas de produto e de apresentação de estratégias de
comercialização, junto ao comércio varejista e a clientes em
perspectiva. Presta assistência aos revendedores e às equipes de venda,
demonstrando os produtos e fornecendo informações sobre seus manuseios e
usos. Em caso de solicitação de informações específicas de
aplicabilidade de um produto, pode fazer encaminhamento das questões
para a área técnica da empresa produtora. Pode auxiliar o vendedor nas
vendas do produto promovido. Pode distribuir brindes e oferecer
degustação de produtos. Pode promover produtos em exposições, feiras e
eventos. Elabora e organiza documentação em meio impresso e digital,
para a promoção de vendas junto ao comércio varejista. Prepara
documentação subsidiária para orientação a vendedores dos pontos de
venda, em meio impresso e digital.

\begin{Form}
\textbf{Esta correspondência está adequada?} \\ 
\ChoiceMenu[radio,radiosymbol=\ding{52},bordercolor=0 0.5 0,name=cbovalid_ 521115 ]{}{CBO deve ficar=sim, \hspace{6em} CBO não fica aqui=nao}
\end{Form}

\newpage
\section{ 05002 -- Técnico em Desenvolvimento de Sistemas }

\textbf{Perfil do Curso (CNCT)}

O Técnico em Desenvolvimento de Sistemas será habilitado para:

Desenvolver sistemas computacionais utilizando ambiente de
desenvolvimento. Dimensionar requisitos e funcionalidades do sistema.
Realizar testes funcionais de programas de computador e aplicativos.
Manter registros para análise e refinamento de resultados. Executar
manutenção de programas de computador e suporte técnico. Realizar
modelagem de aplicações computacionais. Codificar aplicações e rotinas
utilizando linguagens de programação específicas. Executar alterações e
manutenções em aplicações e rotinas de acordo com as definições
estabelecidas. Prestar apoio técnico na elaboração da documentação de
sistemas. Realizar prospecções, testes e avaliações de ferramentas e
produtos de desenvolvimento de sistemas.

Para atuação como Técnico em Desenvolvimento de Sistemas, são
fundamentais:

Conhecimentos e saberes relacionados aos processos de planejamento e
execução de projetos computacionais de forma a garantir a entrega de
produtos digitais, análise de softwares, testagem de protótipos, de
acordo com suas finalidades. Conhecimentos e saberes relacionados às
normas técnicas, à liderança de equipes, à solução de problemas técnicos
e à assertividade na comunicação de laudos e análises.

\textbf{CBOs correspondidos e seus perfis (presentes no CNCT, e presentes na predição da IA)}

\textbf{CBO: 317110  ( Desenvolvedor de sistemas de tecnologia da informação (técnico) )}

Desenvolve sistemas e aplicações de tecnologia da informação, elaborando
a interface gráfica e aplicando critérios ergonômicos de navegação.
Monta a estrutura do banco de dados e codifica os programas e
aplicativos, aplicando sistemas de rotinas de segurança. Compila os
programas. Testa aplicativos e programas, elaborando casos de testes.
Gera aplicativos para instalação e gerenciamento dos sistemas. Documenta
os sistemas e aplicações e avalia o desempenho dos produtos
desenvolvidos. Implanta sistemas e aplicações de tecnologia da
informação, instalando os programas e configurando os equipamentos em
que a aplicação será executada. Homologa os sistemas e aplicações
desenvolvidos, avaliando e validando os resultados da implantação.
Elabora material para capacitação de usuários. Avalia os objetivos e as
metas de projetos de sistemas e aplicações, tendo em vista os resultados
obtidos na fase de implantação. Publica o código final no servidor de
sistemas e aplicações. Realiza a manutenção de sistemas e aplicações de
tecnologia da informação, atualizando informações gráficas, textuais e
audiovisuais e convertendo seus códigos fontes para outras linguagens ou
plataformas. Implementa alterações nos sistemas e aplicações e em suas
estruturas de armazenamento de dados, adequando-os aos sistemas e
ambientes operacionais. Monitora a performance dos sistemas e
aplicações, utilizando ferramentas de software, para obter subsídios à
atividade de manutenção. Atualiza a documentação dos sistemas e
aplicações, em função das intervenções de manutenção. Fornece suporte
técnico para o cliente interno. Colabora na seleção dos recursos de
desenvolvimento, especificando máquinas, ferramentas, acessórios e
suprimentos. Participa da seleção das linguagens de programação e
metodologias. Seleciona ferramentas de desenvolvimento. Participa da
definição da nomenclatura padrão. Colabora no planejamento das etapas e
ações de trabalho, participando das definições de atividades, tarefas e
cronograma. Participa de reuniões com a equipe de trabalho ou com o
cliente. Participa das definições de padronizações. Projeta sistemas e
aplicações de tecnologia da informação, identificando a demanda do
cliente, coletando dados e desenvolvendo o leiaute de telas e
relatórios. Elabora o pré-projeto e os projetos conceitual, lógico,
estrutural, físico e gráfico dos sistemas e aplicações, definindo os
requisitos de projeto. Participa das definições dos critérios
ergonômicos de navegação e da interface de comunicação e interatividade.
Dimensiona a vida útil dos sistemas e aplicações, modela sua estrutura
de banco de dados e dimensiona a capacidade de armazenamento de dados
dos sistemas.

\begin{Form}
\textbf{Esta correspondência está adequada?} \\ 
\ChoiceMenu[radio,radiosymbol=\ding{52},bordercolor=0 0.5 0,name=cbovalid_ 317110 ]{}{CBO deve ficar=sim, \hspace{6em} CBO não fica aqui=nao}
\end{Form}

\newpage
\section{ 05003 -- Técnico em Informática }

\textbf{Perfil do Curso (CNCT)}

O Técnico em Informática será habilitado para:

Desenvolver sistemas computacionais utilizando ambiente de
desenvolvimento. Realizar modelagem, desenvolvimento, testes,
implementação e manutenção de sistemas computacionais. Modelar,
construir e realizar manutenção de banco de dados. Executar montagem,
instalação e configuração de equipamentos de informática. Instalar e
configurar sistemas operacionais e aplicativos em equipamentos
computacionais. Realizar manutenção preventiva e corretiva de
equipamentos de informática. Instalar e configurar dispositivos de
acesso à rede e realizar testes de conectividade. Realizar atendimento
help-desk. Operar, instalar, configurar e realizar manutenção em redes
de computadores. Aplicar técnicas de instalação e configuração da rede
física e lógica. Instalar, configurar e administrar sistemas
operacionais em redes de computadores. Executar as rotinas de
monitoramento do ambiente operacional. Identificar e registrar os
desvios e adotar os procedimentos de correção. Executar procedimentos de
segurança, pré-definidos, para ambiente de rede.

Para atuação como Técnico em Informática, são fundamentais:

Conhecimentos e saberes relacionados aos processos de planejamento e
execução de projetos computacionais de forma a garantir a entrega de
produtos digitais, análise de softwares, testagem de protótipos, de
acordo com suas finalidades. Conhecimentos e saberes relacionados às
normas técnicas, à liderança de equipes, à solução de problemas técnicos
e à assertividade na comunicação de laudos e análises. Habilidades
relacionadas à construção de soluções em BI e integrações sistêmicas.

\textbf{CBOs correspondidos e seus perfis (presentes no CNCT, e presentes na predição da IA)}

\textbf{CBO: 317105  ( Desenvolvedor web (técnico) )}

Desenvolve sistemas e aplicações para web, elaborando a interface
gráfica e aplicando critérios ergonômicos de navegação. Monta a
estrutura do banco de dados e codifica os programas e aplicativos,
aplicando sistemas de rotinas de segurança. Compila os programas. Testa
aplicativos e programas, elaborando casos de testes. Gera aplicativos
para instalação e gerenciamento dos sistemas. Documenta os sistemas e
aplicações e avalia o desempenho dos produtos desenvolvidos. Implanta
sistemas e aplicações para web, instalando os programas, adaptando o
conteúdo para mídias interativas e configurando os equipamentos em que a
aplicação será executada. Homologa os sistemas e aplicações
desenvolvidos, avaliando e validando os resultados da implantação.
Elabora material para capacitação de usuários. Avalia os objetivos e as
metas de projetos de sistemas e aplicações, tendo em vista os resultados
obtidos na fase de implantação. Publica o código final no servidor de
sistemas e aplicações. Realiza a manutenção de sistemas e aplicações
para web, atualizando informações gráficas, textuais e audiovisuais e
convertendo seus códigos fontes para outras linguagens ou plataformas.
Implementa alterações nos sistemas e aplicações e em suas estruturas de
armazenamento de dados, adequando-os aos sistemas e ambientes
operacionais. Monitora a performance dos sistemas e aplicações,
utilizando ferramentas de software, para obter subsídios à atividade de
manutenção. Atualiza a documentação dos sistemas e aplicações, em função
das intervenções de manutenção. Fornece suporte técnico para o cliente
interno. Colabora na seleção dos recursos de desenvolvimento,
especificando máquinas, ferramentas, acessórios e suprimentos. Participa
da seleção das linguagens de programação e metodologias. Seleciona
ferramentas de desenvolvimento. Participa da definição da nomenclatura
padrão. Colabora no planejamento das etapas e ações de trabalho,
participando das definições de atividades, tarefas e cronograma.
Participa de reuniões com a equipe de trabalho ou com o cliente.
Participa das definições de padronizações. Projeta sistemas e aplicações
para web, identificando a demanda do cliente, coletando dados e
desenvolvendo o leiaute de telas e relatórios. Elabora o pré-projeto e
os projetos conceitual, lógico, estrutural, físico e gráfico dos
sistemas e aplicações, definindo os requisitos de projeto. Participa das
definições dos critérios ergonômicos de navegação e da interface de
comunicação e interatividade. Dimensiona a vida útil dos sistemas e
aplicações, modela sua estrutura de banco de dados e dimensiona a
capacidade de armazenamento de dados dos sistemas.

\begin{Form}
\textbf{Esta correspondência está adequada?} \\ 
\ChoiceMenu[radio,radiosymbol=\ding{52},bordercolor=0 0.5 0,name=cbovalid_ 317105 ]{}{CBO deve ficar=sim, \hspace{6em} CBO não fica aqui=nao}
\end{Form}

\textbf{CBO: 317110  ( Desenvolvedor de sistemas de tecnologia da informação (técnico) )}

Desenvolve sistemas e aplicações de tecnologia da informação, elaborando
a interface gráfica e aplicando critérios ergonômicos de navegação.
Monta a estrutura do banco de dados e codifica os programas e
aplicativos, aplicando sistemas de rotinas de segurança. Compila os
programas. Testa aplicativos e programas, elaborando casos de testes.
Gera aplicativos para instalação e gerenciamento dos sistemas. Documenta
os sistemas e aplicações e avalia o desempenho dos produtos
desenvolvidos. Implanta sistemas e aplicações de tecnologia da
informação, instalando os programas e configurando os equipamentos em
que a aplicação será executada. Homologa os sistemas e aplicações
desenvolvidos, avaliando e validando os resultados da implantação.
Elabora material para capacitação de usuários. Avalia os objetivos e as
metas de projetos de sistemas e aplicações, tendo em vista os resultados
obtidos na fase de implantação. Publica o código final no servidor de
sistemas e aplicações. Realiza a manutenção de sistemas e aplicações de
tecnologia da informação, atualizando informações gráficas, textuais e
audiovisuais e convertendo seus códigos fontes para outras linguagens ou
plataformas. Implementa alterações nos sistemas e aplicações e em suas
estruturas de armazenamento de dados, adequando-os aos sistemas e
ambientes operacionais. Monitora a performance dos sistemas e
aplicações, utilizando ferramentas de software, para obter subsídios à
atividade de manutenção. Atualiza a documentação dos sistemas e
aplicações, em função das intervenções de manutenção. Fornece suporte
técnico para o cliente interno. Colabora na seleção dos recursos de
desenvolvimento, especificando máquinas, ferramentas, acessórios e
suprimentos. Participa da seleção das linguagens de programação e
metodologias. Seleciona ferramentas de desenvolvimento. Participa da
definição da nomenclatura padrão. Colabora no planejamento das etapas e
ações de trabalho, participando das definições de atividades, tarefas e
cronograma. Participa de reuniões com a equipe de trabalho ou com o
cliente. Participa das definições de padronizações. Projeta sistemas e
aplicações de tecnologia da informação, identificando a demanda do
cliente, coletando dados e desenvolvendo o leiaute de telas e
relatórios. Elabora o pré-projeto e os projetos conceitual, lógico,
estrutural, físico e gráfico dos sistemas e aplicações, definindo os
requisitos de projeto. Participa das definições dos critérios
ergonômicos de navegação e da interface de comunicação e interatividade.
Dimensiona a vida útil dos sistemas e aplicações, modela sua estrutura
de banco de dados e dimensiona a capacidade de armazenamento de dados
dos sistemas.

\begin{Form}
\textbf{Esta correspondência está adequada?} \\ 
\ChoiceMenu[radio,radiosymbol=\ding{52},bordercolor=0 0.5 0,name=cbovalid_ 317110 ]{}{CBO deve ficar=sim, \hspace{6em} CBO não fica aqui=nao}
\end{Form}

\textbf{CBO: 317205  ( Operador de computador )}

Assegura o funcionamento do hardware e do software, inicializando
sistemas e aplicativos e realizando as diversas configurações -- de
hardware, software e sistemas operacionais --, de acordo com os
requisitos do processamento de dados. Desativa sistemas e aplicativos,
conforme a necessidade. Formata equipamentos e instala softwares e/ou
sistemas operacionais. Realiza diagnóstico em hardware, software e redes
de comunicação, identificando falhas, para possibilitar suas correções.
Requisita manutenção preventiva e/ou corretiva de hardware e/ou
software, podendo, também, executar os serviços. Realiza revisão em
hardware. Assegura o funcionamento de equipamento reserva (standby),
acionando-o quando necessário. Repõe suprimentos em equipamentos.
Monitora sistemas e aplicações, notificando problemas de infraestrutura.
Realiza o monitoramento dos recursos de rede e dos dispositivos de
entrada e saída de dados. Monitora os sistemas de banco de dados, a
disponibilidade e a funcionalidade de aplicativos bem como os registros
de erros e o consumo percentual do tempo da unidade central de
processamento (Central Processing Unit -- CPU), de cada tarefa.
Administra o processamento de dados, acompanhando os recursos
disponíveis e criando script (sequência de instruções) para melhoria da
performance. Providencia correção de erros de tarefas, notificando o
setor de suporte. Administra cronograma de atividades planejadas,
schedule (agenda) e escalonamento de tarefas. Presta suporte ao
cliente/usuário, disponibilizando recursos operacionais, reparando
arquivos e reprocessando tarefas de acordo com solicitação. Recupera e
transfere arquivos, programas e relatórios. Atende cliente/usuário,
planejando as soluções para as demandas apresentadas, orientando-o
quanto à utilização de hardware e/ou software e conduzindo a solicitação
de suporte. Administra chamados abertos e gerencia incidentes e
problemas. Administra a segurança das informações, fazendo cópias de
segurança (backup) e guardando-as em local prescrito. Verifica acesso
lógico do usuário, destrói informações sigilosas descartadas e faz
rodízio de mídias. Verifica temperatura, umidade local, equipamentos de
energia, sistema de detecção de incêndio, dentre outras condições
técnicas do ambiente de trabalho. Organiza o cabeamento dos
equipamentos.

\begin{Form}
\textbf{Esta correspondência está adequada?} \\ 
\ChoiceMenu[radio,radiosymbol=\ding{52},bordercolor=0 0.5 0,name=cbovalid_ 317205 ]{}{CBO deve ficar=sim, \hspace{6em} CBO não fica aqui=nao}
\end{Form}

\textbf{CBO: 317210  ( Técnico de suporte ao usuário de tecnologia da informação )}

Presta suporte ao cliente/usuário de tecnologia da informação,
realizando o atendimento inicial e planejando as soluções para os
chamados de acordo com o diagnóstico. Configura contas de
clientes/usuários e orienta-os quanto à utilização de hardware e/ou
software. Conduz a solicitação de suporte. Administra chamados abertos.
Disponibiliza recursos operacionais e repara arquivos. Recupera e
transfere arquivos, programas e relatórios. Registra as ações dos
chamados. Monitora sistemas e aplicações, notificando problemas de
infraestrutura. Realiza o monitoramento dos recursos de rede e dos
dispositivos de entrada e saída de dados. Monitora a disponibilidade e a
funcionalidade de aplicativos. Providencia correção de erros de tarefas,
notificando o setor de suporte. Assegura o funcionamento do hardware e
do software, inicializando sistemas e aplicativos e realizando as
diversas configurações -- de hardware, software e de sistemas
operacionais --, de acordo com os requisitos do processamento de dados.
Desativa sistemas e aplicativos, conforme a necessidade. Formata
equipamentos e instala softwares e/ou sistemas operacionais. Realiza
diagnóstico em hardware, software e redes de comunicação, identificando
falhas, para possibilitar suas correções. Requisita manutenção
preventiva e/ou corretiva de hardware e/ou software, podendo, também,
executar os serviços. Realiza revisão em hardware. Assegura o
funcionamento de equipamento reserva (standby). Repõe suprimentos em
equipamentos. Administra a segurança das informações, realizando
varredura e eliminação de vírus e garantindo o uso, pelo cliente, de
versões homologadas de software. Faz cópias de segurança (backup),
guardando-as em local prescrito. Verifica acesso lógico do usuário.
Verifica temperatura, umidade local e equipamentos de energia, dentre
outras condições técnicas do ambiente de trabalho. Recomenda mudanças na
disposição de equipamentos e organiza seu cabeamento.

\begin{Form}
\textbf{Esta correspondência está adequada?} \\ 
\ChoiceMenu[radio,radiosymbol=\ding{52},bordercolor=0 0.5 0,name=cbovalid_ 317210 ]{}{CBO deve ficar=sim, \hspace{6em} CBO não fica aqui=nao}
\end{Form}

\newpage
\section{ 05004 -- Técnico em Informática para Internet }

\textbf{Perfil do Curso (CNCT)}

O Técnico em Informática para internet será habilitado para:

Planejar e documentar aplicações para Web e dispositivos móveis.
Desenvolver e organizar elementos estruturais e visuais de aplicações
para Web e dispositivos móveis. Monitorar projetos de aplicações para
Web e dispositivos móveis. Estruturar e implementar banco de dados para
aplicações Web. Codificar aplicações para Web e dispositivos móveis.
Publicar e testar aplicações para Web e dispositivos móveis. Documentar
e realizar manutenção de aplicações para Web e dispositivos móveis.

Para atuação como Técnico em Informática para Internet, são
fundamentais:

Conhecimentos e saberes relacionados aos processos de planejamento e
execução de projetos em websites focados na experiência do usuário, na
testagem e análises de produtos web, na liderança de equipe e na ética
profissional.

\textbf{CBOs correspondidos e seus perfis (presentes no CNCT, e presentes na predição da IA)}

\textbf{CBO: 317110  ( Desenvolvedor de sistemas de tecnologia da informação (técnico) )}

Desenvolve sistemas e aplicações de tecnologia da informação, elaborando
a interface gráfica e aplicando critérios ergonômicos de navegação.
Monta a estrutura do banco de dados e codifica os programas e
aplicativos, aplicando sistemas de rotinas de segurança. Compila os
programas. Testa aplicativos e programas, elaborando casos de testes.
Gera aplicativos para instalação e gerenciamento dos sistemas. Documenta
os sistemas e aplicações e avalia o desempenho dos produtos
desenvolvidos. Implanta sistemas e aplicações de tecnologia da
informação, instalando os programas e configurando os equipamentos em
que a aplicação será executada. Homologa os sistemas e aplicações
desenvolvidos, avaliando e validando os resultados da implantação.
Elabora material para capacitação de usuários. Avalia os objetivos e as
metas de projetos de sistemas e aplicações, tendo em vista os resultados
obtidos na fase de implantação. Publica o código final no servidor de
sistemas e aplicações. Realiza a manutenção de sistemas e aplicações de
tecnologia da informação, atualizando informações gráficas, textuais e
audiovisuais e convertendo seus códigos fontes para outras linguagens ou
plataformas. Implementa alterações nos sistemas e aplicações e em suas
estruturas de armazenamento de dados, adequando-os aos sistemas e
ambientes operacionais. Monitora a performance dos sistemas e
aplicações, utilizando ferramentas de software, para obter subsídios à
atividade de manutenção. Atualiza a documentação dos sistemas e
aplicações, em função das intervenções de manutenção. Fornece suporte
técnico para o cliente interno. Colabora na seleção dos recursos de
desenvolvimento, especificando máquinas, ferramentas, acessórios e
suprimentos. Participa da seleção das linguagens de programação e
metodologias. Seleciona ferramentas de desenvolvimento. Participa da
definição da nomenclatura padrão. Colabora no planejamento das etapas e
ações de trabalho, participando das definições de atividades, tarefas e
cronograma. Participa de reuniões com a equipe de trabalho ou com o
cliente. Participa das definições de padronizações. Projeta sistemas e
aplicações de tecnologia da informação, identificando a demanda do
cliente, coletando dados e desenvolvendo o leiaute de telas e
relatórios. Elabora o pré-projeto e os projetos conceitual, lógico,
estrutural, físico e gráfico dos sistemas e aplicações, definindo os
requisitos de projeto. Participa das definições dos critérios
ergonômicos de navegação e da interface de comunicação e interatividade.
Dimensiona a vida útil dos sistemas e aplicações, modela sua estrutura
de banco de dados e dimensiona a capacidade de armazenamento de dados
dos sistemas.

\begin{Form}
\textbf{Esta correspondência está adequada?} \\ 
\ChoiceMenu[radio,radiosymbol=\ding{52},bordercolor=0 0.5 0,name=cbovalid_ 317110 ]{}{CBO deve ficar=sim, \hspace{6em} CBO não fica aqui=nao}
\end{Form}

\newpage
\section{ 05005 -- Técnico em Manutenção e Suporte em Informática }

\textbf{Perfil do Curso (CNCT)}

O Técnico em Manutenção e Suporte em Informática será habilitado para:

Realizar montagem, diagnóstico, manutenção e instalação de computadores.
Instalar e configurar software (sistema operacional e aplicativos) para
desktop e servidores. Realizar instalação e manutenção de redes de
computadores. Realizar manutenção preventiva e corretiva de computadores
e periféricos. Prestar assistência técnica aos usuários em relação à
utilização dos serviços de TI. Auxiliar nas atividades de infraestrutura
de TI, mantendo a disponibilidade de sistemas. Prestar suporte ao
ambiente interno, instalação e configuração de sistemas operacionais,
redes e impressoras. Identificar problemas e/ou dificuldades de acesso e
utilização de aplicações. Acompanhar e avaliar os níveis de serviços
prestados. Analisar a requisição ou problema apresentado, identificando
a complexidade técnica para atuar na solução e direcionar para
atendimento de acordo com nível técnico correspondente. Verificar os
sistemas das requisições e incidentes na fila de atendimento e analisar
a prioridade conforme a urgência de cada caso. Detectar e diagnosticar,
pessoalmente, os sintomas apresentados pelo equipamento de um
solicitante, fisicamente ou virtualmente, verificando as condições de
funcionamento das instalações físicas e do sistema, para tomar as
providências necessárias de acordo com o problema apresentado. Responder
pela organização e controle de peças e equipamentos quando retirados do
estoque, controlando a logística e movimentação deles. Configurar
equipamentos para novos funcionários ou postos de trabalho, registrando
os dados (protocolos de identificação, e-mail, perfil, dispositivos
móveis) no equipamento destinado ao funcionário. Realizar constante
manutenção nos equipamentos, substituindo componentes/periféricos quando
necessário, visando garantir o funcionamento adequado. Recolher
equipamentos usados (que não serão mais utilizados pelos funcionários),
realizar a formatação e substituição de peças, otimizando o hardware
(upgrade) com o objetivo de disponibilizar o equipamento a outro
colaborador. Estabelecer comunicação oral e escrita para agilizar o
trabalho, redigir documentação técnica e organizar o local de trabalho.

Para atuação como Manutenção e Suporte em Informática, são fundamentais:

Conhecimentos e saberes relacionados aos processos de configurações de
dispositivos de informática, resolução de problemas relacionados às
diversas tecnologias. Saberes relacionados às práticas de lideranças de
equipe, de boas práticas de comunicação e de preservação das boas
práticas de uso de tecnologias.

\textbf{CBOs correspondidos e seus perfis (presentes no CNCT, e presentes na predição da IA)}

\textbf{CBO: 313220  ( Técnico em manutenção de equipamentos de informática )}

Conserta e faz manutenção corretiva em equipamentos de informática -
como computadores e seus periféricos, dispositivos de acesso a redes e
fontes chaveadas -, interpretando esquemas elétricos e substituindo
componentes defeituosos. Pode utilizar equipamentos de diagnóstico
eletrônico e fazer uso da internet para acesso remoto. Instala
equipamentos de informática, observando as condições do ambiente,
realizando ajustes e testes. Faz manutenções preventivas e preditivas em
equipamentos de informática, de acordo com planos de manutenção,
presencialmente ou online. Sugere alterações nos processos produtivos,
para melhoria do sistema, integrando dispositivos e sistemas eletrônicos
por meio de redes digitais de comunicação de dados, locais e a
distância. Desenvolve, monta e testa dispositivos de circuitos
eletrônicos para equipamentos de informática, especificando e calculando
custos de componentes. Pode usar microcontroladores, microprocessadores
e outros dispositivos eletrônicos digitais programáveis. Descreve
procedimentos de trabalho, preenche laudos técnicos e formulário de
reposição de peças rejeitadas e emite relatórios técnicos. Elabora
gráficos de resultados. Pode treinar, orientar e avaliar o desempenho de
operadores.

\begin{Form}
\textbf{Esta correspondência está adequada?} \\ 
\ChoiceMenu[radio,radiosymbol=\ding{52},bordercolor=0 0.5 0,name=cbovalid_ 313220 ]{}{CBO deve ficar=sim, \hspace{6em} CBO não fica aqui=nao}
\end{Form}

\textbf{CBO: 317210  ( Técnico de suporte ao usuário de tecnologia da informação )}

Presta suporte ao cliente/usuário de tecnologia da informação,
realizando o atendimento inicial e planejando as soluções para os
chamados de acordo com o diagnóstico. Configura contas de
clientes/usuários e orienta-os quanto à utilização de hardware e/ou
software. Conduz a solicitação de suporte. Administra chamados abertos.
Disponibiliza recursos operacionais e repara arquivos. Recupera e
transfere arquivos, programas e relatórios. Registra as ações dos
chamados. Monitora sistemas e aplicações, notificando problemas de
infraestrutura. Realiza o monitoramento dos recursos de rede e dos
dispositivos de entrada e saída de dados. Monitora a disponibilidade e a
funcionalidade de aplicativos. Providencia correção de erros de tarefas,
notificando o setor de suporte. Assegura o funcionamento do hardware e
do software, inicializando sistemas e aplicativos e realizando as
diversas configurações -- de hardware, software e de sistemas
operacionais --, de acordo com os requisitos do processamento de dados.
Desativa sistemas e aplicativos, conforme a necessidade. Formata
equipamentos e instala softwares e/ou sistemas operacionais. Realiza
diagnóstico em hardware, software e redes de comunicação, identificando
falhas, para possibilitar suas correções. Requisita manutenção
preventiva e/ou corretiva de hardware e/ou software, podendo, também,
executar os serviços. Realiza revisão em hardware. Assegura o
funcionamento de equipamento reserva (standby). Repõe suprimentos em
equipamentos. Administra a segurança das informações, realizando
varredura e eliminação de vírus e garantindo o uso, pelo cliente, de
versões homologadas de software. Faz cópias de segurança (backup),
guardando-as em local prescrito. Verifica acesso lógico do usuário.
Verifica temperatura, umidade local e equipamentos de energia, dentre
outras condições técnicas do ambiente de trabalho. Recomenda mudanças na
disposição de equipamentos e organiza seu cabeamento.

\begin{Form}
\textbf{Esta correspondência está adequada?} \\ 
\ChoiceMenu[radio,radiosymbol=\ding{52},bordercolor=0 0.5 0,name=cbovalid_ 317210 ]{}{CBO deve ficar=sim, \hspace{6em} CBO não fica aqui=nao}
\end{Form}

\newpage
\section{ 05007 -- Técnico em Programação de Jogos Digitais }

\textbf{Perfil do Curso (CNCT)}

O Técnico em Programação de Jogos Digitais será habilitado para:

Planejar o desenvolvimento do jogo digital para multiplataformas.
Planejar as atividades de programação para o desenvolvimento do jogo
digital. Configurar e incorporar elementos multimídia à plataforma de
desenvolvimento. Desenvolver e selecionar os algoritmos e a estrutura de
dados para jogos digitais. Programar e integrar os elementos multimídia
do jogo digital para computadores, consoles e dispositivos móveis.
Programar jogos digitais multiplayer. Realizar testes em jogos digitais.
Realizar manutenção de jogos digitais.

Para atuação como Técnico em Programação de Jogos Digitais, são
fundamentais:

Conhecimentos e saberes relacionados aos processos de produção de
conteúdo e roteirização, bem como à expertise da produção em equipe, aos
métodos de liderança, às boas práticas de comunicação e produção de
análises e pareceres técnicos, sempre garantindo o pleno atendimento dos
prazos, dos critérios de qualidade e do perfil técnico de suas
produções.

\textbf{CBOs correspondidos e seus perfis (presentes no CNCT, e presentes na predição da IA)}

\textbf{CBO: 317120  ( Desenvolvedor de multimídia )}

Desenvolve sistemas e aplicações de multimídia, elaborando a interface
gráfica e aplicando critérios ergonômicos de navegação. Monta a
estrutura do banco de dados e codifica os programas e aplicativos,
aplicando sistemas de rotinas de segurança. Compila os programas. Testa
aplicativos e programas, elaborando casos de testes. Gera aplicativos
para instalação e gerenciamento dos sistemas. Documenta sistemas e
aplicações e avalia o desempenho dos produtos desenvolvidos. Implanta
sistemas e aplicações de multimídia, instalando os programas, adaptando
o conteúdo para mídias interativas e configurando os equipamentos em que
a aplicação será executada. Homologa sistemas e aplicações
desenvolvidos, avaliando e validando os resultados da implantação.
Elabora material para capacitação de usuários. Avalia os objetivos e as
metas de projetos de sistemas e aplicações, tendo em vista os resultados
obtidos na fase de implantação. Publica o código final no servidor de
sistemas e aplicações. Realiza a manutenção de sistemas e aplicações de
multimídia, atualizando informações gráficas, textuais e audiovisuais.
Implementa alterações em sistemas e aplicações e em suas estruturas de
armazenamento de dados, adequando-os aos sistemas e ambientes
operacionais. Monitora a performance de sistemas e aplicações,
utilizando ferramentas de software, para obter subsídios à atividade de
manutenção. Atualiza a documentação de sistemas e aplicações, em função
das intervenções de manutenção. Fornece suporte técnico para o cliente
interno. Colabora na seleção dos recursos de desenvolvimento,
especificando máquinas, ferramentas, acessórios e suprimentos. Participa
da seleção das linguagens de programação e metodologias. Seleciona
ferramentas de desenvolvimento. Participa da definição da nomenclatura
padrão. Colabora no planejamento das etapas e ações de trabalho,
participando das definições de atividades, tarefas e cronograma.
Participa de reuniões com a equipe de trabalho ou com o cliente.
Participa das definições de padronizações. Projeta sistemas e aplicações
de multimídia, identificando a demanda do cliente, coletando dados e
desenvolvendo o leiaute de telas e relatórios. Elabora o pré-projeto e
os projetos conceitual, lógico, estrutural, físico e gráfico dos
sistemas e aplicações, definindo os requisitos de projeto. Participa das
definições dos critérios ergonômicos de navegação e da interface de
comunicação e interatividade. Dimensiona a vida útil dos sistemas e
aplicações.

\begin{Form}
\textbf{Esta correspondência está adequada?} \\ 
\ChoiceMenu[radio,radiosymbol=\ding{52},bordercolor=0 0.5 0,name=cbovalid_ 317120 ]{}{CBO deve ficar=sim, \hspace{6em} CBO não fica aqui=nao}
\end{Form}

\newpage
\section{ 05008 -- Técnico em Redes de Computadores }

\textbf{Perfil do Curso (CNCT)}

O Técnico em Redes de Computadores será habilitado para:

Instalar, configurar e operar sistemas de redes computacionais. Executar
cabeamento de redes industriais e comerciais. Configurar e dimensionar
sistemas de protocolos de redes de comunicação de equipamentos
computacionais e equipamentos de produção industrial e controle
comercial. Monitorar o ambiente de rede e executar as rotinas
pré-estabelecidas de administração de ambiente de TI. Identificar e
corrigir desvios relacionados a recursos de rede, conforme procedimentos
pré-definidos. Operar, realizar testes e homologar recursos de rede,
conforme requisitos pré-definidos. Executar procedimentos de segurança
pré-definidos para ambiente de rede. Instalar, programar, configurar e
customizar os recursos de rede, de acordo com os procedimentos
operacionais e padrões técnicos pré-definidos. Instalar, configurar e
disponibilizar softwares aplicativos e plataformas operacionais em rede
local, de acordo com os procedimentos operacionais e padrões técnicos
pré-definidos. Efetuar o cadastramento e a habilitação de usuários no
ambiente de rede. Prestar assistência técnica e orientar usuários quanto
à utilização dos recursos de rede. Coletar informações e elaborar
relatórios técnicos para acompanhamento e contabilização dos serviços de
rede. Executar a medição dos serviços de rede, verificando o cumprimento
dos níveis de serviços. Verificar a segurança da rede e a transmissão de
dados, como também testar, periodicamente, a vulnerabilidade da rede em
possíveis ataques. Instalar, configurar e atender problemas relacionados
a produtos que se conectam em redes domésticas e corporativas ? Internet
das Coisas (IOT).

Para atuação como Técnico em Redes de Computadores, são fundamentais:

Conhecimentos e saberes relacionados aos processos operacionais de
soluções em rede, em computadores e tecnologias sensíveis ao processo de
controle operacional das redes, bem como aos métodos e práticas de
conectividade interna e externa, sempre garantindo o pleno atendimento
dos prazos, dos critérios de qualidade e do perfil técnico.

\textbf{CBOs correspondidos e seus perfis (presentes no CNCT, e presentes na predição da IA)}

\textbf{CBO: 313220  ( Técnico em manutenção de equipamentos de informática )}

Conserta e faz manutenção corretiva em equipamentos de informática -
como computadores e seus periféricos, dispositivos de acesso a redes e
fontes chaveadas -, interpretando esquemas elétricos e substituindo
componentes defeituosos. Pode utilizar equipamentos de diagnóstico
eletrônico e fazer uso da internet para acesso remoto. Instala
equipamentos de informática, observando as condições do ambiente,
realizando ajustes e testes. Faz manutenções preventivas e preditivas em
equipamentos de informática, de acordo com planos de manutenção,
presencialmente ou online. Sugere alterações nos processos produtivos,
para melhoria do sistema, integrando dispositivos e sistemas eletrônicos
por meio de redes digitais de comunicação de dados, locais e a
distância. Desenvolve, monta e testa dispositivos de circuitos
eletrônicos para equipamentos de informática, especificando e calculando
custos de componentes. Pode usar microcontroladores, microprocessadores
e outros dispositivos eletrônicos digitais programáveis. Descreve
procedimentos de trabalho, preenche laudos técnicos e formulário de
reposição de peças rejeitadas e emite relatórios técnicos. Elabora
gráficos de resultados. Pode treinar, orientar e avaliar o desempenho de
operadores.

\begin{Form}
\textbf{Esta correspondência está adequada?} \\ 
\ChoiceMenu[radio,radiosymbol=\ding{52},bordercolor=0 0.5 0,name=cbovalid_ 313220 ]{}{CBO deve ficar=sim, \hspace{6em} CBO não fica aqui=nao}
\end{Form}

\textbf{CBO: 317210  ( Técnico de suporte ao usuário de tecnologia da informação )}

Presta suporte ao cliente/usuário de tecnologia da informação,
realizando o atendimento inicial e planejando as soluções para os
chamados de acordo com o diagnóstico. Configura contas de
clientes/usuários e orienta-os quanto à utilização de hardware e/ou
software. Conduz a solicitação de suporte. Administra chamados abertos.
Disponibiliza recursos operacionais e repara arquivos. Recupera e
transfere arquivos, programas e relatórios. Registra as ações dos
chamados. Monitora sistemas e aplicações, notificando problemas de
infraestrutura. Realiza o monitoramento dos recursos de rede e dos
dispositivos de entrada e saída de dados. Monitora a disponibilidade e a
funcionalidade de aplicativos. Providencia correção de erros de tarefas,
notificando o setor de suporte. Assegura o funcionamento do hardware e
do software, inicializando sistemas e aplicativos e realizando as
diversas configurações -- de hardware, software e de sistemas
operacionais --, de acordo com os requisitos do processamento de dados.
Desativa sistemas e aplicativos, conforme a necessidade. Formata
equipamentos e instala softwares e/ou sistemas operacionais. Realiza
diagnóstico em hardware, software e redes de comunicação, identificando
falhas, para possibilitar suas correções. Requisita manutenção
preventiva e/ou corretiva de hardware e/ou software, podendo, também,
executar os serviços. Realiza revisão em hardware. Assegura o
funcionamento de equipamento reserva (standby). Repõe suprimentos em
equipamentos. Administra a segurança das informações, realizando
varredura e eliminação de vírus e garantindo o uso, pelo cliente, de
versões homologadas de software. Faz cópias de segurança (backup),
guardando-as em local prescrito. Verifica acesso lógico do usuário.
Verifica temperatura, umidade local e equipamentos de energia, dentre
outras condições técnicas do ambiente de trabalho. Recomenda mudanças na
disposição de equipamentos e organiza seu cabeamento.

\begin{Form}
\textbf{Esta correspondência está adequada?} \\ 
\ChoiceMenu[radio,radiosymbol=\ding{52},bordercolor=0 0.5 0,name=cbovalid_ 317210 ]{}{CBO deve ficar=sim, \hspace{6em} CBO não fica aqui=nao}
\end{Form}

\newpage
\section{ 05009 -- Técnico em Telecomunicações }

\textbf{Perfil do Curso (CNCT)}

O Técnico em Telecomunicações será habilitado para:

Participar na elaboração de projetos de telecomunicações. Instalar,
testar e realizar manutenções preventivas e corretivas em sistemas de
telecomunicações. Configurar equipamentos nas áreas de telefonia,
transmissão e redes de comunicação. Supervisionar tecnicamente processos
e serviços de telecomunicações. Elaborar documentação técnica. Prestar
assistência técnica aos clientes. Realizar programação de softwares
específicos para equipamentos de telecomunicações. Participar na
elaboração da documentação técnica

Para atuação como Técnico em Telecomunicações, são fundamentais:

Conhecimentos e saberes relacionados aos processos técnicos de
telecomunicação cabeada ou de transmissão/tráfego de dados móveis, bem
como às boas práticas de comunicação e de liderança de equipes.

\textbf{CBOs correspondidos e seus perfis (presentes no CNCT, e presentes na predição da IA)}

\textbf{CBO: 313305  ( Técnico de comunicação de dados )}

Instala sistemas de comunicação de dados, certificando-se da adequação
da infraestrutura e do fornecimento de energia, fixando componentes do
sistema, estabelecendo a interligação dos equipamentos via interfaces e
cabeamento e ativando o sistema. Testa sistemas de comunicação de dados,
providenciando configuração e programação, verificando se o
funcionamento está de acordo com as especificações e realizando
medições, aferições e ajustes. Pode substituir componentes, quando
necessário. Presta orientação ao usuário final sobre a utilização do
sistema. Realiza manutenção preventiva e corretiva em redes de
comunicação de dados, executando rotinas de teste, identificando e
corrigindo falhas, reprogramando e configurando programas do sistema.
Pode participar da elaboração de projetos em sistemas de comunicação de
dados, reunindo dados e fazendo avaliações para definição de desenhos
técnicos, plataformas, equipamentos, materiais, mão de obra e custos, em
função das capacidades previstas para a rede. Pode prestar assistência
técnica, orientando o cliente e intermediando sua relação com a empresa.
Realiza supervisão e treinamento de equipes de trabalho, avaliando
desempenho e propondo programas de aperfeiçoamento e atualização
profissional.

\begin{Form}
\textbf{Esta correspondência está adequada?} \\ 
\ChoiceMenu[radio,radiosymbol=\ding{52},bordercolor=0 0.5 0,name=cbovalid_ 313305 ]{}{CBO deve ficar=sim, \hspace{6em} CBO não fica aqui=nao}
\end{Form}

\textbf{CBO: 313310  ( Técnico de rede (telecomunicações) )}

Instala sistemas de redes de telecomunicações, certificando-se da
adequação da infraestrutura e do fornecimento de energia, fixando
componentes do sistema, estabelecendo a interligação dos equipamentos
via interfaces e cabeamento e ativando o sistema. Testa sistemas de
redes de telecomunicações, providenciando configuração e programação,
verificando se o funcionamento está de acordo com as especificações,
realizando medições, aferições e ajustes. Pode substituir componentes,
quando necessário. Orienta usuário final sobre a utilização do sistema.
Realiza manutenção preventiva e corretiva em equipamentos de redes de
telecomunicação, executando rotinas de teste, identificando e corrigindo
falhas, reprogramando e configurando programas do sistema e reparando
equipamentos. Participa da elaboração de projetos em redes de
telecomunicação, reunindo dados e fazendo avaliações para definição de
desenhos técnicos, plataformas, equipamentos, materiais, mão de obra e
custos, em função das capacidades previstas para a rede. Pode prestar
assistência técnica, orientando o cliente e intermediando sua relação
com a empresa. Realiza supervisão e treinamento de equipes de trabalho,
avaliando desempenho e propondo programas de aperfeiçoamento e
atualização profissional.

\begin{Form}
\textbf{Esta correspondência está adequada?} \\ 
\ChoiceMenu[radio,radiosymbol=\ding{52},bordercolor=0 0.5 0,name=cbovalid_ 313310 ]{}{CBO deve ficar=sim, \hspace{6em} CBO não fica aqui=nao}
\end{Form}

\textbf{CBO: 313315  ( Técnico de telecomunicações (telefonia) )}

Instala sistemas de telefonia, certificando-se da adequação da
infraestrutura e do fornecimento de energia, fixando componentes do
sistema, estabelecendo a interligação dos equipamentos via interfaces e
cabeamento e ativando o sistema. Testa sistemas de telefonia,
providenciando configuração e programação, verificando se o
funcionamento está de acordo com as especificações e realizando
medições, aferições e ajustes. Pode substituir componentes, quando
necessário. Orienta o usuário final sobre a utilização do sistema.
Realiza manutenção preventiva e corretiva em equipamentos de telefonia,
executando rotinas de teste, identificando e corrigindo falhas,
reprogramando e configurando programas do sistema e reparando
equipamentos. Participa da elaboração de projetos em telefonia, reunindo
dados e fazendo avaliações para definição de desenhos técnicos,
plataformas, equipamentos, materiais, mão de obra e custos, em função
das capacidades previstas para a rede. Pode prestar assistência técnica,
orientando o cliente e intermediando sua relação com a empresa. Realiza
supervisão e treinamento de equipes de trabalho, avaliando desempenho e
propondo programas de aperfeiçoamento e atualização profissional.

\begin{Form}
\textbf{Esta correspondência está adequada?} \\ 
\ChoiceMenu[radio,radiosymbol=\ding{52},bordercolor=0 0.5 0,name=cbovalid_ 313315 ]{}{CBO deve ficar=sim, \hspace{6em} CBO não fica aqui=nao}
\end{Form}

\textbf{CBO: 313320  ( Técnico de transmissão (telecomunicações) )}

Instala sistemas de transmissão nas telecomunicações, certificando-se da
adequação da infraestrutura e do fornecimento de energia, fixando
componentes do sistema, estabelecendo a interligação dos equipamentos
via interfaces e cabeamento e ativando o sistema. Testa sistemas de
transmissão nas telecomunicações, providenciando configuração e
programação, verificando se o funcionamento está de acordo com
especificações e realizando medições, aferições e ajustes. Pode fazer
substituição de componentes, quando necessário. Orienta o usuário final
sobre a utilização do sistema. Realiza manutenção preventiva e corretiva
em equipamentos de transmissão nas telecomunicações, executando rotinas
de teste, identificando e corrigindo falhas, reprogramando e
configurando programas do sistema e reparando equipamentos. Participa da
elaboração de projetos em sistemas de transmissão nas telecomunicações,
reunindo dados e fazendo avaliações para definição de desenhos técnicos,
plataformas, equipamentos, materiais, mão de obra e custos, em função
das capacidades previstas para a rede. Pode prestar assistência técnica,
orientando o cliente e intermediando sua relação com a empresa. Realiza
supervisão e treinamento de equipes de trabalho, avaliando desempenho e
propondo programas de aperfeiçoamento e atualização profissional.

\begin{Form}
\textbf{Esta correspondência está adequada?} \\ 
\ChoiceMenu[radio,radiosymbol=\ding{52},bordercolor=0 0.5 0,name=cbovalid_ 313320 ]{}{CBO deve ficar=sim, \hspace{6em} CBO não fica aqui=nao}
\end{Form}

\textbf{CBO: 731305  ( Instalador-reparador de equipamentos de comutação em telefonia )}

Planeja a instalação e a reparação de equipamentos de comutação em
telefonia, interpretando ordem de serviço e documentação técnica de
projetos, para selecionar materiais, ferramentas e equipamentos. Instala
equipamentos de comutação em telefonia, montando as ferragens de suporte
para os equipamentos e instalando o cabeamento periférico de comutação.
Monta os equipamentos, conecta seus cabos internos e liga o sistema de
proteção elétrica. Alimenta o sistema com energia elétrica, carrega o
seu software de operação e programa os equipamentos de comutação. Repara
equipamentos de comutação em telefonia, analisando seus circuitos
eletroeletrônicos e monitorando as centrais de comutação. Analisa causas
dos defeitos e substitui componentes, peças e módulos defeituosos,
manuseando ferramentas, instrumentos e equipamentos eletroeletrônicos de
medição, diagnóstico e reparo. Repara peças e módulos e ajusta
equipamentos de comutação. Encaminha componentes, peças ou módulos de
comutação para reparos em laboratório. Identifica falhas nos
entroncamentos de centrais analógicas e digitais. Analisa encaminhamento
de chamadas nas centrais de comutação. Fornece suporte técnico ao
usuário de equipamentos de comutação. Controla resultados de
funcionamento dos equipamentos de comutação em telefonia, manuseando
equipamentos de teste e elaborando planilhas de dados para controle.
Mantém ferramentas e equipamentos limpos, acondicionados e em plenas
condições de uso e funcionamento. Controla o desperdício de material.
Seleciona materiais reutilizáveis e faz o descarte de resíduos de acordo
com as normas ambientais. Protege o meio ambiente contra soluções
tóxicas de baterias. Trabalha com segurança, utilizando equipamentos de
proteção individual e equipamentos de proteção coletiva e prevenindo
acidentes. Protege-se em relação aos agentes tóxicos, gases e energia.

\begin{Form}
\textbf{Esta correspondência está adequada?} \\ 
\ChoiceMenu[radio,radiosymbol=\ding{52},bordercolor=0 0.5 0,name=cbovalid_ 731305 ]{}{CBO deve ficar=sim, \hspace{6em} CBO não fica aqui=nao}
\end{Form}

\newpage
\section{ 06001 -- Técnico em Aeroportuário }

\textbf{Perfil do Curso (CNCT)}

O Técnico Aeroportuário será habilitado para:

Inspecionar e garantir a manutenção de instalações e equipamentos.
Realizar a calibração e reparar sistemas aeroportuários quanto à
segurança e integridade de funcionamento. Executar plano de manutenção
do aeroporto. Controlar e coordenar o trânsito de pessoal e de viaturas
na área operacional. Supervisionar o carregamento e o descarregamento de
aeronaves. Auxiliar na operação de solo e sinalização de aeronaves.
Controlar o combustível de aviação e executar testes e abastecimento de
aeronaves. Verificar e monitorar o estado, o funcionamento e a
utilização de instalações e de unidades operacionais. Realizar
atividades relacionadas às empresas e aos usuários. Controlar os
serviços envolvidos e o acesso às salas de entrada restrita. Controlar a
manutenção e gestão de elevadores e de passarelas/escadas rolantes e
?fingers?. Analisar relatórios operacionais.

Para atuação como Técnico Aeroportuário, são fundamentais:

Conhecimentos e saberes relacionados às atividades de planejamento
associadas à instalação, à manutenção, à operação e à segurança de
sistemas e equipamentos aeroportuários, de acordo com a regulamentação
brasileira de aviação civil e de atividades de calibração. Além disso,
deve prezar pela ética e preservação do meio ambiente, ter espírito
inovador e empreendedor, e ser capaz de supervisionar equipes com o
intuito de solucionar problemas técnicos, trabalhistas e gestão de
conflitos.

\textbf{CBOs correspondidos e seus perfis (presentes no CNCT, e presentes na predição da IA)}

\textbf{CBO: 342535  ( Operador de atendimento aeroviário )}

Levanta informações de planejamento de voo sobre número de passageiros,
carga e bagagem, para definir o número de bilhetes que poderá ser
vendido. Atende clientes e usuários da aviação civil. Vende bilhetes,
faz reserva de voos e providencia acomodações. No processo de embarque,
confere bilhetes e documentação dos passageiros. Examina documento
específico para transporte de objetos atípicos. Etiqueta bagagem. Pode
chamar lista de espera. Insere dados no sistema informativo de voo
(SIV). Encaminha os passageiros à sala de embarque. Solicita descarte de
objetos proibidos e/ou perigosos. Pode solicitar consentimento de
passageiros para busca pessoal. Efetua chamada para embarque, de acordo
com prioridades. Pode acomodar bagagem fora do padrão ou excedente no
porão da aeronave. Informa à tripulação quantidade de pessoas a bordo.
Acondiciona, conforme legislação vigente, objetos especiais declarados.
Identifica passageiros faltantes no embarque, podendo providenciar a
retirada de sua bagagem. Autoriza e acompanha o fechamento das portas da
aeronave e dos porões. Na chegada ao destino, providencia o atendimento
de passageiros especiais -- tais como pessoas doentes e/ou com
dificuldades de locomoção -- para desembarque. Encaminha passageiros em
conexão para embarque no outro voo. Direciona passageiros para sala de
desembarque. Em caso de emergência médica, providencia atendimento.
Guarda bagagem desacompanhada, extraviada ou esquecida, conforme
procedimento padrão. Pode abrir processo de extravio ou danificação de
bagagem. Pode providenciar indenizações ao passageiro, em casos
previstos na legislação. Pode fornecer informações gerais e registrar
reclamações de passageiro. Elabora relatórios. Participa de reuniões.

\begin{Form}
\textbf{Esta correspondência está adequada?} \\ 
\ChoiceMenu[radio,radiosymbol=\ding{52},bordercolor=0 0.5 0,name=cbovalid_ 342535 ]{}{CBO deve ficar=sim, \hspace{6em} CBO não fica aqui=nao}
\end{Form}

\textbf{CBO: 342545  ( Supervisor de empresa aérea em aeroportos )}

Realiza o planejamento operacional das atividades da empresa aérea no
aeroporto, analisando informações, tais como condições meteorológicas,
número de voos previstos, volume de cargas a serem movimentadas e
quantidade de passageiros previstos para embarque e desembarque nas
aeronaves da empresa. Insere dados e faz pesquisas no sistema
informativo de voo (SIV). Coordena e supervisiona equipe de trabalho da
empresa no aeroporto, distribuindo tarefas, orientando a execução das
atividades, avaliando desempenho e realizando programas de treinamento.
Supervisiona o treinamento em serviço do pessoal de segurança,
programado a partir da avaliação de seu desempenho ou da introdução de
novas diretrizes e/ou novos procedimentos. Monitora e avalia o
desempenho de equipe de trabalho vinculada a empresas de serviços
auxiliares de transporte aéreo, contratadas pela empresa aérea.
Supervisiona a organização de cargas para embarque nas aeronaves da
empresa. Acompanha, por sistema informatizado, o check-in realizado pelo
passageiro na Internet ou nos totens de autoatendimento. Monitora os
procedimentos de check-in realizados pelos operadores de atendimento
aeroviário nas instalações do aeroporto. Confere documentação específica
para transporte de objetos atípicos. Monitora a implementação de
procedimentos de triagem de bagagens, guarda da bagagem despachada e
carregamento da aeronave. Supervisiona o processo de desembarque de
cargas e de passageiros das aeronaves da empresa no aeroporto,
solucionando problemas imprevistos. Registra e encaminha providências
para resolver casos de danificação de bagagem de passageiros e de
embalagens de carga. Recebe reclamações, providenciando respostas e/ou
decidindo sobre a forma de eliminar as falhas apontadas. Supervisiona as
atividades da empresa aérea relacionadas à segurança da aviação civil no
aeroporto. Aplica e acompanha procedimentos de segurança de acordo com
programa de prevenção de acidentes aeronáuticos (PPAA). Aciona serviço
de alerta e de operações de emergência para atendimento de aeronaves da
empresa, podendo solicitar isolamento de áreas do aeroporto. Providencia
atendimento de emergências médicas. Monitora transporte de artigos
perigosos e controlados. Notifica comandante sobre existência de
passageiros armados a bordo. Inspeciona áreas restritas de segurança,
identificando e recolhendo objetos que possam causar danos às aeronaves
e aos passageiros. Realiza os procedimentos de controle de qualidade
descritos pelo Programa de Segurança de Empresa Aérea (PSEA). Acompanha
a elaboração, o registro e o encaminhamento de relatórios de incidentes
e acidentes. Mantém contatos com outros órgãos e entidades envolvidos
com a segurança da aviação civil.

\begin{Form}
\textbf{Esta correspondência está adequada?} \\ 
\ChoiceMenu[radio,radiosymbol=\ding{52},bordercolor=0 0.5 0,name=cbovalid_ 342545 ]{}{CBO deve ficar=sim, \hspace{6em} CBO não fica aqui=nao}
\end{Form}

\newpage
\section{ 06002 -- Técnico em Agrimensura }

\textbf{Perfil do Curso (CNCT)}

O Técnico em Agrimensura será habilitado para:

Executar levantamentos geodésicos e topográficos. Utilizar equipamentos
e métodos específicos. Fazer a locação de obras de sistemas de
transporte, civis, industriais e rurais. Delimitar glebas. Identificar
elementos na superfície e pontos de apoio para georreferenciamento e
amarração. Organizar e supervisionar ações de levantamento e mapeamento.
Efetuar aerotriangulação. Restituir fotografias aéreas para a elaboração
de produtos cartográficos em diferentes sistemas de referências e
projeções. Processar e interpretar dados de sensoriamento remoto, fotos
terrestres e fotos aéreas de modo integrado a dados de cartas, mapas e
plantas. Utilizar ferramentas de geoprocessamento. Executar cadastro
técnico multifinalitário. Identificar métodos e equipamentos para a
coleta de dados. Participar do planejamento de loteamentos,
desmembramentos e obras de engenharia. Dar assistência técnica na
compra, venda e utilização de produtos e equipamentos especializados.
Executar levantamentos e coletas de dados espaciais e geométricos

Para atuação como Técnico em Agrimensura, são fundamentais:

Conhecimentos e saberes relacionados à execução de levantamentos
geodésicos e topográficos, a vistorias e arbitramentos relativos à
Agrimensura, com o intuito de permitir a organização fundiária do espaço
rural, incluindo as medições, as demarcações, as divisões, os
mapeamentos, as avaliações e a regulamentação das terras. Compromisso e
ética para assegurar o cumprimento da legislação e das normas técnicas
vigentes. Habilidade de liderança de equipes para solução de problemas
técnicos e trabalhistas e para a gestão de conflitos.

\textbf{CBOs correspondidos e seus perfis (presentes no CNCT, e presentes na predição da IA)}

\textbf{CBO: 312305  ( Técnico em agrimensura )}

Planeja as atividades em agrimensura, definindo escopo e logística,
dimensionando equipes de trabalho, especificando e quantificando
equipamentos, acessórios e materiais. Elabora planilha de custos e
cronograma físico-financeiro. Pode utilizar software de geoprocessamento
e operar softwares de automação topográfica, tais como SNGS-Sistemas de
Navegação Global por Satélite (GNSS-Global Navigation Satellite Systems)
e SIG - Sistema de Informações Geográficas (GIS-Geography Information
System). Pode usar drones (aeronaves remotamente pilotadas), para obter
fotografias aéreas com digitalização em tempo real. Coleta dados para
atualização de plantas cadastrais urbanas e rurais, elaborando
representações gráficas e planta topográfica, em conformidade com normas
da Associação Brasileira de Normas Técnicas (ABNT). Define limites e
confrontações georreferenciadas, aviventando e declinando rumos e
azimutes magnéticos para estabelecimento de relações geométricas em
mapeamentos. Materializa marcos e pontos topográficos, delimitando
glebas e fixando obras de sistema de transporte, linhas de transmissão,
obras civis, industriais e rurais e parcelamento de solos. Elabora
projetos de terraplenagem de pequeno porte, calculando declinação
magnética, convergência meridiana, norte verdadeiro, áreas de terrenos,
volumes para movimento de solo, distâncias, azimutes e coordenadas,
concordâncias vertical e horizontal, curvas de nível por interpolação,
offset e greide, tendo em vista cotas de projeto. Executa levantamentos
geodésicos e topo-hidrográficos, medindo ângulos e distâncias e fazendo
levantamentos altimétricos e planimétricos. Realiza topografias
especiais (industriais, subterrâneas e batimétricas). Realiza operações
geodésicas. Determina coordenadas geográficas e plano-retangulares pelo
Sistema Universal Transverso de Mercator (UTM). Determina norte
verdadeiro e norte magnético, tendo em vista o georreferenciamento mais
adequado para o projeto. Demarca e mede áreas em campo, elaborando
croqui de campo e relatório de demarcação para levantamento cadastral.
Analisa documentos e informações cartográficas, coletando dados
geométricos; interpretando fotos terrestres e aéreas; interpretando
mapas, cartas e plantas; interpretando relevos para implantação de
linhas de exploração; e identificando acidentes geométricos e pontos de
apoio para georreferenciamento e amarração. Elabora documentos
cartográficos, definindo sistema de projeção, escalas e cálculos
cartográficos. Efetua aerotriangulação. Revisa documentos cartográficos,
restituindo e reambulando fotografia aérea. Edita e cria arte final de
documentos cartográficos. Auxilia agrimensor de nível superior na
organização e na supervisão de equipes de trabalho, orientando e
acompanhando o cumprimento de tarefas e realizando programas de
treinamento. Conserva o local das atividades limpo e organizado,
providenciando o acondicionamento de materiais. Mantém equipamentos e
acessórios limpos, organizados e em plenas condições de uso e
funcionamento. Destina adequadamente os resíduos gerados, orientando a
reciclagem de materiais e providenciando descarte de acordo com as
normas ambientais. Zela pela segurança, fiscalizando o uso de
equipamentos de proteção individual e coletiva e atuando na prevenção de
acidentes.

\begin{Form}
\textbf{Esta correspondência está adequada?} \\ 
\ChoiceMenu[radio,radiosymbol=\ding{52},bordercolor=0 0.5 0,name=cbovalid_ 312305 ]{}{CBO deve ficar=sim, \hspace{6em} CBO não fica aqui=nao}
\end{Form}

\textbf{CBO: 312310  ( Técnico em geodesia e cartografia )}

Planeja as atividades em geodésia e cartografia, definido escopo,
metodologia e logística; dimensionando equipes de trabalho;
especificando e quantificando equipamentos, acessórios e materiais.
Seleciona softwares para as atividades, utilizando software de
geoprocessamento e podendo operar softwares de automação topográfica,
tais como SNGS -Sistemas de Navegação Global por Satélite (GNSS-Global
Navigation Satellite Systems) e SIG-Sistema de Informações Geográficas
(GIS-Geography Information System). Pode usar drones (aeronaves
remotamente pilotadas), para obter fotografias aéreas com digitalização
em tempo real. Elabora planilha de custos e cronograma
físico-financeiro. Executa levantamentos geodésicos, medindo ângulos e
distâncias e realizando levantamentos altimétricos e planimétricos.
Determina norte verdadeiro e norte magnético. Determina coordenadas
geográficas e plano-retangulares, pelo Sistema Universal Transverso de
Mercator (UTM). Elabora relatório de demarcação para levantamento
cadastral. Coleta dados para atualização de plantas cadastrais,
efetuando cálculos de áreas de terrenos. Calcula volumes para movimento
de solo e curvas de nível por interpolação. Calcula declinação
magnética, convergência meridiana, norte verdadeiro, distâncias,
azimutes e coordenadas. Elabora representações gráficas, fazendo planta
topográfica, conforme normas da Associação Brasileira de Normas Técnicas
(ABNT). Analisa documentos e informações cartográficas, interpretando
fotos terrestres, fotos aéreas e imagens orbitais; manipulando mapas
analógicos e digitais; e processando dados de campo e dados coletados
por sensores orbitais e aéreos, para produção cartográfica. Elabora
documento cartográfico, definindo sistema de projeção, escalas e
cálculos cartográficos. Efetua aerotriangulação. Cria base cartográfica,
com base na análise de documentos e no processamento das informações
cartográficas. Edita e elabora a arte final do documento cartográfico.
Revisa documentos cartográficos, restituindo e fazendo reambulação de
fotografias aéreas. Auxilia agrimensor, engenheiro cartográfico ou outro
profissional de nível superior na organização e na supervisão das
equipes de trabalho de campo, orientando e acompanhando o cumprimento de
tarefas e realizando programas de treinamento. Conserva o local das
atividades limpo e organizado, providenciando o acondicionamento de
materiais. Mantém equipamentos e acessórios limpos, organizados e em
plenas condições de uso e funcionamento. Destina adequadamente os
resíduos gerados, orientando a reciclagem de materiais e providenciando
descarte de acordo com as normas ambientais. Zela pela segurança,
atuando na prevenção de acidentes e fiscalizando o uso de equipamentos
de proteção individual e coletiva.

\begin{Form}
\textbf{Esta correspondência está adequada?} \\ 
\ChoiceMenu[radio,radiosymbol=\ding{52},bordercolor=0 0.5 0,name=cbovalid_ 312310 ]{}{CBO deve ficar=sim, \hspace{6em} CBO não fica aqui=nao}
\end{Form}

\textbf{CBO: 312315  ( Técnico em hidrografia )}

Planeja as atividades em hidrografia, definindo escopo, metodologia e
logística e especificando equipamentos, acessórios e materiais.
Seleciona softwares para as atividades, operando sistemas de
geoprocessamento e podendo operar softwares de automação topográfica,
tais como SNGS-Sistemas de Navegação Global por Satélite (GNSS-Global
Navigation Satellite Systems) e SIG-Sistema de Informações Geográficas
(GIS-Geography Information System). Executa levantamentos geodésicos e
topo-hidrográficos, medindo ângulos e distâncias e realizando
levantamentos altimétricos e planimétricos. Determina coordenadas
geográficas e plano-retangulares, pelo Sistema Universal Transverso de
Mercator (UTM). Determina norte verdadeiro e norte magnético, tendo em
vista o georreferenciamento mais adequado para o projeto hidrográfico.
Realiza cálculos batimétricos e cálculos de marés. Mede e demarca
superfícies de águas fluviais, lacustres, marinhas e oceânicas, para
elaborar plantas e mapas de canais de navegação, fornecendo elementos
básicos para construção ou reparo de portos, cais, açudes e outras obras
marítimas, fluviais ou lacustres. Coleta dados meteorológicos,
maregráficos, oceanográficos, hidrográficos, por meio da operação de
equipamentos específicos. Interpreta boletins meteorológicos. Identifica
astros e elementos que determinam sua posição na esfera celeste.
Determina a posição de um navio por meio de métodos específicos.
Confecciona cartas náuticas oceânicas e de vias navegáveis interiores.
Elabora documentos cartográficos, definindo sistema de projeção, escalas
e cálculos cartográficos. Cria a base cartográfica e edita os
documentos. Calibra, opera e mantém, sob supervisão, equipamentos de
levantamento hidrográfico, incluindo sistemas de sonar de feixe único,
feixe múltiplo e de varredura lateral; instrumentos de posicionamento;
sensores de movimento, de velocidade do som e equipamentos auxiliares,
como dispositivos de amostragem de sedimentos e sistemas hidráulicos.
Zela pela segurança, atuando na prevenção de acidentes. Destina
adequadamente os resíduos gerados, orientando a reciclagem de materiais
e providenciando descarte de acordo com as normas ambientais.

\begin{Form}
\textbf{Esta correspondência está adequada?} \\ 
\ChoiceMenu[radio,radiosymbol=\ding{52},bordercolor=0 0.5 0,name=cbovalid_ 312315 ]{}{CBO deve ficar=sim, \hspace{6em} CBO não fica aqui=nao}
\end{Form}

\textbf{CBO: 312320  ( Topografo )}

Planeja atividades em topografia, definindo escopo, metodologia e
logística; especificando e quantificando equipamentos e instrumentos de
trabalho; e dimensionando equipe de auxiliares. Seleciona softwares para
as atividades, utilizando software de georreferenciamento e operando
softwares de automação topográfica, tais como SNGS -Sistemas de
Navegação Global por Satélite (GNSS-Global Navigation Satellite Systems)
e SIG-Sistema de Informações Geográficas (GIS-Geography Information
System). Elabora planilha de custos e cronograma físico-financeiro.
Prepara levantamentos topográficos planimétricos e altimétricos,
estudando documentos técnicos; analisando mapas, plantas, títulos de
propriedade, registros e especificações; e calculando as medições a
serem efetuadas. Analisa as características de terrenos para
reconhecimento básico da área programada, identificando pontos de
partida e vias de melhor acesso. Realiza os levantamentos topográficos
das áreas demarcadas em projetos de construção e de exploração de solos.
Determina altitudes, distâncias, ângulos, coordenadas, referências de
níveis e outras características de superfície terrestre, de áreas
subterrâneas e de edifícios, utilizando instrumentos e equipamentos,
tais como teodolitos, níveis, trenas, bússolas e telêmetros. Registra os
dados topográficos levantados, apontando valores identificados e
cálculos numéricos efetuados, tendo em vista o georreferenciamento para
projetos e obras de construção e exploração de solos. Confere e verifica
a precisão de informações e dados coligidos, consultando tabelas,
efetuando cálculos, avaliando as diferenças entre distâncias e
altitudes, aplicando fórmulas de altimetria e declividade para
nivelamentos planimétricos. Prepara documentos cartográficos, elaborando
esboços, plantas e relatórios técnicos, indicando traçados, pontos e
convenções para desenvolvimento de mapas, cartas e projetos. Pode
desenhar plantas detalhadas das áreas levantadas. Auxilia agrimensor ou
outro profissional de nível superior da área na orientação e na
supervisão de auxiliares, especificando as tarefas a serem realizadas,
detalhando o modo de execução e apontando o grau de precisão a ser
observado. Realiza aferição dos instrumentos de levantamentos
topográficos, responsabilizando-se por sua manutenção, guarda e
conservação. Destina adequadamente os resíduos gerados, orientando a
reciclagem de materiais e providenciando descarte de acordo com as
normas ambientais Zela pela segurança, fiscalizando o uso de
equipamentos de proteção individual e coletiva e atuando na prevenção de
acidentes.

\begin{Form}
\textbf{Esta correspondência está adequada?} \\ 
\ChoiceMenu[radio,radiosymbol=\ding{52},bordercolor=0 0.5 0,name=cbovalid_ 312320 ]{}{CBO deve ficar=sim, \hspace{6em} CBO não fica aqui=nao}
\end{Form}

\textbf{CBO: 318110  ( Desenhista técnico (cartografia) )}

Planeja o trabalho relacionado a desenho técnico de cartografia,
propondo alternativas para elaboração do desenho e determinando a
metodologia para sua execução. Seleciona instrumentos e outros meios
necessários à elaboração do desenho. Pode prever uso de sistemas BIM -
Modelagem da Informação da Construção (Building Information Modeling).
Calcula e define o custo do desenho. Estipula prazo de execução. Prepara
o local de trabalho, de acordo com recomendações normativas de
ergonomia, de segurança do trabalho e de saúde ocupacional. Coleta dados
para elaboração do desenho técnico de cartografia, interpretando
desenhos cartográficos existentes -- analógicos ou digitais -- e
consultando informações em arquivos. Interpreta fotos terrestres, fotos
aéreas e imagens orbitais. Utiliza GIS-Sistema de Informações
Geográficas (Geographic Information System) e GPS-Sistema de
Posicionamento Global (Global Positioning System) para coleta de dados.
Pode obter dados por meio de escaneamento a laser ou com uso de drones
em tempo real. Busca informações complementares e organiza dados
coletados. Processa dados para desenho técnico de cartografia,
interpretando informações do conjunto de dados coletados. Realiza
pesquisas complementares -- inclusive via internet --, para obter novos
dados relacionados ao desenho ou ampliar a compreensão sobre o conjunto
de dados à disposição. Processa dados de campo e dados coletados por
sensores orbitais e aéreos para produção cartográfica. Elabora desenhos
técnicos de cartografia -- tais como mapas, plantas e cartas, dentre
outros --, aplicando normas técnicas e realizando visitas técnicas para
rever dados. Utiliza softwares específicos para desenho cartográfico.
Pode utilizar CAD-Desenho Assistido por Computador (Computer Aided
Design), com o software funcionando de forma independente (stand-alone)
e/ou integrado ao GIS-Sistema de Informações Geográficas (Geographic
Information System) ou similar. Utiliza estações de trabalho
computadorizadas, impressoras, plotters e outros periféricos. Pode usar
instrumentos para execução do traçado diretamente em substratos de
papel, quando requerido. Define formatos e escalas e detalha desenhos.
Diagrama pranchas ou folhas de desenho técnico. Legenda plantas e
elabora lista de material relacionada ao conteúdo do desenho. Revisa
desenhos técnicos de cartografia, conferindo escalas, dimensionamentos,
informações descritivas e outros conteúdos. Atualiza desenhos, em função
de mudanças de legislação. Realiza complementação ou modificação de
desenhos, quando necessário. Em situações em que os desenhos passam por
supervisão, realiza as compatibilizações necessárias dos desenhos,
conforme os apontamentos da supervisão. Nas etapas de trabalho, aplica
convenções internacionais para cartografia. Fecha a ordem de serviço do
desenho de cartografia, relatando mudanças de procedimento adotadas ao
longo dos trabalhos; liberando o desenho para arquivo eletrônico ou
mapoteca; enviando originais e/ou cópias do desenho para clientes.
Retorna documentos utilizados para arquivo. Organiza arquivos técnicos
-- físicos e eletrônicos -- relacionados aos serviços com desenhos
técnicos de cartografia, definindo o tipo de arquivo a ser usado e
indexando documentos pertinentes à área. Reúne documentos e organiza
catálogos de fornecedores. Compacta e armazena arquivos digitais.
Conserva o local de trabalho limpo e organizado. Controla luminosidade
do ambiente e verifica condições ergonômicas para uso de monitor de
vídeo, mouse, mesa digitalizadora, prancheta e teclado. Mantém
equipamentos e instrumentos de trabalho em plenas condições de uso e
funcionamento. Requisita serviço de manutenção de computador,
impressoras, plotters e outros periféricos, quando necessário. Destina
adequadamente os resíduos gerados, fazendo reciclagem e executando
descarte de acordo com normas ambientais. Zela pela segurança, atuando
na prevenção de acidentes e usando equipamentos de proteção individual,
sempre que necessário.

\begin{Form}
\textbf{Esta correspondência está adequada?} \\ 
\ChoiceMenu[radio,radiosymbol=\ding{52},bordercolor=0 0.5 0,name=cbovalid_ 318110 ]{}{CBO deve ficar=sim, \hspace{6em} CBO não fica aqui=nao}
\end{Form}

\newpage
\section{ 06004 -- Técnico em Desenho de Construção Civil }

\textbf{Perfil do Curso (CNCT)}

O Técnico em Desenho em Construção Civil será habilitado para:

Elaborar desenhos e detalhamentos de construções prediais, estradas,
obras de saneamento, estruturas, instalações (hidráulicas, elétricas,
telefônicas, de gás liquefeito de petróleo, de ar-condicionado,
preventivas de incêndio) e redes (de esgoto, águas pluviais e de
abastecimento de água), em meio analógico ou digital. Coletar e
processar dados. Planejar a elaboração do projeto. Calcular e definir
custos de desenho. Analisar croquis. Elaborar maquetes.

Para atuação como Técnico em Desenho de Construção Civil, são
fundamentais:

Conhecimentos e saberes relacionados ao planejamento, às técnicas e
processos de produção na construção civil, no desenvolvimento de
projeto, utilizando expressões gráficas para a planificação de sólidos
existentes na construção civil, predial ou de infraestrutura, através de
softwares específicos para desenho digital ou em forma analógica,
respeitando às normas técnicas vigentes. Além disso, deve prezar pela
ética, viabilidade técnico-econômica e preservação do meio ambiente, ter
espírito inovador e empreendedor, e ser capaz de supervisionar equipes
com o intuito de solucionar problemas técnicos e à gestão de conflitos.

\textbf{CBOs correspondidos e seus perfis (presentes no CNCT, e presentes na predição da IA)}

\textbf{CBO: 318005  ( Desenhista técnico )}

Analisa solicitação de desenho técnico, interpretando documentos de
apoio, tais como plantas, projetos, catálogos, croquis e normas. Propõe
ao solicitante alternativas para execução do desenho, chegando a um
acordo sobre os detalhes técnicos finais. Estima tempo para realização
do desenho e define prioridades, conforme cronograma. Programa ações
para realização do desenho técnico, observando características técnicas,
fazendo esboços, estabelecendo formatos e definindo escalas e sistemas
de representação. Elabora desenhos técnicos, representando formas
geométricas em duas e três dimensões. Utiliza computador, programa --
tal como CAD-Desenho Assistido por Computador (Computer Aided Design) -,
impressoras, plotters e outros periféricos, além de fazer uso de
instrumentos para execução do traçado diretamente em substratos de
papel. Constrói o desenho, traçando linhas auxiliares, cotando e
definindo a lista de materiais e componentes. Envia o desenho para
revisão, recebendo a aprovação do solicitante. Realiza cópias de
segurança e disponibiliza o desenho técnico final e/ou suas revisões
para as áreas afins, que providenciarão a entrega ao solicitante.
Conserva o local de trabalho limpo e organizado. Controla graus de
luminosidade do ambiente e verifica as condições ergonômicas para uso de
monitor de vídeo, mouse, mesa digitalizadora, prancheta e teclado.
Mantém equipamentos e instrumentos de trabalho em plenas condições de
uso e funcionamento. Requisita serviço de manutenção de computador,
impressoras, plotters e outros periféricos, quando necessário. Organiza
e arquiva documentos físicos e eletrônicos, tais como comprovantes de
aprovação e cópias de segurança dos desenhos técnicos. Destina
adequadamente os resíduos gerados, fazendo reciclagem e realizando
descarte de acordo com as normas ambientais.

\begin{Form}
\textbf{Esta correspondência está adequada?} \\ 
\ChoiceMenu[radio,radiosymbol=\ding{52},bordercolor=0 0.5 0,name=cbovalid_ 318005 ]{}{CBO deve ficar=sim, \hspace{6em} CBO não fica aqui=nao}
\end{Form}

\textbf{CBO: 318010  ( Desenhista copista )}

Analisa solicitação de cópia de desenho técnico, observando as
características técnicas do desenho original e interpretando documentos
de apoio, tais como originais de plantas, projetos, catálogos, croquis e
normas. Propõe ao solicitante alternativas para execução da cópia,
chegando a acordo em relação aos detalhes técnicos finais. Estima tempo
para realização do trabalho. Elabora cópias de desenhos técnicos,
reproduzindo representações de formas geométricas em duas e três
dimensões. Utiliza computador, programa -- tal como CAD-Desenho
Assistido por Computador (Computer Aided Design) -, mesas
digitalizadoras, impressoras, plotters e outros periféricos, além de
fazer uso de instrumentos para execução do traçado diretamente em
substratos de papel. Constrói a cópia do desenho, traçando linhas
auxiliares e cotando. Envia a cópia para revisão, recebendo a aprovação
do solicitante. Realiza cópias de segurança e disponibiliza o desenho
final copiado e/ou suas revisões para as áreas afins, que providenciarão
a entrega ao solicitante. Conserva o local de trabalho limpo e
organizado. Controla graus de luminosidade do ambiente e verifica as
condições ergonômicas para uso de monitor de vídeo, mouse, mesa
digitalizadora, prancheta e teclado. Mantém equipamentos e instrumentos
de trabalho em plenas condições de uso e funcionamento. Requisita
serviço de manutenção de computador, impressoras, plotters e outros
periféricos, quando necessário. Organiza e arquiva documentos físicos e
eletrônicos, tais como comprovantes de aprovação e uma via de cada
desenho copiado. Destina adequadamente os resíduos gerados, fazendo
reciclagem e realizando descarte de acordo com as normas ambientais.

\begin{Form}
\textbf{Esta correspondência está adequada?} \\ 
\ChoiceMenu[radio,radiosymbol=\ding{52},bordercolor=0 0.5 0,name=cbovalid_ 318010 ]{}{CBO deve ficar=sim, \hspace{6em} CBO não fica aqui=nao}
\end{Form}

\textbf{CBO: 318015  ( Desenhista detalhista )}

Analisa solicitação de detalhamento de desenho técnico, interpretando
original e documentos de apoio, tais como plantas, projetos, catálogos,
croquis e normas. Propõe ao solicitante alternativas para execução do
detalhamento, chegando a um acordo sobre pormenores técnicos finais.
Estima tempo para realização do detalhamento e estabelece prioridades,
conforme cronograma. Programa ações para realização do desenho
detalhado, observando características técnicas, fazendo esboços,
estabelecendo formatos e definindo escalas e sistemas de representação.
Elabora desenhos técnicos detalhados, representando formas geométricas
em duas e três dimensões. Utiliza computador, programa -- tal como
CAD-Desenho Assistido por Computador (Computer Aided Design) -,
impressoras, plotters e outros periféricos, além de fazer uso de
instrumentos para execução do traçado diretamente em substratos de
papel. Constrói o detalhamento do desenho, traçando linhas auxiliares,
cotando e realçando aspectos e características de formas geométricas
complexas e outros elementos a detalhar. Elabora a lista de materiais e
componentes do desenho detalhado. Envia o desenho técnico detalhado para
revisão, recebendo a aprovação do solicitante. Realiza cópias de
segurança e disponibiliza o desenho detalhado final e/ou suas revisões
para as áreas afins, que providenciarão a entrega ao solicitante.
Conserva o local de trabalho limpo e organizado. Controla graus de
luminosidade do ambiente e verifica as condições ergonômicas para uso de
monitor de vídeo, mouse, mesa digitalizadora, prancheta e teclado.
Mantém equipamentos e instrumentos de trabalho em plenas condições de
uso e funcionamento. Requisita serviço de manutenção de computador,
impressoras, plotters e outros periféricos, quando necessário. Organiza
e arquiva documentos físicos e eletrônicos, tais como comprovantes de
aprovação e cópias de segurança dos desenhos técnicos detalhados.
Destina adequadamente os resíduos gerados, fazendo reciclagem e
realizando descarte de acordo com as normas ambientais.

\begin{Form}
\textbf{Esta correspondência está adequada?} \\ 
\ChoiceMenu[radio,radiosymbol=\ding{52},bordercolor=0 0.5 0,name=cbovalid_ 318015 ]{}{CBO deve ficar=sim, \hspace{6em} CBO não fica aqui=nao}
\end{Form}

\textbf{CBO: 318105  ( Desenhista técnico (arquitetura) )}

Planeja o trabalho relacionado ao desenho técnico de arquitetura,
propondo alternativas para elaboração do desenho e determinando a
metodologia para sua execução. Seleciona instrumentos e outros meios
necessários à elaboração do desenho. Pode prever o uso de sistemas BIM -
Modelagem da Informação da Construção (Building Information Modeling).
Calcula e define o custo do desenho. Estipula prazo de execução. Prepara
o local de trabalho, de acordo com recomendações normativas de
ergonomia, de segurança do trabalho e de saúde ocupacional. Coleta dados
para elaboração do desenho técnico de arquitetura, interpretando
projetos existentes, consultando informações em arquivos e fazendo
levantamento de campo. Busca informações complementares e organiza os
dados coletados. Processa dados para desenho técnico de arquitetura,
analisando croquis obtidos por meio dos levantamentos de campo;
interpretando memória de cálculo e outras informações do conjunto de
dados coletados. Realiza pesquisas complementares -- inclusive via
internet --, para obter novos dados relacionados ao desenho ou ampliar a
compreensão sobre o conjunto de dados à disposição. Elabora desenhos
técnicos de arquitetura -- tais como plantas de apresentação e
anteprojeto; planta legal ou de licenciamento e planta humanizada,
dentre outros tipos de desenho arquitetônico --, aplicando normas
técnicas e realizando visitas técnicas para rever dados. Utiliza
softwares específicos para desenho, especialmente CAD-Desenho Assistido
por Computador (Computer Aided Design). Utiliza estações de trabalho
computadorizadas, impressoras, plotters e outros periféricos. Pode fazer
uso de instrumentos para execução do traçado diretamente em substratos
de papel, quando requerido. Define formatos e escalas e detalha
desenhos. Diagrama pranchas ou folhas de desenho técnico. Legenda
plantas e elabora a lista de material relacionada ao conteúdo do
desenho. Pode elaborar desenhos de plantas de execução. Revisa desenhos
técnicos de arquitetura, conferindo cotas, dimensionamentos, informações
descritivas e outros conteúdos. Atualiza desenhos, em função de mudanças
de legislação. Realiza a complementação ou a modificação de desenhos,
quando necessário. Em situações em que os desenhos passam por
supervisão, realiza as compatibilizações necessárias dos desenhos, de
acordo com os apontamentos da supervisão. Fecha a ordem de serviço do
desenho de arquitetura, relatando mudanças de procedimento adotadas ao
longo dos trabalhos; liberando o desenho para arquivo eletrônico ou
mapoteca; e enviando originais e/ou cópias do desenho para clientes.
Retorna documentos utilizados nos trabalhos para arquivo. Elabora
relatório final. Organiza arquivos técnicos -- físicos e eletrônicos --
relacionados aos serviços com desenhos técnicos de arquitetura,
determinando o tipo de arquivo a ser utilizado e indexando documentos
pertinentes à área. Reúne documentos e organiza catálogos de
fornecedores. Compacta e armazena arquivos digitais. Conserva o local de
trabalho limpo e organizado. Controla luminosidade do ambiente e
verifica as condições ergonômicas para uso de monitor de vídeo, mouse,
mesa digitalizadora, prancheta e teclado. Mantém equipamentos e
instrumentos de trabalho em plenas condições de uso e funcionamento.
Requisita serviço de manutenção de computador, impressoras, plotters e
outros periféricos, quando necessário. Destina adequadamente os resíduos
gerados, fazendo reciclagem e realizando descarte de acordo com as
normas ambientais. Zela pela segurança, atuando na prevenção de
acidentes e utilizando equipamentos de proteção individual, sempre que
necessário, especialmente durante os levantamentos de dados em campo.

\begin{Form}
\textbf{Esta correspondência está adequada?} \\ 
\ChoiceMenu[radio,radiosymbol=\ding{52},bordercolor=0 0.5 0,name=cbovalid_ 318105 ]{}{CBO deve ficar=sim, \hspace{6em} CBO não fica aqui=nao}
\end{Form}

\textbf{CBO: 318110  ( Desenhista técnico (cartografia) )}

Planeja o trabalho relacionado a desenho técnico de cartografia,
propondo alternativas para elaboração do desenho e determinando a
metodologia para sua execução. Seleciona instrumentos e outros meios
necessários à elaboração do desenho. Pode prever uso de sistemas BIM -
Modelagem da Informação da Construção (Building Information Modeling).
Calcula e define o custo do desenho. Estipula prazo de execução. Prepara
o local de trabalho, de acordo com recomendações normativas de
ergonomia, de segurança do trabalho e de saúde ocupacional. Coleta dados
para elaboração do desenho técnico de cartografia, interpretando
desenhos cartográficos existentes -- analógicos ou digitais -- e
consultando informações em arquivos. Interpreta fotos terrestres, fotos
aéreas e imagens orbitais. Utiliza GIS-Sistema de Informações
Geográficas (Geographic Information System) e GPS-Sistema de
Posicionamento Global (Global Positioning System) para coleta de dados.
Pode obter dados por meio de escaneamento a laser ou com uso de drones
em tempo real. Busca informações complementares e organiza dados
coletados. Processa dados para desenho técnico de cartografia,
interpretando informações do conjunto de dados coletados. Realiza
pesquisas complementares -- inclusive via internet --, para obter novos
dados relacionados ao desenho ou ampliar a compreensão sobre o conjunto
de dados à disposição. Processa dados de campo e dados coletados por
sensores orbitais e aéreos para produção cartográfica. Elabora desenhos
técnicos de cartografia -- tais como mapas, plantas e cartas, dentre
outros --, aplicando normas técnicas e realizando visitas técnicas para
rever dados. Utiliza softwares específicos para desenho cartográfico.
Pode utilizar CAD-Desenho Assistido por Computador (Computer Aided
Design), com o software funcionando de forma independente (stand-alone)
e/ou integrado ao GIS-Sistema de Informações Geográficas (Geographic
Information System) ou similar. Utiliza estações de trabalho
computadorizadas, impressoras, plotters e outros periféricos. Pode usar
instrumentos para execução do traçado diretamente em substratos de
papel, quando requerido. Define formatos e escalas e detalha desenhos.
Diagrama pranchas ou folhas de desenho técnico. Legenda plantas e
elabora lista de material relacionada ao conteúdo do desenho. Revisa
desenhos técnicos de cartografia, conferindo escalas, dimensionamentos,
informações descritivas e outros conteúdos. Atualiza desenhos, em função
de mudanças de legislação. Realiza complementação ou modificação de
desenhos, quando necessário. Em situações em que os desenhos passam por
supervisão, realiza as compatibilizações necessárias dos desenhos,
conforme os apontamentos da supervisão. Nas etapas de trabalho, aplica
convenções internacionais para cartografia. Fecha a ordem de serviço do
desenho de cartografia, relatando mudanças de procedimento adotadas ao
longo dos trabalhos; liberando o desenho para arquivo eletrônico ou
mapoteca; enviando originais e/ou cópias do desenho para clientes.
Retorna documentos utilizados para arquivo. Organiza arquivos técnicos
-- físicos e eletrônicos -- relacionados aos serviços com desenhos
técnicos de cartografia, definindo o tipo de arquivo a ser usado e
indexando documentos pertinentes à área. Reúne documentos e organiza
catálogos de fornecedores. Compacta e armazena arquivos digitais.
Conserva o local de trabalho limpo e organizado. Controla luminosidade
do ambiente e verifica condições ergonômicas para uso de monitor de
vídeo, mouse, mesa digitalizadora, prancheta e teclado. Mantém
equipamentos e instrumentos de trabalho em plenas condições de uso e
funcionamento. Requisita serviço de manutenção de computador,
impressoras, plotters e outros periféricos, quando necessário. Destina
adequadamente os resíduos gerados, fazendo reciclagem e executando
descarte de acordo com normas ambientais. Zela pela segurança, atuando
na prevenção de acidentes e usando equipamentos de proteção individual,
sempre que necessário.

\begin{Form}
\textbf{Esta correspondência está adequada?} \\ 
\ChoiceMenu[radio,radiosymbol=\ding{52},bordercolor=0 0.5 0,name=cbovalid_ 318110 ]{}{CBO deve ficar=sim, \hspace{6em} CBO não fica aqui=nao}
\end{Form}

\textbf{CBO: 318115  ( Desenhista técnico (construção civil) )}

Planeja o trabalho relacionado ao desenho técnico de construção civil,
propondo alternativas para elaboração do desenho e determinando a
metodologia para sua execução. Seleciona instrumentos e outros meios
necessários à elaboração do desenho. Pode prever uso de sistemas BIM -
Modelagem da Informação da Construção (Building Information Modeling).
Calcula e define o custo do desenho. Estipula prazo de execução. Prepara
o local de trabalho, de acordo com recomendações normativas de
ergonomia, de segurança do trabalho e de saúde ocupacional. Coleta dados
para elaboração do desenho técnico de construção civil, interpretando
projetos existentes, consultando informações em arquivos e fazendo
levantamento de campo. Busca informações complementares e organiza os
dados coletados. Processa dados para desenho técnico de construção
civil, analisando croquis obtidos por meio dos levantamentos de campo; e
interpretando informações do conjunto de dados coletados. Realiza
pesquisas complementares -- inclusive via internet --, para obter novos
dados relacionados ao desenho ou ampliar a compreensão sobre o conjunto
de dados à disposição. Elabora desenhos técnicos de construção civil --
tais como plantas de construções prediais, estradas e estruturas;
desenhos de instalações elétricas, de ar-condicionado e do sistema de
prevenção contra incêndios, dentre outros --, aplicando normas técnicas
e realizando visitas técnicas para rever dados. Utiliza softwares
específicos para desenho, especialmente CAD-Desenho Assistido por
Computador (Computer Aided Design). Utiliza estações de trabalho
computadorizadas, impressoras, plotters e outros periféricos. Pode fazer
uso de instrumentos para execução do traçado diretamente em substratos
de papel, quando requerido. Define formatos e escalas e detalha
desenhos. Diagrama pranchas ou folhas de desenho técnico. Legenda
plantas e elabora a lista de material relacionada ao conteúdo do desenho
técnico. Revisa desenhos técnicos de construção civil, conferindo cotas,
dimensionamentos, informações descritivas e outros conteúdos. Atualiza
desenhos, em função de mudanças de legislação. Realiza a complementação
ou a modificação de desenhos, quando necessário. Em situações em que os
desenhos passam por supervisão, realiza as compatibilizações necessárias
dos desenhos, de acordo com os apontamentos da supervisão. Fecha a ordem
de serviço do desenho de construção civil, relatando mudanças de
procedimento adotadas ao longo dos trabalhos; liberando o desenho para
arquivo eletrônico ou mapoteca; e enviando originais e/ou cópias do
desenho para clientes. Retorna documentos utilizados nos trabalhos para
arquivo e preenche relatório final. Organiza arquivos técnicos --
físicos e eletrônicos -- relacionados aos serviços com desenhos técnicos
de construção civil, determinando o tipo de arquivo a ser utilizado e
indexando documentos pertinentes à área. Reúne documentos e organiza
catálogos de fornecedores. Compacta e armazena arquivos digitais.
Conserva o local de trabalho limpo e organizado. Controla luminosidade
do ambiente e verifica as condições ergonômicas para uso de monitor de
vídeo, mouse, mesa digitalizadora, prancheta e teclado. Mantém
equipamentos e instrumentos de trabalho em plenas condições de uso e
funcionamento. Requisita serviço de manutenção de computador,
impressoras, plotters e outros periféricos, quando necessário. Destina
adequadamente os resíduos gerados, fazendo reciclagem e realizando
descarte de acordo com as normas ambientais. Zela pela segurança,
atuando na prevenção de acidentes e utilizando equipamentos de proteção
individual, sempre que necessário, especialmente durante os
levantamentos de dados em campo.

\begin{Form}
\textbf{Esta correspondência está adequada?} \\ 
\ChoiceMenu[radio,radiosymbol=\ding{52},bordercolor=0 0.5 0,name=cbovalid_ 318115 ]{}{CBO deve ficar=sim, \hspace{6em} CBO não fica aqui=nao}
\end{Form}

\textbf{CBO: 318120  ( Desenhista técnico (instalações hidrossanitárias) )}

Planeja o trabalho relacionado ao desenho técnico de instalações
hidrossanitárias, propondo alternativas para elaboração do desenho e
determinando a metodologia para sua execução. Seleciona instrumentos e
outros meios necessários à elaboração do desenho. Pode prever uso de
sistemas BIM - Modelagem da Informação da Construção (Building
Information Modeling). Calcula e define o custo do desenho. Estipula
prazo de execução. Prepara o local de trabalho, de acordo com
recomendações normativas de ergonomia, de segurança do trabalho e de
saúde ocupacional. Coleta dados para elaboração do desenho técnico de
instalações hidrossanitárias, interpretando projetos existentes,
consultando informações em arquivos e fazendo levantamento de campo.
Busca informações complementares e organiza os dados coletados. Processa
dados para desenho técnico de instalações hidrossanitárias, analisando
croquis obtidos por meio dos levantamentos de campo; interpretando
memória de cálculo e outras informações do conjunto de dados coletados.
Transfere dados da estação local de campo -- para o escritório ou
ambiente onde será executado o desenho --, utilizando ferramentas de
tecnologia da informação requeridas. Calcula dados da caderneta de
campo. Realiza pesquisas complementares -- inclusive via internet --,
para obter novos dados relacionados ao desenho ou ampliar a compreensão
sobre o conjunto de dados à disposição. Elabora desenhos técnicos de
instalações hidrossanitárias -- tais como desenhos de redes de esgoto,
águas pluviais, abastecimento de água e de sistemas de captação,
purificação e distribuição de água --, aplicando normas técnicas e
realizando visitas técnicas para rever dados. Utiliza softwares
específicos para desenho, especialmente CAD-Desenho Assistido por
Computador (Computer Aided Design). Utiliza estações de trabalho
computadorizadas, impressoras, plotters e outros periféricos. Pode fazer
uso de instrumentos para execução do traçado diretamente em substratos
de papel, quando requerido. Define formatos e escalas e detalha
desenhos. Diagrama pranchas ou folhas de desenho técnico. Legenda
plantas e elabora a lista de material relacionada ao conteúdo do
desenho. Revisa desenhos técnicos de instalações hidrossanitárias,
conferindo cotas, dimensionamentos, informações descritivas e outros
conteúdos. Atualiza desenhos, em função de mudanças de legislação.
Realiza a complementação ou a modificação de desenhos, quando
necessário. Em situações em que os desenhos passam por supervisão,
realiza as compatibilizações necessárias dos desenhos, de acordo com os
apontamentos da supervisão. Fecha a ordem de serviço do desenho de
instalações hidrossanitárias, relatando mudanças de procedimento
adotadas ao longo dos trabalhos; liberando o desenho para arquivo
eletrônico ou mapoteca; e enviando originais e/ou cópias do desenho para
clientes. Retorna documentos utilizados para arquivo. Organiza arquivos
técnicos -- físicos e eletrônicos -- relacionados aos serviços com
desenhos técnicos de instalações hidrossanitárias, determinando o tipo
de arquivo a ser utilizado e indexando documentos pertinentes à área.
Reúne documentos e organiza catálogos de fornecedores. Compacta e
armazena arquivos digitais. Conserva o local de trabalho limpo e
organizado. Controla luminosidade do ambiente e verifica as condições
ergonômicas para uso de monitor de vídeo, mouse, mesa digitalizadora,
prancheta e teclado. Mantém equipamentos e instrumentos de trabalho em
plenas condições de uso e funcionamento. Requisita serviço de manutenção
de computador, impressoras, plotters e outros periféricos, quando
necessário. Destina adequadamente os resíduos gerados, fazendo
reciclagem e realizando descarte de acordo com as normas ambientais.
Zela pela segurança, atuando na prevenção de acidentes e utilizando
equipamentos de proteção individual, sempre que necessário,
especialmente durante os levantamentos de dados em campo.

\begin{Form}
\textbf{Esta correspondência está adequada?} \\ 
\ChoiceMenu[radio,radiosymbol=\ding{52},bordercolor=0 0.5 0,name=cbovalid_ 318120 ]{}{CBO deve ficar=sim, \hspace{6em} CBO não fica aqui=nao}
\end{Form}

\textbf{CBO: 318505  ( Desenhista projetista de arquitetura )}

Planeja a elaboração de desenhos de projetos e detalhes arquitetônicos
de edificações e de obras de infraestrutura, avaliando a solicitação,
definindo a metodologia de trabalho, dimensionando a equipe de
desenhistas, propondo o cronograma e orçando os serviços. Seleciona
programas de informática específicos para a elaboração de projetos,
analisa o uso de BIM - Modelagem da Informação da Construção (Building
Information Modeling) e pesquisa a aplicação de novas tecnologias.
Auxilia o arquiteto na realização de estudo de viabilidade. Registra os
dados do local de construção da edificação ou da obra de infraestrutura
e obtém informações básicas junto aos órgãos competentes, para
elaboração de projeto de arquitetura. Acompanha a realização de estudo
preliminar e a preparação de anteprojetos. Realiza plano de massas e
cálculos de projeto. Elabora croquis, faz pré-dimensionamento de
estruturas e instalações, bem como elabora leiaute dos projetos
arquitetônicos, sob supervisão de arquiteto. Desenha plantas baixas,
plantas complementares, plantas de cortes e de elevações, detalhamentos
e respectivas representações gráficas, utilizando computador e programa,
tal como o CAD-Desenho Assistido por Computador (Computer Aided Design).
Auxilia na compatibilização de projetos arquitetônicos, participando de
reuniões de avaliação dos projetos e consultando normas técnicas e
padrões de desenhos arquitetônicos. Colabora na definição das diretrizes
arquitetônicas de projetos. Confere projetos sob supervisão de
arquiteto, podendo solicitar projetos complementares. Auxilia na
compatibilização de projetos executivos. No final da obra, elabora o
projeto ``como construído'' (``as built''), com as mesmas informações
daquele inicialmente aprovado, adicionadas eventuais modificações que
possam ter ocorrido desde a aprovação até o término da obra. Pode usar
programa para visualização de passeio virtual em 3D ao empreendimento,
permitindo uma ampla visualização do projeto. Auxilia arquiteto na
organização e na supervisão de equipes de trabalho - desenhistas e
estagiários -, distribuindo tarefas, fornecendo dados necessários à
elaboração dos desenhos técnicos, orientando e acompanhando o
cumprimento das atividades. Conserva o ambiente de trabalho limpo e
organizado. Controla graus de luminosidade, temperatura e ruídos.
Verifica as condições ergonômicas para uso de mobiliário, monitor de
vídeo de computador, mouse e teclado. Solicita a aquisição de materiais,
para reposição. Requisita serviço para manutenção de computador,
impressoras e outros periféricos, quando necessário. Organiza e arquiva
documentos físicos e eletrônicos, tais como projetos aprovados e
documentação referente aos projetos. Destina adequadamente os resíduos
gerados, fazendo reciclagem e realizando descarte de acordo com as
normas ambientais. Zela pela segurança, atuando na prevenção de
acidentes e utilizando Equipamentos de Proteção Individual (EPI) sempre
que necessário, especialmente durante visita às obras.

\begin{Form}
\textbf{Esta correspondência está adequada?} \\ 
\ChoiceMenu[radio,radiosymbol=\ding{52},bordercolor=0 0.5 0,name=cbovalid_ 318505 ]{}{CBO deve ficar=sim, \hspace{6em} CBO não fica aqui=nao}
\end{Form}

\textbf{CBO: 318510  ( Desenhista projetista de construção civil )}

Planeja a elaboração de desenhos de projetos de construção civil,
avaliando a solicitação, definindo a metodologia de trabalho,
dimensionando a equipe de desenhistas, propondo o cronograma e orçando
os serviços. Seleciona programas de informática específicos para a
elaboração de projetos, analisa o uso de BIM - Modelagem da Informação
da Construção (Building Information Modeling) e pesquisa a aplicação de
novas tecnologias. Auxilia o engenheiro no estudo de viabilidade técnica
do projeto construtivo. Registra os dados do local de construção e obtém
informações básicas junto aos órgãos competentes, para elaboração do
projeto. Elabora croquis, faz pré-dimensionamento de estruturas e
instalações e elabora leiaute dos projetos estruturais e de coberturas,
sob supervisão de engenheiro. Desenha plantas baixas, plantas
complementares, plantas de cortes e de elevações, detalhamentos e
respectivas representações gráficas, utilizando computador e programa,
tal como o CAD-Desenho Assistido por Computador (Computer Aided Design).
Desenha detalhes de projetos de fundação, de vigas de baldrame, de
pilares, de vigas, de lajes e de coberturas. Faz o dimensionamento de
ferragens estruturais. Auxilia na compatibilização de projetos
multidisciplinares - elétrico, eletrônico, de tubulações e hidráulicos
-, participando de reuniões de avaliação do projeto e consultando normas
técnicas e padrões de desenho de construção civil. Colabora na definição
das diretrizes para o conceito estrutural de projetos. Confere projetos
sob supervisão de engenheiros, solicitando projetos complementares.
Auxilia na compatibilização de projetos executivos. No final da obra,
elabora o projeto ``como construído'' (``as built''), com as mesmas
informações daquele inicialmente aprovado, adicionadas eventuais
modificações que possam ter ocorrido desde a aprovação até o término da
obra. Auxilia engenheiro na organização e na supervisão de equipes de
trabalho - desenhistas e estagiários -, distribuindo tarefas, fornecendo
dados necessários à elaboração dos desenhos técnicos de construção
civil, orientando e acompanhando o cumprimento das atividades. Conserva
o ambiente de trabalho limpo e organizado. Controla graus de
luminosidade, temperatura e ruídos. Verifica as condições ergonômicas
para uso de mobiliário, monitor de vídeo de computador, mouse e teclado.
Solicita a aquisição de materiais, para reposição. Requisita serviço
para manutenção de computador, impressoras e outros periféricos, quando
necessário. Organiza e arquiva documentos físicos e eletrônicos, tais
como projetos aprovados e documentação referente aos projetos. Destina
adequadamente os resíduos gerados, fazendo reciclagem e realizando
descarte de acordo com as normas ambientais. Zela pela segurança,
atuando na prevenção de acidentes e utilizando Equipamentos de Proteção
Individual (EPI) sempre que necessário, especialmente durante visita às
obras.

\begin{Form}
\textbf{Esta correspondência está adequada?} \\ 
\ChoiceMenu[radio,radiosymbol=\ding{52},bordercolor=0 0.5 0,name=cbovalid_ 318510 ]{}{CBO deve ficar=sim, \hspace{6em} CBO não fica aqui=nao}
\end{Form}

\newpage
\section{ 06005 -- Técnico em Edificações }

\textbf{Perfil do Curso (CNCT)}

O Técnico em Edificações será habilitado para:

Desenvolver projetos de arquitetura, estrutura, instalações elétricas e
hidrossanitárias de até 80 m2 usando meios físicos ou digitais. Elaborar
orçamentos de obras e serviços. Planejar a execução dos serviços de
construção e manutenção predial. Executar obras e serviços de construção
e manutenção predial. Executar ensaios de materiais de construção, solos
e controle tecnológico. Conduzir planos de qualidade da construção.
Coordenar a execução de serviços de manutenção de equipamentos e
instalações em edificações.

Para atuação como Técnico em Edificações, são fundamentais:

Conhecimentos e saberes relacionados aos processos de planejamento e
construção de edificações de modo a assegurar a saúde e a segurança dos
trabalhadores e dos futuros ocupantes do imóvel. Conhecimentos e saberes
relacionados à sustentabilidade do processo produtivo, às técnicas e
processos de produção na construção civil, às normas técnicas.
Habilidades e competências relacionadas à liderança de equipes, à
solução de problemas técnicos e trabalhistas e à gestão de conflitos.

\textbf{CBOs correspondidos e seus perfis (presentes no CNCT, e presentes na predição da IA)}

\textbf{CBO: 312105  ( Técnico de obras civis )}

Realiza levantamentos topográficos e planialtimétricos. Elabora desenhos
topográficos e desenvolve cálculos. Elabora, dentro dos limites legais,
projetos de obras civis, de estrutura de concreto e metálica, além de
outros de caráter estrutural (instalações elétricas, hidrossanitárias,
facilities e segurança etc.), compatibilizando os projetos para eliminar
interferências. Legaliza projeto de obras, providenciando documentos e
os encaminhando para aprovação pelos órgãos competentes. Controla prazo
e solicita aprovação de vistoria. Providencia encerramento das obras.
Organiza arquivos técnicos. Elabora orçamento da obra e estudos
comparativos de custos, de acordo com cronograma físico-financeiro.
Estabelece cronograma de compras, selecionando fornecedores, negociando
preços, prazos e condições de pagamento. Elabora plano de trabalho,
definindo logística e propondo cronograma de execução. Analisa a
viabilidade econômica e técnica da introdução de inovações tecnológicas
no processo construtivo, tais como robôs, alumínio transparente,
bioconcreto, concreto feito de compósitos recicláveis, novos
pré-moldados, contrapisos autonivelantes, entre outras. Dimensiona
equipes de trabalho e lista máquinas, equipamentos e ferramentas.
Supervisiona o canteiro de obras, racionalizando o trabalho e o uso de
materiais, solucionando problemas de execução, padronizando
procedimentos e acompanhando os resultados dos serviços. Elabora e
analisa relatórios técnicos, tendo em vista o controle tecnológico e
suas aplicações em obras de construção civil. Pode utilizar BIM
(Building Information Modeling ou Modelagem da Informação da
Construção), conjunto de tecnologias e processos integrados que permite
a criação, o uso e a atualização de modelos digitais de uma construção,
potencialmente servindo a todos os participantes do empreendimento,
durante todo o ciclo de vida da construção. Pode prestar assistência
técnica para construção e reforma de obras, realizando diagnósticos,
propondo soluções técnicas e tecnológicas para realização de reparos e
de revitalização ou modernização (``retrofitting'').

\begin{Form}
\textbf{Esta correspondência está adequada?} \\ 
\ChoiceMenu[radio,radiosymbol=\ding{52},bordercolor=0 0.5 0,name=cbovalid_ 312105 ]{}{CBO deve ficar=sim, \hspace{6em} CBO não fica aqui=nao}
\end{Form}

\textbf{CBO: 318105  ( Desenhista técnico (arquitetura) )}

Planeja o trabalho relacionado ao desenho técnico de arquitetura,
propondo alternativas para elaboração do desenho e determinando a
metodologia para sua execução. Seleciona instrumentos e outros meios
necessários à elaboração do desenho. Pode prever o uso de sistemas BIM -
Modelagem da Informação da Construção (Building Information Modeling).
Calcula e define o custo do desenho. Estipula prazo de execução. Prepara
o local de trabalho, de acordo com recomendações normativas de
ergonomia, de segurança do trabalho e de saúde ocupacional. Coleta dados
para elaboração do desenho técnico de arquitetura, interpretando
projetos existentes, consultando informações em arquivos e fazendo
levantamento de campo. Busca informações complementares e organiza os
dados coletados. Processa dados para desenho técnico de arquitetura,
analisando croquis obtidos por meio dos levantamentos de campo;
interpretando memória de cálculo e outras informações do conjunto de
dados coletados. Realiza pesquisas complementares -- inclusive via
internet --, para obter novos dados relacionados ao desenho ou ampliar a
compreensão sobre o conjunto de dados à disposição. Elabora desenhos
técnicos de arquitetura -- tais como plantas de apresentação e
anteprojeto; planta legal ou de licenciamento e planta humanizada,
dentre outros tipos de desenho arquitetônico --, aplicando normas
técnicas e realizando visitas técnicas para rever dados. Utiliza
softwares específicos para desenho, especialmente CAD-Desenho Assistido
por Computador (Computer Aided Design). Utiliza estações de trabalho
computadorizadas, impressoras, plotters e outros periféricos. Pode fazer
uso de instrumentos para execução do traçado diretamente em substratos
de papel, quando requerido. Define formatos e escalas e detalha
desenhos. Diagrama pranchas ou folhas de desenho técnico. Legenda
plantas e elabora a lista de material relacionada ao conteúdo do
desenho. Pode elaborar desenhos de plantas de execução. Revisa desenhos
técnicos de arquitetura, conferindo cotas, dimensionamentos, informações
descritivas e outros conteúdos. Atualiza desenhos, em função de mudanças
de legislação. Realiza a complementação ou a modificação de desenhos,
quando necessário. Em situações em que os desenhos passam por
supervisão, realiza as compatibilizações necessárias dos desenhos, de
acordo com os apontamentos da supervisão. Fecha a ordem de serviço do
desenho de arquitetura, relatando mudanças de procedimento adotadas ao
longo dos trabalhos; liberando o desenho para arquivo eletrônico ou
mapoteca; e enviando originais e/ou cópias do desenho para clientes.
Retorna documentos utilizados nos trabalhos para arquivo. Elabora
relatório final. Organiza arquivos técnicos -- físicos e eletrônicos --
relacionados aos serviços com desenhos técnicos de arquitetura,
determinando o tipo de arquivo a ser utilizado e indexando documentos
pertinentes à área. Reúne documentos e organiza catálogos de
fornecedores. Compacta e armazena arquivos digitais. Conserva o local de
trabalho limpo e organizado. Controla luminosidade do ambiente e
verifica as condições ergonômicas para uso de monitor de vídeo, mouse,
mesa digitalizadora, prancheta e teclado. Mantém equipamentos e
instrumentos de trabalho em plenas condições de uso e funcionamento.
Requisita serviço de manutenção de computador, impressoras, plotters e
outros periféricos, quando necessário. Destina adequadamente os resíduos
gerados, fazendo reciclagem e realizando descarte de acordo com as
normas ambientais. Zela pela segurança, atuando na prevenção de
acidentes e utilizando equipamentos de proteção individual, sempre que
necessário, especialmente durante os levantamentos de dados em campo.

\begin{Form}
\textbf{Esta correspondência está adequada?} \\ 
\ChoiceMenu[radio,radiosymbol=\ding{52},bordercolor=0 0.5 0,name=cbovalid_ 318105 ]{}{CBO deve ficar=sim, \hspace{6em} CBO não fica aqui=nao}
\end{Form}

\textbf{CBO: 514310  ( Auxiliar de manutenção predial )}

Apoia o serviço de manutenção elétrica, acompanhando a verificação do
funcionamento e a identificação de avarias - de equipamentos e
instalações elétricas e de iluminação -, e auxiliando na execução dos
reparos necessários. Apoia a troca de instalação elétrica e de
equipamentos de iluminação. Auxilia na execução de soldagem de objetos.
Apoia o serviço de manutenção hidráulica, acompanhando a verificação do
funcionamento das instalações, limpando e substituindo filtros,
desentupindo ralos, pias e vasos sanitários e executando outras tarefas
que lhe foram atribuídas para conserto das instalações. Acompanha e
auxilia na execução de serviços gerais de carpintaria, marcenaria e
alvenaria, reparando trincas e rachaduras, impermeabilizando
superfícies, recuperando pinturas, repondo cerâmica, consertando móveis,
forros, portas e janelas, reparando divisórias e realizando ajustes
diversos. Utiliza, sob orientação, ferramentas e equipamentos
apropriados para cada tipo de atividade. Limpa e organiza os locais onde
foram realizados os serviços de manutenção, removendo entulhos,
varrendo, eliminando resíduos e manchas. Zela pela conservação dos
materiais, ferramentas e equipamentos utilizados nos serviços de
manutenção. Pode controlar o estoque de materiais, peças, componentes,
ferramentas e equipamentos, verificando quantidades e registrando em
documentos e planilhas sua movimentação.

\begin{Form}
\textbf{Esta correspondência está adequada?} \\ 
\ChoiceMenu[radio,radiosymbol=\ding{52},bordercolor=0 0.5 0,name=cbovalid_ 514310 ]{}{CBO deve ficar=sim, \hspace{6em} CBO não fica aqui=nao}
\end{Form}

\newpage
\section{ 06006 -- Técnico em Estradas }

\textbf{Perfil do Curso (CNCT)}

O Técnico em Estradas será habilitado para:

Executar estudos topográfico, hidrológico e geotécnico. Executar
levantamento topográfico, geoprocessamento, hidrológico e geotécnico.
Auxiliar e supervisionar a elaboração de projetos BIM de topografia, de
geoprocessamento, geométrico, de terraplenagem, de drenagem, de
geotecnia, de pavimentação e de sinalização rodoviária. Auxiliar e
supervisionar a construção, o gerenciamento, a manutenção e a
conservação de vias rodoviárias e urbanas. Implementar ações para
melhoria da produtividade de máquinas e equipamentos. Supervisionar e
executar ensaios de solos, agregados, misturas betuminosas e concretos.
Elaborar orçamento, medição e controle de custos.

Para atuação como Técnico em Estradas, são fundamentais:

Conhecimentos e saberes relacionados aos processos de planejamento e
construção no âmbito da infraestrutura de modo a assegurar a saúde e a
segurança dos trabalhadores e dos futuros usuários. Conhecimentos e
saberes relacionados à sustentabilidade do processo produtivo, às
técnicas e processos de produção na construção civil, às normas
técnicas, à liderança de equipes, à solução de problemas técnicos e
trabalhistas e à gestão de conflitos.

\textbf{CBOs correspondidos e seus perfis (presentes no CNCT, e presentes na predição da IA)}

\textbf{CBO: 312205  ( Técnico de estradas )}

Elabora anteprojeto, realizando levantamento topográfico, demarcando
áreas e detalhando memorial descritivo. Solicita projetos
complementares, podendo fiscalizar sua elaboração e participar da sua
análise final. Pesquisa e analisa novas tecnologias e novos materiais,
para incorporá-los no anteprojeto. Analisa a inclusão de ferramentas
para gestão e controle do projeto, em especial BIM - Modelagem da
Informação da Construção (Building Information Modeling). Submete o
anteprojeto à aprovação do engenheiro responsável pela obra. Participa
da elaboração de proposta de preservação do meio ambiente no entorno das
obras, indicando ações preventivas e serviços de recuperação de áreas
degradadas. Participa da elaboração de planos de segurança das áreas de
mananciais. Planeja o trabalho para implantação de projetos de
construção ou de manutenção de estradas, dimensionando recursos humanos,
relacionando máquinas e equipamentos e especificando materiais. Orça a
obra, calculando custos diretos e indiretos. Elabora cronograma
físico-financeiro. Contata e cadastra fornecedores, avaliando perfis.
Participa do desenvolvimento de materiais e produtos alternativos.
Acompanha a execução dos serviços de construção ou manutenção de
estradas, verificando o cumprimento do cronograma, as condições técnicas
do local, a adequação do tamanho das equipes e o uso de máquinas e
equipamentos, em conformidade com o projeto. Registra no documento
``como construído'' (``as built'') - com representações técnicas, tais
como plantas, cortes e fachadas - todas as alterações e modificações
promovidas durante a construção, em relação ao projeto inicial. Auxilia
o engenheiro na supervisão de equipes, elaborando escalas de horários de
trabalho, treinando equipes, acompanhando o cumprimento de tarefas e
fiscalizando a aplicação de normas e regulamentos. Realiza trabalhos em
laboratório para controle de qualidade, coletando amostras e realizando
ensaios. Elabora planilhas dos ensaios realizados e emite relatórios.
Pode elaborar planta cadastral da região e montar banco de dados com as
informações de solos e bacias hidrográficas. Padroniza procedimentos e
fixa parâmetros técnicos, participando da montagem e revisão de manuais
técnicos e colaborando na elaboração e revisão das normas e
procedimentos. Pode realizar vistorias técnicas, analisando origem da
solicitação, identificando possíveis causas do problema, elaborando
registro -- croqui, fotografias, filmagens e medições -, propondo
soluções para resolução do problema e emitindo parecer técnico.
Acompanha a execução de serviços de manutenção de máquinas e
equipamentos utilizados na obra. Verifica a organização dos ambientes de
trabalho e o descarte de resíduos de acordo com normas ambientais. Zela
pela segurança, prevenindo acidentes, supervisionando o uso de
equipamentos de proteção individual e coletiva e estruturando campanhas
de prevenção e combate a incêndios.

\begin{Form}
\textbf{Esta correspondência está adequada?} \\ 
\ChoiceMenu[radio,radiosymbol=\ding{52},bordercolor=0 0.5 0,name=cbovalid_ 312205 ]{}{CBO deve ficar=sim, \hspace{6em} CBO não fica aqui=nao}
\end{Form}

\textbf{CBO: 312320  ( Topografo )}

Planeja atividades em topografia, definindo escopo, metodologia e
logística; especificando e quantificando equipamentos e instrumentos de
trabalho; e dimensionando equipe de auxiliares. Seleciona softwares para
as atividades, utilizando software de georreferenciamento e operando
softwares de automação topográfica, tais como SNGS -Sistemas de
Navegação Global por Satélite (GNSS-Global Navigation Satellite Systems)
e SIG-Sistema de Informações Geográficas (GIS-Geography Information
System). Elabora planilha de custos e cronograma físico-financeiro.
Prepara levantamentos topográficos planimétricos e altimétricos,
estudando documentos técnicos; analisando mapas, plantas, títulos de
propriedade, registros e especificações; e calculando as medições a
serem efetuadas. Analisa as características de terrenos para
reconhecimento básico da área programada, identificando pontos de
partida e vias de melhor acesso. Realiza os levantamentos topográficos
das áreas demarcadas em projetos de construção e de exploração de solos.
Determina altitudes, distâncias, ângulos, coordenadas, referências de
níveis e outras características de superfície terrestre, de áreas
subterrâneas e de edifícios, utilizando instrumentos e equipamentos,
tais como teodolitos, níveis, trenas, bússolas e telêmetros. Registra os
dados topográficos levantados, apontando valores identificados e
cálculos numéricos efetuados, tendo em vista o georreferenciamento para
projetos e obras de construção e exploração de solos. Confere e verifica
a precisão de informações e dados coligidos, consultando tabelas,
efetuando cálculos, avaliando as diferenças entre distâncias e
altitudes, aplicando fórmulas de altimetria e declividade para
nivelamentos planimétricos. Prepara documentos cartográficos, elaborando
esboços, plantas e relatórios técnicos, indicando traçados, pontos e
convenções para desenvolvimento de mapas, cartas e projetos. Pode
desenhar plantas detalhadas das áreas levantadas. Auxilia agrimensor ou
outro profissional de nível superior da área na orientação e na
supervisão de auxiliares, especificando as tarefas a serem realizadas,
detalhando o modo de execução e apontando o grau de precisão a ser
observado. Realiza aferição dos instrumentos de levantamentos
topográficos, responsabilizando-se por sua manutenção, guarda e
conservação. Destina adequadamente os resíduos gerados, orientando a
reciclagem de materiais e providenciando descarte de acordo com as
normas ambientais Zela pela segurança, fiscalizando o uso de
equipamentos de proteção individual e coletiva e atuando na prevenção de
acidentes.

\begin{Form}
\textbf{Esta correspondência está adequada?} \\ 
\ChoiceMenu[radio,radiosymbol=\ding{52},bordercolor=0 0.5 0,name=cbovalid_ 312320 ]{}{CBO deve ficar=sim, \hspace{6em} CBO não fica aqui=nao}
\end{Form}

\newpage
\section{ 06007 -- Técnico em Geodésia e Cartografia }

\textbf{Perfil do Curso (CNCT)}

O Técnico em Geodésia e Cartografia será habilitado para:

Manipular mapas analógicos e digitais para obtenção de informações
espaciais. Utilizar dados coletados por sensores orbitais e aéreos para
produção cartográfica. Coletar e processar dados para posicionamento
terrestre. Realizar o processamento digital de imagens, sistemas de
informação geográfica e desenhos assistidos por computador. Utilizar
ferramentas de geoprocessamento.

Para atuação como Técnico em Geodésia e Cartografia, são fundamentais:

Conhecimentos e saberes relacionados à elaboração de mapas, plantas,
cartas digitais, memoriais e laudos, obtidos a partir de levantamentos
topográficos e geodésicos e sensoriamento remoto. Compromisso e ética
para garantir o cumprimento da legislação e das normas técnicas
vigentes. Habilidades e saberes relacionados à liderança de equipes, à
solução de problemas técnicos e trabalhistas e à gestão de conflitos

\textbf{CBOs correspondidos e seus perfis (presentes no CNCT, e presentes na predição da IA)}

\textbf{CBO: 312305  ( Técnico em agrimensura )}

Planeja as atividades em agrimensura, definindo escopo e logística,
dimensionando equipes de trabalho, especificando e quantificando
equipamentos, acessórios e materiais. Elabora planilha de custos e
cronograma físico-financeiro. Pode utilizar software de geoprocessamento
e operar softwares de automação topográfica, tais como SNGS-Sistemas de
Navegação Global por Satélite (GNSS-Global Navigation Satellite Systems)
e SIG - Sistema de Informações Geográficas (GIS-Geography Information
System). Pode usar drones (aeronaves remotamente pilotadas), para obter
fotografias aéreas com digitalização em tempo real. Coleta dados para
atualização de plantas cadastrais urbanas e rurais, elaborando
representações gráficas e planta topográfica, em conformidade com normas
da Associação Brasileira de Normas Técnicas (ABNT). Define limites e
confrontações georreferenciadas, aviventando e declinando rumos e
azimutes magnéticos para estabelecimento de relações geométricas em
mapeamentos. Materializa marcos e pontos topográficos, delimitando
glebas e fixando obras de sistema de transporte, linhas de transmissão,
obras civis, industriais e rurais e parcelamento de solos. Elabora
projetos de terraplenagem de pequeno porte, calculando declinação
magnética, convergência meridiana, norte verdadeiro, áreas de terrenos,
volumes para movimento de solo, distâncias, azimutes e coordenadas,
concordâncias vertical e horizontal, curvas de nível por interpolação,
offset e greide, tendo em vista cotas de projeto. Executa levantamentos
geodésicos e topo-hidrográficos, medindo ângulos e distâncias e fazendo
levantamentos altimétricos e planimétricos. Realiza topografias
especiais (industriais, subterrâneas e batimétricas). Realiza operações
geodésicas. Determina coordenadas geográficas e plano-retangulares pelo
Sistema Universal Transverso de Mercator (UTM). Determina norte
verdadeiro e norte magnético, tendo em vista o georreferenciamento mais
adequado para o projeto. Demarca e mede áreas em campo, elaborando
croqui de campo e relatório de demarcação para levantamento cadastral.
Analisa documentos e informações cartográficas, coletando dados
geométricos; interpretando fotos terrestres e aéreas; interpretando
mapas, cartas e plantas; interpretando relevos para implantação de
linhas de exploração; e identificando acidentes geométricos e pontos de
apoio para georreferenciamento e amarração. Elabora documentos
cartográficos, definindo sistema de projeção, escalas e cálculos
cartográficos. Efetua aerotriangulação. Revisa documentos cartográficos,
restituindo e reambulando fotografia aérea. Edita e cria arte final de
documentos cartográficos. Auxilia agrimensor de nível superior na
organização e na supervisão de equipes de trabalho, orientando e
acompanhando o cumprimento de tarefas e realizando programas de
treinamento. Conserva o local das atividades limpo e organizado,
providenciando o acondicionamento de materiais. Mantém equipamentos e
acessórios limpos, organizados e em plenas condições de uso e
funcionamento. Destina adequadamente os resíduos gerados, orientando a
reciclagem de materiais e providenciando descarte de acordo com as
normas ambientais. Zela pela segurança, fiscalizando o uso de
equipamentos de proteção individual e coletiva e atuando na prevenção de
acidentes.

\begin{Form}
\textbf{Esta correspondência está adequada?} \\ 
\ChoiceMenu[radio,radiosymbol=\ding{52},bordercolor=0 0.5 0,name=cbovalid_ 312305 ]{}{CBO deve ficar=sim, \hspace{6em} CBO não fica aqui=nao}
\end{Form}

\textbf{CBO: 312310  ( Técnico em geodesia e cartografia )}

Planeja as atividades em geodésia e cartografia, definido escopo,
metodologia e logística; dimensionando equipes de trabalho;
especificando e quantificando equipamentos, acessórios e materiais.
Seleciona softwares para as atividades, utilizando software de
geoprocessamento e podendo operar softwares de automação topográfica,
tais como SNGS -Sistemas de Navegação Global por Satélite (GNSS-Global
Navigation Satellite Systems) e SIG-Sistema de Informações Geográficas
(GIS-Geography Information System). Pode usar drones (aeronaves
remotamente pilotadas), para obter fotografias aéreas com digitalização
em tempo real. Elabora planilha de custos e cronograma
físico-financeiro. Executa levantamentos geodésicos, medindo ângulos e
distâncias e realizando levantamentos altimétricos e planimétricos.
Determina norte verdadeiro e norte magnético. Determina coordenadas
geográficas e plano-retangulares, pelo Sistema Universal Transverso de
Mercator (UTM). Elabora relatório de demarcação para levantamento
cadastral. Coleta dados para atualização de plantas cadastrais,
efetuando cálculos de áreas de terrenos. Calcula volumes para movimento
de solo e curvas de nível por interpolação. Calcula declinação
magnética, convergência meridiana, norte verdadeiro, distâncias,
azimutes e coordenadas. Elabora representações gráficas, fazendo planta
topográfica, conforme normas da Associação Brasileira de Normas Técnicas
(ABNT). Analisa documentos e informações cartográficas, interpretando
fotos terrestres, fotos aéreas e imagens orbitais; manipulando mapas
analógicos e digitais; e processando dados de campo e dados coletados
por sensores orbitais e aéreos, para produção cartográfica. Elabora
documento cartográfico, definindo sistema de projeção, escalas e
cálculos cartográficos. Efetua aerotriangulação. Cria base cartográfica,
com base na análise de documentos e no processamento das informações
cartográficas. Edita e elabora a arte final do documento cartográfico.
Revisa documentos cartográficos, restituindo e fazendo reambulação de
fotografias aéreas. Auxilia agrimensor, engenheiro cartográfico ou outro
profissional de nível superior na organização e na supervisão das
equipes de trabalho de campo, orientando e acompanhando o cumprimento de
tarefas e realizando programas de treinamento. Conserva o local das
atividades limpo e organizado, providenciando o acondicionamento de
materiais. Mantém equipamentos e acessórios limpos, organizados e em
plenas condições de uso e funcionamento. Destina adequadamente os
resíduos gerados, orientando a reciclagem de materiais e providenciando
descarte de acordo com as normas ambientais. Zela pela segurança,
atuando na prevenção de acidentes e fiscalizando o uso de equipamentos
de proteção individual e coletiva.

\begin{Form}
\textbf{Esta correspondência está adequada?} \\ 
\ChoiceMenu[radio,radiosymbol=\ding{52},bordercolor=0 0.5 0,name=cbovalid_ 312310 ]{}{CBO deve ficar=sim, \hspace{6em} CBO não fica aqui=nao}
\end{Form}

\textbf{CBO: 312315  ( Técnico em hidrografia )}

Planeja as atividades em hidrografia, definindo escopo, metodologia e
logística e especificando equipamentos, acessórios e materiais.
Seleciona softwares para as atividades, operando sistemas de
geoprocessamento e podendo operar softwares de automação topográfica,
tais como SNGS-Sistemas de Navegação Global por Satélite (GNSS-Global
Navigation Satellite Systems) e SIG-Sistema de Informações Geográficas
(GIS-Geography Information System). Executa levantamentos geodésicos e
topo-hidrográficos, medindo ângulos e distâncias e realizando
levantamentos altimétricos e planimétricos. Determina coordenadas
geográficas e plano-retangulares, pelo Sistema Universal Transverso de
Mercator (UTM). Determina norte verdadeiro e norte magnético, tendo em
vista o georreferenciamento mais adequado para o projeto hidrográfico.
Realiza cálculos batimétricos e cálculos de marés. Mede e demarca
superfícies de águas fluviais, lacustres, marinhas e oceânicas, para
elaborar plantas e mapas de canais de navegação, fornecendo elementos
básicos para construção ou reparo de portos, cais, açudes e outras obras
marítimas, fluviais ou lacustres. Coleta dados meteorológicos,
maregráficos, oceanográficos, hidrográficos, por meio da operação de
equipamentos específicos. Interpreta boletins meteorológicos. Identifica
astros e elementos que determinam sua posição na esfera celeste.
Determina a posição de um navio por meio de métodos específicos.
Confecciona cartas náuticas oceânicas e de vias navegáveis interiores.
Elabora documentos cartográficos, definindo sistema de projeção, escalas
e cálculos cartográficos. Cria a base cartográfica e edita os
documentos. Calibra, opera e mantém, sob supervisão, equipamentos de
levantamento hidrográfico, incluindo sistemas de sonar de feixe único,
feixe múltiplo e de varredura lateral; instrumentos de posicionamento;
sensores de movimento, de velocidade do som e equipamentos auxiliares,
como dispositivos de amostragem de sedimentos e sistemas hidráulicos.
Zela pela segurança, atuando na prevenção de acidentes. Destina
adequadamente os resíduos gerados, orientando a reciclagem de materiais
e providenciando descarte de acordo com as normas ambientais.

\begin{Form}
\textbf{Esta correspondência está adequada?} \\ 
\ChoiceMenu[radio,radiosymbol=\ding{52},bordercolor=0 0.5 0,name=cbovalid_ 312315 ]{}{CBO deve ficar=sim, \hspace{6em} CBO não fica aqui=nao}
\end{Form}

\textbf{CBO: 318110  ( Desenhista técnico (cartografia) )}

Planeja o trabalho relacionado a desenho técnico de cartografia,
propondo alternativas para elaboração do desenho e determinando a
metodologia para sua execução. Seleciona instrumentos e outros meios
necessários à elaboração do desenho. Pode prever uso de sistemas BIM -
Modelagem da Informação da Construção (Building Information Modeling).
Calcula e define o custo do desenho. Estipula prazo de execução. Prepara
o local de trabalho, de acordo com recomendações normativas de
ergonomia, de segurança do trabalho e de saúde ocupacional. Coleta dados
para elaboração do desenho técnico de cartografia, interpretando
desenhos cartográficos existentes -- analógicos ou digitais -- e
consultando informações em arquivos. Interpreta fotos terrestres, fotos
aéreas e imagens orbitais. Utiliza GIS-Sistema de Informações
Geográficas (Geographic Information System) e GPS-Sistema de
Posicionamento Global (Global Positioning System) para coleta de dados.
Pode obter dados por meio de escaneamento a laser ou com uso de drones
em tempo real. Busca informações complementares e organiza dados
coletados. Processa dados para desenho técnico de cartografia,
interpretando informações do conjunto de dados coletados. Realiza
pesquisas complementares -- inclusive via internet --, para obter novos
dados relacionados ao desenho ou ampliar a compreensão sobre o conjunto
de dados à disposição. Processa dados de campo e dados coletados por
sensores orbitais e aéreos para produção cartográfica. Elabora desenhos
técnicos de cartografia -- tais como mapas, plantas e cartas, dentre
outros --, aplicando normas técnicas e realizando visitas técnicas para
rever dados. Utiliza softwares específicos para desenho cartográfico.
Pode utilizar CAD-Desenho Assistido por Computador (Computer Aided
Design), com o software funcionando de forma independente (stand-alone)
e/ou integrado ao GIS-Sistema de Informações Geográficas (Geographic
Information System) ou similar. Utiliza estações de trabalho
computadorizadas, impressoras, plotters e outros periféricos. Pode usar
instrumentos para execução do traçado diretamente em substratos de
papel, quando requerido. Define formatos e escalas e detalha desenhos.
Diagrama pranchas ou folhas de desenho técnico. Legenda plantas e
elabora lista de material relacionada ao conteúdo do desenho. Revisa
desenhos técnicos de cartografia, conferindo escalas, dimensionamentos,
informações descritivas e outros conteúdos. Atualiza desenhos, em função
de mudanças de legislação. Realiza complementação ou modificação de
desenhos, quando necessário. Em situações em que os desenhos passam por
supervisão, realiza as compatibilizações necessárias dos desenhos,
conforme os apontamentos da supervisão. Nas etapas de trabalho, aplica
convenções internacionais para cartografia. Fecha a ordem de serviço do
desenho de cartografia, relatando mudanças de procedimento adotadas ao
longo dos trabalhos; liberando o desenho para arquivo eletrônico ou
mapoteca; enviando originais e/ou cópias do desenho para clientes.
Retorna documentos utilizados para arquivo. Organiza arquivos técnicos
-- físicos e eletrônicos -- relacionados aos serviços com desenhos
técnicos de cartografia, definindo o tipo de arquivo a ser usado e
indexando documentos pertinentes à área. Reúne documentos e organiza
catálogos de fornecedores. Compacta e armazena arquivos digitais.
Conserva o local de trabalho limpo e organizado. Controla luminosidade
do ambiente e verifica condições ergonômicas para uso de monitor de
vídeo, mouse, mesa digitalizadora, prancheta e teclado. Mantém
equipamentos e instrumentos de trabalho em plenas condições de uso e
funcionamento. Requisita serviço de manutenção de computador,
impressoras, plotters e outros periféricos, quando necessário. Destina
adequadamente os resíduos gerados, fazendo reciclagem e executando
descarte de acordo com normas ambientais. Zela pela segurança, atuando
na prevenção de acidentes e usando equipamentos de proteção individual,
sempre que necessário.

\begin{Form}
\textbf{Esta correspondência está adequada?} \\ 
\ChoiceMenu[radio,radiosymbol=\ding{52},bordercolor=0 0.5 0,name=cbovalid_ 318110 ]{}{CBO deve ficar=sim, \hspace{6em} CBO não fica aqui=nao}
\end{Form}

\newpage
\section{ 06008 -- Técnico em Geoprocessamento }

\textbf{Perfil do Curso (CNCT)}

O Técnico em Geoprocessamento será habilitado para:

Executar levantamentos e coletas de dados espaciais. Implantar projetos
de sistemas de transporte, de obras civis, industriais e rurais.
Elaborar produtos cartográficos a partir de fotos terrestres, aéreas e
imagens de satélite. Analisar dados espaciais. Utilizar ferramentas de
geoprocessamento. realizar a modelagem de dados espaciais Definir
consultas relacionadas aos fenômenos mapeados para geração de relatórios
e mapas temáticos. Prestar assistência técnica na compra, venda e
utilização de equipamentos especializados. Coordenar e supervisionar a
execução de serviços técnicos. Realizar perícias técnicas. Organizar e
supervisionar levantamento e mapeamento.

Para atuação como Técnico em Geoprocessamento, são fundamentais:

Conhecimentos e saberes relacionados à obtenção de dados espaciais e
cadastrais, de mapeamento da superfície terrestre, auxiliando em
atividades nas áreas de cartografia, fotogrametria, sensoriamento remoto
e sistema de informação geográfica. Espírito de liderança, compromisso e
ética para garantir o cumprimento da legislação e das normas técnicas
vigentes. Conhecimentos relacionados à solução de problemas técnicos e
trabalhistas e à gestão de conflitos.

\textbf{CBOs correspondidos e seus perfis (presentes no CNCT, e presentes na predição da IA)}

\textbf{CBO: 312305  ( Técnico em agrimensura )}

Planeja as atividades em agrimensura, definindo escopo e logística,
dimensionando equipes de trabalho, especificando e quantificando
equipamentos, acessórios e materiais. Elabora planilha de custos e
cronograma físico-financeiro. Pode utilizar software de geoprocessamento
e operar softwares de automação topográfica, tais como SNGS-Sistemas de
Navegação Global por Satélite (GNSS-Global Navigation Satellite Systems)
e SIG - Sistema de Informações Geográficas (GIS-Geography Information
System). Pode usar drones (aeronaves remotamente pilotadas), para obter
fotografias aéreas com digitalização em tempo real. Coleta dados para
atualização de plantas cadastrais urbanas e rurais, elaborando
representações gráficas e planta topográfica, em conformidade com normas
da Associação Brasileira de Normas Técnicas (ABNT). Define limites e
confrontações georreferenciadas, aviventando e declinando rumos e
azimutes magnéticos para estabelecimento de relações geométricas em
mapeamentos. Materializa marcos e pontos topográficos, delimitando
glebas e fixando obras de sistema de transporte, linhas de transmissão,
obras civis, industriais e rurais e parcelamento de solos. Elabora
projetos de terraplenagem de pequeno porte, calculando declinação
magnética, convergência meridiana, norte verdadeiro, áreas de terrenos,
volumes para movimento de solo, distâncias, azimutes e coordenadas,
concordâncias vertical e horizontal, curvas de nível por interpolação,
offset e greide, tendo em vista cotas de projeto. Executa levantamentos
geodésicos e topo-hidrográficos, medindo ângulos e distâncias e fazendo
levantamentos altimétricos e planimétricos. Realiza topografias
especiais (industriais, subterrâneas e batimétricas). Realiza operações
geodésicas. Determina coordenadas geográficas e plano-retangulares pelo
Sistema Universal Transverso de Mercator (UTM). Determina norte
verdadeiro e norte magnético, tendo em vista o georreferenciamento mais
adequado para o projeto. Demarca e mede áreas em campo, elaborando
croqui de campo e relatório de demarcação para levantamento cadastral.
Analisa documentos e informações cartográficas, coletando dados
geométricos; interpretando fotos terrestres e aéreas; interpretando
mapas, cartas e plantas; interpretando relevos para implantação de
linhas de exploração; e identificando acidentes geométricos e pontos de
apoio para georreferenciamento e amarração. Elabora documentos
cartográficos, definindo sistema de projeção, escalas e cálculos
cartográficos. Efetua aerotriangulação. Revisa documentos cartográficos,
restituindo e reambulando fotografia aérea. Edita e cria arte final de
documentos cartográficos. Auxilia agrimensor de nível superior na
organização e na supervisão de equipes de trabalho, orientando e
acompanhando o cumprimento de tarefas e realizando programas de
treinamento. Conserva o local das atividades limpo e organizado,
providenciando o acondicionamento de materiais. Mantém equipamentos e
acessórios limpos, organizados e em plenas condições de uso e
funcionamento. Destina adequadamente os resíduos gerados, orientando a
reciclagem de materiais e providenciando descarte de acordo com as
normas ambientais. Zela pela segurança, fiscalizando o uso de
equipamentos de proteção individual e coletiva e atuando na prevenção de
acidentes.

\begin{Form}
\textbf{Esta correspondência está adequada?} \\ 
\ChoiceMenu[radio,radiosymbol=\ding{52},bordercolor=0 0.5 0,name=cbovalid_ 312305 ]{}{CBO deve ficar=sim, \hspace{6em} CBO não fica aqui=nao}
\end{Form}

\textbf{CBO: 312310  ( Técnico em geodesia e cartografia )}

Planeja as atividades em geodésia e cartografia, definido escopo,
metodologia e logística; dimensionando equipes de trabalho;
especificando e quantificando equipamentos, acessórios e materiais.
Seleciona softwares para as atividades, utilizando software de
geoprocessamento e podendo operar softwares de automação topográfica,
tais como SNGS -Sistemas de Navegação Global por Satélite (GNSS-Global
Navigation Satellite Systems) e SIG-Sistema de Informações Geográficas
(GIS-Geography Information System). Pode usar drones (aeronaves
remotamente pilotadas), para obter fotografias aéreas com digitalização
em tempo real. Elabora planilha de custos e cronograma
físico-financeiro. Executa levantamentos geodésicos, medindo ângulos e
distâncias e realizando levantamentos altimétricos e planimétricos.
Determina norte verdadeiro e norte magnético. Determina coordenadas
geográficas e plano-retangulares, pelo Sistema Universal Transverso de
Mercator (UTM). Elabora relatório de demarcação para levantamento
cadastral. Coleta dados para atualização de plantas cadastrais,
efetuando cálculos de áreas de terrenos. Calcula volumes para movimento
de solo e curvas de nível por interpolação. Calcula declinação
magnética, convergência meridiana, norte verdadeiro, distâncias,
azimutes e coordenadas. Elabora representações gráficas, fazendo planta
topográfica, conforme normas da Associação Brasileira de Normas Técnicas
(ABNT). Analisa documentos e informações cartográficas, interpretando
fotos terrestres, fotos aéreas e imagens orbitais; manipulando mapas
analógicos e digitais; e processando dados de campo e dados coletados
por sensores orbitais e aéreos, para produção cartográfica. Elabora
documento cartográfico, definindo sistema de projeção, escalas e
cálculos cartográficos. Efetua aerotriangulação. Cria base cartográfica,
com base na análise de documentos e no processamento das informações
cartográficas. Edita e elabora a arte final do documento cartográfico.
Revisa documentos cartográficos, restituindo e fazendo reambulação de
fotografias aéreas. Auxilia agrimensor, engenheiro cartográfico ou outro
profissional de nível superior na organização e na supervisão das
equipes de trabalho de campo, orientando e acompanhando o cumprimento de
tarefas e realizando programas de treinamento. Conserva o local das
atividades limpo e organizado, providenciando o acondicionamento de
materiais. Mantém equipamentos e acessórios limpos, organizados e em
plenas condições de uso e funcionamento. Destina adequadamente os
resíduos gerados, orientando a reciclagem de materiais e providenciando
descarte de acordo com as normas ambientais. Zela pela segurança,
atuando na prevenção de acidentes e fiscalizando o uso de equipamentos
de proteção individual e coletiva.

\begin{Form}
\textbf{Esta correspondência está adequada?} \\ 
\ChoiceMenu[radio,radiosymbol=\ding{52},bordercolor=0 0.5 0,name=cbovalid_ 312310 ]{}{CBO deve ficar=sim, \hspace{6em} CBO não fica aqui=nao}
\end{Form}

\textbf{CBO: 312315  ( Técnico em hidrografia )}

Planeja as atividades em hidrografia, definindo escopo, metodologia e
logística e especificando equipamentos, acessórios e materiais.
Seleciona softwares para as atividades, operando sistemas de
geoprocessamento e podendo operar softwares de automação topográfica,
tais como SNGS-Sistemas de Navegação Global por Satélite (GNSS-Global
Navigation Satellite Systems) e SIG-Sistema de Informações Geográficas
(GIS-Geography Information System). Executa levantamentos geodésicos e
topo-hidrográficos, medindo ângulos e distâncias e realizando
levantamentos altimétricos e planimétricos. Determina coordenadas
geográficas e plano-retangulares, pelo Sistema Universal Transverso de
Mercator (UTM). Determina norte verdadeiro e norte magnético, tendo em
vista o georreferenciamento mais adequado para o projeto hidrográfico.
Realiza cálculos batimétricos e cálculos de marés. Mede e demarca
superfícies de águas fluviais, lacustres, marinhas e oceânicas, para
elaborar plantas e mapas de canais de navegação, fornecendo elementos
básicos para construção ou reparo de portos, cais, açudes e outras obras
marítimas, fluviais ou lacustres. Coleta dados meteorológicos,
maregráficos, oceanográficos, hidrográficos, por meio da operação de
equipamentos específicos. Interpreta boletins meteorológicos. Identifica
astros e elementos que determinam sua posição na esfera celeste.
Determina a posição de um navio por meio de métodos específicos.
Confecciona cartas náuticas oceânicas e de vias navegáveis interiores.
Elabora documentos cartográficos, definindo sistema de projeção, escalas
e cálculos cartográficos. Cria a base cartográfica e edita os
documentos. Calibra, opera e mantém, sob supervisão, equipamentos de
levantamento hidrográfico, incluindo sistemas de sonar de feixe único,
feixe múltiplo e de varredura lateral; instrumentos de posicionamento;
sensores de movimento, de velocidade do som e equipamentos auxiliares,
como dispositivos de amostragem de sedimentos e sistemas hidráulicos.
Zela pela segurança, atuando na prevenção de acidentes. Destina
adequadamente os resíduos gerados, orientando a reciclagem de materiais
e providenciando descarte de acordo com as normas ambientais.

\begin{Form}
\textbf{Esta correspondência está adequada?} \\ 
\ChoiceMenu[radio,radiosymbol=\ding{52},bordercolor=0 0.5 0,name=cbovalid_ 312315 ]{}{CBO deve ficar=sim, \hspace{6em} CBO não fica aqui=nao}
\end{Form}

\textbf{CBO: 312320  ( Topografo )}

Planeja atividades em topografia, definindo escopo, metodologia e
logística; especificando e quantificando equipamentos e instrumentos de
trabalho; e dimensionando equipe de auxiliares. Seleciona softwares para
as atividades, utilizando software de georreferenciamento e operando
softwares de automação topográfica, tais como SNGS -Sistemas de
Navegação Global por Satélite (GNSS-Global Navigation Satellite Systems)
e SIG-Sistema de Informações Geográficas (GIS-Geography Information
System). Elabora planilha de custos e cronograma físico-financeiro.
Prepara levantamentos topográficos planimétricos e altimétricos,
estudando documentos técnicos; analisando mapas, plantas, títulos de
propriedade, registros e especificações; e calculando as medições a
serem efetuadas. Analisa as características de terrenos para
reconhecimento básico da área programada, identificando pontos de
partida e vias de melhor acesso. Realiza os levantamentos topográficos
das áreas demarcadas em projetos de construção e de exploração de solos.
Determina altitudes, distâncias, ângulos, coordenadas, referências de
níveis e outras características de superfície terrestre, de áreas
subterrâneas e de edifícios, utilizando instrumentos e equipamentos,
tais como teodolitos, níveis, trenas, bússolas e telêmetros. Registra os
dados topográficos levantados, apontando valores identificados e
cálculos numéricos efetuados, tendo em vista o georreferenciamento para
projetos e obras de construção e exploração de solos. Confere e verifica
a precisão de informações e dados coligidos, consultando tabelas,
efetuando cálculos, avaliando as diferenças entre distâncias e
altitudes, aplicando fórmulas de altimetria e declividade para
nivelamentos planimétricos. Prepara documentos cartográficos, elaborando
esboços, plantas e relatórios técnicos, indicando traçados, pontos e
convenções para desenvolvimento de mapas, cartas e projetos. Pode
desenhar plantas detalhadas das áreas levantadas. Auxilia agrimensor ou
outro profissional de nível superior da área na orientação e na
supervisão de auxiliares, especificando as tarefas a serem realizadas,
detalhando o modo de execução e apontando o grau de precisão a ser
observado. Realiza aferição dos instrumentos de levantamentos
topográficos, responsabilizando-se por sua manutenção, guarda e
conservação. Destina adequadamente os resíduos gerados, orientando a
reciclagem de materiais e providenciando descarte de acordo com as
normas ambientais Zela pela segurança, fiscalizando o uso de
equipamentos de proteção individual e coletiva e atuando na prevenção de
acidentes.

\begin{Form}
\textbf{Esta correspondência está adequada?} \\ 
\ChoiceMenu[radio,radiosymbol=\ding{52},bordercolor=0 0.5 0,name=cbovalid_ 312320 ]{}{CBO deve ficar=sim, \hspace{6em} CBO não fica aqui=nao}
\end{Form}

\textbf{CBO: 318110  ( Desenhista técnico (cartografia) )}

Planeja o trabalho relacionado a desenho técnico de cartografia,
propondo alternativas para elaboração do desenho e determinando a
metodologia para sua execução. Seleciona instrumentos e outros meios
necessários à elaboração do desenho. Pode prever uso de sistemas BIM -
Modelagem da Informação da Construção (Building Information Modeling).
Calcula e define o custo do desenho. Estipula prazo de execução. Prepara
o local de trabalho, de acordo com recomendações normativas de
ergonomia, de segurança do trabalho e de saúde ocupacional. Coleta dados
para elaboração do desenho técnico de cartografia, interpretando
desenhos cartográficos existentes -- analógicos ou digitais -- e
consultando informações em arquivos. Interpreta fotos terrestres, fotos
aéreas e imagens orbitais. Utiliza GIS-Sistema de Informações
Geográficas (Geographic Information System) e GPS-Sistema de
Posicionamento Global (Global Positioning System) para coleta de dados.
Pode obter dados por meio de escaneamento a laser ou com uso de drones
em tempo real. Busca informações complementares e organiza dados
coletados. Processa dados para desenho técnico de cartografia,
interpretando informações do conjunto de dados coletados. Realiza
pesquisas complementares -- inclusive via internet --, para obter novos
dados relacionados ao desenho ou ampliar a compreensão sobre o conjunto
de dados à disposição. Processa dados de campo e dados coletados por
sensores orbitais e aéreos para produção cartográfica. Elabora desenhos
técnicos de cartografia -- tais como mapas, plantas e cartas, dentre
outros --, aplicando normas técnicas e realizando visitas técnicas para
rever dados. Utiliza softwares específicos para desenho cartográfico.
Pode utilizar CAD-Desenho Assistido por Computador (Computer Aided
Design), com o software funcionando de forma independente (stand-alone)
e/ou integrado ao GIS-Sistema de Informações Geográficas (Geographic
Information System) ou similar. Utiliza estações de trabalho
computadorizadas, impressoras, plotters e outros periféricos. Pode usar
instrumentos para execução do traçado diretamente em substratos de
papel, quando requerido. Define formatos e escalas e detalha desenhos.
Diagrama pranchas ou folhas de desenho técnico. Legenda plantas e
elabora lista de material relacionada ao conteúdo do desenho. Revisa
desenhos técnicos de cartografia, conferindo escalas, dimensionamentos,
informações descritivas e outros conteúdos. Atualiza desenhos, em função
de mudanças de legislação. Realiza complementação ou modificação de
desenhos, quando necessário. Em situações em que os desenhos passam por
supervisão, realiza as compatibilizações necessárias dos desenhos,
conforme os apontamentos da supervisão. Nas etapas de trabalho, aplica
convenções internacionais para cartografia. Fecha a ordem de serviço do
desenho de cartografia, relatando mudanças de procedimento adotadas ao
longo dos trabalhos; liberando o desenho para arquivo eletrônico ou
mapoteca; enviando originais e/ou cópias do desenho para clientes.
Retorna documentos utilizados para arquivo. Organiza arquivos técnicos
-- físicos e eletrônicos -- relacionados aos serviços com desenhos
técnicos de cartografia, definindo o tipo de arquivo a ser usado e
indexando documentos pertinentes à área. Reúne documentos e organiza
catálogos de fornecedores. Compacta e armazena arquivos digitais.
Conserva o local de trabalho limpo e organizado. Controla luminosidade
do ambiente e verifica condições ergonômicas para uso de monitor de
vídeo, mouse, mesa digitalizadora, prancheta e teclado. Mantém
equipamentos e instrumentos de trabalho em plenas condições de uso e
funcionamento. Requisita serviço de manutenção de computador,
impressoras, plotters e outros periféricos, quando necessário. Destina
adequadamente os resíduos gerados, fazendo reciclagem e executando
descarte de acordo com normas ambientais. Zela pela segurança, atuando
na prevenção de acidentes e usando equipamentos de proteção individual,
sempre que necessário.

\begin{Form}
\textbf{Esta correspondência está adequada?} \\ 
\ChoiceMenu[radio,radiosymbol=\ding{52},bordercolor=0 0.5 0,name=cbovalid_ 318110 ]{}{CBO deve ficar=sim, \hspace{6em} CBO não fica aqui=nao}
\end{Form}

\newpage
\section{ 06011 -- Técnico em Portos }

\textbf{Perfil do Curso (CNCT)}

O Técnico em Portos será habilitado para:

Desenvolver atividades de gerenciamento, monitoramento, supervisão,
programação e controle em operações portuárias diversas. Controlar,
programar e coordenar operações de transportes em geral, inclusive o
transporte de cargas perigosas. Prestar suporte técnico em atividades de
armazenagem de cargas, inclusive armazenagem de cargas perigosas.
Supervisionar operações de embarque, transbordo e desembarque de cargas
entre os diversos modos de transporte. Prestar suporte técnico para o
agenciamento de embarcações. Encaminhar procedimentos de importação e
exportação. Verificar as condições de segurança dos meios de
transportes, equipamentos utilizados e das cargas. Programar e
supervisionar a manutenção de equipamentos eletromecânicos de operação
portuária. Verificar e inspecionar a eficiência operacional de
equipamentos e veículos. Interpretar, elaborar e preparar a documentação
necessária ao desembaraço aduaneiro de cargas. Atender clientes internos
e externos. Elaborar a cotação de preços de serviços de transporte,
inclusive transporte multimodal. Identificar e programar rotas de
transporte de cargas. Utilizar tecnologias aplicadas ao processo de
gestão da informação sobre condições do transporte e da carga.

Para atuação como Técnico em Portos, são fundamentais:

Conhecimentos e saberes relacionados aos processos de desembaraço
aduaneiro de cargas, transporte terrestre de contêineres, operações
logísticas, transporte e armazenagem de mercadorias perigosas,
sistemáticas de importação e exportação, operações de
embarque/desembarque de navios e logística de armazéns. Comprometimento
com as questões ambientais, sociais e de desenvolvimento tecnológico e
para a solução de problemas e busca por inovações tecnológicas.
Conhecimentos relacionados à liderança de equipes, à solução de
problemas técnicos e trabalhistas e à gestão de conflitos.

\textbf{CBOs correspondidos e seus perfis (presentes no CNCT, e presentes na predição da IA)}

\textbf{CBO: 342210  ( Despachante aduaneiro )}

Executa atividades relacionadas ao despacho aduaneiro de mercadorias -
inclusive bagagem de viajante - na importação, na exportação e na
internação, realizadas por qualquer via de transporte. Elabora e/ou
supervisiona a elaboração de documentos relativos ao despacho aduaneiro.
Providencia credenciamento de cliente junto à Receita Federal. Analisa
cartas de crédito. Elabora fatura proforma e fatura comercial. Prepara
lista de embarque. Executa e/ou supervisiona a classificação de
mercadorias, analisando a documentação técnica das mercadorias;
verificando suas funções de uso e aplicação; e observando seu material
constitutivo. Analisa necessidade de anuência de órgãos pertinentes e,
quando necessário, consulta a Receita Federal sobre dúvidas de
classificação. Analisa enquadramento legal das alíquotas de impostos de
importação e exportação. Enquadra mercadorias na Tarifa Externa Comum
(TEC), na Nomenclatura da Associação Latino-americana de Integração
(Naladi) e na Tabela de Incidência do Imposto sobre Produtos
Industrializados (TIPI). Preenche Certificado de Origem; Conhecimento de
Transporte internacional; Documento de Arrecadação de Receitas Federais
(DARF); guia de Imposto sobre Operações relativas à Circulação de
Mercadorias e sobre Prestações de Serviços de Transporte Interestadual e
Intermunicipal e de Comunicação (ICMS); além de outros formulários
necessários às operações de comércio exterior. Realiza e/ou supervisiona
a atividade de entrada de dados no Sistema Integrado de Comércio
Exterior (Siscomex) e sistemas auxiliares. Coleta dados de operações e
registra licenças e declarações de importação e de despacho de
exportação. Efetua registros no módulo Registro de Operações Financeiras
do Sistema Registro Declaratório Eletrônico (RDE-ROF) e de atos
concessórios de regime aduaneiro especial (``Drawback''). Faz Declaração
de Trânsito Aduaneiro (DTA). Efetua acompanhamento da tramitação de
documentos relacionados ao despacho aduaneiro. Solicita deferimento de
licenças de importação e de declarações retificadoras. Subscreve
documentos, inclusive termos de responsabilidade. Faz recebimento de
intimações, de notificações, de autos de infração, de despachos, de
decisões e de outros atos e termos processuais relacionados com o
procedimento de despacho aduaneiro. Realiza ou supervisiona o
recebimento de mercadorias e bagagens. Paga impostos e taxas. Pode
apresentar desistência de vistoria aduaneira. Executa acompanhamento de
vistoria da mercadoria na conferência aduaneira, observando a
conferência física e a retirada de amostra de mercadoria para
assistência técnica e perícia. Avalia exigências da receita federal e
demais órgãos pertinentes, presta esclarecimentos e apresenta
documentação. Retira comprovantes de importação e exportação junto à
receita federal. Presta assessoria a importadores e exportadores,
identificando suas necessidades; informando sobre legislação pertinente;
propondo alternativas de logística de transportes; e avaliando
benefícios fiscais. Calcula custos de operação. Pleiteia benefícios
fiscais. Monitora processos administrativos e informa a cliente sobre
tramitação de processos. Realiza contratação de serviços de terceiros,
supervisionando a cotação de preços de armazenagem e preços de fretes
nacionais e internacionais; orçando seguros; e negociando preços e
prazos. Confere cobrança de terceiros por serviços prestados e
providencia o pagamento. Supervisiona equipe de trabalho, distribuindo
tarefas, avaliando desempenho e desenvolvendo treinamentos. Monitora a
limpeza e a organização do ambiente de trabalho. Controla a aplicação de
procedimentos da empresa para cumprimento de normas ambientais,
inspecionando o reaproveitamento de sobras de material e o descarte de
resíduos.

\begin{Form}
\textbf{Esta correspondência está adequada?} \\ 
\ChoiceMenu[radio,radiosymbol=\ding{52},bordercolor=0 0.5 0,name=cbovalid_ 342210 ]{}{CBO deve ficar=sim, \hspace{6em} CBO não fica aqui=nao}
\end{Form}

\textbf{CBO: 342605  ( Chefe de estação portuária )}

Planeja as atividades de estações portuárias (instalações portuárias -
localizadas dentro ou fora da área do porto organizado - utilizadas em
movimentação de passageiros e/ou de cargas, no transporte aquaviário),
levando em conta os espaços disponíveis, os fluxos, os equipamentos e os
demais recursos instalados. Supervisiona recebimento e embarque de
cargas e passageiros nas estações portuárias, examinando a previsão de
operações de embarque e desembarque e analisando os relatórios das
embarcações. Supervisiona o recebimento de documentação e o processo de
liberação de cargas e passageiros. Supervisiona a equipe de operação.
Organiza distribuição de cargas e passageiros, selecionando tipos e
tamanhos de embarcações de acordo com volume e peso de cargas e
quantidade de passageiros; estipulando destinos para as embarcações; e
calculando horários de saída e de chegada das viagens. Pode definir a
utilização de navios especializados de acordo com o tipo de carga. Pode
organizar a distribuição de cargas com uso de robôs. Coordena serviços
de embarcação em estações portuárias, distribuindo tripulação para
embarcações; determinando horários para as operações; e definindo tipos
de embarcação de acordo com a demanda por serviços. Determina o
acionamento de motores para dar início à movimentação de embarcações.
Providencia serviços de auxílio à tripulação e serviços de auxílio às
operações com embarcações. Define prioridades de atracação de
embarcações, de acordo com a programação das operações, tendo em vista
maior agilidade operacional, com redução do período em que a embarcação
permanece atracada. Nas operações de atracação, desatracação, embarque e
desembarque nas estações portuárias, distribui tarefas a serem cumpridas
na embarcação e/ou na estação; e estabelece ponte de embarque e
desembarque. Monitora o embarque e o desembarque de passageiros. Presta
assessoramento nas ocasiões de visitas de representantes da capitania
dos portos, apresentando as instalações da estação portuária e
disponibilizando as informações solicitadas. Colabora com setor de
tecnologia da empresa no desenvolvimento e na implantação de recursos
para a área de segurança da estação portuária, tais como sistemas de
localização em tempo real, análise inteligente de vídeo, câmeras
térmicas e reconhecimento facial, dentre outros. Participa na
implantação de recursos tecnológicos - tais como ``QR Code'' (sistema de
resposta rápida com código de barras) -, para facilitar embarques e
desembarques em cruzeiros internacionais e outras modalidades de
serviço. Estuda e implanta soluções para redução de perdas e
melhoramento contínuo da produtividade e da segurança, nas operações da
estação portuária. Pode fazer uso de sistemas do tipo ``Big Data''
(grandes conjuntos de dados), analisando informações e empregando
técnicas e critérios definidos, a fim de prever demandas e tendências
que interferem nas atividades das estações portuárias. Prepara a equipe
de trabalho, avaliando desempenho; identificando necessidades de
treinamento; e viabilizando a participação em programas de treinamento.
Elabora documentos administrativos relacionados às atividades da estação
portuária, tais como listas de equipamentos de embarcação, planilhas de
estatísticas operacionais, dentre outros. Encaminha a documentação
elaborada para setores da empresa. Mantém iniciativas integradoras entre
a estação portuária e seu entorno, considerando aspectos urbanísticos,
ambientais e de negócios, bem como a integração com outros modais de
transporte. Zela pela preservação ambiental, colaborando -- no âmbito da
estação portuária -- para a redução de emissões no transporte
aquaviário. Trabalha com segurança, usando equipamentos de proteção
individual e aplicando princípios básicos de ergonomia. Desenvolve
campanhas de prevenção de acidentes e de proteção do meio ambiente.

\begin{Form}
\textbf{Esta correspondência está adequada?} \\ 
\ChoiceMenu[radio,radiosymbol=\ding{52},bordercolor=0 0.5 0,name=cbovalid_ 342605 ]{}{CBO deve ficar=sim, \hspace{6em} CBO não fica aqui=nao}
\end{Form}

\textbf{CBO: 342610  ( Supervisor de operações portuárias )}

Planeja as atividades, considerando espaços disponíveis e fluxos de
pessoas e de cargas nos locais de trabalho. Participa em reuniões com
representantes de autoridades portuárias, agentes e operadores, a fim de
obter subsídios para o planejamento das operações. Apresenta, para as
autoridades portuárias, a previsão de carga e descarga e informa o
período previsto das operações. Informa, ao setor financeiro da empresa,
valores para pagamentos de taxas portuárias. Supervisiona o recebimento
de documentação e o processo de liberação de cargas. Verifica
documentação de exportação e de importação de mercadorias. Controla
distribuição de cargas, monitorando a organização das cargas segundo
peso; por tipo; e por embarcação. Confere a separação de tipos de
contêineres por porto de destino. Seleciona tipos e tamanhos de
embarcações e calcula horários de saída e de chegada no destino. Mapeia
infraestrutura para recebimento de cargas. Utiliza sistemas de automação
na distribuição de cargas e nas demais operações portuárias. Programa
atracação de embarcações, definindo prioridades; agendando atracação de
embarcações; e providenciando equipe de trabalhadores. Organiza serviços
de auxílio à tripulação e serviços de auxílio às operações com
embarcações. Compatibiliza o plano de carga e descarga com os relatórios
operacionais. Avisa autoridades sobre a chegada de embarcações. Examina
previsão de embarque e desembarque, analisando relatórios das
embarcações. Define equipamentos para carga e descarga. Controla o fluxo
de caminhões no porto e nas suas áreas de acesso. Monitora a
movimentação de cargas em embarcação, estabelecendo ponte de embarque e
desembarque; distribuindo tarefas nas instalações portuárias e na
embarcação; definindo equipamentos para carga e descarga; monitorando
desapeação e desembarque de carga; e controlando serviços de embarque e
peação de carga. Colabora nas iniciativas de melhoria da infraestrutura
da logística portuária, apresentando sugestões, fazendo recomendações e
participando no desenvolvimento de soluções. Analisa relatórios das
operações portuárias, utilizando critérios estabelecidos pela empresa, a
fim de promover o aperfeiçoamento permanente das operações. Colabora com
a empresa no desenvolvimento e na implantação de soluções para a área da
segurança nas operações portuárias, tais como sistemas de localização em
tempo real e análise inteligente de vídeo. Colabora em outros
desenvolvimentos que objetivam o aumento da produtividade nos portos,
como aplicações de inteligência artificial e algoritmos de atracação,
dentre outros. Presta assessoramento nas ocasiões de visitas -- de
representantes da capitania dos portos, da receita federal, da polícia
federal e da vigilância sanitária --, apresentando as instalações
portuárias e disponibilizando as informações solicitadas. Relaciona-se
com clientes, apresentando as vantagens operacionais dos serviços sob
sua responsabilidade. Prepara a equipe de trabalho, avaliando seu
desempenho, levantando necessidades de treinamento e viabilizando a
participação dos trabalhadores em programas de treinamento. Elabora
documentos relacionados às operações portuárias, tais como lista de
equipamento das embarcações, relatórios com características das
embarcações, lista de sequenciamento de carga e descarga e relatório de
recebimento de carga. Elabora planilhas de estatísticas operacionais e
planilhas de custo das operações portuárias. Preenche livros de
ocorrências. Encaminha documentação para diversos órgãos, tais como
setores da empresa e autoridades portuárias. Participa dos esforços em
prol da desburocratização nas operações portuárias. Trabalha com
segurança, usando equipamentos de proteção individual e aplicando
princípios básicos de ergonomia. Desenvolve campanhas de prevenção de
acidentes e de proteção do meio ambiente.

\begin{Form}
\textbf{Esta correspondência está adequada?} \\ 
\ChoiceMenu[radio,radiosymbol=\ding{52},bordercolor=0 0.5 0,name=cbovalid_ 342610 ]{}{CBO deve ficar=sim, \hspace{6em} CBO não fica aqui=nao}
\end{Form}

\textbf{CBO: 414215  ( Conferente de carga e descarga )}

Registra dados de produção -- como horas-homens trabalhadas,
horas-máquinas, entre outros --, acompanha pesagem dos produtos e
realiza a separação da produção, para organizar a logística de cargas.
Realiza a cubagem de cargas e dos espaços destinados à carga. Registra
tempo de carga e descarga, para controlar o transporte. Elabora plano de
carga e organiza listas de volumes por espaço. Verifica documentação de
carga e descarga. Monitora processos, requerendo expedição da carga
gerada em processo de produção; coordenando fluxos de carga;
determinando local de descarga. Intermedeia atividades de operações de
navios, em função das atividades de carga e descarga. Realiza
conferência de mercadorias e outros materiais, nas operações de carga
(expedição) e descarga, utilizando checklist. Verifica tipo, peso,
número de lote, marca, contramarca, procedência, destino e manifesto da
carga. Realiza inspeção visual, conta volumes, inspeciona lacres,
assiste pesagem da carga e apura casos de faltas e acréscimos. Comunica
ocorrências e solicita providências. Acompanha uso de equipamentos de
terceiros. Encaminha documentos e repassa informações ao superior e a
outros setores, para relatar ocorrências e acidentes de trabalho. Emite
termo de avaria. Redige memorandos. Documenta atividades, preenchendo
relatórios, impressos, formulários, guias, boletins e ficha de liberação
de lotes. Organiza, encaminha e arquiva a documentação relacionada aos
trabalhos do setor. Lidera equipe, distribuindo tarefas e acompanhando o
desempenho de cada profissional. Repassa informações sobre resultados
almejados à equipe de trabalho, esclarecendo dúvidas. Fiscaliza uso de
equipamentos de proteção individual e distribui máquinas e equipamentos
para realização dos trabalhos. Orienta os trabalhadores quanto à
segurança do trabalho, participando do Diálogo Diário de Segurança
(DDS).

\begin{Form}
\textbf{Esta correspondência está adequada?} \\ 
\ChoiceMenu[radio,radiosymbol=\ding{52},bordercolor=0 0.5 0,name=cbovalid_ 414215 ]{}{CBO deve ficar=sim, \hspace{6em} CBO não fica aqui=nao}
\end{Form}

\newpage
\section{ 06012 -- Técnico em Saneamento }

\textbf{Perfil do Curso (CNCT)}

O Técnico em Saneamento será habilitado para:

Coordenar projetos e obras de aterros sanitários. Supervisionar a
disposição e reciclagem de resíduos em unidades de compostagem.
Desenvolver, coordenar e executar projetos de obras de sistemas e
estação de tratamento de águas (captação, transporte, tratamento e
distribuição) e de esgotos (coleta, transporte, tratamento e disposição
final). Executar e fiscalizar obras de drenagem urbana. Realizar a
manutenção de equipamentos e redes. Estruturar o serviço de coleta de
resíduos sólidos das obras. Controlar os procedimentos de preservação do
meio ambiente. Fiscalizar atividades e obras. Realizar vistorias,
inspeções e análises técnicas de projetos, obras e processos. Promover a
educação sanitária e ambiental. Elaborar orçamentos de obras e serviços
de saneamento básico.

Para atuação como Técnico em Saneamento, são fundamentais:

Conhecimentos e saberes relacionados às atividades de planejamento e
elaboração de projetos, associados à operação e manutenção de sistemas e
estação de tratamento de águas e esgoto, resíduos sólidos e drenagem
urbana, respeitando as normas de higiene, saúde e segurança no trabalho.
Além disso, deve prezar pela ética, viabilidade técnico-econômica e
preservação do meio ambiente, ter espírito inovador e empreendedor, e
ser capaz de supervisionar equipes com o intuito de solucionar problemas
técnicos e trabalhistas.

\textbf{CBOs correspondidos e seus perfis (presentes no CNCT, e presentes na predição da IA)}

\textbf{CBO: 312210  ( Técnico de saneamento )}

Elabora anteprojeto de infraestrutura para construção de estações de
tratamento de água e esgoto, realizando levantamento topográfico,
demarcando áreas e detalhando memorial descritivo. Solicita a
contratação de serviços terceirizados para elaboração de projetos
complementares, podendo fiscalizar a execução dos serviços e participar
da sua avaliação final. Pesquisa e analisa a incorporação no anteprojeto
de novos materiais, de novas tecnologias e de ferramentas de gestão e
controle, em especial da metodologia BIM - Modelagem da Informação da
Construção (Building Information Modeling). Submete o anteprojeto à
aprovação do engenheiro responsável pela obra. Planeja o trabalho para
construção de estações de tratamento de água e esgoto, dimensionando
recursos humanos, especificando materiais e selecionando máquinas e
equipamentos. Orça a obra, calculando custos diretos e indiretos.
Elabora cronograma físico-financeiro. Contata e cadastra fornecedores,
avaliando perfis. Participa do desenvolvimento de materiais e produtos
alternativos. Controla a execução da obra, verificando o cumprimento do
cronograma, as condições técnicas do local, a adequação do tamanho das
equipes e o uso de máquinas e equipamentos. Registra no documento ``como
construído'' (``as built'') - representações técnicas, tais como
plantas, corte e fachadas - todas as modificações realizadas durante a
construção, em relação ao projeto inicial. Auxilia o engenheiro na
supervisão de equipes durante a construção de estações de tratamento de
água e esgoto, elaborando escalas de horários de trabalho, treinando
equipes, acompanhando o cumprimento de tarefas e fiscalizando a
aplicação de normas e regulamentos. Realiza trabalhos em laboratório
para controle de qualidade, coletando amostras e realizando ensaios.
Elabora planilhas dos ensaios realizados e emite relatórios. Padroniza
procedimentos e fixa parâmetros técnicos, participando da montagem e
revisão de manuais técnicos e colaborando na elaboração e revisão das
normas e procedimentos. Pode elaborar planta cadastral da região e
montar banco de dados com as informações de solos e bacias
hidrográficas. Pode realizar vistorias técnicas, analisando origem da
solicitação, identificando possíveis causas do problema, elaborando
registro -- croqui, fotografias, filmagens e medições -, propondo
soluções para resolução do problema e emitindo parecer técnico. Pode
coordenar projetos e obras de aterros sanitários e fiscalizar ou
supervisionar a execução de obras de drenagem urbana. Controla os
procedimentos de preservação do meio ambiente. Propõe ações preventivas,
supervisiona serviços de recuperação de áreas degradadas e participa da
elaboração de planos de segurança das áreas de mananciais. Pode
implantar projetos de reflorestamento e recuperação de matas. Efetua
vistorias e perícias técnicas circundantes ao ecossistema local,
coletando dados referentes à fauna e flora, identificando fontes de
poluição e sugerindo formas de minimização de impactos ambientais.
Planeja e supervisiona a execução de campanha educativa para coleta de
resíduos sólidos. Emite relatório final da campanha. Estrutura o serviço
de coleta dos resíduos sólidos das obras. Realiza diagnóstico da
realidade local no entorno da construção de estações de tratamento de
água e esgoto, identificando as características dos resíduos, elaborando
roteiro de coleta e/ou varrição e acompanhando a coleta. Destina
adequadamente os resíduos coletados, reaproveitando materiais e
realizando descarte de acordo com as normas ambientais. Acompanha a
execução de serviços de manutenção de máquinas e equipamentos utilizados
na obra. Zela pela segurança, prevenindo acidentes, supervisionando o
uso de equipamentos de proteção individual e coletiva e estruturando
campanhas de prevenção e combate a incêndios.

\begin{Form}
\textbf{Esta correspondência está adequada?} \\ 
\ChoiceMenu[radio,radiosymbol=\ding{52},bordercolor=0 0.5 0,name=cbovalid_ 312210 ]{}{CBO deve ficar=sim, \hspace{6em} CBO não fica aqui=nao}
\end{Form}

\newpage
\section{ 06014 -- Técnico em Trânsito }

\textbf{Perfil do Curso (CNCT)}

O Técnico em Trânsito será habilitado para:

Realizar procedimentos de gestão, planejamento, fiscalização e operação
do trânsito. Promover a educação e a segurança do trânsito. Organizar a
operação do tráfego urbano. Organizar o controle da manutenção de
equipamentos de tráfego, o monitoramento do trânsito e das vias
públicas, a fiscalização de trânsito e de veículos. Supervisionar o
cumprimento da legislação referente ao trânsito de veículos. Realizar
pesquisas e tratamentos estatísticos de tráfego. Supervisionar operações
de tráfego. Realizar estudos e implantação de melhorias para o trânsito
nas vias rurais, nas cidades e em regiões metropolitanas.

Para atuação como Técnico em Trânsito, são fundamentais:

Conhecimentos e saberes relacionados aos processos de gestão,
planejamento, projeto, implantação e operação de atividades relacionadas
ao transporte de passageiros, à logística e ao transporte de cargas.
Conhecimentos e saberes relacionados à educação e à segurança no
trânsito; à engenharia de tráfego, assegurando a mobilidade urbana de
acordo com as técnicas e processos; e às normas técnicas. Habilidades e
competências relacionadas à liderança de equipes, à solução de problemas
técnicos e trabalhistas e à gestão de conflitos.

\textbf{CBOs correspondidos e seus perfis (presentes no CNCT, e presentes na predição da IA)}

\textbf{CBO: 342305  ( Chefe de serviço de transporte rodoviário (passageiros e cargas) )}

Planeja e controla a disponibilidade da frota para atendimento à demanda
de serviços, determinando a escala de utilização do conjunto de veículos
de transporte rodoviário da empresa. Analisa demanda do mercado para, se
necessário, redimensionar a frota, remanejar os veículos e contratar
serviços de terceiros. Estabelece a programação de viagens, levando em
conta os percursos e as jornadas de trabalho dos motoristas e
auxiliares. Supervisiona o embarque de carga e de passageiros,
monitorando horários de entrada e saída de veículos. Inspeciona
baldeação de bagagem. Emite documento fiscal da carga e libera recursos
financeiros para viagem. Pode entrevistar os motoristas antes da viagem,
antecipando e resolvendo eventuais problemas. Monitora os veículos ao
longo do percurso - inclusive por meio de sistemas de rastreamento -,
para informar-se sobre as ocorrências. Atende acidentes na estrada,
deslocando-se ao local do acidente. Providencia socorro às vítimas.
Aciona polícia e perícia técnica. Relata à empresa o número e o destino
das vítimas e a relação de passageiros, acionando o setor de serviço
social. Informa sobre as circunstâncias do acidente ao departamento
jurídico da empresa. Fornece ao cliente as informações relacionadas ao
acidente. Aciona a seguradora, fotografa o local do acidente e
providencia remoção do veículo do local. Quando possível, atua para dar
continuidade ao serviço, assistindo passageiros não vitimados,
recolhendo seus pertences pessoais e providenciando veículo, motorista e
auxiliares para o prosseguimento da viagem. Realiza o acerto de contas
de viagens com os motoristas, conferindo notas fiscais, controlando o
consumo de combustível e verificando a planilha de despesas de viagens.
Fecha boletins de caixa e encaminha prestação de contas ao setor
financeiro da empresa. Controla custos das viagens. Supervisiona
motoristas e auxiliares de viagem, conferindo sua documentação e
providenciando treinamento. Orienta motoristas e auxiliares,
advertindo-os, se necessário, verbalmente e por escrito. Atende
clientes, registrando e tratando suas reclamações e sugestões. Encaminha
veículos para serviços de manutenção.

\begin{Form}
\textbf{Esta correspondência está adequada?} \\ 
\ChoiceMenu[radio,radiosymbol=\ding{52},bordercolor=0 0.5 0,name=cbovalid_ 342305 ]{}{CBO deve ficar=sim, \hspace{6em} CBO não fica aqui=nao}
\end{Form}

\textbf{CBO: 342310  ( Inspetor de serviços de transportes rodoviários (passageiros e cargas) )}

Inspeciona o embarque, a trajetória e o desembarque de passageiros e
cargas, verificando o cumprimento dos horários previstos, observando a
conduta dos motoristas e registrando as ocorrências durante o percurso.
Na impossibilidade do cumprimento da programação de uma viagem,
encaminha soluções, podendo remanejar veículos e substituir motoristas e
auxiliares. Inspeciona o trabalho dos motoristas e auxiliares de viagem,
examinando seu desempenho. Confere a documentação dos motoristas e
atualiza seu histórico. Entrevista os motoristas antes da viagem,
antecipando e resolvendo eventuais problemas. Orienta motoristas e
auxiliares, podendo propor e supervisionar treinamentos para melhoria de
seu desempenho. Pode advertir motoristas e auxiliares, verbalmente e por
escrito, pelos resultados constatados na medição de teor alcoólico e na
análise de tacógrafo ou pela observação de problemas de conduta. Atende
acidentes na estrada, deslocando-se ao local do acidente. Aciona a
polícia e a perícia técnica, podendo preencher o boletim de ocorrência.
Providencia socorro às vítimas. Notifica a seguradora, apresentando a
sequência de fatos em torno da ocorrência, fotografando o local do
acidente e providenciando remoção do veículo do local. Levanta e fornece
informações relacionadas ao acidente à empresa e ao cliente. Aciona o
setor de serviço social e o departamento jurídico da empresa, informando
sobre as circunstâncias do acidente e relatando o número e a condição
das vítimas. Quando possível, atua para dar continuidade à prestação do
serviço, assistindo passageiros não vitimados, recolhendo seus pertences
pessoais e providenciando veículo, motorista e auxiliares para seguir
viagem. Colabora nas atividades relacionadas ao controle da frota,
checando disponibilidade de veículos e motoristas. Verifica condições de
limpeza e higiene dos veículos e, quando necessário, os encaminha para
manutenção. Controla despesas de viagens e envia prestações de contas ao
setor financeiro da empresa. Atende clientes, registrando e encaminhando
suas reclamações e sugestões.

\begin{Form}
\textbf{Esta correspondência está adequada?} \\ 
\ChoiceMenu[radio,radiosymbol=\ding{52},bordercolor=0 0.5 0,name=cbovalid_ 342310 ]{}{CBO deve ficar=sim, \hspace{6em} CBO não fica aqui=nao}
\end{Form}

\newpage
\section{ 06015 -- Técnico em Transporte Aquaviário }

\textbf{Perfil do Curso (CNCT)}

O Técnico em Transporte Aquaviário será habilitado para:

Operar, coordenar e fiscalizar o transporte aquaviário de pessoas e de
cargas. Operar movimentação em terminal, em logística e em navegação.
Coordenar e fiscalizar atividades de prestação de serviços de transporte
aquaviário.

Para atuação como Técnico Aquaviário, são fundamentais:

Conhecimentos e saberes relacionados ao desenvolvimento de operação,
coordenação, supervisão e fiscalização do transporte aquaviário,
incluindo movimentação em terminal, logística e navegação e de apoio
portuário, primando por um elevado grau de responsabilidade social e
ambiental. Aptidão para o trabalho em equipe. Ética, iniciativa,
responsabilidade, respeito às normas técnicas e de segurança, com
reconhecida competência técnica, política.

\textbf{CBOs correspondidos e seus perfis (presentes no CNCT, e presentes na predição da IA)}

\textbf{CBO: 342605  ( Chefe de estação portuária )}

Planeja as atividades de estações portuárias (instalações portuárias -
localizadas dentro ou fora da área do porto organizado - utilizadas em
movimentação de passageiros e/ou de cargas, no transporte aquaviário),
levando em conta os espaços disponíveis, os fluxos, os equipamentos e os
demais recursos instalados. Supervisiona recebimento e embarque de
cargas e passageiros nas estações portuárias, examinando a previsão de
operações de embarque e desembarque e analisando os relatórios das
embarcações. Supervisiona o recebimento de documentação e o processo de
liberação de cargas e passageiros. Supervisiona a equipe de operação.
Organiza distribuição de cargas e passageiros, selecionando tipos e
tamanhos de embarcações de acordo com volume e peso de cargas e
quantidade de passageiros; estipulando destinos para as embarcações; e
calculando horários de saída e de chegada das viagens. Pode definir a
utilização de navios especializados de acordo com o tipo de carga. Pode
organizar a distribuição de cargas com uso de robôs. Coordena serviços
de embarcação em estações portuárias, distribuindo tripulação para
embarcações; determinando horários para as operações; e definindo tipos
de embarcação de acordo com a demanda por serviços. Determina o
acionamento de motores para dar início à movimentação de embarcações.
Providencia serviços de auxílio à tripulação e serviços de auxílio às
operações com embarcações. Define prioridades de atracação de
embarcações, de acordo com a programação das operações, tendo em vista
maior agilidade operacional, com redução do período em que a embarcação
permanece atracada. Nas operações de atracação, desatracação, embarque e
desembarque nas estações portuárias, distribui tarefas a serem cumpridas
na embarcação e/ou na estação; e estabelece ponte de embarque e
desembarque. Monitora o embarque e o desembarque de passageiros. Presta
assessoramento nas ocasiões de visitas de representantes da capitania
dos portos, apresentando as instalações da estação portuária e
disponibilizando as informações solicitadas. Colabora com setor de
tecnologia da empresa no desenvolvimento e na implantação de recursos
para a área de segurança da estação portuária, tais como sistemas de
localização em tempo real, análise inteligente de vídeo, câmeras
térmicas e reconhecimento facial, dentre outros. Participa na
implantação de recursos tecnológicos - tais como ``QR Code'' (sistema de
resposta rápida com código de barras) -, para facilitar embarques e
desembarques em cruzeiros internacionais e outras modalidades de
serviço. Estuda e implanta soluções para redução de perdas e
melhoramento contínuo da produtividade e da segurança, nas operações da
estação portuária. Pode fazer uso de sistemas do tipo ``Big Data''
(grandes conjuntos de dados), analisando informações e empregando
técnicas e critérios definidos, a fim de prever demandas e tendências
que interferem nas atividades das estações portuárias. Prepara a equipe
de trabalho, avaliando desempenho; identificando necessidades de
treinamento; e viabilizando a participação em programas de treinamento.
Elabora documentos administrativos relacionados às atividades da estação
portuária, tais como listas de equipamentos de embarcação, planilhas de
estatísticas operacionais, dentre outros. Encaminha a documentação
elaborada para setores da empresa. Mantém iniciativas integradoras entre
a estação portuária e seu entorno, considerando aspectos urbanísticos,
ambientais e de negócios, bem como a integração com outros modais de
transporte. Zela pela preservação ambiental, colaborando -- no âmbito da
estação portuária -- para a redução de emissões no transporte
aquaviário. Trabalha com segurança, usando equipamentos de proteção
individual e aplicando princípios básicos de ergonomia. Desenvolve
campanhas de prevenção de acidentes e de proteção do meio ambiente.

\begin{Form}
\textbf{Esta correspondência está adequada?} \\ 
\ChoiceMenu[radio,radiosymbol=\ding{52},bordercolor=0 0.5 0,name=cbovalid_ 342605 ]{}{CBO deve ficar=sim, \hspace{6em} CBO não fica aqui=nao}
\end{Form}

\textbf{CBO: 342610  ( Supervisor de operações portuárias )}

Planeja as atividades, considerando espaços disponíveis e fluxos de
pessoas e de cargas nos locais de trabalho. Participa em reuniões com
representantes de autoridades portuárias, agentes e operadores, a fim de
obter subsídios para o planejamento das operações. Apresenta, para as
autoridades portuárias, a previsão de carga e descarga e informa o
período previsto das operações. Informa, ao setor financeiro da empresa,
valores para pagamentos de taxas portuárias. Supervisiona o recebimento
de documentação e o processo de liberação de cargas. Verifica
documentação de exportação e de importação de mercadorias. Controla
distribuição de cargas, monitorando a organização das cargas segundo
peso; por tipo; e por embarcação. Confere a separação de tipos de
contêineres por porto de destino. Seleciona tipos e tamanhos de
embarcações e calcula horários de saída e de chegada no destino. Mapeia
infraestrutura para recebimento de cargas. Utiliza sistemas de automação
na distribuição de cargas e nas demais operações portuárias. Programa
atracação de embarcações, definindo prioridades; agendando atracação de
embarcações; e providenciando equipe de trabalhadores. Organiza serviços
de auxílio à tripulação e serviços de auxílio às operações com
embarcações. Compatibiliza o plano de carga e descarga com os relatórios
operacionais. Avisa autoridades sobre a chegada de embarcações. Examina
previsão de embarque e desembarque, analisando relatórios das
embarcações. Define equipamentos para carga e descarga. Controla o fluxo
de caminhões no porto e nas suas áreas de acesso. Monitora a
movimentação de cargas em embarcação, estabelecendo ponte de embarque e
desembarque; distribuindo tarefas nas instalações portuárias e na
embarcação; definindo equipamentos para carga e descarga; monitorando
desapeação e desembarque de carga; e controlando serviços de embarque e
peação de carga. Colabora nas iniciativas de melhoria da infraestrutura
da logística portuária, apresentando sugestões, fazendo recomendações e
participando no desenvolvimento de soluções. Analisa relatórios das
operações portuárias, utilizando critérios estabelecidos pela empresa, a
fim de promover o aperfeiçoamento permanente das operações. Colabora com
a empresa no desenvolvimento e na implantação de soluções para a área da
segurança nas operações portuárias, tais como sistemas de localização em
tempo real e análise inteligente de vídeo. Colabora em outros
desenvolvimentos que objetivam o aumento da produtividade nos portos,
como aplicações de inteligência artificial e algoritmos de atracação,
dentre outros. Presta assessoramento nas ocasiões de visitas -- de
representantes da capitania dos portos, da receita federal, da polícia
federal e da vigilância sanitária --, apresentando as instalações
portuárias e disponibilizando as informações solicitadas. Relaciona-se
com clientes, apresentando as vantagens operacionais dos serviços sob
sua responsabilidade. Prepara a equipe de trabalho, avaliando seu
desempenho, levantando necessidades de treinamento e viabilizando a
participação dos trabalhadores em programas de treinamento. Elabora
documentos relacionados às operações portuárias, tais como lista de
equipamento das embarcações, relatórios com características das
embarcações, lista de sequenciamento de carga e descarga e relatório de
recebimento de carga. Elabora planilhas de estatísticas operacionais e
planilhas de custo das operações portuárias. Preenche livros de
ocorrências. Encaminha documentação para diversos órgãos, tais como
setores da empresa e autoridades portuárias. Participa dos esforços em
prol da desburocratização nas operações portuárias. Trabalha com
segurança, usando equipamentos de proteção individual e aplicando
princípios básicos de ergonomia. Desenvolve campanhas de prevenção de
acidentes e de proteção do meio ambiente.

\begin{Form}
\textbf{Esta correspondência está adequada?} \\ 
\ChoiceMenu[radio,radiosymbol=\ding{52},bordercolor=0 0.5 0,name=cbovalid_ 342610 ]{}{CBO deve ficar=sim, \hspace{6em} CBO não fica aqui=nao}
\end{Form}

\newpage
\section{ 06016 -- Técnico em Transporte de Cargas }

\textbf{Perfil do Curso (CNCT)}

O Técnico em Transporte de Cargas será habilitado para:

Planejar, executar, coordenar, controlar e fiscalizar as operações de
transporte de cargas. Realizar o acondicionamento e o movimento de
cargas. Realizar o controle de custos e o apoio à gestão operacional.
Coordenar processos de acondicionamento, embalagem e movimentação de
cargas em diferentes modais de transportes. Organizar sistemas de
informação, documentações e arquivos. Colaborar na definição e
negociação de tarifas e na definição e controle de custos de
transportes. Coordenar e fiscalizar atividades de prestação de serviços
no transporte de cargas. Selecionar fornecedores de veículos,
componentes e serviços.

Para atuação como Técnico em Transporte de Carga, são fundamentais:

Conhecimentos e saberes relacionados ao controle de processos de
acondicionamento, embalagem e movimentação de cargas, determinando o
sistema de transportes e de frota, considerando os modais, roteirização
e composição de custos de frete e de negociação. Compromisso e ética com
as questões sociais e de desenvolvimento tecnológico. Formação e
habilidades para a solução de problemas inerentes ao processo produtivo
e à busca de inovações tecnológicas, à liderança de equipes, à solução
de problemas técnicos e trabalhistas e à gestão de conflitos.

\textbf{CBOs correspondidos e seus perfis (presentes no CNCT, e presentes na predição da IA)}

\textbf{CBO: 342305  ( Chefe de serviço de transporte rodoviário (passageiros e cargas) )}

Planeja e controla a disponibilidade da frota para atendimento à demanda
de serviços, determinando a escala de utilização do conjunto de veículos
de transporte rodoviário da empresa. Analisa demanda do mercado para, se
necessário, redimensionar a frota, remanejar os veículos e contratar
serviços de terceiros. Estabelece a programação de viagens, levando em
conta os percursos e as jornadas de trabalho dos motoristas e
auxiliares. Supervisiona o embarque de carga e de passageiros,
monitorando horários de entrada e saída de veículos. Inspeciona
baldeação de bagagem. Emite documento fiscal da carga e libera recursos
financeiros para viagem. Pode entrevistar os motoristas antes da viagem,
antecipando e resolvendo eventuais problemas. Monitora os veículos ao
longo do percurso - inclusive por meio de sistemas de rastreamento -,
para informar-se sobre as ocorrências. Atende acidentes na estrada,
deslocando-se ao local do acidente. Providencia socorro às vítimas.
Aciona polícia e perícia técnica. Relata à empresa o número e o destino
das vítimas e a relação de passageiros, acionando o setor de serviço
social. Informa sobre as circunstâncias do acidente ao departamento
jurídico da empresa. Fornece ao cliente as informações relacionadas ao
acidente. Aciona a seguradora, fotografa o local do acidente e
providencia remoção do veículo do local. Quando possível, atua para dar
continuidade ao serviço, assistindo passageiros não vitimados,
recolhendo seus pertences pessoais e providenciando veículo, motorista e
auxiliares para o prosseguimento da viagem. Realiza o acerto de contas
de viagens com os motoristas, conferindo notas fiscais, controlando o
consumo de combustível e verificando a planilha de despesas de viagens.
Fecha boletins de caixa e encaminha prestação de contas ao setor
financeiro da empresa. Controla custos das viagens. Supervisiona
motoristas e auxiliares de viagem, conferindo sua documentação e
providenciando treinamento. Orienta motoristas e auxiliares,
advertindo-os, se necessário, verbalmente e por escrito. Atende
clientes, registrando e tratando suas reclamações e sugestões. Encaminha
veículos para serviços de manutenção.

\begin{Form}
\textbf{Esta correspondência está adequada?} \\ 
\ChoiceMenu[radio,radiosymbol=\ding{52},bordercolor=0 0.5 0,name=cbovalid_ 342305 ]{}{CBO deve ficar=sim, \hspace{6em} CBO não fica aqui=nao}
\end{Form}

\textbf{CBO: 342310  ( Inspetor de serviços de transportes rodoviários (passageiros e cargas) )}

Inspeciona o embarque, a trajetória e o desembarque de passageiros e
cargas, verificando o cumprimento dos horários previstos, observando a
conduta dos motoristas e registrando as ocorrências durante o percurso.
Na impossibilidade do cumprimento da programação de uma viagem,
encaminha soluções, podendo remanejar veículos e substituir motoristas e
auxiliares. Inspeciona o trabalho dos motoristas e auxiliares de viagem,
examinando seu desempenho. Confere a documentação dos motoristas e
atualiza seu histórico. Entrevista os motoristas antes da viagem,
antecipando e resolvendo eventuais problemas. Orienta motoristas e
auxiliares, podendo propor e supervisionar treinamentos para melhoria de
seu desempenho. Pode advertir motoristas e auxiliares, verbalmente e por
escrito, pelos resultados constatados na medição de teor alcoólico e na
análise de tacógrafo ou pela observação de problemas de conduta. Atende
acidentes na estrada, deslocando-se ao local do acidente. Aciona a
polícia e a perícia técnica, podendo preencher o boletim de ocorrência.
Providencia socorro às vítimas. Notifica a seguradora, apresentando a
sequência de fatos em torno da ocorrência, fotografando o local do
acidente e providenciando remoção do veículo do local. Levanta e fornece
informações relacionadas ao acidente à empresa e ao cliente. Aciona o
setor de serviço social e o departamento jurídico da empresa, informando
sobre as circunstâncias do acidente e relatando o número e a condição
das vítimas. Quando possível, atua para dar continuidade à prestação do
serviço, assistindo passageiros não vitimados, recolhendo seus pertences
pessoais e providenciando veículo, motorista e auxiliares para seguir
viagem. Colabora nas atividades relacionadas ao controle da frota,
checando disponibilidade de veículos e motoristas. Verifica condições de
limpeza e higiene dos veículos e, quando necessário, os encaminha para
manutenção. Controla despesas de viagens e envia prestações de contas ao
setor financeiro da empresa. Atende clientes, registrando e encaminhando
suas reclamações e sugestões.

\begin{Form}
\textbf{Esta correspondência está adequada?} \\ 
\ChoiceMenu[radio,radiosymbol=\ding{52},bordercolor=0 0.5 0,name=cbovalid_ 342310 ]{}{CBO deve ficar=sim, \hspace{6em} CBO não fica aqui=nao}
\end{Form}

\newpage
\section{ 06017 -- Técnico em Transporte Metroferroviário }

\textbf{Perfil do Curso (CNCT)}

O Técnico em Transporte Metroferroviário será habilitado para:

Operar, coordenar e fiscalizar o transporte metroferroviário. Coordenar
a circulação de veículos metroferroviários. Trabalhar no planejamento,
na execução e no controle de atividades ligadas às operações de pátios e
terminais, veículos, sinalização e equipamentos do transporte
metroferroviário. Coordenar a circulação de veículos metroferroviários.
Administrar e controlar as atividades de pátios e terminais. Controlar e
programar os horários de circulação de trens. Operar equipamentos e
sistemas de sinalização, telecomunicações e bordo utilizados nos
sistemas metroferroviários. Manobrar equipamentos e veículos
metroferroviários. Preencher relatórios, planilhas, documentos de
despacho, diário operacional e boletins de ocorrência.

Para atuação como Técnico em Transporte Metroferroviário, são
fundamentais:

Conhecimentos e saberes relacionados ao controle e operação de veículos
em pátios e vias ferroviárias, de modo a assegurar a saúde e a segurança
dos trabalhadores e dos usuários do transporte metroferroviário. Além de
ter o comprometimento com as questões sociais e de desenvolvimento
tecnológico, proatividade, formação para a solução de problemas e a
busca de inovações tecnológicas. Habilidades e saberes relacionados à
liderança de equipes, à solução de problemas técnicos e trabalhistas e à
gestão de conflitos.

\textbf{CBOs correspondidos e seus perfis (presentes no CNCT, e presentes na predição da IA)}

\textbf{CBO: 342110  ( Operador de transporte multimodal )}

Programa, controla e coordena operações de transportes em geral - em
diferentes modais de transporte -, acompanhando as operações de
embarque, transbordo e desembarque de carga junto ao regime de trânsito
aduaneiro. Supervisiona as condições de segurança e de manutenção dos
meios de transportes e equipamentos utilizados. Controla a eficiência
operacional do armazenamento e do transporte da carga. Remaneja
equipamentos e veículos, quando necessário. Controla recursos
financeiros e insumos, elaborando planilhas de custos e receitas dos
serviços prestados. Pesquisa mercado, analisando a viabilidade de novos
serviços, coletando dados de clientes e fornecedores potenciais e
identificando rotas de transporte e respectiva estrutura de serviço.
Define parcerias e fornecedores de serviços com critérios técnicos,
contratando serviços de transporte e negociando preços e prazos de
pagamento. Elabora ou auxilia na elaboração de documentação - como
contratos, registros legais de transporte de cargas do comércio
doméstico ou do exterior - necessária ao desembaraço de cargas. Registra
dados em sistemas e alimenta sistemas informacionais internos. Atende
clientes, classificando-os por demanda, identificando suas necessidades,
elaborando propostas de preços e serviços, disponibilizando informações
da carga e verificando a sua satisfação. Treina e orienta funcionários e
estagiários, avaliando os seus desempenhos e analisando as necessidades
de treinamento da equipe.

\begin{Form}
\textbf{Esta correspondência está adequada?} \\ 
\ChoiceMenu[radio,radiosymbol=\ding{52},bordercolor=0 0.5 0,name=cbovalid_ 342110 ]{}{CBO deve ficar=sim, \hspace{6em} CBO não fica aqui=nao}
\end{Form}

\textbf{CBO: 783110  ( Manobrador )}

Planeja a execução de manobras de trens, examinando a programação;
verificando os tipos de vagões e seus destinos; e identificando bitolas
e linhas. Confere freios manuais e calços. Engata mangueiras e abre
torneira angular de ar. Orienta os maquinistas via rádio ou por sinais
manuais ou luminosos. Prepara o acoplamento de vagões de carga e carros
de passageiros, examinando a bitola das linhas e aproximando a máquina.
Verifica a posição, alinha e conecta engate no trem. Confirma engate e
conecta plugues ou jumpers entre vagões e locomotivas. Verifica calço do
trem. Opera aparelho de mudança de via (AMV), conferindo e destravando a
posição da chave. Aciona chave manual e verifica ponta da agulha. Trava
a chave na posição desejada. Abre o aparelho de mudança de via, limpando
e lubrificando a chave. Efetua o estacionamento do trem dentro do marco,
acionando freio de emergência ou de serviço. Aciona freio manual dos
vagões. Calça o trem. Recarrega freio, se necessário. Comunica ao Centro
de Controle Operacional (CCO) o recolhimento de trens e a situação do
pátio. Inspeciona os veículos ferroviários, vistoriando o estado físico
dos vagões e verificando o fechamento de suas portas ou seus lacres.
Executa pequenos reparos nos vagões. Limpa as locomotivas,
abastecendo-as com areia. Comunica anormalidades à área de manutenção.
Prepara o desacoplamento de vagões e carros, aplicando freio de serviço
de emergência no veículo. Fecha a torneira de ar e desengata as
mangueiras, prendendo-as no suporte. Desacopla plugues ou jumpers e
aciona a alavanca de corte do engate. Trabalha com segurança, utilizando
equipamentos de proteção individual; fazendo sinais manuais ou luminosos
de alerta; verificando os sistemas de freios e a timoneira; e retirando
obstáculos dos trilhos. Orienta as pessoas sobre as áreas de risco,
utilizando rádio transceptor.

\begin{Form}
\textbf{Esta correspondência está adequada?} \\ 
\ChoiceMenu[radio,radiosymbol=\ding{52},bordercolor=0 0.5 0,name=cbovalid_ 783110 ]{}{CBO deve ficar=sim, \hspace{6em} CBO não fica aqui=nao}
\end{Form}

\newpage
\section{ 06018 -- Técnico em Transporte Rodoviário }

\textbf{Perfil do Curso (CNCT)}

O Técnico em Transporte Rodoviário será habilitado para:

Organizar e controlar as operações e executar a logística de tráfego
rodoviário. Planejar, operacionalizar e executar a logística do
transporte de passageiros. Administrar e controlar a frota de veículos
no transporte rodoviário de cargas e passageiros. Executar a operação,
comercialização e manutenção de equipamentos. Planejar a armazenagem e o
processo de expedição das empresas e centros de distribuições. Planejar
e executar a distribuição de pessoal e cargas. Coordenar ações de
intermodalidade de transportes. Identificar as características da malha
viária e os diversos tipos de veículos transportadores. Aplicar a
legislação de trânsito de veículos e de transporte de passageiros.
Preparar e gerenciar a documentação necessária para operações de
transportes.

Para atuação como Técnico em Transporte Rodoviário, são fundamentais:

Conhecimentos e saberes relacionados ao controle de operações e
logística de transportes rodoviários, de passageiros e de cargas, de
modo a assegurar a saúde e a segurança dos trabalhadores e usuários.
Compromisso e ética com as questões ambientais, sociais e de
desenvolvimento tecnológico. Habilidades para a solução de problemas.
Conhecimentos e habilidades relacionados à busca de inovações
tecnológicas, à liderança de equipes, à solução de problemas técnicos e
trabalhistas e à gestão de conflitos.

\textbf{CBOs correspondidos e seus perfis (presentes no CNCT, e presentes na predição da IA)}

\textbf{CBO: 342305  ( Chefe de serviço de transporte rodoviário (passageiros e cargas) )}

Planeja e controla a disponibilidade da frota para atendimento à demanda
de serviços, determinando a escala de utilização do conjunto de veículos
de transporte rodoviário da empresa. Analisa demanda do mercado para, se
necessário, redimensionar a frota, remanejar os veículos e contratar
serviços de terceiros. Estabelece a programação de viagens, levando em
conta os percursos e as jornadas de trabalho dos motoristas e
auxiliares. Supervisiona o embarque de carga e de passageiros,
monitorando horários de entrada e saída de veículos. Inspeciona
baldeação de bagagem. Emite documento fiscal da carga e libera recursos
financeiros para viagem. Pode entrevistar os motoristas antes da viagem,
antecipando e resolvendo eventuais problemas. Monitora os veículos ao
longo do percurso - inclusive por meio de sistemas de rastreamento -,
para informar-se sobre as ocorrências. Atende acidentes na estrada,
deslocando-se ao local do acidente. Providencia socorro às vítimas.
Aciona polícia e perícia técnica. Relata à empresa o número e o destino
das vítimas e a relação de passageiros, acionando o setor de serviço
social. Informa sobre as circunstâncias do acidente ao departamento
jurídico da empresa. Fornece ao cliente as informações relacionadas ao
acidente. Aciona a seguradora, fotografa o local do acidente e
providencia remoção do veículo do local. Quando possível, atua para dar
continuidade ao serviço, assistindo passageiros não vitimados,
recolhendo seus pertences pessoais e providenciando veículo, motorista e
auxiliares para o prosseguimento da viagem. Realiza o acerto de contas
de viagens com os motoristas, conferindo notas fiscais, controlando o
consumo de combustível e verificando a planilha de despesas de viagens.
Fecha boletins de caixa e encaminha prestação de contas ao setor
financeiro da empresa. Controla custos das viagens. Supervisiona
motoristas e auxiliares de viagem, conferindo sua documentação e
providenciando treinamento. Orienta motoristas e auxiliares,
advertindo-os, se necessário, verbalmente e por escrito. Atende
clientes, registrando e tratando suas reclamações e sugestões. Encaminha
veículos para serviços de manutenção.

\begin{Form}
\textbf{Esta correspondência está adequada?} \\ 
\ChoiceMenu[radio,radiosymbol=\ding{52},bordercolor=0 0.5 0,name=cbovalid_ 342305 ]{}{CBO deve ficar=sim, \hspace{6em} CBO não fica aqui=nao}
\end{Form}

\newpage
\section{ 07010 -- Técnico em Guarda e Segurança }

\textbf{Perfil do Curso (CNCT)}

O Técnico em Guarda e Segurança será habilitado, de acordo com cada
Força, para:

Executar serviços de segurança das instalações militares. Executar
atividades de implantação de sistemas de vigilância. Operar equipamentos
de comunicação dotados de Medidas de Proteção Eletrônica. Empregar a
doutrina de Infantaria da Aeronáutica nas ações de Força Aérea. Aplicar
técnicas e procedimentos de autodefesa de superfície de instalações
militares, utilizando os recursos de Polícia, Autodefesa Antiaérea,
Proteção da Força e Operações Especiais. Participar de operações
militares de ajuda humanitária no Brasil e no exterior.

\textbf{CBOs correspondidos e seus perfis (presentes no CNCT, e presentes na predição da IA)}

\textbf{CBO: 517310  ( Agente de segurança )}

Fiscaliza e controla a movimentação e a permanência de pessoas e
veículos, em área de acesso livre ou restrito, como medida de segurança.
Recepciona pessoas, identificando-as e consultando a pessoa a ser
visitada, conforme normas ou regulamentos, para permitir ou impedir o
acesso às dependências de empresa ou instituição. Aborda, encaminha e
acompanha visitante, controlando a sua movimentação. Auxilia idosos e
pessoas com deficiência. Recepciona autoridades e requisita transporte.
Controla e acompanha a entrega de bens e mercadorias dentro do local de
trabalho, verificando a documentação e impedindo a entrada de objetos
ilícitos ou irregulares. Realiza rondas no prédio e imediações,
monitorizadas ou a pé, inspecionado dependências, examinando portas,
janelas, portões e alarmes, para constatar eventuais anormalidades.
Resgata pessoas acidentadas, prestando-lhes primeiros socorros e
providenciando socorro médico. Desenergiza linha metro-ferroviária para
proteção do usuário, em situação de emergência. Pode combater pequenos
incêndios. Fotografa ocorrências. Comunica à chefia imediata qualquer
irregularidade ocorrida durante seu plantão, para que sejam tomadas as
devidas providências. Comunica, à autoridade competente, as ocorrências
de focos de incêndio, de uso e tráfico de substâncias tóxicas, de
tentativa de furto, de atos obscenos, de vandalismo e de outros delitos.
Na execução de tarefas pertinentes à área de atuação, pode utilizar
equipamentos e programas de informática. Comunica-se via
radiocomunicador, telefone ou outro equipamento e opera equipamento de
vigilância eletrônica, utilizando-se de gestos, sinais e códigos de
segurança. Pode prestar informações ao público sobre normas de
segurança.

\begin{Form}
\textbf{Esta correspondência está adequada?} \\ 
\ChoiceMenu[radio,radiosymbol=\ding{52},bordercolor=0 0.5 0,name=cbovalid_ 517310 ]{}{CBO deve ficar=sim, \hspace{6em} CBO não fica aqui=nao}
\end{Form}

\newpage
\section{ 07015 -- Técnico em Mecânica de Aeronaves  }

\textbf{Perfil do Curso (CNCT)}

O Técnico em Mecânica de Aeronaves será habilitado para:

Atuar na manutenção de aeronaves e em seus equipamentos. Executar
inspeções em motores de aviões e helicópteros e nos seguintes sistemas:
hélice, hidráulica, pneumática, combustível, comandos de voo em aviões e
helicópteros, radiocomunicação, radionavegação e combate à corrosão,
conforme especificações e normas técnicas. Interpretar manuais técnicos
das diferentes aeronaves e equipamentos. Atuar como mecânico de voo.
Exercer tarefas alinhadas às atividades de estrutura e pintura de
aeronaves e de eletricidade e instrumentos aeronáuticos.

\textbf{CBOs correspondidos e seus perfis (presentes no CNCT, e presentes na predição da IA)}

\textbf{CBO: 314310  ( Técnico mecânico (aeronaves) )}

Prepara ensaios de fabricação de equipamentos, estruturas e sistemas
mecânicos de aviões e helicópteros, verificando cronograma do processo
de fabricação e interpretando desenhos de projeto, para elaborar croquis
e fabricar ou adaptar peças, ferramentas e dispositivos a serem
utilizados nos ensaios. Realiza ensaios de fabricação de equipamentos,
estruturas e sistemas mecânicos de aviões e helicópteros, abrangendo
fuselagem, montantes, naceles, capotas de motor, carenagens, rotores e
outras superfícies aerodinâmicas de célula, trens de aterrissagem,
sistemas de freio e demais equipamentos mecânicos da aeronave, incluindo
seus acessórios e controles. Identifica necessidade de ajustes
dimensionais, adapta peças e acessórios e torna modificações adequadas
às características da aeronave. Participa dos testes finais de ensaios e
montagem. Prepara a montagem de componentes em aviões e helicópteros,
interpretando desenhos técnicos, verificando a disponibilidade de itens
no estoque para as montagens requeridas e analisando condições do local
de instalação e dos componentes, para realizar o preparo do local.
Executa, seguindo as instruções de trabalho, a montagem de componentes
em bancada e os instala em sistemas da aeronave de acordo com o projeto,
utilizando ferramentas, instrumentos e equipamentos. Elimina
interferências prejudiciais à segurança, assegurando ajustes entre peças
e destas com a estrutura. Realiza testes funcionais de aviões e
helicópteros, cumprindo procedimentos, seguindo parâmetros estabelecidos
pelo fabricante e utilizando equipamentos específicos e calibrados.
Executa testes iniciais para confirmação dos problemas apresentados,
testa sistemas operacionais mecânicos, de estrutura da aeronave e de
aviônica aplicada ao controle de equipamentos e sistemas mecânicos.
Simula processos, realiza análise dimensional por etapas de montagem e
participa da execução de provas de inclinação da aeronave. Libera
processos para itens de segurança. Prepara a execução de manutenção de
estruturas, equipamentos e sistemas mecânicos de aviões e helicópteros,
identificando sua forma de utilização, seus defeitos e suas falhas.
Executa os reparos, utilizando ferramentas e equipamentos apropriados,
substituindo e regulando componentes e utilizando técnicas de proteção
química e de acabamento de superfície. Nivela a linha de eixo do motor
da aeronave. Presta assessoria técnica interna e externa, treinando
equipes de trabalho, acompanhando o processo de montagem e manutenção,
provendo suporte técnico a equipes de trabalho e clientes, realizando
atendimento de emergência e fazendo visita técnica. Analisa resultados
da assessoria prestada, com a finalidade de assegurar a satisfação dos
clientes. Participa do programa de qualidade da empresa, seguindo as
normas internas e as definidas por organismos externos competentes,
utilizando ferramentas da qualidade, seguindo procedimentos de serviço e
identificando pontos críticos para propor melhorias, tendo em vista o
alcance de metas de qualidade. Fornece informações técnicas para
execução do controle estatístico do processo. Elabora documentação
técnica. Faz orçamento de serviços, com estimativa de tempo e de custo.
Emite parecer técnico. Elabora plano de ação referente às
não-conformidades detectadas em produtos e processos. Elabora relatório
técnico e de campo, identificando os componentes instalados e
registrando a não-conformidade relacionada ao componente substituído ou
reparado. Documenta serviços concluídos, fornecendo registros para
liberação das fases de processos e do programa de qualidade aos
clientes. Emite documento para liberação final do produto. Conserva o
local de trabalho limpo e organizado. Descarta resíduos seguindo normas
ambientais. Trabalha com segurança, utilizando Equipamentos de Proteção
Individual (EPI) e Equipamentos de Proteção Coletiva (EPC). Participa de
programas de prevenção de acidentes, colaborando nas ações e propondo
melhorias nas condições de segurança.

\begin{Form}
\textbf{Esta correspondência está adequada?} \\ 
\ChoiceMenu[radio,radiosymbol=\ding{52},bordercolor=0 0.5 0,name=cbovalid_ 314310 ]{}{CBO deve ficar=sim, \hspace{6em} CBO não fica aqui=nao}
\end{Form}

\textbf{CBO: 914105  ( Mecânico de manutenção de aeronaves, em geral )}

Planeja o trabalho, verificando ordens de serviço. Interpreta manuais,
desenhos técnicos mecânicos e elétricos de motopropulsores, estruturas,
sistemas diversificados da aeronave e componentes aviônicos, de acordo
com o tipo -- aviões e helicópteros, de grande, médio e pequeno porte --
de aeronave. Analisa ``diretivas de aeronavegabilidade'' -- notificações
ou avisos de que existem deficiências de segurança conhecidas,
associadas às aeronaves, que podem estar relacionadas à manutenção --,
para cumprir diretrizes de segurança aeronáutica. Utiliza sistemas de
informação de apoio à manutenção. Identifica novos materiais utilizados
na construção de peças de motores e de estruturas. Repara motores
convencionais e motores a reação (a jato), removendo-os, corrigindo
defeitos, instalando acessórios, executando ensaios não destrutivos e
realizando revisão geral. Balanceia e instala motores. Coleta amostras
de fluidos para análise laboratorial. Pode reparar aeronaves com
motorização elétrica. Repara estruturas de aeronaves, inspecionando-as a
fim de detectar defeitos e verificando fixação de rebites. Repara
chapeamento, aplica materiais compósitos e trata estruturas contra
corrosão. Alinha superfícies em relação à fuselagem. Repara sistema de
combustível, examinando tipo de combustível, verificando nível e pressão
do sistema e realizando destanqueamento (retirada de combustível dos
tanques). Remove componentes; repara tanques; descontaminha o sistema;
substitui componentes; e instala tanques. Realiza revisão geral de
componentes e abastece os tanques. Realiza manutenção do sistema de trem
de pouso, removendo rodas, pneus, trens de pouso, amortecedores,
atuadores e sistema de freios; revisando sistema de emergência e sistema
de freio automático; montando e ajustando componentes. Lubrifica
componentes mecânicos e realiza testes operacionais em componentes do
trem de pouso. Executa manutenção de sistemas elétricos e eletrônicos,
verificando acionamento elétrico dos diversos sistemas; reparando e
instalando componentes. Faz manutenção do sistema de comandos de voo,
inspecionando funcionamento de componentes, removendo comandos e
balanceando superfícies de comando. Regula e instala comandos de voo.
Realiza manutenção do sistema interior de aeronaves, substituindo
componentes de poltronas e revisando estado geral de poltronas e de
cintos de segurança. Repara componentes do sistema da copa (``galley'').
Conserta sistemas de som, vídeo, iluminação e comunicação. Substitui
pisos e trilhos e repara componentes das toaletes. Executa manutenção em
sistemas diversificados, cumprindo rotinas de inspeção de manutenção e
executando serviços de substituição de peças ou equipamentos. Substitui
peças dos sistemas de oxigênio; ar-condicionado; pressurização; água
potável; e de detritos. Substitui equipamentos de emergência. Substitui
peças do sistema antichuva e antigelo, bem como do sistema de detecção e
extinção de fogo. Pode realizar manutenção de sistemas hidráulicos.
Repara sistemas de hélices e de rotores de helicópteros, removendo-os e
substituindo engrenagens, selos mecânicos e rolamentos. Balanceia e
monta os sistemas. Ajusta conjunto de sincronismo de rotores. Faz testes
de componentes da aeronave, de acordo com procedimentos específicos, a
fim de comprovar seu funcionamento. Testa funcionamento de motores; de
sistemas de alimentação e geração de energia elétrica; e de sistema
antiderrapagem. Realiza testes de freio do rotor de helicópteros e
testes de funcionamento de rotores e hélices de helicópteros. Preenche
relatórios de avarias e fichas de inspeção. Mantém o local de trabalho
limpo e organizado. Descarta resíduos seguindo normas ambientais.
Trabalha com segurança, aplicando conceitos de ergonomia e utilizando
Equipamentos de Proteção Individual (EPI). Sinaliza áreas de manutenção
e delimita áreas de risco. Opera equipamentos de extinção de incêndio.
Orienta balizamento de aeronaves. Realiza aterramento elétrico da
aeronave em manutenção.

\begin{Form}
\textbf{Esta correspondência está adequada?} \\ 
\ChoiceMenu[radio,radiosymbol=\ding{52},bordercolor=0 0.5 0,name=cbovalid_ 914105 ]{}{CBO deve ficar=sim, \hspace{6em} CBO não fica aqui=nao}
\end{Form}

\textbf{CBO: 953105  ( Eletricista de instalações (aeronaves) )}

Planeja serviços de instalação e de manutenção eletroeletrônica em
componentes e equipamentos auxiliares de aeronaves, interpretando ordens
de serviços, elaborando orçamentos e estabelecendo cronograma de
atividades. Estima prazos e necessidade de mão de obra para execução dos
serviços. Faz a especificação de materiais e componentes
eletroeletrônicos. Seleciona ferramentas, equipamentos e materiais.
Instala componentes e equipamentos eletroeletrônicos auxiliares em
aeronaves, interpretando esquemas eletroeletrônicos, elaborando leiaute
das instalações e selecionando equipamentos, instrumentos e ferramentas
apropriados aos serviços. Instala a infraestrutura de suporte -- calhas,
suportes, eletrodutos, conduítes e outros elementos --, de acordo com
especificações. Fixa -- manualmente -- chicotes, cabeamento em geral,
equipamentos e acessórios. Conecta cabos aos equipamentos e acessórios.
Libera equipamentos para teste de funcionamento. Testa o funcionamento
de equipamentos e sistemas para operação. Pode elaborar ou corrigir
esquemas eletroeletrônicos das instalações. Realiza medições elétricas
de tensão, de resistência e de corrente; verifica anomalias nas
instalações -- tais como ruídos, vibrações, vazamentos, correntes
elétricas de fuga, entre outras --; e realiza o acabamento das
instalações eletroeletrônicas de componentes e equipamentos auxiliares
de aeronaves, observando especificações técnicas e seguindo padrões de
qualidade. Realiza manutenção eletroeletrônica -- preventiva, preditiva
e corretiva -- em componentes e equipamentos auxiliares de aeronaves,
listando os bens incluídos nas rotinas de manutenção, elaborando
cronograma de manutenção preventiva e realizando a inspeção visual de
manutenção de máquinas e equipamentos, para coleta de informações e
planejamento do processo. Lista itens de verificação dos sistemas
eletroeletrônicos. Limpa equipamentos e o local de trabalho. Realiza o
diagnóstico de defeitos, utilizando ferramentas, instrumentos ou
aparelhos de diagnóstico, para preparar as intervenções de manutenção.
Desmonta componentes eletroeletrônicos, para permitir a realização de
reparação ou de ajustes nos componentes. Monta os componentes após as
intervenções de manutenção. Substitui componentes eletroeletrônicos em
aeronaves, quando necessário. Faz ajuste em peças de suporte ou
relacionadas às instalações elétricas das aeronaves. Lubrifica
componentes. Simula funcionamento de componentes e de equipamentos
eletroeletrônicos. Busca manter-se atualizado no que se refere aos
desenvolvimentos em aviônica, interpretando literatura técnica
especializada e informações de fabricantes, a fim de realizar -- com
qualidade -- os serviços de instalação e de manutenção eletroeletrônica
em componentes e equipamentos auxiliares de aeronaves. Elabora
documentação relacionada aos trabalhos de instalação e de manutenção
eletroeletrônica em aeronaves, emitindo laudos técnicos e registrando
ocorrências e anomalias. Elabora relatórios de serviço e comunicados de
indisponibilidade de equipamentos. Preenche ordens de serviço e
requisições de materiais. Cumpre normas de segurança, meio ambiente e
saúde no trabalho, utilizando Equipamentos de Proteção Individual (EPI)
e Equipamentos de Proteção Coletiva (EPC); solicitando inspeções nos
locais de trabalho; e sinalizando equipamentos ou locais de trabalho
para isolamento. Interpreta laudos de segurança, meio ambiente e saúde
no trabalho. Deposita resíduos em locais apropriados.

\begin{Form}
\textbf{Esta correspondência está adequada?} \\ 
\ChoiceMenu[radio,radiosymbol=\ding{52},bordercolor=0 0.5 0,name=cbovalid_ 953105 ]{}{CBO deve ficar=sim, \hspace{6em} CBO não fica aqui=nao}
\end{Form}

\newpage
\section{ 08001 -- Técnico em Agroindústria }

\textbf{Perfil do Curso (CNCT)}

O Técnico em Agroindústria será habilitado para:

Aplicar tecnologias voltadas à conservação e ao processamento das
matérias-primas de origem animal e vegetal nas agroindústrias e extensão
rural. Realizar a implantação, a execução e a avaliação de programas
preventivos de segurança do trabalho, de gestão de resíduos, de
diminuição do impacto ambiental e de higienização e sanitização da
produção agroindustrial. Realizar análises laboratoriais de alimentos.
Desenvolver técnicas mercadológicas de produtos e insumos para a
agroindústria e promover a inovação tecnológica.

Para atuação como Técnico em Agroindústria, são fundamentais:

Conhecimentos e saberes relacionados às tecnologias de processamento e
conservação de alimentos, aos programas de controle de qualidade, à
gestão de resíduos, à proatividade, à liderança, à capacidade de
trabalho em equipes e à inovação tecnológica.

\textbf{CBOs correspondidos e seus perfis (presentes no CNCT, e presentes na predição da IA)}

\textbf{CBO: 818110  ( Auxiliar de laboratório de análises físico-químicas )}

Auxilia na realização de análises físico-químicas no laboratório,
interpretando as ordens de serviço. Planeja as atividades, selecionando
método de análise, programando as etapas de trabalho e relacionando
materiais, equipamentos e instrumentos necessários. Verifica o
suprimento de materiais. Efetua cálculos, preenchendo fichas e
formulários. Prepara vidrarias para serem utilizadas em análises
físico-químicas, lavando-as, secando-as e esterilizando-as, a fim de
evitar contaminações. Identifica e armazena as vidrarias. Prepara as
soluções para as análises laboratoriais, selecionando vidrarias e
equipamentos e pesando ou medindo o volume dos reagentes. Mistura os
reagentes, para homogeneizar a solução preparada. Padroniza as soluções
e mede os parâmetros necessários para a sua caracterização. Armazena as
soluções, rotulando os frascos com as informações para a sua
identificação, seguindo as boas práticas de laboratório. Interpreta
manual de operações e prepara equipamentos de medição e ensaios,
identificando tensão elétrica. Seleciona e monta acessórios. Verifica a
aferição dos equipamentos. Coleta amostras para análise laboratorial,
interpretando as instruções do plano de amostragem. Etiqueta as amostras
e registra os seus dados, realizando os procedimentos para preservação
das características das amostras. Prepara as amostras para a análise,
organizando a bancada de trabalho, separando reagentes e soluções,
selecionando meios de cultura e outros insumos. Realiza análises,
interpretando dados. Emite laudos. Mantém arquivos físicos e/ou digitais
de fichas, formulários e laudos. Organiza literaturas técnicas, para
consulta. Conserva o ambiente limpo, higienizado e organizado.
Acondiciona e etiqueta materiais químicos e amostras para armazenamento.
Executa o descarte dos resíduos de acordo com as normas ambientais.
Limpa e organiza acessórios e instrumentos de trabalho. Solicita serviço
de manutenção de máquinas e equipamentos, quando necessário. Trabalha
com segurança, prevenindo acidentes e usando equipamentos de proteção
individual. Presta atenção aos rótulos de reagentes e obedece às
especificidades neles descritas.

\begin{Form}
\textbf{Esta correspondência está adequada?} \\ 
\ChoiceMenu[radio,radiosymbol=\ding{52},bordercolor=0 0.5 0,name=cbovalid_ 818110 ]{}{CBO deve ficar=sim, \hspace{6em} CBO não fica aqui=nao}
\end{Form}

\newpage
\section{ 08002 -- Técnico em Alimentos }

\textbf{Perfil do Curso (CNCT)}

O Técnico em Alimentos será habilitado para:

Coordenar, conduzir, dirigir e executar o processamento e a conservação
de matérias-primas, ingredientes, produtos e subprodutos da indústria
alimentícia e de bebidas, da agroindústria e do comércio de alimentos.
Realizar análises físico-químicas, microbiológicas e sensoriais de
controle de processos. Implantar e coordenar procedimentos de segurança
de alimentos em programas de garantia e controle da qualidade.
Supervisionar a instalação e a manutenção de equipamentos, controlando e
corrigindo desvios nos processos manuais, automatizados e indústria 4.0.
Aplicar soluções tecnológicas para aumentar a produtividade e
desenvolver produtos e processos. Responsabilizar-se pela elaboração e
execução de projetos. Promover assistência técnica na compra, venda e
utilização de produtos, equipamentos e maquinários.

Para atuação como Técnico em Alimentos, são fundamentais:

Conhecimentos e saberes relacionados ao processamento e à conservação de
matérias-primas, produtos e subprodutos da indústria alimentícia e de
bebidas. Conhecimentos e saberes relacionados às transformações
químicas, bioquímicas e físicas dos alimentos, à realização de análises
laboratoriais e sensoriais, à gestão de sistemas de controle, garantia
da qualidade e segurança de alimentos, à visão global dos processos de
produção manual, automatizado e indústria 4.0, à responsabilidade
técnica, às normas técnicas. Conhecimentos e saberes relacionados à
liderança de equipes e tomada de decisões, à capacidade de adaptação a
novos ambientes e situações, à atitude profissional, à postura ética, à
proatividade, à solução de problemas técnicos e trabalhistas e à gestão
de conflitos.

\textbf{CBOs correspondidos e seus perfis (presentes no CNCT, e presentes na predição da IA)}

\textbf{CBO: 325205  ( Técnico de alimentos )}

Planeja as rotinas de trabalho, interpretando as ordens de serviço;
elaborando o cronograma de produção; especificando e calculando os
materiais e os insumos; selecionando os procedimentos para cada
atividade; e verificando a disponibilidade dos equipamentos, materiais,
insumos e gêneros alimentícios. Supervisiona os processos de produção,
recebendo insumos, produtos e gêneros alimentícios; acompanhando
pré-preparo e preparo de alimentos; monitorando processos -- tais como
de trituração, pasteurização, mistura, fermentação e cocção -; e
assegurando as condições operacionais do processo produtivo. Corrige
desvios nos processos manuais e automatizados. Controla o tempo de
produção. Controla o peso e a dimensão do produto. Verifica as condições
da embalagem do produto final. Acompanha a distribuição de produtos até
os pontos de venda. Avalia a qualidade de produtos acabados, coletando
amostras dos produtos e realizando análises microbiológicas, sensoriais
e físico-químicas de matérias-primas, insumos e produtos intermediários
e acabados. Realiza testes de desempenho de matérias-primas, insumos e
produtos intermediários e acabados. Controla a qualidade nas etapas de
produção, identificando e corrigindo pontos críticos de controle e
assegurando as condições higiênico-sanitárias. Verifica as condições de
armazenamento, controlando a data de vencimento dos produtos e
acompanhando o controle integrado de pragas e vetores. Participa, sob
supervisão, de pesquisas e de projetos para melhoria, adequação e
desenvolvimento de produtos, acompanhando as necessidades do mercado e
participando na elaboração da formulação e na testagem dos novos
produtos. Assessora a implementação das mudanças aprovadas. Participa na
elaboração da informação nutricional, que consta na rotulagem do
produto. Orienta clientes internos e externos sobre aspectos
nutricionais dos produtos. Coordena e supervisiona equipes de trabalho,
participando na seleção de pessoal, orientando as atividades,
remanejando pessoal, avaliando os resultados de desempenho e realizando
treinamento de rotinas operacionais. Controla as condições do ambiente
de trabalho, monitorando a temperatura e a umidade do ar. Acompanha os
procedimentos de sanitização das instalações. Verifica condições de
segurança ambiental e controla os pontos críticos e o uso de
equipamentos de proteção individual. Monitora o funcionamento de
máquinas, equipamentos e instrumentos, solicitando serviço de
manutenção, quando necessário. Controla desperdícios e destina
adequadamente os resíduos gerados, fazendo reciclagem e realizando
descarte de acordo com as normas ambientais.

\begin{Form}
\textbf{Esta correspondência está adequada?} \\ 
\ChoiceMenu[radio,radiosymbol=\ding{52},bordercolor=0 0.5 0,name=cbovalid_ 325205 ]{}{CBO deve ficar=sim, \hspace{6em} CBO não fica aqui=nao}
\end{Form}

\textbf{CBO: 818110  ( Auxiliar de laboratório de análises físico-químicas )}

Auxilia na realização de análises físico-químicas no laboratório,
interpretando as ordens de serviço. Planeja as atividades, selecionando
método de análise, programando as etapas de trabalho e relacionando
materiais, equipamentos e instrumentos necessários. Verifica o
suprimento de materiais. Efetua cálculos, preenchendo fichas e
formulários. Prepara vidrarias para serem utilizadas em análises
físico-químicas, lavando-as, secando-as e esterilizando-as, a fim de
evitar contaminações. Identifica e armazena as vidrarias. Prepara as
soluções para as análises laboratoriais, selecionando vidrarias e
equipamentos e pesando ou medindo o volume dos reagentes. Mistura os
reagentes, para homogeneizar a solução preparada. Padroniza as soluções
e mede os parâmetros necessários para a sua caracterização. Armazena as
soluções, rotulando os frascos com as informações para a sua
identificação, seguindo as boas práticas de laboratório. Interpreta
manual de operações e prepara equipamentos de medição e ensaios,
identificando tensão elétrica. Seleciona e monta acessórios. Verifica a
aferição dos equipamentos. Coleta amostras para análise laboratorial,
interpretando as instruções do plano de amostragem. Etiqueta as amostras
e registra os seus dados, realizando os procedimentos para preservação
das características das amostras. Prepara as amostras para a análise,
organizando a bancada de trabalho, separando reagentes e soluções,
selecionando meios de cultura e outros insumos. Realiza análises,
interpretando dados. Emite laudos. Mantém arquivos físicos e/ou digitais
de fichas, formulários e laudos. Organiza literaturas técnicas, para
consulta. Conserva o ambiente limpo, higienizado e organizado.
Acondiciona e etiqueta materiais químicos e amostras para armazenamento.
Executa o descarte dos resíduos de acordo com as normas ambientais.
Limpa e organiza acessórios e instrumentos de trabalho. Solicita serviço
de manutenção de máquinas e equipamentos, quando necessário. Trabalha
com segurança, prevenindo acidentes e usando equipamentos de proteção
individual. Presta atenção aos rótulos de reagentes e obedece às
especificidades neles descritas.

\begin{Form}
\textbf{Esta correspondência está adequada?} \\ 
\ChoiceMenu[radio,radiosymbol=\ding{52},bordercolor=0 0.5 0,name=cbovalid_ 818110 ]{}{CBO deve ficar=sim, \hspace{6em} CBO não fica aqui=nao}
\end{Form}

\newpage
\section{ 08003 -- Técnico em Cervejaria }

\textbf{Perfil do Curso (CNCT)}

O Técnico em Cervejaria será habilitado para:

Coordenar e supervisionar atividades de fabricação de cervejas.
Desenvolver técnicas mercadológicas de produtos e insumos para a
indústria cervejeira. Coordenar a aquisição e a manutenção de
equipamentos e insumos. Promover a inovação tecnológica e o
desenvolvimento de novos produtos. Controlar e corrigir desvios nos
processos manuais e automatizados. Utilizar boas práticas de fabricação,
de rotulagem e de identificação de embalagens adequadas. Efetuar o
controle de qualidade. Realizar análises químicas, físicas, biológicas e
sensoriais. Planejar e executar o processo de trabalho, de comércio e de
venda de cervejas. Supervisionar o tratamento e o destino adequado de
resíduos e efluentes. Operar software para controle do processo
cervejeiro.

Para atuação como Técnico em Cervejaria, são fundamentais:

Conhecimentos e saberes relacionados aos processos manuais e
automatizados de produção de cerveja e à operacionalização de software
para controle do processo. Conhecimentos e saberes relacionados à gestão
de resíduos da indústria cervejeira, à gestão na indústria cervejeira,
ao empreendedorismo e à inovação. Conhecimentos e saberes relacionados
às normas técnicas de qualidade da indústria cervejeira e da cerveja.
Conhecimentos e saberes relacionados à execução de análises de qualidade
da cerveja. Conhecimentos e saberes relacionados à solução de problemas,
à criatividade e ao trabalho em equipe.

\textbf{CBOs correspondidos e seus perfis (presentes no CNCT, e presentes na predição da IA)}

\textbf{CBO: 841710  ( Filtrador de cerveja )}

Prepara-se para a execução da filtragem de cerveja, interpretando ordem
de fabricação. Seleciona os equipamentos de filtragem. Realiza a limpeza
interna e externa dos equipamentos selecionados. Acompanha os processos
de moagem e maceração do malte e separação do mosto, para dar início ao
seu trabalho. Executa filtragem do mosto, abastecendo tanque com placa
perfurada que funciona como leito filtrante, retendo no fundo a fração
sólida - o bagaço do malte - e deixando passar a fração líquida,
bombeada ao tanque de fervura. Depois da fervura, efetua a segunda
filtração do mosto, em tanque circular, onde o mosto entra
tangencialmente em alta velocidade. Submete o mosto a um processo de
decantação hidrodinâmica, separando as proteínas e outras partículas por
efeito centrífugo. Retira o resíduo sólido, denominado trub grosso.
Encaminha o mosto limpo para as próximas etapas de produção
(resfriamento, fermentação e maturação). Na etapa pós-maturação,
participa na execução da clarificação da cerveja em filtros de terra
diatomácea, seja do tipo filtro de placas horizontais; filtro de placas
verticais; filtro de placa e suporte; ou filtro de vela. Monitora, em
qualquer tipo de filtro, as partículas em suspensão da cerveja ficarem
retidas nos interstícios, entre as partículas de terra. Faz a limpeza do
filtro, quando a camada filtrante atingir espessura preestabelecida.
Controla o tempo do processo, a temperatura e a pressão. Realiza o
bombeamento da cerveja filtrada para os processos finais (pasteurização
e envase) de fabricação. Nas atividades de filtragem, controla os
parâmetros dos processos. Pode operar sistema de filtragem de mistura
totalmente automatizado e em rede. Preenche documentos e formulários com
registros dos processos de filtragem. Conserva o local de trabalho limpo
e higienizado. Executa a manutenção básica (ou de primeiro nível) de
equipamentos, fazendo sua limpeza, efetuando lubrificação e realizando
pequenos reparos. Monitora o funcionamento dos equipamentos, para
requisitar serviços de manutenção corretiva quando detectar falhas e
defeitos. Contribui na redução dos impactos ambientais do processo
produtivo, seguindo procedimentos de reaproveitamento, como a separação,
na primeira etapa de filtragem, da parte sólida retida (bagaço de
malte), que pode ser utilizada para a fabricação de ração animal, quando
acrescida de outros componentes. Trabalha com segurança, utilizando
equipamentos de proteção individual.

\begin{Form}
\textbf{Esta correspondência está adequada?} \\ 
\ChoiceMenu[radio,radiosymbol=\ding{52},bordercolor=0 0.5 0,name=cbovalid_ 841710 ]{}{CBO deve ficar=sim, \hspace{6em} CBO não fica aqui=nao}
\end{Form}

\newpage
\section{ 08004 -- Técnico em Confeitaria }

\textbf{Perfil do Curso (CNCT)}

O Técnico em Confeitaria será habilitado para:

Planejar e desenvolver produções de confeitaria, de forma artesanal ou
industrializada, utilizando equipamentos, utensílios e tecnologias
aplicadas aos processos, conforme as boas práticas de manipulação de
alimentos. Elaborar e padronizar fichas técnicas e de controle
operacional para diferentes tipos de operação, promovendo a inovação e o
desenvolvimento de novos produtos e processos. Utilizar as boas práticas
na manipulação de alimentos, rotulagem e identificação da embalagem
adequada. Efetuar controle de qualidade, de estoque, de custos e de
consumo. Utilizar técnicas mercadológicas de produtos e insumos para a
confeitaria. Planejar e executar a aquisição e a manutenção preventiva
de equipamentos.

Para atuação como Técnico em Confeitaria, são fundamentais:

Conhecimentos e saberes relacionados às tecnologias de processamento e
conservação de produções de confeitaria, à implementação de análises
laboratoriais e sensoriais, à gestão dos sistemas de controle e garantia
da qualidade e da segurança de alimentos, à visão global dos processos
de produção manual e automatizado, à responsabilidade técnica, às normas
técnicas, à gestão sustentável, à proatividade, à liderança, à
capacidade de trabalho em equipes e à inovação tecnológica.

\textbf{CBOs correspondidos e seus perfis (presentes no CNCT, e presentes na predição da IA)}

\textbf{CBO: 848310  ( Confeiteiro )}

Prepara as atividades, de acordo com ficha técnica de produção e
programação da produção. Seleciona utensílios e instrumentos, de acordo
com as especificações do processo. Inspeciona as condições operacionais
de máquinas e equipamentos. Interpreta receitas e procedimentos de
execução, verificando as proporções de ingredientes requeridas. Faz a
seleção dos ingredientes e o preparo de massas, recheios, coberturas,
cremes, merengues, caldas, geleias, compotas e outros produtos,
aplicando técnicas de confeitaria. Elabora acabamentos e adereços de
confeitaria, considerando valores estéticos, finalidade e receituário.
Utiliza bicos e saco de confeitar, considerando pressão, agilidade,
movimento e padronização. Decora produções da confeitaria considerando
técnicas de modelagem à base de açúcar e técnicas com chocolate. Pode
preparar produtos sem lactose, sem glúten, diet, light e com outras
adequações, para atender a pessoas em dietas ou com restrições
alimentares. Pode operar máquinas para embalar e etiquetar os produtos.
Identifica, na etiqueta, o produto, as datas de fabricação e validade e
outras informações previstas em normas. Pode armazenar produtos,
controlando o prazo de validade. Controla a produção, fazendo os
registros necessários. Conserva o local de trabalho limpo, higienizado e
organizado. Limpa, organiza e higieniza instrumentos, utensílios,
máquinas e equipamentos. Controla a reposição de matérias-primas e
insumos, examinando a programação da produção e o estoque disponível.
Recebe e armazena as mercadorias, controlando seu estado de conservação.
Destina adequadamente os resíduos gerados, realizando descarte de acordo
com as normas ambientais. Verifica o funcionamento das máquinas e dos
equipamentos, requisitando a necessária manutenção.

\begin{Form}
\textbf{Esta correspondência está adequada?} \\ 
\ChoiceMenu[radio,radiosymbol=\ding{52},bordercolor=0 0.5 0,name=cbovalid_ 848310 ]{}{CBO deve ficar=sim, \hspace{6em} CBO não fica aqui=nao}
\end{Form}

\newpage
\section{ 08006 -- Técnico em Panificação }

\textbf{Perfil do Curso (CNCT)}

O Técnico em Panificação será habilitado para:

Planejar e desenvolver a elaboração artesanal ou industrializada de
pães, massas, salgados e similares, utilizando equipamentos, utensílios
e tecnologias aplicadas aos processos, conforme as boas práticas de
manipulação de alimentos. Elaborar e padronizar fichas técnicas e de
controle operacional para diferentes tipos de operação, promovendo a
inovação e o desenvolvimento de novos produtos e processos. Utilizar as
boas práticas na manipulação de alimentos, rotulagem e identificação da
embalagem adequada. Efetuar controle de qualidade, de estoque, de custos
e de consumo. Utilizar técnicas mercadológicas de produtos e insumos
para a panificação.\\
Planejar e executar a aquisição e a manutenção preventiva de
equipamentos.

Para atuação como Técnico em Panificação, são fundamentais:

Conhecimentos e saberes relacionados às tecnologias de processamento e
conservação de produções de panificação, à implementação de análises
laboratoriais e sensoriais, à gestão dos sistemas de controle e garantia
da qualidade e da segurança de alimentos, à visão global dos processos
de produção manual e automatizado, à responsabilidade técnica, às normas
técnicas, à gestão sustentável, à proatividade, à liderança, à
capacidade de trabalho em equipes e à inovação tecnológica.

\textbf{CBOs correspondidos e seus perfis (presentes no CNCT, e presentes na predição da IA)}

\textbf{CBO: 513505  ( Auxiliar nos serviços de alimentação )}

Consulta ficha técnica, examinando os ingredientes e o modo de preparo
de cada prato. Prepara o mise em place, de acordo com o cardápio e com
as receitas definidas. Verifica a disponibilidade e a qualidade dos
gêneros alimentícios. Confere o prazo de validade na rotulagem dos
produtos embalados. Separa e organiza materiais e utensílios. Descongela
alimentos, escolhe grãos e cereais e mói ingredientes. Limpa, higieniza
e corta legumes, verduras e frutas. Pré-prepara carnes, aves e peixes
para cozimento, cortando-os, limpando-os, pesando-os e separando-os de
acordo com porções solicitadas. Prepara molhos, fundos e caldos. Prepara
caldas, recheios e coberturas. Auxilia na montagem dos pratos, dispondo
guarnição, molhos, caldos e outros alimentos. Ajuda na finalização e na
decoração dos pratos. Pode ajudar a manter abastecido o balcão de buffet
para autoatendimento (self-service), transportando alimentos quentes,
frios e sobremesas, repondo pratos e temperos e controlando a
temperatura dos alimentos. Pode coletar e acondicionar amostras de
alimentos para análise. Realiza a higienização do ambiente, das
bancadas, das louças, dos utensílios e dos equipamentos da cozinha em
geral. Auxilia no controle de estoque de alimentos, de modo a atender ao
cardápio pré-definido. Etiqueta e embala gêneros alimentícios. Comunica
necessidade de aquisições e confere os gêneros alimentícios recebidos.
Zela pela conservação dos alimentos estocados, providenciando as
condições necessárias para evitar deterioração e perdas.

\begin{Form}
\textbf{Esta correspondência está adequada?} \\ 
\ChoiceMenu[radio,radiosymbol=\ding{52},bordercolor=0 0.5 0,name=cbovalid_ 513505 ]{}{CBO deve ficar=sim, \hspace{6em} CBO não fica aqui=nao}
\end{Form}

\textbf{CBO: 848305  ( Padeiro )}

Prepara as atividades, de acordo com ficha técnica de produção e
programação da produção. Seleciona utensílios, instrumentos,
matérias-primas e insumos, de acordo com as especificações do processo e
dos produtos. Inspeciona as condições operacionais de máquinas e
equipamentos. Interpreta as receitas de pães e os procedimentos de
execução, determinando as proporções requeridas. Mede e pesa farinha,
sal e outros ingredientes para a massa. Mede e pesa ingredientes para
recheios e coberturas. Opera e controla tempo e velocidade das máquinas,
para mistura dos ingredientes e para homogeneizar a massa. Opera
equipamentos para cilindrar e dividir a massa. Pesa as porções da massa.
Pode modelar a massa. Pode preparar e colocar recheio na massa. Arruma
as porções de massa em fôrmas, colocando-as no forno. Assa os pães,
controlando os parâmetros do forno. Controla o resfriamento e desenforna
os pães. Inspeciona o produto, podendo prová-lo. Pode colocar coberturas
ou confeitar os pães assados. Pode operar máquinas de fatiar os pães.
Além de pães, pode produzir macarrão, biscoitos, bolos, tortas e outros
produtos similares. Pode operar máquinas para embalar e etiquetar os
produtos. Pode armazenar os produtos, controlando o prazo de validade.
Controla a produção, fazendo os registros necessários. Conserva o local
de trabalho limpo, higienizado e organizado. Limpa, organiza e higieniza
instrumentos, utensílios, máquinas e equipamentos. Destina adequadamente
os resíduos gerados, realizando descarte de acordo com as normas
ambientais. Controla a reposição de matérias-primas e insumos,
examinando a programação da produção e o estoque disponível. Recebe e
armazena as mercadorias, controlando seu estado de conservação. Verifica
o funcionamento das máquinas e dos equipamentos, requisitando a
necessária manutenção.

\begin{Form}
\textbf{Esta correspondência está adequada?} \\ 
\ChoiceMenu[radio,radiosymbol=\ding{52},bordercolor=0 0.5 0,name=cbovalid_ 848305 ]{}{CBO deve ficar=sim, \hspace{6em} CBO não fica aqui=nao}
\end{Form}

\newpage
\section{ 08007 -- Técnico em Viticultura e Enologia }

\textbf{Perfil do Curso (CNCT)}

O Técnico em Viticultura e Enologia será habilitado para:

Desenvolver e controlar os processos de cultivo da uva. Selecionar
variedades de uvas para elaboração de vinho e de seus derivados.
Supervisionar os processos de elaboração de vinho e de seus derivados.
Realizar ensaios físicos e análises químicas, sensoriais e
microbiológicas. Controlar e corrigir desvios nos processos manuais e
automatizados. Utilizar boas práticas de fabricação, rotulagem e
identificação da embalagem adequada. Promover a inovação tecnológica e o
desenvolvimento de novos produtos. Promover o tratamento e o destino
adequado de resíduos e efluentes. Operar equipamentos e efetuar o
controle de qualidade. Utilizar técnicas mercadológicas de produtos e
insumos. Planejar e executar o processo de trabalho, de comércio e de
venda de vinhos. Controlar estoques e custos. Realizar aquisição e
manutenção de equipamentos. Prestar assistência técnica na aplicação de
produtos e de serviços.

Para a atuação como Técnico em Viticultura e Enologia, são fundamentais:

Conhecimentos e saberes relacionados aos sistemas de produção de uva, de
vinho e de derivados da uva. Conhecimentos e saberes relacionados à
gestão de resíduos da indústria enológica. Conhecimentos e saberes
relacionados à gestão na indústria enológica, ao empreendedorismo e à
inovação. Conhecimentos e saberes relacionados às normas técnicas de
qualidade da indústria enológica e seus produtos. Conhecimentos e
saberes relacionados à execução de análises de qualidade da uva e do
vinho. Conhecimentos e saberes relacionados à solução de problemas, à
criatividade e ao trabalho em equipe.

\textbf{CBOs correspondidos e seus perfis (presentes no CNCT, e presentes na predição da IA)}

\textbf{CBO: 325005  ( Enólogo )}

Planeja a produção de uva e a fabricação de vinhos e derivados de uva e
vinho, considerando os recursos disponíveis e a demanda do mercado
consumidor. Seleciona máquinas, equipamentos, ferramentas e
instrumentos. Elabora orçamento. Pode adotar novas tecnologias, como
monitoramento da produção em tempo real com uso de sensores,
melhoramento genético de cultivares de uva, entre outras. Presta
informações sobre os processos de produção de uvas, orientando os
viticultores quanto aos aspectos técnicos para formar vinhedos de melhor
produtividade e maior qualidade. Orienta o cultivo de videiras,
recomendando variedades de uvas testadas e aprovadas; definindo o
espaçamento a ser adotado no plantio; e indicando os tratos culturais -
em relação à poda, aos tratamentos fitossanitários, entre outros -- a
serem realizados. Mostra a forma de avaliação do grau de maturação das
uvas, para definição do momento de colheita e dos lotes a serem
colhidos. Organiza processos de elaboração de vinhos e de derivados de
uva e vinho, definindo os tipos de produtos e os respectivos processos
de produção. Indica insumos para cada etapa do processo. Orienta a
equipe de trabalho quanto à higienização de equipamentos e utensílios.
Controla a qualidade de matérias-primas e demais insumos. Avalia o
estado fitossanitário das uvas recebidas e analisa a qualidade de
fermentos e complexos enzimáticos, clarificantes, auxiliares de
filtração, estabilizantes e conservantes. Avalia a qualidade dos
vasilhames e fechamentos. Analisa a qualidade de rótulos, adesivos e
caixas. Monitora os processos de vinificação, tais como recebimento das
uvas, extração do suco, fermentação, entre outros. Controla as etapas de
engarrafamento - filtração, enchimento, fechamento, rotulagem e
acondicionamento - de vinhos e de derivados de uva e vinho. Orienta a
elaboração de vinhos espumantes, vinhos compostos e licorosos;
conhaques, grapas e outros destilados; vinagres; e sucos concentrados de
uva. Promove a execução de análise sensorial - visual, olfativa e
degustativa - durante os processos de elaboração; nos produtos
finalizados; e em produtos concorrentes. Orienta a realização de
análises de controle físico, químico, bioquímico e microbiológico nas
fases dos processos, observando os critérios de qualidade especificados.
Avalia resultados de controles analíticos realizados nos laboratórios.
Desenvolve atividades de pesquisa, fazendo experiências com novos
insumos enológicos; testando novos equipamentos e sistemas
automatizados; e verificando as melhorias decorrentes do uso de novos
processos de elaboração de vinhos e derivados de uva e vinho. Mantém-se
atualizado em relação às inovações em viticultura e enologia. Desenvolve
atividades de divulgação de produtos, promovendo visitação às adegas
produtoras de vinho, ministrando cursos e proferindo palestras e
conferências. Coordena ações para o cumprimento de normas legais,
orientando o registro de produtos, o registro da empresa e a elaboração
de relatórios de produção, estoques e comercialização para o órgão
fiscalizador. Pode executar perícias exigidas em processos judiciais a
título de prova e contraprova. Presta suporte técnico a clientes
internos e externos, fornecendo subsídios técnicos para cálculos de
custo e de rentabilidade; assessorando a área comercial no atendimento a
clientes; executando serviços técnicos em órgãos oficiais de pesquisa e
fiscalização do setor vitivinícola; e realizando consultoria técnica à
área industrial de fabricação de vinhos. Pode supervisionar equipe de
trabalho, avaliando desempenho e desenvolvendo treinamentos. Observa a
organização do local de trabalho. Prepara plano de manutenção de
máquinas e equipamentos. Contribui para a redução dos impactos
ambientais, promovendo o tratamento e o destino adequado de resíduos e
efluentes. Zela pela segurança, prevenindo situações de risco. Utiliza e
monitora o uso de equipamentos de proteção individual. Presta primeiros
socorros.

\begin{Form}
\textbf{Esta correspondência está adequada?} \\ 
\ChoiceMenu[radio,radiosymbol=\ding{52},bordercolor=0 0.5 0,name=cbovalid_ 325005 ]{}{CBO deve ficar=sim, \hspace{6em} CBO não fica aqui=nao}
\end{Form}

\newpage
\section{ 09001 -- Técnico em Artes Circenses }

\textbf{Perfil do Curso (CNCT)}

O Técnico em Artes Circenses será habilitado para:

Criar, desenvolver e executar apresentações circenses em espaços de
circo, teatro, estúdio de televisão, públicos e culturais. Utilizar
técnicas artísticas e corporais de acrobacia aérea e de solo,
equilibrismo, malabarismo, antipodismo, ilusionismo, comicidade, canto,
dança e pantomima. Organizar e supervisionar a estrutura, a montagem e o
funcionamento do circo e dos equipamentos. Administrar, produzir e
divulgar espetáculos.

Para atuação como Técnico em Artes Circenses, são fundamentais:

Conhecimentos interdisciplinares relacionados aos processos de criação,
envolvendo pesquisa, idealização, planejamento, execução técnica,
fruição e recepção estética. Competências comunicativas e empreendedoras
voltadas à proposição de projetos, ao coletivo, à gestão, à solução de
problemas e à resiliência, entre outras competências socioemocionais.

\textbf{CBOs correspondidos e seus perfis (presentes no CNCT, e presentes na predição da IA)}

\textbf{CBO: 376205  ( Acrobata )}

Cria novo número para apresentações de acrobacia de solo, pesquisando
uso de materiais, aparelhos e tecnologias. Pesquisa possibilidades no
uso das cores, em luzes, figurino e materiais. Pesquisa movimentos
corporais, possibilidades de expressão artística e formas de comunicação
com o público. Faz intercâmbio de informações com profissionais do
circo, com escolas de circo e com profissionais de outras áreas
artísticas, pessoalmente, usando Internet ou outras formas de
comunicação. Pode misturar números existentes, criando outros. Pode
criar números de ginástica acrobática, praticada ao som de música, por
pares ou grupos, que demonstram equilíbrio, flexibilidade e agilidade na
execução. Confecciona, monta e adapta equipamentos, aparelhos e
materiais -- tais como argolas, cavalos de madeira, camas elásticas,
entre outros - necessários ao desenvolvimento do número circense,
realizando as adequações necessárias conforme o seu biótipo e as
condições do local das apresentações; e aplicando novas tecnologias
disponíveis. Pode criar, confeccionar e montar novo aparelho para o
número. Produz o número circense, definindo coreografia e selecionando
música. Cria guarda-roupa e maquiagem. Seleciona equipamentos de
segurança. Ensaia o número, repetindo-o para aperfeiçoar técnicas de
expressão corporal e assimilar os tempos na realização das acrobacias.
Prepara a forma de sua entrada para iniciar a apresentação. Incorpora
equipamentos de segurança no número, durante o ensaio. Prepara material,
aparelho e objetos para a apresentação. Colabora na divulgação do número
e do espetáculo, providenciando material impresso, dando entrevistas e
utilizando jornal, carros de som, redes sociais e outras formas de
comunicação. Prepara corpo e mente para apresentação do número,
alongando e aquecendo o corpo e realizando exercícios de concentração.
Apresenta o número, interagindo com o público; desenvolvendo formas de
comunicação para conquistar a empatia das pessoas; e usando técnicas e
estratégias para valorização do seu trabalho. Lida com imprevistos,
durante a apresentação, de forma criativa. Interage com os demais
profissionais envolvidos no espetáculo, combinando códigos para informar
imprevistos e ajustando estratégias para lidar com incidentes, acidentes
ou outros acontecimentos. Trabalha em equipe, estabelecendo vínculos de
confiança com os colegas. Pode vender o número ou espetáculo, avaliando
custos para fazer preço do trabalho e investigando o valor da atividade
circense no mercado. Pode ministrar aulas, promover oficinas e oferecer
cursos ligados às artes acrobáticas, para públicos de diferentes faixas
etárias. Seleciona métodos e estratégias para transmissão de
conhecimentos e realização da prática profissional. Motiva o aluno e
estimula o seu desenvolvimento físico. Transmite a ética circense.
Conserva equipamentos, aparelhos e materiais, usados nas acrobacias,
limpos e organizados. Zela pela sua segurança, pela segurança da equipe
e pela segurança do público, realizando as atividades em ambiente
seguro, contando com pessoal de apoio capacitado e fazendo uso de
equipamentos de segurança. Pode prestar primeiros socorros.

\begin{Form}
\textbf{Esta correspondência está adequada?} \\ 
\ChoiceMenu[radio,radiosymbol=\ding{52},bordercolor=0 0.5 0,name=cbovalid_ 376205 ]{}{CBO deve ficar=sim, \hspace{6em} CBO não fica aqui=nao}
\end{Form}

\textbf{CBO: 376210  ( Artista aéreo )}

Cria novo número para apresentações de acrobacia aéreas, pesquisando uso
de materiais, aparelhos e tecnologias. Pesquisa possibilidades no uso
das cores, em luzes, figurino e materiais. Pesquisa movimentos
corporais, possibilidades de expressão artística e formas de comunicação
com o público. Faz intercâmbio de informações com profissionais do
circo, com escolas de circo e com profissionais de outras áreas
artísticas, pessoalmente, usando Internet ou outras formas de
comunicação. Pode fazer fusão de números existentes, criando outros.
Confecciona, monta e adapta equipamentos, aparelhos e materiais -- tais
como lira (aro de ferro suspenso), mastro chinês, tecido acrobático,
corda indiana, entre outros - necessários ao desenvolvimento do número
circense, aplicando novas tecnologias disponíveis; e realizando
adequações necessárias conforme o seu biótipo e os biótipos dos artistas
coadjuvantes e de acordo com as condições do local das apresentações.
Produz o número circense, definindo coreografia e selecionando música.
Cria guarda-roupa e maquiagem. Seleciona equipamentos de segurança.
Ensaia o número, repetindo-o para aperfeiçoar técnicas de expressão
corporal e assimilar os tempos na realização das acrobacias aéreas.
Prepara a forma de sua entrada para iniciar a apresentação. Incorpora
equipamentos de segurança no número, durante o ensaio. Colabora na
divulgação do número e do espetáculo, providenciando material impresso,
dando entrevistas e utilizando jornal, carros de som, redes sociais e
outras formas de comunicação. Prepara corpo e mente para apresentação do
número, alongando e aquecendo o corpo e realizando exercícios de
concentração. Apresenta o número, interagindo com o público;
desenvolvendo formas de comunicação para conquistar a empatia das
pessoas; e usando técnicas e estratégias para valorização do seu
trabalho. Executa os fundamentos técnicos e princípios artísticos que
norteiam as acrobacias aéreas em dupla ou coletivas. Lida com
imprevistos, durante a apresentação, de forma criativa. Interage com os
demais profissionais envolvidos no espetáculo, combinando códigos para
informar imprevistos e ajustando estratégias para lidar com incidentes,
acidentes ou outros acontecimentos. Trabalha em equipe, estabelecendo
vínculos de confiança com os colegas. Pode vender o número ou
espetáculo, avaliando custos para fazer preço do trabalho e investigando
o valor da atividade circense no mercado. Pode ministrar aulas, promover
oficinas e oferecer cursos ligados às artes acrobáticas aéreas, para
públicos de diferentes faixas etárias. Seleciona métodos e estratégias
para transmissão de conhecimentos e realização da prática profissional.
Motiva o aluno e estimula o seu desenvolvimento físico. Transmite a
ética circense. Conserva equipamentos, aparelhos e materiais, usados nas
acrobacias aéreas, limpos e organizados. Zela pela sua segurança, pela
segurança da equipe e pela segurança do público, realizando as
atividades em ambiente seguro, contando com pessoal de apoio capacitado
e fazendo uso de equipamentos de segurança. Pode prestar primeiros
socorros.

\begin{Form}
\textbf{Esta correspondência está adequada?} \\ 
\ChoiceMenu[radio,radiosymbol=\ding{52},bordercolor=0 0.5 0,name=cbovalid_ 376210 ]{}{CBO deve ficar=sim, \hspace{6em} CBO não fica aqui=nao}
\end{Form}

\textbf{CBO: 376215  ( Artista de circo (outros) )}

Auxilia artistas circenses na criação de personagens para apresentações
itinerantes, propondo a utilização e a combinação de técnicas cênicas e
circenses para entreter a plateia. Pode prestar apoio a diferentes
artistas circenses, por possuir noções básicas de diferentes técnicas
artísticas circenses, como dança, mágica, malabarismo, equilibrismo,
acrobacia, caracterizações cômicas, entre outras. Pode atuar como
coadjuvante dos artistas circenses durante as apresentações. Contribui
com a realização de pesquisas para criação de um novo número circense,
sugerindo fontes de pesquisa. Levanta informações sobre materiais,
aparelhos e tecnologias; sobre uso das cores, em luzes, figurino e
materiais; e sobre movimentos corporais e formas de comunicação com o
público. Pesquisa endereços de profissionais do circo, escolas de circo
e profissionais de outras áreas artísticas, usando Internet ou fontes de
informação, para possibilitar intercâmbio de informações entre os
artistas. Auxilia na confecção, montagem e adaptação de equipamentos,
aparelhos e materiais -- tais como malabares, cordas, argolas, tecidos
acrobáticos, trapézios, entre outros - necessários ao desenvolvimento do
número circense. Recebe orientações para fazer adequações nos
equipamentos e nos aparelhos - conforme o biótipo do artista principal e
os biótipos dos artistas coadjuvantes - e adaptações às condições do
local das apresentações. Contribui na produção do número circense,
auxiliando na montagem e na definição de coreografia, música, figurino e
maquiagem. Ensaia o número - quando atua como coadjuvante - repetindo-o
para aperfeiçoar sua interação com o artista principal. Colabora na
divulgação do número e do espetáculo, distribuindo material impresso ou
usando outras estratégias para atrair o público. Prepara-se para
participar no número, alongando e aquecendo o corpo e realizando
exercícios de concentração. Participa da apresentação, seguindo o
roteiro definido pelo artista principal. Ajuda no monitoramento do tempo
de cada parte do número e na disponibilização de aparelhos e materiais
predefinidos para execução da apresentação. Pode atuar como aparador,
servindo de apoio a artista principal em suas execuções no solo ou no
ar. Ajuda na desmontagem de aparelhos, após os números. Interage com os
demais profissionais envolvidos no espetáculo, estabelecendo vínculos de
confiança com os colegas. Pode auxiliar na realização de cursos ligados
às artes circenses, prestando apoio a públicos de diferentes faixas
etárias para execução de exercícios e outras atividades práticas.
Conserva limpos e organizados os equipamentos, aparelhos e materiais
usados em diferentes apresentações. Pode fazer pequenos reparos nos
aparelhos, quando necessário. Zela pela sua segurança, pela segurança
dos artistas principais e pela segurança do público, capacitando-se para
atuar como coadjuvante em novos números. Realiza as atividades em
ambiente seguro e faz uso de equipamentos de segurança. Pode prestar
primeiros socorros.

\begin{Form}
\textbf{Esta correspondência está adequada?} \\ 
\ChoiceMenu[radio,radiosymbol=\ding{52},bordercolor=0 0.5 0,name=cbovalid_ 376215 ]{}{CBO deve ficar=sim, \hspace{6em} CBO não fica aqui=nao}
\end{Form}

\textbf{CBO: 376220  ( Contorcionista )}

Cria novo número para apresentações de contorcionismo, pesquisando a
utilização de diferentes técnicas de flexão e contorção do corpo.
Pesquisa possibilidades no uso das cores, em luzes e figurino. Pesquisa
dança, malabarismo e outras expressões artísticas para mesclar no número
de contorcionismo. Faz intercâmbio de informações com profissionais do
circo, com escolas de circo e com profissionais de outras áreas
artísticas, pessoalmente, usando Internet ou outras formas de
comunicação. Pode fazer fusão de números existentes, criando outros.
Confecciona e monta caixas - de vidro, acrílico ou outro material - de
diferentes tamanhos e outros objetos, para usar durante a apresentação.
Produz o número circense, definindo coreografia e selecionando música.
Cria guarda-roupa e maquiagem. Seleciona equipamentos de segurança.
Ensaia o número, repetindo-o para assimilar os tempos na realização do
contorcionismo e para treinar danças, formação de figuras de animais com
o corpo ou outras formas de expressões corporais que integram o número.
Prepara sua entrada no picadeiro para iniciar a apresentação. Incorpora
equipamentos de segurança no número, durante o ensaio. Colabora na
divulgação do número e do espetáculo, providenciando material impresso,
dando entrevistas e utilizando jornal, carros de som, redes sociais e
outras formas de comunicação. Prepara corpo e mente para apresentação do
número, alongando e aquecendo o corpo e realizando exercícios de
concentração. Apresenta o número, interagindo com o público;
desenvolvendo formas de comunicação para conquistar a empatia das
pessoas; e usando técnicas e estratégias para valorização do seu
trabalho. Lida com imprevistos, durante a apresentação, de forma
criativa. Interage com os demais profissionais envolvidos no espetáculo,
combinando códigos para informar imprevistos e ajustando estratégias
para lidar com incidentes, acidentes ou outros acontecimentos. Trabalha
em equipe, estabelecendo vínculos de confiança com os colegas. Pode
vender o número ou espetáculo, avaliando custos para fazer preço do
trabalho e investigando o valor da atividade circense no mercado. Pode
ministrar aulas, promover oficinas e oferecer cursos ligados ao
contorcionismo, para públicos de diferentes faixas etárias. Seleciona
métodos e estratégias para transmissão de conhecimentos e realização da
prática profissional. Motiva o aluno e estimula suas capacidades
físicas, principalmente a flexibilidade. Transmite a ética circense.
Conserva limpos e organizados objetos - como arco e flechas, caixas,
bolas, entre outros - usados durante os números de contorcionismo. Zela
pela sua segurança, pela segurança da equipe e pela segurança do
público, realizando as atividades em ambiente seguro, contando com
pessoal de apoio capacitado e fazendo uso de equipamentos de segurança.
Pode prestar primeiros socorros.

\begin{Form}
\textbf{Esta correspondência está adequada?} \\ 
\ChoiceMenu[radio,radiosymbol=\ding{52},bordercolor=0 0.5 0,name=cbovalid_ 376220 ]{}{CBO deve ficar=sim, \hspace{6em} CBO não fica aqui=nao}
\end{Form}

\textbf{CBO: 376225  ( Domador de animais (circense) )}

Atua em apresentações circenses com animais -- como elefantes, leões,
tigres, cavalos, entre outros -, utilizando e combinando técnicas
cênicas e circenses. Cuida e zela pela segurança e pela saúde dos
animais, executando - sob orientação de veterinários e técnicos - os
procedimentos necessários e adequados para abrigo, alimentação,
higienização e controle de saúde e bem-estar. Providencia espaços para
que os animais se exercitem. Usa princípios de psicologia animal no
trato diário. Avalia cada animal para determinar seu temperamento e sua
aptidão para o treinamento. Doma e condiciona animais para espetáculos
circenses, aplicando técnicas de amestramento não violentas e usando os
mesmos sinais e comandos para treinar e para atuar nas apresentações.
Leva em conta as peculiaridades dos animais. Pode aproveitar gestos e
ações que cada animal normalmente faz em liberdade -- como o elefante
ficar em pé para pegar as folhas mais novas de uma árvore -- para
condicioná-lo a fazer o mesmo, a partir de um comando, em um novo
número. Faz pesquisas de espetáculos circenses com animais para criar
números. Realiza intercâmbio de informações com outros domadores e com
escolas de circo, pessoalmente, usando Internet ou outras formas de
comunicação. Produz número circense, montando e definindo
sequencialidade, selecionando figurino e materiais a serem usados, bem
como determinando o tempo para realização da apresentação. Confecciona,
monta e adapta aparelhos e materiais -- como jaulas, bolas, cordas,
suportes, arcos, entre outros - necessários ao desenvolvimento seguro do
espetáculo, realizando adequações necessárias considerando o seu biótipo
e o biótipo dos animais, bem como as condições do local da apresentação.
Ensaia o número, repetindo-o para aumentar o condicionamento dos animais
aos sinais e comandos previstos no número. Fala com os animais para
familiarizá-los com vozes humanas. Prepara a forma de sua entrada para
iniciar a apresentação. Incorpora equipamentos de segurança no número,
durante o ensaio. Colabora na divulgação do número e do espetáculo,
dando entrevistas e utilizando jornal, carros de som, redes sociais e
outras formas de comunicação. Prepara-se para realizar o número,
paramentando-se e maquiando-se para a apresentação. Prepara os animais
para o espetáculo. Apresenta o número, usando técnicas e estratégias
para valorização do seu trabalho. Interage com os demais profissionais
envolvidos no espetáculo, combinando códigos para informar imprevistos e
ajustando estratégias para lidar com incidentes, acidentes ou outros
acontecimentos. Trabalha em equipe, estabelecendo vínculos de confiança
com os colegas. Pode vender o número ou espetáculo, avaliando custos
para fazer preço do trabalho e investigando o valor da atividade
circense no mercado. Pode ministrar aulas, promover oficinas e oferecer
cursos ligados à arte de domar animais. Seleciona métodos e estratégias
para transmissão de conhecimentos e realização da prática profissional.
Transmite a ética circense. Conserva limpos e organizados equipamentos,
aparelhos e materiais usados durante as apresentações com animais. Zela
pela sua segurança, pela segurança dos animais e pela segurança do
público, realizando as atividades em ambiente seguro, contando com
pessoal de apoio capacitado e fazendo uso de equipamentos de segurança.
Pode prestar primeiros socorros.

\begin{Form}
\textbf{Esta correspondência está adequada?} \\ 
\ChoiceMenu[radio,radiosymbol=\ding{52},bordercolor=0 0.5 0,name=cbovalid_ 376225 ]{}{CBO deve ficar=sim, \hspace{6em} CBO não fica aqui=nao}
\end{Form}

\textbf{CBO: 376230  ( Equilibrista )}

Cria novo número para apresentações de equilibrismo, pesquisando a
utilização de diferentes técnicas de controle e equilíbrio do corpo e de
objetos. Pesquisa movimentos corporais, possibilidades de expressão
artística e formas de comunicação com o público. Pesquisa possibilidades
no uso das cores, em luzes e figurino. Faz intercâmbio de informações
com profissionais do circo, com escolas de circo e com profissionais de
outras áreas artísticas, pessoalmente, usando Internet ou outras formas
de comunicação. Pode fazer fusão de números existentes, criando outros.
Confecciona, adapta e monta aparelhos e prepara objetos para as
apresentações. Faz equilíbrio andando no arame (baixo, alto, bambo ou
inclinado); no arame à grande altura com maromba (longa vara auxiliar
para manter o equilíbrio); no monociclo; na perna-de-pau; no rola-rola
(prancha de madeira ou de ferro colocada sobre um ou vários cilindros
superpostos); entre outros. Faz equilíbrio com objetos, tais como pires,
argolas, xícaras e cadeiras. Produz o número circense, definindo
coreografia e selecionando música. Cria guarda-roupa e maquiagem.
Seleciona equipamentos de segurança. Ensaia o número, repetindo-o para
aperfeiçoar técnicas de controle e de equilíbrio do corpo e de objetos;
e para assimilar os tempos na realização de cada parte do número. Treina
sua entrada no picadeiro para iniciar a apresentação. Incorpora
equipamentos de segurança no número, durante o ensaio. Prepara aparelhos
e objetos para a apresentação. Colabora na divulgação do número e do
espetáculo, providenciando material impresso, dando entrevistas e
utilizando jornal, carros de som, redes sociais e outras formas de
comunicação. Prepara corpo e mente para apresentação do número,
alongando e aquecendo o corpo e realizando exercícios de concentração.
Apresenta o número, interagindo com o público; desenvolvendo formas de
comunicação para conquistar a empatia das pessoas; e usando técnicas e
estratégias para valorização do seu trabalho. Lida com imprevistos,
durante a apresentação, de forma criativa. Interage com os demais
profissionais envolvidos no espetáculo, combinando códigos para informar
imprevistos e ajustando estratégias para lidar com incidentes, acidentes
ou outros acontecimentos. Trabalha em equipe, estabelecendo vínculos de
confiança com os colegas. Pode vender o número ou espetáculo, avaliando
custos para fazer preço do trabalho e investigando o valor da atividade
circense no mercado. Pode ministrar aulas, promover oficinas e oferecer
cursos relacionados ao equilibrismo, para públicos de diferentes faixas
etárias. Seleciona métodos e estratégias para transmissão de
conhecimentos e realização da prática profissional. Motiva o aluno e
estimula suas capacidades físicas, principalmente a destreza e o
equilíbrio. Transmite a ética circense. Conserva limpos e organizados
aparelhos e objetos usados durante os números de equilibrismo. Zela pela
sua segurança, pela segurança da equipe e pela segurança do público,
realizando as atividades em ambiente seguro, contando com pessoal de
apoio capacitado e fazendo uso de equipamentos de segurança. Pode
prestar primeiros socorros.

\begin{Form}
\textbf{Esta correspondência está adequada?} \\ 
\ChoiceMenu[radio,radiosymbol=\ding{52},bordercolor=0 0.5 0,name=cbovalid_ 376230 ]{}{CBO deve ficar=sim, \hspace{6em} CBO não fica aqui=nao}
\end{Form}

\textbf{CBO: 376235  ( Mágico )}

Cria novo número para apresentações de mágica, ilusionismo e
prestidigitação, pesquisando a utilização de diferentes técnicas cênicas
e circenses, para entreter plateias. Pesquisa movimentos corporais,
expressões faciais, truques e formas de comunicação com o público.
Pesquisa possibilidades no uso de elementos visuais, tais como cores,
luzes e figurino. Faz intercâmbio de informações com profissionais do
circo, com escolas de circo e com profissionais de outras áreas
artísticas, pessoalmente, usando Internet ou outras formas de
comunicação. Pode fazer fusão de números existentes, criando outros.
Confecciona, adapta e monta aparelhos e prepara objetos para as
apresentações, tais como cartola, moedas, caixas com fundo falso,
espelhos, camisas de força, aquários, entre outros. Prepara baú, onde
entra algemado e trancado, para surgir depois sem algemas e substituído
por outro artista algemado. Faz truques com cartas de baralho. Realiza
números com animais domesticados, como pássaros e coelhos. Produz o
número circense, definindo coreografia e selecionando música. Cria
guarda-roupa e maquiagem. Seleciona equipamentos de segurança. Ensaia o
número, repetindo-o para aperfeiçoar arte e prática de mágica e
ilusionismo e técnicas de prestidigitação; e para assimilar os tempos na
realização de cada parte do número. Fala com os animais para
familiarizá-los com vozes humanas. Treina sua entrada no picadeiro para
iniciar a apresentação. Incorpora equipamentos de segurança no número,
durante o ensaio. Colabora na divulgação do número e do espetáculo,
providenciando material impresso, dando entrevistas e utilizando jornal,
carros de som, redes sociais e outras formas de comunicação.
Caracteriza-se como mágico, paramentando-se e maquiando-se para a
apresentação. Realiza exercícios de concentração e aquecimento vocal.
Apresenta o número, interagindo com o público; desenvolvendo formas de
comunicação para conquistar a empatia das pessoas; e usando técnicas e
estratégias para valorização do seu trabalho. Lida com imprevistos,
durante a apresentação, de forma criativa. Interage com os demais
profissionais envolvidos no espetáculo, combinando códigos para informar
imprevistos e ajustando estratégias para lidar com incidentes, acidentes
ou outros acontecimentos. Trabalha em equipe, estabelecendo vínculos de
confiança com os colegas. Pode atuar em teatros e outros espaços
culturais e em meios de comunicação - como televisão e cinema -, com a
finalidade de entreter plateias e audiências. Pode vender o número ou
espetáculo, avaliando custos para fazer preço do trabalho e investigando
o valor da atividade circense no mercado. Pode ministrar aulas, promover
oficinas e oferecer cursos relacionados à mágica, ao ilusionismo e à
prestidigitação, para públicos de diferentes faixas etárias. Seleciona
métodos e estratégias para transmissão de conhecimentos e realização da
prática profissional. Motiva o aluno e estimula sua rapidez dos gestos e
das mãos; sua presença de palco - para prender a atenção e o interesse
do público --; e sua capacidade verbal, para interagir e estimular a
participação da plateia. Transmite a ética circense. Cuida da
alimentação, da higiene, da saúde e do bem-estar dos animais
domesticados que integram seus números. Considera princípios de direitos
e segurança dos animais. Conserva limpos e organizados aparelhos e
objetos usados durante os números de mágica, ilusionismo e
prestidigitação. Zela pela sua segurança, pela segurança da equipe, pela
segurança dos animais e pela segurança do público, realizando as
atividades em ambiente seguro, contando com pessoal de apoio capacitado
e fazendo uso de equipamentos de segurança. Pode prestar primeiros
socorros.

\begin{Form}
\textbf{Esta correspondência está adequada?} \\ 
\ChoiceMenu[radio,radiosymbol=\ding{52},bordercolor=0 0.5 0,name=cbovalid_ 376235 ]{}{CBO deve ficar=sim, \hspace{6em} CBO não fica aqui=nao}
\end{Form}

\textbf{CBO: 376240  ( Malabarista )}

Cria novo número para apresentações de malabarismo, pesquisando a
utilização de diferentes técnicas de controle e equilíbrio do corpo e de
objetos, para entreter plateias. Pesquisa movimentos corporais,
expressões faciais e formas de comunicação com o público. Pesquisa
possibilidades no uso de elementos visuais, tais como cores, luzes e
figurino. Faz intercâmbio de informações com profissionais do circo, com
escolas de circo e com profissionais de outras áreas artísticas,
pessoalmente, usando Internet ou outras formas de comunicação. Pode
fazer fusão de números existentes, criando outros. Prepara aparelhos e
objetos para as apresentações, tais como aros - de madeira ou de
plástico -, argolas, bambolês, bolas, entre outros. Faz malabarismo com
as mãos, equilibrando clavas, chapéus, tochas acesas e outros objetos no
ar. Faz malabarismo de objetos com os pés, geralmente com barris e
cilindros. Executa troca - entre duas ou mais pessoas - com claves ou
outros objetos, podendo realizar truques de lançamento por baixo das
pernas ou por trás do ombro. Produz o número circense, definindo
coreografia e selecionando música. Cria guarda-roupa e maquiagem.
Seleciona equipamentos de segurança. Ensaia o número, repetindo-o para
aperfeiçoar postura corporal e cênica; destreza nos movimentos;
manipulação cênica dos objetos; e forma de relacionamento com o público.
Assimila os tempos na realização de cada parte do número. Treina sua
entrada no picadeiro para iniciar a apresentação. Incorpora equipamentos
de segurança no número, durante o ensaio. Colabora na divulgação do
número e do espetáculo, providenciando material impresso, dando
entrevistas e utilizando jornal, carros de som, redes sociais e outras
formas de comunicação. Prepara corpo e mente para apresentação do
número, alongando e aquecendo o corpo e realizando exercícios de
concentração. Apresenta o número, interagindo com o público;
desenvolvendo formas de comunicação para conquistar a empatia das
pessoas; e usando técnicas e estratégias para valorização do seu
trabalho. Lida com imprevistos, durante a apresentação, de forma
criativa. Interage com os demais profissionais envolvidos no espetáculo,
combinando códigos para informar imprevistos e ajustando estratégias
para lidar com incidentes, acidentes ou outros acontecimentos. Trabalha
em equipe, estabelecendo vínculos de confiança com os colegas. Pode
vender o número ou espetáculo, avaliando custos para fazer preço do
trabalho e investigando o valor da atividade circense no mercado. Pode
ministrar aulas, promover oficinas e oferecer cursos relacionados ao
malabarismo, para públicos de diferentes faixas etárias. Seleciona
métodos e estratégias para transmissão de conhecimentos e realização da
prática profissional. Desenvolve no aluno a capacidade de dominar
fundamentos técnicos e princípios artísticos do malabarismo e a
capacidade verbal, para interagir com a plateia. Transmite a ética
circense. Conserva limpos e organizados os aparelhos e objetos usados
durante os números de malabarismo. Zela pela sua segurança, pela
segurança da equipe e pela segurança do público, realizando as
atividades em ambiente seguro, contando com pessoal de apoio capacitado
e fazendo uso de equipamentos de segurança. Pode prestar primeiros
socorros.

\begin{Form}
\textbf{Esta correspondência está adequada?} \\ 
\ChoiceMenu[radio,radiosymbol=\ding{52},bordercolor=0 0.5 0,name=cbovalid_ 376240 ]{}{CBO deve ficar=sim, \hspace{6em} CBO não fica aqui=nao}
\end{Form}

\textbf{CBO: 376245  ( Palhaço )}

Cria personagens para apresentações de números cômicos e pantomimas.
Inventa número, pesquisando a utilização de diferentes técnicas cênicas
e circenses. Pesquisa movimentos corporais, expressões faciais e formas
de comunicação com o público. Pesquisa dança e outras expressões
artísticas para usar nas apresentações. Pesquisa possibilidades no uso
de elementos visuais, tais como cores, luzes, figurino e maquilagem. Faz
intercâmbio de informações com profissionais do circo, com escolas de
circo e com profissionais de outras áreas artísticas, pessoalmente,
usando Internet ou outras formas de comunicação. Pode fazer fusão de
números existentes, criando outros. Confecciona, monta e adapta
aparelhos e objetos para o desenvolvimento do número - como caixotes de
vidro, instrumentos musicais, bolas, argolas, entre outros -, realizando
as adequações necessárias - conforme o seu biótipo e os biótipos dos
artistas eventualmente participantes - e adaptações às condições do
local das apresentações. Usa noções básicas de diferentes técnicas
artísticas circenses -- como malabarismo, equilibrismo, acrobacia, entre
outras - em seu número. Pode fazer apresentações em que se manifesta só
por gestos, expressões corporais ou fisionômicas para transmitir
sentimentos, pensamentos e ideias. Pode usar dança, canto e outras
expressões artísticas. Produz o número circense, definindo coreografia e
selecionando música. Cria guarda-roupa e maquiagem. Seleciona
equipamentos de segurança. Ensaia o número, repetindo-o para aperfeiçoar
arte e prática de apresentar ações e expressões faciais ou corporais
para divertir ou comover o público; assimilar os tempos na realização de
cada parte do número; e aprimorar a interação com demais artistas,
quando atua com outros palhaços ou com coadjuvantes. Treina sua entrada
no picadeiro para iniciar a apresentação. Incorpora equipamentos de
segurança no número, durante o ensaio. Colabora na divulgação do número
e do espetáculo, providenciando material impresso, dando entrevistas e
utilizando jornal, carros de som, redes sociais e outras formas de
comunicação. Prepara corpo e mente para desenvolvimento do espetáculo,
alongando e aquecendo o corpo, realizando exercícios de concentração e
fazendo aquecimento vocal. Caracteriza-se como personagem excêntrico -
pintando o rosto, colocando nariz vermelho e usando chapéu, roupas
folgadas com colarinho e sapatos enormes -, mímico, vagabundo ou tipo de
criação própria, paramentando-se e maquiando-se para a apresentação.
Apresenta o número, interagindo com o público; desenvolvendo formas de
comunicação para conquistar a empatia das pessoas; e usando técnicas e
estratégias para valorização do seu trabalho. Lida com imprevistos,
durante a apresentação, de forma criativa. Interage com os demais
profissionais envolvidos no espetáculo, combinando códigos para informar
imprevistos e ajustando estratégias para lidar com incidentes, acidentes
ou outros acontecimentos. Trabalha em equipe, estabelecendo vínculos de
confiança com os colegas. Pode vender o número ou espetáculo, avaliando
custos para fazer preço do trabalho e investigando o valor da atividade
circense no mercado. Pode ministrar aulas, promover oficinas e oferecer
cursos relacionados ao ofício de palhaço, para públicos de diferentes
faixas etárias. Seleciona métodos e estratégias para transmissão de
conhecimentos e realização da prática profissional. Motiva o aluno e
estimula sua presença de palco e sua capacidade verbal, para interagir e
estimular a participação da plateia. Transmite a ética circense.
Conserva limpos e organizados aparelhos e objetos usados durante as
apresentações. Zela pela sua segurança, pela segurança da equipe e pela
segurança do público, realizando as atividades em ambiente seguro,
contando com pessoal de apoio capacitado e fazendo uso de equipamentos
de segurança. Pode prestar primeiros socorros.

\begin{Form}
\textbf{Esta correspondência está adequada?} \\ 
\ChoiceMenu[radio,radiosymbol=\ding{52},bordercolor=0 0.5 0,name=cbovalid_ 376245 ]{}{CBO deve ficar=sim, \hspace{6em} CBO não fica aqui=nao}
\end{Form}

\textbf{CBO: 376250  ( Titeriteiro )}

Cria personagens para bonecos - representando animais, pessoas e/ou
objetos animados -, que se movem por cordéis ou engonços, em espetáculos
circenses. Inventa número, pesquisando a utilização de diferentes
técnicas cênicas e circenses. Pesquisa movimentos antropomórficos e
formas de comunicação com o público. Pesquisa dança e outras expressões
artísticas para utilizar nas apresentações. Pesquisa possibilidades no
uso de elementos visuais, tais como cores e luzes. Faz intercâmbio de
informações com profissionais do circo, com escolas de circo e com
profissionais de outras áreas artísticas, pessoalmente, usando Internet
ou outras formas de comunicação. Pode fazer fusão de números existentes,
criando outros. Confecciona boneco manipulável, assegurando sua
capacidade de movimentar-se e dando-lhe forma, cor e expressões,
conforme características do personagem que representa. Confecciona e
monta cenário para atuação dos bonecos. Pode fazer a manipulação de
forma oculta -- ficando atrás de cortina ou cenário -- ou exposto,
representando personagem - como vendedor de loja de brinquedo, professor
que apresenta uma história, entre outros -, mas sempre mantendo os
bonecos como foco principal para a plateia. Produz o número circense,
definindo coreografia, selecionando música e criando guarda-roupa para
os bonecos e, quando for o caso, para si próprio, enquanto personagem.
Define os movimentos a fazer para dar vida aos bonecos. Pode atuar com
outros titeriteiros ou com coadjuvantes. Seleciona equipamentos de
segurança. Ensaia o número, repetindo-o na busca do movimento perfeito
para os bonecos; no aprimoramento de sua atuação como personagem; na
assimilação dos tempos na realização de cada parte do número; no
aperfeiçoamento da interação entre seu personagem e os bonecos; e na
melhoria contínua da sua interação com demais artistas, quando atua com
outros titeriteiros ou com coadjuvantes. Treina sua entrada no picadeiro
para iniciar a apresentação. Incorpora equipamentos de segurança no
número, durante o ensaio. Colabora na divulgação do número e do
espetáculo, providenciando material impresso, dando entrevistas e
utilizando jornal, carros de som, redes sociais e outras formas de
comunicação. Prepara corpo e mente para desenvolvimento do espetáculo,
alongando e aquecendo o corpo e realizando exercícios de concentração.
Apresenta o número, interagindo com o público; desenvolvendo formas de
comunicação para conquistar a empatia das pessoas; e usando técnicas e
estratégias para valorização do seu trabalho. Lida com imprevistos,
durante a apresentação, de forma criativa. Interage com os demais
profissionais envolvidos no espetáculo, combinando códigos para informar
imprevistos e ajustando estratégias para lidar com incidentes, acidentes
ou outros acontecimentos. Trabalha em equipe, estabelecendo vínculos de
confiança com os colegas. Pode atuar em teatros -- especialmente em
peças voltadas ao público infantil - e em outros espaços culturais. Pode
vender o número ou espetáculo, avaliando custos para fazer preço do
trabalho e investigando o valor da atividade circense no mercado. Pode
ministrar aulas, promover oficinas e oferecer cursos relacionados à
manipulação de bonecos, para públicos de diferentes faixas etárias.
Seleciona métodos e estratégias para transmissão de conhecimentos e
realização da prática profissional. Motiva o aluno, estimula sua
criatividade para criar números e desenvolve sua capacidade de manipular
bonecos. Transmite a ética circense. Conserva e faz a manutenção dos
bonecos, limpando-os, trocando componentes, verificando as suas
articulações, executando pintura e efetuando reparos, quando necessário.
Zela pela sua segurança, pela segurança da equipe e pela segurança do
público, realizando as atividades em ambiente seguro, contando com
pessoal de apoio capacitado e fazendo uso de equipamentos de segurança.
Pode prestar primeiros socorros.

\begin{Form}
\textbf{Esta correspondência está adequada?} \\ 
\ChoiceMenu[radio,radiosymbol=\ding{52},bordercolor=0 0.5 0,name=cbovalid_ 376250 ]{}{CBO deve ficar=sim, \hspace{6em} CBO não fica aqui=nao}
\end{Form}

\textbf{CBO: 376255  ( Trapezista )}

Cria novo número para apresentações de acrobacias aéreas em trapézio,
pesquisando uso de materiais, aparelhos e tecnologias. Pesquisa
possibilidades no uso das cores, em luzes, figurino e materiais.
Pesquisa movimentos corporais, possibilidades de expressão artística e
formas de comunicação com o público. Faz intercâmbio de informações com
profissionais do circo, com escolas de circo e com profissionais de
outras áreas artísticas, pessoalmente, usando Internet ou outras formas
de comunicação. Pode fazer fusão de números existentes, criando outros.
Utiliza diferentes tipos de trapézio, tais como trapézio fixo (barra de
ferro suspensa por duas cordas, que estão fixas em estrutura no alto),
trapézio múltiplo (construído para suportar duas ou mais pessoas, com
barra metálica maior que a do trapézio simples e suspensa por três,
quatro ou mais cordas), trapézio de voo (composto por uma estrutura
metálica, em que estão pendurados um trapézio simples e um trapézio para
o aparador, além de uma plataforma suspensa, de onde partem e chegam os
artistas que voam e que fazem o balanço), e outros aparelhos similares.
Verifica a fixação do seu aparelho. Monta aparelhos e objetos para o
desenvolvimento seguro do espetáculo, tal como testeira (rede auxiliar
de proteção colocada nas extremidades do trapézio) e cama elástica no
solo, para a segurança dos trapezistas e coadjuvantes. Adapta aparelhos
e equipamentos, considerando seu biótipo; os biótipos de outros
trapezistas ou de coadjuvantes; e as condições do local das
apresentações. Produz o número circense, definindo coreografia e
selecionando música. Cria guarda-roupa e maquiagem. Orienta o trabalho
do iluminador e do músico, indicando o uso de cores e sons para criar
suspense e ressaltar o ápice da apresentação. Ensaia o número,
repetindo-o para aperfeiçoar técnicas e expressões corporais; promover
melhoria contínua da sua interação com demais trapezistas ou
coadjuvantes; e assimilar os tempos na realização das acrobacias no
trapézio. Prepara a forma de sua entrada para iniciar a apresentação.
Incorpora equipamentos de segurança no número, durante o ensaio.
Colabora na divulgação do número e do espetáculo, providenciando
material impresso, dando entrevistas e utilizando jornal, carros de som,
redes sociais e outras formas de comunicação. Prepara corpo e mente para
apresentação do número, alongando e aquecendo o corpo e realizando
exercícios de concentração. Apresenta o número, interagindo com o
público; desenvolvendo formas de comunicação para conquistar a empatia
das pessoas; e usando técnicas e estratégias para valorização do seu
trabalho. Adota os fundamentos técnicos e princípios artísticos que
norteiam as acrobacias nos trapézios, em atuação individual ou em grupo.
Lida com imprevistos, durante a apresentação, de forma criativa.
Interage com os demais profissionais envolvidos no espetáculo,
combinando códigos para informar imprevistos e ajustando estratégias
para lidar com incidentes, acidentes ou outros acontecimentos. Trabalha
em equipe, estabelecendo vínculos de confiança com os colegas. Pode
vender o número ou espetáculo, avaliando custos para fazer preço do
trabalho e investigando o valor da atividade circense no mercado. Pode
ministrar aulas, promover oficinas e oferecer cursos ligados às artes
acrobáticas em trapézio, para públicos de diferentes faixas etárias.
Seleciona métodos e estratégias para transmissão de conhecimentos e
realização da prática profissional. Motiva o aluno e estimula o seu
desenvolvimento físico e a sua destreza. Transmite a ética circense.
Conserva aparelhos e objetos, usados nas acrobacias em trapézios, limpos
e organizados. Verifica a instalação do aparelho e a vida útil de
faixas, roldanas e cordas. Zela pela sua segurança, pela segurança da
equipe e pela segurança do público, realizando as atividades em ambiente
seguro, contando com pessoal de apoio capacitado e fazendo uso de
equipamentos de segurança. Pode prestar primeiros socorros.

\begin{Form}
\textbf{Esta correspondência está adequada?} \\ 
\ChoiceMenu[radio,radiosymbol=\ding{52},bordercolor=0 0.5 0,name=cbovalid_ 376255 ]{}{CBO deve ficar=sim, \hspace{6em} CBO não fica aqui=nao}
\end{Form}

\textbf{CBO: 376325  ( Apresentador de circo )}

Participa no planejamento do itinerário do circo, estudando localidades
onde há terrenos para instalação e estimativa de público, para gerar
receita suficiente para cobrir despesas. Discute suas propostas com
demais integrantes do circo, normalmente integrantes de famílias que se
dedicam à arte circense. Divulga a chegada do circo na localidade,
juntando-se a mágicos, trapezistas, malabaristas, acrobatas,
bilheteiros, vendedores de pipoca e outros componentes da trupe. Usa
linguagem acessível para, em megafone ou microfone, apresentar as
atrações, o local e os horários do espetáculo, atraindo especialmente
crianças. Prepara-se para realizar apresentação do espetáculo de circo,
estudando a programação a ser apresentada. Acompanha os ensaios dos
diferentes artistas, para verificar o que deve ser ressaltado na
apresentação de cada número. Sugere alteração na ordem tradicionalmente
adotada para os números, ajudando a montar o roteiro do espetáculo.
Orienta o trabalho do iluminador e dos músicos da banda, indicando o uso
de cores e sons para aumentar o entusiasmo do público, em cada atração
anunciada. Prepara-se física e mentalmente para realizar a apresentação
do espetáculo de circo, selecionando vestimenta e adereços; controlando
respiração; realizando exercícios de concentração; e aquecendo a voz.
Pode levar em conta, na sua fala, características da plateia e
especificidades do local ou da região. Inicia a apresentação do
espetáculo do circo, ajustando sua postura em cena; saudando o público;
dando explicações introdutórias; realizando pausas para a plateia
assimilar sua fala; e anunciando a primeira atração. Pode improvisar, se
ocorrer situação nova. Conduz o espetáculo, atento ao roteiro
predefinido. Interage com o público, estabelecendo relações de empatia,
criando expectativas e estimulando os aplausos para cada número
anunciado. Controla o tempo da performance de acrobata, contorcionista,
domador de animais, equilibrista, mágico, trapezista, palhaço e outros
artistas. Transmite informações e comunicados relevantes ao
desenvolvimento do espetáculo. Faz adequação de sua postura física e
entonação de voz à informação a ser transmitida. Informa ao público
sobre ocorrências durante o espetáculo, tais como alertas sobre criança
perdida, estacionamento irregular de veículo de pessoa que está na
plateia, entre outras. Faz contato - por meio do olhar - com a plateia,
para perceber suas reações durante o espetáculo. Despede-se da plateia,
no final, solicitando a divulgação do espetáculo para formação de novas
plateias. Pode atuar em circos que não se desloquem para outras
localidades, ficando fixos em parques temáticos, praças, lonas culturais
ou outros locais. Pode atuar em outras configurações de circo,
instalados em espaços abertos ou fechados, com artistas normalmente
formados em cursos de artes circenses.

\begin{Form}
\textbf{Esta correspondência está adequada?} \\ 
\ChoiceMenu[radio,radiosymbol=\ding{52},bordercolor=0 0.5 0,name=cbovalid_ 376325 ]{}{CBO deve ficar=sim, \hspace{6em} CBO não fica aqui=nao}
\end{Form}

\newpage
\section{ 09003 -- Técnico em Artesanato }

\textbf{Perfil do Curso (CNCT)}

O Técnico em Artesanato será habilitado para:

Criar e produzir trabalhos artesanais de peças decorativas e
utilitárias, com materiais diversos e recursos naturais. Explorar a
riqueza e o repertório cultural existentes. Comercializar produtos
artesanais no varejo e no atacado. Gerenciar negócios na perspectiva do
associativismo e do cooperativismo. Selecionar técnicas de tratamento,
preparação e transformação de matérias- primas. Respeitar e valorizar o
traço e a diversidade cultural da região.

Para atuação como Técnico em Artesanato, são fundamentais:

Conhecimentos interdisciplinares relacionados aos processos de criação,
envolvendo pesquisa, idealização, planejamento, execução técnica,
fruição e recepção estética. Competências comunicativas e empreendedoras
voltadas à proposição de projetos, ao coletivo, à gestão, à solução de
problemas e à resiliência, entre outras competências socioemocionais.

\textbf{CBOs correspondidos e seus perfis (presentes no CNCT, e presentes na predição da IA)}

\textbf{CBO: 752105  ( Artesão modelador (vidros) )}

Planeja as atividades de sopro e de moldagem de vidros e cristais,
interpretando as ordens de serviço. Seleciona moldes, equipamentos e
ferramentas de trabalho e verifica suas condições de uso. Requisita
matérias-primas e a quantidade necessária de vidro fundido para as
atividades previstas. Analisa as informações decorrentes da troca de
turnos e as mudanças de processos de produção. Estabelece a sequência de
atividades e distribui tarefas à equipe de trabalho. Controla etapas do
plano de produção, orientando a equipe para seu cumprimento. Modela
objetos de vidro e de cristal manualmente, coletando o vidro fundido com
a cana de sopro; conferindo a viscosidade do vidro; e pré-conformando o
vidro. Dá forma à peça, artesanalmente, por sopro e com o auxílio de
moldes e ferramentas. Aplica detalhes à peça, inclusive com vidros
coloridos. Destaca a peça da cana de sopro. Faz o acabamento. Leva a
peça para o forno de recozimento. Controla a qualidade do produto
acabado e do processo de produção, de acordo com as especificações
contidas na ordem de produção. Controla o peso e as dimensões das peças.
Identifica defeitos e não conformidades nos produtos. Segrega os
produtos não conformes, corrigindo as anomalias do processo de produção
que estejam acarretando defeitos nos produtos. Emite relatórios, com
descrição de defeitos e informações sobre ocorrências do processo.
Registra dados do processo. Conserva o local de trabalho limpo e
organizado. Mantém moldes, ferramentas e equipamentos limpos,
organizados e acondicionados. Comunica-se com setor de manutenção,
descrevendo falhas em equipamentos e ferramentas e solicitando serviços
de reparação. Segrega cacos de vidro, destinando-os para a reciclagem no
próprio processo de fabricação. Monitora o descarte e a destinação de
resíduos de acordo com normas ambientais. Trabalha com segurança,
utilizando equipamentos de proteção individual. Avalia condições de
segurança dos equipamentos e do local de trabalho, informando sobre
condições de riscos. Transporta matérias-primas e produtos de vidro e
cristal de acordo com especificações de segurança.

\begin{Form}
\textbf{Esta correspondência está adequada?} \\ 
\ChoiceMenu[radio,radiosymbol=\ding{52},bordercolor=0 0.5 0,name=cbovalid_ 752105 ]{}{CBO deve ficar=sim, \hspace{6em} CBO não fica aqui=nao}
\end{Form}

\textbf{CBO: 791105  ( Artesão bordador )}

Planeja a criação de produtos bordados artesanais, para atender às
demandas de clientes. Pesquisa tendências e características culturais.
Faz esboço e define especificações do produto. Analisa viabilidade
técnica do projeto, testando materiais, processos e produtos. Seleciona
equipamentos, ferramentas manuais e instrumentos. Pesquisa e seleciona
matérias-primas. Verifica a disponibilidade das matérias-primas
selecionadas, adquirindo as faltantes. Prepara as matérias-primas,
podendo fazer sua lavagem, seu desfiamento, sua fiação e seu tingimento.
Confecciona produtos artesanais, organizando e medindo materiais. Risca,
corta, dobra e costura materiais para confecção de peças de vestuário,
de cama, mesa e banho, dentre outras. Tece, trama e entrelaça materiais
para confecção de artigos como bolsas e adereços. Monta as peças ou os
artigos. Cria o desenho para o bordado, definindo temas e utilizando
cores que marcam sua identidade como artesão e podem retratar a cultura
local ou regional. Faz a adequação do tamanho do desenho à peça ou ao
artigo. Risca o desenho, transferindo os traços para a peça ou o artigo
a bordar. Define a técnica de bordar a ser aplicada, como richelieu,
bordado livre, ponto cruz, vagonite, entre outras. Seleciona uma ou mais
agulhas de bordado, de acordo com a linha a ser usada e com a técnica de
bordado a ser aplicada. Seleciona a linha -- de algodão, lã, seda,
dentre outras -- para bordar. Pode utilizar fitas de cetim ou de seda,
além das linhas, para realização do bordado. Executa o bordado de acordo
com desenhos, aplicando a técnica selecionada. Pode acrescentar
elementos decorativos - tais como rendas, miçangas, lantejoulas, pérolas
ou pedrarias - para complementar o bordado. Pode usar dedal ou dedeira
para proteger o dedo que empurra a agulha, durante a execução do
bordado. Pode utilizar bastidor, fixando o tecido para facilitar a
realização dos pontos do bordado. Realiza acabamento da peça ou do
artigo bordado, fazendo arremates. Identifica possíveis imprecisões,
efetuando correções. Passa peça ou artigo bordado, eliminando dobras ou
amassados. Pode executar engomagem. Pode emoldurar peças bordadas.
Prende etiqueta e acondiciona os produtos, preparando-os para a
comercialização. Verifica a qualidade do processo e do produto,
inspecionando o bordado em relação a cores; formas; dimensões;
disposição e posição na peça ou no artigo; contornos; detalhes da trama;
ausência de fios soltos; e perfeição da trama no avesso. Comercializa
produtos bordados, desenvolvendo embalagens, calculando custo dos
produtos e estabelecendo preços de venda. Pesquisa mercado consumidor.
Divulga produtos, podendo fazer uso inclusive de ``sites'' e redes
sociais. Participa de feiras e eventos, expondo produtos. Aplica
estratégias de venda, negocia valores e prazos com cliente e embala
produtos vendidos. Gerencia o próprio negócio, organizando produção,
controlando estoques e administrando orçamento. Define o leiaute de
oficina ou ateliê. Estabelece parcerias e seleciona fornecedores. Pode
capacitar artesãos. Pode vincular-se a uma associação de classe. Realiza
manutenção básica de equipamentos, fazendo sua limpeza e lubrificação.
Providencia serviços de manutenção corretiva. Conserva o local de
trabalho limpo e organizado. Mantém ferramentas, instrumentos e
acessórios limpos, organizados, acondicionados e em plenas condições de
uso e funcionamento. Evita desperdícios, separando sobras de materiais
para reaproveitamento. Realiza descarte de resíduos, de maneira
ambientalmente adequada. Segue normas de segurança no trabalho e
orientações de ergonomia.

\begin{Form}
\textbf{Esta correspondência está adequada?} \\ 
\ChoiceMenu[radio,radiosymbol=\ding{52},bordercolor=0 0.5 0,name=cbovalid_ 791105 ]{}{CBO deve ficar=sim, \hspace{6em} CBO não fica aqui=nao}
\end{Form}

\textbf{CBO: 791110  ( Artesão ceramista )}

Planeja a criação de produtos artesanais de cerâmica - mantendo estilo
que marca sua identidade como artesão e pode retratar a cultura local ou
regional -, para atender às demandas de clientes. Pesquisa tendências e
características culturais. Faz esboço de peças utilitárias -- como
tigelas, panelas, xícaras e outros utensílios -- ou peças decorativas,
tais como figuras típicas regionais (como cangaceiros e rendeiras),
religiosas, da natureza (como flores e animais), dentre outras. Define
especificações do produto. Analisa viabilidade técnica do projeto,
testando materiais, processos e produtos. Seleciona equipamentos,
ferramentas manuais e instrumentos. Pesquisa e seleciona matéria-prima.
Coleta ou adquire a matéria-prima selecionada. Pode reciclar materiais.
Prepara a matéria-prima, misturando argilas de diversos tipos e cores -
para atingir coloração personalizada e de acordo com a peça a ser
produzida -; e adicionando água para formar a massa. Trata a massa,
refinando-a; retirando as impurezas; amassando-a e sovando-a; e
deixando-a totalmente livre de bolhas. Prepara moldes e formas para
confecção da peça. Confecciona a peça artesanal de cerâmica, modelando a
massa de argila com aplicação de técnicas de modelagem livre ou com uso
de torno manual ou elétrico. Monta a peça, ajustando as dimensões,
perfurando e colando partes da peça. Realiza a secagem da peça,
estabelecendo o tempo necessário. Leva a peça ao forno para a primeira
queima -- denominada queima de biscoito --, para que fique resistente e
porosa para receber esmalte cerâmico. Cria desenho ou padrão para
esmaltação da peça cerâmica, definindo temáticas e mistura de cores.
Prepara a peça para esmaltação, limpando sujeira superficial; lixando
protuberância ou irregularidade; e polindo. Realiza esmaltação da peça,
mergulhando-a no esmalte ou usando pincel para fazer a aplicação, de
acordo com desenho ou padrão. Leva a peça novamente ao forno. Controla a
temperatura do novo processo de queima. Finaliza a peça, realizando
limpeza e fazendo acabamento. Verifica a qualidade do processo e do
produto, considerando parâmetros operacionais, estéticos e de utilidade
do produto. Inspeciona a peça, considerando formas, cores, dimensões e
ausência de trincas ou rachaduras. Analisa a resistência do material a
impactos e a altas temperaturas. Embala e identifica as peças,
preparando-as para a comercialização. Comercializa produtos artesanais
de cerâmica, desenvolvendo embalagens, calculando custo dos produtos e
estabelecendo preços de venda. Pesquisa mercado consumidor. Divulga
produtos, podendo fazer uso inclusive de ``sites'' e redes sociais.
Participa de feiras e eventos, expondo produtos. Aplica estratégias de
venda, negocia valores e prazos com cliente e embala produtos vendidos.
Gerencia o próprio negócio, organizando produção, controlando estoques e
administrando orçamento. Define o leiaute de oficina ou ateliê.
Estabelece parcerias e seleciona fornecedores. Pode capacitar artesãos.
Pode vincular-se a uma associação de classe. Realiza manutenção básica
de equipamentos, fazendo sua limpeza e lubrificação. Providencia
serviços de manutenção corretiva. Conserva o local de trabalho limpo e
organizado. Mantém ferramentas, instrumentos e acessórios limpos,
organizados, acondicionados e em plenas condições de uso e
funcionamento. Evita desperdícios, separando sobras de materiais para
reaproveitamento. Realiza descarte de resíduos, de maneira
ambientalmente adequada. Trabalha com segurança, prevenindo acidentes e
incêndios. Utiliza equipamentos de proteção individual, principalmente
em ambientes sujeitos a poeiras e altas temperaturas.

\begin{Form}
\textbf{Esta correspondência está adequada?} \\ 
\ChoiceMenu[radio,radiosymbol=\ding{52},bordercolor=0 0.5 0,name=cbovalid_ 791110 ]{}{CBO deve ficar=sim, \hspace{6em} CBO não fica aqui=nao}
\end{Form}

\textbf{CBO: 791115  ( Artesão com material reciclável )}

Planeja a criação de produtos artesanais de materiais reciclados -
mantendo estilo que marca sua identidade como artesão -, para atender às
demandas de clientes. Elabora esboço e define especificações do produto.
Analisa viabilidade técnica do projeto, testando materiais, processos e
produtos. Seleciona equipamentos, ferramentas manuais e instrumentos.
Pesquisa e seleciona materiais recicláveis. Reaproveita materiais - tais
como metal, vidro, papelão, madeira, cortiça, tecido, entre outros --
selecionados para confecção do produto. Prepara materiais
reaproveitados, limpando-os e fazendo sua lavagem e secagem. Pode fazer
cozimento, desfiamento, amaciamento, tingimento e impermeabilização de
materiais. Pode enrolar e lixar materiais. Confecciona a peça artesanal
- utilitária ou decorativa -, realizando operações de acordo com o
produto a ser confeccionado e o tipo de material reaproveitado. Pode
riscar, cortar, soldar, fundir, trefilar e laminar metais; pintar
vidros; classificar, misturar em água, retirar impurezas, prensar e
secar papéis; entalhar madeira; esculpir cortiça; tecer, trançar e
entrelaçar fios; medir, cortar, modelar e costurar tecidos; entre outros
processos. Finaliza o produto, fazendo vitrificação; esmaltação;
emborrachamento; queima; incrustação de elemento decorativo; ou outro
tipo de acabamento. Executa ajustes, correções e eliminação de rebarbas
ou amassados. Embala e identifica o produto, preparando-o para a
comercialização. Verifica a qualidade do processo e do produto,
considerando parâmetros operacionais, estéticos e de utilidade.
Inspeciona o produto, considerando formas, cores, dimensões e fixação
das partes. Examina ausência de trincas, amassados, rachaduras e outras
possíveis imperfeições. Analisa a resistência do material a impactos, a
variações térmicas e a outros fenômenos físicos ou químicos.
Comercializa produtos artesanais de materiais reciclados, desenvolvendo
embalagens, calculando custo dos produtos e estabelecendo preços de
venda. Pesquisa mercado consumidor. Divulga produtos, podendo fazer uso
inclusive de ``sites'' e redes sociais. Participa de feiras e eventos,
expondo produtos. Aplica estratégias de venda, negocia valores e prazos
com cliente e embala produtos vendidos. Gerencia o próprio negócio,
organizando produção, controlando estoques e administrando orçamento.
Define o leiaute de oficina ou ateliê. Estabelece parcerias e seleciona
fornecedores. Pode capacitar artesãos. Pode vincular-se a uma associação
de classe. Realiza manutenção básica de equipamentos, fazendo sua
limpeza e lubrificação. Providencia serviços de manutenção corretiva.
Conserva o local de trabalho limpo e organizado. Mantém ferramentas,
instrumentos e acessórios limpos, organizados, acondicionados e em
plenas condições de uso e funcionamento. Realiza descarte de resíduos,
de maneira ambientalmente adequada. Trabalha com segurança, usando
equipamentos de proteção individual. Identifica e elimina condições
inseguras no local de trabalho, prevenindo acidentes. Observa princípios
básicos de ergonomia.

\begin{Form}
\textbf{Esta correspondência está adequada?} \\ 
\ChoiceMenu[radio,radiosymbol=\ding{52},bordercolor=0 0.5 0,name=cbovalid_ 791115 ]{}{CBO deve ficar=sim, \hspace{6em} CBO não fica aqui=nao}
\end{Form}

\textbf{CBO: 791120  ( Artesão confeccionador de biojoias e ecojoias )}

Planeja a criação de biojoias (com uso de recursos naturais, como
sementes, fibras, madeira, conchas, pedras, entre outros) e ecojoias
(com garrafas feitas de PET-Polietileno Tereftalato; tecidos; metais não
ferrosos; couro; entre outros materiais descartados, que passam por
processo -- ``upcycling'' - de reutilização criativa), para atendimento
às demandas de clientes. Elabora esboços e define especificações do
produto. Analisa viabilidade técnica do projeto, testando materiais,
processos e produtos. Seleciona equipamentos, ferramentas manuais e
instrumentos. Pesquisa e seleciona matérias-primas, coletando ou
adquirindo os recursos naturais e/ou materiais descartados selecionados.
Prepara matérias-primas, limpando-as e fazendo sua lavagem e secagem.
Pode escovar, tingir, lixar e impermeabilizar matéria-prima. Pode
utilizar estufa para a secagem e câmara de luz para a eliminação de
fungos de sementes e outros recursos naturais. Aplica técnicas e
processos de produção artesanal na confecção de biojoias e ecojoias.
Define materiais, componentes e acessórios de acordo com desenhos e
especificações do projeto. Mede e corta materiais. Pode soldar, fundir,
trefilar e laminar metais reutilizados; perfurar e pintar sementes e
conchas; polir pedras; entalhar madeiras; esculpir cortiça; dobrar e
gravar couro descartado; entre outros processos. Monta a peça,
atribuindo ao objeto características artísticas pessoais e da cultura
local ou regional. Finaliza o produto, fazendo gravações; efetuando
incrustação de elemento decorativo; e realizando outros tipos de
acabamento. Executa ajustes e correções na peça. Verifica a qualidade do
processo e do produto, considerando parâmetros operacionais e estéticos.
Inspeciona a peça, verificando formas, cores, dimensões e fixação das
partes. Examina ausência de trincas, amassados e outras possíveis
imperfeições. Analisa a resistência do material a impactos e a outros
fenômenos físicos ou químicos. Coloca etiqueta no produto, preparando-o
para a comercialização. Pode identificar biojoia e ecojoia com
logomarca. Comercializa biojoias e ecojoias, desenvolvendo embalagens --
normalmente criativas e com apelo ecológico -; calculando custo dos
produtos; e estabelecendo preços de venda. Pesquisa mercado consumidor.
Divulga produtos, podendo fazer uso inclusive de ``sites'' e redes
sociais. Participa de feiras e eventos -- especialmente os relacionados
à sustentabilidade -, expondo produtos. Aplica estratégias de venda,
negocia valores e prazos com cliente e embala produtos vendidos.
Gerencia o próprio negócio, organizando produção, controlando estoques e
administrando orçamento. Define o leiaute de oficina ou ateliê.
Estabelece parcerias e seleciona fornecedores. Pode capacitar artesãos.
Pode vincular-se a uma associação de classe. Realiza manutenção básica
de equipamentos, fazendo sua limpeza e lubrificação. Providencia
serviços de manutenção corretiva. Conserva o local de trabalho limpo e
organizado. Mantém ferramentas, instrumentos e acessórios limpos,
organizados, acondicionados e em plenas condições de uso e
funcionamento. Zela pela preservação ambiental, coletando e extraindo
materiais da natureza - dentro de parâmetros de responsabilidade com o
meio ambiente - e reutilizando materiais descartados, para confecção de
um novo produto. Realiza descarte de resíduos, de maneira ambientalmente
adequada. Trabalha com segurança, usando equipamentos de proteção
individual. Identifica e elimina condições inseguras no local de
trabalho, prevenindo acidentes. Observa princípios básicos de ergonomia.

\begin{Form}
\textbf{Esta correspondência está adequada?} \\ 
\ChoiceMenu[radio,radiosymbol=\ding{52},bordercolor=0 0.5 0,name=cbovalid_ 791120 ]{}{CBO deve ficar=sim, \hspace{6em} CBO não fica aqui=nao}
\end{Form}

\textbf{CBO: 791125  ( Artesão do couro )}

Planeja a criação de produtos artesanais de couro -- tais como bolsas;
carteiras; porta-moedas; cintos; selas, arreios e outros artefatos de
montaria; peças de vestuário; e adornos -, para atender às demandas de
clientes. Pesquisa tendências e características culturais. Elabora
esboços e define especificações do produto. Analisa viabilidade técnica
do projeto, testando materiais, processos e produtos. Seleciona
equipamentos, ferramentas e instrumentos. Pesquisa e seleciona
matérias-primas. Verifica a disponibilidade das matérias-primas
selecionadas, adquirindo as faltantes. Prepara as matérias-primas,
limpando-as e fazendo sua lavagem e secagem. Pode fazer cozimento,
curtimento, amaciamento, tingimento e impermeabilização de materiais.
Pode escovar, enrolar e lixar materiais. Aplica técnicas e processos de
produção artesanal na confecção das partes e do produto de couro. Pode
medir, modelar, riscar, cortar e costurar materiais. Pode perfurar,
martelar, dobrar, colar e trançar materiais. Monta peças. Pode fazer
pinturas, executar gravações e confeccionar bordados nas peças montadas.
Finaliza o produto, fazendo emborrachamento; executando enceramento;
efetuando incrustação de elemento decorativo; ou realizando outro tipo
de acabamento. Executa arremates. Verifica a qualidade do processo e do
produto, considerando parâmetros operacionais e estéticos. Inspeciona a
peça, verificando formas, cores e dimensões; conferindo fixação das
partes; e observando demais detalhes. Analisa a resistência do material
a impactos e a outros fenômenos físicos ou químicos. Comercializa
produtos artesanais de couro, desenvolvendo embalagens, calculando custo
dos produtos e estabelecendo preços de venda. Pesquisa mercado
consumidor. Divulga produtos, podendo fazer uso inclusive de ``sites'' e
redes sociais. Participa de feiras e eventos, expondo produtos. Aplica
estratégias de venda; negocia valores e prazos com cliente; e embala
produtos vendidos. Pode consertar, recuperar ou restaurar peças e
objetos de couro. Gerencia o próprio negócio, organizando produção,
controlando estoques e administrando orçamento. Define o leiaute de
oficina ou ateliê. Estabelece parcerias e seleciona fornecedores. Pode
capacitar artesãos. Pode vincular-se a uma associação de classe. Realiza
manutenção básica de equipamentos, fazendo sua limpeza e lubrificação.
Providencia serviços de manutenção corretiva. Conserva o local de
trabalho limpo e organizado. Mantém ferramentas, instrumentos e
acessórios limpos, organizados, acondicionados e em plenas condições de
uso e funcionamento. Evita desperdícios, separando sobras de materiais
para reaproveitamento. Realiza descarte de resíduos, de maneira
ambientalmente adequada. Trabalha com segurança, usando equipamentos de
proteção individual. Identifica e elimina condições inseguras no local
de trabalho, prevenindo acidentes. Observa princípios básicos de
ergonomia.

\begin{Form}
\textbf{Esta correspondência está adequada?} \\ 
\ChoiceMenu[radio,radiosymbol=\ding{52},bordercolor=0 0.5 0,name=cbovalid_ 791125 ]{}{CBO deve ficar=sim, \hspace{6em} CBO não fica aqui=nao}
\end{Form}

\textbf{CBO: 791130  ( Artesão escultor )}

Planeja a confecção de peças de escultura, para atendimento às demandas
de clientes. Elabora desenhos, definindo as características técnicas e
estéticas do objeto a ser criado. Analisa viabilidade técnica do
projeto, testando materiais, processos e produtos. Seleciona
equipamentos, ferramentas (como formão, rebarbadora de ângulos e
estecas) e instrumentos (tais como cinzel e espátulas). Pesquisa e
seleciona matérias-primas para confecção da escultura, coletando ou
adquirindo argila, madeira, gesso, pedras -- como mármore, alabastro,
granito, arenito e calcário --, bronze, látex, entre outros. Define
elementos de união e fixação. Prepara matérias-primas, limpando-as. Pode
cozer, tingir e lixar matéria-prima. Prepara massas. Pode defumar látex.
Aplica técnicas artesanais para transformar a matéria-prima em objeto de
arte, atribuindo-lhe características artísticas pessoais. Confecciona
escultura, criando moldes, tirando medidas, riscando materiais e fazendo
modelagem. Usa ferramentas e instrumentos para cinzelar, talhar e
desbastar materiais. Faz aglomeração de partículas. Cola ou solda
materiais. Funde, trefila e lamina metais. Monta peças. Finaliza formas
ou figuras tridimensionais, tais como bustos, animais, símbolos, imagens
religiosas, objetos decorativos, entre outras. Pinta e grava materiais.
Incrusta e engasta elementos decorativos. Verifica a presença de
possíveis imperfeições, efetuando correções e ajustes. Controla a
qualidade do processo e do produto, considerando parâmetros
operacionais, estéticos e de precisão dos detalhes. Inspeciona a peça,
considerando formas, cores e dimensões. Examina ausência de trincas ou
rachaduras. Analisa a resistência do material a impactos, a intempéries
e a efeitos físicos e químicos. Comercializa esculturas, desenvolvendo
embalagens; calculando custo dos produtos; e estabelecendo preços de
venda. Pesquisa mercado consumidor e divulga produtos. Participa de
feiras e eventos, expondo produtos. Aplica estratégias de venda, negocia
valores e prazos com cliente e embala produtos vendidos. Gerencia o
próprio negócio, organizando armazenagem de materiais e produtos,
controlando estoques e administrando orçamento. Define o leiaute de
oficina ou ateliê. Estabelece parcerias e seleciona fornecedores. Pode
capacitar artesãos. Pode vincular-se a uma associação de classe. Realiza
manutenção básica de equipamentos, fazendo sua limpeza e lubrificação.
Providencia serviços de manutenção corretiva. Conserva o local de
trabalho limpo e organizado. Mantém ferramentas, instrumentos e
acessórios limpos, organizados, acondicionados e em plenas condições de
uso e funcionamento. Zela pela preservação ambiental, realizando
reaproveitamento de sobras de materiais. Faz descarte de resíduos de
maneira ambientalmente adequada. Trabalha com segurança, usando
equipamentos de proteção individual, tais como protetor auricular,
respirador e óculos de segurança. Identifica e elimina condições
inseguras no local de trabalho, prevenindo acidentes. Observa princípios
básicos de ergonomia.

\begin{Form}
\textbf{Esta correspondência está adequada?} \\ 
\ChoiceMenu[radio,radiosymbol=\ding{52},bordercolor=0 0.5 0,name=cbovalid_ 791130 ]{}{CBO deve ficar=sim, \hspace{6em} CBO não fica aqui=nao}
\end{Form}

\textbf{CBO: 791135  ( Artesão moveleiro (exceto reciclado) )}

Planeja a produção artesanal de peças utilitárias e decorativas e
móveis, para atendimento às demandas de clientes. Elabora desenhos,
definindo as características técnicas e estéticas dos objetos a serem
criados. Analisa viabilidade técnica do projeto, testando materiais,
processos e produtos. Seleciona equipamentos, ferramentas e
instrumentos. Pesquisa e seleciona matérias-primas para confecção de
peças e móveis. Verifica a disponibilidade das matérias-primas
selecionadas, adquirindo as faltantes. Define elementos de união e
fixação. Prepara matérias-primas, limpando-as. Pode cozer, desfibrar,
amaciar e lixar materiais. Aplica técnicas artesanais para transformar
madeira em peças e móveis, com estilo que marca sua identidade como
artesão. Confecciona peças e móveis, medindo, moldando, modelando,
riscando e cortando materiais. Pode perfurar, colar e dobrar. Pode
trançar, tramar e entrelaçar materiais. Monta peças e móveis. Finaliza
peças e móveis, lixando, polindo e colocando acessórios. Pode entalhar
madeira. Pode pintar e envernizar. Verifica a presença de possíveis
imperfeições, efetuando correções e ajustes. Controla a qualidade do
processo e do produto, considerando parâmetros operacionais, estéticos e
de precisão dos detalhes. Inspeciona os produtos, considerando formas,
cores e dimensões. Examina ausência de trincas ou rachaduras. Analisa a
resistência do material a impactos, a intempéries e a efeitos físicos e
químicos. Coloca etiqueta em peças e móveis, preparando-os para a
comercialização. Pode usar logomarca, para identificação dos produtos.
Comercializa peças utilitárias e decorativas e móveis - feitos
artesanalmente em madeira -, calculando custo dos produtos; e
estabelecendo preços de venda. Pesquisa mercado consumidor e divulga
produtos. Participa de feiras e eventos, expondo produtos. Aplica
estratégias de venda, negocia valores e prazos com cliente e embala
produtos vendidos. Gerencia o próprio negócio, organizando produção;
controlando estoques; e administrando orçamento. Define o leiaute de
oficina ou ateliê. Pode estabelecer parceria com ``designer'', para
troca de conhecimentos e experiências. Seleciona fornecedores. Pode
capacitar artesãos. Pode vincular-se a uma associação de classe. Realiza
manutenção básica de equipamentos, fazendo sua limpeza e lubrificação.
Providencia serviços de manutenção corretiva. Conserva o local de
trabalho limpo e organizado. Mantém ferramentas, instrumentos e
acessórios limpos, organizados, acondicionados e em plenas condições de
uso e funcionamento. Zela pela preservação ambiental, realizando
reaproveitamento de sobras de materiais. Faz descarte de resíduos de
maneira ambientalmente adequada. Trabalha com segurança, usando
equipamentos de proteção individual, tais como luvas, protetor auricular
e óculos de segurança. Identifica e elimina condições inseguras no local
de trabalho, prevenindo acidentes. Observa princípios básicos de
ergonomia.

\begin{Form}
\textbf{Esta correspondência está adequada?} \\ 
\ChoiceMenu[radio,radiosymbol=\ding{52},bordercolor=0 0.5 0,name=cbovalid_ 791135 ]{}{CBO deve ficar=sim, \hspace{6em} CBO não fica aqui=nao}
\end{Form}

\textbf{CBO: 791140  ( Artesão tecelão )}

Planeja a confecção de peça de tecido, para atendimento às demandas de
cliente. Pesquisa tendências e características culturais. Elabora
esboços e define especificações do produto. Analisa viabilidade técnica
do projeto, testando materiais, processos e produtos. Seleciona
equipamentos, ferramentas e instrumentos. Pesquisa e seleciona
matérias-primas. Pode coletar matérias-primas ou adquiri-las,
verificando sua procedência. Prepara as matérias-primas, limpando-as e
fazendo sua lavagem e secagem. Pode desfolhar, desfibrar e desfiar
matéria-prima. Pode executar cozimento, amaciamento, tingimento e
escovação de materiais. Pode cardar, fiar e enrolar matéria-prima.
Confecciona tecidos artesanais, aplicando técnicas e métodos para
realizar operações de tecelagem. Executa operações referentes à
preparação da urdidura, ajustando parâmetros, tensão dos fios,
comprimento, largura e centralização. Coloca a urdidura no tear. Opera
tear, passando o fio da trama entre os fios da urdidura, para tecer a
peça de tecido. Pode usar diferentes cores, emendando fios de acordo com
a padronagem definida. Finaliza o produto, fazendo arremates. Pode
confeccionar bordados, para enfeitar o tecido. Pode aplicar adornos e
acessórios. Pode montar peças -- como bolsas, cortinas e tapetes -,
cortando e costurando tecidos, de acordo com esboços e modelos. Controla
a qualidade do processo e do produto, verificando precisão da padronagem
e da textura; formas; cores; dimensões; e união das partes. Inspeciona
padrão de caimento, vincos e pregas. Examina imperfeições, efetuando
ajustes ou reparos. Analisa a resistência do produto confeccionado à
lavagem, à exposição solar e a demais fenômenos físicos ou químicos.
Coloca etiqueta nos tecidos confeccionados, preparando-os para a
comercialização. Pode usar logomarca, para identificação dos produtos.
Comercializa peças artesanais de tecido, desenvolvendo embalagens,
calculando custo dos produtos e estabelecendo preços de venda. Pesquisa
mercado consumidor. Divulga produtos, podendo fazer uso inclusive de
``sites'' e redes sociais. Participa de feiras e eventos, expondo
produtos. Aplica estratégias de venda; negocia valores e prazos com
cliente; e embala os produtos vendidos. Gerencia o próprio negócio,
organizando produção, controlando estoques e administrando orçamento.
Define o leiaute de oficina ou ateliê. Estabelece parcerias e seleciona
fornecedores. Pode capacitar artesãos. Pode vincular-se a uma associação
de classe. Realiza manutenção básica de equipamentos, fazendo sua
limpeza e lubrificação. Providencia serviços de manutenção corretiva.
Conserva o local de trabalho limpo e organizado. Mantém ferramentas,
instrumentos e acessórios limpos, organizados, acondicionados e em
plenas condições de uso e funcionamento. Evita desperdícios, separando
sobras de materiais para reaproveitamento. Realiza descarte de resíduos,
de maneira ambientalmente adequada. Trabalha com segurança, usando
equipamentos de proteção individual. Identifica e elimina condições
inseguras no local de trabalho, prevenindo acidentes. Observa princípios
básicos de ergonomia.

\begin{Form}
\textbf{Esta correspondência está adequada?} \\ 
\ChoiceMenu[radio,radiosymbol=\ding{52},bordercolor=0 0.5 0,name=cbovalid_ 791140 ]{}{CBO deve ficar=sim, \hspace{6em} CBO não fica aqui=nao}
\end{Form}

\textbf{CBO: 791145  ( Artesão trançador )}

Planeja a confecção de produtos artesanais trançados - utilitários ou
decorativos - para atendimento às demandas de cliente. Pesquisa
tendências e características culturais. Elabora esboços e define
especificações do produto. Analisa viabilidade técnica do projeto,
testando materiais, processos e produtos. Seleciona equipamentos,
ferramentas e instrumentos. Pesquisa e seleciona matérias-primas, como
vime, capim dourado, bambu, palha, tiras de tecidos, tiras de couro cru,
pelos da crina de cavalo, entre outras. Define acessórios, guarnições,
elementos de união e fixação e outros materiais necessários. Pode
coletar matérias-primas ou adquiri-las, verificando sua procedência.
Prepara as matérias-primas, limpando-as e fazendo sua lavagem e secagem.
Pode desfolhar, desfibrar, desfiar e enrolar matérias-primas. Pode
executar amaciamento, tingimento e escovação. Confecciona produtos
artesanais trançados - como redes, bolsas, peneiras, rédeas para cavalo,
chapéus, cestas, entre outros -, de acordo com desenhos, especificações
e modelos. Mede e corta materiais. Trança, trama e entrelaça, podendo
usar diferentes cores de materiais, conforme padronagem definida.
Realiza operações com as peças trançadas, dobrando e perfurando
materiais; costurando ou colando para unir partes; entre outras. Monta o
produto, atribuindo ao objeto características artísticas pessoais e/ou
da cultura local ou regional. Finaliza o produto, fazendo arremates.
Pode confeccionar bordados, para enfeitar o produto. Pode aplicar
adornos e acessórios. Pode passar e engomar. Controla a qualidade do
processo e do produto, verificando precisão de detalhes e da trama do
trançado. Inspeciona formas; cores; dimensões; e união das partes.
Examina imperfeições, efetuando ajustes ou reparos. Analisa a
resistência do produto confeccionado à umidade, à exposição solar e a
outros fenômenos físicos ou químicos. Coloca etiquetas nas peças
confeccionadas, preparando-as para a comercialização. Pode usar
logomarca, para identificação dos produtos. Adota procedimentos para
acondicionamento, limpeza e conservação de peças artesanais trançadas,
para evitar avarias, deformação e demais efeitos externos. Comercializa
produtos artesanais trançados, desenvolvendo embalagens, calculando
custo dos produtos e estabelecendo preços de venda. Pesquisa mercado
consumidor. Divulga produtos, podendo fazer uso inclusive de ``sites'' e
redes sociais. Participa de feiras e eventos, expondo produtos. Aplica
estratégias de venda; negocia valores e prazos com cliente; e embala os
produtos vendidos. Pode consertar, recuperar ou restaurar peças
artesanais trançadas. Gerencia o próprio negócio, organizando produção,
controlando estoques e administrando orçamento. Define o leiaute de
oficina ou ateliê. Estabelece parcerias e seleciona fornecedores. Pode
capacitar artesãos. Pode vincular-se a uma associação de classe. Realiza
manutenção básica de equipamentos, fazendo sua limpeza e lubrificação.
Providencia serviços de manutenção corretiva. Conserva o local de
trabalho limpo e organizado. Mantém ferramentas, instrumentos e
acessórios limpos, organizados, acondicionados e em plenas condições de
uso e funcionamento. Evita desperdícios, separando sobras de materiais
para reaproveitamento. Realiza descarte de resíduos, de maneira
ambientalmente adequada. Segue normas e procedimentos de saúde e
segurança no trabalho. Identifica e elimina condições inseguras no local
de trabalho, prevenindo acidentes. Observa princípios básicos de
ergonomia.

\begin{Form}
\textbf{Esta correspondência está adequada?} \\ 
\ChoiceMenu[radio,radiosymbol=\ding{52},bordercolor=0 0.5 0,name=cbovalid_ 791145 ]{}{CBO deve ficar=sim, \hspace{6em} CBO não fica aqui=nao}
\end{Form}

\textbf{CBO: 791150  ( Artesão crocheteiro )}

Planeja a confecção de produtos artesanais de crochê, para atendimento
às demandas de cliente. Pesquisa tendências e características culturais.
Seleciona equipamentos, ferramentas e instrumentos. Interpreta gráficos,
simbologia de crochê e descritivo de receitas. Realiza pesquisas sobre
tipos diferenciados de crochê -- como o crochê tunisiano, o crochê
peruano e o crochê de grampo -; diferentes modelos de peças; e novas
combinações de cores. Define o produto artesanal de crochê, elaborando
esboços e detalhando especificações. Atribui, ao produto a ser
confeccionado, características artísticas pessoais. Pode incorporar, ao
produto, características típicas da cultura local ou regional. Prepara a
confecção da peça, escolhendo fios - como de algodão, de lã, de
barbante, dentre outros -; definindo o tipo e as cores de fios;
selecionando o tipo e o tamanho da agulha de gancho, com base na
espessura do fio; e estabelecendo os tipos de ponto -- correntinha
(ponto base); ponto alto; ponto baixo; ponto baixíssimo; ponto pipoca;
ponto zigue-zague; entre outros -- a serem utilizados. Faz uma amostra
com o fio, as cores e os pontos escolhidos, para examinar o resultado,
antes de começar a peça. Confecciona vestido, blusa, cachecol, gorro,
colcha, almofada, tapete ou outro produto artesanal de crochê. Entrelaça
fio na agulha, iniciando a correntinha de base, fazendo as carreiras e
compondo a trama. Finaliza a peça, fazendo arremates. Costura as
diversas partes que compõem o produto. Aplica adornos e acessórios. Pode
engomar a peça. Controla a qualidade do processo e do produto,
verificando a precisão e a abertura dos pontos, da trama e das costuras;
observando os espaços entre os pontos; e conferindo as formas, as cores
e as dimensões. Verifica a união e o agrupamento das partes. Examina
imperfeições, efetuando ajustes ou reparos. Analisa a resistência do
produto confeccionado à umidade e à exposição solar. Coloca etiquetas
nos produtos confeccionados, preparando-os para a comercialização. Pode
usar logomarca, para identificação dos produtos. Adota procedimentos
para acondicionamento, limpeza e conservação de peças artesanais de
crochê, para evitar avarias e deformação. Comercializa produtos
artesanais de crochê, desenvolvendo embalagens, calculando custo dos
produtos e estabelecendo preços de venda. Pesquisa mercado consumidor.
Divulga produtos, podendo fazer uso inclusive de ``sites'' e redes
sociais. Participa de feiras e eventos, expondo produtos. Aplica
estratégias de venda; negocia valores e prazos com cliente; e embala os
produtos vendidos. Gerencia o próprio negócio, organizando produção,
controlando estoques e administrando orçamento. Define o leiaute de
oficina ou ateliê. Estabelece parcerias e seleciona fornecedores. Pode
capacitar artesãos. Pode vincular-se a uma associação de classe. Realiza
manutenção básica de equipamentos, fazendo sua limpeza e lubrificação.
Providencia serviços de manutenção corretiva. Mantém a limpeza e a
organização do local de trabalho. Organiza e acondiciona materiais,
ferramentas e instrumentos. Evita desperdícios, separando sobras de
materiais para reaproveitamento. Realiza descarte de resíduos, de
maneira ambientalmente adequada. Segue normas e procedimentos de saúde e
segurança no trabalho. Identifica e elimina condições inseguras no local
de trabalho, prevenindo acidentes. Reduz o impacto de repetição de
movimentos com as mãos por tempo prolongado -- que pode gerar câimbras,
dores e lesões --, fazendo pausas regulares durante o trabalho;
realizando exercício de alongamento das mãos; e utilizando agulhas com
cabo ergonômico (formato que se encaixa melhor nas mãos), para diminuir
esforço na execução dos pontos. Observa princípios básicos de ergonomia.

\begin{Form}
\textbf{Esta correspondência está adequada?} \\ 
\ChoiceMenu[radio,radiosymbol=\ding{52},bordercolor=0 0.5 0,name=cbovalid_ 791150 ]{}{CBO deve ficar=sim, \hspace{6em} CBO não fica aqui=nao}
\end{Form}

\textbf{CBO: 791155  ( Artesão tricoteiro )}

Planeja a confecção de produtos artesanais de tricô, para atendimento às
demandas de cliente. Pesquisa tendências e características culturais.
Seleciona equipamentos, ferramentas e instrumentos. Interpreta
simbologia, diagramas e receitas de tricô. Realiza pesquisas sobre novos
modelos, novas texturas e colorações diferenciadas de peças em tricô.
Define o produto artesanal de tricô, elaborando esboços e detalhando
especificações. Atribui, ao produto a ser confeccionado, características
artísticas pessoais. Pode incorporar, ao produto, características
típicas da cultura local ou regional. Prepara a confecção da peça,
escolhendo linhas ou fios - como de algodão, lã, acrílico, dentre outros
-; e definindo suas cores. Seleciona agulhas (de madeira, alumínio,
plástico ou outro material), definindo seu estilo (circular ou reta),
verificando seu comprimento (conforme tamanho da peça a tricotar) e
estabelecendo seu diâmetro (com base no uso de linha mais grossa ou mais
fina e no volume desejado para a peça). Escolhe tipo(s) de ponto - tais
como ponto tricô, ponto meia, ponto mate simples, ponto mate duplo,
ponto arroz, ponto barra, entre outros -, de acordo com a peça a ser
tricotada. Faz uma amostra com a linha, as cores e os pontos escolhidos,
para examinar o resultado, antes de tricotar a peça. Confecciona blusa,
manta, casaco, gorro, cachecol, meias, sapatinhos de bebê ou outro
produto artesanal de tricô. Coloca o número calculado de pontos - de
acordo com a dimensão e formato da peça - na agulha; entrelaça a linha
na agulha, conforme tipo(s) de ponto escolhido(s); e faz as carreiras
para compor a trama. Executa emendas de linhas, trocando as cores
conforme modelo e padronagem da peça. Controla as dimensões da peça.
Finaliza a peça, fazendo arremates. Costura as diversas partes que
compõem o produto. Pode aplicar adornos e acessórios. Controla a
qualidade do processo e do produto confeccionado, verificando a
precisão, o tamanho e a abertura dos pontos, da trama e das costuras.
Observa os espaços entre os pontos. Inspeciona as formas, as cores e a
dimensão da peça. Verifica a união e o agrupamento das partes. Examina
imperfeições, efetuando ajustes ou reparos. Analisa a resistência do
produto confeccionado à umidade, à exposição solar e a demais efeitos
externos. Coloca etiquetas nos produtos confeccionados, preparando-os
para a comercialização. Pode usar logomarca, para identificação dos
produtos. Adota procedimentos para acondicionamento, limpeza e
conservação de peças artesanais de tricô, para evitar avarias e
deformação. Comercializa produtos artesanais de tricô, desenvolvendo
embalagens, calculando custo dos produtos e estabelecendo preços de
venda. Pesquisa mercado consumidor. Divulga produtos, podendo fazer uso
inclusive de ``sites'' e redes sociais. Participa de feiras e eventos,
expondo produtos. Aplica estratégias de venda; negocia valores e prazos
com cliente; e embala os produtos vendidos. Gerencia o próprio negócio,
organizando produção, controlando estoques e administrando orçamento.
Define o leiaute de oficina ou ateliê. Estabelece parcerias e seleciona
fornecedores. Pode capacitar artesãos. Pode vincular-se a uma associação
de classe. Realiza manutenção básica de equipamentos, fazendo sua
limpeza e lubrificação. Providencia serviços de manutenção corretiva.
Mantém a limpeza e a organização do local de trabalho. Organiza e
acondiciona materiais, ferramentas e instrumentos. Evita desperdícios,
separando sobras de materiais para reaproveitamento. Realiza descarte de
resíduos, de maneira ambientalmente adequada. Segue normas e
procedimentos de saúde e segurança no trabalho. Identifica e elimina
condições inseguras no local de trabalho, prevenindo acidentes. Observa
princípios básicos de ergonomia.

\begin{Form}
\textbf{Esta correspondência está adequada?} \\ 
\ChoiceMenu[radio,radiosymbol=\ding{52},bordercolor=0 0.5 0,name=cbovalid_ 791155 ]{}{CBO deve ficar=sim, \hspace{6em} CBO não fica aqui=nao}
\end{Form}

\textbf{CBO: 791160  ( Artesão rendeiro )}

Planeja a confecção artesanal de produtos de renda (como toalhas de
mesa, vestidos de noiva, colchas, entre outros) e de peças de renda a
serem aplicadas em produtos de tecido (como gola e punhos de blusa;
barra de toalha; entre outras), para atendimento às demandas de cliente.
Pesquisa tendências e características culturais. Seleciona equipamentos,
ferramentas, instrumentos e acessórios. Define o produto ou a peça de
renda, elaborando desenhos e detalhando especificações. Incorpora, à
renda a ser confeccionada, características artísticas pessoais. Pode
agregar, à renda, características típicas da cultura local ou regional,
que são transmitidas de geração em geração. Prepara a confecção manual
da peça ou do produto, selecionando os tipos de renda, tais como
renascença, irlandesa, nhanduti, frivoleté, labirinto, macramê, entre
outros. Escolhe fios - como de algodão, linho, seda, dentre outros -;
definindo sua espessura e suas cores. Seleciona o tipo e o tamanho da
agulha e define naveta (espécie de lançadeira) a ser utilizada. Risca
desenho ou diagrama em papel manteiga, papel cartão ou em outro
material, para orientar a confecção da renda. Confecciona a renda,
normalmente usando os processos de conduzir o fio com agulhas; trançar o
fio com bilros -- instrumentos de madeira ou metal -; e usar naveta ou
agulhas para fazer nós com fios. Segue padrões e desenhos para compor a
trama e a padronagem das rendas, tecendo, trançando, laçando e
entrelaçando fios; e desfiando tecido, a fim de separar os fios em
pequenos grupos e entrelaçá-los por nós. Executa acabamento, costurando
as diversas partes que compõem um produto de renda ou preparando peças
de renda para serem aplicadas como partes de um produto de tecido. Faz
arremates na renda. Pode aplicar adornos e acessórios. Pode engomar peça
de renda. Controla a qualidade do processo e do produto, verificando a
sinuosidade dos desenhos e a confecção da trama da renda. Inspeciona as
formas, as cores e a dimensão da peça ou do produto. Verifica a união e
o agrupamento das partes. Identifica imperfeições, efetuando ajustes ou
reparos. Analisa a resistência da peça ou do produto à umidade, à
exposição solar e a outros efeitos externos. Coloca etiquetas nas rendas
confeccionadas, preparando-as para a comercialização. Pode usar
logomarca, para identificação dos produtos. Adota procedimentos para
acondicionamento, limpeza e conservação de peças e produtos artesanais
de renda, para evitar avarias e deformação. Comercializa peças e
produtos artesanais de renda, desenvolvendo embalagens, calculando
custos e estabelecendo preços de venda. Pesquisa mercado consumidor.
Divulga peças e produtos, podendo fazer uso inclusive de ``sites'' e
redes sociais. Participa de feiras e eventos, expondo as rendas. Aplica
estratégias de venda; negocia valores e prazos com cliente; e embala
peças e produtos. Gerencia o próprio negócio, organizando produção,
controlando estoques e administrando orçamento. Define o leiaute de
oficina ou ateliê. Estabelece parcerias e seleciona fornecedores. Pode
capacitar artesãos. Pode vincular-se a uma associação de classe. Realiza
manutenção básica de equipamentos, fazendo sua limpeza e lubrificação.
Providencia serviços de manutenção corretiva. Mantém a limpeza e a
organização do local de trabalho. Organiza e acondiciona materiais,
ferramentas, instrumentos e acessórios. Evita desperdícios, separando
sobras de materiais para reaproveitamento. Realiza descarte de resíduos,
de maneira ambientalmente adequada. Segue normas e procedimentos de
saúde e segurança no trabalho. Identifica e elimina condições inseguras
no local de trabalho, prevenindo acidentes. Observa princípios básicos
de ergonomia.

\begin{Form}
\textbf{Esta correspondência está adequada?} \\ 
\ChoiceMenu[radio,radiosymbol=\ding{52},bordercolor=0 0.5 0,name=cbovalid_ 791160 ]{}{CBO deve ficar=sim, \hspace{6em} CBO não fica aqui=nao}
\end{Form}

\newpage
\section{ 09005 -- Técnico em Cenografia }

\textbf{Perfil do Curso (CNCT)}

O Técnico em Cenografia será habilitado para:

Criar e desenvolver projeto de ambientes cenográficos mediante
especificações em desenhos técnicos, croquis, plantas e maquetes.
Selecionar os materiais e equipamentos, de acordo com as especificações
técnicas do projeto cenográfico. Executar a produção. Supervisionar a
construção dos cenários.

Para atuação como Técnico em Cenografia, são fundamentais:

Conhecimentos interdisciplinares relacionados aos processos de criação,
envolvendo pesquisa, idealização, planejamento, execução técnica,
fruição e recepção estética. Competências comunicativas e empreendedoras
voltadas à proposição de projetos, ao coletivo, à gestão, à solução de
problemas e à resiliência, entre outras competências socioemocionais.

\textbf{CBOs correspondidos e seus perfis (presentes no CNCT, e presentes na predição da IA)}

\textbf{CBO: 374205  ( Cenotécnico (cinema, vídeo, televisão, teatro e espetáculos) )}

Planeja as atividades de construção de cenário, mobiliários e adereços,
analisando projeto elaborado por cenógrafo para produção cultural.
Estima tempo de execução e define etapas de trabalho. Dimensiona equipes
de trabalho, verificando a necessidade de contratar serviços de
terceiros. Seleciona ferramentas, máquinas, equipamentos e acessórios.
Faz orçamento do serviço. Submete plano de trabalho à aprovação do
diretor da produção cultural. Realiza adequação do projeto cenográfico
às medidas dos espaços cênicos. Representa graficamente projeto de
ambientes cenográficos. Pode utilizar ``software'' para elaboração de
cenário virtual. Executa pesquisa de vários materiais e objetos, para
escolher quais se adequam melhor ao projeto cenográfico. Faz a
especificação de materiais e objetos para aquisição. Supervisiona a
construção de cenário, de mobiliários e de adereços, orientando a
execução de trabalhos de carpintaria, serralheria, mecânica, modelagem,
pintura e vidraçaria. Coordena a produção de elementos cenográficos,
acompanhando os trabalhos de escultura, em fibras, de costura, de
textura e de cordame. Examina a qualidade do acabamento de cenário,
mobiliário e adereços. Cria e confecciona mecanismos de efeitos
especiais. Supervisiona acondicionamento e transporte de cenário, de
mobiliários e de adereços. Coordena a equipe de montagem do cenário,
considerando o projeto e o espaço cênico. Orienta os trabalhos de
posicionamento do cenário em espaço cênico; de adaptação do cenário ao
tamanho do espaço; de montagem de peças de cenário; e de fixação de
cenário móvel em maquinaria. Acompanha os serviços de decoração do
cenário; de afixação de adereços e objetos em cena; e de montagem dos
efeitos especiais. Compatibiliza cenário e adereços à luz. Confere
urdimento e condições físicas de palcos e ``set'' de filmagem. Organiza
trânsito de coxias. Roteiriza entrada e saída de cenário. Supervisiona a
operação de varas elétricas e guinchos automatizados ou robotizados.
Afina vestimentas de palcos. Pode acompanhar a execução do programa ou
evento cultural, para providenciar reparos de danos em cenário e
adereços; retoques em cenário; ou outras atividades relacionadas à
manutenção de peças cenográficas. Coordena a desmontagem do espaço
cenográfico, o transporte das peças e a organização das partes do
cenário para armazenagem. Elabora relatórios técnicos referentes à
adaptação e à execução de projeto cenográfico. Administra pessoal,
participando do processo de seleção, distribuindo tarefas e controlando
absenteísmo. Supervisiona o trabalho da equipe, avaliando desempenho e
levantando as necessidades de treinamento. Viabiliza a realização de
programas de treinamento. Orienta a equipe para conservar equipamentos,
ferramentas, instrumentos e acessórios limpos, acondicionados e em
plenas condições de uso e funcionamento. Requisita serviço de manutenção
corretiva de equipamentos, quando necessário. Zela pela segurança,
utilizando e monitorando o uso de equipamentos de proteção individual.
Identifica situações de risco e colabora para sua eliminação.

\begin{Form}
\textbf{Esta correspondência está adequada?} \\ 
\ChoiceMenu[radio,radiosymbol=\ding{52},bordercolor=0 0.5 0,name=cbovalid_ 374205 ]{}{CBO deve ficar=sim, \hspace{6em} CBO não fica aqui=nao}
\end{Form}

\textbf{CBO: 374210  ( Maquinista de cinema e vídeo )}

Planeja atividades, analisando o projeto de cenografia e examinando as
adaptações feitas por cenotécnico para adequação do projeto ao espaço
cênico. Estabelece as etapas do seu trabalho. Seleciona equipamentos
analógicos, automatizados e/ou robotizados. Confere materiais
necessários e requisita os itens faltantes. Pode orçar serviços. Atua na
preparação de equipamentos para filmagem cinematográfica e gravação de
vídeo, realizando instalação de suporte para deslocamento e movimentação
de câmeras; prestando apoio a eletricistas na montagem dos parques de
luz e em locação adaptada; operando varas elétricas (guinchos); armando
telas difusoras; e confeccionando mecanismos de efeitos especiais. Monta
cenário e suas peças, posicionando o cenário em espaço cênico; fixando
cenário móvel em maquinaria; e afixando adereços e objetos. Executa
trabalhos de cordame. Monta efeitos especiais. Compatibiliza cenário e
adereços à luz. Afina vestimentas de palcos. Confere urdimento e
condições físicas de ``set'' de filmagem. Examina o cenário montado,
verificando a disposição e o funcionamento de seus componentes, para
certificar-se de sua qualidade e sua segurança. Submete o cenário
finalizado à apreciação de cenotécnico e da direção da produção de filme
ou vídeo, para realizar eventuais ajustes ou melhorias. Durante ensaios,
gravação e filmagem, opera equipamentos - como gruas, carrinhos sobre
trilhos, entre outros - que deslocam e movimentam as câmeras. Manobra
rebatedores, opera cortinas, movimenta cenário sobre rodas, opera varas
cenográficas e opera efeitos especiais, de acordo com roteiro
predefinido de filmagem ou gravação. Pode auxiliar na elaboração de
relatórios técnicos referentes à montagem de cenário e às demais
atividades, relatando ocorrências durante os trabalhos. Recolhe
materiais - após a filmagem ou gravação de vídeo -, organizando-os para
guarda em depósito. Conserva equipamentos, ferramentas, instrumentos e
acessórios limpos, acondicionados e em plenas condições de uso e
funcionamento. Requisita serviço de manutenção corretiva de
equipamentos, quando necessário. Zela pela segurança, utilizando
equipamentos de proteção individual e atuando na prevenção de acidentes.
Identifica situações de risco e colabora para sua eliminação. Combate
princípios de incêndio. Pode prestar primeiros socorros.

\begin{Form}
\textbf{Esta correspondência está adequada?} \\ 
\ChoiceMenu[radio,radiosymbol=\ding{52},bordercolor=0 0.5 0,name=cbovalid_ 374210 ]{}{CBO deve ficar=sim, \hspace{6em} CBO não fica aqui=nao}
\end{Form}

\textbf{CBO: 374215  ( Maquinista de teatro e espetáculos )}

Planeja atividades, analisando o projeto de cenografia e examinando as
adaptações feitas por cenotécnico para adequação do projeto ao espaço
cênico. Estabelece as etapas do seu trabalho. Seleciona equipamentos
analógicos, automatizados e/ou robotizados. Confere materiais
necessários e requisita os itens faltantes. Pode orçar serviços. Atua na
operação de maquinaria para preparação de cenários de produções teatrais
e espetáculos, identificando cenários em relação a varas; fixando
manobras em cenários; e afinando vestimentas de palcos. Desce as varas
de luz, contrapesa refletores instalados e posiciona a vara de
iluminação. Confecciona mecanismos de efeitos especiais. Monta cenário e
suas peças, posiciona o cenário em espaço cênico e afixa adereços e
objetos. Fixa os cenários aéreos e móveis na maquinaria. Executa
trabalhos de cordame. Monta efeitos especiais. Compatibiliza cenário e
adereços à luz. Confere urdimento e observa condições físicas do palco.
Examina o cenário montado, verificando a disposição e o funcionamento de
seus componentes, para certificar-se de sua qualidade e sua segurança.
Submete o cenário finalizado à apreciação de cenotécnico e da direção da
produção de peça teatral ou espetáculo, para realizar eventuais ajustes
ou melhorias. Participa de ensaios de peça teatral ou espetáculo,
marcando altura do cenário aéreo; posicionando objetos de cena;
ajustando a sincronia de palco, luz e som; e verificando os movimentos
entre as coxias. Durante peça de teatro ou espetáculo, libera a sala
para o público e orienta técnicos de palco. Movimenta cenários aéreos e
térreos; cortinas de cena; contrapeso; cordas da maquinaria; alçapões; e
outros elementos cenográficos. Opera efeitos especiais. Pode operar
sistemas de varas, motorizadas, automatizadas ou robotizadas. Organiza
entradas e saídas de cenários e a movimentação por trás dos cenários.
Após término das apresentações de peça de teatro ou espetáculo, desmonta
cenário aéreo; retira adereços; suspende varas cênicas; e desmonta
cenário térreo (em palco) e piso. Armazena itens desmontados e
retirados. Conserva equipamentos, ferramentas, instrumentos e acessórios
limpos, acondicionados e em plenas condições de uso e funcionamento. Faz
a manutenção da maquinaria usada nos cenários. Pode auxiliar na
elaboração de relatórios técnicos referentes à montagem de cenário e às
demais atividades, relatando ocorrências durante os trabalhos. Zela pela
segurança, utilizando equipamentos de proteção individual e atuando na
prevenção de acidentes. Identifica situações de risco e colabora para
sua eliminação. Combate princípios de incêndio. Pode prestar primeiros
socorros.

\begin{Form}
\textbf{Esta correspondência está adequada?} \\ 
\ChoiceMenu[radio,radiosymbol=\ding{52},bordercolor=0 0.5 0,name=cbovalid_ 374215 ]{}{CBO deve ficar=sim, \hspace{6em} CBO não fica aqui=nao}
\end{Form}

\newpage
\section{ 09007 -- Técnico em Conservação e Restauro }

\textbf{Perfil do Curso (CNCT)}

O Técnico em Conservação e Restauro será habilitado para:

Realizar ações de conservação preventiva em acervos e bens culturais, de
acordo com a especificidade técnica: documentos, bens culturais e obras
de arte. Auxiliar e realizar processos de conservação e restauro de
peças do patrimônio histórico e cultural de acordo com cada
especificidade, com base nas diretrizes que orientam as teorias de
restauro vigentes. Auxiliar na análise e no diagnóstico de espaços
expositivos, acervos e reservas técnicas quanto à incidência de pragas,
fungos e outros fatores de risco, bem como orientar a adequada
iluminação, temperatura e umidade em acervos, espaços expositivos e
reservas técnicas. Auxiliar e realizar o processo de conservação e
restauro aplicando técnicas materiais e intervindo de acordo com a
necessidade em bens culturais, obras de arte, acervos e documentos de
acordo com a especificidade: documentos, obras em papel, mobiliário,
artes decorativas, escultura, pintura, gravura, têxteis.

Para atuação como Técnico em Conservação e Restauro, são fundamentais:

Conhecimentos interdisciplinares relacionados aos processos de criação,
envolvendo pesquisa, idealização, planejamento, execução técnica,
fruição e recepção estética. Competências comunicativas e empreendedoras
voltadas à proposição de projetos, ao coletivo, à gestão, à solução de
problemas e à resiliência, entre outras competências socioemocionais.

\textbf{CBOs correspondidos e seus perfis (presentes no CNCT, e presentes na predição da IA)}

\textbf{CBO: 768710  ( Restaurador de livros )}

Planeja o trabalho, analisando o material a ser restaurado e
identificando as fibras que formam o papel da obra a ser tratada.
Realiza o desmonte do documento e a numeração de cadernos. Faz
mapeamento dos danos nos cadernos. Pode representar graficamente os
danos do documento. Realiza teste de solubilidade no suporte (papel),
verificando se a tinta usada é solúvel em água. Após confirmar que não é
solúvel em água, executa banho de água deionizada - procedimento de
desacidificação -, para remover do papel as sujidades. Passa, ainda, o
material por banho com hidróxido de cálcio, para aumentar sua
alcalinidade. Caso o documento não possa passar por tratamento aquoso,
aplica outros métodos, como a ``limpeza'' com pó de borracha. Faz a
recuperação de papel deteriorado, calculando a gramatura das folhas;
executando reenfibragem, com uso da MOP-Máquina Obturadora de Papel; e
passando a área recuperada por processo de encolagem, com uso de cola
metilcelulose, para reforçar o material. Envolve o documento em feltro
de lã de carneiro ou outro processo similar, para tirar o excesso de
umidade e unir as fibras. Passa o documento pela prensa, pela secadora e
pelo refilamento. Faz a recuperação de lombadas, capas e guardas. Pode
aplicar outros métodos de restauração, a partir da análise inicial dos
danos mapeados em documento, pesquisando -- em Internet, instituição
especializada em restauração ou em outras fontes - casos similares.
Realiza controle de qualidade do trabalho executado, avaliando os
resultados da restauração. Pode realizar ou orientar processo de
encadernação e gravação em capas e lombadas, para restaurar sua
aparência original. Pode atuar em preservação e conservação de livros e
outros documentos, aplicando técnicas de higienização do acervo;
orientando o armazenamento adequado; desenvolvendo ações de
conscientização de usuários sobre o manuseio correto de obras; e
controlando fatores ambientais, tais como temperatura, umidade relativa
do ar, radiação da luz e presença de agentes biológicos -- como insetos,
roedores e fungos -- que levam à deterioração de acervos. Documenta as
ações realizadas, elaborando relatórios. Conserva o local de trabalho
limpo e organizado. Mantém instrumentos e acessórios limpos, organizados
e acondicionados. Monitora o funcionamento de máquinas e equipamentos,
requisitando serviço de manutenção corretiva, no caso de identificação
de falhas. Controla desperdício de material, reaproveitando sobras e
fazendo descarte de resíduos de acordo com normas ambientais. Trabalha
com segurança, prevenindo acidentes, mantendo postura ergonomicamente
correta e utilizando equipamentos de proteção individual.

\begin{Form}
\textbf{Esta correspondência está adequada?} \\ 
\ChoiceMenu[radio,radiosymbol=\ding{52},bordercolor=0 0.5 0,name=cbovalid_ 768710 ]{}{CBO deve ficar=sim, \hspace{6em} CBO não fica aqui=nao}
\end{Form}

\newpage
\section{ 09008 -- Técnico em Dança }

\textbf{Perfil do Curso (CNCT)}

O Técnico em Dança será habilitado para:

Criar e interpretar coreografias e performances diversas. Desenvolver
práticas corporais, técnicas de criação e formação em dança a partir de
matrizes plurais de conhecimento em dança, corpo e cultura. Realizar
investigações de dança em interface com outras linguagens artísticas, em
conexão com saberes tradicionais, populares, urbanos e digitais.\\
Elaborar e executar ações relacionadas à arte e à dança em projetos
socioculturais.

Para atuação como Técnico em Dança, são fundamentais:

Conhecimentos interdisciplinares relacionados aos processos de criação,
envolvendo pesquisa, idealização, planejamento, execução técnica,
fruição e recepção estética. Competências comunicativas e empreendedoras
voltadas à proposição de projetos, ao coletivo, à gestão, à solução de
problemas e à resiliência, entre outras competências socioemocionais.

\textbf{CBOs correspondidos e seus perfis (presentes no CNCT, e presentes na predição da IA)}

\textbf{CBO: 376105  ( Dançarino tradicional )}

Pesquisa danças tradicionais ou folclóricas como manifestações culturais
- produzidas espontaneamente em diferentes regiões, criadas em rituais
mágicos e religiosos, surgidas em brincadeiras ou jogos infantis e
apresentadas em festejos e em outros eventos - transmitidas de geração
em geração. Vivencia a dança nesses ambientes onde foram criadas,
perpetuadas ou recuperadas. Estuda as danças, analisando seus passos,
movimentos e gestos; identificando sua pluralidade rítmica; apreciando
sua música; e examinando seus figurinos. Realiza ensaios, preparando o
corpo para dançar ao toque dos tambores e ao som de loas, saudações,
ladainhas e cantos. Avalia o nível de dificuldade dos passos, movimentos
e gestos. Introduz pequenas modificações na dança, experimentando novos
movimentos; explorando possibilidades de movimentos com objetos; e
aprimorando a coreografia. Atua nas músicas dançadas, fazendo mixagem e
trabalhando o tempo ritual. Dirige ensaios, coordenando movimentos e
expressão corporal. Interpreta as danças dramáticas, as danças rituais e
as danças como representações de histórias, respeitando as tradições.
Utiliza instrumentos musicais, figurinos, adereços e objetos adequados e
próprios à dança. Executa a dança conforme coreografia e ritmos
tradicionais. Dá aulas relacionadas às danças tradicionais, adequando
métodos de ensino às diferentes faixas etárias dos alunos e transmitindo
conhecimentos sobre os mestres das tradições e sobre a história das
danças. Ministra aulas para profissionalizar dançarinos e professores,
realizando oficinas, seminários, apresentações, exposições temáticas e
debates. Promove cursos para educadores, terapeutas e profissionais de
outras áreas. Insere o acervo cultural da dança tradicional em
diferentes contextos. Aplica a dança tradicional em projetos sociais,
empresariais e institucionais; em terapias; em contextos pedagógicos; e
em ações culturais. Administra atividades relacionadas à dança.
Supervisiona equipe de trabalho, atribuindo funções a cada um de seus
membros, avaliando seu desempenho e promovendo seu desenvolvimento
profissional. Gerencia a produção artística e a produção executiva.
Realiza eventos e consulta leis de incentivo à cultura para apresentar
projetos, com a finalidade de captar recursos. Registra dados coletados
e sistematizados sobre danças tradicionais ou folclóricas. Organiza e
arquiva documentos físicos e eletrônicos, vídeos e outros registros
gravados.

\begin{Form}
\textbf{Esta correspondência está adequada?} \\ 
\ChoiceMenu[radio,radiosymbol=\ding{52},bordercolor=0 0.5 0,name=cbovalid_ 376105 ]{}{CBO deve ficar=sim, \hspace{6em} CBO não fica aqui=nao}
\end{Form}

\textbf{CBO: 376110  ( Dançarino popular )}

Pesquisa danças populares como manifestações culturais que surgem em
diferentes regiões ou segmentos sociais, em festejos e em ambientes
específicos de dança. Examina passos, gestos, movimentos e figurinos.
Pesquisa movimentos para o aparecimento de novos ritmos. Vivencia a
dança. Estuda as danças populares, analisando seus passos, movimentos e
gestos; identificando sua pluralidade rítmica; apreciando sua música; e
examinando seus figurinos. Realiza ensaios, preparando o corpo através
de alongamento, aquecimento e outras técnicas específicas. Avalia o
nível de dificuldade dos passos, movimentos e gestos. Introduz pequenas
modificações na dança, experimentando novos movimentos; explorando
possibilidades de movimentos com objetos; e aprimorando a coreografia.
Atua nas músicas, fazendo mixagem. Faz a adequação da dança aos novos
ritmos e trabalha o tempo cênico. Dirige ensaios, coordenando movimentos
e expressão corporal. Interpreta as danças com objetivos cênicos,
reinterpreta as danças populares e interpreta as danças como
representações de histórias. Utiliza instrumentos musicais, figurinos,
adereços e objetos adequados e próprios à dança. Executa a dança
conforme coreografia e ritmos populares. Cria espetáculos, concebendo
coreografias a partir do desenvolvimento de objetos cênicos, cenários,
figurinos e trilhas sonoras. Elabora roteiros. Faz a fusão da dança com
outras linguagens artísticas e tecnológicas. Dá aulas relacionadas às
danças populares, desenvolvendo propostas pedagógicas e adequando
métodos de ensino às diferentes faixas etárias dos alunos. Transmite
conhecimentos sobre os mestres e sobre a história das danças populares.
Ministra aulas para profissionalizar dançarinos e professores,
realizando oficinas, seminários, apresentações, exposições temáticas e
debates. Promove cursos para educadores, terapeutas e profissionais de
outras áreas. Insere o acervo cultural da dança popular em diferentes
contextos. Aplica a dança popular em projetos sociais, empresariais e
institucionais; em terapias; em contextos pedagógicos; e em ações
culturais. Administra atividades relacionadas à dança. Supervisiona
equipe de trabalho, atribuindo funções a cada um de seus membros,
avaliando seu desempenho e promovendo seu desenvolvimento profissional.
Gerencia a produção artística e a produção executiva. Realiza eventos e
consulta leis de incentivo à cultura para apresentar projetos, com a
finalidade de captar recursos. Registra dados coletados e sistematizados
sobre danças populares. Organiza e arquiva documentos físicos e
eletrônicos, vídeos e outros registros gravados.

\begin{Form}
\textbf{Esta correspondência está adequada?} \\ 
\ChoiceMenu[radio,radiosymbol=\ding{52},bordercolor=0 0.5 0,name=cbovalid_ 376110 ]{}{CBO deve ficar=sim, \hspace{6em} CBO não fica aqui=nao}
\end{Form}

\newpage
\section{ 09009 -- Técnico em Design de Calçados }

\textbf{Perfil do Curso (CNCT)}

O Técnico em Design de Calçados será habilitado para:

Realizar pesquisa de tendências de moda, comportamento e mercado. Propor
materiais e componentes. Coordenar projetos de inovação estética,
funcional e tecnológica. Executar peças-piloto e processos de
fabricação.

Para atuação como Técnico em Design de Calçados, são fundamentais:

Conhecimentos interdisciplinares relacionados aos processos de criação,
envolvendo pesquisa, idealização, planejamento, execução técnica,
fruição e recepção estética. Competências comunicativas e empreendedoras
voltadas à proposição de projetos, ao coletivo, à gestão, à solução de
problemas e à resiliência, entre outras competências socioemocionais.

\textbf{CBOs correspondidos e seus perfis (presentes no CNCT, e presentes na predição da IA)}

\textbf{CBO: 318815  ( Modelista de calçados )}

Planeja a confecção de moldes e o desenvolvimento de protótipos de
calçados, examinando projeto do designer; verificando tendências de
mercado; consultando boletins técnicos; avaliando estilos de design,
concepções e produtos de concorrentes; e coletando informações
tecnológicas em feiras de calçados. Confecciona moldes de calçados,
interpretando modelos e desenhos; elaborando desenhos planificados;
traçando, cortando e codificando moldes. Utiliza prancheta e softwares,
tais como CAD-Desenho Assistido por Computador (Computer Aided Design) e
CAM--Manufatura Auxiliada por Computador (Computer Aided Manufacturing).
Marca referências em moldes, com a finalidade de orientar a produção de
calçados. Desenvolve protótipos de calçados, interpretando normas
técnicas para produção; elaborando ficha técnica; especificando
aviamentos e acessórios; ajustando moldes e gabaritos; e determinando a
posição de ornamentos, detalhes e acessórios. Produz e testa protótipo
de calçado. Requisita testes laboratoriais do protótipo. Realiza
ampliações e reduções das dimensões do produto por escala. Avalia
materiais para aquisição com vistas à produção de calçados ergonômicos,
verificando a resistência dos materiais e avaliando sua composição e
acabamento. Avalia a costurabilidade e testa a fixação de cores em
couros. Interpreta resultados de testes laboratoriais, elaborando
relatórios de testagem de materiais. Avalia a relação entre custo e
benefício.

\begin{Form}
\textbf{Esta correspondência está adequada?} \\ 
\ChoiceMenu[radio,radiosymbol=\ding{52},bordercolor=0 0.5 0,name=cbovalid_ 318815 ]{}{CBO deve ficar=sim, \hspace{6em} CBO não fica aqui=nao}
\end{Form}

\newpage
\section{ 09011 -- Técnico em Design de Interiores }

\textbf{Perfil do Curso (CNCT)}

O Técnico em Design de Interiores será habilitado para:

Criar, desenvolver e viabilizar a execução de projetos de interiores
residenciais, comerciais, de vitrines e exposições. Orientar e
desenvolver projetos com base em ergonomia e desenho universal,
conforto, saúde e bem-estar. Desenvolver esboços, perspectivas e
desenhos técnicos. Planejar e organizar o espaço, com base nos estudos
ergonômicos, estéticos e funcionais. Identificar elementos básicos para
a concepção projetual dos espaços internos habitados. Representar os
elementos de projeto no espaço bi e tridimensional. Aplicar métodos de
representação gráfica. Prestar assistência técnica no estudo e
desenvolvimento de projetos e pesquisas tecnológicas, voltadas às
atividades da área. Orientar e coordenar a execução dos serviços de
manutenção de equipamentos e instalações de ambientes e mobiliários
fixos. Reformar ambientes sem alteração estrutural.

Para atuação como Técnico em Design de Interiores, são fundamentais:

Conhecimentos interdisciplinares relacionados aos processos de criação,
envolvendo pesquisa, idealização, planejamento, execução técnica,
fruição e recepção estética. Competências comunicativas e empreendedoras
voltadas à proposição de projetos, ao coletivo, à gestão, à solução de
problemas e à resiliência, entre outras competências socioemocionais.

\textbf{CBOs correspondidos e seus perfis (presentes no CNCT, e presentes na predição da IA)}

\textbf{CBO: 318015  ( Desenhista detalhista )}

Analisa solicitação de detalhamento de desenho técnico, interpretando
original e documentos de apoio, tais como plantas, projetos, catálogos,
croquis e normas. Propõe ao solicitante alternativas para execução do
detalhamento, chegando a um acordo sobre pormenores técnicos finais.
Estima tempo para realização do detalhamento e estabelece prioridades,
conforme cronograma. Programa ações para realização do desenho
detalhado, observando características técnicas, fazendo esboços,
estabelecendo formatos e definindo escalas e sistemas de representação.
Elabora desenhos técnicos detalhados, representando formas geométricas
em duas e três dimensões. Utiliza computador, programa -- tal como
CAD-Desenho Assistido por Computador (Computer Aided Design) -,
impressoras, plotters e outros periféricos, além de fazer uso de
instrumentos para execução do traçado diretamente em substratos de
papel. Constrói o detalhamento do desenho, traçando linhas auxiliares,
cotando e realçando aspectos e características de formas geométricas
complexas e outros elementos a detalhar. Elabora a lista de materiais e
componentes do desenho detalhado. Envia o desenho técnico detalhado para
revisão, recebendo a aprovação do solicitante. Realiza cópias de
segurança e disponibiliza o desenho detalhado final e/ou suas revisões
para as áreas afins, que providenciarão a entrega ao solicitante.
Conserva o local de trabalho limpo e organizado. Controla graus de
luminosidade do ambiente e verifica as condições ergonômicas para uso de
monitor de vídeo, mouse, mesa digitalizadora, prancheta e teclado.
Mantém equipamentos e instrumentos de trabalho em plenas condições de
uso e funcionamento. Requisita serviço de manutenção de computador,
impressoras, plotters e outros periféricos, quando necessário. Organiza
e arquiva documentos físicos e eletrônicos, tais como comprovantes de
aprovação e cópias de segurança dos desenhos técnicos detalhados.
Destina adequadamente os resíduos gerados, fazendo reciclagem e
realizando descarte de acordo com as normas ambientais.

\begin{Form}
\textbf{Esta correspondência está adequada?} \\ 
\ChoiceMenu[radio,radiosymbol=\ding{52},bordercolor=0 0.5 0,name=cbovalid_ 318015 ]{}{CBO deve ficar=sim, \hspace{6em} CBO não fica aqui=nao}
\end{Form}

\textbf{CBO: 375105  ( Designer de interiores )}

Entrevista cliente, identificando suas expectativas e necessidades em
relação ao projeto e definindo seu limite orçamentário. Verifica os usos
e as funções dos espaços a serem trabalhados. Estabelece o contrato de
trabalho. Realiza levantamento métrico e fotográfico do espaço, para
identificar elementos básicos para a concepção do projeto. Desenvolve
esboços, perspectivas e desenhos. Concebe o design de interiores - sem
interferência na estrutura da edificação -, elaborando planta de
distribuição dos espaços internos, considerando a circulação e
distribuindo os volumes no espaço. Planeja e organiza o uso e a ocupação
dos espaços, otimizando o conforto, a estética, a saúde e a segurança do
usuário e atuando de acordo com as normas técnicas de acessibilidade, de
ergonomia e de conforto luminoso, térmico e acústico. Concebe
programação visual para espaços comerciais. Apresenta o estudo
preliminar ao cliente, para aprovação. Elabora o projeto executivo,
especificando os materiais a serem utilizados, escolhendo escala
cromática para o ambiente e alocando pontos de iluminação, pontos de
ar-condicionado, pontos elétricos, pontos hidráulicos, pontos para
informática e pontos de telefonia. Representa graficamente as soluções
para o ambiente, utilizando desenho manual ou recursos tecnológicos,
como CAD-Projeto Assistido por Computador (Computer Aided Design). Pode
fazer projetos em 2D e 3D, maquetes eletrônicas e exposições por meio de
passeios virtuais. Apresenta a proposta e o orçamento do projeto ao
cliente. Após autorização da execução do projeto pelo cliente, elabora o
cronograma, seleciona fornecedores e contrata serviços especializados.
Coordena as equipes para execução de instalações, montagens, reparos,
pintura e outros serviços, em reformas de ambientes. Acompanha os
trabalhos, avalia os resultados e providencia os ajustes necessários.
Pode orientar a realização dos serviços de manutenção de equipamentos,
instalações de ambientes e mobiliários fixos. Presta assistência técnica
ao cliente na compra de mobiliários, acessórios, materiais e
revestimentos, compatibilizando necessidades e expectativas do cliente
com as suas condições socioeconômicas e culturais. Pode colaborar na
criação de peças especiais, móveis e produtos. Realiza a entrega dos
serviços finalizados ao cliente, orientando-o na melhor forma de
utilização de mobiliários e acessórios. Pode organizar portfólio com os
resultados de seus trabalhos. Pode coletar, tratar e analisar dados de
natureza técnica, para elaboração de laudos ou relatórios técnicos.
Realiza pesquisas nos mercados nacional e internacional e participa de
exposições, mostras e feiras, para conhecer inovações em design de
interiores. Compatibiliza os seus projetos com as exigências legais
relacionadas à segurança contra incêndio, à saúde e ao meio ambiente.
Observa e orienta prestadores de serviços na prevenção de acidentes, no
manuseio correto de cargas e na utilização de equipamentos de proteção
individual.

\begin{Form}
\textbf{Esta correspondência está adequada?} \\ 
\ChoiceMenu[radio,radiosymbol=\ding{52},bordercolor=0 0.5 0,name=cbovalid_ 375105 ]{}{CBO deve ficar=sim, \hspace{6em} CBO não fica aqui=nao}
\end{Form}

\textbf{CBO: 375110  ( Designer de vitrines )}

Realiza estudo preliminar, levantando as expectativas do cliente, os
limites orçamentários e o prazo para execução do projeto. Analisa o
perfil do público consumidor e as características do produto a ser
exposto. Realiza levantamento fotográfico e calcula as medidas do espaço
disponível, para identificar os elementos básicos para a concepção do
projeto da vitrine. Analisa a possibilidade de aproveitamento das
instalações pré-existentes. Concebe o projeto da vitrine, com base no
limite de custo e no prazo definidos. Desenha o croqui da vitrine,
utilizando técnicas de desenho de observação e perspectiva. Apresenta a
versão preliminar do projeto de vitrine ao cliente, podendo realizar
ajustes por ele solicitados. Elabora o projeto executivo, selecionando
elementos decorativos, cenográficos, mobiliários, suportes e
expositores. Escolhe a escala cromática. Aloca pontos elétricos e pontos
de iluminação. Representa graficamente as soluções para a vitrine,
utilizando desenho manual ou recursos tecnológicos, como CAD-Projeto
Assistido por Computador (Computer Aided Design). Faz o orçamento para
execução do projeto, a ser submetido à aprovação do cliente. Após
aprovação, consulta carteira de fornecedores, para aquisição de produtos
e serviços. Realiza a montagem da vitrine, de acordo com os princípios e
as técnicas de composição e exposição de produtos. Utiliza recursos de
iluminação, considerando os efeitos de destaque da decoração, dos
materiais promocionais e dos produtos expostos na vitrine. Avalia os
resultados do projeto, analisando atratividade da vitrine junto ao
público, variação de vendas do produto exposto e satisfação do cliente.
Realiza pesquisas nos mercados nacional e internacional e participa de
exposições, mostras e feiras, para conhecer inovações em vitrinismo.
Zela pela segurança, prevenindo acidentes e utilizando equipamentos de
proteção individual.

\begin{Form}
\textbf{Esta correspondência está adequada?} \\ 
\ChoiceMenu[radio,radiosymbol=\ding{52},bordercolor=0 0.5 0,name=cbovalid_ 375110 ]{}{CBO deve ficar=sim, \hspace{6em} CBO não fica aqui=nao}
\end{Form}

\textbf{CBO: 375115  ( Visual merchandiser )}

Entrevista o responsável pelo ponto de venda para identificar suas
expectativas, diagnosticar os problemas enfrentados pelo estabelecimento
ou pela marca, estabelecer limites orçamentários do projeto e estimar
prazo para execução do projeto. Avalia o espaço físico e os aspectos
impactantes de sua ocupação, como os arquitetônicos, ergonômicos e de
segurança. Analisa as características e o comportamento do consumidor,
verifica o tipo de produto em venda e pesquisa os fatores que
influenciam na decisão de compra do consumidor. Elabora projeto para o
ponto de venda, levando em consideração as estratégias de marketing e
concebendo a programação visual para espaços. Projeta estratégia
criativa de merchandising visual, definindo a decoração de interiores e
selecionando as cores a serem usadas no espaço; criando dispositivos
para a apresentação de informações (displays) com visual atraente e
engajante; combinando arte e criatividade na produção de vitrines, para
valorização dos produtos expostos. Pode sugerir eventuais modificações
no projeto arquitetônico. Projeta a organização do espaço de modo a
facilitar a circulação dos clientes. Planeja climatização agradável e
boa iluminação. Elabora proposta de estímulos auditivos -- como músicas
adequadas ao perfil dos clientes -- e oferta de petiscos e bebidas --
como água e café --, para tornar a permanência no espaço uma experiência
agradável. Estabelece ações estratégicas para atrair a atenção do
cliente para a política de descontos. Prevê dispositivos e sistemas -
com base em nuvem e novas tecnologias -, para dar maior segurança aos
clientes e ao espaço. Apresenta a proposta e o orçamento ao cliente,
para aprovação. Interage com fornecedores, para aquisição de produtos e
serviços para o projeto. Executa o projeto, coordenando a implementação
das mudanças propostas. Promove programa de treinamento para a equipe de
vendas, para oferecer ao cliente o melhor atendimento possível. Avalia
os resultados da implementação do projeto, verificando a fidelização dos
clientes, analisando o volume de vendas e levantando a satisfação dos
clientes. Pode fazer ajustes no projeto inicial, a partir da avaliação
dos resultados. Faz manutenção programada do projeto. Participa de
palestras e cursos, conhecendo propostas inovadoras em visual
merchandising, para manter processo de melhoria contínua dos resultados
no ponto de vendas.

\begin{Form}
\textbf{Esta correspondência está adequada?} \\ 
\ChoiceMenu[radio,radiosymbol=\ding{52},bordercolor=0 0.5 0,name=cbovalid_ 375115 ]{}{CBO deve ficar=sim, \hspace{6em} CBO não fica aqui=nao}
\end{Form}

\textbf{CBO: 375120  ( Decorador de eventos )}

Entrevista o cliente para identificar suas expectativas e preferências,
verificar os limites orçamentários do projeto, definir a data do evento
e estimar a quantidade de convidados. Apresenta seu portfólio de
trabalhos ao cliente, com opções de estilo de decoração. Elabora estudos
preliminares, fazendo levantamento fotográfico e das medidas no local do
evento e verificando a viabilidade de adequar elementos já existentes ao
espaço. Elabora planta de distribuição dos espaços internos e planeja a
circulação dos convidados. Concebe a programação visual para o espaço do
evento. Planeja a distribuição do mobiliário. Propõe os tipos de
arranjos florais. Interage com outros profissionais contratados para o
evento -- tais como cenógrafo, apresentador e organizador de evento -,
para compatibilização das propostas de trabalho. Apresenta o estudo
preliminar ao cliente, realizando os ajustes por ele solicitados.
Elabora o projeto executivo, especificando os materiais a serem
utilizados e escolhendo a escala cromática. Define a técnica de
iluminação. Relaciona os móveis a serem utilizados e sua localização no
ambiente, levando em consideração a praticidade e o conforto dos
convidados. Representa graficamente o projeto de decoração, utilizando
desenho manual ou recursos tecnológicos, como CAD-Projeto Assistido por
Computador (Computer Aided Design). Apresenta o projeto executivo, o
cronograma e o orçamento ao cliente, para aprovação. Seleciona
fornecedores, para aquisição de produtos e serviços necessários à
execução do projeto. Coordena a montagem do evento, de acordo com o
projeto aprovado. Providencia atrativos sensoriais no ambiente. Cumpre o
cronograma, acompanhando o trabalho da equipe e controlando a entrega
dos fornecedores. Cria soluções para imprevistos durante a montagem, com
base em conhecimentos sobre materiais e técnicas de decoração. Acompanha
a execução orçamentária. Após o evento, acompanha a desmontagem da
decoração, preservando o local como antes do início do projeto. Avalia
os resultados do projeto, analisando a reação dos convidados durante o
evento e verificando a satisfação do cliente. Realiza pesquisas nos
mercados nacional e internacional e participa de exposições, mostras e
feiras, para conhecer inovações em relação à decoração de eventos. Zela
pela segurança, prevenindo acidentes durante a execução do evento e nos
processos de montagem e desmontagem da decoração.

\begin{Form}
\textbf{Esta correspondência está adequada?} \\ 
\ChoiceMenu[radio,radiosymbol=\ding{52},bordercolor=0 0.5 0,name=cbovalid_ 375120 ]{}{CBO deve ficar=sim, \hspace{6em} CBO não fica aqui=nao}
\end{Form}

\newpage
\section{ 09014 -- Técnico em Design de Móveis }

\textbf{Perfil do Curso (CNCT)}

O Técnico em Design de Móveis será habilitado para:

Desenvolver esboços, perspectivas e desenhos normatizados de móveis.
Realizar estudos volumétricos e maquetes convencionais e eletrônicas.
Aplicar aspectos ergonômicos ao projeto. Pesquisar e definir materiais,
ferragens e acessórios. Elaborar documentação técnica normatizada.
Acompanhar a execução de protótipos ou peças-piloto. Aplicar os
conceitos de sustentabilidade ao desenvolvimento de móveis.

Para atuação como Técnico em Design de Móveis, são fundamentais:

Conhecimentos interdisciplinares relacionados aos processos de criação,
envolvendo pesquisa, idealização, planejamento, execução técnica,
fruição e recepção estética. Competências comunicativas e empreendedoras
voltadas à proposição de projetos, ao coletivo, à gestão, à solução de
problemas e à resiliência, entre outras competências socioemocionais.

\textbf{CBOs correspondidos e seus perfis (presentes no CNCT, e presentes na predição da IA)}

\textbf{CBO: 318425  ( Desenhista técnico (mobiliário) )}

Prepara as atividades, organizando o posto de trabalho de desenho,
inspecionando equipamentos e selecionando ferramentas, instrumentos e
materiais. Interpreta solicitação de desenho, reunindo e organizando
informações pertinentes. Consulta revistas, catálogos e outras
publicações técnicas. Desenvolve esboços manualmente e com auxílio de
recursos computacionais. Define os meios de representação gráfica,
empregando recursos analógicos e/ou digitais adequados à demanda.
Planeja a elaboração de desenho de móveis, ordenando o trabalho por
etapas e fixando prazo de entrega. Define escalas e estabelece o formato
de apresentação, impresso e/ou eletrônico. Consulta normas técnicas para
especificar as características do desenho. Desenvolve o desenho,
combinando elementos de imagens, conforme a necessidade, para melhor
orientar a produção do mobiliário. Realiza representação gráfica pelo
método manual, fazendo uso de instrumentos para execução do traçado
diretamente em substratos de papel, bem como utilizando computadores,
programas integrados de CAD-Desenho Assistido por Computador (Computer
Aided Design), ilustração e tratamento de imagens. Codifica o desenho e
relaciona as especificações técnicas envolvidas em sua produção. Confere
as especificações em face do trabalho realizado e examina o atendimento
às normas técnicas de representação gráfica, para imprimir e requisitar
a aprovação do desenho. Realiza as correções indicadas pelo solicitante,
registra e arquiva o desenho aprovado. Executa o acabamento final do
desenho, indicando as características de materiais de acabamento. Obtém
a aprovação final. Confecciona a matriz do desenho, preenche a legenda
do desenho e tira as cópias de segurança (backup). Conserva o local de
trabalho limpo e organizado. Controla graus de luminosidade do ambiente,
temperatura e ruídos. Verifica as condições ergonômicas para uso de
mobiliário, monitor de vídeo de computador, mouse e teclado. Requisita a
aquisição de materiais e instrumentos de trabalho, para reposição.
Conserva equipamentos, ferramentas e instrumentos de trabalho em plenas
condições de uso e funcionamento. Requisita serviço para manutenção de
computador, impressoras e outros periféricos, quando necessário.
Organiza e arquiva documentos físicos e eletrônicos, tais como
comprovantes de aprovação e cópias de segurança dos desenhos aprovados.
Destina adequadamente os resíduos gerados, fazendo reciclagem e
realizando descarte de acordo com as normas ambientais.

\begin{Form}
\textbf{Esta correspondência está adequada?} \\ 
\ChoiceMenu[radio,radiosymbol=\ding{52},bordercolor=0 0.5 0,name=cbovalid_ 318425 ]{}{CBO deve ficar=sim, \hspace{6em} CBO não fica aqui=nao}
\end{Form}

\textbf{CBO: 318805  ( Projetista de móveis )}

Planeja projetos e protótipos de móveis para produção sob encomenda ou
em série, com base em tendências de mercado; consultando área de
produção; consultando boletins técnicos; avaliando estilos de design,
concepções e produtos de concorrentes; e coletando informações
tecnológicas em feiras do mobiliário. Desenvolve projetos de móveis,
analisando locais de instalação de móveis sob medida e interpretando
desenhos e modelos. Elabora desenhos e gabaritos para a produção de
móveis, utilizando prancheta e software -- tal como o CAD-Desenho
Assistido por Computador (Computer Aided Design) -; dimensionando
componentes; especificando madeiras, derivados e acessórios; e
especificando materiais para acabamento, como tintas e vernizes. Pode
utilizar software de CAM--Manufatura Auxiliada por Computador (Computer
Aided Manufacturing), gerando programas de CNC--Comando Numérico
Computadorizado, para fabricação de mobiliário e confecção de maquetes.
Pode gerar maquetes eletrônicas de móveis, em três dimensões, a partir
de desenhos de projetos em software, tal como CAD-Desenho Assistido por
Computador (Computer Aided Design). Desenvolve protótipos de móveis,
interpretando normas técnicas para produção; elaborando ficha técnica;
especificando ferragens e acessórios; e determinando a posição de
ornamentos, detalhes e acessórios. Orienta a confecção de protótipos e
peças-piloto de móveis, realizando e requisitando testes de
laboratórios; ajustando gabaritos; e realizando ampliações e reduções
das dimensões do móvel por escala. Testa características do protótipo do
móvel, como as qualidades de segurança e manuseio, a eficiência com que
ele pode ser produzido e as maneiras de usá-lo e mantê-lo. Avalia
materiais para aquisição com vistas à produção de mobiliário,
verificando a resistência dos materiais; examinando sua adequação para
produzir móveis ergonômicos e funcionais; e avaliando sua composição e
acabamento. Interpreta resultados de testes laboratoriais, elaborando
relatórios. Avalia a relação entre custo e benefício de materiais.
Refina projetos, usando modelos de trabalho em conformidade com as
especificações do cliente, limitações de produção ou alterações de
tendências de estilo. Avalia a viabilidade do projeto, com base em
fatores como estilo, segurança, função, características do mercado,
facilidade de manutenção, orçamento e custos da produção. Pode
supervisionar e orientar equipe de trabalho na realização de
prototipagem e testes, distribuindo tarefas, coordenando e avaliando a
execução das atividades.

\begin{Form}
\textbf{Esta correspondência está adequada?} \\ 
\ChoiceMenu[radio,radiosymbol=\ding{52},bordercolor=0 0.5 0,name=cbovalid_ 318805 ]{}{CBO deve ficar=sim, \hspace{6em} CBO não fica aqui=nao}
\end{Form}

\textbf{CBO: 771105  ( Marceneiro )}

Planeja o trabalho, interpretando projetos, desenhos e especificações,
para a confecção e para a restauração de móveis e artefatos de madeira e
derivados. Seleciona matéria-prima, insumos e componentes, máquinas,
equipamentos, ferramentas manuais, ferramentas elétricas portáteis e
instrumentos de medição e controle. Elabora croqui ou esboço e realiza
cálculos para dimensionamento do produto, conforme solicitação do
cliente e local de instalação. Pode utilizar software de apoio para
desenhar o projeto completo em computador, especificando madeiras,
acessórios, ferragens e acabamentos. Confecciona gabaritos ou moldes
para execução das peças de madeira e derivados. Define roteiro para as
atividades de confecção ou restauração. Pode elaborar orçamento. Prepara
máquinas e ferramentas, afiando ferramentas e ajustando parâmetros, tais
como de posicionamento e geometria de corte, velocidade e pressão da
ferramenta, entre outros. Confecciona móveis e artefatos de madeira e
derivados, regulando as máquinas para obter o produto conforme o
projeto. Faz a traçagem na madeira, derivados e outros materiais,
observando o sentido dos veios e definindo o posicionamento e a
geometria do corte. Usina as partes, definindo a sequência de operações
a serem aplicadas para cortar, tornear, serrar, furar, aplainar, lixar,
fresar e laminar. Pode programar e utilizar máquinas-ferramenta com
comando numérico computadorizado (CNC) e fazer uso de programas para
usinagem das peças integrados a ferramentas de software de CAD/CAM.
Monta móveis e artefatos de madeira e derivados, definindo a sequência
de montagem. Realiza operações de união de peças para prensar, unir com
cavilha e/ou cola, unir com dispositivos, entre outras. Aplica cola -
com coladeira de borda ou sob pressão - ou massa. Utiliza elementos de
fixação, acessórios, ferragens e dispositivos de montagem. Realiza
operações de acabamento em móveis e artefatos de madeira e derivados,
aplicando tintas e vernizes, revestindo e polindo, apondo apliques e
lâminas, colocando ferragens para reajuste e regulando o funcionamento
das partes móveis, de acordo com o projeto. Pode realizar teste de
viscosidade nos produtos utilizados para acabamento da superfície.
Restaura móveis e artefatos de madeira e derivados, avaliando produtos
para determinar a viabilidade de restauração, confeccionando peças para
reposição, substituindo peças danificadas, reapertando elementos de
fixação e preparando o produto para acabamento. Verifica peças e
produtos, para garantir a qualidade dos móveis e artefatos
confeccionados. Avalia a matéria-prima, a resistência, as formas e
dimensões, os encaixes, a união e a fixação das partes, as condições do
acabamento da superfície e da aplicação de ferragens, acessórios e
dispositivos. Compara as características do produto final com os
requisitos do cliente ou do projeto, providenciando ajustes, quando
necessários. Pode desmontar o produto confeccionado, para posterior
entrega. Prepara a entrega dos móveis e dos artefatos confeccionados sob
medida ou restaurados, embalando separadamente os acessórios e os
elementos principais, definindo as ferramentas a serem utilizadas para
fazer o transporte e a montagem do produto no local da instalação e
apondo os arremates finais. Organiza o local de trabalho. Prepara os
locais para carga, descarga e armazenamento de materiais. Realiza
limpeza e conservação dos recursos materiais empregados. Mantém os
materiais, instrumentos e acessórios de trabalho organizados,
acondicionados e em plenas condições de uso e funcionamento. Controla
desperdícios de material. Separa, identifica e classifica resíduos --
sólidos e líquidos -, para descarte, reúso e reciclagem. Realiza as
atividades utilizando equipamentos de proteção individual e coletiva.
Preenche fichas de controle de processo, ficha técnica do produto,
requisição de materiais, entre outros documentos.

\begin{Form}
\textbf{Esta correspondência está adequada?} \\ 
\ChoiceMenu[radio,radiosymbol=\ding{52},bordercolor=0 0.5 0,name=cbovalid_ 771105 ]{}{CBO deve ficar=sim, \hspace{6em} CBO não fica aqui=nao}
\end{Form}

\textbf{CBO: 771110  ( Modelador de madeira )}

Planeja o trabalho, verificando ordens de serviço e de produção e
interpretando projetos, desenhos e especificações técnicas para a
confecção e restauração de matrizes de madeira e derivados. Seleciona
matéria-prima, insumos e componentes, instrumentos de medição e
controle, ferramentas manuais, ferramentas elétricas portáteis, máquinas
e equipamentos, para a execução do trabalho. Define roteiro para as
atividades de confecção e restauração. Pode elaborar orçamento. Elabora
croqui ou esboço e realiza cálculos para dimensionamento do produto,
conforme documentos de solicitação. Pode utilizar software de apoio para
desenhar o projeto completo em computador, especificando os materiais e
as operações para a confecção do produto. Confecciona gabaritos ou
moldes para execução das peças de madeira e derivados, de acordo com
especificações. Prepara máquinas e ferramentas. Afia ferramentas e
ajusta parâmetros operacionais, tais como de posicionamento e geometria
de corte, velocidade e pressão da ferramenta, entre outros. Confecciona
matrizes de madeira e derivados, aplicando processos de usinagem,
programando, operando e regulando as máquinas para obter o produto
conforme o projeto. Realiza cálculos e traça os contornos da peça,
utilizando instrumentos de desenho e traçagem e ajustando o
posicionamento e a geometria do corte. Usina as partes, definindo a
sequência de operações a serem aplicadas para cortar, serrar, furar,
aplainar, tornear, lixar, fresar e laminar. Ajusta os elementos e aplica
materiais e dispositivos de fixação. Pode programar e utilizar
máquinas-ferramenta com comando numérico computadorizado (CNC) e
programas para usinagem das peças integrados a ferramentas de software
de CAD/CAM. Monta matrizes de madeira e derivados utilizando elementos
de fixação e colocando acessórios e ferragens para reajuste. Coloca
arremates finais. Verifica as peças e os produtos confeccionados, para
garantir a qualidade das matrizes de madeira e derivados, avaliando a
matéria-prima, a resistência física, aferindo as formas e as dimensões e
as condições do acabamento. Compara as características do produto final
com os requisitos do cliente ou do projeto, providenciando ajustes,
quando necessários. Desmonta o produto confeccionado, no todo ou em
parte, para posterior entrega. Prepara a expedição de matrizes de
madeira e derivados conforme demanda do cliente, embalando separadamente
os acessórios e os elementos principais; definindo as ferramentas a
serem utilizadas para fazer o transporte e a entrega do produto; e
apondo os arremates finais. Restaura matrizes de madeira e derivados,
avaliando produtos para determinar a viabilidade de restauração,
confeccionando peças para reposição, substituindo peças danificadas,
reapertando elementos de fixação e preparando o produto para acabamento.
Organiza o local de trabalho, para o desenvolvimento das atividades.
Mantém os materiais, instrumentos e acessórios de trabalho organizados e
acondicionados. Mantém as máquinas, os equipamentos e as ferramentas em
plenas condições de uso e funcionamento. Realiza a limpeza e a
conservação dos recursos materiais. Controla desperdícios de material.
Separa, identifica e classifica resíduos -- sólidos e líquidos -, para
descarte, reúso e reciclagem. Preenche documentação técnica relacionada
aos processos, tais como ordens de serviço, fichas de controle de
processo, ficha técnica do produto, requisição de materiais, entre
outros. Interpreta normas, manuais técnicos e catálogos. Utiliza
equipamentos de proteção individual e coletiva.

\begin{Form}
\textbf{Esta correspondência está adequada?} \\ 
\ChoiceMenu[radio,radiosymbol=\ding{52},bordercolor=0 0.5 0,name=cbovalid_ 771110 ]{}{CBO deve ficar=sim, \hspace{6em} CBO não fica aqui=nao}
\end{Form}

\textbf{CBO: 771115  ( Maquetista na marcenaria )}

Planeja o trabalho, verificando ordens de serviço e de produção e
interpretando projetos, desenhos e especificações técnicas para a
confecção e a restauração de maquetes de madeira e derivados. Define
roteiro para as atividades. Seleciona matéria-prima, insumos e
componentes, ferramentas manuais, ferramentas elétricas portáteis,
instrumentos de medição e controle, acessórios, máquinas e equipamentos.
Pode elaborar orçamento. Elabora croqui ou esboço, realizando cálculos
para dimensionamento do produto, conforme documentos de solicitação.
Pode utilizar software para desenhar ou acessar o projeto completo em
computador, especificando os materiais e as operações para confecção e
acabamento do produto. Confecciona gabaritos ou moldes para execução das
peças para maquetes. Prepara máquinas e ferramentas. Afia ferramentas e
ajusta parâmetros operacionais, tais como de posicionamento e geometria
de corte, velocidade e pressão da ferramenta, entre outros. Confecciona
maquetes de madeira e derivados, aplicando processos de usinagem,
programando, operando e regulando as máquinas para obter o produto
conforme o projeto. Realiza cálculos e traça os contornos das partes,
utilizando instrumentos de desenho e traçagem e ajustando o
posicionamento e a geometria do corte. Usina peças, definindo a
sequência de operações a serem aplicadas para cortar, serrar, furar,
aplainar, tornear, lixar, fresar e laminar. Ajusta os elementos e aplica
materiais e dispositivos de fixação. Pode programar e utilizar
máquinas-ferramenta com comando numérico computadorizado (CNC) e
programas para usinagem das peças integrados a ferramentas de software
de CAD/CAM. Monta as maquetes, definindo a sequência de montagem e
realizando operações de união de peças -- tais como unir com cavilha
e/ou cola, unir com dispositivos, entre outras - para prensar. Utiliza
elementos de fixação, acessórios, ferragens e dispositivos de montagem.
Realiza operações de acabamento das maquetes, aplicando tintas e
vernizes; revestindo e polindo; aplicando cola - com coladeira de borda
ou sob pressão - ou massa para montagem; apondo apliques e lâminas;
colocando ferragens para reajuste; e regulando o funcionamento das
partes móveis, de acordo com o projeto. Pode realizar teste de
viscosidade nos produtos utilizados para acabamento da superfície.
Verifica peças e produtos confeccionados, para garantir a qualidade das
maquetes, avaliando a matéria-prima, a resistência física, as formas e
dimensões, os encaixes, a união e fixação das partes, as condições do
acabamento da superfície e da aplicação de ferragens, acessórios e
dispositivos. Compara as características do produto final com os
requisitos do cliente ou do projeto, providenciando ajustes, quando
necessários. Desmonta o produto confeccionado, para posterior entrega.
Restaura maquetes de madeira e derivados, avaliando produtos para
determinar a viabilidade de restauração, confeccionando peças para
reposição, substituindo peças danificadas, reapertando elementos de
fixação e preparando o produto para acabamento. Prepara a entrega das
maquetes confeccionadas sob medida ou restauradas, embalando
separadamente os acessórios e os elementos principais; definindo as
ferramentas a serem utilizadas para fazer o transporte e a montagem do
produto no local da instalação; e apondo os arremates finais. Organiza o
local de trabalho. Mantém os materiais, instrumentos e acessórios de
trabalho organizados e acondicionados. Mantém as máquinas, os
equipamentos e as ferramentas em plenas condições de uso e
funcionamento. Realiza a limpeza e a conservação dos recursos materiais.
Controla desperdícios de material. Separa, identifica e classifica
resíduos -- sólidos e líquidos -, para descarte, reúso e reciclagem.
Preenche fichas de controle de processo, ficha técnica do produto,
requisição de materiais, entre outros documentos. Interpreta normas,
manuais técnicos e catálogos. Utiliza equipamentos de proteção
individual e coletiva.

\begin{Form}
\textbf{Esta correspondência está adequada?} \\ 
\ChoiceMenu[radio,radiosymbol=\ding{52},bordercolor=0 0.5 0,name=cbovalid_ 771115 ]{}{CBO deve ficar=sim, \hspace{6em} CBO não fica aqui=nao}
\end{Form}

\newpage
\section{ 09015 -- Técnico em Design Gráfico }

\textbf{Perfil do Curso (CNCT)}

O Técnico em Design Gráfico será habilitado para:

Coordenar e executar projetos de comunicação visual de mídias impressas
e digitais seguindo padrões e normas técnicas, referentes à propriedade
intelectual, acessibilidade, usabilidade e sustentabilidade. Aplicar e
implementar sinalizações, ilustrações, tipografias, animações e
usabilidade de projetos de design gráfico. Analisar, interpretar e
propor a produção da identidade visual das peças gráficas. Criar,
controlar, organizar e armazenar arquivos e materiais de produção
gráfica e digital.

Para atuação como Técnico em Design Gráfico, são fundamentais:

Conhecimentos interdisciplinares relacionados aos processos de criação,
envolvendo pesquisa, idealização, planejamento, execução técnica,
fruição e recepção estética. Competências comunicativas e empreendedoras
voltadas à proposição de projetos, ao coletivo, à gestão, à solução de
problemas e à resiliência, entre outras competências socioemocionais.

\textbf{CBOs correspondidos e seus perfis (presentes no CNCT, e presentes na predição da IA)}

\textbf{CBO: 318405  ( Desenhista técnico (artes gráficas) )}

Prepara as atividades, organizando o posto de trabalho de desenho,
inspecionando equipamentos e selecionando ferramentas, instrumentos e
materiais. Interpreta solicitação de desenho, reunindo e organizando
informações pertinentes. Consulta revistas, catálogos e outras
publicações técnicas. Desenvolve esboços manualmente e com auxílio de
recursos computacionais. Define os meios de representação gráfica,
empregando recursos analógicos e/ou digitais adequados à demanda.
Planeja a elaboração do desenho, ordenando o trabalho por etapas e
fixando prazo de entrega. Define escalas e estabelece o formato da
apresentação. Consulta normas técnicas de desenho e de artes gráficas,
para especificar características do desenho. Elabora desenhos de
pré-impressão, combinando figuras, ilustrações, gráficos, textos e
outras reproduções visíveis, para compor a imagem a ser impressa.
Realiza representações gráficas pelo método manual, fazendo uso de
instrumentos para execução do traçado diretamente em substratos de
papel, bem como utilizando computadores, programas integrados de
CAD-Desenho Assistido por Computador (Computer Aided Design),
ilustração, tratamento de imagens e editoração gráfica. Codifica os
desenhos e relaciona as especificações técnicas envolvidas em sua
produção. Confere as especificações em face do trabalho realizado e
examina o atendimento às normas técnicas de representação gráfica, para
imprimir e requisitar a aprovação dos desenhos. Realiza as correções
indicadas pelo solicitante. Registra e arquiva o desenho aprovado.
Executa o acabamento final dos desenhos, indicando as características de
materiais de acabamento. Obtém a aprovação final. Confecciona a matriz
do desenho, preenche a legenda e tira as cópias de segurança (backup).
Conserva o local de trabalho limpo e organizado. Controla graus de
luminosidade do ambiente, temperatura e ruídos. Verifica as condições
ergonômicas para uso de mobiliário, monitor de vídeo de computador,
mouse e teclado. Requisita a aquisição de materiais e instrumentos de
trabalho, para reposição. Conserva equipamentos, ferramentas e
instrumentos de trabalho em plenas condições de uso e funcionamento.
Requisita serviço para manutenção de computador, impressoras, plotters e
outros periféricos, quando necessário. Organiza e arquiva documentos
físicos e eletrônicos, tais como comprovantes de aprovação e cópias de
segurança dos desenhos aprovados. Destina adequadamente os resíduos
gerados, fazendo reciclagem e realizando descarte de acordo com as
normas ambientais.

\begin{Form}
\textbf{Esta correspondência está adequada?} \\ 
\ChoiceMenu[radio,radiosymbol=\ding{52},bordercolor=0 0.5 0,name=cbovalid_ 318405 ]{}{CBO deve ficar=sim, \hspace{6em} CBO não fica aqui=nao}
\end{Form}

\textbf{CBO: 318410  ( Desenhista técnico (ilustrações artísticas) )}

Prepara as atividades, organizando o posto de trabalho de desenho,
inspecionando equipamentos e selecionando ferramentas, instrumentos e
materiais. Interpreta solicitação de desenho, reunindo e organizando as
informações pertinentes. Consulta revistas, catálogos e outras
publicações técnicas. Desenvolve esboços manualmente e com auxílio de
recursos computacionais. Define os meios de representação gráfica,
empregando recursos analógicos e/ou digitais adequados à demanda.
Planeja a elaboração do desenho artístico para ilustração, ordenando o
trabalho por etapas e fixando prazo de entrega. Define escalas e
estabelece o formato de apresentação em folhetos, catálogos, manuais,
documentação técnica em geral e textos de outros trabalhos, impressos
e/ou eletrônicos. Consulta normas técnicas para especificar
características do desenho. Desenvolve o desenho, combinando elementos
de imagens, para compor a ilustração artística a ser inserida nos
trabalhos ou utilizada à parte. Realiza representação gráfica pelo
método manual, fazendo uso de instrumentos para execução do traçado
diretamente em substratos de papel, bem como utilizando computador,
programas integrados de CAD-Desenho Assistido por Computador (Computer
Aided Design), ilustração, tratamento de imagens e editoração gráfica.
Codifica o desenho e relaciona as especificações técnicas envolvidas em
sua produção. Confere as especificações em face do trabalho realizado e
examina o atendimento às normas técnicas de representação gráfica, para
imprimir e requisitar a aprovação do desenho. Realiza as correções
indicadas pelo solicitante, registra e arquiva o desenho aprovado.
Executa o acabamento final do desenho, indicando as características de
materiais de acabamento. Obtém a aprovação final. Confecciona a matriz
da ilustração artística, preenche a legenda do desenho e tira as cópias
de segurança (back-up). Conserva o local de trabalho limpo e organizado.
Controla graus de luminosidade do ambiente, temperatura e ruídos.
Verifica as condições ergonômicas para uso de mobiliário, monitor de
vídeo de computador, mouse e teclado. Requisita a aquisição de materiais
e instrumentos de trabalho, para reposição. Conserva equipamentos,
ferramentas e instrumentos de trabalho em plenas condições de uso e
funcionamento. Requisita serviço para manutenção de computador,
impressoras, plotters e outros periféricos, quando necessário. Organiza
e arquiva documentos físicos e eletrônicos, tais como comprovantes de
aprovação e cópias de segurança dos desenhos aprovados. Destina
adequadamente os resíduos gerados, fazendo reciclagem e realizando
descarte de acordo com as normas ambientais.

\begin{Form}
\textbf{Esta correspondência está adequada?} \\ 
\ChoiceMenu[radio,radiosymbol=\ding{52},bordercolor=0 0.5 0,name=cbovalid_ 318410 ]{}{CBO deve ficar=sim, \hspace{6em} CBO não fica aqui=nao}
\end{Form}

\textbf{CBO: 318415  ( Desenhista técnico (ilustrações técnicas) )}

Prepara as atividades, organizando o posto de trabalho de desenho,
inspecionando equipamentos e selecionando ferramentas, instrumentos e
materiais. Interpreta solicitação de desenho, reunindo e organizando as
informações pertinentes. Consulta revistas, catálogos e outras
publicações técnicas. Desenvolve esboços manualmente e com auxílio de
recursos computacionais. Define os meios de representação gráfica,
empregando recursos analógicos e/ou digitais adequados à demanda.
Planeja a elaboração do desenho técnico para ilustração, ordenando o
trabalho por etapas e fixando prazo de entrega. Define escalas e
estabelece o formato de apresentação em folhetos, catálogos, manuais,
documentação técnica em geral e outros trabalhos, impressos e/ou
eletrônicos. Consulta normas técnicas para especificar características
do desenho. Desenvolve o desenho, combinando elementos de imagens, para
compor a ilustração técnica a ser inserida nos trabalhos ou utilizada à
parte. Realiza representação gráfica pelo método manual, fazendo uso de
instrumentos para execução do traçado diretamente em substratos de
papel, bem como utilizando computadores, programas integrados de
CAD-Desenho Assistido por Computador (Computer Aided Design),
ilustração, tratamento de imagens e editoração gráfica. Codifica o
desenho e relaciona as especificações técnicas envolvidas em sua
produção. Confere as especificações em face do trabalho realizado e
examina o atendimento às normas técnicas de representação gráfica, para
imprimir e requisitar a aprovação do desenho. Realiza as correções
indicadas pelo solicitante, registra e arquiva o desenho aprovado.
Executa o acabamento final do desenho, indicando as características de
materiais de acabamento. Obtém a aprovação final. Confecciona a matriz
da ilustração técnica, preenche a legenda do desenho e tira as cópias de
segurança (back-up). Conserva o local de trabalho limpo e organizado.
Controla graus de luminosidade do ambiente, temperatura e ruídos.
Verifica as condições ergonômicas para uso de mobiliário, monitor de
vídeo de computador, mouse e teclado. Requisita a aquisição de materiais
e instrumentos de trabalho, para reposição. Conserva equipamentos,
ferramentas e instrumentos de trabalho em plenas condições de uso e
funcionamento. Requisita serviço para manutenção de computador,
impressoras, plotters e outros periféricos, quando necessário. Organiza
e arquiva documentos físicos e eletrônicos, tais como comprovantes de
aprovação e cópias de segurança dos desenhos aprovados. Destina
adequadamente os resíduos gerados, fazendo reciclagem e realizando
descarte de acordo com as normas ambientais.

\begin{Form}
\textbf{Esta correspondência está adequada?} \\ 
\ChoiceMenu[radio,radiosymbol=\ding{52},bordercolor=0 0.5 0,name=cbovalid_ 318415 ]{}{CBO deve ficar=sim, \hspace{6em} CBO não fica aqui=nao}
\end{Form}

\textbf{CBO: 318420  ( Desenhista técnico (indústria têxtil) )}

Prepara as atividades, organizando o posto de trabalho de desenho,
inspecionando equipamentos e selecionando ferramentas, instrumentos e
materiais. Interpreta solicitação de desenho e projeto de estamparia,
texturas e cores, reunindo e organizando informações pertinentes.
Consulta revistas, catálogos e outras publicações técnicas. Desenvolve
esboços manualmente e com auxílio de recursos computacionais. Define os
meios de representação gráfica, empregando recursos analógicos e/ou
digitais adequados à demanda. Planeja a elaboração de desenho e projeto
de estamparia, texturas e cores, ordenando o trabalho por etapas e
fixando prazo de entrega. Define escalas e estabelece o formato de
apresentação, impresso e/ou eletrônico. Consulta normas técnicas para
especificar características do desenho. Desenvolve o desenho, combinando
elementos de imagens, para compor a estampa dos tecidos. Realiza
representação gráfica pelo método manual, fazendo uso de instrumentos
para execução do traçado diretamente em substratos de papel, bem como
utilizando computadores, programas integrados de CAD-Desenho Assistido
por Computador (Computer Aided Design), ilustração e tratamento de
imagens. Codifica o desenho e relaciona as especificações técnicas
envolvidas em sua produção. Confere as especificações em face do
trabalho realizado e examina o atendimento às normas técnicas de
representação gráfica, para imprimir e requisitar a aprovação do
desenho. Realiza as correções indicadas pelo solicitante, registra e
arquiva o desenho aprovado. Executa o acabamento final do desenho,
indicando as características de materiais de acabamento. Obtém a
aprovação final. Confecciona a matriz do desenho da estampa, preenche a
legenda do desenho e tira cópias de segurança (back-up). Conserva o
local de trabalho limpo e organizado. Controla graus de luminosidade do
ambiente, temperatura e ruídos. Verifica as condições ergonômicas para
uso de mobiliário, monitor de vídeo de computador, mouse e teclado.
Requisita a aquisição de materiais e instrumentos de trabalho, para
reposição. Conserva equipamentos, ferramentas e instrumentos de trabalho
em plenas condições de uso e funcionamento. Requisita serviço para
manutenção de computador, impressoras e outros periféricos, quando
necessário. Organiza e arquiva documentos físicos e eletrônicos, tais
como comprovantes de aprovação e cópias de segurança dos desenhos e
projetos de estamparia aprovados. Destina adequadamente os resíduos
gerados, fazendo reciclagem e realizando descarte de acordo com as
normas ambientais.

\begin{Form}
\textbf{Esta correspondência está adequada?} \\ 
\ChoiceMenu[radio,radiosymbol=\ding{52},bordercolor=0 0.5 0,name=cbovalid_ 318420 ]{}{CBO deve ficar=sim, \hspace{6em} CBO não fica aqui=nao}
\end{Form}

\textbf{CBO: 318425  ( Desenhista técnico (mobiliário) )}

Prepara as atividades, organizando o posto de trabalho de desenho,
inspecionando equipamentos e selecionando ferramentas, instrumentos e
materiais. Interpreta solicitação de desenho, reunindo e organizando
informações pertinentes. Consulta revistas, catálogos e outras
publicações técnicas. Desenvolve esboços manualmente e com auxílio de
recursos computacionais. Define os meios de representação gráfica,
empregando recursos analógicos e/ou digitais adequados à demanda.
Planeja a elaboração de desenho de móveis, ordenando o trabalho por
etapas e fixando prazo de entrega. Define escalas e estabelece o formato
de apresentação, impresso e/ou eletrônico. Consulta normas técnicas para
especificar as características do desenho. Desenvolve o desenho,
combinando elementos de imagens, conforme a necessidade, para melhor
orientar a produção do mobiliário. Realiza representação gráfica pelo
método manual, fazendo uso de instrumentos para execução do traçado
diretamente em substratos de papel, bem como utilizando computadores,
programas integrados de CAD-Desenho Assistido por Computador (Computer
Aided Design), ilustração e tratamento de imagens. Codifica o desenho e
relaciona as especificações técnicas envolvidas em sua produção. Confere
as especificações em face do trabalho realizado e examina o atendimento
às normas técnicas de representação gráfica, para imprimir e requisitar
a aprovação do desenho. Realiza as correções indicadas pelo solicitante,
registra e arquiva o desenho aprovado. Executa o acabamento final do
desenho, indicando as características de materiais de acabamento. Obtém
a aprovação final. Confecciona a matriz do desenho, preenche a legenda
do desenho e tira as cópias de segurança (backup). Conserva o local de
trabalho limpo e organizado. Controla graus de luminosidade do ambiente,
temperatura e ruídos. Verifica as condições ergonômicas para uso de
mobiliário, monitor de vídeo de computador, mouse e teclado. Requisita a
aquisição de materiais e instrumentos de trabalho, para reposição.
Conserva equipamentos, ferramentas e instrumentos de trabalho em plenas
condições de uso e funcionamento. Requisita serviço para manutenção de
computador, impressoras e outros periféricos, quando necessário.
Organiza e arquiva documentos físicos e eletrônicos, tais como
comprovantes de aprovação e cópias de segurança dos desenhos aprovados.
Destina adequadamente os resíduos gerados, fazendo reciclagem e
realizando descarte de acordo com as normas ambientais.

\begin{Form}
\textbf{Esta correspondência está adequada?} \\ 
\ChoiceMenu[radio,radiosymbol=\ding{52},bordercolor=0 0.5 0,name=cbovalid_ 318425 ]{}{CBO deve ficar=sim, \hspace{6em} CBO não fica aqui=nao}
\end{Form}

\textbf{CBO: 318430  ( Desenhista técnico de embalagens, maquetes e leiautes )}

Prepara as atividades, organizando o posto de trabalho de desenho,
inspecionando equipamentos e selecionando ferramentas, instrumentos e
materiais. Interpreta solicitação de desenhos, reunindo e organizando
informações pertinentes. Consulta revistas, catálogos e outras
publicações técnicas. Desenvolve esboços manualmente e com auxílio de
recursos computacionais. Define os meios de representação gráfica,
empregando recursos analógicos e/ou digitais adequados à demanda.
Planeja a elaboração de desenhos de embalagens, maquete e leiaute,
ordenando o trabalho por etapas e fixando prazo de entrega. Define
escalas e estabelece o formato de apresentação, impresso e/ou
eletrônico. Consulta normas técnicas para especificar as características
dos desenhos. Desenvolve os desenhos, combinando elementos de imagens,
conforme a necessidade, para melhor orientar a fabricação de embalagens;
guiar a montagem e a reprodução eletrônica de maquetes; e definir a
disposição de elementos com referência em leiautes. Pode aplicar
recursos de animação computadorizada à maquete eletrônica. Realiza
representações gráficas pelo método manual, fazendo uso de instrumentos
para execução do traçado diretamente em substratos de papel, bem como
utilizando computadores, programas integrados de CAD-Desenho Assistido
por Computador (Computer Aided Design), de ilustração, de tratamento de
imagens e de animação computadorizada. Codifica os desenhos e relaciona
as especificações técnicas envolvidas em sua produção. Confere as
especificações em face do trabalho realizado e examina o atendimento às
normas técnicas de representação gráfica, para imprimir e requisitar a
aprovação dos desenhos. Realiza as correções indicadas pelo solicitante,
registra e arquiva os desenhos aprovados. Executa o acabamento final dos
desenhos, indicando as características de materiais de acabamento. Obtém
a aprovação final. Confecciona as matrizes dos desenhos, preenche a
legenda dos desenhos e tira as cópias de segurança (backup). Conserva o
local de trabalho limpo e organizado. Controla graus de luminosidade do
ambiente, temperatura e ruídos. Verifica as condições ergonômicas para
uso de mobiliário, monitor de vídeo de computador, mouse e teclado.
Requisita a aquisição de materiais e instrumentos de trabalho, para
reposição. Conserva equipamentos, ferramentas e instrumentos de trabalho
em plenas condições de uso e funcionamento. Requisita serviço para
manutenção de computador, impressoras e outros periféricos, quando
necessário. Organiza e arquiva documentos físicos e eletrônicos, tais
como comprovantes de aprovação e cópias de segurança dos desenhos
aprovados. Destina adequadamente os resíduos gerados, fazendo reciclagem
e realizando descarte de acordo com as normas ambientais.

\begin{Form}
\textbf{Esta correspondência está adequada?} \\ 
\ChoiceMenu[radio,radiosymbol=\ding{52},bordercolor=0 0.5 0,name=cbovalid_ 318430 ]{}{CBO deve ficar=sim, \hspace{6em} CBO não fica aqui=nao}
\end{Form}

\textbf{CBO: 371305  ( Técnico em programação visual )}

Planeja o serviço de programação visual, analisando ordem de serviço e
identificando necessidades de cliente. Estabelece prazo para elaboração
do projeto. Faz estimativa de custos, para aprovação de recursos que
poderão ser utilizados no projeto. Prepara especificação de recursos
materiais. Examina as condições dos equipamentos, testando seu
funcionamento. Confere a disponibilidade de materiais, requisitando os
faltantes. Seleciona ``softwares'' específicos para elaboração de
projeto. Prepara esboços e leiautes de amostras, para visualizar a
posição de imagens, blocos de textos e outros elementos de página.
Seleciona estilos e tamanho de fontes. Localiza imagens ou fotos de
produtos. Pesquisa conceitos de ``design'' a serem aplicados. Estabelece
padrões de referência de qualidade. Elabora o projeto de programação
visual de materiais -- publicitários, jornalísticos, entre outros -- a
serem impressos, aplicando elementos criativos e estéticos. Desenvolve
elementos de identidade visual. Cria ilustrações e desenha logotipos.
Efetua a editoração de textos e o tratamento de imagens, digitando,
diagramando, ilustrando e formatando. Elabora tabelas e gráficos. Em
publicações editoriais, define o tamanho da lombada e confecciona
boneco, miolo e capa. Prepara originais, verificando a qualidade, a
quantidade e a organização de textos e imagens. Confecciona prova
digital ou peça-piloto. Analisa a viabilidade econômica do projeto.
Solicita, a cliente, a aprovação da arte-final. Pode operar máquinas e
equipamentos ou executar processos eletrônicos, em testes de
processamento e reprodução de textos e imagens. Faz ajustes técnicos --
tais como de resolução, de espessura de linhas, de extensão das imagens,
dentre outros - da programação visual, antes do processo de impressão.
Realiza o fechamento do arquivo com a arte-final. Encaminha, junto com o
arquivo, ficha de impressão, que detalha informações técnicas - tais
como tipo e gramatura do papel; número de páginas; quantidade de
impresso; entre outras - para a produção gráfica. Presta assistência
técnica para as diversas áreas da empresa, fornecendo informações
técnicas. Avalia resultados da produção realizada junto a clientes.
Analisa reclamações e informa a clientes sobre possibilidades e
restrições de materiais e processos. Busca alternativas para melhorias
em tecnologias e matérias-primas. Realiza visitas técnicas a empresas e
fornecedores, para coleta de informações e análise de tecnologias
disponíveis. Consulta normas e fontes de informações técnicas. Avalia
tendências de mercado. Pesquisa ``sites'' especializados da internet e
verifica lançamentos em feiras nacionais e internacionais da indústria
gráfica, para manter-se atualizado em relação às inovações e à
introdução de novas tecnologias e de novos materiais na sua área de
atuação. Coordena equipe de trabalho, definindo atribuições e
assegurando o cumprimento de normas e procedimentos no trabalho. Pode
supervisionar o trabalho de equipe, avaliando desempenho e realizando
programas de treinamento. Monitora a organização, a classificação e o
arquivo de documentos físicos e digitais. Verifica a limpeza e a
organização do ambiente de trabalho. Orienta a equipe para conservar
instrumentos limpos, acondicionados e em plenas condições de uso e
funcionamento. Mantém máquinas e equipamentos em condições operacionais,
providenciando manutenção de primeiro nível. Requisita serviços de
manutenção corretiva de máquinas e equipamentos. Controla a aplicação de
procedimentos da empresa para cumprimento das normas ambientais.
Inspeciona o reaproveitamento de sobras de material e o descarte de
resíduos. Zela pela segurança, utilizando e monitorando o uso de
equipamentos de proteção individual. Identifica situações de risco e
colabora para sua eliminação.

\begin{Form}
\textbf{Esta correspondência está adequada?} \\ 
\ChoiceMenu[radio,radiosymbol=\ding{52},bordercolor=0 0.5 0,name=cbovalid_ 371305 ]{}{CBO deve ficar=sim, \hspace{6em} CBO não fica aqui=nao}
\end{Form}

\textbf{CBO: 371310  ( Técnico gráfico )}

Participa no planejamento de atividades, elaborando programação para
desenvolvimento da produção gráfica. Organiza arranjo físico em função
do programa de produção e analisa fluxograma. Realiza análise de custos,
efetuando cálculos para fechamento do preços. Elabora orçamentos. Faz
desenvolvimento de produtos gráficos, elaborando projeto de confecção e
edição. Realiza processo criativo de concepção de produto. Elabora
especificação de matérias-primas e outros insumos para a produção.
Define tecnologias, sistemas e processos para confecção dos produtos.
Orienta o fechamento de arquivo -- com resultado da fase de
pré-impressão - em aplicativos, para encaminhamento ao processo de
impressão. Analisa viabilidade técnica e econômica do projeto. Coordena
o processo de impressão e presta informações técnicas - tais como tipo e
gramatura do papel a ser impresso -, de acordo com o projeto gráfico.
Pode operar o processo de impressão e reprodução de produto gráfico.
Estabelece o processo de pós-impressão, para realização de cortes
especiais, relevos e outras operações. Seleciona equipamentos de
acabamento editorial e cartotécnico. Pode ajustar e operar equipamentos
de pós-impressão. Analisa e avalia as características do produto
confeccionado. Monitora estoque do produto acabado. Controla a qualidade
do processo e do produto, estabelecendo padrões de referência. Analisa
variáveis do processo, utilizando metodologia de análise de problemas.
Audita processos a partir de normas estabelecidas. Define plano de
amostragem e realiza ensaios em amostras de matérias-primas e produtos.
Emite relatório de não conformidade do produto. Realiza controle dos
processos gráficos, alimentando bancos de dados, monitorando o andamento
das atividades e calculando os índices de produção. Busca alternativas
para melhorias em tecnologias e matérias-primas. Realiza visitas
técnicas a empresas e fornecedores, para coleta de informações e análise
de tecnologias disponíveis. Consulta normas e fontes de informações
técnicas. Avalia tendências de mercado. Pesquisa ``sites''
especializados da internet e acompanha lançamentos em feiras nacionais e
internacionais da indústria gráfica. Implanta novas tecnologias,
estimando custos de implantação e de operação e orientando fornecedores
para adequação de insumos às tecnologias implantadas. Define
características técnicas de equipamentos e promove melhoria
(``upgrade'') em seu desempenho. Presta assistência técnica, orientando
clientes sobre possibilidades e restrições de materiais e processos.
Propõe soluções corretivas no local de ocorrência de problemas. Fornece
informações técnicas para as diversas áreas da empresa e elabora manuais
de procedimentos. Emite relatório técnico. Coordena equipe de trabalho,
definindo atribuições e assegurando o cumprimento de normas e
procedimentos no trabalho. Pode prestar orientação no uso de
``softwares'' e na operação de máquinas e equipamentos. Pode
supervisionar o trabalho de equipe, avaliando desempenho e levantando
necessidades de treinamento. Realiza programas de treinamento. Verifica
a limpeza e a organização do ambiente de trabalho. Orienta a equipe para
conservar instrumentos limpos, acondicionados e em plenas condições de
uso e funcionamento. Mantém máquinas e equipamentos em condições
operacionais, providenciando manutenção de primeiro nível. Requisita
serviço de manutenção corretiva de máquinas e equipamentos, quando
necessário. Controla a aplicação de procedimentos da empresa para
cumprimento das normas ambientais. Inspeciona o reaproveitamento de
sobras de material e o descarte de resíduos. Zela pela segurança,
utilizando e monitorando o uso de equipamentos de proteção individual.
Identifica situações de risco e colabora para sua eliminação.

\begin{Form}
\textbf{Esta correspondência está adequada?} \\ 
\ChoiceMenu[radio,radiosymbol=\ding{52},bordercolor=0 0.5 0,name=cbovalid_ 371310 ]{}{CBO deve ficar=sim, \hspace{6em} CBO não fica aqui=nao}
\end{Form}

\newpage
\section{ 09018 -- Técnico em Fabricação de Instrumentos Musicais }

\textbf{Perfil do Curso (CNCT)}

O Técnico em Fabricação de Instrumentos Musicais será habilitado para:

Projetar instrumentos musicais. Distinguir acústicas de materiais para a
fabricação dos instrumentos musicais. Preparar matérias-primas para
confecção dos instrumentos. Confeccionar componentes dos instrumentos.
Realizar acabamentos, montar, afinar, consertar e vender instrumentos
musicais.

Para atuação como Técnico em Fabricação de Instrumentos Musicais, são
fundamentais:

Conhecimentos interdisciplinares relacionados aos processos de criação,
envolvendo pesquisa, idealização, planejamento, execução técnica,
fruição e recepção estética. Competências comunicativas e empreendedoras
voltadas à proposição de projetos, ao coletivo, à gestão, à solução de
problemas e à resiliência, entre outras competências socioemocionais.

\textbf{CBOs correspondidos e seus perfis (presentes no CNCT, e presentes na predição da IA)}

\textbf{CBO: 742105  ( Afinador de instrumentos musicais )}

Planeja as atividades, definindo as etapas e a sequência de trabalho, de
acordo com o tipo de instrumento musical a ser afinado. Seleciona
equipamentos, ferramentas, instrumentos, componentes e acessórios
necessários, considerando requisição de trabalho, especificações,
manuais, laudos e pareceres técnicos. Testa o instrumento musical, para
avaliar a qualidade do som e localizar defeitos. Revisa componentes
mecânicos do instrumento musical e detecta vazão de ar -- no caso de
instrumento de sopro -, durante a análise de problemas. Opera diapasão,
comparando seu som com o produzido pelo instrumento musical, para
detectar as dissonâncias. Faz análise, ainda, comparando os sons por
audição. Realiza afinação de instrumentos musicais, de sopro, de cordas
e de percussão, utilizando afinadores analógicos, digitais ou
automatizados. Regula tensores de instrumento. Executa limpeza de
orifícios, componentes e lâminas de madeira, metal e sintéticas. Pode
utilizar softwares e aplicativos. Analisa a geração do som, utilizando
tabela de acordes e seguindo especificações e parâmetros técnicos.
Avalia a qualidade e a precisão da afinação realizada, aplicando testes
auditivos e visuais e utilizando instrumentos e equipamentos. Examina e
ajusta parâmetros de frequência do som, altura, intensidade, timbre,
articulação, duração, entre outros. Define procedimentos para
conservação dos instrumentos musicais. Pode providenciar laudos
técnicos. Acondiciona o instrumento musical, tendo em vista manter sua
integridade física. Pode reparar instrumentos, inspecionando-os para
localizar defeitos, desmontando-os, ajustando as partes e substituindo
as peças desgastadas ou defeituosas. Remonta os instrumentos após o
reparo. Pode instalar, montar e afinar instrumentos em locais de
espetáculo. Para não comprometer a performance do artista, leva em
consideração, no trabalho de afinação, as características da orquestra
em que o músico vai tocar. Considera, ainda, a temperatura, a umidade do
ar e o calor das luzes do local de espetáculo, que podem interferir nas
peças metálicas. Pode desmontar os instrumentos após apresentações
itinerantes. Elabora documentos técnicos, tais como relatórios de
serviços, manuais, tabelas, especificações, laudos e fichas de controle
de afinações e de reparos. Mantém o local de trabalho limpo e
organizado. Conserva as ferramentas, componentes, acessórios e
instrumentos de trabalho limpos, organizados, acondicionados e em plenas
condições de uso e funcionamento. Providencia os serviços de manutenção
de equipamentos. Controla desperdícios de materiais, separando sobras
para reaproveitamento. Identifica resíduos, providenciando sua
segregação e destinação de acordo com normas ambientais. Trabalha com
segurança, utilizando equipamentos de proteção individual e atuando na
prevenção de acidentes.

\begin{Form}
\textbf{Esta correspondência está adequada?} \\ 
\ChoiceMenu[radio,radiosymbol=\ding{52},bordercolor=0 0.5 0,name=cbovalid_ 742105 ]{}{CBO deve ficar=sim, \hspace{6em} CBO não fica aqui=nao}
\end{Form}

\textbf{CBO: 742110  ( Confeccionador de acordeão )}

Planeja a produção de acordeão, em série ou sob encomenda, programando
as etapas e avaliando o tempo de fabricação. Verifica as medidas e as
matérias-primas a serem usadas. Seleciona máquinas, equipamentos,
ferramentas, acessórios e instrumentos de medição. Calcula custos. Pode
esboçar croqui e confeccionar protótipo. Prepara matérias-primas para
confecção do acordeão, definindo os tipos de cortes dos materiais.
Classifica a madeira e faz sua secagem em estufa. Executa o tingimento
das lâminas de madeira. Confecciona componentes do acordeão, construindo
moldes; manuseando gabaritos; utilizando ferramentas e instrumentos de
medição; regulando e operando máquinas; e colando materiais. Executa
lixamento primário em superfície e solda componentes metálicos.
Confecciona acordeão, montando a estrutura composta por uma armação de
madeira (castelo ou cavalete), fole e palhetas. Encaixa, une, fixa e
ajusta as demais partes, componentes e acessórios - tais como teclado,
alavancas de válvula, registro de melodia, trava de fole, baixos,
registros de baixo, botão de saída de ar e correias de sustentação -,
para executar a montagem do acordeão. Realiza a confecção de acordo com
ordem de produção, desenhos, protótipos, gabaritos, especificações e
manuais técnicos. Utiliza, no processo de produção, máquinas e
equipamentos mecânicos e elétricos, além de instrumentos e ferramentas.
Pode utilizar máquinas e equipamentos digitais e automatizados. Realiza
ajustes de montagem e acabamento, regulando componentes
eletroeletrônicos; aplicando banhos eletrostáticos; usando tintas para
revestimento; jateando com areia e polindo superfícies metálicas. Coloca
etiquetas identificadoras e instala proteção contra agentes externos.
Realiza afinação do acordeão, utilizando afinadores analógicos, digitais
ou automatizados. Opera diapasão, comparando seu som com o produzido
pelo acordeão, para detectar as dissonâncias. Faz análise, ainda,
comparando os sons por audição. Usa captadores de som e equipamentos
mecânicos, eletrônicos e de informática, para examinar a frequência do
som, altura, intensidade, timbre, articulação, entre outros itens.
Revisa componentes mecânicos, verifica vazamento de ar e regula
tensores. Limpa componentes e lâminas de madeira, de metal e sintéticas.
Avalia a geração do som, utilizando tabela de acordes. Avalia a
qualidade do acordeão, aplicando testes sensoriais e visuais e usando
instrumentos e equipamentos, para verificar e ajustar parâmetros
dimensionais, estruturais, funcionais e acústicos. Pode providenciar
laudos técnicos. Acondiciona o acordeão, tendo em vista manter a
integridade física das estruturas, partes e componentes. Pode realizar a
venda de acordeão, pesquisando mercados e divulgando o produto. Calcula
o valor de venda e estabelece o mecanismo de cobrança. Após efetivar a
venda, seleciona os meios de transporte para entrega. Analisa a
satisfação do cliente. Pode executar conserto e restauração de acordeão,
avaliando as informações do cliente, identificando a viabilidade do
serviço e orçando os custos. Compra componentes e constrói peças de
reposição. Desmonta o acordeão e substitui partes e componentes
desgastados ou defeituosos. Após realização do serviço, monta o
acordeão. Prepara documentos técnicos, participando da elaboração de
manuais técnicos e preenchendo relatórios. Mantém o local de trabalho
limpo e organizado. Conserva ferramentas, acessórios e instrumentos de
trabalho limpos, organizados, acondicionados e em plenas condições de
uso e funcionamento. Providencia serviços de manutenção de equipamentos.
Controla desperdícios de materiais, realizando reaproveitamento.
Identifica resíduos, providenciando sua destinação de acordo com normas
ambientais. Trabalha com segurança, utilizando equipamentos de proteção
individual e atuando na prevenção de acidentes.

\begin{Form}
\textbf{Esta correspondência está adequada?} \\ 
\ChoiceMenu[radio,radiosymbol=\ding{52},bordercolor=0 0.5 0,name=cbovalid_ 742110 ]{}{CBO deve ficar=sim, \hspace{6em} CBO não fica aqui=nao}
\end{Form}

\textbf{CBO: 742115  ( Confeccionador de instrumentos de corda )}

Planeja a produção de instrumentos musicais de corda, em série ou sob
encomenda, considerando o tipo - instrumentos cordófonos dedilhados
eletrificados, dedilhados acústicos ou instrumentos cordófonos a arco
acústicos -- para programar as etapas e avaliar o tempo de fabricação.
Verifica a técnica a ser aplicada no processo de produção, confere as
matérias-primas a serem usadas e seleciona máquinas, equipamentos,
ferramentas, acessórios e instrumentos de medição, de acordo com ordem
de produção, especificações, fichas técnicas e manuais. Pode calcular
custos. Prepara a madeira, o metal e outros materiais para construção
dos componentes utilizados na confecção de instrumentos de corda.
Desdobra e classifica madeira. Seca madeira na estufa. Prensa e tinge
lâminas de madeira. Confecciona componentes, construindo moldes;
manuseando gabaritos; utilizando ferramentas e instrumentos de medição;
regulando e operando máquinas; e colando materiais. Calandra madeiras e
metais, executa lixamento primário em superfície e solda componentes
metálicos. Confecciona variados instrumentos de corda - tais como
violão, guitarra, cavaquinho, violino, violoncelo e harpa -, de acordo
com ordem de produção, desenhos, protótipos e manuais técnicos. Monta a
estrutura (corpo do instrumento de corda), composta de uma caixa de
ressonância. Encaixa, une, fixa e ajusta partes, componentes e
acessórios - como braço, cravelhas, cordas, tensor, roseta, pestana,
entre outros -, para realizar a montagem do instrumento de corda.
Utiliza, no processo de produção, máquinas e equipamentos mecânicos e
elétricos, instrumentos e ferramentas. Pode utilizar máquinas e
equipamentos digitais e automatizados. Realiza ajustes de montagem e
acabamento, encaixando e regulando componentes eletroeletrônicos;
verificando a altura, a espessura e o comprimento das cordas, para
fixá-las de forma a manter a distância predeterminada entre elas; e
unindo os conjuntos de elementos. Aplica tintas, vernizes, óleos, ceras
e outros materiais e produtos para revestimento. Coloca etiquetas
identificadoras e instala proteção contra agentes externos. Realiza
afinação do instrumento de corda, aumentando ou diminuindo a tensão das
cordas. Faz comparação dos sons por audição e pode usar diapasão normal
ou digital, software e aplicativos, para identificar dissonâncias e
afinar o instrumento. Avalia a qualidade do instrumento de corda,
aplicando testes sensoriais e visuais e usando instrumentos de medição e
equipamentos, para verificar e ajustar parâmetros dimensionais,
estruturais, funcionais e acústicos. Pode providenciar laudos técnicos.
Faz o acondicionamento do instrumento de corda produzido. Pode realizar
a venda de instrumento de corda, pesquisando mercados e divulgando o
produto. Calcula o valor de venda e estabelece o mecanismo de cobrança.
Após realizar a venda, seleciona os meios de transporte para entrega.
Analisa a satisfação do cliente. Pode realizar conserto e restauração de
instrumentos de corda, avaliando as informações do cliente,
identificando a viabilidade do serviço e orçando os custos. Compra
componentes e constrói peças de reposição. Desmonta o instrumento e
substitui partes e componentes desgastados ou defeituosos. Após
realização do serviço, monta o instrumento de corda. Prepara documentos
técnicos, participando da elaboração de manuais técnicos e preenchendo
relatórios. Mantém o local de trabalho limpo e organizado. Conserva
ferramentas, acessórios e instrumentos de trabalho limpos, organizados,
acondicionados e em plenas condições de uso e funcionamento. Providencia
serviços de manutenção de equipamentos. Controla desperdícios de
materiais, realizando reaproveitamento. Identifica resíduos,
providenciando sua destinação de acordo com normas ambientais. Trabalha
com segurança, utilizando equipamentos de proteção individual e atuando
na prevenção de acidentes.

\begin{Form}
\textbf{Esta correspondência está adequada?} \\ 
\ChoiceMenu[radio,radiosymbol=\ding{52},bordercolor=0 0.5 0,name=cbovalid_ 742115 ]{}{CBO deve ficar=sim, \hspace{6em} CBO não fica aqui=nao}
\end{Form}

\textbf{CBO: 742120  ( Confeccionador de instrumentos de percussão (pele, couro ou plástico) )}

Planeja a produção de instrumentos musicais de percussão, em série ou
sob encomenda, programando as etapas e avaliando o tempo de fabricação.
Verifica a técnica a ser aplicada no processo de produção, confere as
matérias-primas a serem usadas e seleciona máquinas, equipamentos,
ferramentas, acessórios e instrumentos de medição e controle, de acordo
com ordem de produção, especificações, fichas técnicas e manuais
técnicos. Pode calcular custos. Prepara as matérias-primas a serem
utilizadas na confecção de instrumentos de percussão. Distingue
acústicas dos materiais, analisando as propriedades de ressonância e
identificando diferentes timbres. Classifica madeira e faz seu secamento
na estufa. Tinge lâminas de madeira. Confecciona componentes,
construindo moldes; manuseando gabaritos; utilizando ferramentas e
instrumentos de medição e controle; regulando e operando máquinas; e
colando materiais. Calandra madeiras e metais, executa lixamento
primário em superfície e solda componentes metálicos. Confecciona
instrumentos de percussão -- de impacto, atrito ou agitação -, tais como
tambor, surdo, pandeiro, bateria, triângulo, marimba, chocalho, cuíca,
entre outros. Executa a montagem do instrumento, montando a estrutura e
encaixando, unindo e fixando partes, componentes e acessórios. Realiza
operações - manuais, mecânicas e automatizadas - de traçagem, corte,
usinagem, recozimento, prensagem, furação, lixamento, rebitagem, entre
outras. Realiza acabamento nos instrumentos musicais, preparando
superfícies; pintando ou encerando; revestindo instrumentos com
celuloide, acetato ou outros materiais; jateando com areia as
superfícies dos instrumentos metálicos e peles sintéticas; e polindo
superfícies. Executa a confecção de acordo com ordem de produção,
desenhos, protótipos, especificações e manuais técnicos. Utiliza, no
processo de produção, máquinas e equipamentos mecânicos e elétricos,
instrumentos e ferramentas. Pode utilizar máquinas e equipamentos
digitais e automatizados. Aplica técnicas de afinação do instrumento de
percussão. Coloca etiquetas identificadoras nos instrumentos produzidos
e aplica procedimentos de proteção contra agentes externos. Avalia a
qualidade dos instrumentos de percussão, aplicando testes sensoriais e
visuais e usando instrumentos de medição e equipamentos, para verificar
e ajustar parâmetros dimensionais, estruturais, funcionais e acústicos.
Pode providenciar laudos técnicos. Pode tocar o instrumento para avaliar
a qualidade do som, além de identificar e corrigir defeitos. Acondiciona
os instrumentos de percussão, tendo em vista manter a integridade física
das estruturas, partes e componentes. Pode realizar a venda de
instrumento de percussão, pesquisando mercados e divulgando o produto.
Calcula o valor de venda e estabelece o mecanismo de cobrança. Após
realizar a venda, seleciona os meios de transporte para entrega. Analisa
a satisfação do cliente. Pode realizar conserto e restauração de
instrumentos de percussão, avaliando as informações do cliente,
identificando a viabilidade do serviço e orçando os custos. Compra
componentes e constrói peças de reposição. Desmonta o instrumento e
substitui partes e componentes desgastados ou defeituosos. Após
realização do serviço, monta o instrumento de percussão. Prepara
documentos técnicos, participando da elaboração de manuais técnicos e
preenchendo relatórios. Mantém o local de trabalho limpo e organizado.
Conserva ferramentas, acessórios e instrumentos de trabalho limpos,
organizados, acondicionados e em plenas condições de uso e
funcionamento. Providencia serviços de manutenção de equipamentos.
Controla desperdícios de materiais, realizando reaproveitamento.
Identifica resíduos e providencia sua destinação de acordo com normas
ambientais. Trabalha com segurança, utilizando equipamentos de proteção
individual e atuando na prevenção de acidentes.

\begin{Form}
\textbf{Esta correspondência está adequada?} \\ 
\ChoiceMenu[radio,radiosymbol=\ding{52},bordercolor=0 0.5 0,name=cbovalid_ 742120 ]{}{CBO deve ficar=sim, \hspace{6em} CBO não fica aqui=nao}
\end{Form}

\textbf{CBO: 742125  ( Confeccionador de instrumentos de sopro (madeira) )}

Planeja a produção de instrumentos de sopro de madeira, assim
classificados pelo seu funcionamento em que, na maior parte dos casos, o
sopro do músico sobre uma palheta - lâmina flexível simples ou dupla,
geralmente feita de madeira ou bambu - faz com que ela vibre e produza o
som desejado. Programa as etapas e avalia o tempo de produção. Verifica
a técnica a ser aplicada no processo de produção, confere as
matérias-primas a serem usadas e seleciona máquinas, equipamentos,
ferramentas, acessórios e instrumentos de medição, de acordo com ordem
de produção, especificações, fichas técnicas e manuais técnicos. Pode
calcular custos. Prepara as matérias-primas a serem utilizadas na
confecção de instrumentos de sopro de madeira. Distingue acústicas dos
materiais, analisando as propriedades de ressonância e identificando os
diferentes timbres. Define os tipos de cortes dos materiais. Classifica
madeira e faz seu secamento na estufa. Tinge lâminas de madeira.
Confecciona componentes, construindo moldes; manuseando gabaritos;
utilizando ferramentas e instrumentos de medição; regulando e operando
máquinas; e colando materiais. Executa lixamento primário em superfície
e solda componentes metálicos. Confecciona instrumentos de sopro de
madeira, tais como flautas, oboé, clarinetes e fagotes. Monta a
estrutura (corpo do instrumento), composta por um tubo cilíndrico ou
cônico. Encaixa, une, fixa e ajusta partes, componentes e acessórios -
como chaves (ou orifícios), embocadura, barrilete, campânula, palheta
simples ou dupla, entre outros -, para realizar a montagem do
instrumento. Ajusta os conjuntos de elementos. Realiza acabamento nos
instrumentos musicais, preparando superfícies; aplicando materiais e
produtos para revestimento; e fazendo polimento. Executa a confecção de
instrumentos de sopro de madeira de acordo com ordem de produção,
desenhos, protótipos, especificações e manuais técnicos. Utiliza, no
processo de produção, máquinas e equipamentos mecânicos e elétricos,
instrumentos e ferramentas. Pode utilizar máquinas e equipamentos
digitais ou automatizados. Aplica técnicas de afinação dos instrumentos
de sopro de madeira. Avalia a qualidade dos instrumentos de sopro de
madeira, aplicando testes sensoriais e visuais e usando instrumentos de
medição e equipamentos, para verificar e ajustar parâmetros
dimensionais, geométricos, estruturais, funcionais e acústicos. Pode
providenciar laudos técnicos. Pode tocar o instrumento para avaliar a
qualidade do som, além de identificar e corrigir defeitos. Coloca
etiquetas identificadoras e embala os instrumentos de sopro de madeira
produzidos. Aplica procedimentos de proteção contra agentes externos.
Pode realizar a venda de instrumentos de sopro de madeira, pesquisando
mercados e divulgando o produto. Calcula o valor de venda e estabelece o
mecanismo de cobrança. Após realizar a venda, seleciona os meios de
transporte para entrega. Analisa a satisfação do cliente. Pode realizar
conserto e restauração de instrumento de sopro de madeira, avaliando as
informações do cliente, identificando a viabilidade do serviço e orçando
os custos. Compra componentes e constrói peças de reposição. Desmonta o
instrumento e substitui partes e componentes desgastados ou defeituosos.
Após realização do serviço, monta o instrumento de sopro de madeira.
Prepara documentos técnicos, participando da elaboração de manuais
técnicos e preenchendo relatórios. Mantém o local de trabalho limpo e
organizado. Conserva ferramentas, acessórios e instrumentos de trabalho
limpos, organizados, acondicionados e em plenas condições de uso e
funcionamento. Providencia serviços de manutenção de equipamentos.
Controla desperdícios de materiais, realizando reaproveitamento.
Identifica resíduos e providencia sua destinação de acordo com normas
ambientais. Trabalha com segurança, utilizando equipamentos de proteção
individual e atuando na prevenção de acidentes.

\begin{Form}
\textbf{Esta correspondência está adequada?} \\ 
\ChoiceMenu[radio,radiosymbol=\ding{52},bordercolor=0 0.5 0,name=cbovalid_ 742125 ]{}{CBO deve ficar=sim, \hspace{6em} CBO não fica aqui=nao}
\end{Form}

\textbf{CBO: 742130  ( Confeccionador de instrumentos de sopro (metal) )}

Planeja a produção de instrumentos de sopro de metal, em série ou sob
encomenda, programando as etapas e avaliando o tempo para confecção.
Verifica a técnica a ser aplicada no processo de produção, confere os
materiais a serem usados e seleciona máquinas, equipamentos,
ferramentas, acessórios e instrumentos de medição, de acordo com ordem
de produção, especificações, fichas técnicas e manuais técnicos. Pode
calcular custos. Prepara os materiais a serem utilizados na confecção de
instrumentos de sopro de metal. Distingue acústicas dos materiais,
analisando as propriedades de ressonância e identificando os diferentes
timbres. Define os tipos de cortes dos materiais. Classifica metais e
faz seu recozimento. Confecciona componentes, construindo moldes;
manuseando gabaritos; utilizando ferramentas e instrumentos de medição;
regulando e operando máquinas; e colando materiais. Calandra metais,
executa lixamento primário em superfície, modela peças por meio de
repuxo e solda componentes metálicos. Confecciona instrumentos de sopro
de metal - tais como trompete, trombone, tuba, trompa e corneta de
pistões -, de acordo com ordem de produção, desenhos, protótipos,
especificações e manuais técnicos. Monta a estrutura (corpo do
instrumento) composta de tubo, com espessura e comprimento
predeterminados. Encaixa, une, fixa e ajusta as demais partes,
componentes e acessórios -- como boquilha, braçadeira, campânula,
chaves, entre outros -, para realizar a montagem do instrumento. Ajusta
os conjuntos de elementos. Realiza acabamento do instrumento de sopro de
metal, preparando superfícies, aplicando banhos eletrostáticos; pintando
ou fazendo aplicação de laca, níquel ou prata; jateando com areia as
superfícies; e fazendo polimento. Utiliza, no processo de produção,
máquinas e equipamentos mecânicos e elétricos, instrumentos e
ferramentas. Pode utilizar máquinas e equipamentos digitais ou
automatizados. Aplica técnicas de afinação dos instrumentos de sopro de
metal. Avalia a qualidade dos instrumentos de sopro de metal, aplicando
testes sensoriais e visuais e usando instrumentos de medição e
equipamentos, para verificar e ajustar parâmetros - dimensionais,
geométricos, estruturais, funcionais e acústicos -, tais como
nivelamento da campânula; regulagem e vedação das sapatilhas; entre
outros. Pode providenciar laudos técnicos. Pode tocar o instrumento para
avaliar a qualidade do som, além de identificar e corrigir defeitos.
Coloca etiquetas identificadoras e embala os instrumentos de sopro de
metal produzidos. Aplica procedimentos de proteção contra agentes
externos. Pode realizar a venda de instrumentos de sopro de metal,
pesquisando mercados e divulgando o produto. Calcula o valor de venda e
estabelece o mecanismo de cobrança. Após efetuar a venda, seleciona os
meios de transporte para entrega. Analisa a satisfação do cliente. Pode
realizar conserto e restauração de instrumento de sopro de metal,
avaliando as informações do cliente, identificando a viabilidade do
serviço e orçando os custos. Compra componentes e constrói peças de
reposição. Desmonta o instrumento, substitui partes e componentes
desgastados ou defeituosos; desamassa a superfície; e ajusta as partes.
Após realização do serviço, monta o instrumento de sopro de metal.
Prepara documentos técnicos, participando da elaboração de manuais
técnicos e preenchendo relatórios. Mantém o local de trabalho limpo e
organizado. Conserva ferramentas, acessórios e instrumentos de trabalho
limpos, organizados, acondicionados e em plenas condições de uso e
funcionamento. Providencia serviços de manutenção de equipamentos.
Controla desperdícios de materiais, realizando reaproveitamento.
Identifica resíduos e providencia sua destinação de acordo com normas
ambientais. Trabalha com segurança, utilizando equipamentos de proteção
individual e atuando na prevenção de acidentes.

\begin{Form}
\textbf{Esta correspondência está adequada?} \\ 
\ChoiceMenu[radio,radiosymbol=\ding{52},bordercolor=0 0.5 0,name=cbovalid_ 742130 ]{}{CBO deve ficar=sim, \hspace{6em} CBO não fica aqui=nao}
\end{Form}

\textbf{CBO: 742135  ( Confeccionador de órgão )}

Planeja a produção de órgãos musicais, em série ou sob encomenda,
programando as etapas e avaliando o tempo para confecção. Verifica a
técnica a ser aplicada no processo de produção, confere as
matérias-primas a serem usadas e seleciona máquinas, equipamentos,
ferramentas, acessórios e instrumentos de medição, de acordo com ordem
de produção, especificações, fichas técnicas e manuais técnicos. Pode
calcular custos. Prepara as matérias-primas a serem utilizadas na
confecção de órgãos musicais. Distingue acústicas dos materiais,
analisando as propriedades de ressonância e identificando os diferentes
timbres. Define os tipos de cortes dos materiais. Classifica madeira e
faz seu secamento em estufa. Tinge lâminas de madeira. Confecciona
componentes, construindo moldes; manuseando gabaritos; utilizando
ferramentas e instrumentos de medição; regulando e operando máquinas; e
colando materiais. Realiza operações referentes a processos mecânicos,
eletroeletrônicos, pneumáticos, de marcenaria e de soldagem. Aplica
princípios básicos de acústica e estética. Confecciona órgãos musicais,
de acordo com ordem de produção, desenhos, protótipos, especificações e
manuais técnicos. Monta a estrutura, composta de caixa e do console
(mesa de comando). Encaixa, une, fixa, acopla e ajusta partes,
componentes e acessórios - como tubos, teclado, pedais, ventilador,
fole, canais de vento, entre outros -, para a montagem do órgão.
Verifica diâmetro e comprimento das partes internas e externas. Ajusta
os conjuntos de elementos. Realiza acabamento dos órgãos musicais,
preparando superfícies; aplicando materiais e produtos para
revestimento; e colocando ou pintando elementos decorativos. Utiliza, no
processo de produção, máquinas e equipamentos mecânicos, eletrônicos e
elétricos, além de instrumentos e ferramentas. Pode utilizar máquinas e
equipamentos digitais ou automatizados. Aplica técnicas para afinação
dos órgãos musicais. Avalia a qualidade dos órgãos musicais, aplicando
testes sensoriais e visuais e usando instrumentos de medição e
equipamentos, para verificar e ajustar parâmetros dimensionais,
geométricos, estruturais, funcionais e acústicos. Pode providenciar
laudos técnicos. Pode tocar o órgão para avaliar a qualidade do som,
além de identificar e corrigir defeitos. Coloca etiquetas
identificadoras e aplica procedimentos de proteção contra agentes
externos. Expõe relação de procedimentos para conservação preventiva e
estratégias de contenção de processos degenerativos causados por fatores
ambientais, tais como temperatura, umidade, radiação luminosa e
poluentes atmosféricos. Embala os órgãos musicais produzidos. Pode
realizar a venda de órgãos musicais, pesquisando mercados e divulgando o
produto. Calcula o valor de venda e estabelece o mecanismo de cobrança.
Após efetuar a venda, seleciona os meios de transporte para entrega.
Analisa a satisfação do cliente. Pode realizar conserto e restauração de
órgão musical, avaliando as informações do cliente, identificando a
viabilidade do serviço e orçando os custos. Compra componentes e
constrói peças de reposição. Desmonta o órgão, substitui partes e
componentes desgastados ou defeituosos, e ajusta as partes. Após
realização do serviço, monta o órgão musical. Prepara documentos
técnicos, participando da elaboração de manuais técnicos e preenchendo
relatórios. Mantém o local de trabalho limpo e organizado. Conserva
ferramentas, acessórios e instrumentos de trabalho limpos, organizados,
acondicionados e em plenas condições de uso e funcionamento. Providencia
serviços de manutenção de equipamentos. Controla desperdícios de
materiais, realizando reaproveitamento. Identifica resíduos e
providencia sua destinação de acordo com normas ambientais. Trabalha com
segurança, utilizando equipamentos de proteção individual e atuando na
prevenção de acidentes.

\begin{Form}
\textbf{Esta correspondência está adequada?} \\ 
\ChoiceMenu[radio,radiosymbol=\ding{52},bordercolor=0 0.5 0,name=cbovalid_ 742135 ]{}{CBO deve ficar=sim, \hspace{6em} CBO não fica aqui=nao}
\end{Form}

\textbf{CBO: 742140  ( Confeccionador de piano )}

Planeja a produção de pianos acústicos verticais ou de cauda, em série
ou sob encomenda, programando as etapas e avaliando o tempo para
confecção. Verifica a técnica a ser aplicada; confere matérias-primas; e
seleciona máquinas, equipamentos, ferramentas, acessórios e instrumentos
de medição, de acordo com ordem de produção, especificações, fichas
técnicas e manuais técnicos. Pode calcular custos. Prepara as
matérias-primas a serem utilizadas na confecção dos pianos. Distingue
acústicas dos materiais, analisando as propriedades de ressonância e
identificando os diferentes timbres. Desdobra e classifica madeira. Seca
madeira em estufa. Prensa e tinge lâminas de madeira. Confecciona
componentes, construindo moldes; manuseando gabaritos; utilizando
ferramentas e instrumentos de medição; regulando e operando máquinas; e
colando materiais. Executa lixamento primário em superfície e solda
componentes metálicos. Confecciona piano, de acordo com ordem de
produção, desenhos, protótipos, especificações e manuais técnicos. Monta
a mesa de som e a estrutura de metal (corpo do instrumento musical).
Instala, encaixa, une, fixa e ajusta partes, componentes e acessórios,
como pedais; cordas de fios de aço revestidos com cobre; abafadores;
cravelha, entre outros. Monta o teclado, com teclas brancas, feitas de
madeira e cobertas com marfim ou material plástico, e teclas pretas,
revestidas geralmente de ébano. Monta o mecanismo de ação, com martelos
de madeira revestidos com placas de feltro ou similar, pinos metálicos e
outros materiais. Verifica fixação, aperto, diâmetro e comprimento das
partes internas e externas. Ajusta os conjuntos de elementos. Realiza
acabamento do piano, preparando superfícies e aplicando materiais e
produtos para revestimento. Utiliza, no processo de produção, máquinas e
equipamentos mecânicos e eletrônicos, além de instrumentos e
ferramentas. Pode utilizar máquinas e equipamentos digitais ou
automatizados. Aplica técnicas para afinação do piano, processo de
ajuste das cordas para ter a intensidade correta para produzir a escala
musical. Avalia a qualidade dos pianos, aplicando testes sensoriais e
visuais e usando instrumentos de medição e equipamentos, para verificar
e ajustar parâmetros dimensionais, geométricos, estruturais, funcionais
e acústicos. Pode providenciar laudos técnicos. Coloca etiquetas
identificadoras. Aplica procedimentos de proteção contra agentes
externos. Expõe relação de procedimentos para conservação do piano e
contenção de processos degenerativos causados por fatores ambientais,
como temperatura, umidade e radiação luminosa. Embala os pianos
produzidos. Pode realizar a venda de pianos, pesquisando mercados e
divulgando o produto. Calcula o valor de venda e estabelece o mecanismo
de cobrança. Após efetuar a venda, seleciona os meios de transporte para
entrega. Analisa a satisfação do cliente. Pode realizar conserto e
restauração de piano acústico, avaliando as informações do cliente,
identificando a viabilidade do serviço e orçando os custos. Compra
componentes e constrói peças de reposição. Desmonta o piano, substitui
partes e componentes desgastados ou defeituosos, e ajusta as partes.
Após realização do serviço, monta o piano. Pode confeccionar piano
digital - sem cordas e que gera o som de forma eletrônica -, executando
diferente processo de produção em comparação ao piano acústico. Prepara
documentos técnicos, participando da elaboração de manuais técnicos e
preenchendo relatórios. Mantém o local de trabalho limpo e organizado.
Conserva ferramentas, acessórios e instrumentos de trabalho limpos,
organizados, acondicionados e em plenas condições de uso e
funcionamento. Providencia serviços de manutenção de equipamentos.
Controla desperdícios de materiais, realizando reaproveitamento.
Identifica resíduos e providencia sua destinação de acordo com normas
ambientais. Trabalha com segurança, utilizando equipamentos de proteção
individual e atuando na prevenção de acidentes.

\begin{Form}
\textbf{Esta correspondência está adequada?} \\ 
\ChoiceMenu[radio,radiosymbol=\ding{52},bordercolor=0 0.5 0,name=cbovalid_ 742140 ]{}{CBO deve ficar=sim, \hspace{6em} CBO não fica aqui=nao}
\end{Form}

\newpage
\section{ 09022 -- Técnico em Modelagem do Vestuário }

\textbf{Perfil do Curso (CNCT)}

O Técnico em Modelagem do Vestuário será habilitado para:

Criar e desenvolver projetos de moda. Utilizar técnicas de modelagem bi
e tridimensionais. Elaborar desenhos e fichas técnicas. Representar
graficamente peças de vestuário planificadas. Utilizar ferramentas da
computação gráfica para moda. Supervisionar produção de peça-piloto e
produção em série. Avaliar a vestibilidade e a viabilidade técnica do
produto.

Para atuação como Técnico em Modelagem do Vestuário, são fundamentais:

Conhecimentos interdisciplinares relacionados aos processos de criação,
envolvendo pesquisa, idealização, planejamento, execução técnica,
fruição e recepção estética. Competências comunicacionais e
empreendedoras voltadas à proposição de projetos, ao coletivo, à gestão,
à solução de problemas e à resiliência, entre outras competências
socioemocionais.

\textbf{CBOs correspondidos e seus perfis (presentes no CNCT, e presentes na predição da IA)}

\textbf{CBO: 318810  ( Modelista de roupas )}

Planeja a confecção de moldes e o desenvolvimento de protótipos de
roupas, examinando a coleção do estilista; verificando tendências de
mercado; consultando boletins técnicos; avaliando estilos de design,
concepções e produtos de concorrentes; e coletando informações
tecnológicas em feiras e em desfiles de vestuário. Confecciona moldes de
roupas, interpretando modelos e desenhos; elaborando desenhos
planificados de roupas; traçando, cortando e codificando moldes. Utiliza
prancheta e softwares, tais como CAD-Desenho Assistido por Computador
(Computer Aided Design) e CAM--Manufatura Auxiliada por Computador
(Computer Aided Manufacturing). Marca referências em moldes e elabora
tabelas de medidas de roupas, para orientar a produção do vestuário.
Desenvolve protótipo de roupa, interpretando normas técnicas para
produção, elaborando ficha técnica, especificando aviamentos e
acessórios, ajustando moldes e gabaritos e determinando a posição de
ornamentos, detalhes e acessórios. Confecciona e testa peça-piloto
(protótipo). Requisita testes laboratoriais do protótipo. Realiza
ampliações e reduções das dimensões do produto por escala. Avalia
materiais para aquisição com vistas à produção de vestuário, verificando
a resistência dos materiais e analisando sua composição e acabamento.
Avalia a costurabilidade de tecidos e couros. Avalia a relação entre
custo e benefício de materiais. Interpreta resultados de testes
laboratoriais de materiais e elabora relatório.

\begin{Form}
\textbf{Esta correspondência está adequada?} \\ 
\ChoiceMenu[radio,radiosymbol=\ding{52},bordercolor=0 0.5 0,name=cbovalid_ 318810 ]{}{CBO deve ficar=sim, \hspace{6em} CBO não fica aqui=nao}
\end{Form}

\newpage
\section{ 09023 -- Técnico em Multimídia }

\textbf{Perfil do Curso (CNCT)}

O Técnico em Multimídia será habilitado para:

Desenvolver comunicação visual em meios eletrônicos, interfaces
interativas, publicações digitais, animações 2D e 3D, jogos eletrônicos,
web sites, web TV, TV digital e conteúdo audiovisual. Organizar e
preparar arquivos digitais para aplicações web e multimídia, animações e
games. Aplicar técnicas de tratamento de imagens estáticas e em
movimento que compõem estruturas de navegação em mídias digitais.
Executar atualização de páginas web e portais.

Para atuação como Técnico em Multimídia, são fundamentais:

Conhecimentos interdisciplinares relacionados aos processos de criação,
envolvendo pesquisa, idealização, planejamento, execução técnica,
fruição e recepção estética. Competências comunicativas e empreendedoras
voltadas à proposição de projetos, ao coletivo, à gestão, à solução de
problemas e à resiliência, entre outras competências socioemocionais.

\textbf{CBOs correspondidos e seus perfis (presentes no CNCT, e presentes na predição da IA)}

\textbf{CBO: 317105  ( Desenvolvedor web (técnico) )}

Desenvolve sistemas e aplicações para web, elaborando a interface
gráfica e aplicando critérios ergonômicos de navegação. Monta a
estrutura do banco de dados e codifica os programas e aplicativos,
aplicando sistemas de rotinas de segurança. Compila os programas. Testa
aplicativos e programas, elaborando casos de testes. Gera aplicativos
para instalação e gerenciamento dos sistemas. Documenta os sistemas e
aplicações e avalia o desempenho dos produtos desenvolvidos. Implanta
sistemas e aplicações para web, instalando os programas, adaptando o
conteúdo para mídias interativas e configurando os equipamentos em que a
aplicação será executada. Homologa os sistemas e aplicações
desenvolvidos, avaliando e validando os resultados da implantação.
Elabora material para capacitação de usuários. Avalia os objetivos e as
metas de projetos de sistemas e aplicações, tendo em vista os resultados
obtidos na fase de implantação. Publica o código final no servidor de
sistemas e aplicações. Realiza a manutenção de sistemas e aplicações
para web, atualizando informações gráficas, textuais e audiovisuais e
convertendo seus códigos fontes para outras linguagens ou plataformas.
Implementa alterações nos sistemas e aplicações e em suas estruturas de
armazenamento de dados, adequando-os aos sistemas e ambientes
operacionais. Monitora a performance dos sistemas e aplicações,
utilizando ferramentas de software, para obter subsídios à atividade de
manutenção. Atualiza a documentação dos sistemas e aplicações, em função
das intervenções de manutenção. Fornece suporte técnico para o cliente
interno. Colabora na seleção dos recursos de desenvolvimento,
especificando máquinas, ferramentas, acessórios e suprimentos. Participa
da seleção das linguagens de programação e metodologias. Seleciona
ferramentas de desenvolvimento. Participa da definição da nomenclatura
padrão. Colabora no planejamento das etapas e ações de trabalho,
participando das definições de atividades, tarefas e cronograma.
Participa de reuniões com a equipe de trabalho ou com o cliente.
Participa das definições de padronizações. Projeta sistemas e aplicações
para web, identificando a demanda do cliente, coletando dados e
desenvolvendo o leiaute de telas e relatórios. Elabora o pré-projeto e
os projetos conceitual, lógico, estrutural, físico e gráfico dos
sistemas e aplicações, definindo os requisitos de projeto. Participa das
definições dos critérios ergonômicos de navegação e da interface de
comunicação e interatividade. Dimensiona a vida útil dos sistemas e
aplicações, modela sua estrutura de banco de dados e dimensiona a
capacidade de armazenamento de dados dos sistemas.

\begin{Form}
\textbf{Esta correspondência está adequada?} \\ 
\ChoiceMenu[radio,radiosymbol=\ding{52},bordercolor=0 0.5 0,name=cbovalid_ 317105 ]{}{CBO deve ficar=sim, \hspace{6em} CBO não fica aqui=nao}
\end{Form}

\textbf{CBO: 317120  ( Desenvolvedor de multimídia )}

Desenvolve sistemas e aplicações de multimídia, elaborando a interface
gráfica e aplicando critérios ergonômicos de navegação. Monta a
estrutura do banco de dados e codifica os programas e aplicativos,
aplicando sistemas de rotinas de segurança. Compila os programas. Testa
aplicativos e programas, elaborando casos de testes. Gera aplicativos
para instalação e gerenciamento dos sistemas. Documenta sistemas e
aplicações e avalia o desempenho dos produtos desenvolvidos. Implanta
sistemas e aplicações de multimídia, instalando os programas, adaptando
o conteúdo para mídias interativas e configurando os equipamentos em que
a aplicação será executada. Homologa sistemas e aplicações
desenvolvidos, avaliando e validando os resultados da implantação.
Elabora material para capacitação de usuários. Avalia os objetivos e as
metas de projetos de sistemas e aplicações, tendo em vista os resultados
obtidos na fase de implantação. Publica o código final no servidor de
sistemas e aplicações. Realiza a manutenção de sistemas e aplicações de
multimídia, atualizando informações gráficas, textuais e audiovisuais.
Implementa alterações em sistemas e aplicações e em suas estruturas de
armazenamento de dados, adequando-os aos sistemas e ambientes
operacionais. Monitora a performance de sistemas e aplicações,
utilizando ferramentas de software, para obter subsídios à atividade de
manutenção. Atualiza a documentação de sistemas e aplicações, em função
das intervenções de manutenção. Fornece suporte técnico para o cliente
interno. Colabora na seleção dos recursos de desenvolvimento,
especificando máquinas, ferramentas, acessórios e suprimentos. Participa
da seleção das linguagens de programação e metodologias. Seleciona
ferramentas de desenvolvimento. Participa da definição da nomenclatura
padrão. Colabora no planejamento das etapas e ações de trabalho,
participando das definições de atividades, tarefas e cronograma.
Participa de reuniões com a equipe de trabalho ou com o cliente.
Participa das definições de padronizações. Projeta sistemas e aplicações
de multimídia, identificando a demanda do cliente, coletando dados e
desenvolvendo o leiaute de telas e relatórios. Elabora o pré-projeto e
os projetos conceitual, lógico, estrutural, físico e gráfico dos
sistemas e aplicações, definindo os requisitos de projeto. Participa das
definições dos critérios ergonômicos de navegação e da interface de
comunicação e interatividade. Dimensiona a vida útil dos sistemas e
aplicações.

\begin{Form}
\textbf{Esta correspondência está adequada?} \\ 
\ChoiceMenu[radio,radiosymbol=\ding{52},bordercolor=0 0.5 0,name=cbovalid_ 317120 ]{}{CBO deve ficar=sim, \hspace{6em} CBO não fica aqui=nao}
\end{Form}

\newpage
\section{ 09024 -- Técnico em Museologia }

\textbf{Perfil do Curso (CNCT)}

O Técnico em Museologia será habilitado para:

Promover a difusão dos bens culturais sob tutela de instituições
museológicas, culturais e afins. Organizar exposições de diferentes
naturezas e duração. Auxiliar e realizar planejamento e gerenciamento de
acervos e espaços expositivos.\\
Organizar e gerenciar bases de dados e arquivos de documentos de
natureza histórica, artística, literária, científica no âmbito da
cultura material e imaterial, bem como auxiliar e realizar pesquisas
específicas. Auxiliar e realizar gerenciamento de espaços expositivos,
bem como produção e montagem de exposições. Orientar a seleção de bens
culturais para fins de preservação e auxiliar nas ações preventivas e de
restauro. Auxiliar e realizar planejamento e gerenciamento de reserva
técnica. Desenvolver ações educativas para o público e oferecer produtos
culturais.

Para atuação como Técnico em Museologia, são fundamentais:

Conhecimentos interdisciplinares relacionados aos processos de criação,
envolvendo pesquisa, idealização, planejamento, execução técnica,
fruição e recepção estética. Competências comunicativas e empreendedoras
voltadas à proposição de projetos, ao coletivo, à gestão, à solução de
problemas e à resiliência, entre outras competências socioemocionais.

\textbf{CBOs correspondidos e seus perfis (presentes no CNCT, e presentes na predição da IA)}

\textbf{CBO: 371205  ( Colecionador de selos e moedas )}

Planeja o trabalho, programando as atividades. Seleciona materiais,
instrumentos e acessórios -- tais como filigranoscópio, lupas,
luminárias com lente de aumento, embalagens de acrílico para
acondicionar moedas e selos, entre outros - para uso nas atividades
programadas, verificando sua disponibilidade e realizando a aquisição
dos faltantes. Levanta a oferta de selos e moedas e realiza pesquisa
para avaliação da qualidade de cada peça. Considera, na pesquisa, o
valor estimado de mercado; autenticidade; limpeza; integridade e estado
de conservação -- conforme tabelas de classificação -; data de impressão
e de cunho; raridade; peso e diâmetro; entre outros aspectos. Leva em
conta, ainda, os contextos histórico, cultural e social que marcam o
momento da tiragem e da emissão. Adquire selos e moedas, examinando a
peça, consultando catálogos nacionais e internacionais e efetuando a
triagem. Acondiciona selos e moedas, utilizando pastas compartimentadas,
envelopes, fichários ou outras embalagens, para a organização, a
conservação e o transporte do acervo. Pode realizar publicação do acervo
com visualização de moedas e selos por realidade virtual. Conserva o
acervo, estabelecendo as normas para realização desse trabalho. Faz a
higienização e acondiciona as moedas e os selos nas embalagens.
Implementa ações de segurança das peças. Fiscaliza o estado do acervo.
Assessora outras instituições na conservação de acervo. Elabora catálogo
do acervo de selos e moedas, criando fichas de catalogação e
informatização. Numera, toma medidas e fotografa as peças. Produz
registro de localização de cada selo e moeda. Elabora laudo técnico de
catalogação. Auxilia na classificação e no tombamento do acervo. Avalia
as peças para efeitos de empréstimos. Utiliza, nas diversas atividades
de montagem e controle do acervo, métodos e recursos de informatização.
Realiza atividades administrativas - tais como elaboração de
correspondência, registro de empréstimos e compra de materiais e
acessórios -, para gestão, divulgação e registros do acervo. Atualiza e
arquiva documentos -- físicos e digitais - relativos às peças do acervo.
Organiza exposição de coleções, colaborando no planejamento da sua
logística e em sua programação visual. Seleciona as peças do acervo para
exposição. Monitora as condições físicas da exposição, escolhendo o
material para a montagem, confeccionando ``passe-partout''; adequando a
iluminação ao ambiente; construindo vitrines; e controlando as condições
ambientais. Supervisiona o translado do acervo. Coloca as peças nos
locais determinados para exposição, sinalizando o circuito. Identifica
as peças com etiquetas. Subsidia, com informações, a criação de
catálogos. Verifica textos elucidativos do acervo. Divulga o evento na
mídia e alimenta ``website''. Acompanha a realização da exposição,
executando a monitoria. Controla a desmontagem da exposição, conferindo
a saída de cada peça. Comercializa selos e moedas, efetuando cadastro de
clientes e ordenando estoques. Leiloa as peças, presencialmente ou por
meios eletrônicos. Organiza catálogos de selos e moedas para
participação em leilão ou feiras. Normatiza os valores das peças.
Realiza intercâmbios profissionais com colecionadores, filatelistas e
numismatas. Pesquisa documentações técnicas, consultando bibliografias e
catálogos. Mantém limpos e organizados o espaço de catalogação e
pesquisa e o ambiente reservado para conservação de selos e moedas. Zela
pela preservação, funcionamento e integridade dos recursos de trabalho.
Controla desperdícios, separando sobras de materiais para
reaproveitamento. Identifica resíduos, providenciando sua segregação e
destinação, de acordo com normas ambientais. Trabalha com segurança,
usando equipamentos de proteção individual, tais como máscaras e luvas
para manuseio de moedas e selos.

\begin{Form}
\textbf{Esta correspondência está adequada?} \\ 
\ChoiceMenu[radio,radiosymbol=\ding{52},bordercolor=0 0.5 0,name=cbovalid_ 371205 ]{}{CBO deve ficar=sim, \hspace{6em} CBO não fica aqui=nao}
\end{Form}

\textbf{CBO: 371210  ( Técnico em museologia )}

Planeja as atividades anuais, participando na elaboração de projetos de
exposições -- itinerantes, temporárias e permanentes - e programando
atividades de manutenção de acervos em museus, centros culturais e
instituições similares. Seleciona equipamentos, instrumentos e materiais
para as atividades. Participa da seleção de objetos e documentos,
realizando pesquisas e avaliando valor cultural, autenticidade,
raridade, integridade e conservação, para fins de incorporação ao
acervo. Cataloga o acervo, elaborando e atualizando fichas de
catalogação e informatização. Auxilia na classificação de peças. Numera
as peças e registra peso e medidas de objetos. Fotografa cada peça e
assinala sua localização no acervo. Auxilia no tombamento do acervo.
Produz laudo técnico de catalogação. Avalia peças para efetivar
empréstimos. Realiza ações de preservação e auxilia na tomada de
decisões sobre a conservação do acervo. Participa na elaboração de
normas sobre a conservação do acervo, monitorando seu cumprimento.
Fiscaliza o estado do acervo. Monitora os serviços de higienização e
acondicionamento de peças. Auxilia na restauração de obras e elabora
laudos técnicos de conservação e restauro. Implementa ações de segurança
e medidas de primeiros socorros às obras. Assessora outras instituições
na conservação do acervo. Participa no planejamento, na organização e no
gerenciamento de reserva técnica, ambiente - com rígido controle de
climatização e constante higienização -- onde ficam as peças do acervo
que não estão em exposição. Controla entrada e saída de peças e
inspeciona o estado de conservação do acervo, na reserva técnica.
Participa na produção e na montagem de exposição de coleções,
colaborando no planejamento da logística e na elaboração da programação
visual. Seleciona peças do acervo para exposição. Escolhe o material
adequado à montagem, acompanha a implantação do sistema de iluminação e
controla as condições ambientais. Supervisiona o translado do acervo e
cria área de reserva técnica de obras em trânsito. Monitora a montagem
das obras, identificando-as com etiquetas. Sinaliza o circuito da
exposição. Subsidia com informações a criação de catálogos e indica
textos elucidativos do acervo. Colabora nas atividades de manutenção da
exposição. Confere a saída das obras, durante a desmontagem da
exposição. Desenvolve ações educativas pré-determinadas. Executa
treinamento do pessoal de apoio, para manuseio de peças, e da equipe de
vigilância das ações. Orienta estagiários. Realiza monitoria, prestando
informações ao público. Pesquisa documentação técnica, realizando
intercâmbios profissionais e consultando bibliografia e catálogos.
Visita museus e exposições para consulta. Entrevista artistas e
familiares sobre obras. Organiza e gerencia bases de dados e arquivos de
documentos -- físicos e virtuais - de natureza histórica, artística,
literária, científica, no âmbito da cultura material e imaterial.
Coordena equipe de trabalho, participando de processo de seleção e de
administração de pessoal. Supervisiona o trabalho de equipe, avaliando
desempenho e levantando necessidade de treinamento. Realiza programas de
treinamento para a equipe. Executa atividades administrativas,
elaborando relatório das atividades e das ocorrências; apoiando na
redação dos termos de aquisição do acervo; e elaborando correspondência.
Registra empréstimos de obras e elabora estatísticas de visitação e
consultas. Preenche requisições para aquisições e solicitações de
serviços de manutenção de equipamentos e instrumentos. Presta conta da
administração de recursos financeiros de destinação específica. Mantém
ambientes de trabalho limpos e organizados. Zela pela preservação e
funcionamento dos recursos de trabalho. Controla desperdícios,
incentivando o reaproveitamento. Identifica resíduos, providenciando
destinação, de acordo com normas ambientais. Trabalha com segurança,
usando equipamentos de proteção individual, tais como máscaras e luvas
para manuseio de obras.

\begin{Form}
\textbf{Esta correspondência está adequada?} \\ 
\ChoiceMenu[radio,radiosymbol=\ding{52},bordercolor=0 0.5 0,name=cbovalid_ 371210 ]{}{CBO deve ficar=sim, \hspace{6em} CBO não fica aqui=nao}
\end{Form}

\newpage
\section{ 09030 -- Técnico em Produção de Áudio e Vídeo }

\textbf{Perfil do Curso (CNCT)}

O Técnico em Produção de Áudio e Vídeo será habilitado para:

Captar imagens e sons. Realizar ambientação e operação de equipamentos
por intermédio de recursos e linguagens. Investigar a utilização de
tecnologias de tratamento acústico, de imagem, luminosidade e animação.
Preparar material audiovisual. Elaborar fichas técnicas, mapas de
programação, distribuição, veiculação de produtos e serviços de
comunicação.

Para atuação como Técnico em Produção de Áudio e Vídeo, são
fundamentais:

Conhecimentos interdisciplinares relacionados aos processos de criação,
envolvendo pesquisa, idealização, planejamento, execução técnica,
fruição e recepção estética. Competências comunicacionais e
empreendedoras voltadas à proposição de projetos, ao coletivo, à gestão,
à solução de problemas e à resiliência, entre outras competências
socioemocionais.

\textbf{CBOs correspondidos e seus perfis (presentes no CNCT, e presentes na predição da IA)}

\textbf{CBO: 372115  ( Operador de câmera de televisão )}

Planeja o processo de captação de imagens em movimento, participando na
elaboração de cronograma das atividades. Interage com equipes de
produção televisiva, para levantar informações dos produtos audiovisuais
em que atuará. Seleciona equipamentos - analógicos e digitais -,
``softwares'' e materiais. Realiza leitura técnica e interpreta projeto
visual desenvolvido pelo diretor de fotografia para produto audiovisual,
analisando desde o conceito fotográfico até a estratégia de execução.
Verifica a tradução de conceito fotográfico em imagens a serem captadas.
Utiliza suporte para fazer movimentos de câmera, em testes fotográficos.
Aplica diferentes técnicas de captação com a câmera, de forma
complementar à iluminação, para verificação do comportamento de
materiais cenográficos e de figurino. Monta câmera e acessórios para
captação de imagem de acordo com as características da produção
audiovisual. Participa de ensaios prévios à captação de imagens. Prepara
e opera câmera para captação de imagens, posicionando o equipamento de
acordo com orientação do diretor de fotografia. Utiliza equipamentos
periféricos conforme possibilidades e necessidades de captação e da
narrativa audiovisual. Regula as funções e o menu da câmera. Realiza
enquadramentos e ajustes de foco. Busca captar nas cenas, além da ação,
as emoções dos personagens. Acompanha movimentos, aplicando técnicas de
operação de câmera e composição de imagem. Propõe soluções de captação
de imagem em produções externas, de acordo com pauta e roteiro. Monitora
sinal de vídeo, na câmera ou em equipamento externo, considerando a
fotografia e a qualidade do sinal. Examina a avaliação feita pelo
diretor de fotografia em relação aos resultados da captação de imagens.
Considera a análise feita -- relacionada a enquadramentos, proporções
espaciais e efeitos de luzes e cores -- como contribuição para promover
processo de melhoria contínua em sua atuação. Conserva os diferentes
tipos de câmera e os acessórios limpos, acondicionados e em plenas
condições de uso e funcionamento. Requisita serviço de manutenção
corretiva, ao detectar falhas ou defeitos durante o uso de câmera.
Mantém-se atualizado em relação ao desenvolvimento técnico e às
inovações na área de captação de imagens em movimento de produções
audiovisuais. Zela pela segurança, utilizando equipamentos de proteção
individual. Identifica situações de risco e colabora para sua
eliminação.

\begin{Form}
\textbf{Esta correspondência está adequada?} \\ 
\ChoiceMenu[radio,radiosymbol=\ding{52},bordercolor=0 0.5 0,name=cbovalid_ 372115 ]{}{CBO deve ficar=sim, \hspace{6em} CBO não fica aqui=nao}
\end{Form}

\textbf{CBO: 373105  ( Operador de mídia audiovisual )}

Prepara o trabalho de operação de mídia audiovisual, verificando a grade
de programação, interpretando roteiro e/ou ``script'' de programa e
acompanhando a instalação de equipamentos e acessórios. Liga e testa os
equipamentos, detectando problemas e acionando manutenção corretiva para
solucioná-los. Confere o funcionamento e a qualidade técnica de mesas
e/ou consoles de som analógico e digital, automatizados ou não. Prepara
ambientes internos e externos para sonorização, captação e gravação para
mídias audiovisuais, organizando o ambiente, selecionando os materiais
adequados e considerando os aspectos ligados à propagação ou ao
isolamento acústico de determinados locais. Faz tratamento do áudio,
selecionando trilhas sonoras e executando captação, mixagem e edição de
áudios, tais como música, vinheta, comercial e chamada promocional.
Nivela modulação. Realiza a gravação. Sonoriza vídeos, vinhetas e
comerciais. Edita programas, adequando formato do material para
veiculação. Arquiva conteúdo final. Manipula áudio e vídeo, configurando
equipamentos audiovisuais; organizando conteúdo de exibição, edição e
gravação; gerando e gravando conteúdo. Opera e monitora sistemas,
equipamentos e softwares para transmissão e apresentação de produções
artísticas, culturais e jornalísticas em mídias audiovisuais que
utilizam diferentes sistemas de transmissão e distribuição de sinais -
radiodifusão, cabo, satélite ou rede mundial de computadores. Verifica o
cumprimento de horários e a sequencialidade da disposição de conteúdos
na grade de transmissão e acompanha a adequada inserção de conteúdos
publicitários e propagandísticos. Executa transmissões em rede,
considerando a configuração dos equipamentos, os sinais, as referências
e a sincronia. Zela pela qualidade técnica, monitorando a qualidade do
sinal, do áudio e do vídeo e verificando os sinais de entrada e de
saída. Aciona, quando necessário, sistema de contingência e/ou
redundância. Interage com os demais profissionais técnicos e artísticos
das transmissões em sistemas de exibição. Instrui auxiliares de
iluminação, áudio e vídeo. Compila informações e documentação técnica.
Relata ocorrências. Auxilia na elaboração de pareceres e de relatórios
técnicos. Realiza manutenção de primeiro nível em equipamentos e
providencia serviços de manutenção corretiva para problemas mais
complexos.

\begin{Form}
\textbf{Esta correspondência está adequada?} \\ 
\ChoiceMenu[radio,radiosymbol=\ding{52},bordercolor=0 0.5 0,name=cbovalid_ 373105 ]{}{CBO deve ficar=sim, \hspace{6em} CBO não fica aqui=nao}
\end{Form}

\textbf{CBO: 373220  ( Supervisor técnico operacional de sistemas de televisão e produtores de vídeo )}

Planeja as atividades técnicas e operacionais relativas à televisão e à
produção audiovisual, participando da definição de metas e orçando
despesas operacionais. Participa da conceituação de projetos e elabora
indicadores de desempenho setorial. Participa de reuniões pré e
pós-demanda, negociando recursos técnico-operacionais e financeiros com
áreas pertinentes. Administra ciclo de vida dos equipamentos e sistemas,
programando revisões preventivas. Supervisiona atividades técnicas e
operacionais, controlando o atingimento de metas. Define a especificação
dos recursos técnicos. Verifica as condições de uso e de conservação dos
equipamentos. Providencia atendimento de solicitação de serviços de
suporte e manutenção corretiva, quando necessário. Solicita linhas de
transmissão de áudio, vídeo, dados e segmento espacial. Acompanha a
execução da programação, monitorando a qualidade do sinal. Apura as
ocorrências durante o processo de transmissão. Comanda acionamento de
sistemas de redundância, sempre que necessário. Monitora o tráfego e a
distribuição de mídias. Coordena avaliações técnicas e homologações de
novos sistemas. Presta suporte para realização de eventos, solicitando
roteiro de gravação e de transmissão para análise. Realiza visita
técnica ao local do evento e providencia recursos de logística. Contrata
fornecimento de energia elétrica e serviços de links dedicados para o
evento. Avalia a operabilidade no cenário. Supervisiona a interligação
de sinais de terceiros com a produção do evento e checa as canalizações
de sinais. Repassa informações sobre condições técnicas e operacionais
do evento para áreas afins. Encaminha solicitação de segurança pública
e/ou particular. Pesquisa tendências tecnológicas, interagindo com
fabricantes e fornecedores para conhecer as inovações na área. Alerta
sua chefia sobre a necessidade de substituição de equipamentos e outros
recursos obsoletos. Administra recursos humanos, participando da
definição do perfil da equipe, gerenciando composição e movimentação de
pessoal, solicitando abertura de vagas e controlando assiduidade.
Programa folgas e férias da equipe. Supervisiona a equipe de trabalho,
avaliando o desempenho e levantando necessidade de treinamento.
Viabiliza a capacitação e o desenvolvimento da equipe. Atua em parceria
com áreas afins. Consolida relatórios de ocorrências e redige
documentos. Zela pelo cumprimento das normas de segurança do trabalho e
da legislação.

\begin{Form}
\textbf{Esta correspondência está adequada?} \\ 
\ChoiceMenu[radio,radiosymbol=\ding{52},bordercolor=0 0.5 0,name=cbovalid_ 373220 ]{}{CBO deve ficar=sim, \hspace{6em} CBO não fica aqui=nao}
\end{Form}

\textbf{CBO: 374105  ( Técnico em gravação de áudio )}

Prepara local para gravação de áudio, em estúdio de gravação ou em
eventos ao vivo. Define necessidades técnicas do evento; examina
dimensões físicas do local; verifica condições de infraestrutura para
acesso e instalação de equipamentos; identifica parâmetros acústicos do
local; e faz a especificação de equipamentos, tais como microfones,
transdutores, cabos e conectores. Elabora documento -- rider técnico -,
relacionando as características técnicas dos equipamentos de
sonorização. Colhe informações sobre evento de sonorização e gravação.
Orça a gravação de áudio. Acondiciona os equipamentos e supervisiona o
seu transporte ao local da gravação. Limpa e monta equipamentos de áudio
e acessórios. Confere o sistema de energia elétrica e aterramento.
Instala transdutores e conecta sinais de áudio. Testa o funcionamento de
equipamentos e substitui os que apresentam defeitos. Confere o
funcionamento do sistema de áudio e checa os níveis de interferência em
sistema. Testa a gravação de vozes ou de algum instrumento musical com
microfones diferentes, na busca da melhor sonoridade. Orienta
microfonista e assistente, indicando o melhor microfone para determinada
voz ou determinado instrumento. Configura sistemas de sonorização e
avalia características - timbre, formato e tipo - da fonte sonora.
Seleciona e posiciona os transdutores eletroacústicos. Opera sistema de
sonorização, misturando e equilibrando sinais de fontes de áudio.
Analisa sinais de áudio. Examina o ambiente de gravação e escolhe o meio
de registro (mídia). Alinha sistemas de gravação e define padrões de
sincronismo. Prepara o sistema de monitoramento para gravação. Capta
sinais de áudio para sistema de gravação. Opera sistema de gravação,
gravando sinais em meio de registro (mídia). Monitora sinais gravados.
Mapeia pistas de gravação. Trata os registros sonoros, selecionando-os,
editando-os e equalizando-os. Controla dinâmica de registros sonoros.
Pode usar efeitos sonoros especiais e simulação de diferentes ambientes,
como um estádio ou um teatro. Gera master. Marca pontos específicos de
faixas, rolos e arquivos digitais. Avalia a qualidade do produto sonoro.
Realiza o arquivamento de conteúdo audiofônico. Pode utilizar
equipamentos multimídias, plataformas digitais e diferentes softwares de
captação, gravação, edição, tratamento e arquivamento de produto sonoro.
Pode preparar e operar equipamentos e softwares de captação e gravação
de áudio em diferentes mídias, conforme roteiro, gênero e formato
pré-estabelecidos, considerando especificidades ligadas aos canais de
transmissão e distribuição -- radiodifusão, cabo, satélite ou rede
mundial de computadores. Documenta as rotinas de trabalho, registrando
ocorrências da produção, relatando medidas de contingenciamento e
comunicando ajustes aos diferentes setores envolvidos nos processos de
trabalho. Auxilia na elaboração de pareceres e relatórios técnicos de
qualidade e segurança. Realiza manutenção de primeiro nível em
equipamentos, executando pequenos reparos. Providencia serviço de
manutenção de equipamentos, para solucionar problemas mais complexos.
Trabalha com segurança, providenciando a fixação adequada dos
equipamentos em paredes, tetos, treliças, pórticos, postes e outros
suportes; examinando a localização e o posicionamento de equipamentos,
de modo a não obstruir saídas e passagens de emergência; verificando a
disponibilidade de dispositivos de extinção de incêndio adequados aos
equipamentos utilizados no sistema de áudio; e prevenindo acidentes, de
modo geral.

\begin{Form}
\textbf{Esta correspondência está adequada?} \\ 
\ChoiceMenu[radio,radiosymbol=\ding{52},bordercolor=0 0.5 0,name=cbovalid_ 374105 ]{}{CBO deve ficar=sim, \hspace{6em} CBO não fica aqui=nao}
\end{Form}

\textbf{CBO: 374110  ( Técnico em instalação de equipamentos de áudio )}

Planeja a instalação de equipamentos de áudio, analisando projeto;
verificando as dimensões físicas do local do evento; examinando as
condições de infraestrutura para acesso e instalação de equipamentos; e
identificando os parâmetros acústicos do local. Estabelece as
necessidades técnicas, a partir das informações levantadas sobre o
evento de sonorização e gravação. Define ângulos de cobertura, níveis de
pressão sonora e resposta de frequência. Faz a especificação dos
equipamentos, tais como microfones, transdutores, cabos e conectores.
Elabora documento -- rider técnico -, relacionando as características
técnicas dos equipamentos de sonorização. Gera desenhos e diagramas de
instalação, incluindo atualizações referentes às alterações realizadas.
Orça o serviço de instalação de equipamentos para evento. Acondiciona
equipamentos, para transportá-los ao local do evento. Supervisiona o
transporte, o descarregamento e a limpeza dos equipamentos e acessórios.
Monta e testa os equipamentos de áudio e acessórios - tais como mesas
e/ou console de som analógico ou digital, caixas de som, monitores de
áudio, amplificadores, entre outros -, substituindo os que apresentam
defeitos. Executa o projeto de instalação de equipamentos, conferindo
sistema de energia elétrica e aterramento. Conecta sistema de caixas
acústicas, instala transdutores e conecta sinais de áudio. Confere o
funcionamento de sistema e a uniformidade de distribuição de áudio.
Checa os níveis de interferência em sistema. Configura e alinha sistema
de sonorização. Avalia as características - timbre, formato e tipo - de
fonte sonora. Seleciona e posiciona transdutores eletroacústicos. No
sistema de gravação, ajusta as estruturas de ganho de sistema e prepara
o sistema de monitoramento para gravação. Pode instalar, testar e
configurar equipamentos e softwares de captação e gravação de conteúdo
audiofônico para diversas mídias, considerando especificidades ligadas
aos canais de transmissão e distribuição -- radiodifusão, cabo, satélite
e rede mundial de computadores -- e respeitando o projeto previamente
estabelecido. Desconecta sistema de caixas acústicas, desinstala
transdutores e desconecta sinais de áudio, no final do evento. Prepara
os equipamentos para retirada do local do evento e para transporte.
Documenta as rotinas de trabalho, registrando ocorrências da produção,
relatando medidas de contingenciamento e comunicando ajustes aos
diferentes setores envolvidos nos processos de trabalho. Informa aos
profissionais, envolvidos em sonorização e gravação de áudio, sobre os
processos realizados de instalação, teste e configuração de
equipamentos. Auxilia na elaboração de pareceres e relatórios técnicos
de qualidade e segurança. Realiza manutenção de primeiro nível em
equipamentos e acessórios, executando pequenos reparos. Providencia
serviço de manutenção de equipamentos, para solucionar problemas mais
complexos. Trabalha com segurança, realizando o dimensionamento adequado
da alimentação elétrica; providenciando a fixação segura dos
equipamentos em paredes, tetos, treliças, pórticos e outros suportes;
examinando a localização e o posicionamento de equipamentos, de modo a
não obstruir saídas e passagens de emergência; verificando a
disponibilidade de dispositivos de extinção de incêndio adequados aos
equipamentos utilizados no sistema de áudio; e prevenindo acidentes, de
modo geral.

\begin{Form}
\textbf{Esta correspondência está adequada?} \\ 
\ChoiceMenu[radio,radiosymbol=\ding{52},bordercolor=0 0.5 0,name=cbovalid_ 374110 ]{}{CBO deve ficar=sim, \hspace{6em} CBO não fica aqui=nao}
\end{Form}

\textbf{CBO: 374135  ( Projetista de sistemas de áudio )}

Elabora projeto de sistemas de sonorização e de gravação, levantando
informações para definir necessidades técnicas. Verifica dimensões
físicas e identifica parâmetros acústicos do local. Examina as condições
de infraestrutura para acesso e instalação de equipamentos. Define
ângulos de cobertura, níveis de pressão sonora e resposta de frequência.
Faz a especificação de equipamentos, microfones, transdutores, cabos e
conectores. Define a demanda de energia elétrica dos sistemas. Elabora
documento -- rider técnico -, relacionando as características técnicas
dos equipamentos de sonorização. Gera desenhos e diagramas de
instalação, incluindo atualizações referentes às alterações realizadas.
Orça o serviço de elaboração e implementação do projeto. Acompanha a
implementação do projeto, conferindo a uniformidade de distribuição de
áudio após a instalação dos equipamentos e acessórios. Configura e
alinha os sistemas de sonorização. Avalia as características - timbre,
formato e tipo - de fonte sonora e posiciona os transdutores
eletroacústicos. Analisa os sinais de áudio e, se adequados, os
distribui para outros sistemas. Ajusta as estruturas de ganho de sistema
e define padrões de sincronismo, no sistema de gravação. Monitora
sistemas, equipamentos e softwares de sonorização, captação e gravação
de sinais sonoros para diversas mídias e/ou plataformas. Avalia o
desempenho dos sistemas, realizando ajustes, quando necessário.
Documenta o serviço realizado, organizando informações, registrando
ocorrências e relatando interação com os setores envolvidos nos
processos. Elabora relatório sobre a implementação do projeto e sobre os
resultados obtidos. Realiza manutenção de primeiro nível em equipamentos
e acessórios, executando pequenos reparos. Providencia serviço de
manutenção de equipamentos, para solucionar problemas mais complexos.
Trabalha com segurança, prevenindo acidentes.

\begin{Form}
\textbf{Esta correspondência está adequada?} \\ 
\ChoiceMenu[radio,radiosymbol=\ding{52},bordercolor=0 0.5 0,name=cbovalid_ 374135 ]{}{CBO deve ficar=sim, \hspace{6em} CBO não fica aqui=nao}
\end{Form}

\textbf{CBO: 374210  ( Maquinista de cinema e vídeo )}

Planeja atividades, analisando o projeto de cenografia e examinando as
adaptações feitas por cenotécnico para adequação do projeto ao espaço
cênico. Estabelece as etapas do seu trabalho. Seleciona equipamentos
analógicos, automatizados e/ou robotizados. Confere materiais
necessários e requisita os itens faltantes. Pode orçar serviços. Atua na
preparação de equipamentos para filmagem cinematográfica e gravação de
vídeo, realizando instalação de suporte para deslocamento e movimentação
de câmeras; prestando apoio a eletricistas na montagem dos parques de
luz e em locação adaptada; operando varas elétricas (guinchos); armando
telas difusoras; e confeccionando mecanismos de efeitos especiais. Monta
cenário e suas peças, posicionando o cenário em espaço cênico; fixando
cenário móvel em maquinaria; e afixando adereços e objetos. Executa
trabalhos de cordame. Monta efeitos especiais. Compatibiliza cenário e
adereços à luz. Afina vestimentas de palcos. Confere urdimento e
condições físicas de ``set'' de filmagem. Examina o cenário montado,
verificando a disposição e o funcionamento de seus componentes, para
certificar-se de sua qualidade e sua segurança. Submete o cenário
finalizado à apreciação de cenotécnico e da direção da produção de filme
ou vídeo, para realizar eventuais ajustes ou melhorias. Durante ensaios,
gravação e filmagem, opera equipamentos - como gruas, carrinhos sobre
trilhos, entre outros - que deslocam e movimentam as câmeras. Manobra
rebatedores, opera cortinas, movimenta cenário sobre rodas, opera varas
cenográficas e opera efeitos especiais, de acordo com roteiro
predefinido de filmagem ou gravação. Pode auxiliar na elaboração de
relatórios técnicos referentes à montagem de cenário e às demais
atividades, relatando ocorrências durante os trabalhos. Recolhe
materiais - após a filmagem ou gravação de vídeo -, organizando-os para
guarda em depósito. Conserva equipamentos, ferramentas, instrumentos e
acessórios limpos, acondicionados e em plenas condições de uso e
funcionamento. Requisita serviço de manutenção corretiva de
equipamentos, quando necessário. Zela pela segurança, utilizando
equipamentos de proteção individual e atuando na prevenção de acidentes.
Identifica situações de risco e colabora para sua eliminação. Combate
princípios de incêndio. Pode prestar primeiros socorros.

\begin{Form}
\textbf{Esta correspondência está adequada?} \\ 
\ChoiceMenu[radio,radiosymbol=\ding{52},bordercolor=0 0.5 0,name=cbovalid_ 374210 ]{}{CBO deve ficar=sim, \hspace{6em} CBO não fica aqui=nao}
\end{Form}

\textbf{CBO: 374405  ( Editor de mídia audiovisual )}

Participa na definição do produto, interagindo com os demais
profissionais envolvidos na produção audiovisual, discutindo a pauta
proposta e sugerindo definição de narrativa. Toma conhecimento do
material bruto disponível -- tais como música, foto, vídeo e som -,
analisando sua relação com o roteiro preestabelecido. Requisita imagens
disponíveis em arquivo e solicita captação de novas imagens. Organiza os
materiais audiovisuais, elaborando índice de conteúdo gravado. Estrutura
narrativas audiovisuais, conforme roteiro, selecionando imagens e sons
diretos. Ordena narrativas de mídias audiovisuais. Modula tempos
narrativos. Define os efeitos visuais e o corte final de mídias
audiovisuais. Cria efeitos especiais e aplica aqueles predefinidos em
softwares. Cria e adiciona artes gráficas. Corrige cores, brilho e
contraste. Faz adequação dos formatos de artes gráficas. Aplica filtros.
Analisa a qualidade de imagem e som. Realiza o tratamento de imagens,
quanto a contraste, brilho, saturação, tons, cores e nitidez. Executa o
tratamento de sons, em relação a frequência, modulação, compressão,
equalização, efeitos, amplitude, fase, timbre, altura e intensidade.
Edita imagens e áudio, conforme roteiro preestabelecido, montando mídias
audiovisuais em sistema de edição. Corta e funde imagens. Corta e faz a
mixagem de áudio. Sincroniza e faz o sequenciamento de sons e imagens.
Confecciona e insere créditos. Revisa a edição final. Elabora e registra
a lista de decisões de edição. Exporta mídias audiovisuais. Envia a
mídia gravada para o setor de qualidade. Finaliza o material
audiovisual, encaminhando-o aos setores de pós-produção. Nas etapas de
trabalho, opera equipamentos e softwares de captação, gravação,
tratamento, edição, finalização e arquivamento de conteúdo audiovisual
para diferentes plataformas e meios de comunicação, considerando
especificidades relacionadas aos canais de transmissão e distribuição,
tais como radiodifusão, cabo, satélite ou rede mundial de computadores.
Segue rotinas de trabalho, registrando ocorrências e medidas de
contingenciamento e comunicando intercorrências e ajustes aos diferentes
setores envolvidos nos processos de trabalho. Participa na elaboração de
pareceres e relatórios técnicos de qualidade e segurança. Compila
informações e documentação técnica e mantém registros de produção.
Realiza manutenção de primeiro nível em equipamentos. Requisita serviços
de manutenção para reparos mais complexos.

\begin{Form}
\textbf{Esta correspondência está adequada?} \\ 
\ChoiceMenu[radio,radiosymbol=\ding{52},bordercolor=0 0.5 0,name=cbovalid_ 374405 ]{}{CBO deve ficar=sim, \hspace{6em} CBO não fica aqui=nao}
\end{Form}

\textbf{CBO: 374415  ( Finalizador de vídeo )}

Participa em reuniões sobre o produto audiovisual, tomando conhecimento
do roteiro, do gênero e do formato preestabelecidos e interagindo com os
demais profissionais envolvidos na produção e na edição do vídeo.
Processa ajustes finos de imagem e som, revisando a edição final.
Captura, faz o sequenciamento, corta e funde imagens. Examina a
qualidade final da imagem e do som e verifica o sincronismo do som com
imagem. Exporta mídias audiovisuais. Atua na definição, na criação e no
aprimoramento de efeitos especiais. Lista as cenas para aplicação de
efeitos visuais, orientando sua filmagem. Corrige cores, brilho e
contraste. Distorce, compõe, recorta, multiplica, inverte, congela,
acelera e desacelera imagens, em função do efeito especial. Aplica
filtros. Insere ilustrações, logotipos, gráficos e animações. Cria e
aplica caracteres. Cria, adiciona e adequa formatos de artes gráficas.
Confecciona e insere créditos, inclusive usando barra de créditos
animadas. Opera equipamentos e softwares no processo de finalização de
conteúdo audiovisual, adequando o vídeo a diferentes plataformas e meios
de comunicação e considerando especificidades ligadas aos canais de
transmissão e distribuição, tais como radiodifusão, cabo, satélite ou
rede mundial de computadores. Segue rotinas de trabalho, registrando
ocorrências e medidas de contingenciamento e comunicando intercorrências
e ajustes aos diferentes setores envolvidos nos processos de trabalho.
Participa na elaboração de pareceres e relatórios técnicos de qualidade
e segurança. Compila informações e documentação técnica e mantém
registros da etapa de pós-produção. Realiza manutenção de primeiro nível
em equipamentos. Requisita serviços de manutenção para reparos mais
complexos.

\begin{Form}
\textbf{Esta correspondência está adequada?} \\ 
\ChoiceMenu[radio,radiosymbol=\ding{52},bordercolor=0 0.5 0,name=cbovalid_ 374415 ]{}{CBO deve ficar=sim, \hspace{6em} CBO não fica aqui=nao}
\end{Form}

\textbf{CBO: 766120  ( Editor de texto e imagem )}

Planeja o serviço de editoração de texto e imagem, analisando a ordem de
serviço. Requisita materiais, seleciona softwares e verifica as
condições dos equipamentos. Faz o orçamento do serviço. Elabora projeto,
para a execução da editoração de acordo com a solicitação do cliente.
Efetua a editoração dos textos, selecionando tipos, serifas e tamanho de
fonte, de acordo com o projeto visual. Digita e posiciona o texto.
Diagrama, ilustra e formata os textos. Executa a editoração de imagens
para a confecção de produtos gráficos, operando escâneres, câmeras e
equipamentos e utilizando catálogo de cores. Organiza material e
ilustrações. Realiza cópias em escala, para reduções e ampliações.
Analisa imagens e fotografias para verificar imperfeições. Revisa e
executa o controle de qualidade de textos e imagens, fazendo o
gerenciamento de cores, utilizando equipamentos de medição de cor - como
densitômetros e colorímetros -, efetuando a separação de cores dos
originais e o fechamento dos arquivos, de acordo com o projeto gráfico.
Avalia se o produto está de acordo com as especificações da solicitação
do cliente e se atende às características do processo de impressão.
Confecciona boneco - com miolo e capa - e prova digital, para a
aprovação do projeto pelo cliente. Organiza a documentação referente ao
projeto, arquivando - em meio digital - a solicitação inicial; a cópia
do serviço de editoração de texto e imagem realizado; e o documento de
aprovação do projeto pelo cliente. Executa a limpeza e a conservação dos
equipamentos. Mantém o local de trabalho limpo e organizado. Faz o
descarte de resíduos de acordo com as normas ambientais.

\begin{Form}
\textbf{Esta correspondência está adequada?} \\ 
\ChoiceMenu[radio,radiosymbol=\ding{52},bordercolor=0 0.5 0,name=cbovalid_ 766120 ]{}{CBO deve ficar=sim, \hspace{6em} CBO não fica aqui=nao}
\end{Form}

\newpage
\section{ 09033 -- Técnico em Rádio e Televisão }

\textbf{Perfil do Curso (CNCT)}

O Técnico em Rádio e Televisão será habilitado para:

Executar a produção e veiculação de programas radiofônicos e
televisivos. Realizar seleção musical, montagem de filmes, videotapes,
trilhas, vinhetas, jingles, spots, roteiros, locução e aplicação de
efeitos especiais. Operar equipamentos analógicos e digitais de estúdio
de gravação.

Para atuação como Técnico em Rádio e Televisão, são fundamentais:

Conhecimentos interdisciplinares relacionados aos processos de criação,
envolvendo pesquisa, idealização, planejamento, execução técnica,
fruição e recepção estética. Competências comunicativas e empreendedoras
voltadas à proposição de projetos, ao coletivo, à gestão, à solução de
problemas e à resiliência, entre outras competências socioemocionais.

\textbf{CBOs correspondidos e seus perfis (presentes no CNCT, e presentes na predição da IA)}

\textbf{CBO: 373105  ( Operador de mídia audiovisual )}

Prepara o trabalho de operação de mídia audiovisual, verificando a grade
de programação, interpretando roteiro e/ou ``script'' de programa e
acompanhando a instalação de equipamentos e acessórios. Liga e testa os
equipamentos, detectando problemas e acionando manutenção corretiva para
solucioná-los. Confere o funcionamento e a qualidade técnica de mesas
e/ou consoles de som analógico e digital, automatizados ou não. Prepara
ambientes internos e externos para sonorização, captação e gravação para
mídias audiovisuais, organizando o ambiente, selecionando os materiais
adequados e considerando os aspectos ligados à propagação ou ao
isolamento acústico de determinados locais. Faz tratamento do áudio,
selecionando trilhas sonoras e executando captação, mixagem e edição de
áudios, tais como música, vinheta, comercial e chamada promocional.
Nivela modulação. Realiza a gravação. Sonoriza vídeos, vinhetas e
comerciais. Edita programas, adequando formato do material para
veiculação. Arquiva conteúdo final. Manipula áudio e vídeo, configurando
equipamentos audiovisuais; organizando conteúdo de exibição, edição e
gravação; gerando e gravando conteúdo. Opera e monitora sistemas,
equipamentos e softwares para transmissão e apresentação de produções
artísticas, culturais e jornalísticas em mídias audiovisuais que
utilizam diferentes sistemas de transmissão e distribuição de sinais -
radiodifusão, cabo, satélite ou rede mundial de computadores. Verifica o
cumprimento de horários e a sequencialidade da disposição de conteúdos
na grade de transmissão e acompanha a adequada inserção de conteúdos
publicitários e propagandísticos. Executa transmissões em rede,
considerando a configuração dos equipamentos, os sinais, as referências
e a sincronia. Zela pela qualidade técnica, monitorando a qualidade do
sinal, do áudio e do vídeo e verificando os sinais de entrada e de
saída. Aciona, quando necessário, sistema de contingência e/ou
redundância. Interage com os demais profissionais técnicos e artísticos
das transmissões em sistemas de exibição. Instrui auxiliares de
iluminação, áudio e vídeo. Compila informações e documentação técnica.
Relata ocorrências. Auxilia na elaboração de pareceres e de relatórios
técnicos. Realiza manutenção de primeiro nível em equipamentos e
providencia serviços de manutenção corretiva para problemas mais
complexos.

\begin{Form}
\textbf{Esta correspondência está adequada?} \\ 
\ChoiceMenu[radio,radiosymbol=\ding{52},bordercolor=0 0.5 0,name=cbovalid_ 373105 ]{}{CBO deve ficar=sim, \hspace{6em} CBO não fica aqui=nao}
\end{Form}

\textbf{CBO: 373220  ( Supervisor técnico operacional de sistemas de televisão e produtores de vídeo )}

Planeja as atividades técnicas e operacionais relativas à televisão e à
produção audiovisual, participando da definição de metas e orçando
despesas operacionais. Participa da conceituação de projetos e elabora
indicadores de desempenho setorial. Participa de reuniões pré e
pós-demanda, negociando recursos técnico-operacionais e financeiros com
áreas pertinentes. Administra ciclo de vida dos equipamentos e sistemas,
programando revisões preventivas. Supervisiona atividades técnicas e
operacionais, controlando o atingimento de metas. Define a especificação
dos recursos técnicos. Verifica as condições de uso e de conservação dos
equipamentos. Providencia atendimento de solicitação de serviços de
suporte e manutenção corretiva, quando necessário. Solicita linhas de
transmissão de áudio, vídeo, dados e segmento espacial. Acompanha a
execução da programação, monitorando a qualidade do sinal. Apura as
ocorrências durante o processo de transmissão. Comanda acionamento de
sistemas de redundância, sempre que necessário. Monitora o tráfego e a
distribuição de mídias. Coordena avaliações técnicas e homologações de
novos sistemas. Presta suporte para realização de eventos, solicitando
roteiro de gravação e de transmissão para análise. Realiza visita
técnica ao local do evento e providencia recursos de logística. Contrata
fornecimento de energia elétrica e serviços de links dedicados para o
evento. Avalia a operabilidade no cenário. Supervisiona a interligação
de sinais de terceiros com a produção do evento e checa as canalizações
de sinais. Repassa informações sobre condições técnicas e operacionais
do evento para áreas afins. Encaminha solicitação de segurança pública
e/ou particular. Pesquisa tendências tecnológicas, interagindo com
fabricantes e fornecedores para conhecer as inovações na área. Alerta
sua chefia sobre a necessidade de substituição de equipamentos e outros
recursos obsoletos. Administra recursos humanos, participando da
definição do perfil da equipe, gerenciando composição e movimentação de
pessoal, solicitando abertura de vagas e controlando assiduidade.
Programa folgas e férias da equipe. Supervisiona a equipe de trabalho,
avaliando o desempenho e levantando necessidade de treinamento.
Viabiliza a capacitação e o desenvolvimento da equipe. Atua em parceria
com áreas afins. Consolida relatórios de ocorrências e redige
documentos. Zela pelo cumprimento das normas de segurança do trabalho e
da legislação.

\begin{Form}
\textbf{Esta correspondência está adequada?} \\ 
\ChoiceMenu[radio,radiosymbol=\ding{52},bordercolor=0 0.5 0,name=cbovalid_ 373220 ]{}{CBO deve ficar=sim, \hspace{6em} CBO não fica aqui=nao}
\end{Form}

\textbf{CBO: 374105  ( Técnico em gravação de áudio )}

Prepara local para gravação de áudio, em estúdio de gravação ou em
eventos ao vivo. Define necessidades técnicas do evento; examina
dimensões físicas do local; verifica condições de infraestrutura para
acesso e instalação de equipamentos; identifica parâmetros acústicos do
local; e faz a especificação de equipamentos, tais como microfones,
transdutores, cabos e conectores. Elabora documento -- rider técnico -,
relacionando as características técnicas dos equipamentos de
sonorização. Colhe informações sobre evento de sonorização e gravação.
Orça a gravação de áudio. Acondiciona os equipamentos e supervisiona o
seu transporte ao local da gravação. Limpa e monta equipamentos de áudio
e acessórios. Confere o sistema de energia elétrica e aterramento.
Instala transdutores e conecta sinais de áudio. Testa o funcionamento de
equipamentos e substitui os que apresentam defeitos. Confere o
funcionamento do sistema de áudio e checa os níveis de interferência em
sistema. Testa a gravação de vozes ou de algum instrumento musical com
microfones diferentes, na busca da melhor sonoridade. Orienta
microfonista e assistente, indicando o melhor microfone para determinada
voz ou determinado instrumento. Configura sistemas de sonorização e
avalia características - timbre, formato e tipo - da fonte sonora.
Seleciona e posiciona os transdutores eletroacústicos. Opera sistema de
sonorização, misturando e equilibrando sinais de fontes de áudio.
Analisa sinais de áudio. Examina o ambiente de gravação e escolhe o meio
de registro (mídia). Alinha sistemas de gravação e define padrões de
sincronismo. Prepara o sistema de monitoramento para gravação. Capta
sinais de áudio para sistema de gravação. Opera sistema de gravação,
gravando sinais em meio de registro (mídia). Monitora sinais gravados.
Mapeia pistas de gravação. Trata os registros sonoros, selecionando-os,
editando-os e equalizando-os. Controla dinâmica de registros sonoros.
Pode usar efeitos sonoros especiais e simulação de diferentes ambientes,
como um estádio ou um teatro. Gera master. Marca pontos específicos de
faixas, rolos e arquivos digitais. Avalia a qualidade do produto sonoro.
Realiza o arquivamento de conteúdo audiofônico. Pode utilizar
equipamentos multimídias, plataformas digitais e diferentes softwares de
captação, gravação, edição, tratamento e arquivamento de produto sonoro.
Pode preparar e operar equipamentos e softwares de captação e gravação
de áudio em diferentes mídias, conforme roteiro, gênero e formato
pré-estabelecidos, considerando especificidades ligadas aos canais de
transmissão e distribuição -- radiodifusão, cabo, satélite ou rede
mundial de computadores. Documenta as rotinas de trabalho, registrando
ocorrências da produção, relatando medidas de contingenciamento e
comunicando ajustes aos diferentes setores envolvidos nos processos de
trabalho. Auxilia na elaboração de pareceres e relatórios técnicos de
qualidade e segurança. Realiza manutenção de primeiro nível em
equipamentos, executando pequenos reparos. Providencia serviço de
manutenção de equipamentos, para solucionar problemas mais complexos.
Trabalha com segurança, providenciando a fixação adequada dos
equipamentos em paredes, tetos, treliças, pórticos, postes e outros
suportes; examinando a localização e o posicionamento de equipamentos,
de modo a não obstruir saídas e passagens de emergência; verificando a
disponibilidade de dispositivos de extinção de incêndio adequados aos
equipamentos utilizados no sistema de áudio; e prevenindo acidentes, de
modo geral.

\begin{Form}
\textbf{Esta correspondência está adequada?} \\ 
\ChoiceMenu[radio,radiosymbol=\ding{52},bordercolor=0 0.5 0,name=cbovalid_ 374105 ]{}{CBO deve ficar=sim, \hspace{6em} CBO não fica aqui=nao}
\end{Form}

\textbf{CBO: 374125  ( Técnico em sonorização )}

Planeja o sistema de sonorização, levando em consideração as condições
do estúdio ou do local do evento. Verifica as dimensões físicas do
espaço e examina as condições de infraestrutura para acesso e instalação
de equipamentos. Identifica os parâmetros acústicos do local. Estabelece
a demanda de energia elétrica do sistema. Define necessidades técnicas
do evento, ângulos de cobertura, níveis de pressão sonora e resposta de
frequência. Faz a especificação de equipamentos, microfones,
transdutores, cabos e conectores. Elabora documento -- rider técnico -,
relacionando as características técnicas dos equipamentos de
sonorização. Orça o serviço de sonorização de evento. Acondiciona os
equipamentos de áudio e supervisiona o seu transporte ao estúdio ou
local do evento. Limpa e monta equipamentos de áudio e acessórios.
Realiza a instalação de equipamentos, conferindo o sistema de energia
elétrica e aterramento, conectando sistema de caixas acústicas,
instalando transdutores e conectando sinais de áudio. Testa o
funcionamento de equipamentos e substitui os que apresentam defeitos.
Confere o funcionamento do sistema de áudio e checa os níveis de
interferência em sistema. Confere a uniformidade de distribuição de
áudio. Configura e alinha sistemas de sonorização. Avalia as
características - timbre, formato e tipo - de fonte sonora. Seleciona e
posiciona os transdutores eletroacústicos. Analisa os sinais de áudio.
Equilibra e mistura sinais de fontes de áudio. Distribui sinais de áudio
para outros sistemas. Seleciona e opera instrumentos de monitoramento de
mau funcionamento dos sistemas e equipamentos, acionando, quando
necessário, sistemas de contingência e/ou redundância. Pode selecionar e
operar softwares e ferramentas digitais, que proporcionam automação no
processo de sonorização, em diferentes plataformas e/ou eventos. Ao
final da prestação de serviços de sonorização, desconecta sistema de
caixas acústicas, desinstala transdutores e desconecta sinais de áudio.
Prepara os equipamentos para retirada do estúdio ou local do evento e
para transporte. Documenta o serviço realizado, organizando informações,
registrando ocorrências e relatando interação com profissionais
envolvidos no evento, como os responsáveis por iluminação e cenografia.
Elabora relatório sobre a implementação do projeto de sonorização e
sobre os resultados obtidos. Realiza manutenção de primeiro nível em
equipamentos e acessórios, executando pequenos reparos. Providencia
serviço de manutenção de equipamentos, para solucionar problemas mais
complexos. Trabalha com segurança, realizando o dimensionamento adequado
da alimentação elétrica; providenciando a fixação segura dos
equipamentos em paredes, tetos, treliças, pórticos e outros suportes;
examinando a localização e o posicionamento de equipamentos, de modo a
não obstruir saídas e passagens de emergência; verificando a
disponibilidade de dispositivos de extinção de incêndio adequados aos
equipamentos utilizados no sistema de áudio; e prevenindo acidentes, de
modo geral.

\begin{Form}
\textbf{Esta correspondência está adequada?} \\ 
\ChoiceMenu[radio,radiosymbol=\ding{52},bordercolor=0 0.5 0,name=cbovalid_ 374125 ]{}{CBO deve ficar=sim, \hspace{6em} CBO não fica aqui=nao}
\end{Form}

\textbf{CBO: 374130  ( Técnico em mixagem de áudio )}

Planeja o serviço de mixagem, programando as etapas a serem realizadas e
estabelecendo cronograma. Seleciona equipamentos analógicos ou digitais.
Pode escolher software de mixagem, entre as estações de trabalho com
áudio digital (Digital Audio Workstations - DAW) disponíveis. Faz o
orçamento do serviço. Recebe arquivos de áudio após o trabalho de
gravação e após os ajustes, as correções e os acertos realizados na fase
de edição. Analisa os diversos elementos -- como vozes, instrumentos
musicais e ruídos da platéia -, gravados ao vivo ou em estúdio. Concebe
o trabalho de mixagem a ser realizado, utilizando técnicas e trabalho
criativo para definir a forma de unir, nivelar e equilibrar os
elementos, para transmitir emoção e mensagem de uma música e para obter
um produto sonoro que agrade ao ouvinte. Prepara o trabalho de mixagem,
atribuindo um nome adequado a cada faixa e organizando os elementos --
tais como grupos de instrumentos musicais similares - por cor. Define a
ordem para trabalhar cada grupo de elementos, podendo iniciar pela
bateria e percussão, seguir com os instrumentos de corda, continuar com
teclas, até chegar aos vocais, mantendo todos equilibrados, com níveis
de volume aproximados e com ambientações semelhantes. Realiza a limpeza
e a remoção de sons indesejados. Executa a mixagem de áudio, processo de
unir faixas gravadas e editadas, ``misturando-as'' por vários processos
-- tais como equalização, compressão e panorama -, utilizando
processadores físicos (aparelhos) ou virtuais. Efetua a equalização,
aumentando ou diminuindo as frequências específicas de sons graves,
médios e agudos. Executa a compressão, reduzindo a diferença entre o
menor e o maior nível, volume e amplitude dos sinais de áudio. Observa o
panorama (panning), espalhando os sons, posicionando-os adequadamente na
mix, à esquerda, à direita e no meio. Pode realizar outros processos de
mixagem, a partir da análise do objetivo a ser atingido. Finaliza o
processo de mixagem, avaliando o conjunto obtido. Documenta as rotinas
de trabalho, organizando informações, registrando ocorrências da
produção, relatando medidas de contingenciamento e comunicando ajustes
aos diferentes setores envolvidos nos processos de trabalho. Auxilia na
elaboração de pareceres e relatórios técnicos de qualidade e segurança.
Realiza manutenção de primeiro nível em equipamentos e acessórios,
executando pequenos reparos. Providencia serviço de manutenção de
equipamentos, para solucionar problemas mais complexos. Trabalha com
segurança, prevenindo acidentes.

\begin{Form}
\textbf{Esta correspondência está adequada?} \\ 
\ChoiceMenu[radio,radiosymbol=\ding{52},bordercolor=0 0.5 0,name=cbovalid_ 374130 ]{}{CBO deve ficar=sim, \hspace{6em} CBO não fica aqui=nao}
\end{Form}

\textbf{CBO: 374145  ( DJ (disc jockey) )}

Prepara-se para atuar como disc-jockey em evento, verificando o tipo de
produção - cultural, político, religioso, entre outros - e identificando
o público-alvo. Examina as dimensões físicas e identifica parâmetros
acústicos no local do evento, para adequar volume e equalização para o
ambiente. Examina condições de infraestrutura para acesso e instalação
de equipamentos. Faz a especificação de equipamentos, microfones,
transdutores, cabos e conectores. Define demanda de energia elétrica do
sistema. Orça sua participação no evento. Confere o funcionamento e
configura diferentes equipamentos, como mesa, controladores, fones,
``mixers'', ``samplers'', microfones, caixas de som, retornos, entre
outros. Realiza o teste de som (sound check). Pode selecionar e operar
softwares e ferramentas digitais que proporcionam automação da
programação audiofônica, em diferentes plataformas. Escolhe estilos
musicais e pesquisa músicas. Seleciona e organiza o repertório,
elaborando a programação. Edita registros sonoros. Aplica técnicas de
discotecagem, como equalização, mixagem e ajuste de BPM (batidas por
minuto). Reproduz diferentes composições musicais, gravadas ou
produzidas instantaneamente, em diferentes plataformas tecnológicas -
radiodifusão, cabo, satélite ou rede mundial de computadores - e/ou
eventos. Seleciona e opera instrumentos de monitoramento de mau
funcionamento dos sistemas de transmissão e de equipamentos, acionando,
quando necessário, sistemas de contingência e/ou redundância. Implementa
roteiros de programação, garantindo o cumprimento de horários e
sequencialidade da disposição de conteúdos musicais e publicitários em
grades de transmissão ou eventos, zelando pela qualidade técnica e
lógica de transmissão. Pode atuar como apresentador ou mestre de
cerimônia. Pesquisa e cataloga diferentes estilos e estruturas musicais,
considerando fontes de busca e repositórios físicos e digitais.
Mantém-se atualizado sobre os lançamentos e tendências musicais.
Pesquisa novas tecnologias de reprodução de áudio. Documenta o trabalho
executado e registra ocorrências, considerando normas técnicas e legais
que regem o setor, como legislações que regulam publicidade e
propaganda, classificação indicativa, direitos autorais, direitos da
infância e da adolescência, entre outras. Pode elaborar e emitir
pareceres e relatórios. Interage com profissionais que atuam em eventos,
como os responsáveis por iluminação e vídeo. Realiza manutenção de
primeiro nível em equipamentos e acessórios, executando pequenos
reparos. Providencia serviço de manutenção de equipamentos, para
solucionar problemas mais complexos. Trabalha com segurança, prevenindo
acidentes.

\begin{Form}
\textbf{Esta correspondência está adequada?} \\ 
\ChoiceMenu[radio,radiosymbol=\ding{52},bordercolor=0 0.5 0,name=cbovalid_ 374145 ]{}{CBO deve ficar=sim, \hspace{6em} CBO não fica aqui=nao}
\end{Form}

\newpage
\section{ 09035 -- Técnico em Teatro }

\textbf{Perfil do Curso (CNCT)}

O Técnico em Teatro será habilitado para:

Estudar e investigar práticas e métodos do processo de criação teatral
na contemporaneidade, sem perder de vista as perspectivas históricas,
sociais e culturais das artes cênicas locais e mundiais. Atuar
profissionalmente e de maneira interdisciplinar no campo das artes do
palco ? cenografia e figurinos, dramaturgia, direção teatral,
iluminação, sonoplastia e técnicas de palco (cenotécnica). Atuar
abrangendo perspectivas desde o drama ao humor, do teatro
infanto-juvenil ao adulto, do teatro brasileiro ao internacional. Criar
cenas, situações, personagens e figuras, com os procedimentos técnicos,
estéticos e éticos que envolvem o trabalho do atuante no teatro e no
audiovisual. Atuar em diferentes modos da produção em artes cênicas,
tais como teatros de grupo, solos, performances e musical. Reconhecer os
diversos campos da representação artística e da perfomatividade,
considerando as práticas performativas identitárias, as diversidades
culturais e artísticas brasileiras: ameríndias, africanas e europeias.
Criar e produzir pensamento crítico sobre as relações do artista com o
público, dentro da esfera das produções destinadas aos espectadores
infantis, juvenis e adultos. Conhecer os mecanismos que envolvem o
desenvolvimento artístico e cultural nas produções das artes cênicas na
contemporaneidade.

Para atuação como Técnico em Teatro, são fundamentais:

Conhecimentos interdisciplinares relacionados aos processos de criação,
envolvendo pesquisa, idealização, planejamento, execução técnica,
fruição e recepção estética. Competências comunicativas e empreendedoras
voltadas à proposição de projetos, ao coletivo, à gestão, à solução de
problemas e à resiliência, entre outras competências socioemocionais.

\textbf{CBOs correspondidos e seus perfis (presentes no CNCT, e presentes na predição da IA)}

\textbf{CBO: 374205  ( Cenotécnico (cinema, vídeo, televisão, teatro e espetáculos) )}

Planeja as atividades de construção de cenário, mobiliários e adereços,
analisando projeto elaborado por cenógrafo para produção cultural.
Estima tempo de execução e define etapas de trabalho. Dimensiona equipes
de trabalho, verificando a necessidade de contratar serviços de
terceiros. Seleciona ferramentas, máquinas, equipamentos e acessórios.
Faz orçamento do serviço. Submete plano de trabalho à aprovação do
diretor da produção cultural. Realiza adequação do projeto cenográfico
às medidas dos espaços cênicos. Representa graficamente projeto de
ambientes cenográficos. Pode utilizar ``software'' para elaboração de
cenário virtual. Executa pesquisa de vários materiais e objetos, para
escolher quais se adequam melhor ao projeto cenográfico. Faz a
especificação de materiais e objetos para aquisição. Supervisiona a
construção de cenário, de mobiliários e de adereços, orientando a
execução de trabalhos de carpintaria, serralheria, mecânica, modelagem,
pintura e vidraçaria. Coordena a produção de elementos cenográficos,
acompanhando os trabalhos de escultura, em fibras, de costura, de
textura e de cordame. Examina a qualidade do acabamento de cenário,
mobiliário e adereços. Cria e confecciona mecanismos de efeitos
especiais. Supervisiona acondicionamento e transporte de cenário, de
mobiliários e de adereços. Coordena a equipe de montagem do cenário,
considerando o projeto e o espaço cênico. Orienta os trabalhos de
posicionamento do cenário em espaço cênico; de adaptação do cenário ao
tamanho do espaço; de montagem de peças de cenário; e de fixação de
cenário móvel em maquinaria. Acompanha os serviços de decoração do
cenário; de afixação de adereços e objetos em cena; e de montagem dos
efeitos especiais. Compatibiliza cenário e adereços à luz. Confere
urdimento e condições físicas de palcos e ``set'' de filmagem. Organiza
trânsito de coxias. Roteiriza entrada e saída de cenário. Supervisiona a
operação de varas elétricas e guinchos automatizados ou robotizados.
Afina vestimentas de palcos. Pode acompanhar a execução do programa ou
evento cultural, para providenciar reparos de danos em cenário e
adereços; retoques em cenário; ou outras atividades relacionadas à
manutenção de peças cenográficas. Coordena a desmontagem do espaço
cenográfico, o transporte das peças e a organização das partes do
cenário para armazenagem. Elabora relatórios técnicos referentes à
adaptação e à execução de projeto cenográfico. Administra pessoal,
participando do processo de seleção, distribuindo tarefas e controlando
absenteísmo. Supervisiona o trabalho da equipe, avaliando desempenho e
levantando as necessidades de treinamento. Viabiliza a realização de
programas de treinamento. Orienta a equipe para conservar equipamentos,
ferramentas, instrumentos e acessórios limpos, acondicionados e em
plenas condições de uso e funcionamento. Requisita serviço de manutenção
corretiva de equipamentos, quando necessário. Zela pela segurança,
utilizando e monitorando o uso de equipamentos de proteção individual.
Identifica situações de risco e colabora para sua eliminação.

\begin{Form}
\textbf{Esta correspondência está adequada?} \\ 
\ChoiceMenu[radio,radiosymbol=\ding{52},bordercolor=0 0.5 0,name=cbovalid_ 374205 ]{}{CBO deve ficar=sim, \hspace{6em} CBO não fica aqui=nao}
\end{Form}

\textbf{CBO: 374210  ( Maquinista de cinema e vídeo )}

Planeja atividades, analisando o projeto de cenografia e examinando as
adaptações feitas por cenotécnico para adequação do projeto ao espaço
cênico. Estabelece as etapas do seu trabalho. Seleciona equipamentos
analógicos, automatizados e/ou robotizados. Confere materiais
necessários e requisita os itens faltantes. Pode orçar serviços. Atua na
preparação de equipamentos para filmagem cinematográfica e gravação de
vídeo, realizando instalação de suporte para deslocamento e movimentação
de câmeras; prestando apoio a eletricistas na montagem dos parques de
luz e em locação adaptada; operando varas elétricas (guinchos); armando
telas difusoras; e confeccionando mecanismos de efeitos especiais. Monta
cenário e suas peças, posicionando o cenário em espaço cênico; fixando
cenário móvel em maquinaria; e afixando adereços e objetos. Executa
trabalhos de cordame. Monta efeitos especiais. Compatibiliza cenário e
adereços à luz. Afina vestimentas de palcos. Confere urdimento e
condições físicas de ``set'' de filmagem. Examina o cenário montado,
verificando a disposição e o funcionamento de seus componentes, para
certificar-se de sua qualidade e sua segurança. Submete o cenário
finalizado à apreciação de cenotécnico e da direção da produção de filme
ou vídeo, para realizar eventuais ajustes ou melhorias. Durante ensaios,
gravação e filmagem, opera equipamentos - como gruas, carrinhos sobre
trilhos, entre outros - que deslocam e movimentam as câmeras. Manobra
rebatedores, opera cortinas, movimenta cenário sobre rodas, opera varas
cenográficas e opera efeitos especiais, de acordo com roteiro
predefinido de filmagem ou gravação. Pode auxiliar na elaboração de
relatórios técnicos referentes à montagem de cenário e às demais
atividades, relatando ocorrências durante os trabalhos. Recolhe
materiais - após a filmagem ou gravação de vídeo -, organizando-os para
guarda em depósito. Conserva equipamentos, ferramentas, instrumentos e
acessórios limpos, acondicionados e em plenas condições de uso e
funcionamento. Requisita serviço de manutenção corretiva de
equipamentos, quando necessário. Zela pela segurança, utilizando
equipamentos de proteção individual e atuando na prevenção de acidentes.
Identifica situações de risco e colabora para sua eliminação. Combate
princípios de incêndio. Pode prestar primeiros socorros.

\begin{Form}
\textbf{Esta correspondência está adequada?} \\ 
\ChoiceMenu[radio,radiosymbol=\ding{52},bordercolor=0 0.5 0,name=cbovalid_ 374210 ]{}{CBO deve ficar=sim, \hspace{6em} CBO não fica aqui=nao}
\end{Form}

\textbf{CBO: 374215  ( Maquinista de teatro e espetáculos )}

Planeja atividades, analisando o projeto de cenografia e examinando as
adaptações feitas por cenotécnico para adequação do projeto ao espaço
cênico. Estabelece as etapas do seu trabalho. Seleciona equipamentos
analógicos, automatizados e/ou robotizados. Confere materiais
necessários e requisita os itens faltantes. Pode orçar serviços. Atua na
operação de maquinaria para preparação de cenários de produções teatrais
e espetáculos, identificando cenários em relação a varas; fixando
manobras em cenários; e afinando vestimentas de palcos. Desce as varas
de luz, contrapesa refletores instalados e posiciona a vara de
iluminação. Confecciona mecanismos de efeitos especiais. Monta cenário e
suas peças, posiciona o cenário em espaço cênico e afixa adereços e
objetos. Fixa os cenários aéreos e móveis na maquinaria. Executa
trabalhos de cordame. Monta efeitos especiais. Compatibiliza cenário e
adereços à luz. Confere urdimento e observa condições físicas do palco.
Examina o cenário montado, verificando a disposição e o funcionamento de
seus componentes, para certificar-se de sua qualidade e sua segurança.
Submete o cenário finalizado à apreciação de cenotécnico e da direção da
produção de peça teatral ou espetáculo, para realizar eventuais ajustes
ou melhorias. Participa de ensaios de peça teatral ou espetáculo,
marcando altura do cenário aéreo; posicionando objetos de cena;
ajustando a sincronia de palco, luz e som; e verificando os movimentos
entre as coxias. Durante peça de teatro ou espetáculo, libera a sala
para o público e orienta técnicos de palco. Movimenta cenários aéreos e
térreos; cortinas de cena; contrapeso; cordas da maquinaria; alçapões; e
outros elementos cenográficos. Opera efeitos especiais. Pode operar
sistemas de varas, motorizadas, automatizadas ou robotizadas. Organiza
entradas e saídas de cenários e a movimentação por trás dos cenários.
Após término das apresentações de peça de teatro ou espetáculo, desmonta
cenário aéreo; retira adereços; suspende varas cênicas; e desmonta
cenário térreo (em palco) e piso. Armazena itens desmontados e
retirados. Conserva equipamentos, ferramentas, instrumentos e acessórios
limpos, acondicionados e em plenas condições de uso e funcionamento. Faz
a manutenção da maquinaria usada nos cenários. Pode auxiliar na
elaboração de relatórios técnicos referentes à montagem de cenário e às
demais atividades, relatando ocorrências durante os trabalhos. Zela pela
segurança, utilizando equipamentos de proteção individual e atuando na
prevenção de acidentes. Identifica situações de risco e colabora para
sua eliminação. Combate princípios de incêndio. Pode prestar primeiros
socorros.

\begin{Form}
\textbf{Esta correspondência está adequada?} \\ 
\ChoiceMenu[radio,radiosymbol=\ding{52},bordercolor=0 0.5 0,name=cbovalid_ 374215 ]{}{CBO deve ficar=sim, \hspace{6em} CBO não fica aqui=nao}
\end{Form}

\newpage
\section{ 10001 -- Técnico em Açúcar e Álcool }

\textbf{Perfil do Curso (CNCT)}

O Técnico em Açúcar e Álcool será habilitado para:

Controlar e supervisionar operações de processos tecnológicos da
produção de açúcar e álcool e subprodutos. Realizar análises
físico-químicas e microbiológicas de matérias-primas e produtos dos
processos de industrialização da cana-de-açúcar para o controle de
qualidade. Desenvolver produtos e processos em açúcar e álcool.
Colaborar com equipe multidisciplinar nas fases de colheita, transporte,
moagem, industrialização e distribuição dos produtos e subprodutos.
Operar etapas e movimentação de materiais e insumos relacionados à área.

Para atuação como Técnico em Açúcar e Álcool, são fundamentais:

Conhecimentos e saberes relacionados aos processos de planejamento e
operação das atribuições da área, de modo a assegurar a saúde e a
segurança dos trabalhadores e dos futuros usuários e operadores de
empresas em processos de transformação em açúcar e álcool. Conhecimentos
e saberes relacionados à sustentabilidade do processo produtivo, às
normas e relatórios técnicos, à legislação da área, às novas tecnologias
relacionadas à indústria 4.0, à liderança de equipes, à solução de
problemas técnicos e à gestão de conflitos.

\textbf{CBOs correspondidos e seus perfis (presentes no CNCT, e presentes na predição da IA)}

\textbf{CBO: 311105  ( Técnico químico )}

Realiza análises químicas, físicas e microbiológicas, observando
critérios de confiabilidade metrológica e seguindo procedimentos e
normas técnicas e de qualidade. Coleta e prepara amostras. Prepara
reagentes. Registra os resultados das análises, assumindo
responsabilidade técnica por eles. Opera e controla máquinas,
equipamentos e plantas de processamento químico, interpretando manuais
técnicos e procedimentos de produção. Monitora o funcionamento de
máquinas, equipamentos e instrumentos de medição e controle. Mantém
máquinas e equipamentos em condições de uso, realizando regulagens,
quando necessárias. Supervisiona processos de produção, definindo e
coordenando equipes de trabalho, elaborando o cronograma de produção,
definindo o fluxo do processo, monitorando os parâmetros do processo e
emitindo relatórios de produção. Participa do desenvolvimento de
produtos e processos, testando insumos e matérias-primas, elaborando
formulações para produção, testando novos produtos, definindo o processo
de fabricação de produtos, fazendo a adequação dos produtos às
necessidades dos clientes e validando o produto desenvolvido. Colabora
na implantação ou na reestruturação das instalações industriais,
definindo leiautes e fluxogramas de produção, acompanhando montagem e
instalação de equipamentos e testando máquinas e equipamentos. Presta
assistência técnica, identificando necessidades dos clientes,
diagnosticando problemas técnicos e propondo soluções e melhorias de
produtos e processos. Participa da implantação de programas de
qualidade, analisando indicadores, aplicando ferramentas da qualidade,
implementando ações preventivas e corretivas e participando de
auditorias. Elabora documentação técnica, emitindo laudos e redigindo
procedimentos e relatórios. Controla o estoque de matérias-primas,
insumos e produtos, requisitando compra para reposição. Conserva o local
de trabalho limpo e organizado. Mantém, limpa e esteriliza os
instrumentos e os equipamentos de laboratório. Faz o descarte de
resíduos seguindo normas ambientais.

\begin{Form}
\textbf{Esta correspondência está adequada?} \\ 
\ChoiceMenu[radio,radiosymbol=\ding{52},bordercolor=0 0.5 0,name=cbovalid_ 311105 ]{}{CBO deve ficar=sim, \hspace{6em} CBO não fica aqui=nao}
\end{Form}

\textbf{CBO: 311110  ( Técnico de celulose e papel )}

Planeja os processos de obtenção da celulose e de fabricação de papel,
elaborando o cronograma de produção e definindo o fluxo do processo.
Opera e controla máquinas, equipamentos e plantas industriais de
obtenção de celulose e de fabricação de papel, interpretando manuais
técnicos e procedimentos de produção. Monitora o funcionamento de
máquinas, equipamentos e instrumentos de medição e controle. Mantém
máquinas e equipamentos em condições de uso, realizando regulagens,
quando necessárias. Coordena e supervisiona os processos de obtenção da
celulose e de fabricação de papel, definindo e coordenando equipes de
trabalho, monitorando os parâmetros do processo e emitindo relatórios de
produção. Realiza ensaios e análises químicas, físicas e físico-químicas
de matérias-primas, insumos e produtos, observando critérios de
confiabilidade metrológica e seguindo procedimentos e normas técnicas e
de qualidade. Coleta e prepara amostras. Prepara reagentes. Registra os
resultados, assumindo responsabilidade técnica por eles. Participa do
desenvolvimento de produtos e processos, testando insumos e
matérias-primas, elaborando formulações para produção, testando novos
produtos, definindo o processo de fabricação de produtos, fazendo a
adequação dos produtos às necessidades dos clientes e validando o
produto desenvolvido. Participa da implantação de programas de
qualidade, analisando indicadores, aplicando ferramentas da qualidade e
implementando ações preventivas e corretivas. Participa de auditorias.
Colabora na implantação ou na reestruturação das instalações
industriais, definindo leiautes e fluxogramas de produção, acompanhando
montagem e instalação de equipamentos e testando máquinas e
equipamentos. Treina pessoal operacional, identificando necessidades de
treinamento, elaborando programa de treinamento e o material necessário
à sua realização. Avalia o desempenho do trabalhador após o treinamento.
Presta assistência técnica, identificando necessidades dos clientes,
diagnosticando problemas técnicos e propondo soluções e melhorias de
produtos e processos. Pode trabalhar em todas as etapas da rota do
papel, incluindo clonagem, plantio e manejo florestal. Elabora
documentação técnica, emitindo laudos e redigindo procedimentos e
relatórios. Conserva o local de trabalho limpo e organizado. Faz o
descarte de resíduos seguindo normas ambientais.

\begin{Form}
\textbf{Esta correspondência está adequada?} \\ 
\ChoiceMenu[radio,radiosymbol=\ding{52},bordercolor=0 0.5 0,name=cbovalid_ 311110 ]{}{CBO deve ficar=sim, \hspace{6em} CBO não fica aqui=nao}
\end{Form}

\newpage
\section{ 10004 -- Técnico em Biotecnologia }

\textbf{Perfil do Curso (CNCT)}

O Técnico em Biotecnologia será habilitado para:

Executar atividades laboratoriais de biotecnologia e biociências.
Controlar e monitorar processos industriais e laboratoriais da sua área.
Preparar materiais, meios de cultura, soluções e reagentes. Analisar
substâncias e materiais biológicos. Cultivar in vivo e in vitro
microrganismos, células e tecidos animais e vegetais. Auxiliar em
pesquisas de melhoramento genético. Realizar o preparo de amostras dos
tecidos animais e vegetais. Extrair, replicar e quanti?car biomoléculas.
Realizar a produção de imunobiológicos, vacinas, diluentes, kits de
diagnóstico. Operar a criação e manejo de animais de experimentação.
Controlar a qualidade de matérias-primas, insumos e produtos

Para atuação como Técnico em Biotecnologia, são fundamentais:

Conhecimentos e saberes relacionados aos processos de planejamento e
operação das atribuições da área, de modo a assegurar a saúde e a
segurança dos trabalhadores e dos futuros usuários e operadores de
empresas em processos de transformação biotecnológica. Conhecimentos e
saberes relacionados à sustentabilidade do processo produtivo, às normas
e relatórios técnicos, à legislação da área, às novas tecnologias
relacionadas à indústria 4.0, à liderança de equipes, à solução de
problemas técnicos e à gestão de conflitos.

\textbf{CBOs correspondidos e seus perfis (presentes no CNCT, e presentes na predição da IA)}

\textbf{CBO: 325305  ( Técnico em biotecnologia )}

Prepara ambiente de trabalho, mapeando áreas de coleta e experimento;
selecionando equipamentos e utensílios; e montando utensílios de
laboratórios. Confere protocolos e procedimentos. Providencia materiais
aplicados à biotecnologia, coletando materiais para pesquisa;
selecionando e colocando identificação nos materiais biológico e
químico; descontaminando e esterilizando materiais. Faz armazenamento e
descarte dos materiais aplicados à biotecnologia de acordo com normas
técnicas e de biossegurança. Pode prestar apoio técnico à criação de
animais, para sua utilização em experimentos de pesquisa científica.
Cuida dos animais e monitora sua reprodução, sob a supervisão de
profissional de nível superior. Mensura taxa de natalidade e mortalidade
dos animais. Pode abater os animais em condições que envolvam, de acordo
com as espécies, o mínimo possível de sofrimento físico ou mental.
Controla as condições do laboratório para realização de experimento.
Confere condições ambientais, examinando luz, temperatura, umidade e
radiação. Dispõe os equipamentos, verificando seu funcionamento, fazendo
sua calibragem e realizando sua adequação ao experimento. Opera
equipamentos de laboratórios. Pode executar manutenção (de até segundo
nível) de equipamentos. Prepara meios de cultura e soluções, calculando
e pesando reagentes; misturando substâncias para produção de inóculos;
diluindo e concentrando soluções. Desenvolve culturas ``in vivo'' e ``in
vitro'' e marcadores moleculares, para estudos da fisiologia e da
bioquímica de células e para detecção da presença de doenças e
acompanhamento de sua evolução. Cultiva tecidos animal e vegetal para
multiplicação celular. Cultiva e inocula microrganismos. Macera tecidos
animais e vegetais. Realiza extração, replicação, sequenciamento e
quantificação de DNA-Ácido desoxirribonucleico (Deoxyribonucleic acid).
Analisa substâncias e compostos biológicos, medindo volume,
concentração, potencial hidrogeniônico e condutividade de soluções; e
determinando umidade de substâncias e compostos biológicos. Realiza
análises químicas e microbiológicas, atuando na produção de
imunobiológicos (medicamentos, soros e vacinas) e nos processos
biotecnológicos de indústrias - química, de alimentos, de medicamentos,
de cosméticos e de produtos veterinários -, sob supervisão de
profissional de nível superior. Controla a qualidade de matérias-primas
e produtos utilizados. Mantém-se atualizado em relação às pesquisas e
inovações na área de biotecnologia. Acompanha a introdução da
automatização e do uso de robótica nos processos e o desenvolvimento de
novos materiais e produtos. Monitora a limpeza do local de trabalho.
Orienta a aplicação de técnicas de higienização. Aplica procedimentos
para cumprimento de normas ambientais, promovendo o tratamento
estabelecido e o destino adequado de materiais biológicos e de outros
resíduos descartados. Atua seguindo as normas de segurança e de
biossegurança, submetendo-se a exames periódicos; vacinando-se contra
agentes patogênicos; e usando equipamentos de proteção individual.
Participa em ações para prevenção de acidentes, colaborando na
eliminação de situações de risco. Atua no combate a pequenos focos de
fogo, manuseando extintores de incêndio. Aplica cuidados de primeiros
socorros.

\begin{Form}
\textbf{Esta correspondência está adequada?} \\ 
\ChoiceMenu[radio,radiosymbol=\ding{52},bordercolor=0 0.5 0,name=cbovalid_ 325305 ]{}{CBO deve ficar=sim, \hspace{6em} CBO não fica aqui=nao}
\end{Form}

\textbf{CBO: 325310  ( Técnico em imunobiológicos )}

Prepara ambiente de trabalho, mapeando áreas de coleta e experimento;
selecionando equipamentos e utensílios; e montando utensílios de
laboratórios. Confere protocolos e procedimentos. Providencia materiais
aplicados à biotecnologia, coletando materiais para pesquisa;
selecionando e colocando identificação nos materiais biológico e
químico; descontaminando e esterilizando materiais. Faz armazenamento e
descarte dos materiais aplicados à biotecnologia de acordo com normas
técnicas e de biossegurança. Pode prestar apoio técnico à criação de
animais, para sua utilização em experimentos de pesquisa científica.
Cuida dos animais e monitora sua reprodução, sob a supervisão de
profissional de nível superior. Mensura taxas de natalidade e
mortalidade dos animais. Pode abater os animais em condições que
envolvam, de acordo com as espécies, o mínimo possível de sofrimento
físico ou mental. Controla as condições do laboratório para realização
de experimento. Confere condições ambientais, examinando luz,
temperatura, umidade e radiação. Dispõe os equipamentos, verificando seu
funcionamento, fazendo sua calibragem e realizando sua adequação ao
experimento. Opera equipamentos de laboratórios. Pode executar
manutenção (de até segundo nível) em equipamentos. Prepara meios de
cultura e soluções, calculando e pesando reagentes; misturando
substâncias para produção de inóculos; diluindo e concentrando soluções.
Desenvolve culturas ``in vivo'' e ``in vitro'' e marcadores moleculares,
para estudos da fisiologia e da bioquímica de células e para detecção da
presença de doenças e acompanhamento de sua evolução. Cultiva tecidos
animal e vegetal para multiplicação celular. Cultiva e inocula
microrganismos. Macera tecidos animais e vegetais. Faz extração,
replicação, sequenciamento e quantificação de DNA-Ácido
desoxirribonucleico (Deoxyribonucleic acid). Realiza manipulação de
RNA-Ácido ribonucleico (Ribonucleic acid). Analisa substâncias e
compostos biológicos, medindo volume, concentração, potencial
hidrogeniônico e condutividade de soluções; e determinando umidade de
substâncias e compostos biológicos. Realiza análises químicas e
microbiológicas, atuando na produção de imunobiológicos (medicamentos,
soros, vacinas e kits de diagnóstico) e no estudo de doenças infecciosas
causadas por microrganismos - tais como bactérias e vírus -, sob
supervisão de profissional de nível superior. Auxilia em pesquisas de
melhoramento genético. Pode colaborar em atividades de perícia criminal
e investigação genética. Controla a qualidade de matérias-primas e
produtos utilizados. Mantém-se atualizado em relação às pesquisas e
inovações na área de imunobiologia. Colabora em estudos para implantação
de inovações, tais como os ``órgãos-em-chip'', capazes de servir como
``simuladores'' de diferentes órgãos do corpo humano, permitindo que
pesquisadores monitorem a jornada de um corpo externo - como bactéria e
droga farmacêutica - com maior segurança e controle, em substituição ao
uso de animais em experimentos. Monitora a limpeza do local de trabalho.
Orienta a aplicação de técnicas de higienização. Aplica procedimentos
para cumprimento de normas ambientais, promovendo o tratamento
estabelecido e o destino adequado de materiais biológicos e de outros
resíduos descartados. Atua seguindo as normas de segurança e de
biossegurança, submetendo-se a exames periódicos; vacinando-se contra
agentes patogênicos; e usando equipamentos de proteção individual.
Participa em ações para prevenção de acidentes, colaborando na
eliminação de situações de risco. Atua no combate a pequenos focos de
fogo, manuseando extintores de incêndio. Aplica cuidados de primeiros
socorros.

\begin{Form}
\textbf{Esta correspondência está adequada?} \\ 
\ChoiceMenu[radio,radiosymbol=\ding{52},bordercolor=0 0.5 0,name=cbovalid_ 325310 ]{}{CBO deve ficar=sim, \hspace{6em} CBO não fica aqui=nao}
\end{Form}

\newpage
\section{ 10005 -- Técnico em Calçados }

\textbf{Perfil do Curso (CNCT)}

O Técnico em Calçados será habilitado para:

Executar e supervisionar operações relativas à fabricação de calçados,
preparação, corte, costura (pesponto), montagem e acabamento. Auxiliar
no planejamento, programação e controle dos processos de produção.
Controlar a qualidade de matérias-primas, insumos e produtos. Acompanhar
tendências de mercado na área de calçados. Verificar materiais
alternartivos para confecção de calçados.

Para atuação como Técnico em Calçados, são fundamentais:

Conhecimentos e saberes relacionados aos processos de planejamento e
operação das atribuições da área, de modo a assegurar a saúde e a
segurança dos trabalhadores e dos futuros usuários e operadores de
empresas em processos de transformação em calçados. Conhecimentos e
saberes relacionados à sustentabilidade do processo produtivo, às normas
e relatórios técnicos, à legislação da área, às novas tecnologias
relacionadas à indústria 4.0, à liderança de equipes, à solução de
problemas técnicos e à gestão de conflitos.

\textbf{CBOs correspondidos e seus perfis (presentes no CNCT, e presentes na predição da IA)}

\textbf{CBO: 318815  ( Modelista de calçados )}

Planeja a confecção de moldes e o desenvolvimento de protótipos de
calçados, examinando projeto do designer; verificando tendências de
mercado; consultando boletins técnicos; avaliando estilos de design,
concepções e produtos de concorrentes; e coletando informações
tecnológicas em feiras de calçados. Confecciona moldes de calçados,
interpretando modelos e desenhos; elaborando desenhos planificados;
traçando, cortando e codificando moldes. Utiliza prancheta e softwares,
tais como CAD-Desenho Assistido por Computador (Computer Aided Design) e
CAM--Manufatura Auxiliada por Computador (Computer Aided Manufacturing).
Marca referências em moldes, com a finalidade de orientar a produção de
calçados. Desenvolve protótipos de calçados, interpretando normas
técnicas para produção; elaborando ficha técnica; especificando
aviamentos e acessórios; ajustando moldes e gabaritos; e determinando a
posição de ornamentos, detalhes e acessórios. Produz e testa protótipo
de calçado. Requisita testes laboratoriais do protótipo. Realiza
ampliações e reduções das dimensões do produto por escala. Avalia
materiais para aquisição com vistas à produção de calçados ergonômicos,
verificando a resistência dos materiais e avaliando sua composição e
acabamento. Avalia a costurabilidade e testa a fixação de cores em
couros. Interpreta resultados de testes laboratoriais, elaborando
relatórios de testagem de materiais. Avalia a relação entre custo e
benefício.

\begin{Form}
\textbf{Esta correspondência está adequada?} \\ 
\ChoiceMenu[radio,radiosymbol=\ding{52},bordercolor=0 0.5 0,name=cbovalid_ 318815 ]{}{CBO deve ficar=sim, \hspace{6em} CBO não fica aqui=nao}
\end{Form}

\textbf{CBO: 319105  ( Técnico em calçados e artefatos de couro )}

Participa no planejamento dos processos de produção de calçados e de
artefatos de couro, considerando as metas preestabelecidas. Propõe
leiaute da produção. Seleciona máquinas e equipamentos. Monitora
estoques de insumos. Auxilia no dimensionamento do quadro de pessoal
necessário nos vários setores de produção. Participa na elaboração da
programação e do cronograma de atividades de produção. Prepara a
produção de calçados e de artefatos de couro, interpretando ordem de
produção. Verifica as especificações e interpreta os desenhos referentes
aos modelos a serem produzidos. Coordena a produção de protótipos,
fornecendo as informações necessárias; monitorando o processo de
produção; e testando o protótipo produzido. Pode utilizar impressão 3D
na confecção dos protótipos. Supervisiona operações relativas à produção
de calçados e de artefatos de couro, orientando corte, costura
(pesponto), montagem e acabamento. Controla a produção, medindo tempo
para desenvolvimento das etapas; implementando ações corretivas nos
processos; e verificando as condições de funcionamento de máquinas e
equipamentos. Controla estoques de produtos acabados. Controla a
qualidade dos processos e produtos. Mede indicadores de desempenho,
avaliando produtividade. Identifica perdas na produção. Elabora
estatísticas de produção. Pode orientar o uso de ``softwares'', como os
de modelagem e os relacionados a outras etapas da produção. Pode fazer
uso de ``softwares'' e sistemas integrados de gestão empresarial para
controle de produção; controle de estoques; e gerenciamento de custos.
Verifica a necessidade de melhoria nos processos de trabalho e propõe
alternativa. Descreve o processo produtivo proposto, definindo o fluxo e
estabelecendo o tempo-padrão das operações de produção. Desenvolve novos
produtos, analisando tendências de mercado e considerando o
público-alvo. Pesquisa novas matérias-primas, identificando as
apropriadas para os produtos. Elabora especificações das dimensões e das
características dos produtos. Elabora ficha técnica dos novos produtos.
Desenvolve amostras e faz sua apresentação. Indica sistemas de embalagem
dos produtos e de acondicionamento da produção. Analisa viabilidade
técnica e econômica dos novos produtos. Verifica a concorrência e propõe
ações de marketing para os novos produtos. Mantém-se atualizado em
relação às inovações na área de calçados e artefatos de couro. Avalia
inovações, como novas máquinas; sistemas de produção automatizados;
sistemas a laser; entre outras. Identifica fornecedores, testando
equipamentos e materiais oferecidos e comparando preços, prazos de
entrega e formas de pagamento. Levanta necessidades de treinamento de
fornecedores e desenvolve programação. Coordena equipe de trabalho,
participando do processo de seleção e de administração de pessoal.
Distribui tarefas, controla absenteísmo e programa férias da equipe.
Supervisiona o trabalho da equipe, avaliando desempenho e levantando as
necessidades de treinamento. Viabiliza a realização de programas de
treinamento. Participa da administração financeira, acompanhando a
execução do orçamento e calculando custo de produto. Documenta o plano
de manutenção preventiva e preditiva de máquinas e equipamentos. Mantém
máquinas e equipamentos em condições operacionais, providenciando
manutenção de primeiro nível. Requisita serviços de manutenção corretiva
de máquinas e equipamentos. Monitora a limpeza e a organização do local
de trabalho. Orienta equipe para conservar ferramentas e instrumentos
limpos, acondicionados e em plenas condições de uso e funcionamento.
Contribui para a redução dos impactos ambientais do processo produtivo,
orientando o reaproveitamento de sobras e monitorando a aplicação das
normas ambientais no descarte de resíduos. Zela pela segurança,
identificando situações de risco e colaborando para sua eliminação.
Elabora procedimentos em conformidade com as normas de segurança.
Utiliza e monitora o uso de equipamentos de proteção individual.

\begin{Form}
\textbf{Esta correspondência está adequada?} \\ 
\ChoiceMenu[radio,radiosymbol=\ding{52},bordercolor=0 0.5 0,name=cbovalid_ 319105 ]{}{CBO deve ficar=sim, \hspace{6em} CBO não fica aqui=nao}
\end{Form}

\textbf{CBO: 319110  ( Técnico em confecções do vestuário )}

Participa no planejamento de produção de confecções em indústria de
vestuário, considerando as metas preestabelecidas. Propõe leiaute da
produção. Seleciona máquinas e equipamentos. Monitora estoques de
insumos. Auxilia no dimensionamento do quadro de pessoal necessário nos
vários setores de produção. Participa na elaboração da programação e do
cronograma de atividades de produção. Prepara a produção de confecções,
interpretando ordem de produção. Verifica especificações de materiais.
Lê e interpreta desenhos e projeto de confecções. Coordena a produção de
protótipos, fornecendo as informações necessárias; monitorando o
processo de produção; e testando o protótipo produzido. Supervisiona
operações relativas à produção de confecções, orientando corte, costura,
montagem e acabamento. Controla a produção, medindo tempo para
desenvolvimento das etapas; implementando ações corretivas nos
processos; e verificando as condições de funcionamento de máquinas e
equipamentos. Controla estoques de produtos acabados. Realiza controle
de qualidade dos processos e produtos. Mede indicadores de desempenho,
avaliando produtividade. Identifica perdas na produção. Elabora
estatísticas de produção. Pode orientar o uso de máquinas automáticas de
enfesto, a utilização de ``software'' de modelagem, entre outras novas
tecnologias aplicadas ao processo de produção. Pode fazer uso de
``softwares'' e sistemas integrados de gestão empresarial para controle
de produção; controle de estoques; e gerenciamento de custos. Verifica a
necessidade de melhoria nos processos de trabalho e propõe alternativa.
Descreve o processo produtivo proposto, definindo o fluxo e
estabelecendo o tempo-padrão das operações de produção. Desenvolve nova
coleção de confecções, considerando o público-alvo e analisando
tendências da moda. Pesquisa novas matérias-primas, identificando as
apropriadas para as confecções. Detalha as características dos produtos.
Elabora ficha técnica das novas confecções. Desenvolve amostras e faz
sua apresentação. Analisa e define a melhor sequência de montagem dos
produtos, considerando as características da matéria-prima. Indica
sistemas de embalagem dos produtos e de acondicionamento da produção.
Analisa viabilidade técnica e econômica dos novos produtos. Verifica a
concorrência e propõe ações de marketing para as novas confecções.
Identifica fornecedores, testando equipamentos e materiais oferecidos e
comparando preços, prazos de entrega e formas de pagamento. Levanta
necessidades de treinamento de fornecedores e desenvolve programação.
Mantém-se atualizado em relação às inovações na área, como máquinas
robotizadas de costura, tecidos inteligentes, entre outras. Coordena
equipe de trabalho, participando do processo de seleção e de
administração de pessoal. Distribui tarefas, controla absenteísmo e
programa férias da equipe. Supervisiona o trabalho da equipe, avaliando
desempenho e levantando as necessidades de treinamento. Viabiliza a
realização de programas de treinamento. Participa da administração
financeira, acompanhando a execução do orçamento e calculando custo de
produto. Documenta o plano de manutenção preventiva e preditiva de
máquinas e equipamentos. Mantém máquinas e equipamentos em condições
operacionais, providenciando manutenção de primeiro nível. Requisita
serviços de manutenção corretiva de máquinas e equipamentos. Monitora a
limpeza e a organização do local de trabalho. Orienta equipe para
conservar ferramentas e instrumentos limpos, acondicionados e em plenas
condições de uso e funcionamento. Contribui para a redução dos impactos
ambientais do processo produtivo, orientando o reaproveitamento de
sobras e monitorando a aplicação das normas ambientais no descarte de
resíduos. Zela pela segurança, identificando situações de risco e
colaborando para sua eliminação. Elabora procedimentos em conformidade
com as normas de segurança. Utiliza e monitora o uso de equipamentos de
proteção individual.

\begin{Form}
\textbf{Esta correspondência está adequada?} \\ 
\ChoiceMenu[radio,radiosymbol=\ding{52},bordercolor=0 0.5 0,name=cbovalid_ 319110 ]{}{CBO deve ficar=sim, \hspace{6em} CBO não fica aqui=nao}
\end{Form}

\newpage
\section{ 10006 -- Técnico em Celulose e Papel }

\textbf{Perfil do Curso (CNCT)}

O Técnico em Celulose e Papel será habilitado para:

Controlar processos de obtenção da celulose e de fabricação de papel.
Realizar ensaios e análises químicas, físicas e físico-químicas de
matérias-primas e produtos seguindo normas e procedimentos técnicos.
Planejar, executar e supervisionar os processos de secagem e corte na
produção de papel.

Para atuação como Técnico em Celulose e Papel, são fundamentais:

Conhecimentos e saberes relacionados aos processos de planejamento e
operação das atribuições da área, de modo a assegurar a saúde e a
segurança dos trabalhadores e dos futuros usuários e operadores de
empresas em processos de obtenção da celulose e de fabricação de papel.
Conhecimentos e saberes relacionados à sustentabilidade do processo
produtivo, às normas e relatórios técnicos, à legislação da área, às
novas tecnologias relacionadas à indústria 4.0, à liderança de equipes,
à solução de problemas técnicos e à gestão de conflitos.

\textbf{CBOs correspondidos e seus perfis (presentes no CNCT, e presentes na predição da IA)}

\textbf{CBO: 311105  ( Técnico químico )}

Realiza análises químicas, físicas e microbiológicas, observando
critérios de confiabilidade metrológica e seguindo procedimentos e
normas técnicas e de qualidade. Coleta e prepara amostras. Prepara
reagentes. Registra os resultados das análises, assumindo
responsabilidade técnica por eles. Opera e controla máquinas,
equipamentos e plantas de processamento químico, interpretando manuais
técnicos e procedimentos de produção. Monitora o funcionamento de
máquinas, equipamentos e instrumentos de medição e controle. Mantém
máquinas e equipamentos em condições de uso, realizando regulagens,
quando necessárias. Supervisiona processos de produção, definindo e
coordenando equipes de trabalho, elaborando o cronograma de produção,
definindo o fluxo do processo, monitorando os parâmetros do processo e
emitindo relatórios de produção. Participa do desenvolvimento de
produtos e processos, testando insumos e matérias-primas, elaborando
formulações para produção, testando novos produtos, definindo o processo
de fabricação de produtos, fazendo a adequação dos produtos às
necessidades dos clientes e validando o produto desenvolvido. Colabora
na implantação ou na reestruturação das instalações industriais,
definindo leiautes e fluxogramas de produção, acompanhando montagem e
instalação de equipamentos e testando máquinas e equipamentos. Presta
assistência técnica, identificando necessidades dos clientes,
diagnosticando problemas técnicos e propondo soluções e melhorias de
produtos e processos. Participa da implantação de programas de
qualidade, analisando indicadores, aplicando ferramentas da qualidade,
implementando ações preventivas e corretivas e participando de
auditorias. Elabora documentação técnica, emitindo laudos e redigindo
procedimentos e relatórios. Controla o estoque de matérias-primas,
insumos e produtos, requisitando compra para reposição. Conserva o local
de trabalho limpo e organizado. Mantém, limpa e esteriliza os
instrumentos e os equipamentos de laboratório. Faz o descarte de
resíduos seguindo normas ambientais.

\begin{Form}
\textbf{Esta correspondência está adequada?} \\ 
\ChoiceMenu[radio,radiosymbol=\ding{52},bordercolor=0 0.5 0,name=cbovalid_ 311105 ]{}{CBO deve ficar=sim, \hspace{6em} CBO não fica aqui=nao}
\end{Form}

\textbf{CBO: 311110  ( Técnico de celulose e papel )}

Planeja os processos de obtenção da celulose e de fabricação de papel,
elaborando o cronograma de produção e definindo o fluxo do processo.
Opera e controla máquinas, equipamentos e plantas industriais de
obtenção de celulose e de fabricação de papel, interpretando manuais
técnicos e procedimentos de produção. Monitora o funcionamento de
máquinas, equipamentos e instrumentos de medição e controle. Mantém
máquinas e equipamentos em condições de uso, realizando regulagens,
quando necessárias. Coordena e supervisiona os processos de obtenção da
celulose e de fabricação de papel, definindo e coordenando equipes de
trabalho, monitorando os parâmetros do processo e emitindo relatórios de
produção. Realiza ensaios e análises químicas, físicas e físico-químicas
de matérias-primas, insumos e produtos, observando critérios de
confiabilidade metrológica e seguindo procedimentos e normas técnicas e
de qualidade. Coleta e prepara amostras. Prepara reagentes. Registra os
resultados, assumindo responsabilidade técnica por eles. Participa do
desenvolvimento de produtos e processos, testando insumos e
matérias-primas, elaborando formulações para produção, testando novos
produtos, definindo o processo de fabricação de produtos, fazendo a
adequação dos produtos às necessidades dos clientes e validando o
produto desenvolvido. Participa da implantação de programas de
qualidade, analisando indicadores, aplicando ferramentas da qualidade e
implementando ações preventivas e corretivas. Participa de auditorias.
Colabora na implantação ou na reestruturação das instalações
industriais, definindo leiautes e fluxogramas de produção, acompanhando
montagem e instalação de equipamentos e testando máquinas e
equipamentos. Treina pessoal operacional, identificando necessidades de
treinamento, elaborando programa de treinamento e o material necessário
à sua realização. Avalia o desempenho do trabalhador após o treinamento.
Presta assistência técnica, identificando necessidades dos clientes,
diagnosticando problemas técnicos e propondo soluções e melhorias de
produtos e processos. Pode trabalhar em todas as etapas da rota do
papel, incluindo clonagem, plantio e manejo florestal. Elabora
documentação técnica, emitindo laudos e redigindo procedimentos e
relatórios. Conserva o local de trabalho limpo e organizado. Faz o
descarte de resíduos seguindo normas ambientais.

\begin{Form}
\textbf{Esta correspondência está adequada?} \\ 
\ChoiceMenu[radio,radiosymbol=\ding{52},bordercolor=0 0.5 0,name=cbovalid_ 311110 ]{}{CBO deve ficar=sim, \hspace{6em} CBO não fica aqui=nao}
\end{Form}

\textbf{CBO: 311705  ( Colorista de papel )}

Recebe, inspeciona e testa corantes, pigmentos e demais materiais
utilizados na coloração de papéis, para garantir a conformidade com as
especificações. Desenvolve cores em laboratórios, selecionando, pesando
e misturando pigmentos, corantes e materiais auxiliares, para preparar
tinturas de acordo com especificações. Coleta amostras de materiais ou
produtos intermediários e acabados para testes laboratoriais. Testa a
qualidade de tinturas aplicadas em papel, comparando a cor com os
padrões estabelecidos - seja visualmente ou usando equipamentos como
colorímetros ou espectrofotômetros -, ajustando e padronizando cores.
Desenvolve cores e elabora cartelas de amostras de cores de acordo com
tendências e em conformidade com padrões internacionais. Controla a
qualidade do processo de tingimento e impressão, verificando a
viscosidade e o pH de tintas pastosas e examinando a densidade e a
espessura de camada de tinta nas impressões. Pode operar e manter
máquinas utilizadas no processo de coloração de papel. Registra os dados
operacionais ou de produção em formulários especificados. Conserva local
de trabalho limpo e organizado. Armazena corantes, pigmentos e demais
materiais. Destina adequadamente os resíduos gerados, reciclando papel e
realizando descarte de acordo com as normas ambientais. Zela pela
segurança, prevenindo acidentes, eliminando condições inseguras e
utilizando equipamentos de proteção individual.

\begin{Form}
\textbf{Esta correspondência está adequada?} \\ 
\ChoiceMenu[radio,radiosymbol=\ding{52},bordercolor=0 0.5 0,name=cbovalid_ 311705 ]{}{CBO deve ficar=sim, \hspace{6em} CBO não fica aqui=nao}
\end{Form}

\newpage
\section{ 10007 -- Técnico em Cerâmica }

\textbf{Perfil do Curso (CNCT)}

O Técnico em Cerâmica será habilitado para:

Planejar, coordenar e supervisionar etapas de produção de materiais
cerâmicos. Operar e controlar linhas de produção de produtos cerâmicos.
Utilizar máquinas, equipamentos e instrumentos da indústria cerâmica.
Manipular e caracterizar matérias-primas e insumos na indústria
cerâmica. Desenvolver melhorias no processo produtivo e programar a
produção. Realizar ensaios físico-químicos para o controle de qualidade
da matéria-prima e do produto acabado.\\
Controlar estoques de produtos acabados.

Para atuação como Técnico em Cerâmica, são fundamentais:

Conhecimentos e saberes relacionados aos processos de planejamento e
operação das atribuições da área, de modo a assegurar a saúde e a
segurança dos trabalhadores e dos futuros usuários e operadores de
empresas em processos de transformação em cerâmica. Conhecimentos e
saberes relacionados à sustentabilidade do processo produtivo, às normas
e relatórios técnicos, à legislação da área, às novas tecnologias
relacionadas à indústria 4.0, à liderança de equipes, à solução de
problemas técnicos e à gestão de conflitos.

\textbf{CBOs correspondidos e seus perfis (presentes no CNCT, e presentes na predição da IA)}

\textbf{CBO: 311305  ( Técnico em materiais, produtos cerâmicos e vidros )}

Elabora programação de produção de materiais cerâmicos e vidros,
avaliando informações do mercado e interagindo com os setores do
estabelecimento industrial para definir a capacidade de produção.
Elabora cronograma de produção, programando consumo de matéria-prima.
Pode realizar ajustes na programação da produção. Calcula composição do
produto e define quantidade de matérias-primas para a produção. Testa
novas formulações, avaliando o material formulado. Avalia custos e
benefícios das matérias-primas e qualifica fornecedores. Supervisiona o
processo de produção de materiais cerâmicos e vidros, considerando metas
e organizando o fluxo de produção. Controla o processo de produção,
coordenando equipes de trabalho, monitorando o funcionamento de máquinas
e equipamentos, checando variáveis no processo e inspecionando a
qualidade do produto intermediário e acabado da produção. Pode operar
máquinas e equipamentos industriais, interpretando manuais. Prepara e
regula máquinas e equipamentos. Executa controle do produto. Realiza
ensaios físico-químicos para o controle de qualidade da matéria-prima e
do produto acabado, coletando e preparando amostras, mantendo
equipamentos calibrados e operando instrumentos de medição e controle.
Executa ensaios e testes em matérias-primas e produtos, aplicando normas
técnicas. Interpreta laudos de análises físico-químicas. Desenvolve
melhorias no processo de produção, avaliando as falhas existentes,
analisando dados estatísticos, criando dispositivos para otimização da
produção e adaptando leiautes. Identifica novas tecnologias que trazem
melhorias e eficiência ao processo, justificando a necessidade de
aquisição de novos sistemas, equipamentos e materiais. Desenvolve
produto, interpretando desenhos e projetos. Define as características do
produto, dimensionando-o de acordo com normas técnicas. Adequa o produto
ao processo produtivo ou adequa o processo de produção ao produto
desenvolvido. Realiza prova industrial e estima o custo do produto.
Mantém-se atualizado na sua área tecnológica de atuação. Presta
assistência técnica, identificando defeitos no produto, interpretando
dados da produção, solicitando ensaios de produtos acabados e
solucionando os problemas técnicos, para atender clientes internos ou
externos. Planeja visitas técnicas, examinando a aplicação do produto e
interpretando termos de garantia. Sugere substituição do produto ou
indenização para cliente. Propõe melhorias no produto. Registra
informações técnicas e administrativas, descrevendo procedimentos
operacionais e elaborando ficha técnica do produto. Emite laudos
técnicos. Elabora relatórios, incluindo gráficos e fluxogramas.
Administra equipe de trabalho, participando da seleção do pessoal;
distribuindo tarefas; definindo período de férias; e estabelecendo
trocas de pessoal. Divulga normas e procedimentos para a equipe.
Supervisiona equipe de trabalho, avaliando seu desempenho e
identificando a necessidade de treinamento. Planeja o desenvolvimento da
equipe, definindo conteúdo programático, preparando material didático e
organizando o local da programação. Realiza treinamentos operacionais e
profere palestras técnicas, avaliando o desempenho dos participantes.
Analisa o impacto do programa de desenvolvimento no desempenho da
equipe. Mantém os equipamentos em condições operacionais, monitorando
seu funcionamento e providenciando manutenção corretiva e preventiva.
Inspeciona a limpeza e a organização do local de trabalho. Examina a
limpeza, o acondicionamento e a conservação de ferramentas e
instrumentos de trabalho. Controla parâmetros de poluição ambiental.
Monitora o descarte de resíduos, verificando a conformidade em relação
às normas ambientais. Zela pela segurança, identificando situações de
risco e participando na elaboração, na implementação e na avaliação de
ações corretivas. Orienta e supervisiona o uso de equipamentos de
proteção individual e coletiva.

\begin{Form}
\textbf{Esta correspondência está adequada?} \\ 
\ChoiceMenu[radio,radiosymbol=\ding{52},bordercolor=0 0.5 0,name=cbovalid_ 311305 ]{}{CBO deve ficar=sim, \hspace{6em} CBO não fica aqui=nao}
\end{Form}

\textbf{CBO: 752305  ( Ceramista )}

Prepara as massas cerâmicas, interpretando formulações, dosando
componentes da mistura e carregando moinhos e misturadores. Realiza a
moagem e a mistura do material, descarregando os moinhos, realizando a
filtro-prensagem e a extrusão de massas pastosas. Opera o atomizador
para a preparação de massas em pó, armazenando as massas preparadas.
Verifica as características das massas, controlando resíduos,
viscosidade e densidade das massas líquidas; realizando análise
granulométrica; e controlando a umidade das massas pastosas e secas.
Prepara a modelagem e a conformação de peças, selecionando e instalando
moldes e formas; e abastecendo prensas, moldes e tornos com massa
cerâmica. Ajusta prensas e tornos segundo as especificações do processo.
Molda, torneia ou prensa a massa cerâmica. Monitora o processo de
moldagem e conformação das peças, controlando a densidade aparente, a
pressão de compactação, a umidade da massa cerâmica e a temperatura de
operação do equipamento. Verifica as dimensões e o peso das peças
cerâmicas. Monitora o volume de produção. Executa o acabamento das peças
cerâmicas, rebarbando, polindo, esponjando e colando partes da peça.
Ajusta as dimensões e o esquadro antes da secagem. Realiza a secagem das
peças, operando o secador, controlando a curva de secagem e verificando
a umidade residual após o processo. Executa a queima das peças, operando
o forno de acordo com curvas de queima padronizadas. Controla a
qualidade da queima das peças cerâmicas. Prepara tintas, esmaltes e
vernizes, interpretando formulações; dosando os componentes da mistura;
realizando a moagem dos componentes da mistura do esmalte; misturando
componentes para tintas e vernizes; e descarregando os moinhos e os
misturadores. Controla as características dos esmaltes, tintas e
vernizes, medindo a densidade e a viscosidade, testando e comparando com
padrão de cores e efeitos superficiais. Realiza ajustes, quando
necessários. Aplica tintas, esmaltes e vernizes por processo mecanizado
ou manual, abastecendo a linha de esmaltação; controlando a viscosidade
e a densidade de tintas, esmaltes e vernizes; e controlando as camadas
de aplicação. Verifica o acabamento do esmalte aplicado, corrigindo
defeitos quando possível. Decora as peças manualmente ou por processo
mecanizado, aplicando decalques, filetes, barrados e imprimindo a marca
do fabricante. Verifica a qualidade das peças acabadas, classificando-as
em função dos possíveis defeitos existentes e selecionando-as por
tonalidade, dimensões e sons. Embala as peças para destiná-las ao
mercado. Pode utilizar insumos de base nanotecnológica, que alteram os
procedimentos de trabalho. Faz reaproveitamento - reciclagem - de
materiais. Realiza o descarte de resíduos de acordo com as normas
ambientais. Mantém ferramentas e instrumentos de trabalho limpos,
organizados, acondicionados e em plenas condições de uso e
funcionamento. Trabalha com segurança, prevenindo incêndios e acidentes.
Utiliza equipamentos de proteção individual e coletiva, em ambientes
sujeitos a poeiras, ruído e altas temperaturas.

\begin{Form}
\textbf{Esta correspondência está adequada?} \\ 
\ChoiceMenu[radio,radiosymbol=\ding{52},bordercolor=0 0.5 0,name=cbovalid_ 752305 ]{}{CBO deve ficar=sim, \hspace{6em} CBO não fica aqui=nao}
\end{Form}

\newpage
\section{ 10008 -- Técnico em Construção Naval }

\textbf{Perfil do Curso (CNCT)}

O Técnico em Construção Naval será habilitado para:

Realizar ensaios e testes e montar componentes na fabricação e
manutenção naval. Desenvolver projetos de construção naval. Controlar e
inspecionar os processos de construção em plantas navais. Apoiar a
coordenação de construção de embarcações e estruturas hidroviárias.
Realizar manutenção e operação de sistemas de navegação. Selecionar
materiais a serem empregados na construção naval. Analisar custos
operacionais na área. Testar a velocidade e a segurança de barcos e
navios. Montar e organizar estaleiros. Operar sistemas de logística para
controle do frete, do armazenamento e da distribuição de cargas. Emitir
laudos técnicos e fazer vistorias nas companhias de navegação dentro de
suas atribuições técnicas.

Para atuação como Técnico em Construção Naval, são fundamentais:

Conhecimentos e saberes relacionados aos processos de planejamento e
operação das atribuições da área, de modo a assegurar a saúde e a
segurança dos trabalhadores e dos futuros usuários e operadores de
empresas em processos de transformação em construção naval.
Conhecimentos e saberes relacionados à sustentabilidade do processo
produtivo, às normas e relatórios técnicos, à legislação da área, às
novas tecnologias relacionadas à indústria 4.0, à liderança de equipes,
à solução de problemas técnicos e à gestão de conflitos.

\textbf{CBOs correspondidos e seus perfis (presentes no CNCT, e presentes na predição da IA)}

\textbf{CBO: 318215  ( Desenhista técnico naval )}

Planeja o trabalho, levando em consideração o cronograma geral da área
ou do projeto -- ao qual se vinculam as atividades de elaboração do
desenho técnico naval --, estabelecendo as prioridades e estimando o
tempo de execução do desenho. Define os meios computacionais e/ou
manuais a serem empregados nos trabalhos e elabora o cronograma
específico de execução do desenho. Coleta os dados necessários à
elaboração do desenho técnico naval, estudando as características do
projeto; identificando interfaces com áreas envolvidas no projeto; e
levantando informações básicas e informações complementares junto às
áreas de interface, a fim de atender aos requisitos de projeto em todas
as etapas de elaboração do desenho. Consulta tabelas, gráficos, normas
técnicas e outras fontes. Identifica o material -- das peças ou dos
conjuntos que integram os artefatos mecânicos de embarcações a serem
desenhados -- e identifica a aplicabilidade do projeto. Identifica
inovações tecnológicas na área naval e seus reflexos nos elementos de
composição do desenho. Elabora croquis, com o auxílio de instrumentos de
medição e materiais de desenho, para fazer estudo prévio do desenho
definitivo. Elabora desenhos técnicos detalhados de estruturas,
componentes, equipamentos, sistemas e outros artefatos mecânicos de
embarcações, utilizando programas computacionais de apoio e/ou
instrumentos manuais e materiais de desenho. Utiliza programas do tipo
CAD-Desenho Assistido por Computador (Computer Aided Design) e os
recursos -- de hardware e de software -- associados a essas ferramentas.
Pode executar o desenho técnico naval a partir de informações da
representação de artefatos mecânicos de embarcações em três dimensões
(3D), utilizando ferramentas computacionais integradas de representação
em 3D e de simulações. Na execução do desenho, aplica normas para
desenho técnico e realiza os cálculos necessários à execução do desenho,
tais como cálculo das dimensões, das superfícies, do volume, do peso e
do centro de gravidade de cada peça. Pode utilizar recursos de software
para auxílio ao cálculo. Efetua cálculos para estabelecer relações entre
as diferentes partes do artefato mecânico de embarcação desenhado,
dimensionando os elementos parciais em escalas adequadas. Executa a
representação de todos os elementos do desenho com precisão técnica,
clareza e nitidez. Detalha desenho de conjunto e desenho de montagem.
Elabora lista de peças de conjunto. Confere o desenho e submete o esboço
elaborado à apreciação do autor do projeto, consultando-o sobre
possíveis correções ou alterações, para efetuar os ajustes necessários.
Elabora o desenho definitivo, registrando comentários, instruções gerais
e informações sobre materiais, processos e técnicas de fabricação e
montagem. Faz cópias do desenho, por meio de plotters, impressoras
eletrônicas ou máquinas copiadoras de diferentes tecnologias. Acompanha
fabricação e montagem dos artefatos mecânicos de embarcações desenhados,
assistindo o executante na fabricação e conferindo as peças, os
conjuntos e a montagem em relação ao desenho. Pode detectar erros,
propor modos de corrigi-los e sugerir melhorias. Elabora relatórios de
acompanhamento dos trabalhos. Pode alterar desenhos, em função de
mudanças de projeto, adequação às normas ou outras demandas. Pode
elaborar desenhos para manuais de referência e publicações técnicas a
fim de descrever a operação e a manutenção de artefatos mecânicos de
embarcações. Identifica os desenhos para arquivo. Organiza os arquivos
dos desenhos em substratos de papel ou similares e em meios eletrônicos.
Controla a distribuição de cópias dos desenhos. Conserva o local de
trabalho limpo e organizado. Controla luminosidade do ambiente e
verifica as condições ergonômicas. Destina adequadamente os resíduos
gerados. Mantém equipamentos e instrumentos de trabalho em plenas
condições de uso e funcionamento. Zela pela segurança, atuando na
prevenção de acidentes e utilizando equipamentos de proteção individual.

\begin{Form}
\textbf{Esta correspondência está adequada?} \\ 
\ChoiceMenu[radio,radiosymbol=\ding{52},bordercolor=0 0.5 0,name=cbovalid_ 318215 ]{}{CBO deve ficar=sim, \hspace{6em} CBO não fica aqui=nao}
\end{Form}

\newpage
\section{ 10009 -- Técnico em Curtimento }

\textbf{Perfil do Curso (CNCT)}

O Técnico em Curtimento será habilitado para:

Planejar, coordenar e supervisionar o processo de bene?ciamento de
couros e peles. Selecionar e executar análises laboratoriais para o
controle de qualidade. Realizar operações e processos de curtimento,
recurtimento, matização, pré-acabamento e acabamento. Desenvolver
melhorias no processo produtivo e programar a produção. Realizar ensaios
físico-químicos para o controle de qualidade da matéria-prima e do
produto acabado. Controlar estoques de produtos acabados.

Para atuação como Técnico em Curtimento, são fundamentais:

Conhecimentos e saberes relacionados aos processos de planejamento e
operação das atribuições da área, de modo a assegurar a saúde e a
segurança dos trabalhadores e dos futuros usuários e operadores de
empresas em processos de transformação em curtimento. Conhecimentos e
saberes relacionados à sustentabilidade do processo produtivo, às normas
e relatórios técnicos, à legislação da área, às novas tecnologias
relacionadas à indústria 4.0, à liderança de equipes, à solução de
problemas técnicos e à gestão de conflitos.

\textbf{CBOs correspondidos e seus perfis (presentes no CNCT, e presentes na predição da IA)}

\textbf{CBO: 311115  ( Técnico em curtimento )}

Opera e controla máquinas e equipamentos empregados nas operações e
processos de curtimento de pele, recurtimento, matização, pré-acabamento
e acabamento de couro, conforme manuais técnicos e procedimentos de
produção. Monitora o funcionamento de máquinas, equipamentos e
instrumentos de medição e controle. Mantém máquinas e equipamentos em
condições de uso, realizando regulagens, quando necessárias. Realiza
análises químicas e físicas, utilizando métodos analíticos e
instrumentais, seguindo procedimentos e normas técnicas e de qualidade e
observando critérios de confiabilidade metrológica. Coleta e prepara
amostras. Prepara reagentes. Registra os resultados, assumindo
responsabilidade técnica por eles. Supervisiona processos de produção,
definindo e coordenando equipes de trabalho, elaborando o cronograma de
produção, definindo o fluxo do processo, monitorando os parâmetros do
processo e emitindo relatórios de produção. Participa do desenvolvimento
de produtos e processos, testando insumos e matérias primas, elaborando
formulações para produção, testando novos produtos, definindo o processo
de fabricação de produtos, fazendo a adequação dos produtos às
necessidades dos clientes e validando o produto desenvolvido. Participa
da implantação de programas de qualidade, analisando indicadores,
aplicando ferramentas da qualidade e implementando ações preventivas e
corretivas. Participa de auditorias. Colabora na implantação ou
reestruturação das instalações industriais, definindo leiautes e
fluxogramas de produção, acompanhando a montagem e instalação de
equipamentos e testando máquinas e equipamentos. Presta assistência
técnica, identificando necessidades dos clientes, diagnosticando
problemas técnicos e propondo soluções e melhorias de produtos e
processos. Elabora documentação técnica, emitindo laudos e redigindo
procedimentos e relatórios. Conserva o local de trabalho limpo e
organizado. Realiza tratamento de resíduos sólidos e efluentes gerados
no curtimento e no beneficiamento, para aproveitamento econômico -- tal
como aplicação na produção de biogás e de fertilizante agrícola -- e
diminuição do impacto ambiental.

\begin{Form}
\textbf{Esta correspondência está adequada?} \\ 
\ChoiceMenu[radio,radiosymbol=\ding{52},bordercolor=0 0.5 0,name=cbovalid_ 311115 ]{}{CBO deve ficar=sim, \hspace{6em} CBO não fica aqui=nao}
\end{Form}

\textbf{CBO: 760205  ( Supervisor de curtimento )}

Planeja o processo de produção, analisando pedidos e ordens de serviços
e examinando a viabilidade de sua execução. Negocia metas de produção.
Define métodos e processos de produção; determina fluxo de produção; e
estabelece itens de controle de processo. Elabora fórmulas para o
processamento de couros. Administra recursos para a produção,
programando e controlando estoques de insumos - peles e produtos
químicos -; dimensionando equipe de trabalho; e controlando a
disponibilidade, o leiaute e a manutenção de máquinas e equipamentos,
para garantir o perfeito funcionamento e a continuidade do processo.
Dimensiona capacidade de produção com os recursos disponíveis. Analisa
especificações de peles e produtos químicos, verificando a possibilidade
de utilização de insumos menos agressivos ao ambiente - biotecnológicos
ou novos insumos -, que possam atribuir, ao produto, as características
exigidas atualmente pelo mercado. Coordena o desenvolvimento de amostras
de couro, avaliando a solicitação. Faz especificação das características
de amostras. Define e controla processo de fabricação de amostras. Testa
amostras e requisita sua aprovação. Controla os processos mecânicos ou
automatizados de transformação de pele em couro, desde a limpeza da pele
até a estabilização e o acabamento do couro. Determina padrões de
produção e orienta fluxo e movimentação de materiais. Determina e
monitora os padrões de qualidade do processo e da produção. Avalia
índice de produtos não conformes e identifica causas de não
conformidades. Administra os resultados da produção, analisando
registros e relatórios. Controla o volume e os custos da produção.
Identifica pontos críticos da produção. Soluciona problemas pós-vendas.
Supervisiona equipe de trabalho, selecionando pessoal, distribuindo
tarefas, controlando as horas trabalhadas e programando férias e folgas
dos trabalhadores. Monitora o cumprimento das normas administrativas e
de segurança do trabalho pela equipe. Avalia o desempenho profissional
dos trabalhadores, identificando as necessidades de treinamento. Treina
a equipe. Elabora documentação técnica, preparando pareceres técnicos e
relatórios. Elabora cronogramas, fichas técnicas de produção e planilhas
de custos. Emite requisições de material. Prepara ordens de serviço e
procedimentos operacionais. Elabora materiais para apresentações.
Monitora a limpeza e a organização do ambiente de trabalho. Orienta
equipe para conservar ferramentas e instrumentos limpos, organizados e
acondicionados. Providencia a manutenção preventiva e a manutenção
corretiva de máquinas e equipamentos. Pode aplicar princípios e
estimular práticas de Produção mais Limpa nos processos de curtimento do
couro, diminuindo a produção de resíduos; reutilizando banhos;
diminuindo o consumo de água; e utilizando produtos com menor
toxicidade. Monitora o cumprimento das normas de segurança,
inspecionando a utilização de equipamentos de proteção individual e
coletiva. Identifica condições de risco no trabalho e atua para
eliminá-las. Implementa programas institucionais que visam aumentar a
segurança no trabalho.

\begin{Form}
\textbf{Esta correspondência está adequada?} \\ 
\ChoiceMenu[radio,radiosymbol=\ding{52},bordercolor=0 0.5 0,name=cbovalid_ 760205 ]{}{CBO deve ficar=sim, \hspace{6em} CBO não fica aqui=nao}
\end{Form}

\textbf{CBO: 762205  ( Curtidor (couros e peles) )}

Prepara-se para realizar processo de curtimento de peles de animais de
diferentes origens - tais como bovina, suína ou caprina -, recebendo as
peles, verificando a disponibilidade de produtos químicos e ajustando
instrumentos de medida. Seleciona fulões (cilindros horizontais fechados
- normalmente de madeira, aço inox ou polipropileno -, dotados de
dispositivos para rotação em torno de seu eixo horizontal, e com porta
na lateral para carga e descarga das peles e para adição dos produtos
químicos). Pode trabalhar com equipamentos totalmente automatizados, com
dispositivos de controle de rotação e temperatura. Controla os
parâmetros do processo de curtimento para a transformação de peles em
couros, interpretando a formulação, pesando peles e produtos químicos;
medindo o volume, o pH e a temperatura da água do banho de curtimento e
verificando o tempo do processo de curtimento. Pode utilizar taninos
vegetais ou sintéticos para curtimento de peles em substituição ao
cromo. Realiza o processo de curtimento, colocando as peles, os produtos
químicos e a água no fulão e acionando os comandos eletromecânicos e
mecânicos do equipamento. Controla o tempo de processo e avalia a
absorção dos produtos químicos pelas peles. Pode, se necessário, repetir
as operações indispensáveis, para completar o curtimento. Pode operar
processos automatizados ou modificados para reduzir o consumo de água e
produtos químicos. Descarrega o fulão e acondiciona as peles curtidas
(couros) em paletes, para o transporte para as etapas subsequentes. Pode
usar esteira ou sistema de transporte contínuo para movimentação das
peles curtidas. Aplica os princípios de Produção mais Limpa, monitorando
o consumo de energia elétrica e verificando focos de desperdício, tais
como perda de vapor e isolamento inadequado de tubulações. Monitora o
consumo de água, podendo reutilizar águas de banho ou diminuir a
quantidade de água nos fulões - desde que essas medidas não interfiram
na qualidade do produto - e efetuar a manutenção das linhas de processo
para evitar vazamentos. Registra os dados, preenchendo planilhas de
controle e fichas técnicas, a fim de monitorar a produção e orientar as
etapas subsequentes ao processo de curtimento. Requisita produtos
químicos, instrumentos, ferramentas e outros materiais de consumo.
Participa da execução de manutenção produtiva total dos fulões, para
evitar paradas desnecessárias e diminuir o risco de acidente. Conserva
local de trabalho limpo e organizado. Mantém instrumentos limpos,
organizados e acondicionados. Zela pela segurança, prevenindo acidentes
e utilizando equipamentos de proteção individual e coletiva. Participa
de treinamentos da Comissão Interna de Prevenção de Acidentes (CIPA).
Atua de acordo com regras de segurança da empresa, eliminando situações
de riscos. Pode prestar primeiros socorros.

\begin{Form}
\textbf{Esta correspondência está adequada?} \\ 
\ChoiceMenu[radio,radiosymbol=\ding{52},bordercolor=0 0.5 0,name=cbovalid_ 762205 ]{}{CBO deve ficar=sim, \hspace{6em} CBO não fica aqui=nao}
\end{Form}

\textbf{CBO: 762210  ( Classificador de couros )}

Prepara-se para o processo de classificação de couros, selecionando e
regulando instrumentos de medição. Recebe peles que passaram por
processo de curtimento - transformando-se em couro -- e que foram
divididas em camada superior - lado externo e parte mais nobre - chamada
``flor'' e a camada inferior - lado interno -, ``raspa''. Registra
couros recebidos em sistema informatizado, com sua identificação e
características, como espessura de cada camada, cor e tamanho. Examina
as camadas ``flor'' e separa, em um primeiro grupo, as que apresentam
correções de defeitos causados pela natureza -- tais como bernes e
carrapatos -- e pelo homem; e problemas decorrentes de falhas no
processo de preparação ao curtimento e no processo de curtimento. Com
relação a cada camada ``flor'' desse grupo, registra sua classificação,
de acordo com padrões estabelecidos; a área da camada em que estão
localizados os problemas; e a quantidade de problemas encontrados.
Coloca, em segundo grupo, camadas ``flor'' integrais, sem defeitos
corrigidos e submetidas a processamentos -- na preparação ao curtimento
e no curtimento -- eficientes e sem falhas. Registra a classificação de
cada camada do grupo, segundo padrões estabelecidos. Avalia o índice de
aproveitamento das camadas ``flor''. Coloca as camadas, distribuídas em
lotes por classificação, em paletes, orientando seu encaminhamento para
processos subsequentes de tratamento. Pode usar esteira ou sistema de
transporte contínuo para movimentação de cada lote. Examina as camadas
``raspa'', colocando em um grupo as camadas perfeitas e, em outro, as
que apresentam defeitos corrigidos e/ou problemas decorrentes de falhas
nos processos de tratamento realizados. Com relação a cada camada
``raspa'', registra sua classificação, de acordo com padrões
estabelecidos, e detalha os problemas nela encontrados. Encaminha as
camadas ``raspa'', distribuídas em lotes por classificação, para
processo de comercialização. Envia - por sistema informatizado - as
fichas de controle, com os dados registrados, para a administração da
produção. Pode realizar classificação com a utilização de ``software'',
aumentando a precisão e a uniformização da análise. Conserva local de
trabalho limpo e organizado. Mantém instrumentos limpos, organizados e
acondicionados. Participa de ações que seguem princípios da Produção
mais Limpa, relacionadas à destinação de resíduos. Zela pela segurança,
prevenindo acidentes, eliminando condições inseguras e utilizando
equipamentos de proteção individual.

\begin{Form}
\textbf{Esta correspondência está adequada?} \\ 
\ChoiceMenu[radio,radiosymbol=\ding{52},bordercolor=0 0.5 0,name=cbovalid_ 762210 ]{}{CBO deve ficar=sim, \hspace{6em} CBO não fica aqui=nao}
\end{Form}

\textbf{CBO: 762215  ( Enxugador de couros )}

Prepara-se para enxugar couros, selecionando instrumentos de trabalho.
Registra nas fichas técnicas, em sistema informatizado, as peças de
couro recebidas. Ajusta balança e pesa os couros. Prepara a máquina,
regulando-a para o tipo de couro a ser enxugado. Abastece a máquina com
o couro. Opera a máquina, acionando e monitorando a ação combinada de
cilindros revestidos por material absorvente, que pressiona o couro para
remover o líquido. Retira o couro enxugado da máquina. Mede o grau de
umidade e o peso dos couros enxugados, utilizando instrumentos de
medição. Empilha os couros em paletes, encaminhando-os à próxima etapa
de tratamento. Pode usar esteira ou sistema de transporte contínuo para
movimentação dos couros enxugados. Aplica princípios de Produção mais
Limpa, monitorando o consumo de energia elétrica e encaminhando o
líquido extraído dos couros para recuperação e/ou reuso, conforme
processo estabelecido pela empresa. Registra, em fichas técnicas, os
pesos e os graus de umidade medidos. Encaminha as fichas técnicas, em
sistema informatizado, à administração da produção. Participa da
execução de manutenção produtiva total da máquina, para evitar paradas
no processo de enxugamento. Conserva local de trabalho limpo e
organizado. Mantém instrumentos limpos, organizados e acondicionados.
Zela pela segurança, prevenindo acidentes e utilizando equipamentos de
proteção individual e coletiva. Participa de treinamentos da Comissão
Interna de Prevenção de Acidentes (CIPA). Atua de acordo com regras de
segurança da empresa, eliminando situações de riscos. Pode prestar
primeiros socorros.

\begin{Form}
\textbf{Esta correspondência está adequada?} \\ 
\ChoiceMenu[radio,radiosymbol=\ding{52},bordercolor=0 0.5 0,name=cbovalid_ 762215 ]{}{CBO deve ficar=sim, \hspace{6em} CBO não fica aqui=nao}
\end{Form}

\textbf{CBO: 762220  ( Rebaixador de couros )}

Prepara-se para realizar o rebaixamento de couros, selecionando
instrumentos de trabalho. Registra nas fichas técnicas, em sistema
informatizado, as peças de couro recebidas. Mede a espessura dos couros.
Interpreta especificações recebidas, para tomar conhecimento da
espessura a ser atingida. Prepara a máquina, regulando-a para atingir a
espessura especificada do couro. Abastece a máquina com o couro. Opera a
máquina de rebaixamento, acionando-a para que suas lâminas helicoidais
de corte nivelem e uniformizem a espessura do couro, a fim de atender às
especificações estabelecidas. Monitora a abertura das fibras - feita
durante o processo de rebaixamento -, para uma melhor resposta do couro
durante as etapas seguintes do pós curtimento. Confere a espessura do
couro rebaixado -- que deverá prevalecer até o final do processamento -,
utilizando instrumentos de medição. Empilha os couros em paletes,
encaminhando-os à próxima etapa do tratamento. Pode usar máquinas
automáticas para empilhar couros. Pode utilizar esteira ou sistema de
transporte contínuo para movimentação dos couros, após o processo de
rebaixamento. Registra, em fichas técnicas, as espessuras medidas nas
peças de couro após o processo de rebaixamento. Encaminha as fichas
técnicas, em sistema informatizado, à administração da produção. Aplica
princípios de Produção mais Limpa, recolhendo o farelo de rebaixamento
-- com diferentes propriedades, conforme o tipo de couro -, para
destiná-lo ao reaproveitamento interno e/ou por terceiros, como
subproduto do processamento do couro. Participa da execução de
manutenção produtiva total da máquina, para evitar paradas durante o
processo de rebaixamento. Conserva local de trabalho limpo e organizado.
Mantém instrumentos limpos, organizados e acondicionados. Zela pela
segurança, prevenindo acidentes e utilizando equipamentos de proteção
individual e coletiva. Participa de treinamentos da Comissão Interna de
Prevenção de Acidentes (CIPA). Atua de acordo com regras de segurança da
empresa, eliminando situações de riscos. Pode prestar primeiros
socorros.

\begin{Form}
\textbf{Esta correspondência está adequada?} \\ 
\ChoiceMenu[radio,radiosymbol=\ding{52},bordercolor=0 0.5 0,name=cbovalid_ 762220 ]{}{CBO deve ficar=sim, \hspace{6em} CBO não fica aqui=nao}
\end{Form}

\newpage
\section{ 10011 -- Técnico em Móveis }

\textbf{Perfil do Curso (CNCT)}

O Técnico em Móveis será habilitado para:

Realizar o desenvolvimento, a fabricação e a manutenção de móveis e
esquadrias. Operar máquinas e equipamentos. Selecionar materiais,
insumos e acessórios. Planejar e implementar melhoria nos produtos e
processos relacionados a móveis. Executar regulagem e manutenção
preventiva de máquinas. Coordenar, planejar e supervisionar linhas de
produção de móveis. Controlar estoques de produtos acabados.

Para atuação como Técnico em Móveis, são fundamentais:

Conhecimentos e saberes relacionados aos processos de planejamento e
operação das atribuições da área, de modo a assegurar a saúde e a
segurança dos trabalhadores e dos futuros usuários e operadores de
empresas em processos de transformação em móveis. Conhecimentos e
saberes relacionados à sustentabilidade do processo produtivo, às normas
e relatórios técnicos, à legislação da área, às novas tecnologias
relacionadas à indústria 4.0, à liderança de equipes, à solução de
problemas técnicos e à gestão de conflitos.

\textbf{CBOs correspondidos e seus perfis (presentes no CNCT, e presentes na predição da IA)}

\textbf{CBO: 318425  ( Desenhista técnico (mobiliário) )}

Prepara as atividades, organizando o posto de trabalho de desenho,
inspecionando equipamentos e selecionando ferramentas, instrumentos e
materiais. Interpreta solicitação de desenho, reunindo e organizando
informações pertinentes. Consulta revistas, catálogos e outras
publicações técnicas. Desenvolve esboços manualmente e com auxílio de
recursos computacionais. Define os meios de representação gráfica,
empregando recursos analógicos e/ou digitais adequados à demanda.
Planeja a elaboração de desenho de móveis, ordenando o trabalho por
etapas e fixando prazo de entrega. Define escalas e estabelece o formato
de apresentação, impresso e/ou eletrônico. Consulta normas técnicas para
especificar as características do desenho. Desenvolve o desenho,
combinando elementos de imagens, conforme a necessidade, para melhor
orientar a produção do mobiliário. Realiza representação gráfica pelo
método manual, fazendo uso de instrumentos para execução do traçado
diretamente em substratos de papel, bem como utilizando computadores,
programas integrados de CAD-Desenho Assistido por Computador (Computer
Aided Design), ilustração e tratamento de imagens. Codifica o desenho e
relaciona as especificações técnicas envolvidas em sua produção. Confere
as especificações em face do trabalho realizado e examina o atendimento
às normas técnicas de representação gráfica, para imprimir e requisitar
a aprovação do desenho. Realiza as correções indicadas pelo solicitante,
registra e arquiva o desenho aprovado. Executa o acabamento final do
desenho, indicando as características de materiais de acabamento. Obtém
a aprovação final. Confecciona a matriz do desenho, preenche a legenda
do desenho e tira as cópias de segurança (backup). Conserva o local de
trabalho limpo e organizado. Controla graus de luminosidade do ambiente,
temperatura e ruídos. Verifica as condições ergonômicas para uso de
mobiliário, monitor de vídeo de computador, mouse e teclado. Requisita a
aquisição de materiais e instrumentos de trabalho, para reposição.
Conserva equipamentos, ferramentas e instrumentos de trabalho em plenas
condições de uso e funcionamento. Requisita serviço para manutenção de
computador, impressoras e outros periféricos, quando necessário.
Organiza e arquiva documentos físicos e eletrônicos, tais como
comprovantes de aprovação e cópias de segurança dos desenhos aprovados.
Destina adequadamente os resíduos gerados, fazendo reciclagem e
realizando descarte de acordo com as normas ambientais.

\begin{Form}
\textbf{Esta correspondência está adequada?} \\ 
\ChoiceMenu[radio,radiosymbol=\ding{52},bordercolor=0 0.5 0,name=cbovalid_ 318425 ]{}{CBO deve ficar=sim, \hspace{6em} CBO não fica aqui=nao}
\end{Form}

\textbf{CBO: 318805  ( Projetista de móveis )}

Planeja projetos e protótipos de móveis para produção sob encomenda ou
em série, com base em tendências de mercado; consultando área de
produção; consultando boletins técnicos; avaliando estilos de design,
concepções e produtos de concorrentes; e coletando informações
tecnológicas em feiras do mobiliário. Desenvolve projetos de móveis,
analisando locais de instalação de móveis sob medida e interpretando
desenhos e modelos. Elabora desenhos e gabaritos para a produção de
móveis, utilizando prancheta e software -- tal como o CAD-Desenho
Assistido por Computador (Computer Aided Design) -; dimensionando
componentes; especificando madeiras, derivados e acessórios; e
especificando materiais para acabamento, como tintas e vernizes. Pode
utilizar software de CAM--Manufatura Auxiliada por Computador (Computer
Aided Manufacturing), gerando programas de CNC--Comando Numérico
Computadorizado, para fabricação de mobiliário e confecção de maquetes.
Pode gerar maquetes eletrônicas de móveis, em três dimensões, a partir
de desenhos de projetos em software, tal como CAD-Desenho Assistido por
Computador (Computer Aided Design). Desenvolve protótipos de móveis,
interpretando normas técnicas para produção; elaborando ficha técnica;
especificando ferragens e acessórios; e determinando a posição de
ornamentos, detalhes e acessórios. Orienta a confecção de protótipos e
peças-piloto de móveis, realizando e requisitando testes de
laboratórios; ajustando gabaritos; e realizando ampliações e reduções
das dimensões do móvel por escala. Testa características do protótipo do
móvel, como as qualidades de segurança e manuseio, a eficiência com que
ele pode ser produzido e as maneiras de usá-lo e mantê-lo. Avalia
materiais para aquisição com vistas à produção de mobiliário,
verificando a resistência dos materiais; examinando sua adequação para
produzir móveis ergonômicos e funcionais; e avaliando sua composição e
acabamento. Interpreta resultados de testes laboratoriais, elaborando
relatórios. Avalia a relação entre custo e benefício de materiais.
Refina projetos, usando modelos de trabalho em conformidade com as
especificações do cliente, limitações de produção ou alterações de
tendências de estilo. Avalia a viabilidade do projeto, com base em
fatores como estilo, segurança, função, características do mercado,
facilidade de manutenção, orçamento e custos da produção. Pode
supervisionar e orientar equipe de trabalho na realização de
prototipagem e testes, distribuindo tarefas, coordenando e avaliando a
execução das atividades.

\begin{Form}
\textbf{Esta correspondência está adequada?} \\ 
\ChoiceMenu[radio,radiosymbol=\ding{52},bordercolor=0 0.5 0,name=cbovalid_ 318805 ]{}{CBO deve ficar=sim, \hspace{6em} CBO não fica aqui=nao}
\end{Form}

\textbf{CBO: 319205  ( Técnico do mobiliário )}

Planeja as atividades de produção de móveis e esquadrias, em série ou
sob encomenda, analisando tendências do mercado e demandas do cliente.
Elabora croquis e orienta a construção de protótipos. Define as etapas e
a sequência de trabalho, de acordo com o produto. Desenha esquemas para
a etapa de montagem. Propõe o desenvolvimento de embalagens para
proteção, acondicionamento e transporte. Seleciona máquinas,
equipamentos, ferramentas, instrumentos, matéria-prima e demais
materiais e insumos, considerando especificações técnicas, projeto e
ordem de produção. Avalia a qualidade e a relação custo-benefício de
materiais, acessórios e insumos. Estima custo do produto. Pode utilizar
o software CAD-Desenho Assistido por Computador (Computer Aided Design)
na elaboração do projeto. Pode fazer adequações em projeto, considerando
os recursos de trabalho disponíveis. Prepara e regula parâmetros
operacionais de máquinas e equipamentos, inclusive os automatizados.
Confecciona e ajusta gabaritos e moldes. Elabora programas para
operações de máquinas-ferramentas com comando numérico computadorizado
(CNC). Utiliza o software CAM--Manufatura Auxiliada por Computador
(Computer Aided Manufacturing), para que as máquinas-ferramentas
convertam o projeto -- com modelo e montagem - criado no software CAD em
peças de móveis e esquadrias. Monitora o processo de produção, regulando
máquinas e ferramentas; efetuando cálculos técnicos; controlando e
ajustando parâmetros operacionais e de medição; inspecionando o
desempenho de máquinas e equipamentos; e controlando o tempo das
operações e as demais variáveis do processo. Pode operar máquinas e
equipamentos, quando necessário. Controla as fases de corte, de montagem
dos produtos e de fixação de acessórios, considerando ordem de produção,
desenhos e projeto. Monitora a execução do acabamento de peças, móveis e
esquadrias, medindo viscosidade, gramatura, brilho, aderência e dureza
de tintas e vernizes. Inspeciona embalagens e monitora armazenamento e
transporte dos produtos. Controla estoques de produtos acabados.
Verifica a qualidade do produto e do processo, conferindo geometria e
dimensões, aplicando testes de resistência e interpretando laudos de
análises químicas e físicas de materiais e produtos. Pode aplicar
ferramentas da qualidade. Propõe e implementa melhorias no processo,
identificando falhas e apresentando soluções. Pesquisa novas tecnologias
e novos métodos de produção. Testa novas tecnologias e materiais. Estuda
adaptação de leiaute em função do processo de produção. Sugere aquisição
de máquinas e equipamentos. Presta assistência técnica, fornecendo
suporte ao cliente. Identifica defeitos no produto e ajusta ou substitui
peças e acessórios. Pode solicitar, à empresa, a substituição do produto
para o cliente. Coordena equipe de trabalho, participando de processo de
seleção e de administração de pessoal. Supervisiona o trabalho da
equipe, avaliando desempenho e levantando necessidade de treinamento.
Viabiliza a realização de programas de treinamento. Elabora e preenche
documentos, tais como relatórios, pareceres, descrição de procedimentos,
memorial descritivo, fichas de controle e dados de desempenho de
máquinas. Monitora a limpeza e a organização dos locais de trabalho.
Zela pela preservação, funcionamento e integridade dos recursos de
trabalho. Programa manutenção preventiva e acompanha a realização de
manutenção corretiva em máquinas, equipamentos, instrumentos e
ferramentas. Controla desperdícios de material, monitorando o
reaproveitamento. Identifica resíduos, providenciando sua segregação e
destinação, de acordo com normas ambientais. Trabalha com segurança,
utilizando equipamentos de proteção individual; identificando e
solucionando condições inseguras; e sinalizando áreas de risco.
Inspeciona e aciona os dispositivos de segurança dos sistemas. Fiscaliza
a utilização de equipamentos de proteção individual e coletiva pelas
equipes de trabalho.

\begin{Form}
\textbf{Esta correspondência está adequada?} \\ 
\ChoiceMenu[radio,radiosymbol=\ding{52},bordercolor=0 0.5 0,name=cbovalid_ 319205 ]{}{CBO deve ficar=sim, \hspace{6em} CBO não fica aqui=nao}
\end{Form}

\newpage
\section{ 10013 -- Técnico em Petróleo e Gás }

\textbf{Perfil do Curso (CNCT)}

O Técnico em Petróleo e Gás será habilitado para:

Operar, controlar, coordenar e monitorar processos de produção e re?no
de petróleo e gás. Programar e planejar a manutenção de máquinas e
equipamentos relacionados ao seu processo. Realizar amostragens e
caracterizações de petróleo, gás natural e derivados. Realizar
procedimento de controle de qualidade de matérias-primas, insumos e
produtos. Analisar dados estatísticos do processo produtivo e
interpretar laudos de análises químicas. Comprar e estocar
matérias-primas, produtos e insumos. Controlar estoques de produtos
acabados.

Para atuação como Técnico em Petróleo e Gás, são fundamentais:

Conhecimentos e saberes relacionados aos processos de planejamento e
operação das atribuições da área, de modo a assegurar a saúde e a
segurança dos trabalhadores e dos futuros usuários e operadores de
empresas em processos de transformação em petróleo e gás. Conhecimentos
e saberes relacionados à sustentabilidade do processo produtivo, às
normas e relatórios técnicos, à legislação da área, às novas tecnologias
relacionadas à indústria 4.0, à liderança de equipes, à solução de
problemas técnicos e à gestão de conflitos.

\textbf{CBOs correspondidos e seus perfis (presentes no CNCT, e presentes na predição da IA)}

\textbf{CBO: 301115  ( Técnico químico de petróleo )}

Organiza as condições para a realização de análises laboratoriais
químicas, físicas e físico-químicas, provendo laboratório dos materiais
de consumo, preparando materiais e equipamentos para ensaio,
selecionando substâncias reagentes, preparando e padronizando soluções
para análise. Realiza amostragem de materiais, aplicando métodos de
amostragem. Coleta, prepara e preserva amostra conforme especificações,
normas e procedimentos. Efetua descarte ou reaproveitamento da amostra,
conforme procedimentos estabelecidos. Pode receber amostras para
análise. Realiza análises laboratoriais químicas, físicas e
físico-químicas, por meio da execução de ensaios para determinar
propriedades do material analisado - como poder calorífico, acidez,
volume, peso específico, salinidade e outras -, interpretando normas
técnicas de ensaios e especificações e definindo metodologia de análise.
Faz a análise, registrando e informando os resultados. Participa no
desenvolvimento de novos produtos, levantando suas características,
auxiliando na sua especificação, realizando testes e monitorando os
resultados. Participa na avaliação de novos fornecedores. Participa no
planejamento e na execução dos projetos de perfuração, completação,
manutenção e segurança de poços de petróleo, bem como do controle das
condicionantes ambientais. Colabora no desenvolvimento de metodologias e
análises, pesquisando normas e métodos atualizados de análise, tal como
a possibilidade de automação de ensaios com uso de instrumentação
computadorizada, redes e ambientes gráficos de programação. Testa,
otimiza, valida e implementa as metodologias de análise. Elabora, testa,
revisa e padroniza procedimentos. Elabora instruções de trabalho.
Controla a qualidade da matéria-prima e de produtos acabados e em
processo, aplicando métodos estatísticos e de metrologia, além de
métodos e técnicas normalizadas de análises e ensaios. Interpreta
resultados e emite relatórios. Garante a calibração dos equipamentos,
monitorando validade de calibração, solicitando manutenção e reparo nos
equipamentos e selecionando prestadores de serviços de calibração. Pode
efetuar calibração de equipamentos em laboratórios credenciados,
aplicando métodos específicos; definindo tipo de padrão para calibração;
registrando dados; aplicando normas e critérios de aceitação de
calibração; interpretando resultados em relação ao padrão; e realizando
manutenção e reparo nos equipamentos. Participa no sistema da qualidade
da empresa, atuando no processo de melhoria contínua; atendendo aos
procedimentos definidos pelo sistema de garantia da qualidade;
atualizando normas de procedimentos internos de análise, ensaio,
processos e produtos; colaborando nas auditorias internas e externas da
qualidade; e monitorando a qualidade dos fornecedores. Zela pela
segurança, utilizando Equipamentos de Proteção Individual (EPI), atuando
na prevenção de acidentes e manuseando os materiais de análise com base
em normas de segurança. Mantém a organização, a limpeza e a higiene no
local de trabalho. Conduz análises para auxiliar no controle de emissões
do processo, aplicando procedimentos de descarte e segregação de
resíduos de laboratório; pesquisando métodos de recuperação, reciclagem
e reaproveitamento; e otimizando métodos de tratamento de resíduos
industriais, para minimizar impactos ambientais indesejáveis.

\begin{Form}
\textbf{Esta correspondência está adequada?} \\ 
\ChoiceMenu[radio,radiosymbol=\ding{52},bordercolor=0 0.5 0,name=cbovalid_ 301115 ]{}{CBO deve ficar=sim, \hspace{6em} CBO não fica aqui=nao}
\end{Form}

\textbf{CBO: 316310  ( Técnico de mineração (óleo e petróleo) )}

Participa no planejamento de atividades de exploração de petróleo e gás
natural - na terra (``onshore'') e no mar (``offshore'') -, programando
etapas da produção. Auxilia na seleção de equipe e de equipamentos; na
escolha de aplicativos a serem utilizados durante pesquisa e extração; e
na definição dos demais recursos para execução das atividades. Participa
no estabelecimento de processos de trabalho. Contribui na elaboração de
plano de custos operacionais. Participa na elaboração de projetos
ambientais. Assessora na seleção de áreas para prospecção e pesquisa de
petróleo e gás natural, interpretando fotografias aéreas; efetuando
levantamento topográfico; levantando perfis e mapeamento geológicos;
fotodocumentando rochas; e medindo direção e mergulho de rochas
aflorantes. Pode ajudar em prospecções geológicas e geofísicas, que
visem localizar jazida de petróleo ou gás natural e calcular seu
potencial. Direciona poços de petróleo conforme coordenadas de direção e
inclinação. Coordena a execução de drenagem de jazidas e controla a
produção de hidrocarbonetos. Auxilia no controle de custos e materiais.
Supervisiona sondagem e perfilagem, posicionando poços artesianos e
furos de sondagem e determinando a porosidade e a permeabilidade em
plugues testemunhos. Examina poços de sondagem, utilizando técnicas e
instrumentos específicos, para determinar pressões, temperaturas,
natureza das camadas encontradas e outros elementos. Presta apoio no
controle técnico da exploração de petróleo ou de gás natural, elaborando
padrões técnicos, procedimentos e instruções operacionais de trabalho.
Transmite instruções de caráter técnico ao pessoal de execução dos
processos e acompanha os trabalhos de perfuração de poços, para garantir
a observância dos prazos, das normas de segurança e dos padrões
estabelecidos. Monitora o fluxo de petróleo, controlando a pressão, para
mantê-lo dentro dos níveis estabelecidos e evitar danos ao material ou
aumento excessivo da temperatura do petróleo. Fiscaliza poços em
produção; tanques de armazenagem; e instalações de transporte por
tubulações. Discute condições de produção com outros setores. Coordena a
coleta de amostras, tais como de sedimentos de corrente; de solo, rocha,
hidrocarbonetos e água; de furos de sondagem; e de poços de pesquisa,
trincheiras e galerias. Prepara as amostras, homogeneizando-as,
quarteando-as e confeccionando lâminas delgadas. Prepara amostras de
petróleo nos tanques de armazenamento, eliminando resíduos e água, para
permitir a sua análise. Descreve características físico-químicas e
litológicas de amostras. Identifica, acondiciona e encaminha amostras
para realização de testes e análises. Colabora no tratamento de dados de
pesquisa, prospecção e extração de petróleo e gás, processando e
analisando dados; confeccionando mapas e seções geológicas de prospecção
e pesquisa georreferenciados; gerando e editando banco de dados de
perfis elétricos, acústicos, radioativos e geológicos. Cria modelo de
bloco baseado em modelo geológico. Aperfeiçoa métodos de extração, para
aplicação de procedimentos mais adequados na exploração de petróleo e de
gás natural. Calcula produção de poços, galerias, trincheiras e
sondagens. Auxilia na estimativa de vida útil de reservatório de
petróleo. Supervisiona a operação e a manutenção de equipamentos. Pode
supervisionar equipes de trabalho, distribuindo tarefas, avaliando
desempenho e desenvolvendo treinamentos. Controla a organização do
ambiente de trabalho e a conservação de ferramentas e instrumentos.
Monitora aplicação de procedimentos da empresa para cumprimento de
normas ambientais. Fiscaliza emissão de poluentes e nível de
assoreamento de rios. Zela pela saúde e segurança no trabalho,
monitorando o uso de equipamentos de proteção individual e a realização
de exames periódicos de saúde. Identifica áreas de risco de poços de
petróleo, propondo análise dos fatores de risco.

\begin{Form}
\textbf{Esta correspondência está adequada?} \\ 
\ChoiceMenu[radio,radiosymbol=\ding{52},bordercolor=0 0.5 0,name=cbovalid_ 316310 ]{}{CBO deve ficar=sim, \hspace{6em} CBO não fica aqui=nao}
\end{Form}

\textbf{CBO: 316325  ( Técnico de produção em refino de petróleo )}

Participa no planejamento de atividades de refino de petróleo e de
movimentação de produtos nas unidades de processamento, armazenamento e
expedição. Auxilia na seleção de equipe; na escolha de equipamentos; e
na definição de processos de trabalho. Participa na elaboração de
programas de manutenção de equipamentos e instalações. Contribui na
elaboração de plano de custos operacionais. Auxilia na elaboração de
projetos de locação de oleodutos e gasodutos. Participa na elaboração de
projetos ambientais. Participa na elaboração de padrões técnicos,
procedimentos e instruções operacionais de trabalho, discutindo
condições de produção com outros setores. Analisa, em conjunto com a
supervisão técnica, o cronograma, as especificações, os resultados de
testes e as recomendações de laboratório, a fim de definir os controles
do equipamento para produzir a quantidade prevista de produtos com a
qualidade preestabelecida. Supervisiona o processo de refino do
petróleo, operando painéis de controle; regulando as variáveis do
processo - como temperatura e pressão -; e direcionando a taxa de fluxo
do produto, de acordo com o cronograma do processo. Monitora os
processos de planta de reagente, de purificação e de beneficiamento,
inspecionando a alimentação da planta e o sistema de rejeito e
verificando o funcionamento dos medidores de entrada e saída de
produtos. Dosa produto de acordo com as especificações. Analisa a
qualidade dos produtos. Controla a movimentação de produtos nas unidades
de processamento, armazenamento e expedição, usando o conhecimento das
interconexões e capacidades do sistema. Monitora o escoamento dos
produtos derivados do petróleo por dutos. Coordena a coleta de amostras
de produtos estocados e embarcados, por meio de válvulas de sangria ou
pelo recolhimento de óleo nos tanques. Registra e compila dados
operacionais, leituras de instrumentos e resultados de análises
laboratoriais. Realiza análise dos dados. Supervisiona a operação e a
manutenção de equipamentos. Pode supervisionar equipe de trabalho,
distribuindo tarefas, avaliando desempenho e desenvolvendo treinamentos.
Monitora a organização do ambiente de trabalho e a conservação de
ferramentas e instrumentos limpos, acondicionados e em plenas condições
de uso e funcionamento. Controla a aplicação de procedimentos da empresa
para cumprimento de normas ambientais. Fiscaliza a emissão de poluentes.
Inspeciona a quantidade de rejeitos, acompanhando testes de produtos
biodegradáveis. Representa departamentos em inspeção ambiental. Zela
pela saúde e segurança da equipe de trabalho, monitorando o uso de
equipamentos de proteção individual e a realização de exames periódicos
de saúde. Identifica áreas de risco, tomando medidas para seu isolamento
e para análise dos fatores de risco. Registra ocorrências de acidentes
de trabalho.

\begin{Form}
\textbf{Esta correspondência está adequada?} \\ 
\ChoiceMenu[radio,radiosymbol=\ding{52},bordercolor=0 0.5 0,name=cbovalid_ 316325 ]{}{CBO deve ficar=sim, \hspace{6em} CBO não fica aqui=nao}
\end{Form}

\newpage
\section{ 10014 -- Técnico em Petroquímica }

\textbf{Perfil do Curso (CNCT)}

O Técnico em Petroquímica será habilitado para:

Coletar amostras e realizar análises químicas e físico-químicas de
processos e produtos petroquímicos. Avaliar e controlar a qualidade de
matérias-primas, insumos e produtos. Controlar estoques de produtos
acabados. Operar, monitorar e controlar processos químicos,
petroquímicos e de refino de petróleo.

Para atuação como Técnico em Petroquímica, são fundamentais:

Conhecimentos e saberes relacionados aos processos de planejamento e
operação das atribuições da área, de modo a assegurar a saúde e a
segurança dos trabalhadores e dos futuros usuários e operadores de
empresas em processos de transformação em petroquímica. Conhecimentos e
saberes relacionados à sustentabilidade do processo produtivo, às normas
e relatórios técnicos, à legislação da área, às novas tecnologias
relacionadas à indústria 4.0, à liderança de equipes, à solução de
problemas técnicos e à gestão de conflitos.

\textbf{CBOs correspondidos e seus perfis (presentes no CNCT, e presentes na predição da IA)}

\textbf{CBO: 311205  ( Técnico em petroquímica )}

Programa atividades de produção petroquímica, identificando tarefas e
ações necessárias ao cumprimento do planejamento. Define meios e
recursos para o cumprimento de metas de produção e de custos,
monitorando estoques de materiais e de insumos e programando parada e
partida de equipamentos. Elabora, controla e revisa cronograma de
atividades operacionais. Programa serviços complementares, tais como
pinturas e limpeza. Elabora plano de análise contingencial para os casos
de quebra de equipamentos e paradas não programadas. Interpreta conjunto
de documentos técnicos - especificações, desenhos, dentre outros --
referentes à execução das atividades. Coordena o processo de produção
petroquímica, consultando relatórios e ocorrências do turno anterior
para dar sequência às atividades. Monitora a realização da produção,
analisando as variáveis do processo. Controla a execução de manobras
operacionais, efetuando ações de ajustes e solucionando problemas.
Inspeciona serviços contratados vinculados à área de produção, podendo
requisitar outros, complementares. Pode interagir com clientes internos
e externos durante a realização de manobras, operando recursos de
comunicação. Controla a qualidade, com o propósito de evitar produtos
fora de especificação. Ordena insumos - intermediários e finais - e
produtos, rastreando-os para identificar falhas. Coleta amostras de
produtos, realizando ensaios qualitativos, quantitativos e instrumentais
para identificar produtos não conformes, com falhas e fora de padrão.
Monitora o cumprimento do plano de análises laboratoriais, interpretando
laudos de análises químicas. Analisa dados estatísticos do processo
produtivo e controla índices técnicos relacionados à produção e aos
custos. Elabora relatórios e boletins de ocorrências, podendo apresentar
propostas para o desenvolvimento de novos produtos, aplicação de métodos
alternativos e implantação de novos sistemas. Administra a equipe de
trabalho, participando da seleção do pessoal; distribuindo tarefas;
definindo período de férias; e estabelecendo trocas de pessoal. Divulga
normas e procedimentos para a equipe. Elabora instruções operacionais
para a execução dos trabalhos. Supervisiona a equipe de trabalho,
avaliando seu desempenho e identificando a necessidade de treinamento.
Planeja o desenvolvimento da equipe, preparando material didático e
instrucional e selecionando técnica de avaliação dos programas.
Desenvolve programação de treinamento no ambiente de trabalho e avalia
seu impacto no desempenho da equipe. Propõe à equipe discussões sobre
problemas operacionais e instrui cada integrante sobre a importância da
utilização de equipamentos de proteção individual e coletiva. Mantém os
equipamentos em condições operacionais, monitorando seu funcionamento e
providenciando serviços de manutenção de primeiro nível. Preenche ordem
de serviço, relacionando os equipamentos que necessitam serviços de
manutenção corretiva. Faz especificação de materiais, ferramentas,
instrumentos e equipamentos para aquisição e reposição de estoque.
Inspeciona a limpeza e a organização do local de trabalho. Examina a
limpeza, o acondicionamento e a conservação de ferramentas e
instrumentos. Encaminha instrumentos para manutenção, aferindo-os após a
realização do serviço. Monitora o descarte de resíduos, efluentes e
gases conforme normas ambientais. Zela pela segurança, fazendo uso dos
equipamentos de proteção individual. Elabora procedimentos em
conformidade com as normas de segurança, realizando sua divulgação e
controlando seu cumprimento pelas equipes de trabalho. Identifica e
analisa situações de risco em conjunto com outros profissionais,
participando na elaboração, na implementação e na avaliação de ações
corretivas. Emite relatório de acidentes e incidentes, fazendo a
investigação de cada ocorrência em conjunto com profissionais das áreas
envolvidas.

\begin{Form}
\textbf{Esta correspondência está adequada?} \\ 
\ChoiceMenu[radio,radiosymbol=\ding{52},bordercolor=0 0.5 0,name=cbovalid_ 311205 ]{}{CBO deve ficar=sim, \hspace{6em} CBO não fica aqui=nao}
\end{Form}

\textbf{CBO: 316325  ( Técnico de produção em refino de petróleo )}

Participa no planejamento de atividades de refino de petróleo e de
movimentação de produtos nas unidades de processamento, armazenamento e
expedição. Auxilia na seleção de equipe; na escolha de equipamentos; e
na definição de processos de trabalho. Participa na elaboração de
programas de manutenção de equipamentos e instalações. Contribui na
elaboração de plano de custos operacionais. Auxilia na elaboração de
projetos de locação de oleodutos e gasodutos. Participa na elaboração de
projetos ambientais. Participa na elaboração de padrões técnicos,
procedimentos e instruções operacionais de trabalho, discutindo
condições de produção com outros setores. Analisa, em conjunto com a
supervisão técnica, o cronograma, as especificações, os resultados de
testes e as recomendações de laboratório, a fim de definir os controles
do equipamento para produzir a quantidade prevista de produtos com a
qualidade preestabelecida. Supervisiona o processo de refino do
petróleo, operando painéis de controle; regulando as variáveis do
processo - como temperatura e pressão -; e direcionando a taxa de fluxo
do produto, de acordo com o cronograma do processo. Monitora os
processos de planta de reagente, de purificação e de beneficiamento,
inspecionando a alimentação da planta e o sistema de rejeito e
verificando o funcionamento dos medidores de entrada e saída de
produtos. Dosa produto de acordo com as especificações. Analisa a
qualidade dos produtos. Controla a movimentação de produtos nas unidades
de processamento, armazenamento e expedição, usando o conhecimento das
interconexões e capacidades do sistema. Monitora o escoamento dos
produtos derivados do petróleo por dutos. Coordena a coleta de amostras
de produtos estocados e embarcados, por meio de válvulas de sangria ou
pelo recolhimento de óleo nos tanques. Registra e compila dados
operacionais, leituras de instrumentos e resultados de análises
laboratoriais. Realiza análise dos dados. Supervisiona a operação e a
manutenção de equipamentos. Pode supervisionar equipe de trabalho,
distribuindo tarefas, avaliando desempenho e desenvolvendo treinamentos.
Monitora a organização do ambiente de trabalho e a conservação de
ferramentas e instrumentos limpos, acondicionados e em plenas condições
de uso e funcionamento. Controla a aplicação de procedimentos da empresa
para cumprimento de normas ambientais. Fiscaliza a emissão de poluentes.
Inspeciona a quantidade de rejeitos, acompanhando testes de produtos
biodegradáveis. Representa departamentos em inspeção ambiental. Zela
pela saúde e segurança da equipe de trabalho, monitorando o uso de
equipamentos de proteção individual e a realização de exames periódicos
de saúde. Identifica áreas de risco, tomando medidas para seu isolamento
e para análise dos fatores de risco. Registra ocorrências de acidentes
de trabalho.

\begin{Form}
\textbf{Esta correspondência está adequada?} \\ 
\ChoiceMenu[radio,radiosymbol=\ding{52},bordercolor=0 0.5 0,name=cbovalid_ 316325 ]{}{CBO deve ficar=sim, \hspace{6em} CBO não fica aqui=nao}
\end{Form}

\newpage
\section{ 10015 -- Técnico em Planejamento e Controle da Produção }

\textbf{Perfil do Curso (CNCT)}

O Técnico em Planejamento e Controle da Produção será habilitado para:

Empregar métodos de planejamento, programação e controle na produção
industrial, preservando os requisitos de qualidade e de consumo, de
acordo com normas, padrões e especificações dos produtos. Monitorar os
insumos e suprimentos necessários de produção, analisando os estoques de
materiais e as dinâmicas de reabastecimento com base no just in time.
Reconhecer plano mestre de produção e planejamento de capacidade de uma
linha de produção. Utilizar tecnologias para administrar os recursos
fabris e melhorar a eficiência dos processos produtivos. Avaliar
indicadores estratégicos de produção quanto ao atendimento dos objetivos
organizacionais e para identificação de causas de falhas e desvios.
Reconhecer as técnicas de controle da produção utilizadas pela filosofia
Lean Manufacturing.

Para atuação como Técnico em Planejamento e Controle da Produção, são
fundamentais:

Conhecimentos e saberes relacionados aos processos de planejamento e
operação das atribuições da área, de modo a assegurar a saúde e a
segurança dos trabalhadores e dos futuros usuários e operadores de
indústrias e serviços do setor da logística. Conhecimentos e saberes
relacionados à sustentabilidade do processo produtivo, às normas e
relatórios técnicos, às novas tecnologias relacionadas à indústria 4.0,
à liderança de equipes, à solução de problemas técnicos e à gestão de
conflitos.

\textbf{CBOs correspondidos e seus perfis (presentes no CNCT, e presentes na predição da IA)}

\textbf{CBO: 391125  ( Técnico de planejamento de produção )}

Planeja a produção, a partir da análise da demanda e da capacidade de
produção. Realiza a programação da produção, analisando ordens de
produção, determinando prioridades, definindo cronograma e estabelecendo
a sequência de produção. Estabelece parâmetros para controle dos
processos de produção. Controla o desenvolvimento dos planos de
produção, com base em parâmetros de controle e monitoramento do fluxo
produtivo. Identifica desvios e define ações corretivas. Propõe
melhorias no processo de produção. Define estoque de segurança de
suprimentos necessários para atender às demandas de produção.
Providencia entrega e distribuição de suprimentos para agilizar o fluxo
de materiais e cumprir os cronogramas de produção. Pode revisar os
cronogramas de produção por falta de trabalhadores, atraso de entrega de
material ou outros problemas. Trata as informações do ciclo
planejamento-produção, elaborando gráficos e relatórios de controle.
Disponibiliza e atualiza informações. Interage com profissionais
envolvidos nas diversas etapas de produção, com clientes e com
fornecedores, para coordenar as atividades de produção, para resolver
reclamações ou para eliminar atrasos. Auxilia na definição das
instruções de trabalho.

\begin{Form}
\textbf{Esta correspondência está adequada?} \\ 
\ChoiceMenu[radio,radiosymbol=\ding{52},bordercolor=0 0.5 0,name=cbovalid_ 391125 ]{}{CBO deve ficar=sim, \hspace{6em} CBO não fica aqui=nao}
\end{Form}

\newpage
\section{ 10016 -- Técnico em Plásticos }

\textbf{Perfil do Curso (CNCT)}

O Técnico em Plásticos será habilitado para:

Planejar, operar, controlar, coordenar e monitorar o processo de
fabricação de produtos de plástico e de reciclagem. Supervisionar a
aquisição de matéria-prima e controlar a qualidade do produto acabado.
Realizar ensaios físicos. Identi?car a composição do material de
produtos acabados. Elaborar o dimensionamento das necessidades da
instalação industrial. Controlar estoques de produtos acabados.

Para atuação como Técnico em Plásticos, são fundamentais:

Conhecimentos e saberes relacionados aos processos de planejamento e
operação das atribuições da área, de modo a assegurar a saúde e a
segurança dos trabalhadores e dos futuros usuários e operadores de
empresas em processos de transformação em polímeros. Conhecimentos e
saberes relacionados à sustentabilidade do processo produtivo, às normas
e relatórios técnicos, à legislação da área, às novas tecnologias
relacionadas à indústria 4.0, à liderança de equipes, à solução de
problemas técnicos e à gestão de conflitos.

\textbf{CBOs correspondidos e seus perfis (presentes no CNCT, e presentes na predição da IA)}

\textbf{CBO: 311405  ( Técnico em borracha )}

Elabora programação da fabricação de produtos de borracha, interagindo
com setores do estabelecimento industrial para avaliar e definir a
capacidade de produção. Define o cronograma de produção e programa o
consumo de matéria-prima. Realiza ajustes na programação da produção,
quando necessário. Supervisiona o processo de fabricação de produtos de
borracha natural e sintética, organizando o fluxo de produção e
controlando o processo de mistura das matérias-primas. Monitora
parâmetros do processo e soluciona problemas da produção. Aplica
ferramentas estatísticas de controle de processo. Otimiza parâmetros de
processo e etapas de produção. Inspeciona a qualidade dos produtos
intermediários e acabados. Controla estoques de produtos acabados. Pode
preparar, regular e operar máquinas e equipamentos industriais --
inclusive para demonstração e orientação aos operadores - na fabricação
de produtos de borracha. Realiza ensaios físico-químicos, coletando e
preparando amostras; mantendo equipamentos calibrados; operando
instrumentos de medição e controle; e executando ensaios em
matérias-primas e produtos intermediários e acabados. Supervisiona
equipe de trabalho, avaliando seu desempenho e identificando a
necessidade de treinamento. Desenvolve programação de treinamento e
avalia o desempenho dos participantes. Desenvolve produto, interpretando
normas técnicas, desenhos e projetos. Define as características do
produto, adequando-o às necessidades do cliente e ao processo produtivo.
Identifica materiais e aplicações. Desenvolve protótipo. Avalia a
viabilidade técnico-econômica do produto, otimizando custos. Desenvolve
processos de produção, definindo etapas de produção; avaliando condições
de utilização de equipamentos e ferramental; aplicando ferramentas de
melhoria e desenvolvimento de processo; e simulando processos,
utilizando softwares de informática. Define arranjo físico dos
equipamentos e executa testes iniciais (``tryout'') de ferramentas.
Elabora documentação sobre parâmetros de processo. Treina operadores de
processo. Pode projetar ferramentas e dispositivos, em função do
processo e do produto. Elabora ou interpreta desenho técnico e seleciona
materiais para construção da ferramenta e dispositivo. Aplica tecnologia
de processos de usinagem. Usa instrumentos de medição e aplica
princípios de termologia. Desenvolve fornecedores de produtos e
serviços, avaliando seu potencial técnico, seu sistema produtivo e seu
sistema de controle de qualidade; e prestando-lhes orientação e apoio
técnico. Acompanha o desempenho de atendimento dos fornecedores,
estabelecendo prazo de entrega e verificando seu cumprimento. Presta
assistência técnica a clientes, avaliando suas solicitações e planejando
visitas técnicas, para solucionar os problemas identificados. Treina
clientes no uso dos produtos adquiridos. Realiza acompanhamento técnico
do produto no campo, coletando informações sobre problemas apontados
pelos clientes. Elabora relatórios técnicos, propondo melhorias no
produto a partir das informações obtidas. Participa das ações do sistema
de gestão da qualidade, elaborando e implementando procedimentos e
aplicando ferramentas do sistema da qualidade. Levanta e registra dados
de avaliação do nível de satisfação do cliente. Controla documentação do
sistema. Participa das ações do sistema de gestão ambiental, aplicando
os procedimentos para atendimento às normas ambientais. Colabora com a
avaliação dos impactos ambientais da atividade de produção. Define
materiais para reciclagem. Monitora o descarte de resíduos, de acordo
com diretrizes de preservação ambiental. Mantém-se atualizado na sua
área tecnológica de atuação. Monitora a organização e a limpeza do local
de trabalho. Providencia serviços de manutenção corretiva e preventiva
de máquinas e equipamentos. Zela pela segurança no trabalho, utilizando
e orientando o uso de equipamentos de proteção coletiva e individual.
Participa da elaboração do mapa de risco. Pode prestar primeiros
socorros.

\begin{Form}
\textbf{Esta correspondência está adequada?} \\ 
\ChoiceMenu[radio,radiosymbol=\ding{52},bordercolor=0 0.5 0,name=cbovalid_ 311405 ]{}{CBO deve ficar=sim, \hspace{6em} CBO não fica aqui=nao}
\end{Form}

\textbf{CBO: 311410  ( Técnico em plástico )}

Elabora programação da fabricação de produtos de plástico, interagindo
com setores do estabelecimento industrial para avaliar e definir a
capacidade de produção. Estabelece o cronograma de produção e programa o
consumo de matéria-prima. Realiza ajustes na programação da produção,
quando necessário. Supervisiona o processo da fabricação de produtos de
plástico, organizando o fluxo de produção e controlando o processo de
mistura das matérias-primas. Monitora parâmetros do processo e soluciona
problemas de produção. Aplica ferramentas estatísticas de controle de
processo. Otimiza parâmetros de processo e etapas de produção.
Inspeciona a qualidade dos produtos intermediários e acabados. Controla
estoques de produtos acabados. Pode preparar, regular e operar máquinas
e equipamentos industriais - inclusive para demonstração e orientação
aos operadores -, na fabricação e na reciclagem de produtos plásticos.
Realiza ensaios físico-químicos, coletando e preparando amostras;
mantendo equipamentos calibrados; operando instrumentos de medição e
controle; e executando ensaios em matérias-primas e produtos
intermediários e acabados. Supervisiona equipe de trabalho, avaliando
seu desempenho e identificando a necessidade de treinamento. Desenvolve
programação de treinamento e avalia o desempenho dos participantes.
Desenvolve produto, interpretando normas técnicas, desenhos e projetos.
Define as características do produto, adequando-o às necessidades do
cliente e ao processo produtivo. Identifica materiais e aplicações.
Desenvolve protótipos. Avalia a viabilidade técnico-econômica do
produto, otimizando custos. Desenvolve processos de produção, definindo
etapas de produção; avaliando condições de utilização de equipamentos e
ferramental; aplicando ferramentas de melhoria e desenvolvimento de
processo; e simulando processos, com uso de softwares. Define arranjo
físico dos equipamentos e executa testes iniciais (``tryout'') de
ferramentas. Elabora documentação sobre os parâmetros de processo.
Treina operadores de processo. Pode projetar ferramentas e dispositivos,
em função do processo e do produto. Elabora ou interpreta desenho
técnico e seleciona materiais para construção da ferramenta e
dispositivo. Aplica tecnologia de processos de usinagem. Usa
instrumentos de medição e aplica princípios de termologia. Desenvolve
fornecedores de produtos e serviços, avaliando seu potencial técnico,
seu sistema produtivo e seu sistema de controle de qualidade; e
prestando-lhes orientação e apoio técnico. Acompanha o desempenho dos
fornecedores, estabelecendo prazo de entrega e verificando seu
cumprimento. Presta assistência técnica a clientes, avaliando suas
solicitações e planejando visitas técnicas, para solucionar os problemas
identificados. Treina clientes no uso dos produtos adquiridos. Realiza
acompanhamento técnico do produto no campo, coletando informações sobre
problemas apontados pelos clientes. Elabora relatórios técnicos,
propondo melhorias no produto a partir das informações obtidas.
Participa das ações do sistema de gestão da qualidade, elaborando e
implementando procedimentos e aplicando ferramentas do sistema da
qualidade. Levanta e registra dados de avaliação do nível de satisfação
do cliente. Controla documentação do sistema. Participa das ações do
sistema de gestão ambiental, aplicando os procedimentos para atendimento
às normas ambientais. Colabora com a avaliação dos impactos ambientais
da atividade de produção. Define materiais para reciclagem. Monitora o
descarte de resíduos, de acordo com diretrizes de preservação ambiental.
Mantém-se atualizado na sua área tecnológica de atuação. Monitora a
organização e a limpeza do local de trabalho. Providencia serviços de
manutenção corretiva e preventiva de máquinas e equipamentos. Zela pela
segurança no trabalho, utilizando e orientando o uso de equipamentos de
proteção coletiva e individual. Participa da elaboração do mapa de
risco. Pode prestar primeiros socorros.

\begin{Form}
\textbf{Esta correspondência está adequada?} \\ 
\ChoiceMenu[radio,radiosymbol=\ding{52},bordercolor=0 0.5 0,name=cbovalid_ 311410 ]{}{CBO deve ficar=sim, \hspace{6em} CBO não fica aqui=nao}
\end{Form}

\newpage
\section{ 10017 -- Técnico em Processamento da Madeira }

\textbf{Perfil do Curso (CNCT)}

O Técnico em Processamento da Madeira será habilitado para:

Realizar processos de tratamento da madeira. Analisar e elaborar
programas de secagem e preservação da madeira. Operar máquinas de
usinagem de madeira para a execução de projetos. Utilizar técnicas de
acabamento e montagem de produtos de madeira. Programar e controlar a
produção atendendo as normas e padrões técnicos, de qualidade, saúde e
segurança e de meio ambiente. Elaborar documentação técnica de
processos.

Para atuação como Técnico em Processamento da Madeira, são fundamentais:

Conhecimentos e saberes relacionados aos processos de planejamento e
operação, de modo a assegurar a saúde e a segurança dos trabalhadores e
dos futuros usuários e operadores de empresas de obtenção e
processamento de madeira. Conhecimentos e saberes relacionados: à
sustentabilidade do processo produtivo; às normas e relatórios técnicos;
às leis; às novas tecnologias relacionadas à indústria 4.0; à liderança
de equipes; à solução de problemas técnicos; e à gestão de conflitos.

\textbf{CBOs correspondidos e seus perfis (presentes no CNCT, e presentes na predição da IA)}

\textbf{CBO: 321205  ( Técnico em madeira )}

Planeja atividades de colheita florestal e uso racional da madeira.
Elabora planilha de custos e cronograma de execução física e operacional
de projetos. Dimensiona o quadro de pessoal e prevê infraestrutura de
instalações e equipamentos para as atividades. Administra unidades de
colheita florestal e de processamento de produtos madeireiros,
projetando o volume de extração de madeira e de fabricação de produtos e
subprodutos. Estabelece procedimentos e padrões de qualidade, para a
extração e a produção. Supervisiona a execução de atividades florestais,
monitorando a correta identificação das espécies florestais e orientando
a colheita de madeira. Controla a cubagem e a etiquetagem da origem dos
lotes de madeira bruta, seguindo as normas legais. Controla os
parâmetros de secagem e tratamento da madeira -- produção intermediária
- em estufas ou ao ar livre. Supervisiona a execução de projetos de
fabricação de produtos e subprodutos de madeira, coordenando processos
de usinagem -- ou operando máquinas-ferramenta - e utilizando técnicas
de montagem e acabamento de produtos. Pode usar máquinas automatizadas
ou que incorporem outras novas tecnologias. Controla qualidade da
produção. Utiliza galpões e outras instalações para armazenar madeira
bruta e produtos madeireiros. Providencia acessos e vias para escoamento
da produção, definindo a logística de carregamento e transporte de
madeira bruta e de produtos madeireiros. Coordena equipes de trabalho,
participando de processo de seleção e de administração de pessoal.
Distribui tarefas e programa férias das equipes. Controla conflitos nas
equipes. Supervisiona equipes de trabalho, avaliando desempenho e
levantando as necessidades de treinamento. Viabiliza a realização de
programas de aperfeiçoamento e atualização profissional. Orienta o uso
de máquinas e implementos específicos para a atividade florestal e para
o processamento da madeira. Participa da administração dos custos e
recursos, controlando estoques e monitorando a execução financeira das
atividades. Realiza extensão florestal, sistematizando informações
socioeconômicas da comunidade; demonstrando viabilidade econômica de
produtos e métodos de produção; e assessorando criação de cooperativas.
Ministra programas de treinamento e palestras para pessoas de
comunidades e empresas, orientando-as sobre utilização racional de
recursos naturais renováveis e sobre legislação ambiental. Pode atuar em
pesquisa, colaborando na realização de experimentos técnico-científicos.
Pode participar de levantamentos fitossociológicos e fitossanitários.
Elabora documentação técnica -- tais como relatórios de atividades
operacionais - e mantém arquivos -- físicos e digitais - de inventários
florestais, cadastros, documentos fiscais, entre outros documentos.
Providencia regularizações e autorizações de instâncias superiores e de
órgãos competentes, com relação aos certificados de origem. Monitora a
limpeza e a organização dos locais de trabalho. Zela pela manutenção de
máquinas e equipamentos, elaborando plano de manutenção preventiva e
preditiva, providenciando manutenção de primeiro nível e solicitando
serviços de manutenção corretiva. Trabalha com segurança, fazendo uso
dos equipamentos de proteção individual. Fiscaliza a utilização de
equipamentos de proteção individual e coletiva pelas equipes de
trabalho. Sinaliza áreas de risco, controlando o fluxo de pessoas.
Providencia meios para controle de incêndios. Presta primeiros socorros
e providencia atendimento médico de urgência, quando necessário.

\begin{Form}
\textbf{Esta correspondência está adequada?} \\ 
\ChoiceMenu[radio,radiosymbol=\ding{52},bordercolor=0 0.5 0,name=cbovalid_ 321205 ]{}{CBO deve ficar=sim, \hspace{6em} CBO não fica aqui=nao}
\end{Form}

\newpage
\section{ 10018 -- Técnico em Processos Gráficos }

\textbf{Perfil do Curso (CNCT)}

O Técnico em Processos Gráficos será habilitado para:

Planejar, coordenar, controlar e realizar serviços de produção grá?ca.
Preparar matrizes de impressão. Reconhecer os diferentes processos de
impressão. Operar e controlar os diferentes processos de impressão de
acordo com o projeto gráfico. Ajustar e operar máquinas de
pós-impressão. Controlar a qualidade do material impresso. Analisar e
avaliar as características de matérias-primas e dos produtos. Controlar
estoques de produtos acabados.

Para atuação como Técnico em Processos Gráficos, são fundamentais:

Conhecimentos e saberes relacionados aos processos de planejamento e
operação das atribuições da área, de modo a assegurar a saúde e a
segurança dos trabalhadores e dos futuros usuários e operadores de
empresas em processos de transformação em processos gráficos.
Conhecimentos e saberes relacionados à sustentabilidade do processo
produtivo, às normas e relatórios técnicos, à legislação da área, às
novas tecnologias relacionadas à indústria 4.0, à liderança de equipes,
à solução de problemas técnicos e à gestão de conflitos.

\textbf{CBOs correspondidos e seus perfis (presentes no CNCT, e presentes na predição da IA)}

\textbf{CBO: 318405  ( Desenhista técnico (artes gráficas) )}

Prepara as atividades, organizando o posto de trabalho de desenho,
inspecionando equipamentos e selecionando ferramentas, instrumentos e
materiais. Interpreta solicitação de desenho, reunindo e organizando
informações pertinentes. Consulta revistas, catálogos e outras
publicações técnicas. Desenvolve esboços manualmente e com auxílio de
recursos computacionais. Define os meios de representação gráfica,
empregando recursos analógicos e/ou digitais adequados à demanda.
Planeja a elaboração do desenho, ordenando o trabalho por etapas e
fixando prazo de entrega. Define escalas e estabelece o formato da
apresentação. Consulta normas técnicas de desenho e de artes gráficas,
para especificar características do desenho. Elabora desenhos de
pré-impressão, combinando figuras, ilustrações, gráficos, textos e
outras reproduções visíveis, para compor a imagem a ser impressa.
Realiza representações gráficas pelo método manual, fazendo uso de
instrumentos para execução do traçado diretamente em substratos de
papel, bem como utilizando computadores, programas integrados de
CAD-Desenho Assistido por Computador (Computer Aided Design),
ilustração, tratamento de imagens e editoração gráfica. Codifica os
desenhos e relaciona as especificações técnicas envolvidas em sua
produção. Confere as especificações em face do trabalho realizado e
examina o atendimento às normas técnicas de representação gráfica, para
imprimir e requisitar a aprovação dos desenhos. Realiza as correções
indicadas pelo solicitante. Registra e arquiva o desenho aprovado.
Executa o acabamento final dos desenhos, indicando as características de
materiais de acabamento. Obtém a aprovação final. Confecciona a matriz
do desenho, preenche a legenda e tira as cópias de segurança (backup).
Conserva o local de trabalho limpo e organizado. Controla graus de
luminosidade do ambiente, temperatura e ruídos. Verifica as condições
ergonômicas para uso de mobiliário, monitor de vídeo de computador,
mouse e teclado. Requisita a aquisição de materiais e instrumentos de
trabalho, para reposição. Conserva equipamentos, ferramentas e
instrumentos de trabalho em plenas condições de uso e funcionamento.
Requisita serviço para manutenção de computador, impressoras, plotters e
outros periféricos, quando necessário. Organiza e arquiva documentos
físicos e eletrônicos, tais como comprovantes de aprovação e cópias de
segurança dos desenhos aprovados. Destina adequadamente os resíduos
gerados, fazendo reciclagem e realizando descarte de acordo com as
normas ambientais.

\begin{Form}
\textbf{Esta correspondência está adequada?} \\ 
\ChoiceMenu[radio,radiosymbol=\ding{52},bordercolor=0 0.5 0,name=cbovalid_ 318405 ]{}{CBO deve ficar=sim, \hspace{6em} CBO não fica aqui=nao}
\end{Form}

\textbf{CBO: 371305  ( Técnico em programação visual )}

Planeja o serviço de programação visual, analisando ordem de serviço e
identificando necessidades de cliente. Estabelece prazo para elaboração
do projeto. Faz estimativa de custos, para aprovação de recursos que
poderão ser utilizados no projeto. Prepara especificação de recursos
materiais. Examina as condições dos equipamentos, testando seu
funcionamento. Confere a disponibilidade de materiais, requisitando os
faltantes. Seleciona ``softwares'' específicos para elaboração de
projeto. Prepara esboços e leiautes de amostras, para visualizar a
posição de imagens, blocos de textos e outros elementos de página.
Seleciona estilos e tamanho de fontes. Localiza imagens ou fotos de
produtos. Pesquisa conceitos de ``design'' a serem aplicados. Estabelece
padrões de referência de qualidade. Elabora o projeto de programação
visual de materiais -- publicitários, jornalísticos, entre outros -- a
serem impressos, aplicando elementos criativos e estéticos. Desenvolve
elementos de identidade visual. Cria ilustrações e desenha logotipos.
Efetua a editoração de textos e o tratamento de imagens, digitando,
diagramando, ilustrando e formatando. Elabora tabelas e gráficos. Em
publicações editoriais, define o tamanho da lombada e confecciona
boneco, miolo e capa. Prepara originais, verificando a qualidade, a
quantidade e a organização de textos e imagens. Confecciona prova
digital ou peça-piloto. Analisa a viabilidade econômica do projeto.
Solicita, a cliente, a aprovação da arte-final. Pode operar máquinas e
equipamentos ou executar processos eletrônicos, em testes de
processamento e reprodução de textos e imagens. Faz ajustes técnicos --
tais como de resolução, de espessura de linhas, de extensão das imagens,
dentre outros - da programação visual, antes do processo de impressão.
Realiza o fechamento do arquivo com a arte-final. Encaminha, junto com o
arquivo, ficha de impressão, que detalha informações técnicas - tais
como tipo e gramatura do papel; número de páginas; quantidade de
impresso; entre outras - para a produção gráfica. Presta assistência
técnica para as diversas áreas da empresa, fornecendo informações
técnicas. Avalia resultados da produção realizada junto a clientes.
Analisa reclamações e informa a clientes sobre possibilidades e
restrições de materiais e processos. Busca alternativas para melhorias
em tecnologias e matérias-primas. Realiza visitas técnicas a empresas e
fornecedores, para coleta de informações e análise de tecnologias
disponíveis. Consulta normas e fontes de informações técnicas. Avalia
tendências de mercado. Pesquisa ``sites'' especializados da internet e
verifica lançamentos em feiras nacionais e internacionais da indústria
gráfica, para manter-se atualizado em relação às inovações e à
introdução de novas tecnologias e de novos materiais na sua área de
atuação. Coordena equipe de trabalho, definindo atribuições e
assegurando o cumprimento de normas e procedimentos no trabalho. Pode
supervisionar o trabalho de equipe, avaliando desempenho e realizando
programas de treinamento. Monitora a organização, a classificação e o
arquivo de documentos físicos e digitais. Verifica a limpeza e a
organização do ambiente de trabalho. Orienta a equipe para conservar
instrumentos limpos, acondicionados e em plenas condições de uso e
funcionamento. Mantém máquinas e equipamentos em condições operacionais,
providenciando manutenção de primeiro nível. Requisita serviços de
manutenção corretiva de máquinas e equipamentos. Controla a aplicação de
procedimentos da empresa para cumprimento das normas ambientais.
Inspeciona o reaproveitamento de sobras de material e o descarte de
resíduos. Zela pela segurança, utilizando e monitorando o uso de
equipamentos de proteção individual. Identifica situações de risco e
colabora para sua eliminação.

\begin{Form}
\textbf{Esta correspondência está adequada?} \\ 
\ChoiceMenu[radio,radiosymbol=\ding{52},bordercolor=0 0.5 0,name=cbovalid_ 371305 ]{}{CBO deve ficar=sim, \hspace{6em} CBO não fica aqui=nao}
\end{Form}

\textbf{CBO: 371310  ( Técnico gráfico )}

Participa no planejamento de atividades, elaborando programação para
desenvolvimento da produção gráfica. Organiza arranjo físico em função
do programa de produção e analisa fluxograma. Realiza análise de custos,
efetuando cálculos para fechamento do preços. Elabora orçamentos. Faz
desenvolvimento de produtos gráficos, elaborando projeto de confecção e
edição. Realiza processo criativo de concepção de produto. Elabora
especificação de matérias-primas e outros insumos para a produção.
Define tecnologias, sistemas e processos para confecção dos produtos.
Orienta o fechamento de arquivo -- com resultado da fase de
pré-impressão - em aplicativos, para encaminhamento ao processo de
impressão. Analisa viabilidade técnica e econômica do projeto. Coordena
o processo de impressão e presta informações técnicas - tais como tipo e
gramatura do papel a ser impresso -, de acordo com o projeto gráfico.
Pode operar o processo de impressão e reprodução de produto gráfico.
Estabelece o processo de pós-impressão, para realização de cortes
especiais, relevos e outras operações. Seleciona equipamentos de
acabamento editorial e cartotécnico. Pode ajustar e operar equipamentos
de pós-impressão. Analisa e avalia as características do produto
confeccionado. Monitora estoque do produto acabado. Controla a qualidade
do processo e do produto, estabelecendo padrões de referência. Analisa
variáveis do processo, utilizando metodologia de análise de problemas.
Audita processos a partir de normas estabelecidas. Define plano de
amostragem e realiza ensaios em amostras de matérias-primas e produtos.
Emite relatório de não conformidade do produto. Realiza controle dos
processos gráficos, alimentando bancos de dados, monitorando o andamento
das atividades e calculando os índices de produção. Busca alternativas
para melhorias em tecnologias e matérias-primas. Realiza visitas
técnicas a empresas e fornecedores, para coleta de informações e análise
de tecnologias disponíveis. Consulta normas e fontes de informações
técnicas. Avalia tendências de mercado. Pesquisa ``sites''
especializados da internet e acompanha lançamentos em feiras nacionais e
internacionais da indústria gráfica. Implanta novas tecnologias,
estimando custos de implantação e de operação e orientando fornecedores
para adequação de insumos às tecnologias implantadas. Define
características técnicas de equipamentos e promove melhoria
(``upgrade'') em seu desempenho. Presta assistência técnica, orientando
clientes sobre possibilidades e restrições de materiais e processos.
Propõe soluções corretivas no local de ocorrência de problemas. Fornece
informações técnicas para as diversas áreas da empresa e elabora manuais
de procedimentos. Emite relatório técnico. Coordena equipe de trabalho,
definindo atribuições e assegurando o cumprimento de normas e
procedimentos no trabalho. Pode prestar orientação no uso de
``softwares'' e na operação de máquinas e equipamentos. Pode
supervisionar o trabalho de equipe, avaliando desempenho e levantando
necessidades de treinamento. Realiza programas de treinamento. Verifica
a limpeza e a organização do ambiente de trabalho. Orienta a equipe para
conservar instrumentos limpos, acondicionados e em plenas condições de
uso e funcionamento. Mantém máquinas e equipamentos em condições
operacionais, providenciando manutenção de primeiro nível. Requisita
serviço de manutenção corretiva de máquinas e equipamentos, quando
necessário. Controla a aplicação de procedimentos da empresa para
cumprimento das normas ambientais. Inspeciona o reaproveitamento de
sobras de material e o descarte de resíduos. Zela pela segurança,
utilizando e monitorando o uso de equipamentos de proteção individual.
Identifica situações de risco e colabora para sua eliminação.

\begin{Form}
\textbf{Esta correspondência está adequada?} \\ 
\ChoiceMenu[radio,radiosymbol=\ding{52},bordercolor=0 0.5 0,name=cbovalid_ 371310 ]{}{CBO deve ficar=sim, \hspace{6em} CBO não fica aqui=nao}
\end{Form}

\textbf{CBO: 760605  ( Supervisor das artes gráficas  (indústria editorial e gráfica) )}

Planeja o processo de produção gráfica, analisando os recursos
disponíveis para atendimento aos pedidos recebidos. Verifica o estoque
de materiais -- tais como papel, tintas, entre outros -; dimensiona a
equipe de trabalho para atender à demanda; identifica capacidade de
máquinas e equipamentos; e analisa custo previsto da produção gráfica,
para negociar metas de produção. Aprova a liberação da produção gráfica.
Supervisiona e controla os processos de produção gráfica, estabelecendo
o fluxo das atividades; definindo a forma de utilização das máquinas e
equipamentos; acompanhando cronograma de execução; e controlando a
logística da saída da produção gráfica. Identifica e elimina
desperdícios de insumos e de tempo, para obter melhor aproveitamento dos
recursos. Pode alimentar bancos de dados para o controle das atividades
de produção. Propõe soluções para problemas do processo produtivo,
efetuando pesquisas na Internet; realizando visitas técnicas a empresas;
e fazendo consultas a fornecedores, para a coleta de informações. Avalia
tendências e analisa a incorporação de novos equipamentos - verificando
sua segurança e estimando custo de implantação e operação - e a
adequação de novos materiais, para aumentar a capacidade de produção.
Analisa os resultados do processo produtivo, avaliando índice de
retrabalho, verificando o cumprimento de prazos e conferindo o
atendimento às metas de produção. Analisa custo dos produtos. Monitora o
padrão de qualidade da produção gráfica, coordenando testes para análise
de insumos e produtos. Avalia satisfação de clientes. Administra
pessoal, participando da seleção de profissionais das artes gráficas
para admissão; distribuindo atribuições; controlando horas trabalhadas;
e programando folga e férias da equipe. Supervisiona o trabalho da
equipe, apresentando as metas a serem cumpridas e orientando a execução
das atividades. Avalia desempenho, identificando necessidades de
treinamento. Promove programas de treinamento, para melhoria na
realização das atividades da produção. Pode instruir a equipe na
atualização dos conhecimentos em relação às novas tecnologias adotadas
pela empresa, tais como a adoção de processos híbridos de impressão -
como a impressão ofsete associada à impressão digital -; a execução do
processo CTP-Computador para chapa (Computer to plate), de transferência
da informação digital de arquivos de textos e imagens para chapas; e a
utilização de sistemas automatizados para a impressão e pós-impressão;
entre outras. Elabora documentos, redigindo relatórios técnicos, fazendo
registros de ocorrências -- tais como falhas de equipamentos e acidentes
de trabalho - e preparando pareceres técnicos. Emite comunicações
internas e requisições de matérias-primas. Prepara e emite ordens de
serviço. Orienta equipe para conservar ferramentas e instrumentos
limpos, organizados e acondicionados. Providencia a manutenção
preventiva e a manutenção corretiva de máquinas e equipamentos. Monitora
a limpeza, a higiene e a organização do ambiente de trabalho. Faz busca
de tecnologias menos agressivas ao ambiente, avaliando tendências e
analisando as opções disponíveis no mercado, para atender à legislação
ambiental. Monitora o cumprimento das normas de segurança, inspecionando
a utilização de equipamentos de proteção individual. Identifica
condições de risco no trabalho e atua para eliminá-las.

\begin{Form}
\textbf{Esta correspondência está adequada?} \\ 
\ChoiceMenu[radio,radiosymbol=\ding{52},bordercolor=0 0.5 0,name=cbovalid_ 760605 ]{}{CBO deve ficar=sim, \hspace{6em} CBO não fica aqui=nao}
\end{Form}

\newpage
\section{ 10019 -- Técnico em Química }

\textbf{Perfil do Curso (CNCT)}

O Técnico em Química será habilitado para:

Operar, controlar e monitorar processos industriais e laboratoriais.
Controlar a qualidade de matérias-primas, insumos e produtos. Realizar
amostragens, análises químicas, físico-químicas e microbiológicas.
Desenvolver produtos e processos. Comprar e estocar matérias-primas,
insumos e produtos. Controlar estoques de produtos acabados. Realizar a
especificação de produtos e processos e a seleção de fornecedores de
produtos químicos.

Para atuação como Técnico em Química, são fundamentais:

Conhecimentos e saberes relacionados aos processos de planejamento e
operação das atribuições da área, de modo a assegurar a saúde e a
segurança dos trabalhadores e dos futuros usuários e operadores de
empresas em processos de transformação em química. Conhecimentos e
saberes relacionados à sustentabilidade do processo produtivo, às normas
e relatórios técnicos, à legislação da área, às novas tecnologias
relacionadas à indústria 4.0, à liderança de equipes, à solução de
problemas técnicos e à gestão de conflitos.

\textbf{CBOs correspondidos e seus perfis (presentes no CNCT, e presentes na predição da IA)}

\textbf{CBO: 301105  ( Técnico de laboratório industrial )}

Organiza as condições para a realização dos ensaios, provendo o
laboratório dos materiais de consumo, preparando materiais e
equipamentos, selecionando substâncias reagentes, preparando e
padronizando soluções para análise. Realiza amostragem de materiais,
aplicando métodos de amostragem. Coleta, prepara e preserva amostra,
conforme especificações, normas e procedimentos. Efetua descarte ou
reaproveitamento da amostra, conforme procedimentos estabelecidos.
Realiza os ensaios físicos, químicos, metalográficos e biológicos,
interpretando especificações e normas técnicas de ensaios e definindo
metodologia de análise. Executa a análise. Monitora e valida resultados
obtidos na análise, em bancada ou analisadores em linha. Registra e
informa os resultados. Participa no desenvolvimento de novos produtos,
levantando suas características, auxiliando na sua especificação,
realizando testes e monitorando os resultados. Organiza a produção em
lote-piloto e escala. Participa na avaliação de novos fornecedores.
Colabora no desenvolvimento de metodologias de análises, pesquisando
normas e métodos atualizados, tal como a possibilidade de automação de
ensaios com uso de instrumentação computadorizada, redes e ambientes
gráficos de programação. Testa, otimiza, valida e implementa as
metodologias de análise. Elabora, testa, revisa e padroniza
procedimentos. Elabora instruções de trabalho. Controla a qualidade da
matéria-prima e de produtos acabados e em processo, aplicando métodos
estatísticos e de metrologia, além de métodos e técnicas normalizadas de
análises e ensaios. Interpreta resultados e emite relatórios de controle
de qualidade. Participa no sistema da qualidade da empresa, atuando no
processo de melhoria contínua; atendendo aos procedimentos definidos
pelo sistema de garantia da qualidade; atualizando normas de
procedimentos internos de análise, ensaio, processos e produtos;
colaborando nas auditorias internas e externas da qualidade; e
monitorando a qualidade dos fornecedores. Garante a calibração dos
equipamentos, monitorando validade de calibração, solicitando manutenção
e reparo nos equipamentos e selecionando prestadores de serviços de
calibração. Pode efetuar calibração de equipamentos em laboratórios de
calibração credenciados, aplicando métodos específicos; definindo tipo
de padrão para calibração; registrando dados de calibração; aplicando
normas e critérios de aceitação de calibração; interpretando resultados
em relação ao padrão; e realizando manutenção e reparo nos equipamentos.
Zela pela segurança, utilizando Equipamentos de Proteção Individual
(EPI), atuando na prevenção de acidentes e manuseando os materiais de
análise com base em normas de segurança. Mantém a organização, a limpeza
e a higiene no local de trabalho. Conduz análises para auxiliar no
controle de emissões do processo, aplicando procedimentos de descarte e
segregação de resíduos de laboratório; pesquisando métodos de
recuperação, reciclagem e reaproveitamento; e otimizando métodos de
tratamento de resíduos industriais, para minimizar impactos ambientais
indesejáveis.

\begin{Form}
\textbf{Esta correspondência está adequada?} \\ 
\ChoiceMenu[radio,radiosymbol=\ding{52},bordercolor=0 0.5 0,name=cbovalid_ 301105 ]{}{CBO deve ficar=sim, \hspace{6em} CBO não fica aqui=nao}
\end{Form}

\textbf{CBO: 301110  ( Técnico de laboratório de análises físico-químicas (materiais de construção) )}

Organiza as condições para a realização de análises laboratoriais
físico-químicas, provendo laboratório com materiais de consumo,
preparando materiais e equipamentos para ensaio, selecionando
substâncias reagentes, preparando e padronizando soluções para análise.
Realiza amostragem de materiais, aplicando métodos de amostragem.
Coleta, prepara e preserva amostra conforme especificações, normas e
procedimentos. Efetua descarte ou reaproveitamento da amostra conforme
procedimentos estabelecidos. Realiza análises laboratoriais
físico-químicas, por meio da execução de ensaios -- tais como de
compactação, cisalhamento, adensamento, permeabilidade, resistência,
equivalência de areia, entre outros -, interpretando normas técnicas de
ensaios e especificações e definindo metodologia de análise. Faz a
análise, registrando e informando os resultados. Participa no
desenvolvimento de novos produtos, levantando suas características,
auxiliando na sua especificação, realizando testes e monitorando os
resultados. Participa na avaliação de novos fornecedores. Colabora no
desenvolvimento de metodologias de análises, pesquisando normas e
métodos atualizados, tal como a possibilidade de automação de ensaios
com uso de instrumentação computadorizada, redes e ambientes gráficos de
programação. Testa, otimiza, valida e implementa as metodologias de
análise. Elabora, testa, revisa e padroniza procedimentos. Elabora
instruções de trabalho. Controla a qualidade da matéria-prima e de
produtos acabados e em processo, aplicando métodos estatísticos e de
metrologia, além de métodos e técnicas normalizadas de análises e
ensaios. Interpreta resultados e emite relatórios. Participa no sistema
da qualidade da empresa, atuando no processo de melhoria contínua;
atendendo aos procedimentos definidos pelo sistema de garantia da
qualidade; atualizando normas de procedimentos internos de análise,
ensaio, processos e produtos; colaborando nas auditorias internas e
externas da qualidade; e monitorando a qualidade dos fornecedores.
Garante a calibração dos equipamentos, monitorando validade de
calibração, solicitando manutenção e reparo nos equipamentos e
selecionando prestadores de serviços de calibração. Pode efetuar
calibração de equipamentos em laboratórios credenciados, aplicando
métodos específicos; definindo tipo de padrão para calibração;
registrando dados; aplicando normas e critérios de aceitação de
calibração; interpretando resultados em relação ao padrão; e realizando
manutenção e reparo nos equipamentos. Zela pela segurança, utilizando
Equipamentos de Proteção Individual (EPI), atuando na prevenção de
acidentes e manuseando os materiais de análise com base em normas de
segurança. Mantém a organização, a limpeza e a higiene no local de
trabalho. Conduz análises para auxiliar no controle de emissões do
processo, aplicando procedimentos de descarte e segregação de resíduos
de laboratório; pesquisando métodos de recuperação, reciclagem e
reaproveitamento; e otimizando métodos de tratamento de resíduos
industriais, para minimizar impactos ambientais indesejáveis.

\begin{Form}
\textbf{Esta correspondência está adequada?} \\ 
\ChoiceMenu[radio,radiosymbol=\ding{52},bordercolor=0 0.5 0,name=cbovalid_ 301110 ]{}{CBO deve ficar=sim, \hspace{6em} CBO não fica aqui=nao}
\end{Form}

\textbf{CBO: 301115  ( Técnico químico de petróleo )}

Organiza as condições para a realização de análises laboratoriais
químicas, físicas e físico-químicas, provendo laboratório dos materiais
de consumo, preparando materiais e equipamentos para ensaio,
selecionando substâncias reagentes, preparando e padronizando soluções
para análise. Realiza amostragem de materiais, aplicando métodos de
amostragem. Coleta, prepara e preserva amostra conforme especificações,
normas e procedimentos. Efetua descarte ou reaproveitamento da amostra,
conforme procedimentos estabelecidos. Pode receber amostras para
análise. Realiza análises laboratoriais químicas, físicas e
físico-químicas, por meio da execução de ensaios para determinar
propriedades do material analisado - como poder calorífico, acidez,
volume, peso específico, salinidade e outras -, interpretando normas
técnicas de ensaios e especificações e definindo metodologia de análise.
Faz a análise, registrando e informando os resultados. Participa no
desenvolvimento de novos produtos, levantando suas características,
auxiliando na sua especificação, realizando testes e monitorando os
resultados. Participa na avaliação de novos fornecedores. Participa no
planejamento e na execução dos projetos de perfuração, completação,
manutenção e segurança de poços de petróleo, bem como do controle das
condicionantes ambientais. Colabora no desenvolvimento de metodologias e
análises, pesquisando normas e métodos atualizados de análise, tal como
a possibilidade de automação de ensaios com uso de instrumentação
computadorizada, redes e ambientes gráficos de programação. Testa,
otimiza, valida e implementa as metodologias de análise. Elabora, testa,
revisa e padroniza procedimentos. Elabora instruções de trabalho.
Controla a qualidade da matéria-prima e de produtos acabados e em
processo, aplicando métodos estatísticos e de metrologia, além de
métodos e técnicas normalizadas de análises e ensaios. Interpreta
resultados e emite relatórios. Garante a calibração dos equipamentos,
monitorando validade de calibração, solicitando manutenção e reparo nos
equipamentos e selecionando prestadores de serviços de calibração. Pode
efetuar calibração de equipamentos em laboratórios credenciados,
aplicando métodos específicos; definindo tipo de padrão para calibração;
registrando dados; aplicando normas e critérios de aceitação de
calibração; interpretando resultados em relação ao padrão; e realizando
manutenção e reparo nos equipamentos. Participa no sistema da qualidade
da empresa, atuando no processo de melhoria contínua; atendendo aos
procedimentos definidos pelo sistema de garantia da qualidade;
atualizando normas de procedimentos internos de análise, ensaio,
processos e produtos; colaborando nas auditorias internas e externas da
qualidade; e monitorando a qualidade dos fornecedores. Zela pela
segurança, utilizando Equipamentos de Proteção Individual (EPI), atuando
na prevenção de acidentes e manuseando os materiais de análise com base
em normas de segurança. Mantém a organização, a limpeza e a higiene no
local de trabalho. Conduz análises para auxiliar no controle de emissões
do processo, aplicando procedimentos de descarte e segregação de
resíduos de laboratório; pesquisando métodos de recuperação, reciclagem
e reaproveitamento; e otimizando métodos de tratamento de resíduos
industriais, para minimizar impactos ambientais indesejáveis.

\begin{Form}
\textbf{Esta correspondência está adequada?} \\ 
\ChoiceMenu[radio,radiosymbol=\ding{52},bordercolor=0 0.5 0,name=cbovalid_ 301115 ]{}{CBO deve ficar=sim, \hspace{6em} CBO não fica aqui=nao}
\end{Form}

\textbf{CBO: 311105  ( Técnico químico )}

Realiza análises químicas, físicas e microbiológicas, observando
critérios de confiabilidade metrológica e seguindo procedimentos e
normas técnicas e de qualidade. Coleta e prepara amostras. Prepara
reagentes. Registra os resultados das análises, assumindo
responsabilidade técnica por eles. Opera e controla máquinas,
equipamentos e plantas de processamento químico, interpretando manuais
técnicos e procedimentos de produção. Monitora o funcionamento de
máquinas, equipamentos e instrumentos de medição e controle. Mantém
máquinas e equipamentos em condições de uso, realizando regulagens,
quando necessárias. Supervisiona processos de produção, definindo e
coordenando equipes de trabalho, elaborando o cronograma de produção,
definindo o fluxo do processo, monitorando os parâmetros do processo e
emitindo relatórios de produção. Participa do desenvolvimento de
produtos e processos, testando insumos e matérias-primas, elaborando
formulações para produção, testando novos produtos, definindo o processo
de fabricação de produtos, fazendo a adequação dos produtos às
necessidades dos clientes e validando o produto desenvolvido. Colabora
na implantação ou na reestruturação das instalações industriais,
definindo leiautes e fluxogramas de produção, acompanhando montagem e
instalação de equipamentos e testando máquinas e equipamentos. Presta
assistência técnica, identificando necessidades dos clientes,
diagnosticando problemas técnicos e propondo soluções e melhorias de
produtos e processos. Participa da implantação de programas de
qualidade, analisando indicadores, aplicando ferramentas da qualidade,
implementando ações preventivas e corretivas e participando de
auditorias. Elabora documentação técnica, emitindo laudos e redigindo
procedimentos e relatórios. Controla o estoque de matérias-primas,
insumos e produtos, requisitando compra para reposição. Conserva o local
de trabalho limpo e organizado. Mantém, limpa e esteriliza os
instrumentos e os equipamentos de laboratório. Faz o descarte de
resíduos seguindo normas ambientais.

\begin{Form}
\textbf{Esta correspondência está adequada?} \\ 
\ChoiceMenu[radio,radiosymbol=\ding{52},bordercolor=0 0.5 0,name=cbovalid_ 311105 ]{}{CBO deve ficar=sim, \hspace{6em} CBO não fica aqui=nao}
\end{Form}

\textbf{CBO: 311205  ( Técnico em petroquímica )}

Programa atividades de produção petroquímica, identificando tarefas e
ações necessárias ao cumprimento do planejamento. Define meios e
recursos para o cumprimento de metas de produção e de custos,
monitorando estoques de materiais e de insumos e programando parada e
partida de equipamentos. Elabora, controla e revisa cronograma de
atividades operacionais. Programa serviços complementares, tais como
pinturas e limpeza. Elabora plano de análise contingencial para os casos
de quebra de equipamentos e paradas não programadas. Interpreta conjunto
de documentos técnicos - especificações, desenhos, dentre outros --
referentes à execução das atividades. Coordena o processo de produção
petroquímica, consultando relatórios e ocorrências do turno anterior
para dar sequência às atividades. Monitora a realização da produção,
analisando as variáveis do processo. Controla a execução de manobras
operacionais, efetuando ações de ajustes e solucionando problemas.
Inspeciona serviços contratados vinculados à área de produção, podendo
requisitar outros, complementares. Pode interagir com clientes internos
e externos durante a realização de manobras, operando recursos de
comunicação. Controla a qualidade, com o propósito de evitar produtos
fora de especificação. Ordena insumos - intermediários e finais - e
produtos, rastreando-os para identificar falhas. Coleta amostras de
produtos, realizando ensaios qualitativos, quantitativos e instrumentais
para identificar produtos não conformes, com falhas e fora de padrão.
Monitora o cumprimento do plano de análises laboratoriais, interpretando
laudos de análises químicas. Analisa dados estatísticos do processo
produtivo e controla índices técnicos relacionados à produção e aos
custos. Elabora relatórios e boletins de ocorrências, podendo apresentar
propostas para o desenvolvimento de novos produtos, aplicação de métodos
alternativos e implantação de novos sistemas. Administra a equipe de
trabalho, participando da seleção do pessoal; distribuindo tarefas;
definindo período de férias; e estabelecendo trocas de pessoal. Divulga
normas e procedimentos para a equipe. Elabora instruções operacionais
para a execução dos trabalhos. Supervisiona a equipe de trabalho,
avaliando seu desempenho e identificando a necessidade de treinamento.
Planeja o desenvolvimento da equipe, preparando material didático e
instrucional e selecionando técnica de avaliação dos programas.
Desenvolve programação de treinamento no ambiente de trabalho e avalia
seu impacto no desempenho da equipe. Propõe à equipe discussões sobre
problemas operacionais e instrui cada integrante sobre a importância da
utilização de equipamentos de proteção individual e coletiva. Mantém os
equipamentos em condições operacionais, monitorando seu funcionamento e
providenciando serviços de manutenção de primeiro nível. Preenche ordem
de serviço, relacionando os equipamentos que necessitam serviços de
manutenção corretiva. Faz especificação de materiais, ferramentas,
instrumentos e equipamentos para aquisição e reposição de estoque.
Inspeciona a limpeza e a organização do local de trabalho. Examina a
limpeza, o acondicionamento e a conservação de ferramentas e
instrumentos. Encaminha instrumentos para manutenção, aferindo-os após a
realização do serviço. Monitora o descarte de resíduos, efluentes e
gases conforme normas ambientais. Zela pela segurança, fazendo uso dos
equipamentos de proteção individual. Elabora procedimentos em
conformidade com as normas de segurança, realizando sua divulgação e
controlando seu cumprimento pelas equipes de trabalho. Identifica e
analisa situações de risco em conjunto com outros profissionais,
participando na elaboração, na implementação e na avaliação de ações
corretivas. Emite relatório de acidentes e incidentes, fazendo a
investigação de cada ocorrência em conjunto com profissionais das áreas
envolvidas.

\begin{Form}
\textbf{Esta correspondência está adequada?} \\ 
\ChoiceMenu[radio,radiosymbol=\ding{52},bordercolor=0 0.5 0,name=cbovalid_ 311205 ]{}{CBO deve ficar=sim, \hspace{6em} CBO não fica aqui=nao}
\end{Form}

\newpage
\section{ 10020 -- Técnico em Têxtil }

\textbf{Perfil do Curso (CNCT)}

O Técnico em Têxtil será habilitado para:

Supervisionar os processos produtivos na cadeia têxtil, da ?ação ao
bene?ciamento. Planejar e controlar as operações nos processos nas áreas
de ?ação, tecelagem e bene?ciamento têxtil. Desenvolver padronagens de
malharia ou tecido plano. Desenvolver produtos e processos de
tinturaria, estamparia e acabamento ?nal. Realizar testes de controle de
qualidade, químicos, físicos e colorimétricos. Analisar laudos técnicos.
Controlar estoques de produtos acabados.

Para atuação como Técnico em Têxtil, são fundamentais:

Conhecimentos e saberes relacionados aos processos de planejamento e
operação das atribuições da área, de modo a assegurar a saúde e a
segurança dos trabalhadores e dos futuros usuários e operadores de
empresas em processos de transformação têxtil. Conhecimentos e saberes
relacionados à sustentabilidade do processo produtivo, às normas e
relatórios técnicos, à legislação da área, às novas tecnologias
relacionadas à indústria 4.0, à liderança de equipes, à solução de
problemas técnicos e à gestão de conflitos.

\textbf{CBOs correspondidos e seus perfis (presentes no CNCT, e presentes na predição da IA)}

\textbf{CBO: 311605  ( Técnico têxtil )}

Planeja os processos produtivos têxteis, prevendo meios e recursos para
o cumprimento das metas de produção. Analisa laudos técnicos. Avalia a
capacidade produtiva dos equipamentos; monitora estoques de materiais e
de insumos; e dimensiona o quadro de pessoal necessário nos vários
setores de produção. Elabora programação e cronograma de atividades
operacionais. Coordena os processos produtivos - de fiação, tecelagem,
malharia e beneficiamento têxtil - analisando e adequando o leiaute e o
fluxo de produção. Identifica e soluciona possíveis ``gargalos'' nos
processos de produção. Controla os processos automatizados, avaliando a
produção e providenciando os ajustes nos processos em função dos desvios
apresentados. Elabora fichas técnicas de produtos e processos de
fabricação. Monitora o cumprimento da programação da produção. Emite
relatórios de produção. Administra os processos de produção,
especificando matérias-primas; determinando e registrando os parâmetros
de processo -- como velocidade e temperatura -; e monitorando os itens
de controle e verificação de processos. Controla estoques de produtos
acabados. Controla a qualidade dos processos e produtos, monitorando
indicadores e utilizando ferramentas da qualidade. Padroniza processos.
Realiza testes de controle de qualidade químico e físico para assegurar
as características e a qualidade dos produtos fabricados. Elabora
sistema de auditoria por amostragem, identificando e tratando não
conformidades. Analisa reclamações, críticas e sugestões dos clientes.
Desenvolve fornecedores. Emite parecer de qualidade. Mantém-se
atualizado em relação às inovações na área têxtil. Verifica a
necessidade de melhoria nos processos de trabalho e propõe alternativa.
Desenvolve o processo proposto, elaborando procedimentos; testando e
padronizando o novo método. Avalia novos equipamentos. Elabora fichas
técnicas do novo processo de fabricação. Identifica possíveis impactos
ambientais. Verifica relação custo-benefício. Desenvolve novos produtos,
analisando tendências de mercado e pesquisando novas matérias-primas.
Elabora ficha técnica de novos produtos e define o fluxo para sua
produção. Desenvolve e apresenta amostras. Analisa viabilidade técnica e
econômica dos novos produtos. Propõe ações de marketing para novos
produtos. Coordena equipes de trabalho, participando de processo de
seleção e de administração de pessoal. Distribui tarefas, controla
absenteísmo e programa férias das equipes. Supervisiona equipes de
trabalho, avaliando desempenho e levantando as necessidades de
treinamento. Viabiliza a realização de programas de aperfeiçoamento e
atualização profissional. Controla conflitos nas equipes. Participa da
administração dos custos e recursos, acompanhando a execução do
orçamento e avaliando o custo dos produtos. Pode utilizar softwares
especializados - como ERP-Planejamento de Recursos Empresariais
(Enterprise Resource Planning) - para maior efetividade e rapidez na
transferência de informações e na tomada de decisões. Monitora a limpeza
e a organização do local de trabalho. Documenta o plano de ação de
manutenção preventiva e preditiva de máquinas e equipamentos. Mantém
máquinas e equipamentos em condições operacionais, providenciando
serviços de manutenção de primeiro nível. Preenche ordem de serviço,
relacionando as máquinas e os equipamentos que necessitam serviços de
manutenção corretiva. Contribui para a redução dos impactos ambientais
do processo produtivo, monitorando a aplicação das normas ambientais no
tratamento de efluentes e no descarte de resíduos. Sensibiliza e
conscientiza as equipes de trabalho quanto à importância da preservação
ambiental. Zela pela segurança, fazendo uso dos equipamentos de proteção
individual. Identifica situações de risco e colabora para sua
eliminação. Elabora procedimentos em conformidade com as normas de
segurança, realizando sua divulgação e controlando seu cumprimento pelas
equipes de trabalho.

\begin{Form}
\textbf{Esta correspondência está adequada?} \\ 
\ChoiceMenu[radio,radiosymbol=\ding{52},bordercolor=0 0.5 0,name=cbovalid_ 311605 ]{}{CBO deve ficar=sim, \hspace{6em} CBO não fica aqui=nao}
\end{Form}

\textbf{CBO: 311610  ( Técnico têxtil (tratamentos químicos) )}

Planeja os processos produtivos têxteis, prevendo meios e recursos para
o cumprimento das metas de produção. Faz especificação de produtos
químicos a serem utilizados em diferentes etapas da produção. Monitora
estoques de materiais e insumos. Elabora cronograma de atividades
operacionais. Coordena os processos produtivos, analisando e adequando o
leiaute e o fluxo de produção. Controla os processos automatizados,
avaliando a produção e providenciando os ajustes nos processos em função
dos desvios apresentados. Elabora fichas técnicas de produtos e
processos de fabricação. Administra os processos de tratamento químico
em fios e tecidos nos processos de produção têxtil, em especial nas
etapas de beneficiamento primário, tingimento, estamparia e acabamento
de tecidos. Monitora a aplicação de produtos químicos - como alvejantes
-, no beneficiamento primário. Controla a aplicação de corantes
adequados a cada tipo de tecido, para tingimento e estamparia. Acompanha
e orienta o processo de lavagem, para tirar os excessos do tingimento e
do alvejamento. Examina os processos que proporcionam melhoria das
características - tais como resistência, brilho, amaciamento, acabamento
antichama, entre outros - do tecido. Emite relatórios de produção.
Controla estoques de produtos acabados. Controla a qualidade dos
processos e produtos, monitorando indicadores e utilizando ferramentas
da qualidade. Padroniza processos. Realiza testes de controle de
qualidade químico e físico para assegurar as características e a
qualidade dos produtos fabricados. Elabora sistema de auditoria por
amostragem, identificando e tratando não conformidades. Analisa
reclamações, críticas e sugestões dos clientes. Desenvolve fornecedores.
Emite parecer de qualidade. Mantém-se atualizado em relação às inovações
na área têxtil. Desenvolve novos métodos de tratamento químico de fios e
tecidos, examinando necessidades de melhoria e identificando - inclusive
entre tecnologias emergentes - a tecnologia adequada ao processo em
análise. Testa e padroniza a aplicação do novo método. Avalia os
resultados e verifica a relação custo-benefício. Elabora procedimentos e
fichas técnicas do novo método de tratamento químico de fios e tecidos.
Desenvolve novos produtos - tais como tecidos inteligentes -,
pesquisando tecnologias e materiais. Elabora ficha técnica de novos
produtos e define o fluxo para sua produção. Propõe ações de marketing e
analisa viabilidade técnica e econômica de novos produtos. Coordena
equipes de trabalho, participando de processo de seleção e de
administração de pessoal. Distribui tarefas, controla absenteísmo e
programa férias das equipes. Supervisiona equipes de trabalho, avaliando
desempenho e levantando as necessidades de treinamento. Realiza
programas de treinamento e avalia resultados. Controla conflitos nas
equipes. Participa da administração dos custos e recursos, acompanhando
a execução do orçamento e avaliando o custo dos produtos. Pode utilizar
softwares especializados - como ERP-Planejamento de Recursos
Empresariais (Enterprise Resource Planning) - para agilizar análise de
informações e aumentar eficácia na tomada de decisões. Monitora a
limpeza e a organização do local de trabalho. Zela pela manutenção de
máquinas e equipamentos, elaborando plano de manutenção preventiva e
preditiva, providenciando manutenção de primeiro nível e solicitando
serviços de manutenção corretiva. Contribui para a redução dos impactos
ambientais do processo produtivo, monitorando a aplicação das normas
ambientais no tratamento de efluentes e no descarte de resíduos.
Sensibiliza e conscientiza as equipes de trabalho quanto à importância
da preservação ambiental. Trabalha com segurança, fazendo uso dos
equipamentos de proteção individual. Identifica situações de risco e
colabora para sua eliminação. Elabora procedimentos em conformidade com
as normas de segurança, realizando sua divulgação e controlando seu
cumprimento pelas equipes de trabalho.

\begin{Form}
\textbf{Esta correspondência está adequada?} \\ 
\ChoiceMenu[radio,radiosymbol=\ding{52},bordercolor=0 0.5 0,name=cbovalid_ 311610 ]{}{CBO deve ficar=sim, \hspace{6em} CBO não fica aqui=nao}
\end{Form}

\textbf{CBO: 311615  ( Técnico têxtil de fiação )}

Planeja os processos produtivos têxteis, prevendo meios e recursos para
o cumprimento das metas de produção. Analisa laudos técnicos. Avalia a
capacidade produtiva dos equipamentos; monitora estoques de materiais e
de insumos; e dimensiona o quadro de pessoal necessário nos vários
setores de produção. Elabora cronograma de atividades operacionais.
Coordena os processos produtivos têxteis, analisando e adequando o
leiaute e o fluxo de produção. Identifica e soluciona possíveis
``gargalos'' nos processos de produção. Controla os processos
automatizados, avaliando a produção e providenciando os ajustes nos
processos em função dos desvios apresentados. Elabora fichas técnicas de
produtos e processos de produção. Atua na área de fiação, efetuando
estudos e cálculos para otimizar processos e ampliar a produtividade.
Realiza programação e controle operacional dos processos nas áreas de
fiação. Monitora o recebimento das matérias-primas - fibras naturais,
artificiais e sintéticas -- e examina os resultados da análise de
amostras realizada no laboratório de fiação. Controla os processos de
mistura de fibras, estiragem, duplicação, escartamento, paralelização e
torção. Monitora e orienta os diferentes tipos de fiações convencionais
e a fiação química. Acompanha a transferência de fios têxteis --
produtos da etapa de fiação -- para os setores de tecelagem e malharia.
Emite relatórios da etapa de fiação. Controla a qualidade dos processos
e produtos, monitorando indicadores e utilizando ferramentas da
qualidade. Realiza testes de controle de qualidade químico e físico para
assegurar as características e a qualidade dos produtos fabricados.
Elabora sistema de auditoria por amostragem, identificando e tratando
não conformidades. Analisa reclamações, críticas e sugestões dos
clientes. Desenvolve fornecedores. Emite parecer de qualidade. Mantém-se
atualizado em relação às inovações na área têxtil. Verifica a
necessidade de melhoria nos processos de trabalho e propõe alternativa.
Desenvolve novo método de fiação, elaborando procedimentos, testando e
padronizando processos. Avalia novos equipamentos. Elabora fichas
técnicas do novo processo de fiação. Identifica possíveis impactos
ambientais. Verifica relação custo-benefício. Desenvolve novos fios
têxteis, analisando tendências de mercado e pesquisando novas
matérias-primas. Elabora ficha técnica de novos produtos e define o
fluxo para sua produção. Analisa viabilidade técnica e econômica dos
novos produtos. Coordena equipes de trabalho, participando de processo
de seleção e de administração de pessoal. Distribui tarefas, controla
absenteísmo e programa férias das equipes. Supervisiona equipes de
trabalho, avaliando desempenho e levantando as necessidades de
treinamento. Viabiliza a realização de programas de aperfeiçoamento e
atualização profissional. Controla conflitos nas equipes. Participa da
administração dos custos e recursos, acompanhando a execução do
orçamento e avaliando o custo dos produtos. Pode utilizar softwares
especializados - como ERP-Planejamento de Recursos Empresariais
(Enterprise Resource Planning) - para agilizar análise de informações e
aumentar eficácia na tomada de decisões. Monitora a limpeza e a
organização do local de trabalho. Zela pela manutenção de máquinas e
equipamentos, elaborando plano de manutenção preventiva e preditiva,
providenciando manutenção de primeiro nível e solicitando serviços de
manutenção corretiva. Contribui para a redução dos impactos ambientais
do processo produtivo, monitorando a aplicação das normas ambientais no
tratamento de efluentes e no descarte de resíduos. Sensibiliza e
conscientiza as equipes de trabalho quanto à importância da preservação
ambiental. Trabalha com segurança, fazendo uso dos equipamentos de
proteção individual. Identifica situações de risco e colabora para sua
eliminação. Elabora procedimentos em conformidade com as normas de
segurança, realizando sua divulgação e controlando seu cumprimento pelas
equipes de trabalho.

\begin{Form}
\textbf{Esta correspondência está adequada?} \\ 
\ChoiceMenu[radio,radiosymbol=\ding{52},bordercolor=0 0.5 0,name=cbovalid_ 311615 ]{}{CBO deve ficar=sim, \hspace{6em} CBO não fica aqui=nao}
\end{Form}

\textbf{CBO: 311620  ( Técnico têxtil de malharia )}

Planeja os processos produtivos têxteis, prevendo meios e recursos para
o cumprimento das metas de produção. Analisa laudos técnicos. Avalia a
capacidade produtiva dos equipamentos; monitora estoques de materiais e
de insumos; e dimensiona o quadro de pessoal necessário nos vários
setores de produção. Elabora cronograma de atividades operacionais.
Coordena os processos produtivos têxteis, analisando e adequando o
leiaute e o fluxo de produção. Identifica e soluciona possíveis
``gargalos'' nos processos de produção. Controla os processos
automatizados, avaliando a produção e providenciando os ajustes nos
processos em função dos desvios apresentados. Elabora fichas técnicas de
produtos e processos de produção. Atua na área de malharia, efetuando
estudos e cálculos para otimizar processos e ampliar a produtividade.
Faz a programação das atividades. Acompanha o recebimento dos fios
têxteis. Controla os processos de malharia de trama - circular e
retilínea -- e malharia de urdume. Monitora o funcionamento de máquinas
-- retilíneas, circulares (de pequeno, médio ou grande diâmetro),
Kettenstuhl e Raschel -, orientando o operador no caso de falhas.
Acompanha a transferência dos tecidos e artigos de malha produzidos para
a área de beneficiamento e acabamento. Emite relatório da produção da
malharia. Controla a qualidade dos processos e produtos, monitorando
indicadores e utilizando ferramentas da qualidade. Realiza testes de
controle de qualidade químico e físico para assegurar as características
e a qualidade dos produtos fabricados. Elabora sistema de auditoria por
amostragem, identificando e tratando não conformidades. Analisa
reclamações, críticas e sugestões dos clientes. Desenvolve fornecedores.
Emite parecer de qualidade. Mantém-se atualizado em relação às inovações
na área têxtil. Desenvolve novos métodos de fabricação de tecidos e
artigos de malha, identificando - inclusive entre tecnologias emergentes
- a tecnologia adequada ao processo em análise. Avalia novos
equipamentos. Elabora procedimentos e fichas técnicas do novo processo
de malharia. Identifica possíveis impactos ambientais. Verifica relação
custo-benefício. Desenvolve novos produtos de malharia, analisando
tendências de mercado e aproveitando a grande variedade de texturas
(ligamentos ou estruturas) dos tecidos de malha. Elabora ficha técnica
de novos produtos e define o fluxo para sua produção. Analisa
viabilidade técnica e econômica dos novos produtos. Coordena equipes de
trabalho, participando de processo de seleção e de administração de
pessoal. Distribui tarefas, controla absenteísmo e programa férias das
equipes. Supervisiona equipes de trabalho, avaliando desempenho e
levantando as necessidades de treinamento. Viabiliza a realização de
programas de aperfeiçoamento e atualização profissional. Controla
conflitos nas equipes. Participa da administração dos custos e recursos,
acompanhando a execução do orçamento e avaliando o custo dos produtos.
Pode utilizar softwares especializados - como ERP-Planejamento de
Recursos Empresariais (Enterprise Resource Planning) - para agilizar
análise de informações e aumentar eficácia na tomada de decisões.
Monitora a limpeza e a organização do local de trabalho. Zela pela
manutenção de máquinas e equipamentos, elaborando plano de manutenção
preventiva e preditiva, providenciando manutenção de primeiro nível e
solicitando serviços de manutenção corretiva. Contribui para a redução
dos impactos ambientais do processo produtivo, monitorando a aplicação
das normas ambientais no tratamento de efluentes e no descarte de
resíduos. Sensibiliza e conscientiza as equipes de trabalho quanto à
importância da preservação ambiental. Trabalha com segurança, fazendo
uso dos equipamentos de proteção individual. Identifica situações de
risco e colabora para sua eliminação. Elabora procedimentos em
conformidade com as normas de segurança, realizando sua divulgação e
controlando seu cumprimento pelas equipes de trabalho.

\begin{Form}
\textbf{Esta correspondência está adequada?} \\ 
\ChoiceMenu[radio,radiosymbol=\ding{52},bordercolor=0 0.5 0,name=cbovalid_ 311620 ]{}{CBO deve ficar=sim, \hspace{6em} CBO não fica aqui=nao}
\end{Form}

\textbf{CBO: 311625  ( Técnico têxtil de tecelagem )}

Planeja os processos produtivos têxteis, prevendo meios e recursos para
o cumprimento das metas de produção. Analisa laudos técnicos. Avalia a
capacidade produtiva dos equipamentos; monitora estoques de materiais e
de insumos; e dimensiona o quadro de pessoal necessário nos vários
setores de produção. Elabora cronograma de atividades operacionais.
Coordena os processos produtivos têxteis, analisando e adequando o
leiaute e o fluxo de produção. Identifica e soluciona possíveis
``gargalos'' nos processos de produção. Controla os processos
automatizados, avaliando a produção e providenciando os ajustes nos
processos em função dos desvios apresentados. Elabora fichas técnicas de
produtos e processos de produção. Atua na área de tecelagem, efetuando
estudos e cálculos para otimizar processos e ampliar a produtividade.
Faz a programação das atividades. Acompanha o recebimento dos fios
têxteis. Controla as operações de preparação da tecelagem e engomagem.
Monitora e orienta a utilização correta dos fios e dos teares, de acordo
com os tipos de tecidos a serem produzidos. Orienta o operador no caso
de falhas no funcionamento do tear. Acompanha a transferência dos
tecidos planos produzidos para a área de beneficiamento e acabamento.
Emite relatório da produção da tecelagem. Controla a qualidade dos
processos e produtos, monitorando indicadores e utilizando ferramentas
da qualidade. Realiza testes de controle de qualidade químico e físico
para assegurar as características e a qualidade dos produtos fabricados.
Elabora sistema de auditoria por amostragem, identificando e tratando
não conformidades. Analisa reclamações, críticas e sugestões dos
clientes. Desenvolve fornecedores. Emite parecer de qualidade. Mantém-se
atualizado em relação às inovações na área têxtil. Desenvolve novos
métodos de produção de tecidos planos, identificando - inclusive entre
tecnologias emergentes - a tecnologia adequada ao processo em análise.
Avalia novos equipamentos. Elabora fichas técnicas e procedimentos
referentes ao novo processo de tecelagem. Identifica possíveis impactos
ambientais. Verifica relação custo-benefício. Desenvolve novos tecidos e
novas padronagens, analisando tendências de mercado. Utiliza softwares
para desenvolvimento de padronagem. Elabora ficha técnica de novos
produtos e define o fluxo para sua produção. Analisa viabilidade técnica
e econômica dos novos produtos. Coordena equipes de trabalho,
participando de processo de seleção e de administração de pessoal.
Distribui tarefas, controla absenteísmo e programa férias das equipes.
Supervisiona equipes de trabalho, avaliando desempenho e levantando as
necessidades de treinamento. Viabiliza a realização de programas de
aperfeiçoamento e atualização profissional. Controla conflitos nas
equipes. Participa da administração dos custos e recursos, acompanhando
a execução do orçamento e avaliando o custo dos produtos. Pode utilizar
softwares especializados - como ERP-Planejamento de Recursos
Empresariais (Enterprise Resource Planning) - para agilizar análise de
informações e aumentar eficácia na tomada de decisões. Monitora a
limpeza e a organização do local de trabalho. Zela pela manutenção de
máquinas e equipamentos, elaborando plano de manutenção preventiva e
preditiva, providenciando manutenção de primeiro nível e solicitando
serviços de manutenção corretiva. Contribui para a redução dos impactos
ambientais do processo produtivo, monitorando a aplicação das normas
ambientais no tratamento de efluentes e no descarte de resíduos.
Sensibiliza e conscientiza as equipes de trabalho quanto à importância
da preservação ambiental. Trabalha com segurança, fazendo uso dos
equipamentos de proteção individual. Identifica situações de risco e
colabora para sua eliminação. Elabora procedimentos em conformidade com
as normas de segurança, realizando sua divulgação e controlando seu
cumprimento pelas equipes de trabalho.

\begin{Form}
\textbf{Esta correspondência está adequada?} \\ 
\ChoiceMenu[radio,radiosymbol=\ding{52},bordercolor=0 0.5 0,name=cbovalid_ 311625 ]{}{CBO deve ficar=sim, \hspace{6em} CBO não fica aqui=nao}
\end{Form}

\textbf{CBO: 311710  ( Colorista têxtil )}

Recebe, inspeciona e testa corantes, pigmentos, ligantes, resinas e
demais materiais utilizados na estamparia e no tingimento de materiais
têxteis, para garantir a conformidade com as especificações. Desenvolve
cores em laboratórios, selecionando, pesando e misturando pigmentos,
corantes e materiais auxiliares para preparar tinturas têxteis, de
acordo com especificações. Calcula e prepara receitas de cores para
produção, a partir de formulações de cores preestabelecidas. Coleta
amostras de materiais ou produtos intermediários e acabados para
realizar testes laboratoriais das receitas. Testa a qualidade de
tinturas aplicadas em materiais têxteis, comparando a cor com os padrões
estabelecidos - seja visualmente ou usando equipamentos como
colorímetros ou espectrofotômetros -, ajustando e padronizando cores.
Desenvolve cartela de cores para produtos têxteis e de vestuários,
pesquisando tendências, selecionando e harmonizando cores. Identifica
cores com códigos, em conformidade com padrões internacionais. Pode
interagir com o designer têxtil, para decidir as cores a serem usadas
para cada padrão de tecido e a combinação de cores em uma coleção para
os segmentos de mercado masculino, feminino e infantil. Analisa as
tendências da indústria da moda. Controla a qualidade de processo de
tingimento e estamparia, examinando as características físicas e
químicas das tinturas e corantes. Pode operar e manter máquinas e
equipamentos utilizados no processo de tingimento e estamparia de
materiais têxteis. Registra os dados operacionais ou de produção em
formulários especificados. Conserva local de trabalho limpo e
organizado. Armazena corantes, pigmentos, ligantes, resinas e demais
materiais. Destina adequadamente os resíduos gerados, reciclando
materiais e realizando descarte de acordo com as normas ambientais. Zela
pela segurança, prevenindo acidentes, eliminando condições inseguras e
utilizando equipamentos de proteção individual.

\begin{Form}
\textbf{Esta correspondência está adequada?} \\ 
\ChoiceMenu[radio,radiosymbol=\ding{52},bordercolor=0 0.5 0,name=cbovalid_ 311710 ]{}{CBO deve ficar=sim, \hspace{6em} CBO não fica aqui=nao}
\end{Form}

\textbf{CBO: 311720  ( Preparador de tintas (fábrica de tecidos) )}

Recebe, inspeciona visualmente a cor, verifica a validade e testa
pigmentos, corantes e produtos auxiliares para tingimentos, estamparias
e impressão. Seleciona, pesa e mistura pigmentos, corantes e materiais
auxiliares para preparar tintas de acordo com especificações. Coleta
amostras de materiais ou produtos intermediários e acabados para testes
laboratoriais. Confere a concentração de corantes e pigmentos no
tingimento, estamparia ou impressão de artigos têxteis. Controla a
qualidade de processo de preparação de tintas. Examina as
características físicas e químicas das tintas, tais como temperatura dos
banhos; pH de banhos e pastas; e viscosidade de tintas e pastas. Opera e
mantém máquinas utilizadas no processo de preparação de tintas. Pode
usar software para calibração de cores e para modelar o comportamento de
camadas tonais sobre camadas impressas. Pode adaptar as técnicas de
preparação de tintas de acordo com novos materiais e máquinas de
estamparia digital de alta precisão. Registra os dados operacionais ou
de produção em formulários especificados ou utilizando software
apropriado. Conserva local de trabalho limpo e organizado. Armazena
pigmentos, corantes e produtos auxiliares. Destina adequadamente os
resíduos gerados, reciclando materiais e realizando descarte de acordo
com as normas ambientais. Zela pela segurança, prevenindo acidentes,
eliminando condições inseguras e utilizando equipamentos de proteção
individual e coletiva.

\begin{Form}
\textbf{Esta correspondência está adequada?} \\ 
\ChoiceMenu[radio,radiosymbol=\ding{52},bordercolor=0 0.5 0,name=cbovalid_ 311720 ]{}{CBO deve ficar=sim, \hspace{6em} CBO não fica aqui=nao}
\end{Form}

\textbf{CBO: 318420  ( Desenhista técnico (indústria têxtil) )}

Prepara as atividades, organizando o posto de trabalho de desenho,
inspecionando equipamentos e selecionando ferramentas, instrumentos e
materiais. Interpreta solicitação de desenho e projeto de estamparia,
texturas e cores, reunindo e organizando informações pertinentes.
Consulta revistas, catálogos e outras publicações técnicas. Desenvolve
esboços manualmente e com auxílio de recursos computacionais. Define os
meios de representação gráfica, empregando recursos analógicos e/ou
digitais adequados à demanda. Planeja a elaboração de desenho e projeto
de estamparia, texturas e cores, ordenando o trabalho por etapas e
fixando prazo de entrega. Define escalas e estabelece o formato de
apresentação, impresso e/ou eletrônico. Consulta normas técnicas para
especificar características do desenho. Desenvolve o desenho, combinando
elementos de imagens, para compor a estampa dos tecidos. Realiza
representação gráfica pelo método manual, fazendo uso de instrumentos
para execução do traçado diretamente em substratos de papel, bem como
utilizando computadores, programas integrados de CAD-Desenho Assistido
por Computador (Computer Aided Design), ilustração e tratamento de
imagens. Codifica o desenho e relaciona as especificações técnicas
envolvidas em sua produção. Confere as especificações em face do
trabalho realizado e examina o atendimento às normas técnicas de
representação gráfica, para imprimir e requisitar a aprovação do
desenho. Realiza as correções indicadas pelo solicitante, registra e
arquiva o desenho aprovado. Executa o acabamento final do desenho,
indicando as características de materiais de acabamento. Obtém a
aprovação final. Confecciona a matriz do desenho da estampa, preenche a
legenda do desenho e tira cópias de segurança (back-up). Conserva o
local de trabalho limpo e organizado. Controla graus de luminosidade do
ambiente, temperatura e ruídos. Verifica as condições ergonômicas para
uso de mobiliário, monitor de vídeo de computador, mouse e teclado.
Requisita a aquisição de materiais e instrumentos de trabalho, para
reposição. Conserva equipamentos, ferramentas e instrumentos de trabalho
em plenas condições de uso e funcionamento. Requisita serviço para
manutenção de computador, impressoras e outros periféricos, quando
necessário. Organiza e arquiva documentos físicos e eletrônicos, tais
como comprovantes de aprovação e cópias de segurança dos desenhos e
projetos de estamparia aprovados. Destina adequadamente os resíduos
gerados, fazendo reciclagem e realizando descarte de acordo com as
normas ambientais.

\begin{Form}
\textbf{Esta correspondência está adequada?} \\ 
\ChoiceMenu[radio,radiosymbol=\ding{52},bordercolor=0 0.5 0,name=cbovalid_ 318420 ]{}{CBO deve ficar=sim, \hspace{6em} CBO não fica aqui=nao}
\end{Form}

\textbf{CBO: 352320  ( Agente fiscal têxtil )}

Fiscaliza produtos têxteis - bem como medidas materializadas e
instrumentos metrológicos empregados em suas medições -, verificando a
existência da informação da composição têxtil do produto. Examina os
aspectos formais, conferindo a existência de todas as informações
obrigatórias, analisando aspectos quantitativos e verificando a
fidelidade de informações. Realiza fiscalizações em operações especiais
(blitz). Coleta produtos têxteis, notifica o fiscalizado e verifica o
cumprimento da notificação. Acompanha liberações. Autua o responsável
por infração. Interdita produtos e instrumentos e os apreende, cautelar
e definitivamente. Realiza ensaios, testes e análises, seguindo a
legislação pertinente aos produtos têxteis. Efetua testes prévios e
seletivos de produtos têxteis, inclusive os pré-medidos. Coleta amostras
nos locais onde se fabricam, armazenam, acondicionam ou vendem fios,
tecidos, confecções ou outros produtos têxteis. Realiza análises
laboratoriais e testes químicos, físicos e colorimétricos de produtos
têxteis. Supervisiona as atividades de fiscalização têxtil, organizando
operações especiais e encaminhando produtos têxteis para ensaios de
laboratório. Supervisiona e coordena equipe, programando e acompanhando
as suas atividades e avaliando os trabalhos executados. Disponibiliza
meios para a execução das atividades e controla equipamentos,
instrumentos e materiais utilizados na fiscalização. Participa de grupos
de estudos. Capacita recursos humanos na área de fiscalização têxtil,
elaborando material didático, preparando aulas, ministrando treinamentos
internos e avaliando o desempenho dos treinandos. Ministra palestras e
colabora na disseminação de conceitos de fiscalização têxtil e
metrologia. Orienta tecnicamente o público -- consumidor e fiscalizado
-- informando sobre procedimentos técnicos e administrativos e
esclarecendo sobre legislação. Participa da elaboração de material de
orientação. Atende denúncias e reclamações. Registra o processo de
fiscalização, emitindo termos, intimações, autos e pareceres técnicos.
Preenche fichas cadastrais e relatórios diários.

\begin{Form}
\textbf{Esta correspondência está adequada?} \\ 
\ChoiceMenu[radio,radiosymbol=\ding{52},bordercolor=0 0.5 0,name=cbovalid_ 352320 ]{}{CBO deve ficar=sim, \hspace{6em} CBO não fica aqui=nao}
\end{Form}

\textbf{CBO: 760105  ( Contramestre de acabamento (indústria têxtil) )}

Planeja as atividades para alcance das metas de produção por produto,
dimensionando a disponibilidade de equipamentos e de materiais,
organizando o fluxo de produção, alocando pessoal e distribuindo
tarefas. Define pontos e itens de controle no processo produtivo,
coletando, interpretando e analisando dados relacionados com a produção.
Garante a qualidade do processo produtivo, participando da elaboração de
procedimentos de produção; controlando a qualidade das matérias-primas;
e acompanhando a inspeção de produtos intermediários e acabados.
Participa de auditorias de processo, avaliando o nível de satisfação dos
clientes; verificando a conformidade aos procedimentos técnicos
adotados; identificando situações de não conformidade; e implementando
ações preventivas e corretivas. Mantém-se atualizado em relação às
tendências do vestuário e às novas tecnologias de acabamento e
beneficiamento têxtil -- tais como inovações nas áreas de nanotecnologia
e biotecnologia, novos materiais, tecnologias sustentáveis, entre outras
--, para apresentar propostas que possam aumentar a eficiência dos
processos, melhorar a qualidade dos produtos e reduzir o impacto
ambiental dos processos produtivos. Participa da otimização dos
processos, praticando gestão participativa, buscando informações que
colaborem com a melhoria dos métodos de trabalho e da organização do
fluxo de trabalho. Testa equipamentos, acessórios e insumos. Administra
os recursos materiais, fazendo a previsão de itens necessários e
controlando sua disponibilidade nos processos de produção. Participa na
administração dos custos e recursos, acompanhando a execução do
orçamento e avaliando o custo do produto. Pode utilizar softwares
especializados - como o Planejamento de Recursos Empresariais
(ERP-Enterprise Resource Planning) - para maior efetividade e rapidez na
transferência de informações entre os vários setores da empresa.
Administra os recursos humanos envolvidos no processo produtivo,
definindo os perfis técnicos necessários, participando da seleção de
pessoal, alocando as pessoas da equipe nos postos de trabalho e
controlando absenteísmo. Programa folgas e férias da equipe.
Supervisiona a equipe de trabalho, avaliando desempenho, identificando o
potencial das pessoas e levantando as necessidades de treinamento.
Viabiliza a realização de programas de aperfeiçoamento e atualização
profissional. Controla conflitos na equipe. Colabora na manutenção dos
equipamentos, participando no seu planejamento e na definição de itens
de controle. Supervisiona a realização da manutenção. Controla resíduos
e desperdícios, visando a redução dos custos de produção e o aumento da
sustentabilidade do processo. Descarta resíduos de acordo com as normas
ambientais. Zela pela segurança das pessoas, assegurando o cumprimento
das normas; identificando e eliminando condições inseguras; orientando e
conscientizando a equipe para o uso de Equipamentos de Proteção
Individual (EPI) e Equipamentos de Proteção Coletiva (EPC).

\begin{Form}
\textbf{Esta correspondência está adequada?} \\ 
\ChoiceMenu[radio,radiosymbol=\ding{52},bordercolor=0 0.5 0,name=cbovalid_ 760105 ]{}{CBO deve ficar=sim, \hspace{6em} CBO não fica aqui=nao}
\end{Form}

\textbf{CBO: 760120  ( Contramestre de tecelagem (indústria têxtil) )}

Planeja as atividades para alcance das metas de produção por produto,
dimensionando a disponibilidade de equipamentos e de materiais,
organizando o fluxo de produção, alocando pessoal e distribuindo
tarefas. Define pontos e itens de controle no processo produtivo,
coletando, interpretando e analisando dados relacionados com a produção.
Garante a qualidade do processo produtivo, participando na elaboração de
procedimentos de produção; controlando a qualidade das matérias-primas;
e acompanhando a inspeção de produtos intermediários e acabados.
Participa de auditorias de processo, avaliando o nível de satisfação dos
clientes; verificando a conformidade aos procedimentos técnicos
adotados; identificando situações de não conformidade; e implementando
ações preventivas e corretivas. Mantém-se atualizado sobre as inovações
relacionadas à tecelagem - incluindo conhecimentos sobre novas
tecnologias, novos materiais, tecnologias sustentáveis, entre outras --,
para apresentar propostas que possam aumentar a eficiência dos
processos, melhorar a qualidade dos produtos e reduzir o impacto
ambiental dos processos produtivos. Participa na otimização dos
processos, praticando gestão participativa, buscando informações que
colaborem com a melhoria dos métodos de trabalho e da organização do
fluxo de trabalho. Testa equipamentos, acessórios e insumos. Administra
os recursos materiais, fazendo a previsão de itens necessários e
controlando sua disponibilidade nos processos de produção. Participa na
administração dos custos e recursos, acompanhando a execução do
orçamento e avaliando o custo do produto. Pode utilizar softwares
especializados - como o Planejamento de Recursos Empresariais
(ERP-Enterprise Resource Planning) - para maior efetividade e rapidez na
transferência de informações entre os vários setores da empresa.
Administra os recursos humanos envolvidos no processo produtivo,
definindo os perfis técnicos necessários, participando da seleção de
pessoal, alocando as pessoas da equipe nos postos de trabalho e
controlando absenteísmo. Programa turnos, folgas e férias da equipe.
Supervisiona a equipe de trabalho, avaliando desempenho, identificando o
potencial das pessoas e levantando as necessidades de treinamento.
Viabiliza a realização de programas de aperfeiçoamento e atualização
profissional. Controla conflitos na equipe. Regula máquinas, solicita
peças de reposição e disponibiliza máquinas para manutenção. Colabora na
manutenção dos equipamentos, participando no seu planejamento e na
definição de itens de controle. Supervisiona a realização da manutenção.
Controla resíduos e desperdícios, visando a redução dos custos de
produção e o aumento da sustentabilidade do processo. Descarta resíduos
de acordo com as normas ambientais. Zela pela segurança das pessoas,
assegurando o cumprimento das normas; identificando e eliminando
condições inseguras; orientando e conscientizando a equipe para o uso de
Equipamentos de Proteção Individual (EPI) e Equipamentos de Proteção
Coletiva (EPC).

\begin{Form}
\textbf{Esta correspondência está adequada?} \\ 
\ChoiceMenu[radio,radiosymbol=\ding{52},bordercolor=0 0.5 0,name=cbovalid_ 760120 ]{}{CBO deve ficar=sim, \hspace{6em} CBO não fica aqui=nao}
\end{Form}

\textbf{CBO: 760125  ( Mestre (indústria têxtil e de confecções) )}

Planeja as atividades de acordo com a programação da produção,
examinando a disponibilidade de equipamentos, adequando os recursos
humanos ao volume de trabalho e verificando o estoque de matérias-primas
e de outros materiais. Planeja a localização dos equipamentos em função
do fluxo de produção. Define a sequência da produção. Administra os
pontos críticos do processo produtivo e o otimiza, eliminando gargalos.
Mantém-se atualizado em relação às inovações nos processos produtivos
têxtil e de confecções, tais como as que surgem nas áreas de
biotecnologia, nanotecnologia, tecnologia da informação -- impressão 3D
e novos softwares - e automação industrial. Analisa o conceito de
fábrica inteligente para a indústria de vestuário, com uso de Internet
das Coisas, robótica colaborativa, realidade aumentada, entre outras
inovações. Avalia tecnologias sustentáveis. Acompanha o lançamento de
novos materiais. Examina a viabilidade e os benefícios da implementação
de cada inovação, para apresentar proposta e/ou elaborar parecer técnico
em projetos da empresa. Controla a qualidade do processo produtivo,
podendo redefinir parâmetros utilizados e implementar medidas
corretivas. Orienta os trabalhadores sobre alcance das metas referentes
à qualidade dos produtos. Controla prazos de entrega. Administra os
recursos materiais, controlando o estoque. Encaminha pedidos aos
fornecedores internos e externos, para assegurar os recursos necessários
à produção. Orienta sobre o armazenamento, o manuseio e o transporte de
materiais. Participa da administração dos custos e recursos,
acompanhando a execução do orçamento e avaliando o custo de cada
produto. Pode utilizar softwares especializados - como o Planejamento de
Recursos Empresariais (ERP-Enterprise Resource Planning) -, para obter
informações que subsidiem suas decisões. Administra as equipes de
trabalho envolvidas no processo produtivo, identificando características
pessoais dos membros, monitorando a qualificação profissional da equipe
e delegando responsabilidades. Entrevista candidatos no processo de
seleção de pessoal. Orienta a equipe nos métodos e processos produtivos.
Coordena reuniões e administra conflitos. Supervisiona equipes de
trabalho nas áreas têxtil e/ou de confecções, levando em consideração as
opiniões emitidas e fazendo a interlocução com níveis superiores. Mantém
as equipes informadas sobre as metas e as prioridades da empresa.
Orienta a equipe quanto às ações de manutenção de primeiro nível dos
equipamentos. Avalia a eficácia dos serviços de manutenção. Documenta a
produção, elaborando relatórios e outros documentos. Monitora a limpeza
e a organização dos ambientes de trabalho. Orienta a equipe sobre a
importância do cumprimento das normas de segurança, solicitando, quando
necessário, apoio da área especializada em segurança industrial.

\begin{Form}
\textbf{Esta correspondência está adequada?} \\ 
\ChoiceMenu[radio,radiosymbol=\ding{52},bordercolor=0 0.5 0,name=cbovalid_ 760125 ]{}{CBO deve ficar=sim, \hspace{6em} CBO não fica aqui=nao}
\end{Form}

\textbf{CBO: 761105  ( Classificador de fibras têxteis )}

Planeja atividades de classificação de fibras têxteis naturais -- de
origem animal, mineral ou vegetal -, a fim de analisar o atendimento aos
requisitos de qualidade para produção de artigos têxteis. Seleciona
máquinas e equipamentos para realização das atividades. Estabelece
etapas e cronograma para execução do trabalho. Monitora a recepção de
fibras naturais para classificação, registrando - em formulário de
sistema informatizado - número de identificação, origem, data e outros
dados de cada fardo ou molho (feixe ou porção) recebido. Realiza a
classificação visual e manual das fibras naturais, examinando a cor das
fibras; a quantidade de impurezas; a existência de fibras mortas; as
contaminações por matérias estranhas; a espessura; a presença de
``neps'' (emaranhados de fibras normalmente ocasionados por fibras
imaturas); entre outros fatores. Extrai amostras de fibras têxteis dos
fardos ou molhos. Registra cada amostra, com seu peso bruto e com o
número de identificação do fardo ou molho do qual foi extraída. Pode
utilizar máquina elétrica portátil no corte de amostras. Embala lotes de
amostras de fibras têxteis e lacra a embalagem, para remessa a
laboratórios credenciados. Registra os dados dos lotes de amostras em
sistema informatizado, para fins de rastreabilidade. Encaminha amostras
de fibras têxteis para realização da classificação tecnológica ou
instrumental, por equipamento - como o Instrumento de Alto Volume para
Teste de Fibra (HVI- High Volume Instrument) -- semiautomático ou
automático. Recebe relatórios emitidos pelo equipamento, com dados sobre
resistência, comprimento, ``micronaire'' (índice determinado por finura
e maturidade da fibra), grau de impureza, entre outras características
das fibras. Pode usar aparelho para medir índice de umidade de amostras.
Analisa resultados da classificação tecnológica ou instrumental,
interpretando-os e comparando-os com os resultados da classificação
visual e manual. Compara resultados com parâmetros de referência.
Considera, na análise, normas técnicas e legislação vigente. Registra,
em sistema informatizado, os resultados. Documenta a classificação das
fibras, elaborando relatório e emitindo laudo e certificado. Organiza e
arquiva - em meios físico e eletrônico -- as informações sobre as fibras
recebidas e sobre a extração de amostras; os resultados da classificação
visual e manual; os resultados da classificação tecnológica ou
instrumental; e as cópias de laudos e outros documentos emitidos.
Conserva o ambiente de trabalho limpo e organizado. Mantém ferramentas,
instrumentos, aparelhos e equipamentos limpos, organizados,
acondicionados e em plenas condições de uso e funcionamento. Realiza
descarte de resíduos, conforme normas ambientais. Trabalha com
segurança, utilizando equipamentos de proteção individual.

\begin{Form}
\textbf{Esta correspondência está adequada?} \\ 
\ChoiceMenu[radio,radiosymbol=\ding{52},bordercolor=0 0.5 0,name=cbovalid_ 761105 ]{}{CBO deve ficar=sim, \hspace{6em} CBO não fica aqui=nao}
\end{Form}

\textbf{CBO: 761110  ( Lavador de lã )}

Planeja atividades de lavagem de lã, a fim de prepará-la para o
processamento industrial têxtil. Interpreta ordem de serviço e
programação de produção. Verifica a disponibilidade de produtos químicos
para realização das atividades. Prepara a solução para lavagem da lã,
empregando as técnicas apropriadas de manipulação de produtos químicos.
Mistura água, detergentes e outros produtos químicos, de acordo com as
quantidades especificadas na ordem de serviço. Homogeneíza a solução.
Prepara máquina para lavagem da lã, regulando-a. Coloca, na máquina, a
solução preparada para lavagem. Controla o processo de lavagem,
monitorando a entrada da lã na máquina; acionando os dispositivos de
comando da máquina; e conferindo a extração de gordura da lã. Monitora
processo de secagem da lã, controlando temperatura no secador. Coleta
amostras de lã, ao finalizar o processo de secagem. Requisita teste da
umidade e do teor graxo das amostras coletadas de lã. Analisa o
resultado dos testes para controlar o grau de umidade e o teor graxo.
Confere o fluxo de produção de lã lavada, comparando-o com a programação
preestabelecida. Limpa a máquina e os equipamentos auxiliares, a fim de
prepará-los para nova operação de lavagem. Conserva o ambiente de
trabalho limpo e organizado, acondicionando os produtos químicos.
Realiza a manutenção básica de máquinas e equipamentos, lubrificando-os,
executando ajustes e fazendo pequenos reparos. Identifica falhas e
defeitos durante o funcionamento de máquinas e equipamentos,
requisitando serviços de manutenção corretiva. Acompanha a realização do
serviço de manutenção e, ao seu término, testa o funcionamento de
máquinas e equipamentos. Descarta resíduos sólidos e efluentes do
processo de lavagem, seguindo os procedimentos estabelecidos pela
empresa para cumprimento de normas ambientais. Trabalha com segurança,
utilizando equipamentos de proteção individual. Avalia condições de
riscos no local de trabalho e realiza ações para prevenir acidentes.
Segue procedimentos de segurança na manipulação de produtos químicos.

\begin{Form}
\textbf{Esta correspondência está adequada?} \\ 
\ChoiceMenu[radio,radiosymbol=\ding{52},bordercolor=0 0.5 0,name=cbovalid_ 761110 ]{}{CBO deve ficar=sim, \hspace{6em} CBO não fica aqui=nao}
\end{Form}

\textbf{CBO: 761810  ( Revisor de fios (produção têxtil) )}

Planeja as atividades de revisão de fios têxteis crus e acabados,
levando em consideração a análise do relatório de produção e o estudo
sobre evolução de casos de não conformidade aos padrões de qualidade.
Define cronograma de ações. Inspeciona - visualmente e com uso de
instrumentos e aparelhos - bobinas de fios têxteis crus e acabados,
identificando e classificando defeitos de acordo com padrões
preestabelecidos. Registra as não conformidades detectadas durante a
inspeção. Executa análise de causas dos defeitos em fios têxteis.
Realiza ações preventivas e corretivas de acordo com os defeitos
identificados, podendo interromper a produção; fazer reunião com equipe
responsável pelas não conformidades para implantação de correções e
melhorias de processos de produção; entre outras providências. Registra
dados para controle estatístico e de qualidade, calculando número de
defeitos em valores estatísticos e emitindo relatórios de não
conformidades para avaliação da produção. Informa, ao encarregado da
produção, a necessidade de reprogramação dos produtos. Certifica lotes
de bobinas de fios têxteis produzidos. Prepara lotes de produção
conforme programação preestabelecida, selecionando bobinas de fios
têxteis para beneficiamento e expedição e separando produtos de
qualidade inferior. Identifica necessidade de realização de treinamento
operacional, avaliando o desempenho dos operadores de fiação. Compara o
desempenho do operador antes do programa de treinamento com o seu índice
de produtividade e qualidade após realizar o treinamento. Apresenta
sugestões para melhoria do plano de treinamento operacional. Mantém-se
atualizado em relação às tendências e às inovações nas áreas de fiação,
acabamento e beneficiamento de fios têxteis. Examina o cumprimento dos
procedimentos de limpeza e organização do ambiente de trabalho.
Incentiva a seleção de resíduos têxteis para reaproveitamento e orienta
o descarte dos demais, de acordo com os procedimentos adotados pela
empresa. Trabalha com segurança, utilizando equipamentos de proteção
individual; identificando condições inseguras e providenciando medidas
para eliminá-las; e prevenindo acidentes.

\begin{Form}
\textbf{Esta correspondência está adequada?} \\ 
\ChoiceMenu[radio,radiosymbol=\ding{52},bordercolor=0 0.5 0,name=cbovalid_ 761810 ]{}{CBO deve ficar=sim, \hspace{6em} CBO não fica aqui=nao}
\end{Form}

\textbf{CBO: 761815  ( Revisor de tecidos acabados )}

Planeja as atividades de revisão de tecidos planos e de malha acabados,
levando em consideração a análise do relatório de produção e o estudo
sobre evolução de casos de não conformidade aos padrões de qualidade.
Define cronograma de ações. Inspeciona - visualmente e com uso de
instrumentos e aparelhos - tecidos planos e de malha acabados,
identificando e classificando defeitos de acordo com padrões
preestabelecidos. Registra as não conformidades detectadas durante a
inspeção. Executa análise de causas dos defeitos em tecidos acabados.
Realiza ações preventivas e corretivas de acordo com os defeitos
identificados, podendo interromper a produção; fazer reunião com equipe
responsável pelas não conformidades para implantação de correções e
melhorias de processos de produção; analisar a possibilidade de remoção
dos defeitos dos tecidos; entre outras providências. Registra dados para
controle estatístico e de qualidade, calculando número de defeitos em
valores estatísticos e emitindo relatórios de não conformidades para
avaliação da produção. Informa, ao encarregado da produção, a
necessidade de reprogramação dos produtos. Certifica lotes de tecidos
acabados produzidos. Prepara lotes de produção conforme programação
preestabelecida, selecionando tecidos planos e de malha acabados para
expedição e separando os produtos de qualidade inferior. Identifica
necessidade de realização de treinamento operacional, avaliando o
desempenho dos operadores de produção. Compara o desempenho do operador
antes do programa de treinamento com o seu índice de produtividade e
qualidade após realizar o treinamento. Apresenta sugestões para melhoria
do plano de treinamento operacional. Mantém-se atualizado em relação às
tendências e às inovações nas áreas de tecelagem, malharia, acabamento e
beneficiamento de produtos têxteis. Examina o cumprimento dos
procedimentos de limpeza e organização do ambiente de trabalho.
Incentiva a seleção de resíduos têxteis para reaproveitamento e orienta
o descarte dos demais, de acordo com os procedimentos adotados pela
empresa. Trabalha com segurança, utilizando equipamentos de proteção
individual; identificando condições inseguras e providenciando medidas
para eliminá-las; e prevenindo acidentes.

\begin{Form}
\textbf{Esta correspondência está adequada?} \\ 
\ChoiceMenu[radio,radiosymbol=\ding{52},bordercolor=0 0.5 0,name=cbovalid_ 761815 ]{}{CBO deve ficar=sim, \hspace{6em} CBO não fica aqui=nao}
\end{Form}

\newpage
\section{ 10021 -- Técnico em Vestuário }

\textbf{Perfil do Curso (CNCT)}

O Técnico em Vestuário será habilitado para:

Supervisionar o processo de confecção do produto conforme padrões de
qualidade. Analisar e de?nir a melhor sequência de montagem do produto,
de acordo com a forma de execução e as características da matéria-prima.
Propor e analisar métodos de trabalho dos processos fabris de vestuário.
Determinar o tempo-padrão das operações e dimensionar recursos
necessários ao atendimento das demandas de clientes. Supervisionar a
utilização de máquinas de costura industrial e equipamentos. Organizar o
fluxo de produção. Monitorar o desempenho da produção. Supervisionar a
execução de plano de manutenção. Controlar estoques de produtos
acabados. Apoiar a equipe de desenvolvimento de produto em função das
características operacionais da produção interna ou externa.

Para atuação como Técnico em Vestuário, são fundamentais:

Conhecimentos e saberes relacionados aos processos de planejamento e
operação das atribuições da área, de modo a assegurar a saúde e a
segurança dos trabalhadores e dos futuros usuários e operadores de
empresas em processos de transformação em vestuário e indústria da moda.
Conhecimentos e saberes relacionados à sustentabilidade do processo
produtivo, às normas e relatórios técnicos, à legislação da área, às
novas tecnologias relacionadas à indústria 4.0, à liderança de equipes,
à solução de problemas técnicos e à gestão de conflitos.

\textbf{CBOs correspondidos e seus perfis (presentes no CNCT, e presentes na predição da IA)}

\textbf{CBO: 319105  ( Técnico em calçados e artefatos de couro )}

Participa no planejamento dos processos de produção de calçados e de
artefatos de couro, considerando as metas preestabelecidas. Propõe
leiaute da produção. Seleciona máquinas e equipamentos. Monitora
estoques de insumos. Auxilia no dimensionamento do quadro de pessoal
necessário nos vários setores de produção. Participa na elaboração da
programação e do cronograma de atividades de produção. Prepara a
produção de calçados e de artefatos de couro, interpretando ordem de
produção. Verifica as especificações e interpreta os desenhos referentes
aos modelos a serem produzidos. Coordena a produção de protótipos,
fornecendo as informações necessárias; monitorando o processo de
produção; e testando o protótipo produzido. Pode utilizar impressão 3D
na confecção dos protótipos. Supervisiona operações relativas à produção
de calçados e de artefatos de couro, orientando corte, costura
(pesponto), montagem e acabamento. Controla a produção, medindo tempo
para desenvolvimento das etapas; implementando ações corretivas nos
processos; e verificando as condições de funcionamento de máquinas e
equipamentos. Controla estoques de produtos acabados. Controla a
qualidade dos processos e produtos. Mede indicadores de desempenho,
avaliando produtividade. Identifica perdas na produção. Elabora
estatísticas de produção. Pode orientar o uso de ``softwares'', como os
de modelagem e os relacionados a outras etapas da produção. Pode fazer
uso de ``softwares'' e sistemas integrados de gestão empresarial para
controle de produção; controle de estoques; e gerenciamento de custos.
Verifica a necessidade de melhoria nos processos de trabalho e propõe
alternativa. Descreve o processo produtivo proposto, definindo o fluxo e
estabelecendo o tempo-padrão das operações de produção. Desenvolve novos
produtos, analisando tendências de mercado e considerando o
público-alvo. Pesquisa novas matérias-primas, identificando as
apropriadas para os produtos. Elabora especificações das dimensões e das
características dos produtos. Elabora ficha técnica dos novos produtos.
Desenvolve amostras e faz sua apresentação. Indica sistemas de embalagem
dos produtos e de acondicionamento da produção. Analisa viabilidade
técnica e econômica dos novos produtos. Verifica a concorrência e propõe
ações de marketing para os novos produtos. Mantém-se atualizado em
relação às inovações na área de calçados e artefatos de couro. Avalia
inovações, como novas máquinas; sistemas de produção automatizados;
sistemas a laser; entre outras. Identifica fornecedores, testando
equipamentos e materiais oferecidos e comparando preços, prazos de
entrega e formas de pagamento. Levanta necessidades de treinamento de
fornecedores e desenvolve programação. Coordena equipe de trabalho,
participando do processo de seleção e de administração de pessoal.
Distribui tarefas, controla absenteísmo e programa férias da equipe.
Supervisiona o trabalho da equipe, avaliando desempenho e levantando as
necessidades de treinamento. Viabiliza a realização de programas de
treinamento. Participa da administração financeira, acompanhando a
execução do orçamento e calculando custo de produto. Documenta o plano
de manutenção preventiva e preditiva de máquinas e equipamentos. Mantém
máquinas e equipamentos em condições operacionais, providenciando
manutenção de primeiro nível. Requisita serviços de manutenção corretiva
de máquinas e equipamentos. Monitora a limpeza e a organização do local
de trabalho. Orienta equipe para conservar ferramentas e instrumentos
limpos, acondicionados e em plenas condições de uso e funcionamento.
Contribui para a redução dos impactos ambientais do processo produtivo,
orientando o reaproveitamento de sobras e monitorando a aplicação das
normas ambientais no descarte de resíduos. Zela pela segurança,
identificando situações de risco e colaborando para sua eliminação.
Elabora procedimentos em conformidade com as normas de segurança.
Utiliza e monitora o uso de equipamentos de proteção individual.

\begin{Form}
\textbf{Esta correspondência está adequada?} \\ 
\ChoiceMenu[radio,radiosymbol=\ding{52},bordercolor=0 0.5 0,name=cbovalid_ 319105 ]{}{CBO deve ficar=sim, \hspace{6em} CBO não fica aqui=nao}
\end{Form}

\textbf{CBO: 319110  ( Técnico em confecções do vestuário )}

Participa no planejamento de produção de confecções em indústria de
vestuário, considerando as metas preestabelecidas. Propõe leiaute da
produção. Seleciona máquinas e equipamentos. Monitora estoques de
insumos. Auxilia no dimensionamento do quadro de pessoal necessário nos
vários setores de produção. Participa na elaboração da programação e do
cronograma de atividades de produção. Prepara a produção de confecções,
interpretando ordem de produção. Verifica especificações de materiais.
Lê e interpreta desenhos e projeto de confecções. Coordena a produção de
protótipos, fornecendo as informações necessárias; monitorando o
processo de produção; e testando o protótipo produzido. Supervisiona
operações relativas à produção de confecções, orientando corte, costura,
montagem e acabamento. Controla a produção, medindo tempo para
desenvolvimento das etapas; implementando ações corretivas nos
processos; e verificando as condições de funcionamento de máquinas e
equipamentos. Controla estoques de produtos acabados. Realiza controle
de qualidade dos processos e produtos. Mede indicadores de desempenho,
avaliando produtividade. Identifica perdas na produção. Elabora
estatísticas de produção. Pode orientar o uso de máquinas automáticas de
enfesto, a utilização de ``software'' de modelagem, entre outras novas
tecnologias aplicadas ao processo de produção. Pode fazer uso de
``softwares'' e sistemas integrados de gestão empresarial para controle
de produção; controle de estoques; e gerenciamento de custos. Verifica a
necessidade de melhoria nos processos de trabalho e propõe alternativa.
Descreve o processo produtivo proposto, definindo o fluxo e
estabelecendo o tempo-padrão das operações de produção. Desenvolve nova
coleção de confecções, considerando o público-alvo e analisando
tendências da moda. Pesquisa novas matérias-primas, identificando as
apropriadas para as confecções. Detalha as características dos produtos.
Elabora ficha técnica das novas confecções. Desenvolve amostras e faz
sua apresentação. Analisa e define a melhor sequência de montagem dos
produtos, considerando as características da matéria-prima. Indica
sistemas de embalagem dos produtos e de acondicionamento da produção.
Analisa viabilidade técnica e econômica dos novos produtos. Verifica a
concorrência e propõe ações de marketing para as novas confecções.
Identifica fornecedores, testando equipamentos e materiais oferecidos e
comparando preços, prazos de entrega e formas de pagamento. Levanta
necessidades de treinamento de fornecedores e desenvolve programação.
Mantém-se atualizado em relação às inovações na área, como máquinas
robotizadas de costura, tecidos inteligentes, entre outras. Coordena
equipe de trabalho, participando do processo de seleção e de
administração de pessoal. Distribui tarefas, controla absenteísmo e
programa férias da equipe. Supervisiona o trabalho da equipe, avaliando
desempenho e levantando as necessidades de treinamento. Viabiliza a
realização de programas de treinamento. Participa da administração
financeira, acompanhando a execução do orçamento e calculando custo de
produto. Documenta o plano de manutenção preventiva e preditiva de
máquinas e equipamentos. Mantém máquinas e equipamentos em condições
operacionais, providenciando manutenção de primeiro nível. Requisita
serviços de manutenção corretiva de máquinas e equipamentos. Monitora a
limpeza e a organização do local de trabalho. Orienta equipe para
conservar ferramentas e instrumentos limpos, acondicionados e em plenas
condições de uso e funcionamento. Contribui para a redução dos impactos
ambientais do processo produtivo, orientando o reaproveitamento de
sobras e monitorando a aplicação das normas ambientais no descarte de
resíduos. Zela pela segurança, identificando situações de risco e
colaborando para sua eliminação. Elabora procedimentos em conformidade
com as normas de segurança. Utiliza e monitora o uso de equipamentos de
proteção individual.

\begin{Form}
\textbf{Esta correspondência está adequada?} \\ 
\ChoiceMenu[radio,radiosymbol=\ding{52},bordercolor=0 0.5 0,name=cbovalid_ 319110 ]{}{CBO deve ficar=sim, \hspace{6em} CBO não fica aqui=nao}
\end{Form}

\newpage
\section{ 11001 -- Técnico em Agricultura }

\textbf{Perfil do Curso (CNCT)}

O Técnico em Agricultura será habilitado para:

Planejar, organizar, dirigir e controlar a produção vegetal de forma
sustentável, analisando as características econômicas, sociais e
ambientais. Elaborar e executar projetos de produção agrícola, aplicando
as Boas Práticas de Produção Agrícola (BPA). Prestar assistência técnica
e assessoria ao estudo e ao desenvolvimento de projetos e pesquisas
tecnológicas, ou aos trabalhos de vistoria, perícia, arbitramento e
consultoria. Elaborar orçamentos, laudos, pareceres, relatórios e
projetos, inclusive de incorporação de novas tecnologias. Prestar
assistência técnica às áreas de crédito rural e agroindustrial, de
topografia na área rural, de impacto ambiental, de paisagismo, de
jardinagem e horticultura, de construção de benfeitorias rurais, de
drenagem e irrigação. Planejar, organizar e monitorar atividades de
exploração e manejo do solo, matas e florestas de acordo com suas
características, com as alternativas de otimização dos fatores
climáticos e seus efeitos no crescimento e desenvolvimento das plantas e
dos animais. Produzir mudas e sementes, em propagação, em cultivos
abertos ou protegidos, em viveiros e em casas de vegetação. Planejar,
organizar e monitorar o processo de aquisição, preparo, conservação e
armazenamento da matéria prima e dos produtos agroindustriais. Aplicar
métodos e programas de melhoramento genético. Prestar assistência
técnica à aplicação, à comercialização, ao manejo de produtos
especializados, à recomendação e à interpretação de análise de solos, à
aplicação de fertilizantes e corretivos nos tratos das culturas.
Identificar os processos simbióticos de absorção, de translocação e os
efeitos alelopáticos entre solo e planta, planejando ações referentes
aos tratos das culturas. Selecionar e aplicar métodos de erradicação e
controle de vetores e pragas, doenças e plantas daninhas. Planejar e
acompanhar a colheita e a pós-colheita. Supervisionar o armazenamento, a
conservação, a comercialização e a industrialização dos produtos
agrícolas. Elaborar, aplicar e monitorar programas profiláticos,
higiênicos e sanitários na produção vegetal e agroindustrial. Implantar
e gerenciar sistemas de controle de qualidade na produção agrícola.
Emitir laudos e documentos de classificação e exercer a fiscalização de
produtos de origem vegetal, animal e agroindustrial. Implantar pomares e
acompanhar seu desenvolvimento até a fase produtiva, emitindo os
respectivos certificados de origem e qualidade de produtos. Treinar e
conduzir equipes nas suas modalidades de atuação profissional. Aplicar
as legislações pertinentes ao processo produtivo e ao meio ambiente.
Aplicar práticas sustentáveis no manejo de conservação do solo e da
água. Identificar e aplicar técnicas mercadológicas para distribuição e
comercialização de produtos agrícolas. Executar a gestão econômica e
financeira da produção agrícola. Administrar e gerenciar propriedades
agrícolas. Realizar procedimentos de desmembramento, parcelamento e
incorporação de imóveis rurais. Operar, manejar e regular máquinas,
implementos e equipamentos agrícolas. Operar veículos aéreos remotamente
pilotados e equipamentos de precisão para monitoramento remoto da
produção agrícola.

Para a atuação como Técnico em Agricultura, são fundamentais:

Conhecimentos e saberes relacionados à produção agrícola, à produção e
ao processamento de alimentos, fitossanidade e proteção ambiental.
Atualização em relação às inovações tecnológicas. Cooperação de forma
construtiva e colaborativa nos trabalhos em equipe, tomada de decisões.
Adoção de senso investigativo, visão sistêmica das atividades e
processos. Capacidade de comunicação e argumentação, autonomia,
proatividade, liderança, respeito às diversidades nos grupos de
trabalho, resiliência frente aos problemas, organização,
responsabilidade, visão crítica, humanística, ética e consciência em
relação ao impacto de sua atuação profissional na sociedade e no
ambiente.

\textbf{CBOs correspondidos e seus perfis (presentes no CNCT, e presentes na predição da IA)}

\textbf{CBO: 321105  ( Técnico agrícola )}

Planeja projeto agrícola, pesquisando mercado consumidor e analisando
viabilidade econômica. Avalia condições da propriedade, verificando sua
infraestrutura - máquinas, equipamentos e instalações -; levantando
dados sobre a área a ser trabalhada -- topografia e extensão --;
examinando as condições edafoclimáticas - solo, clima e água -;
verificando a disponibilidade de mão de obra para as atividades;
levantando a capacitação tecnológica do produtor; e pesquisando o
mercado fornecedor de insumos, materiais, máquinas e equipamentos.
Executa levantamento do custo-benefício para o produtor. Verifica as
condições para implantação de projeto agrícola, coletando amostras para
análise de solo e interpretando os resultados laboratoriais; verificando
a disponibilidade e a qualidade da água a ser utilizada; e coletando
dados meteorológicos. Executa projeto agrícola, definindo cultivares e
rotação de culturas. Elabora planta de construções rurais. Compra
máquinas, equipamentos, insumos e materiais. Supervisiona a colheita e a
pós-colheita das principais culturas, aplicando técnicas de adubação
verde, de cobertura morta e de correção do solo, conforme cada caso.
Orienta os processos de silagem e transporte. Supervisiona trabalhadores
agrícolas, contratando mão de obra, distribuindo tarefas, orientando e
acompanhando a execução de atividades. Fiscaliza produção agrícola,
inspecionando documentação de produtos em trânsito, examinando a
sanidade de produtos agrícolas, enviando amostras de produtos para
análises laboratoriais, fiscalizando aplicação de agrotóxicos e
inspecionando cumprimento de normas e padrões técnicos. Conduz
experimentos de pesquisa, elaborando relatórios, laudos, pareceres,
perícias e avaliações. Administra propriedades rurais, adotando sistema
de produção conforme necessidade do mercado e controlando despesas e
receitas. Administra funcionários da propriedade. Comercializa produção
agrícola. Presta assistência e consultoria técnicas, orientando na
escolha do local para atividade agrícola; aconselhando sobre escolha de
espécies e cultivares; instruindo sobre preparo, correção e conservação
de solo; informando sobre equipamentos, máquinas e implementos;
indicando formas e manejo de irrigação e drenagem; aconselhando sobre
manejo integrado de pragas e doenças; orientando sobre época de plantio,
tratos culturais e colheita; e estimulando o beneficiamento de produtos
agrícolas. Orienta sobre preservação ambiental. Promove extensão e
capacitação rural, organizando reuniões com produtores e orientando
formação de associações e grupos de produtores. Assessora produtores,
orientando-os em cultivo com técnicas convencionais ou alternativas de
plantio, em manejo do solo e em compra e venda de insumos, materiais e
produtos agrícolas. Promove a difusão de tecnologia, orientando a
implantação de agricultura de precisão e o uso de sistemas
automatizados, bem como o desenvolvimento de novos cultivares. Dissemina
o uso de ambientes protegidos de cultivo e o melhoramento de sementes.
Divulga o uso de controle remoto da produção e da qualidade de processos
e produtos. Ministra cursos e programas de treinamentos. Apresenta
resultados de pesquisas em encontros e congressos da área agrícola. Pode
participar do desenvolvimento de tecnologias agrícolas, adequando
equipamentos e instalações conforme necessidade da região e do produtor.
Cria e adapta técnicas convencionais e alternativas de cultivo. Orienta
produtores agrícolas sobre biossegurança, indicando uso racional de
agrotóxicos e esclarecendo sobre destino de embalagens de agrotóxicos.
Informa sobre limpeza e desinfeção de máquinas, equipamentos e
instalações.

\begin{Form}
\textbf{Esta correspondência está adequada?} \\ 
\ChoiceMenu[radio,radiosymbol=\ding{52},bordercolor=0 0.5 0,name=cbovalid_ 321105 ]{}{CBO deve ficar=sim, \hspace{6em} CBO não fica aqui=nao}
\end{Form}

\newpage
\section{ 11002 -- Técnico em Agroecologia }

\textbf{Perfil do Curso (CNCT)}

O Técnico em Agroecologia será habilitado para:

Planejar, organizar, dirigir e controlar a produção agrícola de forma
sustentável, analisando as características econômicas, sociais e
ambientais. Elaborar e executar projetos de sistemas agroecológicos de
produção agropecuária e agroextrativista e sistemas orgânicos de
produção, aplicando as Boas Práticas de Produção Agrícola (BPA).
Planejar, organizar e monitorar atividades de exploração e manejo do
solo, das matas e das florestas de acordo com suas características, com
as alternativas de otimização dos fatores climáticos e seus efeitos no
crescimento e desenvolvimento das plantas e dos animais. Produzir mudas
e sementes, em propagação, em cultivos abertos ou protegidos, em
viveiros e em casas de vegetação. Planejar, organizar e monitorar
programas de nutrição e manejo alimentar em projetos zootécnicos.
Planejar, organizar e monitorar o processo de aquisição, preparo,
conservação e armazenamento da matéria prima e dos produtos
agroindustriais. Elaborar orçamentos, laudos, pareceres, relatórios e
projetos, inclusive de incorporação de novas tecnologias. Orientar
projetos de recomposição florestal em propriedades rurais. Aplicar
métodos e programas de melhoramento genético. Aplicar práticas
sustentáveis no manejo de conservação do solo e da água. Prestar
assistência técnica e assessoria no estudo e desenvolvimento de projetos
e pesquisas tecnológicas, ou nos trabalhos de vistoria, perícia,
arbitramento, consultoria, laudos, pareceres e relatórios técnicos.
Prestar assistência técnica nas áreas de crédito rural e agroindustrial,
topografia na área rural, impacto ambiental, paisagismo, jardinagem e
horticultura, construção de benfeitorias rurais, drenagem e irrigação.
Interpretar a análise de solos e aplicar fertilizantes e corretivos nos
tratos culturais em sistema agroecológico. Identificar os processos
simbióticos de absorção, de translocação e os efeitos alelopáticos entre
solo e planta, planejando ações referentes aos tratos das culturas.
Prestar assistência técnica à aplicação, à comercialização e ao manejo
de produtos especializados. Selecionar e aplicar métodos agroecológicos
de controle de vetores e pragas, doenças e plantas daninhas. Planejar e
acompanhar a colheita e a pós-colheita. Supervisionar o armazenamento, a
conservação, a comercialização e a industrialização dos produtos
agroecológicos. Elaborar, aplicar e monitorar programas profiláticos,
higiênicos e sanitários na produção animal, vegetal e agroindustrial.
Emitir laudos e documentos de classificação e exercer a fiscalização de
produtos agroecológicos de origem vegetal, animal e agroindustrial.
Implantar e gerenciar sistemas de controle de qualidade na produção
agroecológica. Treinar e conduzir equipes nas suas modalidades de
atuação profissional. Aplicar as legislações pertinentes ao processo
produtivo e ao meio ambiente. Executar a gestão econômica e financeira
da produção agroecológica. Administrar e gerenciar propriedades
agroecológicas. Operar e manejar máquinas, implementos e equipamentos
agrícolas, veículos aéreos remotamente pilotados e equipamentos de
precisão para monitoramento remoto inerentes ao sistema de produção
agroecológico. Organizar ações integradas de agricultura familiar. Atuar
na certificação agroecológica.

Para a atuação como Técnico em Agroecologia, são fundamentais:

Desenvolvimento de ações socioambientais para a conservação de recursos
naturais aliados à necessidade econômica em sistemas produtivos locais,
inclusive dos povos tradicionais (quilombolas, indígenas, ribeirinhos,
agricultores familiares). Trabalho em equipe a fim de contribuir e
participar, de forma ética e cidadã, com o coletivo na resolução de
problemas, na sustentabilidade, na proposição de ideias e soluções
ambientais, tecnológicas, políticas, econômicas, sociais e culturais, de
acordo com os princípios e ética profissional. Competências e
habilidades para gerenciar sistemas agroecológicos produtivos e a
organização de ações integradas dos povos tradicionais, organizações não
governamentais, empresas públicas ou privadas ligadas às práticas de
produção agroecológica, com base em princípios éticos, humanísticos,
democráticos, inclusivos, sustentáveis e solidários.

\textbf{CBOs correspondidos e seus perfis (presentes no CNCT, e presentes na predição da IA)}

\textbf{CBO: 321105  ( Técnico agrícola )}

Planeja projeto agrícola, pesquisando mercado consumidor e analisando
viabilidade econômica. Avalia condições da propriedade, verificando sua
infraestrutura - máquinas, equipamentos e instalações -; levantando
dados sobre a área a ser trabalhada -- topografia e extensão --;
examinando as condições edafoclimáticas - solo, clima e água -;
verificando a disponibilidade de mão de obra para as atividades;
levantando a capacitação tecnológica do produtor; e pesquisando o
mercado fornecedor de insumos, materiais, máquinas e equipamentos.
Executa levantamento do custo-benefício para o produtor. Verifica as
condições para implantação de projeto agrícola, coletando amostras para
análise de solo e interpretando os resultados laboratoriais; verificando
a disponibilidade e a qualidade da água a ser utilizada; e coletando
dados meteorológicos. Executa projeto agrícola, definindo cultivares e
rotação de culturas. Elabora planta de construções rurais. Compra
máquinas, equipamentos, insumos e materiais. Supervisiona a colheita e a
pós-colheita das principais culturas, aplicando técnicas de adubação
verde, de cobertura morta e de correção do solo, conforme cada caso.
Orienta os processos de silagem e transporte. Supervisiona trabalhadores
agrícolas, contratando mão de obra, distribuindo tarefas, orientando e
acompanhando a execução de atividades. Fiscaliza produção agrícola,
inspecionando documentação de produtos em trânsito, examinando a
sanidade de produtos agrícolas, enviando amostras de produtos para
análises laboratoriais, fiscalizando aplicação de agrotóxicos e
inspecionando cumprimento de normas e padrões técnicos. Conduz
experimentos de pesquisa, elaborando relatórios, laudos, pareceres,
perícias e avaliações. Administra propriedades rurais, adotando sistema
de produção conforme necessidade do mercado e controlando despesas e
receitas. Administra funcionários da propriedade. Comercializa produção
agrícola. Presta assistência e consultoria técnicas, orientando na
escolha do local para atividade agrícola; aconselhando sobre escolha de
espécies e cultivares; instruindo sobre preparo, correção e conservação
de solo; informando sobre equipamentos, máquinas e implementos;
indicando formas e manejo de irrigação e drenagem; aconselhando sobre
manejo integrado de pragas e doenças; orientando sobre época de plantio,
tratos culturais e colheita; e estimulando o beneficiamento de produtos
agrícolas. Orienta sobre preservação ambiental. Promove extensão e
capacitação rural, organizando reuniões com produtores e orientando
formação de associações e grupos de produtores. Assessora produtores,
orientando-os em cultivo com técnicas convencionais ou alternativas de
plantio, em manejo do solo e em compra e venda de insumos, materiais e
produtos agrícolas. Promove a difusão de tecnologia, orientando a
implantação de agricultura de precisão e o uso de sistemas
automatizados, bem como o desenvolvimento de novos cultivares. Dissemina
o uso de ambientes protegidos de cultivo e o melhoramento de sementes.
Divulga o uso de controle remoto da produção e da qualidade de processos
e produtos. Ministra cursos e programas de treinamentos. Apresenta
resultados de pesquisas em encontros e congressos da área agrícola. Pode
participar do desenvolvimento de tecnologias agrícolas, adequando
equipamentos e instalações conforme necessidade da região e do produtor.
Cria e adapta técnicas convencionais e alternativas de cultivo. Orienta
produtores agrícolas sobre biossegurança, indicando uso racional de
agrotóxicos e esclarecendo sobre destino de embalagens de agrotóxicos.
Informa sobre limpeza e desinfeção de máquinas, equipamentos e
instalações.

\begin{Form}
\textbf{Esta correspondência está adequada?} \\ 
\ChoiceMenu[radio,radiosymbol=\ding{52},bordercolor=0 0.5 0,name=cbovalid_ 321105 ]{}{CBO deve ficar=sim, \hspace{6em} CBO não fica aqui=nao}
\end{Form}

\newpage
\section{ 11004 -- Técnico em Agropecuária }

\textbf{Perfil do Curso (CNCT)}

O Técnico em Agropecuária será habilitado para:

Planejar, organizar, dirigir e controlar a produção agropecuária de
forma sustentável, analisando as características econômicas, sociais e
ambientais. Elaborar, projetar e executar projetos de produção
agropecuária, aplicando as Boas Práticas de Produção Agropecuária (BPA).
Prestar assistência técnica e assessoria ao estudo e ao desenvolvimento
de projetos e pesquisas tecnológicas, ou aos trabalhos de vistoria,
perícia, arbitramento e consultoria. Elaborar orçamentos, laudos,
pareceres, relatórios e projetos, inclusive de incorporação de novas
tecnologias. Prestar assistência técnica às áreas de crédito rural e
agroindustrial, de topografia na área rural, de impacto ambiental, de
construção de benfeitorias rurais, de drenagem e irrigação. Planejar,
organizar e monitorar atividades de exploração e manejo do solo, matas e
florestas de acordo com suas características, alternativas de otimização
dos fatores climáticos e seus efeitos no crescimento e desenvolvimento
das plantas e dos animais. Realizar a produção de mudas e sementes, em
propagação em cultivos abertos ou protegidos, em viveiros e em casas de
vegetação. Planejar, organizar e monitorar programas de nutrição e
manejo alimentar em projetos zootécnicos. Planejar, organizar e
monitorar o processo de aquisição, preparo, conservação e armazenamento
da matéria prima e dos produtos agroindustriais. Orientar projetos de
recomposição florestal em propriedades rurais. Aplicar métodos e
programas de melhoramento genético. Prestar assistência técnica na
aplicação, na comercialização, no manejo de produtos especializados e
insumos (sementes, fertilizantes, defensivos, pastagens, concentrados,
sal mineral, medicamentos e vacinas). Interpretar a análise de solos e
aplicar fertilizantes e corretivos nos tratos culturais. Selecionar e
aplicar métodos de erradicação e controle de vetores e pragas, doenças e
plantas daninhas. Planejar e acompanhar a colheita e a pós-colheita.
Supervisionar o armazenamento, a conservação, a comercialização e a
industrialização dos produtos agropecuários. Elaborar, aplicar e
monitorar programas profiláticos, higiênicos e sanitários na produção
animal, vegetal e agroindustrial. Emitir laudos e documentos de
classificação e exercer a fiscalização de produtos de origem vegetal,
animal e agroindustrial. Implantar e gerenciar sistemas de controle de
qualidade na produção agropecuária. Manejar animais por categoria e
finalidade (criação, reprodução, alimentação e sanidade). Aplicar
técnicas de bem-estar animal na produção agropecuária. Treinar e
conduzir equipes nas suas modalidades de atuação profissional. Aplicar
as legislações pertinentes ao processo produtivo e ao meio ambiente.
Aplicar práticas sustentáveis no manejo de conservação do solo e da
água. Identificar e aplicar técnicas mercadológicas para distribuição e
comercialização de produtos agropecuários e animais. Executar a gestão
econômica e financeira da produção agropecuária. Administrar e gerenciar
propriedades rurais. Realizar procedimentos de desmembramento,
parcelamento e incorporação de imóveis rurais. Operar, manejar e regular
máquinas, implementos e equipamentos agrícolas. Operar veículos aéreos
remotamente pilotados e equipamentos de precisão para monitoramento
remoto da produção agropecuária.

Para a atuação como Técnico em Agropecuária, são fundamentais:

Conhecimentos e saberes relacionados à produção agropecuária, à produção
e ao processamento de alimentos, à fitossanidade e à proteção ambiental.
Atualização em relação às inovações tecnológicas. Cooperação de forma
construtiva e colaborativa nos trabalhos em equipe e tomada de decisões.
Adoção de senso investigativo, visão sistêmica das atividades e
processos, capacidade de comunicação e argumentação, autonomia,
proatividade, liderança, respeito às diversidades nos grupos de
trabalho, resiliência frente aos problemas, organização,
responsabilidade, visão crítica, humanística, ética e consciência em
relação ao impacto de sua atuação profissional na sociedade e no
ambiente.

\textbf{CBOs correspondidos e seus perfis (presentes no CNCT, e presentes na predição da IA)}

\textbf{CBO: 321110  ( Técnico agropecuário )}

Planeja projeto agropecuário, pesquisando mercado consumidor e
analisando viabilidade econômica. Avalia condições da propriedade,
verificando sua infraestrutura; levantando dados sobre a área a ser
trabalhada; verificando a disponibilidade de mão de obra para as
atividades; levantando a capacitação tecnológica do produtor; e
pesquisando o mercado fornecedor de insumos, materiais, máquinas e
equipamentos. Executa levantamento do custo-benefício para o produtor.
Verifica as condições para implantação de projeto agropecuário,
coletando amostras para análise de solo e interpretando os resultados
laboratoriais; verificando a disponibilidade e a qualidade da água a ser
utilizada; e coletando dados meteorológicos. Executa projeto
agropecuário, definindo cultivares, raças e espécies e rotação de
culturas. Elabora planta de construções rurais. Compra máquinas,
equipamentos, insumos e materiais. Supervisiona a colheita e a
pós-colheita das principais culturas, aplicando técnicas de adubação
verde, de cobertura morta e de correção do solo, conforme cada caso.
Maneja reprodução de animais e formula rações. Orienta os processos de
silagem e transporte. Supervisiona trabalhadores em atividades
agropecuárias, contratando mão de obra, distribuindo tarefas, orientando
e acompanhando a execução dos trabalhos. Fiscaliza produção
agropecuária, inspecionando documentação de produtos em trânsito,
fiscalizando vacinação de animais, examinando a sanidade de produtos
agropecuários, enviando amostras de produtos agropecuários para análises
laboratoriais, fiscalizando aplicação de agrotóxicos e inspecionando
cumprimento de normas e padrões técnicos. Conduz experimentos de
pesquisa, elaborando relatórios, laudos, pareceres, perícias e
avaliações. Administra propriedades rurais, adotando sistema de produção
conforme necessidade do mercado e controlando despesas e receitas.
Administra funcionários da propriedade. Comercializa produção
agropecuária. Presta assistência e consultoria técnicas, orientando na
escolha do local para atividade agropecuária; informando sobre
alimentação de animais; orientando sobre manejo do desenvolvimento
animal - cria, recria e terminação -; instruindo sobre preparo, correção
e conservação de solo; informando sobre equipamentos, máquinas e
implementos; indicando formas e manejo de irrigação e drenagem;
aconselhando sobre manejo integrado de pragas e doenças; orientando
sobre época de plantio, tratos culturais e colheita; informando sobre
compra e venda de animais; e estimulando o beneficiamento de produtos
agropecuários. Orienta sobre preservação ambiental. Promove extensão e
capacitação rural, organizando reuniões com produtores e orientando
formação de associações e grupos de produtores. Assessora produtores,
orientando-os em cultivo com técnicas convencionais ou alternativas de
plantio, demonstrando técnicas de manejo de animais e informando sobre
compra e venda de insumos, materiais e produtos agropecuários. Promove a
difusão de tecnologia, informando sobre inseminação artificial e
melhoramento genético de animais; apresentando técnicas de controle de
origem e rastreabilidade animal; orientando a implantação de agricultura
e de pecuária de precisão; informando sobre o uso de sistemas
automatizados; e apresentando o desenvolvimento de novos cultivares.
Dissemina o uso de ambientes protegidos de cultivo e o melhoramento de
sementes. Ministra cursos e programas de treinamentos. Apresenta
resultados de pesquisas em encontros e congressos da área agropecuária.
Pode participar do desenvolvimento de tecnologias agropecuárias,
adequando equipamentos e instalações conforme necessidade da região e do
produtor. Orienta produtores sobre biossegurança, indicando uso racional
de agrotóxicos e esclarecendo sobre destino de embalagens de
agrotóxicos. Orienta sobre destino de animais mortos. Informa sobre
limpeza e desinfeção de máquinas, equipamentos e instalações.

\begin{Form}
\textbf{Esta correspondência está adequada?} \\ 
\ChoiceMenu[radio,radiosymbol=\ding{52},bordercolor=0 0.5 0,name=cbovalid_ 321110 ]{}{CBO deve ficar=sim, \hspace{6em} CBO não fica aqui=nao}
\end{Form}

\newpage
\section{ 11005 -- Técnico em Apicultura }

\textbf{Perfil do Curso (CNCT)}

O Técnico em Apicultura será habilitado para:

Planejar, organizar, dirigir e controlar a produção e criação
sustentável de abelhas, analisando as características econômicas,
sociais e ambientais. Elaborar e executar projetos de produção apícola,
inclusive com a incorporação de novas tecnologias, incluindo a criação
de abelhas nativas sem ferrão. Realizar atividades de implantação,
produção apícola, aquisição e manutenção de equipamentos e manejo da
flora apícola. Prestar assistência técnica e assessoria ao estudo e ao
desenvolvimento de projetos e pesquisas tecnológicas e de consultoria.
Prestar assistência técnica às áreas de crédito rural e agroindustrial.
Realizar produção intensiva e artesanal, controle de qualidade,
identi?cação e avaliação da produção de diferentes espécies de abelhas.
Realizar a criação de abelhas para fins de produção de produtos apícolas
alimentares, cosméticos e farmacológicos e para fins de polinização e
preservação ambiental. Realizar a produção e propagação de abelhas
nativas. Aplicar métodos e programas de reprodução e melhoramento
genético. Aplicar as legislações pertinentes ao processo produtivo e ao
meio ambiente. Aplicar Boas Práticas de Produção Apícola. Realizar o
bene?ciamento e processamento de mel, própolis, geleia real e produtos
da atividade apícola. Implantar e gerenciar sistemas de controle de
qualidade na produção apícola. Elaborar, aplicar e monitorar programas
profiláticos, higiênicos e sanitários na produção apícola. Supervisionar
o armazenamento, a conservação, a comercialização e a industrialização
dos produtos apícolas. Identificar e aplicar técnicas mercadológicas
para distribuição e comercialização de produtos apícolas. Treinar e
conduzir equipes nas suas modalidades de atuação profissional.
Administrar e gerenciar apiários e agroindústrias de produtos apícolas.
Executar a gestão econômica e financeira da produção apícola.

Para a atuação como Técnico em Apicultura, são fundamentais:

Conhecimentos e saberes relacionados aos sistemas de produção, ao
beneficiamento de produtos apícolas, aos sistemas de
gestão/comercialização de apiários, às normas técnicas de segurança.
Atualização em relação às inovações tecnológicas. Cooperação de forma
construtiva e colaborativa nos trabalhos em equipe, tomada de decisões.
Adoção de senso investigativo, visão sistêmica das atividades e
processos. Capacidade de comunicação e argumentação, autonomia,
proatividade, liderança, respeito às diversidades nos grupos de
trabalho, resiliência frente aos problemas, organização,
responsabilidade, visão crítica, humanística, ética e consciência em
relação ao impacto da atuação profissional na sociedade e no ambiente.

\textbf{CBOs correspondidos e seus perfis (presentes no CNCT, e presentes na predição da IA)}

\textbf{CBO: 613405  ( Apicultor )}

Planeja atividades de criação de abelhas, definindo uma ou mais
finalidades da criação -- extrações de mel, cera, própolis, pólen,
geleia real e apitoxina (veneno de abelhas) -; e escolhendo espécie de
abelhas. Define a aplicação de técnicas e procedimentos manuais,
semimecanizados ou totalmente mecanizados, com uso, inclusive, de
sistemas automatizados. Utiliza programas informatizados para
planejamento e controle da criação. Estabelece metas e estima custos.
Seleciona o local para o apiário em terreno plano e em área sombreada.
Adquire ou fabrica colmeia (ou caixa de abelhas), confirmando seu
formato-padrão e as medidas das suas peças. Estabelece a criação
racional de abelhas, coletando colônia natural -- principalmente em
troncos de árvores - e colocando-a na colmeia. Supervisiona atividades
de manejo, monitorando as colmeias do apiário. Orienta procedimentos
para substituição de rainha e confirma a postura da nova, verificando a
presença de ovos e larvas. Acompanha a substituição dos quadros da
colmeia. Faz revisões, quando necessário, para avaliar as condições
gerais das colmeias e a ocorrência de anormalidades, interferindo o
mínimo possível na atividade das abelhas. Planeja alimentação das
abelhas, instalando o apiário próximo à flora apícola, conjunto de
plantas usadas pelas abelhas como fonte de alimento (néctar e pólen).
Pode acrescentar à flora apícola outras espécies de plantas, para que
tenha garantia do florescimento em diferentes épocas do ano. Estabelece
forma de aproveitamento de recursos hídricos disponíveis ou constrói
fonte de água para as abelhas. Pode fornecer à colônia alimentação
suplementar, como pasta ``candy''. Controla e combate pragas e doenças.
Providencia proteção contra formigas em instalações e acondiciona cera
para prevenção de traças. Cuida da situação sanitária das colmeias, para
evitar a disseminação de doenças. Reconhece sintomas de doenças e toma
medidas imediatas, como o isolamento das colmeias atacadas e envio de
amostras de material biológico a laboratórios para análise e diagnóstico
precisos. Providencia a execução de procedimentos de inseminação
artificial de abelhas rainhas. Faz a aquisição de rainhas e zangões
geneticamente melhorados, para aperfeiçoamento do plantel de abelhas.
Supervisiona o manejo do apiário, orientando a colheita dos produtos das
abelhas. Monitora a realização da desoperculação e da centrifugação dos
favos de mel e a execução da filtragem e da decantação do mel. Verifica
coleta, limpeza e secagem do pólen apícola. Controla a coleta da
própolis e da geleia real. Orienta coleta e o reprocessamento da cera.
Confere a aplicação de técnicas apropriadas para extrair o veneno das
abelhas (apitoxina). Monitora a preparação de mel e outros produtos
apícolas para comercialização. Observa a realização de pesagem,
classificação, acondicionamento e armazenagem conforme normas de
sanidade e segurança. Comercializa mel e demais produtos apícolas.
Providencia a expedição e o transporte dos produtos vendidos. Mantém-se
atualizado, pesquisando novas tecnologias na criação de abelhas.
Administra propriedade rural, controlando estoques de insumos;
programando logística de transporte de insumos e produtos; e gerenciando
recursos financeiros. Coordena a montagem da infraestrutura do apiário.
Monitora limpeza, higienização e desinfecção das instalações. Seleciona
equipamentos e acessórios e providencia sua manutenção. Supervisiona
atividades das equipes de trabalho, distribuindo tarefas e avaliando o
desempenho dos trabalhadores. Proporciona programas de treinamento às
equipes. Promove práticas -- tais como arborização e proteção de
nascentes da propriedade - para mitigação de danos ambientais. Zela pela
segurança no trabalho, usando e monitorando o uso, pela equipe, de
equipamentos de proteção individual e de vestimenta apícola adequada e
completa; controlando adoção de protocolos de biossegurança; e
prevenindo acidentes.

\begin{Form}
\textbf{Esta correspondência está adequada?} \\ 
\ChoiceMenu[radio,radiosymbol=\ding{52},bordercolor=0 0.5 0,name=cbovalid_ 613405 ]{}{CBO deve ficar=sim, \hspace{6em} CBO não fica aqui=nao}
\end{Form}

\newpage
\section{ 11006 -- Técnico em Aquicultura }

\textbf{Perfil do Curso (CNCT)}

O Técnico em Aquicultura será habilitado para:

Realizar projetos de implantação e de operação de sistemas de cultivos
aquícolas continentais e marinhos. Elaborar projetos aquícolas,
reconhecer o potencial de áreas geográficas para implantação de
empreendimentos e construções aquícolas. Prestar assistência técnica e
assessoria ao estudo e ao desenvolvimento de projetos e pesquisas
tecnológicas, ou aos trabalhos de vistoria, perícia, arbitramento e
consultoria. Elaborar orçamentos, laudos, pareceres, relatórios e
projetos, inclusive de incorporação de novas tecnologias. Utilizar
tecnologias em sistemas de produção e manejo aquícola. Analisar a
viabilidade técnica e econômica de propostas e projetos aquícolas.
Prestar assistência técnica às áreas de crédito rural e agroindustrial,
de topografia na área rural, de impacto ambiental, de construção de
benfeitorias rurais. Operar equipamentos e métodos qualitativos de
análise de água utilizada em sistemas de cultivo. Reconhecer os aspectos
biológicos, fisiológicos e patológicos das principais espécies de
cultivo e aplicar os princípios de nutrição e de manejo alimentar das
principais espécies cultivadas. Realizar procedimentos para reprodução
das principais espécies de interesse aquícola. Aplicar métodos e
programas de melhoramento genético. Aplicar boas práticas de manipulação
e fabricação, e supervisionar as etapas de conservação, processamento,
beneficiamento e comercialização do pescado. Elaborar, aplicar e
monitorar programas profiláticos, higiênicos e sanitários na produção
aquícola. Implantar e gerenciar sistemas de controle de qualidade na
produção aquícola. Emitir laudos e documentos de classificação e exercer
a fiscalização de produtos de origem aquícola. Treinar e conduzir
equipes nas suas modalidades de atuação profissional. Prevenir situações
de risco à segurança no trabalho. Aplicar as legislações pertinentes ao
processo produtivo e ao meio ambiente. Aplicar práticas sustentáveis no
manejo de conservação do solo e da água. Utilizar equipamentos e
programas para fins topográficos e georreferenciamento. Identificar e
aplicar técnicas mercadológicas para distribuição e comercialização de
produtos aquícolas. Executar a gestão econômica e financeira da produção
aquícolas. Administrar e gerenciar propriedades aquícolas.

Para a atuação como Técnico em Aquicultura, são fundamentais:

Conhecimentos e saberes relacionados a processos de produção e
reprodução de animais aquáticos, como peixes, camarões, mexilhões,
ostras, rãs, entre outros. Conhecimentos relacionados ao monitoramento
da qualidade da água, ao controle sanitário de organismos aquáticos e às
boas práticas de manipulação e beneficiamento do pescado. Conhecimentos
relacionados à gestão de negócios voltados à aquicultura, à legislação
ambiental, ao planejamento de produção, à gestão de projetos e de
processos, ao empreendedorismo, ao mercado e à comercialização do
pescado, à extensão pesqueira, à aquicultura em estabelecimentos rurais,
à aquicultura em águas da União, às políticas públicas para o
desenvolvimento da aquicultura, ao associativismo e ao cooperativismo.
Domínio de uso de tecnologias da informação e bases tecnológicas,
habilidade de comunicação e resolução de situações-problema, habilidade
para o trabalho em equipe e para gestão de conflitos.

\textbf{CBOs correspondidos e seus perfis (presentes no CNCT, e presentes na predição da IA)}

\textbf{CBO: 321305  ( Técnico em piscicultura )}

Elabora projeto de implantação e de operação de sistema de criação de
peixes, de acordo com as características da região e com a aptidão do
ambiente natural. Faz análise da viabilidade técnica e econômica do
projeto. Apresenta as etapas de implantação. Relaciona máquinas,
equipamentos, ferramentas, instrumentos e materiais necessários. Detalha
tamanho de equipe de trabalho e suas atribuições. Elabora orçamento.
Presta orientações técnicas na implantação do projeto de piscicultura.
Seleciona sistema extensivo, semi-intensivo ou intensivo de criação.
Analisa os aspectos biológicos, fisiológicos e patológicos das
principais espécies e examina as condições do local, para recomendar a
espécie de peixes para criação. Orienta a preparação de locais -- tais
como viveiros, tanques escavados e fazendas marinhas -, para a criação
de peixes. Executa levantamento topográfico. Indica a necessidade de
drenagem, colocação de telas, vedação de comportas e outras
providências, conforme infraestrutura predefinida de criação. Acompanha
a montagem de equipamentos e sistemas de aeração. Avalia qualidade da
água utilizada no sistema de criação. Monitora o uso racional da água,
dimensionando sua vazão. Pode recomendar a instalação de viveiros
automatizados. Orienta o processo de reprodução natural ou artificial de
peixes, preferencialmente com espécies e matrizes melhoradas
geneticamente. Monitora a seleção, a mensuração e o acasalamento de
reprodutores, medindo fertilidade e testando grau de maturação de
óvulos. Acompanha a aplicação de técnicas artificiais de reprodução.
Monitora o desenvolvimento dos ovos e das larvas. Pode orientar a
aplicação de técnica de reversão sexual. Organiza o processo de
alimentação dos peixes, calculando seu peso médio e indicando tipo e
quantidade de alimentos. Define horários de alimentação. Demonstra o
processo de preparação de rações. Controla a sanidade dos peixes,
observando o comportamento dos espécimes para identificar eventuais
doenças e irregularidades da criação, como ferimentos ou presença de
fungos, bactérias e parasitas. Indica a necessidade de visita de médico
veterinário ou zootecnista, para ações e prescrições pertinentes.
Monitora a realização dos tratamentos, como isolamento de indivíduos em
quarentena e aplicação dos medicamentos prescritos na ração ou na água.
Orienta o uso de emissão de sinais sonoros; de alarmes automatizados com
sensores; e de outras estratégias, para afastar e coibir predadores.
Demonstra técnicas de anestesiar, insensibilizar ou promover a perda
imediata da consciência dos peixes, para realizar o abate sem
desconforto ou dor. Presta informações técnicas para execução do
beneficiamento dos peixes abatidos, orientando limpeza, descamação,
evisceração, corte, filetagem, fatiamento, moagem e congelamento.
Classifica os peixes. Demonstra a forma de embalar os peixes. Monitora o
transporte e o armazenamento da produção. Orienta a comercialização da
produção. Implanta e gerencia sistemas de controle de qualidade de
processos e produtos. Observa a organização do local de trabalho.
Monitora a aplicação de programas profiláticos, higiênicos e sanitários.
Orienta equipe para conservar ferramentas e instrumentos limpos,
acondicionados e em plenas condições de uso e funcionamento. Prepara
plano de manutenção de máquinas e equipamentos. Elabora relatório
técnico sobre resultados da implantação de projeto de piscicultura. Pode
ministrar palestras. Pode administrar empreendimentos e realizar
pesquisa e extensão, na área de piscicultura. Pode supervisionar equipe
de trabalho, distribuindo tarefas, avaliando desempenho e desenvolvendo
treinamentos. Mantém-se atualizado em relação às inovações na área de
piscicultura. Contribui para a redução dos impactos ambientais do
processo de criação, monitorando a aplicação das normas ambientais. Zela
pela segurança, prevenindo situações de risco. Utiliza e monitora o uso
de equipamentos de proteção individual. Presta primeiros socorros.

\begin{Form}
\textbf{Esta correspondência está adequada?} \\ 
\ChoiceMenu[radio,radiosymbol=\ding{52},bordercolor=0 0.5 0,name=cbovalid_ 321305 ]{}{CBO deve ficar=sim, \hspace{6em} CBO não fica aqui=nao}
\end{Form}

\textbf{CBO: 321310  ( Técnico em carcinicultura )}

Elabora projeto de implantação e de operação de sistema de criação de
camarões marinhos ou de camarões de água doce em cativeiro. Faz análise
da viabilidade técnica e econômica do projeto. Apresenta as etapas de
implantação. Relaciona máquinas, equipamentos, ferramentas, instrumentos
e materiais necessários. Detalha tamanho de equipe de trabalho e suas
atribuições. Elabora orçamento. Presta orientações técnicas na
implantação do projeto de carcinicultura, conforme características da
área, do clima e da localização geográfica. Seleciona sistema extensivo,
semi-intensivo ou intensivo de criação. Recomenda espécie de camarões
para criação. Orienta a construção e a adaptação de instalações - no mar
ou na terra - para a criação de camarões. Executa levantamento
topográfico. Indica necessidade de drenagem, colocação de telas e outras
providências. Orienta a montagem de sistemas de aeração. Pode recomendar
a instalação de viveiros automatizados. Orienta monitoramento constante
de parâmetros da água. Demonstra a forma de coletar amostras e de medir
temperatura, pH, oxigênio, transparência, entre outros parâmetros.
Indica o modo de instalar filtros de purificação e de dimensionar a
vazão da água. Monitora, antes do povoamento do viveiro com as
pós-lavas, a aplicação de calcário e de fertilizantes químicos e
orgânicos, que contribuirão na produção da alimentação natural, formada
por microorganismos, tais como algas, fito e zooplânctons, entre outros.
Além da alimentação natural, orienta o fornecimento de ração balanceada.
Seleciona e recomenda suplementos alimentares. Indica procedimentos
técnicos para aquisição de pós-larvas, para sua transferência ao local
de criação e para sua aclimatação. Pode organizar a reprodução de
camarões, monitorando a seleção, a mensuração e o acasalamento de
reprodutores. Pode orientar a aplicação de técnicas artificiais de
reprodução. Informa sobre técnicas de melhoramento genético. Recomenda o
emprego de tela de proteção durante o abastecimento de água no local da
criação, para impedir a entrada de espécies competidoras e predadoras.
Implanta controle da sanidade dos camarões, com observação de espécimes
e coleta de amostra dirigida - dos que apresentem aspectos menos
saudáveis - de diferentes pontos do viveiro. Orienta a pesagem dos
camarões e imediato envio para a avaliação de tecidos e órgãos, a fim de
identificar enfermidades virais e bacterianas, entre outras doenças e
irregularidades. Recomenda, ainda, o controle de índices zootécnicos -
como taxa de crescimento, peso corporal, entre outros -- indicativos da
saúde dos camarões em cativeiro. Orienta a despesca, a desinfecção e o
abate dos camarões. Organiza o beneficiamento dos camarões abatidos,
monitorando os processos, tais como limpeza, processamento em filés,
classificação, congelamento e embalagem. Monta a logística de
armazenagem e transporte da produção. Implanta e gerencia sistemas de
controle de qualidade de processos e produtos. Observa a organização do
local de trabalho. Monitora a aplicação de programas profiláticos,
higiênicos e sanitários. Orienta a conservação de ferramentas e
instrumentos em plenas condições de funcionamento. Prepara plano de
manutenção de máquinas e equipamentos. Elabora relatório técnico sobre
resultados do projeto de carcinicultura. Pode ministrar palestras. Pode
administrar empreendimentos e realizar pesquisa e extensão, na área de
carcinicultura. Pode supervisionar trabalho de equipe, avaliando
desempenho e desenvolvendo treinamentos. Mantém-se atualizado em relação
às inovações em carcinicultura. Estuda aplicação de novas tecnologias,
como o sistema de criação em bioflocos bacterianos para criação de
camarões marinhos em locais afastados da costa. Contribui para redução
dos impactos ambientais do processo de criação, monitorando a aplicação
das normas ambientais. Zela pela segurança, prevenindo situações de
risco. Utiliza e monitora o uso de equipamentos de proteção individual.
Presta primeiros socorros.

\begin{Form}
\textbf{Esta correspondência está adequada?} \\ 
\ChoiceMenu[radio,radiosymbol=\ding{52},bordercolor=0 0.5 0,name=cbovalid_ 321310 ]{}{CBO deve ficar=sim, \hspace{6em} CBO não fica aqui=nao}
\end{Form}

\textbf{CBO: 321315  ( Técnico em mitilicultura )}

Elabora projeto de implantação e de operação de sistema de criação de
mexilhões no mar. Faz análise da viabilidade técnica e econômica do
projeto. Apresenta as etapas de implantação. Relaciona equipamentos --
como as máquinas de semear, colher e desagregar mexilhões -,
embarcações, ferramentas, instrumentos e materiais necessários. Detalha
tamanho de equipe de trabalho e suas atribuições. Elabora orçamento.
Presta orientações técnicas na implantação do projeto de mitilicultura,
conforme características do ambiente de cultivo e da temperatura do
local. Recomenda espécie para criação. Indica local apropriado para
instalação de cultivo de mexilhões. Seleciona o sistema de criação.
Orienta a seleção e a instalação de boias e cabos. Mostra a importância
do conhecimento do regime de reprodução do mexilhão na região de
cultivo, para definir as épocas de captação de sementes. Indica formas
de captação, detalhando a explotação de bancos naturais; o uso de
coletores artificiais; o aproveitamento de sementes que se assentam
sobre as estruturas de cultivo, entre outras. Recomenda classificação de
sementes por tamanho, para processo de semeadura. Presta informações
técnicas sobre o manejo dos mexilhões. Orienta o controle da densidade
de sementes por metro de corda de cultivo. Monitora os mexilhões
submersos, que filtram a água do mar respirando e se alimentando.
Organiza o controle periódico e contínuo no local, para prevenir a
contaminação dos mexilhões pelo acúmulo de toxinas em seu sistema
digestivo ou tecidos. Indica a necessidade do controle constante nas
áreas de cultivo, testando a qualidade microbiológica da água e
avaliando uma amostra de mexilhões. Indica formas de controle dos
fatores de riscos e ameaças, tais como ataque às sementes por peixes e
outros predadores; roubo e vandalismo; e eventos climáticos extremos.
Orienta a avaliação de uma amostra de mexilhões antes da colheita,
verificando, entre outros fatores, o tamanho alcançado em relação ao
desejável no mercado alvo. Destaca a importância do uso de técnicas de
manuseio adequadas da produção, para garantir melhor qualidade e preço.
Estabelece a logística de armazenagem e de transporte da produção.
Orienta a comercialização da produção. Recomenda que os mexilhões sejam
transportados para a planta de processamento o mais rapidamente
possível, para entrarem na cadeia de frio. Implanta e gerencia sistemas
de controle de qualidade de processos e produtos. Elabora relatório
técnico sobre resultados da implantação de projeto de mitilicultura.
Pode ministrar palestras. Pode orientar a formação de grupos de
mitilicultores - como associações - para, entre outros objetivos,
promover a melhoria do rendimento econômico dos participantes. Pode
elaborar projeto de criação de mexilhões em água doce - lagos, lagoas e
em águas correntes, como riachos e rios - e orientar sua implantação.
Pode formular e implantar projetos de malacocultura, com atividade de
cultivo -- além de mexilhões - de ostras, vieiras e outros moluscos
bivalves. Pode administrar empreendimentos e realizar pesquisa e
extensão, na área de mitilicultura. Pode supervisionar equipe de
trabalho, distribuindo tarefas, avaliando desempenho e desenvolvendo
treinamentos. Mantém-se atualizado em relação às inovações relacionadas
à área de mitilicultura. Seleciona e avalia inovações - como
automatização de processos -- para adoção em novos projetos. Destaca a
importância da organização do local de trabalho e da aplicação de
programas profiláticos, higiênicos e sanitários. Orienta equipe para
conservar ferramentas e instrumentos limpos, acondicionados e em plenas
condições de uso e funcionamento. Prepara plano de manutenção de
embarcações, máquinas e equipamentos. Contribui para a redução dos
impactos ambientais do processo de criação, monitorando a aplicação das
normas ambientais. Zela pela segurança, prevenindo situações de risco.
Utiliza e monitora o uso de equipamentos de proteção individual. Presta
primeiros socorros.

\begin{Form}
\textbf{Esta correspondência está adequada?} \\ 
\ChoiceMenu[radio,radiosymbol=\ding{52},bordercolor=0 0.5 0,name=cbovalid_ 321315 ]{}{CBO deve ficar=sim, \hspace{6em} CBO não fica aqui=nao}
\end{Form}

\textbf{CBO: 321320  ( Técnico em ranicultura )}

Elabora projeto de implantação e de operação de sistema de criação de
rãs em cativeiro. Faz análise da viabilidade técnica e econômica do
projeto. Apresenta as etapas de implantação. Relaciona máquinas,
equipamentos, ferramentas, instrumentos e materiais necessários. Detalha
tamanho de equipe de trabalho e suas atribuições. Elabora orçamento.
Orienta a implantação do projeto de ranicultura, realizada considerando
clima na região; características do terreno; disponibilidade de água; e
distância dos centros de comercialização. Seleciona espécie de rãs para
criação. Indica sistema de criação. Recomenda a instalação de ranário em
terrenos com topografia preferencialmente plana ou pouco acidentada.
Orienta a implantação da estrutura do ranário, com instalações para
todas as fases - reprodução; larvicultura ou girinagem; engorda inicial
ou recria; e engorda final ou terminação - da vida de rãs. Enfatiza a
importância da disponibilidade de água com qualidade física e química
específica. Mostra como são feitos os exames dos principais parâmetros
da água, tais como oxigênio, pH, condutividade elétrica, alcalinidade
total, fósforo, cloretos, ferro, entre outros. Organiza as práticas de
manejo, registrando rãs para controle de seu histórico e seu
desenvolvimento; providenciando a limpeza das baias e a renovação da
água; e controlando a densidade recomendada de animais nas instalações.
Orienta a oferta de alimentos naturais e de ração. Mostra como cuidar do
bem-estar do plantel. Organiza o processo de reprodução natural de
machos e fêmeas aptos ao acasalamento. Principalmente em ranários de
maior porte ou naqueles especializados na produção de imagos para a
venda, mostra como executar a reprodução artificial, controlando a
temperatura, o fotoperíodo e a umidade do local; e submetendo os animais
a tratamento hormonal específico, para obter desovas semelhantes às
naturais. Indica procedimentos para controle da sanidade, como constante
renovação da água e limpeza dos tanques. Orienta o fornecimento de
informações - sobre sinais de possíveis doenças e lesões - ao médico
veterinário, que fará a prescrição e a aplicação de medicamentos, quando
necessário. Indica como fazer a coleta de animais mortos e o seu
transporte a um pequeno forno crematório na propriedade, para
erradicação de quaisquer agentes infecciosos presentes no criatório.
Orienta a seleção de rãs para o abate, respeitando um padrão de tamanho
e peso e verificando a ausência de feridas, traumas, deformações e
malformações. Demonstra como promover a insensibilidade aos animais, de
modo que os procedimentos de abate possam ser realizados sem dor ou
sofrimento. Organiza as etapas de processamento de abate e de preparação
das rãs abatidas para comercialização. Presta informações técnicas sobre
a comercialização da carne congelada de rãs e sobre a venda de
subprodutos - tais como vísceras, pele e gordura - do abate. Implanta e
gerencia sistemas de controle de qualidade de processos e produtos.
Elabora relatório técnico sobre resultados do projeto de ranicultura.
Pode ministrar palestras. Pode administrar empreendimentos e realizar
pesquisa e extensão, na área de ranicultura. Pode supervisionar trabalho
de equipe, avaliando desempenho e desenvolvendo treinamentos. Mantém-se
atualizado em relação às inovações em ranicultura. Seleciona e avalia
inovações - como automatização de processos -- para adoção em novos
projetos. Observa a organização do local de trabalho. Monitora a
aplicação de programas profiláticos, higiênicos e sanitários. Orienta a
conservação de ferramentas e instrumentos em plenas condições de
funcionamento. Prepara plano de manutenção de máquinas e equipamentos.
Contribui para redução dos impactos ambientais do processo de criação,
monitorando a aplicação das normas ambientais. Zela pela segurança,
prevenindo situações de risco. Utiliza e monitora o uso de equipamentos
de proteção individual. Presta primeiros socorros.

\begin{Form}
\textbf{Esta correspondência está adequada?} \\ 
\ChoiceMenu[radio,radiosymbol=\ding{52},bordercolor=0 0.5 0,name=cbovalid_ 321320 ]{}{CBO deve ficar=sim, \hspace{6em} CBO não fica aqui=nao}
\end{Form}

\newpage
\section{ 11007 -- Técnico em Cafeicultura }

\textbf{Perfil do Curso (CNCT)}

O Técnico em Cafeicultura será habilitado para:

Planejar, organizar, dirigir e controlar os processos de implantação,
condução do sistema produtivo cafeeiro, bem como os processos de
colheita, processamento e preparo pós-colheita do café. Prestar
assistência técnica e assessoria ao estudo e ao desenvolvimento de
projetos e pesquisas tecnológicas, ou aos trabalhos de vistoria e
consultoria na cadeia produtiva do café. Elaborar orçamentos, laudos,
pareceres, relatórios e projetos, inclusive de incorporação de novas
tecnologias. Aplicar as legislações pertinentes ao processo produtivo e
ao meio ambiente. Identificar as principais espécies e cultivares de
café. Classificar, beneficiar, industrializar e comercializar produtos
do café. Empregar técnicas para produção de sementes e mudas de
cafeeiro. Realizar o preparo mecânico ou manual do solo e a implantação
da lavoura cafeeira. Aplicar práticas sustentáveis no manejo de
conservação do solo e da água. Recomendar e aplicar os processos de
adubação, correção e conservação do solo para a cultura do café.
Identificar os processos simbióticos de absorção, de translocação e os
efeitos alelopáticos entre solo e planta, planejando ações referentes
aos tratos das culturas. Efetuar tratos culturais em todos os estágios
de desenvolvimento da cultura. Identificar e promover o manejo integrado
de pragas, doenças e plantas espontâneas. Executar levantamentos
topográficos e projetos de geoprocessamento. Conduzir o preparo da área
e definir o sistema de cultivo e manejo. Dimensionar o sistema,
selecionar e manejar projetos de irrigação para a cultura do café.
Dominar os procedimentos para coleta de amostras de solos e folhas,
analisar e interpretar seus os resultados laboratoriais. Planejar a
manutenção de máquinas e equipamentos agrícolas e operá-los. Treinar e
conduzir equipes nas suas modalidades de atuação profissional. Aplicar
Boas Práticas de Produção Agrícola (BPA). Orientar as ações necessárias
à colheita, ao processamento, ao beneficiamento e ao armazenamento do
café. Implantar e gerenciar sistemas de controle de qualidade na
produção cafeeira. Responsabilizar-se pela implantação de cafezais,
acompanhando seu desenvolvimento até a fase produtiva, emitindo os
respectivos certificados de origem e qualidade de produtos.
Supervisionar sistemas de certificação e rastreabilidade na cultura do
café. Administrar a propriedade cafeeira. Executar e planejar a gestão
financeira da propriedade cafeeira. Orientar a gestão de riscos
relacionada à comercialização do café. Prestar assistência técnica às
áreas de crédito rural e agroindustrial. Desenvolver projetos
agroecológicos para a cultura do café. Dominar os procedimentos de
classificação e degustação de café.

Para a atuação como Técnico em Cafeicultura, são fundamentais:

Conhecimentos e saberes relacionados à produção cafeeira. Atualização em
relação às inovações tecnológicas. Cooperação de forma construtiva e
colaborativa nos trabalhos em equipe e na tomada de decisões. Adoção de
senso investigativo, visão sistêmica das atividades e processos,
capacidade de comunicação e argumentação, autonomia, proatividade,
liderança, respeito às diversidades nos grupos de trabalho, resiliência
frente aos problemas, organização, responsabilidade, visão crítica,
humanística, ética e consciência em relação ao impacto da atuação
profissional na sociedade e no ambiente.

\textbf{CBOs correspondidos e seus perfis (presentes no CNCT, e presentes na predição da IA)}

\textbf{CBO: 321105  ( Técnico agrícola )}

Planeja projeto agrícola, pesquisando mercado consumidor e analisando
viabilidade econômica. Avalia condições da propriedade, verificando sua
infraestrutura - máquinas, equipamentos e instalações -; levantando
dados sobre a área a ser trabalhada -- topografia e extensão --;
examinando as condições edafoclimáticas - solo, clima e água -;
verificando a disponibilidade de mão de obra para as atividades;
levantando a capacitação tecnológica do produtor; e pesquisando o
mercado fornecedor de insumos, materiais, máquinas e equipamentos.
Executa levantamento do custo-benefício para o produtor. Verifica as
condições para implantação de projeto agrícola, coletando amostras para
análise de solo e interpretando os resultados laboratoriais; verificando
a disponibilidade e a qualidade da água a ser utilizada; e coletando
dados meteorológicos. Executa projeto agrícola, definindo cultivares e
rotação de culturas. Elabora planta de construções rurais. Compra
máquinas, equipamentos, insumos e materiais. Supervisiona a colheita e a
pós-colheita das principais culturas, aplicando técnicas de adubação
verde, de cobertura morta e de correção do solo, conforme cada caso.
Orienta os processos de silagem e transporte. Supervisiona trabalhadores
agrícolas, contratando mão de obra, distribuindo tarefas, orientando e
acompanhando a execução de atividades. Fiscaliza produção agrícola,
inspecionando documentação de produtos em trânsito, examinando a
sanidade de produtos agrícolas, enviando amostras de produtos para
análises laboratoriais, fiscalizando aplicação de agrotóxicos e
inspecionando cumprimento de normas e padrões técnicos. Conduz
experimentos de pesquisa, elaborando relatórios, laudos, pareceres,
perícias e avaliações. Administra propriedades rurais, adotando sistema
de produção conforme necessidade do mercado e controlando despesas e
receitas. Administra funcionários da propriedade. Comercializa produção
agrícola. Presta assistência e consultoria técnicas, orientando na
escolha do local para atividade agrícola; aconselhando sobre escolha de
espécies e cultivares; instruindo sobre preparo, correção e conservação
de solo; informando sobre equipamentos, máquinas e implementos;
indicando formas e manejo de irrigação e drenagem; aconselhando sobre
manejo integrado de pragas e doenças; orientando sobre época de plantio,
tratos culturais e colheita; e estimulando o beneficiamento de produtos
agrícolas. Orienta sobre preservação ambiental. Promove extensão e
capacitação rural, organizando reuniões com produtores e orientando
formação de associações e grupos de produtores. Assessora produtores,
orientando-os em cultivo com técnicas convencionais ou alternativas de
plantio, em manejo do solo e em compra e venda de insumos, materiais e
produtos agrícolas. Promove a difusão de tecnologia, orientando a
implantação de agricultura de precisão e o uso de sistemas
automatizados, bem como o desenvolvimento de novos cultivares. Dissemina
o uso de ambientes protegidos de cultivo e o melhoramento de sementes.
Divulga o uso de controle remoto da produção e da qualidade de processos
e produtos. Ministra cursos e programas de treinamentos. Apresenta
resultados de pesquisas em encontros e congressos da área agrícola. Pode
participar do desenvolvimento de tecnologias agrícolas, adequando
equipamentos e instalações conforme necessidade da região e do produtor.
Cria e adapta técnicas convencionais e alternativas de cultivo. Orienta
produtores agrícolas sobre biossegurança, indicando uso racional de
agrotóxicos e esclarecendo sobre destino de embalagens de agrotóxicos.
Informa sobre limpeza e desinfeção de máquinas, equipamentos e
instalações.

\begin{Form}
\textbf{Esta correspondência está adequada?} \\ 
\ChoiceMenu[radio,radiosymbol=\ding{52},bordercolor=0 0.5 0,name=cbovalid_ 321105 ]{}{CBO deve ficar=sim, \hspace{6em} CBO não fica aqui=nao}
\end{Form}

\textbf{CBO: 325205  ( Técnico de alimentos )}

Planeja as rotinas de trabalho, interpretando as ordens de serviço;
elaborando o cronograma de produção; especificando e calculando os
materiais e os insumos; selecionando os procedimentos para cada
atividade; e verificando a disponibilidade dos equipamentos, materiais,
insumos e gêneros alimentícios. Supervisiona os processos de produção,
recebendo insumos, produtos e gêneros alimentícios; acompanhando
pré-preparo e preparo de alimentos; monitorando processos -- tais como
de trituração, pasteurização, mistura, fermentação e cocção -; e
assegurando as condições operacionais do processo produtivo. Corrige
desvios nos processos manuais e automatizados. Controla o tempo de
produção. Controla o peso e a dimensão do produto. Verifica as condições
da embalagem do produto final. Acompanha a distribuição de produtos até
os pontos de venda. Avalia a qualidade de produtos acabados, coletando
amostras dos produtos e realizando análises microbiológicas, sensoriais
e físico-químicas de matérias-primas, insumos e produtos intermediários
e acabados. Realiza testes de desempenho de matérias-primas, insumos e
produtos intermediários e acabados. Controla a qualidade nas etapas de
produção, identificando e corrigindo pontos críticos de controle e
assegurando as condições higiênico-sanitárias. Verifica as condições de
armazenamento, controlando a data de vencimento dos produtos e
acompanhando o controle integrado de pragas e vetores. Participa, sob
supervisão, de pesquisas e de projetos para melhoria, adequação e
desenvolvimento de produtos, acompanhando as necessidades do mercado e
participando na elaboração da formulação e na testagem dos novos
produtos. Assessora a implementação das mudanças aprovadas. Participa na
elaboração da informação nutricional, que consta na rotulagem do
produto. Orienta clientes internos e externos sobre aspectos
nutricionais dos produtos. Coordena e supervisiona equipes de trabalho,
participando na seleção de pessoal, orientando as atividades,
remanejando pessoal, avaliando os resultados de desempenho e realizando
treinamento de rotinas operacionais. Controla as condições do ambiente
de trabalho, monitorando a temperatura e a umidade do ar. Acompanha os
procedimentos de sanitização das instalações. Verifica condições de
segurança ambiental e controla os pontos críticos e o uso de
equipamentos de proteção individual. Monitora o funcionamento de
máquinas, equipamentos e instrumentos, solicitando serviço de
manutenção, quando necessário. Controla desperdícios e destina
adequadamente os resíduos gerados, fazendo reciclagem e realizando
descarte de acordo com as normas ambientais.

\begin{Form}
\textbf{Esta correspondência está adequada?} \\ 
\ChoiceMenu[radio,radiosymbol=\ding{52},bordercolor=0 0.5 0,name=cbovalid_ 325205 ]{}{CBO deve ficar=sim, \hspace{6em} CBO não fica aqui=nao}
\end{Form}

\newpage
\section{ 11008 -- Técnico em Florestas }

\textbf{Perfil do Curso (CNCT)}

O Técnico em Florestas será habilitado para:

Planejar, organizar, dirigir e controlar as atividades de preservação,
implantação, conservação e utilização de florestas e produtos de origem
florestal, analisando as características econômicas, sociais e
ambientais. Prestar assistência técnica e assessoria ao estudo e ao
desenvolvimento de projetos e pesquisas tecnológicas, ou aos trabalhos
de vistoria, perícia, arbitramento e consultoria. Elaborar orçamentos,
laudos, pareceres, relatórios e projetos, inclusive de incorporação de
novas tecnologias. Prestar assistência técnica às áreas de crédito rural
e de fomento florestal. Planejar, organizar e monitorar atividades de
exploração e manejo do solo, matas e florestas de acordo com suas
características, alternativas de otimização dos fatores climáticos e
seus efeitos no crescimento e desenvolvimento das plantas e dos animais,
propagação em cultivos abertos ou protegidos, em viveiros e em casas de
vegetação. Aplicar métodos e programas de melhoramento genético
florestal. Prestar assistência técnica à aplicação, à comercialização,
ao manejo de produtos especializados, à recomendação e à interpretação
de análise de solos, à aplicação de fertilizantes e corretivos nos
tratos culturais. Identificar os processos simbióticos de absorção, de
translocação e os efeitos alelopáticos entre solo e planta, planejando
ações referentes aos tratos das culturas. Selecionar e aplicar métodos
de erradicação e controle de vetores e pragas, doenças e plantas
daninhas. Supervisionar a execução de atividades florestais, como a
construção de viveiros florestais, produção de mudas, colheita florestal
com extração e beneficiamento da madeira, manejo de florestas nativas e
comercialização. Supervisionar o armazenamento, a conservação, a
comercialização e a industrialização dos produtos de origem florestal.
Emitir laudos e documentos de classificação, exercer a fiscalização de
produtos de origem florestal. Executar o processo de produção, manejo
sustentável e industrialização dos recursos de origem florestal.
Orientar a prática florestal de menor impacto ambiental. Aplicar as
legislações pertinentes ao processo produtivo e ao meio ambiente.
Aplicar práticas sustentáveis no manejo de conservação do solo e da
água. Orientar e planejar projetos de recomposição florestal em
propriedades rurais. Supervisionar atividades de coleta de dados dentro
dos povoamentos florestais para fins de elaboração de inventários
florestais Utilizar equipamentos e reconhecer os métodos utilizados nas
medições das árvores no campo e aplicações em inventário florestal.
Realizar a coleta, a identificação e a conservação de sementes
florestais. Conhecer os principais produtos florestais derivados da
madeira. Administrar unidades de conservação e de produção florestal.
Fiscalizar e monitorar a fauna e a flora silvestres. Realizar a
identificação botânica de espécies florestais. Elaborar, projetar e
executar projetos de arborização urbana e jardins. Treinar e conduzir
equipes nas suas modalidades de atuação profissional. Identificar e
aplicar técnicas mercadológicas para distribuição e comercialização de
produtos florestais. Executar a gestão econômica e financeira da
produção florestal. Administrar e gerenciar propriedades de produção
florestal. Realizar procedimentos de desmembramento, parcelamento e
incorporação de imóveis rurais. Analisar bancos de dados espaciais e
realizar sensoriamento remoto. Utilizar máquinas e implementos
específicos para a atividade florestal. Utilizar equipamentos e
programas para fins topográficos e georreferenciamento. Realizar
levantamento, coleta, processamento e análise de dados através do
Sistema de Navegação Global por Satélite (GNSS).

Para a atuação como Técnico em Florestas, são fundamentais:

Conhecimentos e saberes relacionados à produção florestal, à tecnologia
da madeira e à proteção ambiental. Atualização em relação às inovações
tecnológicas. Cooperação de forma construtiva e colaborativa nos
trabalhos em equipe e na tomada de decisões. Adoção de senso
investigativo, visão sistêmica das atividades e processos, capacidade de
comunicação e argumentação, autonomia, proatividade, liderança, respeito
às diversidades nos grupos de trabalho, resiliência frente aos
problemas, organização, responsabilidade, visão crítica, humanística,
ética e consciência em relação ao impacto de sua atuação profissional na
sociedade e no ambiente.

\textbf{CBOs correspondidos e seus perfis (presentes no CNCT, e presentes na predição da IA)}

\textbf{CBO: 321205  ( Técnico em madeira )}

Planeja atividades de colheita florestal e uso racional da madeira.
Elabora planilha de custos e cronograma de execução física e operacional
de projetos. Dimensiona o quadro de pessoal e prevê infraestrutura de
instalações e equipamentos para as atividades. Administra unidades de
colheita florestal e de processamento de produtos madeireiros,
projetando o volume de extração de madeira e de fabricação de produtos e
subprodutos. Estabelece procedimentos e padrões de qualidade, para a
extração e a produção. Supervisiona a execução de atividades florestais,
monitorando a correta identificação das espécies florestais e orientando
a colheita de madeira. Controla a cubagem e a etiquetagem da origem dos
lotes de madeira bruta, seguindo as normas legais. Controla os
parâmetros de secagem e tratamento da madeira -- produção intermediária
- em estufas ou ao ar livre. Supervisiona a execução de projetos de
fabricação de produtos e subprodutos de madeira, coordenando processos
de usinagem -- ou operando máquinas-ferramenta - e utilizando técnicas
de montagem e acabamento de produtos. Pode usar máquinas automatizadas
ou que incorporem outras novas tecnologias. Controla qualidade da
produção. Utiliza galpões e outras instalações para armazenar madeira
bruta e produtos madeireiros. Providencia acessos e vias para escoamento
da produção, definindo a logística de carregamento e transporte de
madeira bruta e de produtos madeireiros. Coordena equipes de trabalho,
participando de processo de seleção e de administração de pessoal.
Distribui tarefas e programa férias das equipes. Controla conflitos nas
equipes. Supervisiona equipes de trabalho, avaliando desempenho e
levantando as necessidades de treinamento. Viabiliza a realização de
programas de aperfeiçoamento e atualização profissional. Orienta o uso
de máquinas e implementos específicos para a atividade florestal e para
o processamento da madeira. Participa da administração dos custos e
recursos, controlando estoques e monitorando a execução financeira das
atividades. Realiza extensão florestal, sistematizando informações
socioeconômicas da comunidade; demonstrando viabilidade econômica de
produtos e métodos de produção; e assessorando criação de cooperativas.
Ministra programas de treinamento e palestras para pessoas de
comunidades e empresas, orientando-as sobre utilização racional de
recursos naturais renováveis e sobre legislação ambiental. Pode atuar em
pesquisa, colaborando na realização de experimentos técnico-científicos.
Pode participar de levantamentos fitossociológicos e fitossanitários.
Elabora documentação técnica -- tais como relatórios de atividades
operacionais - e mantém arquivos -- físicos e digitais - de inventários
florestais, cadastros, documentos fiscais, entre outros documentos.
Providencia regularizações e autorizações de instâncias superiores e de
órgãos competentes, com relação aos certificados de origem. Monitora a
limpeza e a organização dos locais de trabalho. Zela pela manutenção de
máquinas e equipamentos, elaborando plano de manutenção preventiva e
preditiva, providenciando manutenção de primeiro nível e solicitando
serviços de manutenção corretiva. Trabalha com segurança, fazendo uso
dos equipamentos de proteção individual. Fiscaliza a utilização de
equipamentos de proteção individual e coletiva pelas equipes de
trabalho. Sinaliza áreas de risco, controlando o fluxo de pessoas.
Providencia meios para controle de incêndios. Presta primeiros socorros
e providencia atendimento médico de urgência, quando necessário.

\begin{Form}
\textbf{Esta correspondência está adequada?} \\ 
\ChoiceMenu[radio,radiosymbol=\ding{52},bordercolor=0 0.5 0,name=cbovalid_ 321205 ]{}{CBO deve ficar=sim, \hspace{6em} CBO não fica aqui=nao}
\end{Form}

\textbf{CBO: 321210  ( Técnico florestal )}

Planeja atividades florestais, dimensionando quadro de pessoal e
prevendo instalações e equipamentos necessários para os trabalhos
programados. Propõe projetos, elaborando cronograma e planilha de custos
para execução. Realiza inventário florestal, auxiliando na definição da
metodologia a ser usada e coordenando coleta, conferência e
processamento de dados a respeito da cobertura vegetal e sobre
características qualitativas e quantitativas de espécies distribuídas na
floresta. Executa levantamentos fitossociológicos. Fiscaliza fauna e
flora, visitando e identificando áreas de intervenções ambientais.
Realiza perícias técnicas para avaliar a extensão de danos ambientais e
localizar os responsáveis pela sua ocorrência. Coíbe tráfico de animais
e extração de espécies vegetais ameaçadas de extinção. Fiscaliza caça e
pesca predatórias. Administra unidades de conservação e de produção,
estabelecendo diretrizes de conservação; coordenando equipe de
vigilância e resgate; e criando placas identificadoras e sinalizadoras.
Efetua levantamentos topográficos e outros estudos, utilizando sistemas
de georreferenciamento, como GPS-Sistema de Posicionamento Global
(Global Position System) ou GNSS-Sistema Global de Navegação por
Satélite (Global Navigation Satellite System). Supervisiona a construção
e a conservação da infraestrutura. Aplica práticas sustentáveis no
manejo de conservação do solo e da água. Realiza a identificação e a
conservação de sementes florestais. Coordena o plantio de espécies
ameaçadas de extinção, orientando a produção de mudas, acompanhando o
crescimento pós-plantio e aplicando métodos de erradicação e controle de
vetores e pragas, doenças e plantas daninhas. Supervisiona manejo de
florestas nativas e comerciais, acompanhando colheita florestal;
monitorando beneficiamento de produtos florestais; e controlando volume
de produção. Executa práticas de extensão florestal, mobilizando
comunidades e lideranças. Orienta o uso de tecnologias. Monta unidades
demonstrativas. Assessora a criação de cooperativas. Atua na preservação
e na conservação ambiental, prestando orientações sobre recuperação de
ecossistemas degradados, sobre legislação ambiental, sobre execução de
projetos técnicos, e sobre utilização racional de recursos naturais
renováveis. Ministra programas de treinamento e profere palestras
referentes à proteção e à manutenção florestal. Participa de pesquisas
florestais, elaborando projetos; selecionando matrizes de espécies
florestais; instalando e acompanhando experimentos; coletando dados e
materiais; e participando de publicação de resultados. Elabora
documentação técnica. Emite laudos técnicos e periciais; autorização de
colheita florestal e queima controlada; relatórios de inventários
florestais; planos de autossuprimento; autos de infração ambiental; e
relatórios de atividades operacionais. Atualiza cadastros. Providencia
renovação de documentos fiscais. Fornece dados técnicos para elaboração
de contratos. Coordena equipes de trabalho, participando de processo de
seleção e de administração de pessoal. Distribui tarefas e programa
férias das equipes. Controla conflitos nas equipes. Supervisiona equipes
de trabalho, avaliando desempenho e levantando as necessidades de
treinamento. Viabiliza a realização de programas de aperfeiçoamento e
atualização profissional. Pode gerenciar propriedades de produção
florestal. Monitora a limpeza e a organização dos locais de trabalho.
Zela pela manutenção de máquinas e equipamentos, providenciando
manutenção de primeiro nível e solicitando serviços de manutenção
corretiva. Trabalha com segurança, fazendo uso dos equipamentos de
proteção individual. Fiscaliza a utilização de equipamentos de proteção
individual e coletiva pelas equipes de trabalho. Sinaliza áreas de
risco, controlando o fluxo de pessoas. Providencia meios para controle
de incêndios. Presta primeiros socorros e providencia atendimento médico
de urgência, quando necessário.

\begin{Form}
\textbf{Esta correspondência está adequada?} \\ 
\ChoiceMenu[radio,radiosymbol=\ding{52},bordercolor=0 0.5 0,name=cbovalid_ 321210 ]{}{CBO deve ficar=sim, \hspace{6em} CBO não fica aqui=nao}
\end{Form}

\newpage
\section{ 11009 -- Técnico em Fruticultura }

\textbf{Perfil do Curso (CNCT)}

O Técnico em Fruticultura será habilitado para:

Planejar, organizar, dirigir e controlar processos de implantação,
condução do sistema produtivo de plantas frutíferas, de forma
sustentável, aplicando as Boas Práticas de Produção Agrícola (BPA).
Supervisionar a colheita e a pós-colheita de frutas, bem como executar
etapas do processo produtivo, desde a produção de sementes e mudas, a
pós-colheita, o controle de qualidade dos processos de produção de
frutas até os serviços de manutenção de instalações. Prestar assistência
técnica e assessoria ao estudo e ao desenvolvimento de projetos e
pesquisas tecnológicas, ou aos trabalhos de vistoria e consultoria na
cadeia produtiva de frutas. Elaborar orçamentos, laudos, pareceres,
relatórios e projetos, inclusive de incorporação de novas tecnologias.
Aplicar métodos e programas de melhoramento genético. Prestar
assistência técnica à aplicação, à comercialização, ao manejo de
produtos especializados, à recomendação e à interpretação de análise de
solos, à aplicação de fertilizantes e corretivos nos tratos das
culturas. Identificar os processos simbióticos de absorção, de
translocação e os efeitos alelopáticos entre solo e planta, planejando
ações referentes aos tratos das culturas. Realizar o preparo mecânico ou
manual do solo e a implantação das culturas. Aplicar práticas
sustentáveis no manejo de conservação do solo e da água. Recomendar e
aplicar os processos de adubação, correção e conservação do solo.
Aplicar as legislações pertinentes ao processo produtivo e ao meio
ambiente. Executar levantamentos topográficos e projetos de
geoprocessamento. Conduzir o preparo da área e definir o sistema de
cultivo e manejo. Dimensionar o sistema, selecionar e manejar projetos
de irrigação. Efetuar tratos culturais em todos os estágios de
desenvolvimento de espécies frutíferas. Identificar e promover o manejo
integrado de pragas, doenças e plantas espontâneas. Dominar os
procedimentos para coleta de amostras de solos, folhas e frutos,
analisar e interpretar os resultados laboratoriais. Planejar a
manutenção de máquinas e equipamentos agrícolas e operá-los. Orientar as
ações necessárias à colheita, à pós-colheita, ao armazenamento, à
conservação, ao processamento, ao beneficiamento e à comercialização de
frutas. Elaborar, aplicar e monitorar programas profiláticos, higiênicos
e sanitários na produção vegetal e agroindustrial. Emitir laudos e
documentos de classificação e exercer a fiscalização e a gerência de
sistemas de controle de qualidade na produção agrícola de frutas.
Responsabilizar-se pela implantação de pomares, acompanhando seu
desenvolvimento até a fase produtiva, emitindo os respectivos
certificados de origem e qualidade de produtos. Treinar e conduzir
equipes nas suas modalidades de atuação profissional. Administrar e
executar planejamento e gestão financeiros da propriedade agrícola.
Orientar na gestão de riscos relacionada à comercialização de frutas.
Prestar assistência técnica às áreas de crédito rural e agroindustrial,
de paisagismo, de construção de benfeitorias rurais, de drenagem e
irrigação. Supervisionar sistemas de certificação e rastreabilidade de
frutas. Realizar práticas convencionais e técnicas de produção orgânica
em cultivo de frutas. Operar veículos aéreos remotamente pilotados e
equipamentos de precisão para monitoramento remoto da produção agrícola.

Para a atuação como Técnico em Fruticultura, são fundamentais:

Conhecimentos e saberes relacionados à produção de espécies frutíferas.
Atualização em relação às inovações tecnológicas. Cooperação de forma
construtiva e colaborativa nos trabalhos em equipe, tomada de decisões.
Adoção de senso investigativo, visão sistêmica das atividades e
processos, capacidade de comunicação e argumentação, autonomia,
proatividade, liderança, respeito às diversidades nos grupos de
trabalho, resiliência frente aos problemas, organização,
responsabilidade, visão crítica, humanística, ética e consciência em
relação ao impacto de sua atuação profissional na sociedade e no
ambiente.

\textbf{CBOs correspondidos e seus perfis (presentes no CNCT, e presentes na predição da IA)}

\textbf{CBO: 321105  ( Técnico agrícola )}

Planeja projeto agrícola, pesquisando mercado consumidor e analisando
viabilidade econômica. Avalia condições da propriedade, verificando sua
infraestrutura - máquinas, equipamentos e instalações -; levantando
dados sobre a área a ser trabalhada -- topografia e extensão --;
examinando as condições edafoclimáticas - solo, clima e água -;
verificando a disponibilidade de mão de obra para as atividades;
levantando a capacitação tecnológica do produtor; e pesquisando o
mercado fornecedor de insumos, materiais, máquinas e equipamentos.
Executa levantamento do custo-benefício para o produtor. Verifica as
condições para implantação de projeto agrícola, coletando amostras para
análise de solo e interpretando os resultados laboratoriais; verificando
a disponibilidade e a qualidade da água a ser utilizada; e coletando
dados meteorológicos. Executa projeto agrícola, definindo cultivares e
rotação de culturas. Elabora planta de construções rurais. Compra
máquinas, equipamentos, insumos e materiais. Supervisiona a colheita e a
pós-colheita das principais culturas, aplicando técnicas de adubação
verde, de cobertura morta e de correção do solo, conforme cada caso.
Orienta os processos de silagem e transporte. Supervisiona trabalhadores
agrícolas, contratando mão de obra, distribuindo tarefas, orientando e
acompanhando a execução de atividades. Fiscaliza produção agrícola,
inspecionando documentação de produtos em trânsito, examinando a
sanidade de produtos agrícolas, enviando amostras de produtos para
análises laboratoriais, fiscalizando aplicação de agrotóxicos e
inspecionando cumprimento de normas e padrões técnicos. Conduz
experimentos de pesquisa, elaborando relatórios, laudos, pareceres,
perícias e avaliações. Administra propriedades rurais, adotando sistema
de produção conforme necessidade do mercado e controlando despesas e
receitas. Administra funcionários da propriedade. Comercializa produção
agrícola. Presta assistência e consultoria técnicas, orientando na
escolha do local para atividade agrícola; aconselhando sobre escolha de
espécies e cultivares; instruindo sobre preparo, correção e conservação
de solo; informando sobre equipamentos, máquinas e implementos;
indicando formas e manejo de irrigação e drenagem; aconselhando sobre
manejo integrado de pragas e doenças; orientando sobre época de plantio,
tratos culturais e colheita; e estimulando o beneficiamento de produtos
agrícolas. Orienta sobre preservação ambiental. Promove extensão e
capacitação rural, organizando reuniões com produtores e orientando
formação de associações e grupos de produtores. Assessora produtores,
orientando-os em cultivo com técnicas convencionais ou alternativas de
plantio, em manejo do solo e em compra e venda de insumos, materiais e
produtos agrícolas. Promove a difusão de tecnologia, orientando a
implantação de agricultura de precisão e o uso de sistemas
automatizados, bem como o desenvolvimento de novos cultivares. Dissemina
o uso de ambientes protegidos de cultivo e o melhoramento de sementes.
Divulga o uso de controle remoto da produção e da qualidade de processos
e produtos. Ministra cursos e programas de treinamentos. Apresenta
resultados de pesquisas em encontros e congressos da área agrícola. Pode
participar do desenvolvimento de tecnologias agrícolas, adequando
equipamentos e instalações conforme necessidade da região e do produtor.
Cria e adapta técnicas convencionais e alternativas de cultivo. Orienta
produtores agrícolas sobre biossegurança, indicando uso racional de
agrotóxicos e esclarecendo sobre destino de embalagens de agrotóxicos.
Informa sobre limpeza e desinfeção de máquinas, equipamentos e
instalações.

\begin{Form}
\textbf{Esta correspondência está adequada?} \\ 
\ChoiceMenu[radio,radiosymbol=\ding{52},bordercolor=0 0.5 0,name=cbovalid_ 321105 ]{}{CBO deve ficar=sim, \hspace{6em} CBO não fica aqui=nao}
\end{Form}

\textbf{CBO: 325205  ( Técnico de alimentos )}

Planeja as rotinas de trabalho, interpretando as ordens de serviço;
elaborando o cronograma de produção; especificando e calculando os
materiais e os insumos; selecionando os procedimentos para cada
atividade; e verificando a disponibilidade dos equipamentos, materiais,
insumos e gêneros alimentícios. Supervisiona os processos de produção,
recebendo insumos, produtos e gêneros alimentícios; acompanhando
pré-preparo e preparo de alimentos; monitorando processos -- tais como
de trituração, pasteurização, mistura, fermentação e cocção -; e
assegurando as condições operacionais do processo produtivo. Corrige
desvios nos processos manuais e automatizados. Controla o tempo de
produção. Controla o peso e a dimensão do produto. Verifica as condições
da embalagem do produto final. Acompanha a distribuição de produtos até
os pontos de venda. Avalia a qualidade de produtos acabados, coletando
amostras dos produtos e realizando análises microbiológicas, sensoriais
e físico-químicas de matérias-primas, insumos e produtos intermediários
e acabados. Realiza testes de desempenho de matérias-primas, insumos e
produtos intermediários e acabados. Controla a qualidade nas etapas de
produção, identificando e corrigindo pontos críticos de controle e
assegurando as condições higiênico-sanitárias. Verifica as condições de
armazenamento, controlando a data de vencimento dos produtos e
acompanhando o controle integrado de pragas e vetores. Participa, sob
supervisão, de pesquisas e de projetos para melhoria, adequação e
desenvolvimento de produtos, acompanhando as necessidades do mercado e
participando na elaboração da formulação e na testagem dos novos
produtos. Assessora a implementação das mudanças aprovadas. Participa na
elaboração da informação nutricional, que consta na rotulagem do
produto. Orienta clientes internos e externos sobre aspectos
nutricionais dos produtos. Coordena e supervisiona equipes de trabalho,
participando na seleção de pessoal, orientando as atividades,
remanejando pessoal, avaliando os resultados de desempenho e realizando
treinamento de rotinas operacionais. Controla as condições do ambiente
de trabalho, monitorando a temperatura e a umidade do ar. Acompanha os
procedimentos de sanitização das instalações. Verifica condições de
segurança ambiental e controla os pontos críticos e o uso de
equipamentos de proteção individual. Monitora o funcionamento de
máquinas, equipamentos e instrumentos, solicitando serviço de
manutenção, quando necessário. Controla desperdícios e destina
adequadamente os resíduos gerados, fazendo reciclagem e realizando
descarte de acordo com as normas ambientais.

\begin{Form}
\textbf{Esta correspondência está adequada?} \\ 
\ChoiceMenu[radio,radiosymbol=\ding{52},bordercolor=0 0.5 0,name=cbovalid_ 325205 ]{}{CBO deve ficar=sim, \hspace{6em} CBO não fica aqui=nao}
\end{Form}

\newpage
\section{ 11010 -- Técnico em Geologia }

\textbf{Perfil do Curso (CNCT)}

O Técnico em Geologia será habilitado para:

Realizar atividades de prospecção, avaliação técnica e econômica,
planejamento, extração e produção referentes aos recursos naturais.
Prestar assistência técnica e assessoria ao estudo e ao desenvolvimento
de projetos e pesquisas tecnológicas, ou aos trabalhos de vistoria,
perícia, arbitramento e consultoria. Elaborar orçamentos, laudos,
pareceres, relatórios e projetos, inclusive de incorporação de novas
tecnologias. Realizar levantamento, coleta, processamento e análise de
dados através do Sistema de Navegação Global por Satélite (GNSS).
Levantar e processar dados obtidos por meio de sensores orbitais,
radares imageadores e aeronaves remotamente pilotadas. Elaborar e
gerenciar dados em Sistema de Informações Geográficas (SIG). Executar
mapeamento geológico e amostragem em superfície e subsolo. Identificar,
qualificar e quantificar ocorrências minerais. Realizar a caracterização
de minérios. Realizar levantamentos topográficos nas atividades de
pesquisa mineral. Operar equipamentos de sondagem, perfuração e pesquisa
mineral.

Para a atuação como Técnico em Geologia, são fundamentais:

Conhecimentos e saberes relacionados aos recursos minerais e à proteção
ambiental. Atualização em relação às inovações tecnológicas. Cooperação
de forma construtiva e colaborativa nos trabalhos em equipe, tomada de
decisões. Adoção de senso investigativo, visão sistêmica das atividades
e processos, capacidade de comunicação e argumentação, autonomia,
proatividade, liderança, respeito às diversidades nos grupos de
trabalho, resiliência frente aos problemas, organização,
responsabilidade, visão crítica, humanística, ética e consciência em
relação ao impacto de sua atuação profissional na sociedade e no
ambiente.

\textbf{CBOs correspondidos e seus perfis (presentes no CNCT, e presentes na predição da IA)}

\textbf{CBO: 316105  ( Técnico em geofísica )}

Realiza estudos geofísicos, analisando a composição e as propriedades
físicas e estruturais do solo e das rochas e investigando fenômenos
magnéticos, elétricos, sísmicos, gravitacionais e térmicos. Localiza
pontos de amostragem para registrar dados geofísicos e dados de
perfilagem de poços e desvios de furos. Processa os dados geofísicos,
fazendo sua compilação, sua conferência, sua avaliação e sua
interpretação. Elabora e atualiza mapas, perfis geográficos e relatórios
técnicos, a fim de contribuir para a execução dos programas de estudos
geofísicos aplicados na verificação da estabilidade dos solos; na
viabilidade de execução de projetos e fundações em terrenos; na análise
dos impactos ambientais da exploração de minérios; e em outras
aplicações similares. Pode utilizar GPR-Radar de penetração no solo
terrestre (Ground Penetration Radar) portátil e métodos geofísicos -
tais como sísmico de reflexão e refração, gravimetria, magnetometria e
radiometria -, para a prospecção do subsolo. Pode participar da
interpretação de dados sísmicos - com o auxílio de computadores e
instrumentos de medição -, construindo modelos tridimensionais
(realidade virtual) para investigação de sistemas subterrâneos e para
simulações de terremotos e de outros eventos que possam ocorrer nesses
sistemas. Participa no planejamento da prospecção de recursos minerais,
auxiliando na elaboração de cronograma de trabalho e na definição de
equipe e equipamentos para a execução das atividades. Fornece suporte
para a prospecção de recursos minerais, auxiliando em processos de
requerimentos e alvarás de áreas de pesquisa de lavra; definindo base
operacional e apoio logístico; e estabelecendo métodos e padrões de
trabalho. Identifica áreas de potencial mineral, verificando vias de
acesso. Consulta mapas aerogeofísicos, relatórios e cartas topográficas,
a fim de subsidiar processos de tomada de decisão. Contata proprietários
superficiários para negociação. Pode supervisionar equipes de trabalho,
distribuindo tarefas, avaliando desempenho e desenvolvendo treinamentos.
Participa do controle de recursos, administrando custos operacionais.
Conserva o local de trabalho limpo, seguro e organizado. Mantém
ferramentas, instrumentos e equipamentos limpos, organizados,
acondicionados e em plenas condições de uso. Contribui para a redução de
impactos ambientais, atuando para o direcionamento de sondagens aos
locais de maior interesse. Sensibiliza e conscientiza as equipes de
trabalho quanto à importância da preservação ambiental. Trabalha com
segurança, usando equipamentos de proteção individual e participando de
ações para prevenção de acidentes. Identifica situações de risco e
propõe soluções para eliminá-las. Pode prestar primeiros socorros.

\begin{Form}
\textbf{Esta correspondência está adequada?} \\ 
\ChoiceMenu[radio,radiosymbol=\ding{52},bordercolor=0 0.5 0,name=cbovalid_ 316105 ]{}{CBO deve ficar=sim, \hspace{6em} CBO não fica aqui=nao}
\end{Form}

\textbf{CBO: 316110  ( Técnico em geologia )}

Programa e executa estudos geológicos, para o conhecimento da
localização e da natureza de formações geológicas. Coleta amostras
geológicas, tais como amostras de sondagem de percussão tracionária; de
sedimentos de corrente; de concentrado de bateia; de rochas; de solos;
de trado; de testemunhos de sondagem; de canal; de plantas e fósseis; e
de água, gases e óleo. Descreve aspectos macroscópicos das amostras.
Organiza e protocola as amostras de acordo com normas legais para
análises laboratoriais. Prepara as amostras minerais para análise,
homogeneizando-as, quarteando-as e secando-as. Calcula as densidades de
amostras. Confecciona lâminas de amostras. Trata amostras, efetuando sua
separação por granulometria; fazendo sua britagem e sua pulverização; e
realizando sua concentração volumetricamente. Monitora os processos de
análise laboratorial. Processa os dados geológicos, fazendo sua
compilação, sua conferência, sua avaliação e sua interpretação. Elabora
e atualiza mapas, perfis geográficos e relatórios técnicos, a fim de
contribuir para a execução dos programas de estudos geológicos. Realiza
mapeamentos geológicos e levantamentos topográficos. Pode interpretar
dados sísmicos, com auxílio de computadores e instrumentos de medição.
Pode utilizar GPR-Radar de penetração no solo terrestre (Ground
Penetration Radar) portátil. Participa no planejamento da prospecção de
recursos minerais, auxiliando na elaboração de cronograma de trabalho e
na definição de equipe e equipamentos para a execução das atividades.
Presta apoio técnico à prospecção de recursos minerais, definindo base
operacional e apoio logístico e estabelecendo métodos e padrões de
trabalho. Identifica áreas de potencial mineral, efetuando levantamento
de publicações técnicas; examinando áreas de pesquisa e vias de acesso;
realizando reconhecimento cartográfico; e fazendo reconhecimento
geológico. Interpreta mapas aerogeofísicos, relatórios e cartas
topográficas; analisa imagens de alta resolução -- capturadas por meio
de satélites e drones - para a investigação e o delineamento de áreas de
interesse; e realiza georreferenciamento, a fim de subsidiar processos
de tomada de decisão. Pode contatar proprietários superficiários para
negociação. Controla a qualidade de frente de lavra, determinando o teor
de minério conforme amostragens. Instala piezômetros. Seleciona e
delimita áreas para a extração de recursos minerais. Fiscaliza o avanço
de lavra. Monitora qualidade de pilhas de estoques e homogeneização.
Gera informações para cálculo de vida útil da mina; para possibilitar a
correta execução da programação estabelecida; e para prevenir acidentes
e eventuais danos às zonas de extração. Participa em estudos de impacto
ambiental, detectando emissão de poluentes no solo, na atmosfera e em
bacias hidrográficas; índices de desmatamento; e índices de erosão, a
fim de conhecer os problemas e minimizar seus efeitos para garantir a
preservação dos ambientes naturais. Participa na elaboração de planos de
recuperação de áreas degradadas. Registra a ocorrência de fauna e flora
específicas. Opera ou ajusta equipamentos ou aparelhos usados para obter
dados geológicos. Pode supervisionar equipes de trabalho, distribuindo
tarefas, avaliando desempenho e desenvolvendo treinamentos. Planeja e
direciona as atividades dos trabalhadores que operam equipamentos para
coletar dados geológicos. Participa do controle de recursos,
administrando custos operacionais. Pode elaborar orçamentos. Conserva o
local de trabalho limpo, seguro e organizado. Mantém ferramentas,
instrumentos e equipamentos limpos, organizados, acondicionados e em
plenas condições de uso. Trabalha com segurança, usando equipamentos de
proteção individual e participando de ações para prevenção de acidentes.
Identifica situações de risco e propõe soluções para eliminá-las. Pode
prestar primeiros socorros.

\begin{Form}
\textbf{Esta correspondência está adequada?} \\ 
\ChoiceMenu[radio,radiosymbol=\ding{52},bordercolor=0 0.5 0,name=cbovalid_ 316110 ]{}{CBO deve ficar=sim, \hspace{6em} CBO não fica aqui=nao}
\end{Form}

\textbf{CBO: 316115  ( Técnico em geoquímica )}

Programa e executa estudos geoquímicos, para delinear os padrões de
distribuição dos elementos químicos em locais preestabelecidos. Realiza
planejamento da amostragem, definindo tipo e número de amostras a serem
coletadas a partir dos objetivos do estudo. Estabelece padrões de
amostras. Pode participar de pesquisas geoquímicas regionais, realizando
viagens de campo para levantamento de dados. Localiza pontos de
amostragem e coleta amostras, tais como sedimentos de corrente; de
concentrado de bateia; de rochas; de solos; de trado; de canal; de
plantas e fósseis; e de água, gases e óleo. Pode utilizar instrumentos
de medição portáteis - eletrônicos, sônicos e nucleares - para a
amostragem em campo. Organiza e protocola as amostras de acordo com
normas legais para análises laboratoriais. Prepara as amostras minerais
para análise, homogeneizando-as, quarteando-as e secando-as. Prepara as
amostras para análise por via úmida. Calcula as densidades de amostras.
Confecciona lâminas de amostras. Trata amostras, efetuando sua separação
por granulometria; fazendo sua britagem e sua pulverização; e realizando
sua concentração volumetricamente. Monitora os processos de análise
laboratorial. Processa os dados geoquímicos, fazendo sua compilação, sua
conferência, sua avaliação e sua interpretação. Realiza tratamento
estatístico básico, tais como elaboração de histogramas e cálculos de
mediana, de valor médio, de valores mínimos e máximos e de desvio
padrão. A partir dos parâmetros estatísticos, desenha as curvas ou
superfícies de tendência para cada elemento químico analisado. Pode
associar tais resultados aos valores constantes de tabelas legais, que
-- quando disponíveis - determinam os patamares de qualidade dos
elementos (como água, solo, entre outros). Elabora mapas geoquímicos,
utilizando software para mostrar os resultados de cada elemento químico
em cada local, formando a apresentação espacial dos dados. Participa em
estudos de avaliação das composições químicas do substrato rochoso, dos
solos e das águas de superfície e de abastecimento público, que visam
disponibilizar dados e informações para pesquisas de novos depósitos
minerais, de fertilidade natural para a agricultura, de fontes de
contaminações de elementos químicos nocivos à saúde humana, à saúde
animal e ao meio ambiente, entre outras. Participa na elaboração de
relatórios técnicos relacionados às análises realizadas. Pode utilizar
sistemas automatizados para promover a padronização, precisão e rapidez
na coleta e no uso de dados geoquímicos. Participa no planejamento da
prospecção de recursos minerais, auxiliando na elaboração de cronograma
de trabalho e na definição de equipe e equipamentos para a execução das
atividades. Presta apoio técnico à prospecção de recursos minerais,
definindo base operacional e apoio logístico e estabelecendo métodos e
padrões de trabalho. Identifica áreas de potencial mineral e verifica
vias de acesso, a fim de subsidiar processos de tomada de decisão. Pode
contatar proprietários superficiários para negociação. Participa em
estudos de impacto ambiental, detectando emissão de poluentes no solo,
na atmosfera e em bacias hidrográficas e verificando índices de
desmatamento, a fim de conhecer problemas e minimizar seus efeitos para
garantir a preservação dos ambientes naturais. Pode supervisionar
equipes de trabalho, distribuindo tarefas, avaliando desempenho e
desenvolvendo treinamentos. Participa do controle de recursos,
administrando custos operacionais. Conserva o local de trabalho limpo,
seguro e organizado. Mantém ferramentas, instrumentos e equipamentos
limpos, organizados, acondicionados e em plenas condições de uso.
Trabalha com segurança, usando equipamentos de proteção individual e
participando de ações para prevenção de acidentes. Identifica situações
de risco e propõe soluções para eliminá-las. Pode prestar primeiros
socorros.

\begin{Form}
\textbf{Esta correspondência está adequada?} \\ 
\ChoiceMenu[radio,radiosymbol=\ding{52},bordercolor=0 0.5 0,name=cbovalid_ 316115 ]{}{CBO deve ficar=sim, \hspace{6em} CBO não fica aqui=nao}
\end{Form}

\textbf{CBO: 316120  ( Técnico em geotecnia )}

Participa no planejamento de projetos geotécnicos, relacionados a
sondagens, fundações, poços, taludes, escavações, estruturas de
contenções, entre outros. Presta apoio na elaboração de cronograma de
trabalho e na definição de equipe e equipamentos para a execução das
atividades. Localiza pontos de amostragem e estabelece padrões de
amostras. Coleta amostras, tais como de rochas; de solo; de trado; de
sondagem de percussão tracionária; de testemunhos de sondagem; e
indeformada. Prepara as amostras, organizando-as e protocolando-as de
acordo com normas legais para as análises laboratoriais. Processa os
dados geológicos-geotécnicos, fazendo sua compilação, sua conferência,
sua avaliação e sua interpretação. Elabora e atualiza mapas e relatórios
técnicos, contribuindo na execução dos programas de prevenção de
desabamentos, desmoronamentos e deslizamentos de encostas; para
preservação de lençóis freáticos; e para outros estudos geotécnicos
realizados em indústrias extrativas e de construção. Aplica
conhecimentos especializados de mecânica dos solos e mecânica das
rochas. Utiliza instrumentos, tais como piezômetros - que medem a
pressão da água em solos e rochas --, inclinômetros e tensiômetros. Na
execução das atividades, pode utilizar CAD-Desenho Assistido por
Computador (Computer Aided Design); GIS-Sistema de Informação Geográfica
(Geographic Information System); BIM-Modelagem da Informação da
Construção (Building Information Modeling), entre outros ``softwares'' e
sistemas. Pode interpretar dados sísmicos com o auxílio de computadores
e de instrumentos de medição. Pode atuar na área de barragens,
acompanhando engenheiros (que atuam na área de geotecnia) e geólogos em
inspeções rotineiras de campo. Utiliza dados coletados de forma remota
por equipamentos -- como veículo aéreo e embarcação não tripulados -,
durante mapeamentos e inspeções. Faz monitoramento da instrumentação das
barragens - em centro de monitoramento geotécnico -, recebendo alertas
de anormalidades; executando análise preliminar para verificar se está
ocorrendo falha dos instrumentos ou de sistemas; e acionando engenheiros
e geólogos para avaliar anomalias nas estruturas. Coleta dados para
emissão de relatórios de inspeção e de monitoramento. Auxilia na
elaboração de relatórios técnicos e na inserção de dados das inspeções
no SIGBM-Sistema Integrado de Gestão de Barragens de Mineração.
Participa na organização e na realização de simulados de emergência de
barragem. Pode participar em estudos para escolha de novos materiais de
reforço - em terrenos, aterros, barreiras, contenções, entre outros - e
em definição de procedimentos para utilização de novos materiais
geossintéticos, como mantas; grelhas e geogrelhas drenantes;
EPS-Poliestireno Expandido (Expanded Polystyrene) associado a
geomembranas; entre outros. Participa em estudos de impacto ambiental,
detectando índices de erosão e identificando outros riscos, a fim de
conhecer problemas e minimizar seus efeitos em obras de barragens, na
escavação de túneis, na compactação de aterros e em outras atividades.
Pode supervisionar equipes de trabalho, distribuindo tarefas, avaliando
desempenho e desenvolvendo treinamentos. Sensibiliza e conscientiza as
equipes de trabalho quanto à importância da preservação ambiental.
Participa do controle de recursos, administrando custos operacionais.
Conserva o local de trabalho limpo, seguro e organizado. Mantém
ferramentas, instrumentos e equipamentos limpos, organizados,
acondicionados e em plenas condições de uso. Trabalha com segurança,
usando equipamentos de proteção individual e participando de ações para
prevenção de acidentes. Identifica situações de risco e propõe soluções
para eliminá-las. Pode prestar primeiros socorros.

\begin{Form}
\textbf{Esta correspondência está adequada?} \\ 
\ChoiceMenu[radio,radiosymbol=\ding{52},bordercolor=0 0.5 0,name=cbovalid_ 316120 ]{}{CBO deve ficar=sim, \hspace{6em} CBO não fica aqui=nao}
\end{Form}

\textbf{CBO: 316320  ( Técnico em pesquisa mineral )}

Participa no planejamento de atividades de prospecção e pesquisa
mineral, programando etapas de trabalho. Auxilia na seleção de equipe;
na escolha de equipamentos, ferramentas e instrumentos; e na definição
de processos de trabalho. Contribui na elaboração de plano de custos
operacionais. Assessora na seleção de áreas para prospecção e pesquisa
de minerais, interpretando fotografias aéreas e analisando imagens de
alta resolução -- capturadas por meio de satélites e drones - para a
investigação e o delineamento de áreas de interesse; efetuando
levantamento topográfico de vias de acesso; fotodocumentando rochas; e
fazendo mapeamento geológico, para subsidiar tomada de decisão.
Participa na elaboração de requerimentos de alvará de pesquisa e lavra.
Notifica superficiários de uso de terreno para prospecção e pesquisa.
Assessora na seleção de métodos de prospecção e de pesquisa a serem
utilizados. Coordena a abertura de vias de acesso para áreas de
pesquisa. Participa no estudo da área de prospecção e pesquisa de
minerais, levantando perfis geológicos e radioativos; medindo direção e
mergulho de rochas aflorantes; e triando microfósseis. Participa na
execução de prospecção e pesquisa de minerais, marcando malhas para
perfuração; locando poços artesianos e furos de sondagem; e
supervisionando sondagem e perfilagem, para confirmar a existência de
depósito de mineral. Coordena a execução de drenagem da jazida.
Determina porosidade e permeabilidade em plugues de testemunhos.
Monitora teor de minério. Coordena a coleta amostras, tais como de
sedimentos de corrente; de solo, rocha e água; de furos de sondagem; e
de poços de pesquisa, trincheiras e galerias. Prepara as amostras,
homogeneizando-as, quarteando-as e confeccionando lâminas delgadas.
Descreve características físico-químicas e litológicas de amostras.
Identifica e acondiciona amostras para testes e análises. Processa e
analisa dados levantados durante prospecção e pesquisa mineral,
confeccionando mapas e seções geológicas de prospecção, pesquisa e
lavra; cubando reserva lavrável e corpo mineralizado; e calculando
produção de poços, galerias, trincheiras. Auxilia na interpretação de
corpo mineralizado. Cria modelo de bloco baseado em modelo geológico.
Participa na apresentação dos resultados dos estudos, estabelecendo a
quantidade de minério a ser extraída; a proposição da tecnologia a ser
aplicada na exploração; e a estimativa de vida útil da mina. Pode
utilizar modelamento geológico 3D - representação em softwares das
rochas e das suas estruturas, feita a partir de dados geológicos
coletados -, para verificar localização e quantidade de minério;
natureza do ambiente onde se encontra o mineral; entre outros
resultados. Pode usar instrumentos de medição eletrônicos, sônicos e
nucleares em laboratório e em campo, durante a realização dos estudos.
Gera e edita banco de dados de perfis elétricos, acústicos, radioativos
e geológicos. Controla custos e materiais utilizados durante as
atividades de prospecção e pesquisa. Participa na elaboração de padrões
técnicos, procedimentos e instruções operacionais de trabalho.
Supervisiona a operação e a manutenção de equipamentos. Pode
supervisionar equipes de trabalho, distribuindo tarefas, avaliando
desempenho e desenvolvendo treinamentos. Monitora a organização do
ambiente de trabalho e a conservação de ferramentas e instrumentos
limpos, acondicionados e em plenas condições de uso e funcionamento.
Controla a aplicação de procedimentos da empresa para cumprimento de
normas ambientais. Zela pela saúde e segurança da equipe de trabalho,
monitorando o uso de equipamentos de proteção individual e a realização
de exames periódicos de saúde. Identifica áreas de risco, tomando
medidas para seu isolamento e para análise dos fatores de risco.
Registra ocorrências de acidentes de trabalho.

\begin{Form}
\textbf{Esta correspondência está adequada?} \\ 
\ChoiceMenu[radio,radiosymbol=\ding{52},bordercolor=0 0.5 0,name=cbovalid_ 316320 ]{}{CBO deve ficar=sim, \hspace{6em} CBO não fica aqui=nao}
\end{Form}

\newpage
\section{ 11011 -- Técnico em Mineração }

\textbf{Perfil do Curso (CNCT)}

O Técnico em Mineração será habilitado para:

Realizar atividades de prospecção, avaliação técnica e econômica,
planejamento, extração e produção referentes aos recursos naturais.
Prestar assistência técnica e assessoria ao estudo e ao desenvolvimento
de projetos e pesquisas tecnológicas, ou aos trabalhos de vistoria,
perícia, arbitramento e consultoria. Elaborar orçamentos, laudos,
pareceres, relatórios e projetos, inclusive de incorporação de novas
tecnologias. Realizar levantamento topográfico, sensoriamento remoto e
geoprocessamento, conforme sua formação profissional. Auxiliar na
caracterização de minérios sob os aspectos físico-químico, mineralógico
e granulométrico. Executar projetos de desmonte, transporte e
carregamento de minérios. Monitorar a estabilidade de rochas em minas
subterrâneas e a céu aberto. Auxiliar no mapeamento geológico e
amostragem em superfície e subsolo. Supervisionar, coordenar e operar
equipamentos de fragmentação, de separação mineral, separação
sólido/líquido, hidrometalúrgicos e de secagem. Supervisionar, coordenar
e operar equipamentos de extração mineral, sondagem, perfuração,
amostragem e transporte. Orientar e coordenar a execução de serviços de
manutenção de equipamentos. Prestar assistência técnica na compra, venda
e utilização de equipamentos especializados.

Para a atuação como Técnico em Mineração, são fundamentais:

Conhecimentos e saberes relacionados aos recursos minerais e à proteção
ambiental. Atualização em relação às inovações tecnológicas. Cooperação
de forma construtiva e colaborativa nos trabalhos em equipe e tomada de
decisões. Adoção de senso investigativo, visão sistêmica das atividades
e processos, capacidade de comunicação e argumentação, autonomia,
proatividade, liderança, respeito às diversidades nos grupos de
trabalho, resiliência frente aos problemas, organização,
responsabilidade, visão crítica, humanística, ética e consciência em
relação ao impacto de sua atuação profissional na sociedade e no
ambiente.

\textbf{CBOs correspondidos e seus perfis (presentes no CNCT, e presentes na predição da IA)}

\textbf{CBO: 316305  ( Técnico de mineração )}

Participa no planejamento de atividades de mineração, programando etapas
da produção de minério. Auxilia na seleção de equipe e de equipamentos;
na escolha de aplicativos de mineração e beneficiamento; e na definição
de demais recursos para execução das atividades. Contribui na elaboração
de plano de custos operacionais. Participa na elaboração de projetos
ambientais. Assessora na seleção de áreas para prospecção e pesquisa de
minerais, interpretando fotografias aéreas e fotodocumentando rochas;
realizando geoprocessamento; efetuando levantamento topográfico de vias
de acesso; e fazendo mapeamento geológico, para subsidiar tomada de
decisão. Participa na elaboração de requerimentos de alvará de pesquisa
e lavra. Presta assessoria ao estudo da área de prospecção e pesquisa de
minerais, levantando perfis geológicos e radioativos; determinando
porosidade e permeabilidade em plugues de testemunhos; e triando
microfósseis. Participa na execução de prospecção e pesquisa de
minerais, locando poços artesianos e furos de sondagem e supervisionando
sondagem e perfilagem. Presta assistência técnica à lavra e à exploração
de jazida, realizando levantamento topográfico de área de lavra e
deposição de estéril; coordenando abertura de vias de acesso; e
controlando produção de minério, estéril e hidrocarbonetos. Monitora o
teor de minério. Acompanha perfuração e desmonte de rocha e verifica
desempenho de explosivos. Participa na elaboração de planos de drenagem
e coordena a execução da drenagem da jazida. Monitora o nível de lençol
freático. Controla o nível de sedimentos de bacia de contenção.
Supervisiona o beneficiamento de minério, definindo variáveis de
controle e monitorando processos. Supervisiona alimentação de planta de
beneficiamento, bombeamento de água e processo de planta de reagente.
Supervisiona processo de cominuição de minério e processos de separação
- por classificação; por desaguamento; por filtragem; gravimétrica; e
magnética - de minério. Supervisiona processo de amostragem de rejeito e
concentrado. Inspeciona sistemas de rejeito. Orienta a dosagem
(blendagem) do produto mineral de acordo com especificações. Analisa a
qualidade do produto. Controla a movimentação da produção final de
minério, inspecionando pátios para evitar contaminações; solicitando
transporte; e controlando embarque e desembarque de minério de vias
férrea, rodoviária, marítima e mineroduto. Controla peso médio líquido
de composições de carregamento. Monitora tempos de carregamentos de
produto mineral. Coordena a coleta amostras, tais como de sedimentos de
corrente; de solo, rocha, hidrocarbonetos e água; de furos de sondagem;
de poços de pesquisa; de frentes de lavra; e de produtos estocados e
embarcados. Prepara as amostras, homogeneizando-as, quarteando-as e
confeccionando lâminas delgadas. Descreve características
físico-químicas e litológicas de amostras. Identifica e acondiciona
amostras para testes e análises. Processa dados de pesquisa mineral,
prospecção e lavra, analisando dados; confeccionando mapas; e gerando
banco de dados de perfis elétricos, acústicos, radioativos e geológicos.
Calcula produção de poços, galerias, trincheiras e sondagens. Auxilia na
estimativa de vida útil de mina e de reservatório de petróleo. Participa
na elaboração de padrões técnicos, procedimentos e instruções
operacionais de trabalho. Supervisiona a operação e a manutenção de
equipamentos. Pode supervisionar equipes de trabalho, distribuindo
tarefas, avaliando desempenho e desenvolvendo treinamentos. Controla a
organização do ambiente de trabalho e a conservação de ferramentas e
instrumentos. Monitora aplicação de procedimentos da empresa para
cumprimento de normas ambientais. Zela pela saúde e segurança no
trabalho, monitorando o uso de equipamentos de proteção individual e a
realização de exames periódicos de saúde. Identifica áreas de risco de
minas e poços de petróleo, tomando medidas para seu isolamento e para
análise dos fatores de risco.

\begin{Form}
\textbf{Esta correspondência está adequada?} \\ 
\ChoiceMenu[radio,radiosymbol=\ding{52},bordercolor=0 0.5 0,name=cbovalid_ 316305 ]{}{CBO deve ficar=sim, \hspace{6em} CBO não fica aqui=nao}
\end{Form}

\textbf{CBO: 316315  ( Técnico em processamento mineral (exceto petróleo) )}

Participa no planejamento de atividades de processamento de minérios e
de controle da movimentação da produção final. Auxilia na seleção de
equipe e de equipamentos; na escolha de aplicativos de beneficiamento; e
na definição de demais recursos para execução das atividades. Participa
na elaboração de programas de manutenção de equipamentos e instalações.
Contribui na elaboração de plano de custos operacionais. Participa na
elaboração de projetos ambientais. Auxilia na elaboração de projetos de
locação de conjunto de correias transportadoras de minérios - em
substituição aos caminhões -, visando, entre outros objetivos, à
diminuição de emissão de poluentes. Supervisiona o beneficiamento de
minério, definindo variáveis de controle; monitorando processos; e
controlando custos e materiais. Supervisiona alimentação de planta de
beneficiamento, controlando a produção de minério - que vai passar pelo
processo de beneficiamento -; monitorando o processo de bombeamento de
água; e acompanhando o processo de planta de reagente. Supervisiona
processos de cominuição de minérios e processos de separação - por
classificação; por desaguamento; por filtragem; gravimétrica; e
magnética - de minérios. Supervisiona processo de amostragem de rejeito
e concentrado. Inspeciona sistemas de rejeito. Monitora teor de minério.
Orienta a dosagem de teores de minério (blendagem) de acordo com
especificações. Analisa a qualidade do produto. Controla a movimentação
da produção final de minério, inspecionando pátios para evitar
contaminações; solicitando transporte; e controlando embarque de minério
- em vias férrea, rodoviária e marítima e em mineroduto - e seu
desembarque. Controla peso médio líquido de composições de carregamento
e desvio padrão de peso em vagões. Monitora tempos de carregamentos de
produto mineral. Verifica demanda de minério de outros departamentos da
empresa. Comercializa produto mineral. Coordena a coleta de amostras de
produtos estocados e embarcados. Realiza processamento e análise de
dados, calculando o rendimento de usina e gerando banco de dados de
perfis elétricos, acústicos, radioativos e geológicos. Participa na
elaboração de padrões técnicos, procedimentos e instruções operacionais
de trabalho. Supervisiona a operação e a manutenção de equipamentos.
Pode utilizar sistemas autônomos para a operação de equipamentos de
forma remota. Pode supervisionar equipe de trabalho, distribuindo
tarefas, avaliando desempenho e desenvolvendo treinamentos. Monitora a
organização do ambiente de trabalho e a conservação de ferramentas e
instrumentos limpos, acondicionados e em plenas condições de uso e
funcionamento. Controla a aplicação de procedimentos da empresa para
cumprimento de normas ambientais. Fiscaliza a emissão de poluentes e o
nível de assoreamento de rios. Representa departamentos em inspeção
ambiental. Zela pela saúde e segurança da equipe de trabalho,
monitorando o uso de equipamentos de proteção individual e a realização
de exames periódicos de saúde. Identifica áreas de risco, tomando
medidas para seu isolamento e para análise dos fatores de risco.
Registra ocorrências de acidentes de trabalho.

\begin{Form}
\textbf{Esta correspondência está adequada?} \\ 
\ChoiceMenu[radio,radiosymbol=\ding{52},bordercolor=0 0.5 0,name=cbovalid_ 316315 ]{}{CBO deve ficar=sim, \hspace{6em} CBO não fica aqui=nao}
\end{Form}

\textbf{CBO: 316320  ( Técnico em pesquisa mineral )}

Participa no planejamento de atividades de prospecção e pesquisa
mineral, programando etapas de trabalho. Auxilia na seleção de equipe;
na escolha de equipamentos, ferramentas e instrumentos; e na definição
de processos de trabalho. Contribui na elaboração de plano de custos
operacionais. Assessora na seleção de áreas para prospecção e pesquisa
de minerais, interpretando fotografias aéreas e analisando imagens de
alta resolução -- capturadas por meio de satélites e drones - para a
investigação e o delineamento de áreas de interesse; efetuando
levantamento topográfico de vias de acesso; fotodocumentando rochas; e
fazendo mapeamento geológico, para subsidiar tomada de decisão.
Participa na elaboração de requerimentos de alvará de pesquisa e lavra.
Notifica superficiários de uso de terreno para prospecção e pesquisa.
Assessora na seleção de métodos de prospecção e de pesquisa a serem
utilizados. Coordena a abertura de vias de acesso para áreas de
pesquisa. Participa no estudo da área de prospecção e pesquisa de
minerais, levantando perfis geológicos e radioativos; medindo direção e
mergulho de rochas aflorantes; e triando microfósseis. Participa na
execução de prospecção e pesquisa de minerais, marcando malhas para
perfuração; locando poços artesianos e furos de sondagem; e
supervisionando sondagem e perfilagem, para confirmar a existência de
depósito de mineral. Coordena a execução de drenagem da jazida.
Determina porosidade e permeabilidade em plugues de testemunhos.
Monitora teor de minério. Coordena a coleta amostras, tais como de
sedimentos de corrente; de solo, rocha e água; de furos de sondagem; e
de poços de pesquisa, trincheiras e galerias. Prepara as amostras,
homogeneizando-as, quarteando-as e confeccionando lâminas delgadas.
Descreve características físico-químicas e litológicas de amostras.
Identifica e acondiciona amostras para testes e análises. Processa e
analisa dados levantados durante prospecção e pesquisa mineral,
confeccionando mapas e seções geológicas de prospecção, pesquisa e
lavra; cubando reserva lavrável e corpo mineralizado; e calculando
produção de poços, galerias, trincheiras. Auxilia na interpretação de
corpo mineralizado. Cria modelo de bloco baseado em modelo geológico.
Participa na apresentação dos resultados dos estudos, estabelecendo a
quantidade de minério a ser extraída; a proposição da tecnologia a ser
aplicada na exploração; e a estimativa de vida útil da mina. Pode
utilizar modelamento geológico 3D - representação em softwares das
rochas e das suas estruturas, feita a partir de dados geológicos
coletados -, para verificar localização e quantidade de minério;
natureza do ambiente onde se encontra o mineral; entre outros
resultados. Pode usar instrumentos de medição eletrônicos, sônicos e
nucleares em laboratório e em campo, durante a realização dos estudos.
Gera e edita banco de dados de perfis elétricos, acústicos, radioativos
e geológicos. Controla custos e materiais utilizados durante as
atividades de prospecção e pesquisa. Participa na elaboração de padrões
técnicos, procedimentos e instruções operacionais de trabalho.
Supervisiona a operação e a manutenção de equipamentos. Pode
supervisionar equipes de trabalho, distribuindo tarefas, avaliando
desempenho e desenvolvendo treinamentos. Monitora a organização do
ambiente de trabalho e a conservação de ferramentas e instrumentos
limpos, acondicionados e em plenas condições de uso e funcionamento.
Controla a aplicação de procedimentos da empresa para cumprimento de
normas ambientais. Zela pela saúde e segurança da equipe de trabalho,
monitorando o uso de equipamentos de proteção individual e a realização
de exames periódicos de saúde. Identifica áreas de risco, tomando
medidas para seu isolamento e para análise dos fatores de risco.
Registra ocorrências de acidentes de trabalho.

\begin{Form}
\textbf{Esta correspondência está adequada?} \\ 
\ChoiceMenu[radio,radiosymbol=\ding{52},bordercolor=0 0.5 0,name=cbovalid_ 316320 ]{}{CBO deve ficar=sim, \hspace{6em} CBO não fica aqui=nao}
\end{Form}

\textbf{CBO: 316330  ( Técnico em planejamento de lavra de minas )}

Participa no planejamento de atividades de lavra de mina, programando
etapas do trabalho. Participa no planejamento da sequência de lavra e
exaustão da mina. Auxilia na seleção de equipe; na escolha de
equipamentos; na seleção e implantação de aplicativos de mineração; e na
definição de processos de trabalho. Participa na elaboração de programas
de manutenção de equipamentos e instalações. Auxilia na elaboração de
projetos de locação de conjunto de correias transportadoras de
minerodutos. Participa na elaboração de projetos ambientais. Contribui
na elaboração de plano de custos operacionais. Participa na elaboração
de requerimentos de alvará de pesquisa e lavra. Presta assistência
técnica à lavra e à exploração de mina, realizando levantamento
topográfico de área de lavra e deposição de estéril; participando na
elaboração de planos de drenagem; coordenando abertura de vias de acesso
à área de lavra; e definindo avanços de lavra, visando à
operacionalização dos planos de lavra de curto e de longo prazo.
Controla a produção de minério e estéril. Monitora o teor de minério.
Acompanha perfuração e desmonte de rocha e verifica desempenho de
explosivos. Monitora pistas de rolamento, movimentação de equipamentos
em mina e tempos de carregamentos de produto mineral. Coordena a
execução da drenagem da jazida. Monitora o nível de lençol freático e o
nível de sedimentos de bacia de contenção. Controla custos e materiais.
Coleta amostras de minerais de frentes de lavra e de produtos da mina
estocados. Processa e analisa dados. Confecciona mapas e seções
geológicas de lavra. Cria modelo de bloco baseado em modelo geológico.
Cuba reserva lavrável e corpo mineralizado. Auxilia na interpretação de
corpo mineralizado. Gera banco de dados de perfis elétricos, acústicos,
radioativos e geológicos. Auxilia na estimativa de vida útil da mina.
Pode utilizar instrumentos de medição em laboratório e em campo e
sistema automatizado para coleta e uso de dados geológicos. Pode fazer
uso de sistema de sensoriamento remoto e de outras tecnologias digitais
para a atualização topográfica periódica. Participa na elaboração de
padrões técnicos, procedimentos e instruções operacionais de trabalho,
discutindo condições de produção com outros setores. Pode supervisionar
equipe de trabalho, distribuindo tarefas, avaliando desempenho e
desenvolvendo treinamentos. Monitora a organização do ambiente de
trabalho e a conservação de ferramentas e instrumentos limpos,
acondicionados e em plenas condições de uso e funcionamento.
Supervisiona a manutenção de equipamentos. Controla a aplicação de
procedimentos da empresa para cumprimento de normas ambientais.
Fiscaliza a emissão de poluentes. Inspeciona a quantidade de rejeitos e
o nível de assoreamento de rios. Representa departamentos em inspeção
ambiental. Zela pela saúde e segurança da equipe de trabalho,
monitorando o uso de equipamentos de proteção individual e a realização
de exames periódicos de saúde. Identifica áreas de risco na mina,
tomando medidas para seu isolamento e para análise dos fatores de risco.
Registra ocorrências de acidentes de trabalho.

\begin{Form}
\textbf{Esta correspondência está adequada?} \\ 
\ChoiceMenu[radio,radiosymbol=\ding{52},bordercolor=0 0.5 0,name=cbovalid_ 316330 ]{}{CBO deve ficar=sim, \hspace{6em} CBO não fica aqui=nao}
\end{Form}

\newpage
\section{ 11013 -- Técnico em Pesca }

\textbf{Perfil do Curso (CNCT)}

O Técnico em Pesca será habilitado para:

Planejar e executar atividades relacionadas à pesca extrativa, a
operações de embarque e desembarque de pescado. Prestar assistência
técnica e assessoria ao estudo e ao desenvolvimento de projetos e
pesquisas tecnológicas ou aos trabalhos de vistoria, perícia,
arbitramento e consultoria. Elaborar orçamentos, laudos, pareceres,
relatórios e projetos, inclusive de incorporação de novas tecnologias.
Analisar a viabilidade técnica e econômica de propostas e projetos
pesqueiros. Prestar assistência técnica às áreas de crédito rural,
agroindustrial e de impacto ambiental. Utilizar procedimentos de armação
para a pesca. Construir e efetuar a manutenção de apetrechos de pesca
(redes, iscas, armadilhas e anzóis). Realizar procedimentos de
beneficiamento e processamento do pescado nas embarcações de pesca e em
frigoríficos. Aplicar boas práticas de manipulação e fabricação,
supervisionar as etapas de conservação, processamento, beneficiamento e
comercialização do pescado. Elaborar, aplicar e monitorar programas
profiláticos, higiênicos e sanitários na cadeia produtiva do pescado
Implantar e gerenciar sistemas de controle de qualidade na produção de
pescado. Emitir laudos e documentos de classificação e exercer a
fiscalização de produtos pesqueiros. Treinar e conduzir equipes nas suas
modalidades de atuação profissional. Prevenir situações de risco à
segurança no trabalho. Aplicar as legislações pertinentes ao processo
produtivo e ao meio ambiente. Identificar e aplicar técnicas
mercadológicas para distribuição e comercialização do pescado. Executar
a gestão econômica e financeira da produção pesqueira. Administrar e
gerenciar empreendimentos pesqueiros. Conduzir embarcações de pesca.
Operar equipamentos como radares, bússolas, barômetros e de Sistema de
Navegação Global por Satélite (GNSS).

Para a atuação como Técnico em Pesca, são fundamentais:

Conhecimentos e saberes relacionados a processos de gestão de negócios
voltados à pesca, à legislação pesqueira brasileira, à legislação
ambiental, à gestão de projetos, à gestão de processos, ao
empreendedorismo, à estatística pesqueira, a políticas públicas para o
desenvolvimento da pesca, à higiene e sanitização pessoal, das
embarcações, das instalações de processamento e do pescado alinhado aos
requisitos e protocolos internacionais. Domínio de uso de tecnologias da
informação e bases tecnológicas. Habilidade de comunicação, resolução de
situações-problema, trabalho em equipe e gestão de conflitos.

\textbf{CBOs correspondidos e seus perfis (presentes no CNCT, e presentes na predição da IA)}

\textbf{CBO: 341220  ( Patrão de pesca de alto-mar )}

Realiza operações de navegação em mar aberto com embarcação de pesca,
traçando a rota de navegação; relatando posição da embarcação e hora
estimada de chegada; e monitorando condições de navegabilidade. Guarnece
o passadiço e o timão. Aciona equipamentos auxiliares e luzes
regulamentares da embarcação. Opera equipamentos de embarcação de pesca
na navegação em mar aberto, tais como equipamentos de comunicação, de
combate a incêndio, de salvatagem e de orientação para posicionamento
geográfico. Opera piloto automático, instrumento ecossondador e sonar.
Atraca e desatraca embarcação de pesca na navegação em mar aberto,
acionando a seção de máquinas; trocando informações com a estação de
apoio; e analisando as condições intervenientes (maré, vento, tráfego de
embarcações, visibilidade e profundidade). Solicita serviços de apoio
portuário. Realiza as manobras necessárias para atracar ou desatracar a
embarcação. Reboca embarcações. Monitora carga e descarga da embarcação
de pesca na navegação em mar aberto, emitindo notificação de prontidão
da embarcação; verificando condições dos porões de carga; classificando
cargas; e elaborando o plano de carregamento. Verifica documentação das
cargas. Designa pessoal para carregamento e descarregamento. Monitora o
calado, a disposição das cargas e as condições de conservação de
pescados. Verifica condições e fixação das cargas. Informa unidade
receptora sobre características das cargas. Supervisiona serviços de
manutenção de embarcação de pesca na navegação em mar aberto,
solicitando serviços de reparo e fiscalizando reparos realizados a
bordo. Fiscaliza condições de conservação da embarcação e condições de
equipamentos de convés. Fiscaliza equipamentos de combate a incêndio e
de salvatagem. Fiscaliza dotação de material para contenção de
poluentes. Fiscaliza compartimentos habitáveis. Administra recursos
materiais e financeiros nas operações com embarcação de pesca na
navegação em mar aberto, requisitando materiais, comprando insumos,
requisitando provisões de alimentos e requisitando combustível,
lubrificantes e água. Administra custeio de bordo. Inventaria materiais.
Efetua pagamento para a tripulação. Supervisiona tripulação de
embarcação de pesca na navegação em mar aberto, divulgando normas e
regulamentos, determinando horário de trabalho, distribuindo tarefas
para guarnição e monitorando atividades. Treina novatos. Aplica
advertências e demais punições. Comanda embarcação de pesca na navegação
em mar aberto, de acordo com normas da autoridade marítima da Diretoria
de Portos e Costas da Marinha do Brasil, que define capacidades das
embarcações e limites em que pode atuar. Registra dados da embarcação de
pesca na navegação em mar aberto, preenchendo rol de tripulação (rol de
equipagem); emitindo documentação de entrada, saída e permanência no
porto; e atualizando cartas e publicações náuticas. Organiza
documentação de embarcação e carga. Escritura diário de navegação.
Preenche livro de carga e documentação de despacho de lixo. Escritura
diário de comunicação e redige atas de reuniões. Preenche mapa de bordo
para o Instituto Brasileiro do Meio Ambiente e dos Recursos Naturais
Renováveis (Ibama). Orienta tripulação sobre temas de segurança do
trabalho, tais como utilização de Equipamentos de Proteção individual
(EPI); situações de emergência; e condições e atos inseguros. Simula
situações adversas para treinamento da tripulação. Divulga informações
sobre saúde e orienta sobre questões ambientais.

\begin{Form}
\textbf{Esta correspondência está adequada?} \\ 
\ChoiceMenu[radio,radiosymbol=\ding{52},bordercolor=0 0.5 0,name=cbovalid_ 341220 ]{}{CBO deve ficar=sim, \hspace{6em} CBO não fica aqui=nao}
\end{Form}

\textbf{CBO: 341225  ( Patrão de pesca na navegação interior )}

Realiza operações de navegação interior com embarcação de pesca,
traçando a rota de navegação; relatando posição da embarcação e hora
estimada de chegada; e monitorando condições de navegabilidade. Guarnece
o passadiço e o timão. Aciona equipamentos auxiliares e luzes
regulamentares da embarcação. Opera equipamentos de embarcação de pesca
na navegação interior, tais como equipamentos de comunicação, de combate
a incêndio, de salvatagem e de orientação para posicionamento
geográfico. Opera piloto automático, instrumento ecossondador e sonar.
Atraca e desatraca embarcação de pesca na navegação interior, acionando
a seção de máquinas; trocando informações com a estação de apoio; e
analisando as condições intervenientes (vento, tráfego de embarcações,
visibilidade e profundidade). Solicita serviços de apoio portuário.
Realiza as manobras necessárias para atracar ou desatracar a embarcação.
Reboca embarcações. Monitora carga e descarga da embarcação de pesca na
navegação interior, emitindo notificação de prontidão da embarcação;
verificando condições dos porões de carga; e classificando cargas.
Verifica documentação das cargas. Designa pessoal para carregamento e
descarregamento. Monitora o calado, a disposição das cargas e as
condições de conservação de pescados. Verifica condições e fixação das
cargas. Informa unidade receptora sobre características das cargas.
Supervisiona serviços de manutenção de embarcação de pesca na navegação
interior, solicitando serviços de reparo e fiscalizando reparos
realizados a bordo. Fiscaliza condições de conservação da embarcação e
condições de equipamentos de convés. Fiscaliza equipamentos de combate a
incêndio e de salvatagem. Fiscaliza dotação de material para contenção
de poluentes. Fiscaliza compartimentos habitáveis. Administra recursos
materiais e financeiros nas operações com embarcação de pesca na
navegação interior, requisitando materiais, comprando insumos,
requisitando provisões de alimentos e requisitando combustível,
lubrificantes e água. Administra custeio de bordo. Inventaria materiais.
Efetua pagamento para a tripulação. Supervisiona tripulação de
embarcação de pesca na navegação interior, divulgando normas e
regulamentos, determinando horário de trabalho, distribuindo tarefas
para guarnição e monitorando atividades. Treina novatos. Aplica
advertências e demais punições. Comanda embarcação de pesca na navegação
interior, de acordo com normas da autoridade marítima da Diretoria de
Portos e Costas da Marinha do Brasil, que define capacidades das
embarcações e limites em que pode atuar. Registra dados da embarcação de
pesca na navegação interior, preenchendo rol de tripulação (rol de
equipagem); emitindo documentação de entrada, saída e permanência no
porto; e atualizando cartas e publicações náuticas. Organiza
documentação de embarcação e carga. Escritura diário de navegação.
Preenche livro de carga e documentação de despacho de lixo. Escritura
diário de comunicação e redige atas de reuniões. Preenche mapa de bordo
para o Instituto Brasileiro do Meio Ambiente e dos Recursos Naturais
Renováveis (Ibama). Orienta tripulação sobre temas de segurança do
trabalho, tais como utilização de Equipamentos de Proteção individual
(EPI); situações de emergência; e condições e atos inseguros. Simula
situações adversas para treinamento da tripulação. Divulga informações
sobre saúde e orienta sobre questões ambientais.

\begin{Form}
\textbf{Esta correspondência está adequada?} \\ 
\ChoiceMenu[radio,radiosymbol=\ding{52},bordercolor=0 0.5 0,name=cbovalid_ 341225 ]{}{CBO deve ficar=sim, \hspace{6em} CBO não fica aqui=nao}
\end{Form}

\newpage
\section{ 11016 -- Técnico em Recursos Pesqueiros }

\textbf{Perfil do Curso (CNCT)}

O Técnico em Recursos Pesqueiros será habilitado para:

Realizar operações do setor pesqueiro com base no manejo e na qualidade
da cadeia produtiva do pescado. Analisar e avaliar os aspectos técnicos,
sociais e econômicos da cadeia produtiva do setor pesqueiro. Prestar
assistência técnica e assessoria ao estudo e ao desenvolvimento de
projetos e pesquisas tecnológicas, ou aos trabalhos de vistoria,
perícia, arbitramento e consultoria. Elaborar orçamentos, laudos,
pareceres, relatórios e projetos, inclusive de incorporação de novas
tecnologias. Analisar a viabilidade técnica e econômica de propostas de
projetos aquícolas e pesqueiros. Prestar assistência técnica às áreas de
crédito rural, agroindustrial e impacto ambiental. Reconhecer os
aspectos fisiológicos, biológicos e ecológicos das espécies cultivadas e
exploradas pela pesca. Planejar e organizar ações de extensão pesqueira.
Planejar, organizar, dirigir e controlar as operações de pesca e de
despesca. Confeccionar, montar e operar apetrechos e equipamentos de
pesca e de aquicultura e dar manutenção adequada. Realizar procedimentos
para reprodução das principais espécies de interesse pesqueiro.
Monitorar e realizar o beneficiamento e o processamento de pescado em
embarcações pesqueiras e frigoríficos de pescado. Aplicar boas práticas
de manipulação e fabricação, supervisionar as etapas de conservação,
processamento, beneficiamento e comercialização do pescado. Elaborar,
aplicar e monitorar programas profiláticos, higiênicos e sanitários na
cadeia produtiva do pescado. Implantar e gerenciar sistemas de controle
de qualidade na produção do pescado. Emitir laudos e documentos de
classificação e exercer a fiscalização de produtos pesqueiros. Treinar e
conduzir equipes nas suas modalidades de atuação profissional. Prevenir
situações de risco à segurança no trabalho. Aplicar as legislações
pertinentes ao processo produtivo e ao meio ambiente. Identificar e
aplicar técnicas mercadológicas para distribuição e comercialização do
pescado. Executar a gestão econômica e financeira da produção pesqueira.
Administrar e gerenciar empreendimentos do setor pesqueiro. Conduzir
embarcações de pesca. Operar equipamentos como radares, bússolas,
barômetros e de Sistema de Navegação Global por Satélite (GNSS).

Para a atuação como Técnico em Recursos Pesqueiros, são fundamentais:

Conhecimentos e saberes relacionados a processos de gestão de negócios
voltados à aquicultura e pesca, à legislação pesqueira brasileira, à
legislação ambiental, ao planejamento de produção, à gestão de projetos,
à gestão de processos, ao empreendedorismo, ao mercado e comercialização
do pescado, à extensão aquícola e pesqueira, à aquicultura em águas da
União, a políticas públicas para o desenvolvimento da aquicultura e da
pesca, ao associativismo e ao cooperativismo. Domínio do uso de
tecnologias da informação.\\
Habilidade de comunicação, resolução de situações-problema, trabalho em
equipe e gestão de conflitos.

\textbf{CBOs correspondidos e seus perfis (presentes no CNCT, e presentes na predição da IA)}

\textbf{CBO: 321305  ( Técnico em piscicultura )}

Elabora projeto de implantação e de operação de sistema de criação de
peixes, de acordo com as características da região e com a aptidão do
ambiente natural. Faz análise da viabilidade técnica e econômica do
projeto. Apresenta as etapas de implantação. Relaciona máquinas,
equipamentos, ferramentas, instrumentos e materiais necessários. Detalha
tamanho de equipe de trabalho e suas atribuições. Elabora orçamento.
Presta orientações técnicas na implantação do projeto de piscicultura.
Seleciona sistema extensivo, semi-intensivo ou intensivo de criação.
Analisa os aspectos biológicos, fisiológicos e patológicos das
principais espécies e examina as condições do local, para recomendar a
espécie de peixes para criação. Orienta a preparação de locais -- tais
como viveiros, tanques escavados e fazendas marinhas -, para a criação
de peixes. Executa levantamento topográfico. Indica a necessidade de
drenagem, colocação de telas, vedação de comportas e outras
providências, conforme infraestrutura predefinida de criação. Acompanha
a montagem de equipamentos e sistemas de aeração. Avalia qualidade da
água utilizada no sistema de criação. Monitora o uso racional da água,
dimensionando sua vazão. Pode recomendar a instalação de viveiros
automatizados. Orienta o processo de reprodução natural ou artificial de
peixes, preferencialmente com espécies e matrizes melhoradas
geneticamente. Monitora a seleção, a mensuração e o acasalamento de
reprodutores, medindo fertilidade e testando grau de maturação de
óvulos. Acompanha a aplicação de técnicas artificiais de reprodução.
Monitora o desenvolvimento dos ovos e das larvas. Pode orientar a
aplicação de técnica de reversão sexual. Organiza o processo de
alimentação dos peixes, calculando seu peso médio e indicando tipo e
quantidade de alimentos. Define horários de alimentação. Demonstra o
processo de preparação de rações. Controla a sanidade dos peixes,
observando o comportamento dos espécimes para identificar eventuais
doenças e irregularidades da criação, como ferimentos ou presença de
fungos, bactérias e parasitas. Indica a necessidade de visita de médico
veterinário ou zootecnista, para ações e prescrições pertinentes.
Monitora a realização dos tratamentos, como isolamento de indivíduos em
quarentena e aplicação dos medicamentos prescritos na ração ou na água.
Orienta o uso de emissão de sinais sonoros; de alarmes automatizados com
sensores; e de outras estratégias, para afastar e coibir predadores.
Demonstra técnicas de anestesiar, insensibilizar ou promover a perda
imediata da consciência dos peixes, para realizar o abate sem
desconforto ou dor. Presta informações técnicas para execução do
beneficiamento dos peixes abatidos, orientando limpeza, descamação,
evisceração, corte, filetagem, fatiamento, moagem e congelamento.
Classifica os peixes. Demonstra a forma de embalar os peixes. Monitora o
transporte e o armazenamento da produção. Orienta a comercialização da
produção. Implanta e gerencia sistemas de controle de qualidade de
processos e produtos. Observa a organização do local de trabalho.
Monitora a aplicação de programas profiláticos, higiênicos e sanitários.
Orienta equipe para conservar ferramentas e instrumentos limpos,
acondicionados e em plenas condições de uso e funcionamento. Prepara
plano de manutenção de máquinas e equipamentos. Elabora relatório
técnico sobre resultados da implantação de projeto de piscicultura. Pode
ministrar palestras. Pode administrar empreendimentos e realizar
pesquisa e extensão, na área de piscicultura. Pode supervisionar equipe
de trabalho, distribuindo tarefas, avaliando desempenho e desenvolvendo
treinamentos. Mantém-se atualizado em relação às inovações na área de
piscicultura. Contribui para a redução dos impactos ambientais do
processo de criação, monitorando a aplicação das normas ambientais. Zela
pela segurança, prevenindo situações de risco. Utiliza e monitora o uso
de equipamentos de proteção individual. Presta primeiros socorros.

\begin{Form}
\textbf{Esta correspondência está adequada?} \\ 
\ChoiceMenu[radio,radiosymbol=\ding{52},bordercolor=0 0.5 0,name=cbovalid_ 321305 ]{}{CBO deve ficar=sim, \hspace{6em} CBO não fica aqui=nao}
\end{Form}

\textbf{CBO: 321315  ( Técnico em mitilicultura )}

Elabora projeto de implantação e de operação de sistema de criação de
mexilhões no mar. Faz análise da viabilidade técnica e econômica do
projeto. Apresenta as etapas de implantação. Relaciona equipamentos --
como as máquinas de semear, colher e desagregar mexilhões -,
embarcações, ferramentas, instrumentos e materiais necessários. Detalha
tamanho de equipe de trabalho e suas atribuições. Elabora orçamento.
Presta orientações técnicas na implantação do projeto de mitilicultura,
conforme características do ambiente de cultivo e da temperatura do
local. Recomenda espécie para criação. Indica local apropriado para
instalação de cultivo de mexilhões. Seleciona o sistema de criação.
Orienta a seleção e a instalação de boias e cabos. Mostra a importância
do conhecimento do regime de reprodução do mexilhão na região de
cultivo, para definir as épocas de captação de sementes. Indica formas
de captação, detalhando a explotação de bancos naturais; o uso de
coletores artificiais; o aproveitamento de sementes que se assentam
sobre as estruturas de cultivo, entre outras. Recomenda classificação de
sementes por tamanho, para processo de semeadura. Presta informações
técnicas sobre o manejo dos mexilhões. Orienta o controle da densidade
de sementes por metro de corda de cultivo. Monitora os mexilhões
submersos, que filtram a água do mar respirando e se alimentando.
Organiza o controle periódico e contínuo no local, para prevenir a
contaminação dos mexilhões pelo acúmulo de toxinas em seu sistema
digestivo ou tecidos. Indica a necessidade do controle constante nas
áreas de cultivo, testando a qualidade microbiológica da água e
avaliando uma amostra de mexilhões. Indica formas de controle dos
fatores de riscos e ameaças, tais como ataque às sementes por peixes e
outros predadores; roubo e vandalismo; e eventos climáticos extremos.
Orienta a avaliação de uma amostra de mexilhões antes da colheita,
verificando, entre outros fatores, o tamanho alcançado em relação ao
desejável no mercado alvo. Destaca a importância do uso de técnicas de
manuseio adequadas da produção, para garantir melhor qualidade e preço.
Estabelece a logística de armazenagem e de transporte da produção.
Orienta a comercialização da produção. Recomenda que os mexilhões sejam
transportados para a planta de processamento o mais rapidamente
possível, para entrarem na cadeia de frio. Implanta e gerencia sistemas
de controle de qualidade de processos e produtos. Elabora relatório
técnico sobre resultados da implantação de projeto de mitilicultura.
Pode ministrar palestras. Pode orientar a formação de grupos de
mitilicultores - como associações - para, entre outros objetivos,
promover a melhoria do rendimento econômico dos participantes. Pode
elaborar projeto de criação de mexilhões em água doce - lagos, lagoas e
em águas correntes, como riachos e rios - e orientar sua implantação.
Pode formular e implantar projetos de malacocultura, com atividade de
cultivo -- além de mexilhões - de ostras, vieiras e outros moluscos
bivalves. Pode administrar empreendimentos e realizar pesquisa e
extensão, na área de mitilicultura. Pode supervisionar equipe de
trabalho, distribuindo tarefas, avaliando desempenho e desenvolvendo
treinamentos. Mantém-se atualizado em relação às inovações relacionadas
à área de mitilicultura. Seleciona e avalia inovações - como
automatização de processos -- para adoção em novos projetos. Destaca a
importância da organização do local de trabalho e da aplicação de
programas profiláticos, higiênicos e sanitários. Orienta equipe para
conservar ferramentas e instrumentos limpos, acondicionados e em plenas
condições de uso e funcionamento. Prepara plano de manutenção de
embarcações, máquinas e equipamentos. Contribui para a redução dos
impactos ambientais do processo de criação, monitorando a aplicação das
normas ambientais. Zela pela segurança, prevenindo situações de risco.
Utiliza e monitora o uso de equipamentos de proteção individual. Presta
primeiros socorros.

\begin{Form}
\textbf{Esta correspondência está adequada?} \\ 
\ChoiceMenu[radio,radiosymbol=\ding{52},bordercolor=0 0.5 0,name=cbovalid_ 321315 ]{}{CBO deve ficar=sim, \hspace{6em} CBO não fica aqui=nao}
\end{Form}

\textbf{CBO: 321320  ( Técnico em ranicultura )}

Elabora projeto de implantação e de operação de sistema de criação de
rãs em cativeiro. Faz análise da viabilidade técnica e econômica do
projeto. Apresenta as etapas de implantação. Relaciona máquinas,
equipamentos, ferramentas, instrumentos e materiais necessários. Detalha
tamanho de equipe de trabalho e suas atribuições. Elabora orçamento.
Orienta a implantação do projeto de ranicultura, realizada considerando
clima na região; características do terreno; disponibilidade de água; e
distância dos centros de comercialização. Seleciona espécie de rãs para
criação. Indica sistema de criação. Recomenda a instalação de ranário em
terrenos com topografia preferencialmente plana ou pouco acidentada.
Orienta a implantação da estrutura do ranário, com instalações para
todas as fases - reprodução; larvicultura ou girinagem; engorda inicial
ou recria; e engorda final ou terminação - da vida de rãs. Enfatiza a
importância da disponibilidade de água com qualidade física e química
específica. Mostra como são feitos os exames dos principais parâmetros
da água, tais como oxigênio, pH, condutividade elétrica, alcalinidade
total, fósforo, cloretos, ferro, entre outros. Organiza as práticas de
manejo, registrando rãs para controle de seu histórico e seu
desenvolvimento; providenciando a limpeza das baias e a renovação da
água; e controlando a densidade recomendada de animais nas instalações.
Orienta a oferta de alimentos naturais e de ração. Mostra como cuidar do
bem-estar do plantel. Organiza o processo de reprodução natural de
machos e fêmeas aptos ao acasalamento. Principalmente em ranários de
maior porte ou naqueles especializados na produção de imagos para a
venda, mostra como executar a reprodução artificial, controlando a
temperatura, o fotoperíodo e a umidade do local; e submetendo os animais
a tratamento hormonal específico, para obter desovas semelhantes às
naturais. Indica procedimentos para controle da sanidade, como constante
renovação da água e limpeza dos tanques. Orienta o fornecimento de
informações - sobre sinais de possíveis doenças e lesões - ao médico
veterinário, que fará a prescrição e a aplicação de medicamentos, quando
necessário. Indica como fazer a coleta de animais mortos e o seu
transporte a um pequeno forno crematório na propriedade, para
erradicação de quaisquer agentes infecciosos presentes no criatório.
Orienta a seleção de rãs para o abate, respeitando um padrão de tamanho
e peso e verificando a ausência de feridas, traumas, deformações e
malformações. Demonstra como promover a insensibilidade aos animais, de
modo que os procedimentos de abate possam ser realizados sem dor ou
sofrimento. Organiza as etapas de processamento de abate e de preparação
das rãs abatidas para comercialização. Presta informações técnicas sobre
a comercialização da carne congelada de rãs e sobre a venda de
subprodutos - tais como vísceras, pele e gordura - do abate. Implanta e
gerencia sistemas de controle de qualidade de processos e produtos.
Elabora relatório técnico sobre resultados do projeto de ranicultura.
Pode ministrar palestras. Pode administrar empreendimentos e realizar
pesquisa e extensão, na área de ranicultura. Pode supervisionar trabalho
de equipe, avaliando desempenho e desenvolvendo treinamentos. Mantém-se
atualizado em relação às inovações em ranicultura. Seleciona e avalia
inovações - como automatização de processos -- para adoção em novos
projetos. Observa a organização do local de trabalho. Monitora a
aplicação de programas profiláticos, higiênicos e sanitários. Orienta a
conservação de ferramentas e instrumentos em plenas condições de
funcionamento. Prepara plano de manutenção de máquinas e equipamentos.
Contribui para redução dos impactos ambientais do processo de criação,
monitorando a aplicação das normas ambientais. Zela pela segurança,
prevenindo situações de risco. Utiliza e monitora o uso de equipamentos
de proteção individual. Presta primeiros socorros.

\begin{Form}
\textbf{Esta correspondência está adequada?} \\ 
\ChoiceMenu[radio,radiosymbol=\ding{52},bordercolor=0 0.5 0,name=cbovalid_ 321320 ]{}{CBO deve ficar=sim, \hspace{6em} CBO não fica aqui=nao}
\end{Form}

\textbf{CBO: 341225  ( Patrão de pesca na navegação interior )}

Realiza operações de navegação interior com embarcação de pesca,
traçando a rota de navegação; relatando posição da embarcação e hora
estimada de chegada; e monitorando condições de navegabilidade. Guarnece
o passadiço e o timão. Aciona equipamentos auxiliares e luzes
regulamentares da embarcação. Opera equipamentos de embarcação de pesca
na navegação interior, tais como equipamentos de comunicação, de combate
a incêndio, de salvatagem e de orientação para posicionamento
geográfico. Opera piloto automático, instrumento ecossondador e sonar.
Atraca e desatraca embarcação de pesca na navegação interior, acionando
a seção de máquinas; trocando informações com a estação de apoio; e
analisando as condições intervenientes (vento, tráfego de embarcações,
visibilidade e profundidade). Solicita serviços de apoio portuário.
Realiza as manobras necessárias para atracar ou desatracar a embarcação.
Reboca embarcações. Monitora carga e descarga da embarcação de pesca na
navegação interior, emitindo notificação de prontidão da embarcação;
verificando condições dos porões de carga; e classificando cargas.
Verifica documentação das cargas. Designa pessoal para carregamento e
descarregamento. Monitora o calado, a disposição das cargas e as
condições de conservação de pescados. Verifica condições e fixação das
cargas. Informa unidade receptora sobre características das cargas.
Supervisiona serviços de manutenção de embarcação de pesca na navegação
interior, solicitando serviços de reparo e fiscalizando reparos
realizados a bordo. Fiscaliza condições de conservação da embarcação e
condições de equipamentos de convés. Fiscaliza equipamentos de combate a
incêndio e de salvatagem. Fiscaliza dotação de material para contenção
de poluentes. Fiscaliza compartimentos habitáveis. Administra recursos
materiais e financeiros nas operações com embarcação de pesca na
navegação interior, requisitando materiais, comprando insumos,
requisitando provisões de alimentos e requisitando combustível,
lubrificantes e água. Administra custeio de bordo. Inventaria materiais.
Efetua pagamento para a tripulação. Supervisiona tripulação de
embarcação de pesca na navegação interior, divulgando normas e
regulamentos, determinando horário de trabalho, distribuindo tarefas
para guarnição e monitorando atividades. Treina novatos. Aplica
advertências e demais punições. Comanda embarcação de pesca na navegação
interior, de acordo com normas da autoridade marítima da Diretoria de
Portos e Costas da Marinha do Brasil, que define capacidades das
embarcações e limites em que pode atuar. Registra dados da embarcação de
pesca na navegação interior, preenchendo rol de tripulação (rol de
equipagem); emitindo documentação de entrada, saída e permanência no
porto; e atualizando cartas e publicações náuticas. Organiza
documentação de embarcação e carga. Escritura diário de navegação.
Preenche livro de carga e documentação de despacho de lixo. Escritura
diário de comunicação e redige atas de reuniões. Preenche mapa de bordo
para o Instituto Brasileiro do Meio Ambiente e dos Recursos Naturais
Renováveis (Ibama). Orienta tripulação sobre temas de segurança do
trabalho, tais como utilização de Equipamentos de Proteção individual
(EPI); situações de emergência; e condições e atos inseguros. Simula
situações adversas para treinamento da tripulação. Divulga informações
sobre saúde e orienta sobre questões ambientais.

\begin{Form}
\textbf{Esta correspondência está adequada?} \\ 
\ChoiceMenu[radio,radiosymbol=\ding{52},bordercolor=0 0.5 0,name=cbovalid_ 341225 ]{}{CBO deve ficar=sim, \hspace{6em} CBO não fica aqui=nao}
\end{Form}

\newpage
\section{ 11017 -- Técnico em Zootecnia }

\textbf{Perfil do Curso (CNCT)}

O Técnico em Zootecnia será habilitado para:

Planejar, organizar, dirigir e controlar a produção e a criação
sustentável de animais domésticos e silvestres, analisando as
características econômicas, sociais e ambientais. Elaborar, projetar e
executar projetos de produção pecuária, inclusive com a incorporação de
novas tecnologias. Prestar assistência técnica e assessoria ao estudo e
ao desenvolvimento de projetos e pesquisas tecnológicas e de
consultoria. Prestar assistência técnica às áreas de crédito rural.
Planejar, organizar e monitorar atividades de produção animal, processo
de aquisição, preparo, conservação e armazenamento de matérias primas e
produtos pecuários. Planejar, organizar e monitorar programas de
nutrição e manejo alimentar em projetos zootécnicos. Elaborar, aplicar e
monitorar programas de manejo preventivo, higiênico, sanitário,
nutricional e de reprodução animal. Realizar procedimentos de
inseminação artificial em animais. Aplicar métodos e programas de
reprodução animal e de melhoramento genético. Implantar e realizar o
manejo das pastagens. Aplicar procedimentos relativos ao preparo e
conservação do solo e da água. Realizar e monitorar a produção de
silagem e forragem. Aplicar técnicas de bem-estar animal na produção
pecuária. Projetar instalações zootécnicas. Prestar assistência técnica
à aplicação, à comercialização e ao manejo de produtos especializados
(sementes, fertilizantes, defensivos, pastagens, concentrados, sal
mineral, medicamentos e vacinas). Elaborar, aplicar e monitorar
programas profiláticos, higiênicos e sanitários na produção animal.
Emitir laudos e documentos de classificação e exercer a fiscalização de
produtos de origem animal. Implantar e gerenciar sistemas de controle de
qualidade na produção pecuária. Supervisionar o armazenamento, a
conservação, a comercialização e a industrialização dos produtos
pecuários. Treinar e conduzir equipes nas suas modalidades de atuação
profissional. Aplicar as legislações pertinentes ao processo produtivo e
ao meio ambiente. Aplicar práticas sustentáveis no manejo de conservação
do solo e da água. Identificar e aplicar técnicas mercadológicas para
distribuição e comercialização de produtos pecuários e animais. Executar
a gestão econômica e financeira da produção pecuária. Administrar e
gerenciar propriedades rurais. Operar e manejar máquinas, implementos,
equipamentos, veículos aéreos remotamente pilotados e equipamentos de
precisão para monitoramento remoto da produção pecuária.

Para a atuação como Técnico em Zootecnia, são fundamentais:

Conhecimentos e saberes relacionados à produção pecuária, à produção e
ao processamento de alimentos de origem animal. Atualização em relação
às inovações tecnológicas. Cooperação de forma construtiva e
colaborativa nos trabalhos em equipe, tomada de decisões. Adoção de
senso investigativo, visão sistêmica das atividades e processos,
capacidade de comunicação e argumentação, autonomia, proatividade,
liderança, respeito às diversidades nos grupos de trabalho, resiliência
frente aos problemas, organização, responsabilidade, visão crítica,
humanística, ética e consciência em relação ao impacto de sua atuação
profissional na sociedade e no ambiente.

\textbf{CBOs correspondidos e seus perfis (presentes no CNCT, e presentes na predição da IA)}

\textbf{CBO: 323105  ( Técnico em pecuária )}

Elabora projeto de implantação e de operação de sistema de criação de
animais - tais como bovinos, suínos, equinos, entre outros --
considerando características da região e particularidades de cada
espécie animal. Faz análise da viabilidade técnica e econômica do
projeto. Apresenta as etapas de implantação. Relaciona máquinas,
equipamentos, ferramentas, instrumentos e materiais necessários. Detalha
tamanho de equipe de trabalho e suas atribuições. Elabora orçamento.
Presta orientações técnicas na implantação do projeto de criação animal,
de acordo com os objetivos da produção. Orienta a escolha correta da
raça, para alcance de maior produtividade e de melhor rentabilidade.
Seleciona sistemas de criação extensiva, intensiva ou semi-intensiva.
Pode selecionar e indicar veículos aéreos remotamente pilotados e
equipamentos de precisão para monitoramento remoto da produção pecuária.
Presta informações técnicas para preparação da infraestrutura, podendo
propor inovações, tais como sistemas automatizados para transporte de
insumos; criadouros automatizados; e armazéns frigorificados. Orienta o
manejo animal racional e humanizado, realizando as atividades -- de
condução, adestramento, imobilização, dentre outras -- de modo a manter
relacionamento harmônico com os animais. Planeja, organiza e monitora
programas de nutrição e manejo alimentar. Orienta o cultivo de
forrageiras e grãos para ensilagem e fenação. Implanta e orienta o
manejo das pastagens. Programa e controla alimentação dos animais,
destacando os tipos de nutrientes imprescindíveis e as quantidades mais
próximas às necessidades dos animais; definindo frequência e horários da
alimentação; estabelecendo tipo e quantidade de rações, alimentos ``in
natura'', água, suplementos e vitaminas; e proporcionando alimentação de
pastagem, com aplicação de técnica de rodízio de pasto. Planeja e
organiza o controle da saúde e do bem-estar dos animais. Orienta a
condução dos processos de reprodução natural e o uso de técnicas
artificiais de reprodução. Monitora o manejo de animais recém-nascidos.
Orienta a observação do comportamento dos animais em todas as etapas da
criação, levantando sinais de possíveis doenças, lesões e traumatismos,
para solicitar visita de médico veterinário. Presta informações sobre
aplicação de métodos e programas de melhoramento genético. Apresenta
técnicas de anestesiar, insensibilizar ou promover a perda imediata da
consciência dos animais, para realizar o abate sem desconforto ou dor.
Presta informações técnicas para execução do beneficiamento dos animais
abatidos, orientando as ações promovidas, tais como classificação de
produtos; escolha de embalagens; e indicação de técnicas de conservação
dos produtos. Estabelece a logística de armazenagem e de transporte da
produção. Orienta a comercialização. Implanta e gerencia sistemas de
controle de qualidade de processos e produtos. Observa a organização do
local de trabalho. Elabora, aplica e monitora programas profiláticos,
higiênicos e sanitários na produção animal. Orienta equipe para
conservar ferramentas e instrumentos limpos, acondicionados e em plenas
condições de uso e funcionamento. Prepara plano de manutenção de
máquinas e equipamentos. Elabora relatório técnico sobre resultados da
implantação de projeto de pecuária. Pode ministrar palestras. Pode
administrar empreendimentos e realizar pesquisa e extensão, na área de
pecuária. Pode supervisionar equipe de trabalho, distribuindo tarefas,
avaliando desempenho e desenvolvendo treinamentos. Mantém-se atualizado
em relação às inovações na área de pecuária. Promove a sustentabilidade
pecuária, indicando a utilização de técnicas que reduzem os danos
causados à natureza. Orienta o uso de tecnologias de baixa emissão de
carbono e incentiva o aproveitamento de sobras de materiais. Zela pela
segurança, prevenindo situações de risco. Utiliza e monitora o uso de
equipamentos de proteção individual. Presta primeiros socorros.

\begin{Form}
\textbf{Esta correspondência está adequada?} \\ 
\ChoiceMenu[radio,radiosymbol=\ding{52},bordercolor=0 0.5 0,name=cbovalid_ 323105 ]{}{CBO deve ficar=sim, \hspace{6em} CBO não fica aqui=nao}
\end{Form}

\newpage
\section{ 12004 -- Técnico em Prevenção e Combate a Incêndios }

\textbf{Perfil do Curso (CNCT)}

O Técnico em Prevenção e Combate a Incêndio será habilitado para:

Executar atividades de prevenção e controle de incêndios e atendimentos
de emergência de resgate técnico, produtos perigosos e ambientais e
atendimento pré-hospitalar de emergências médicas. Analisar situações
que possam oferecer riscos para a vida. Interpretar projetos de proteção
contra incêndio. Analisar e implementar plano de emergência. Elaborar
procedimentos de abandono de áreas. Realizar atendimentos de primeiros
socorros e/ou atendimento pré-hospitalar de emergências médicas.
Selecionar, inspecionar e operar equipamentos e recursos materiais
empregados nos atendimentos às emergências. Realizar atendimentos de
prevenção e controle especializado de incêndio. Realizar atendimentos de
resgate técnico. Realizar atendimentos a emergências com produtos
perigosos. Analisar os principais potenciais de danos ambientais por
consequência de acidentes e/ou incêndios. Analisar os principais
potenciais de perdas de propriedades por consequência de acidentes e/ou
incêndios. Elaborar procedimentos operacionais empregados como padrão
para os atendimentos às emergências. Elaborar procedimentos
administrativos de elaboração de relatórios e gestão de pessoas.
Implantar e coordenar brigadas de incêndio e emergências. Realizar
atividades de ensino de educação continuada.

Para atuação como Técnico em Prevenção e Combate a Incêndio, são
fundamentais:

Conhecimentos e saberes relacionados aos processos produtivos do ramo de
atividade de atuação. Conhecimento das normas técnicas e
regulamentadoras relativas à atividade. Conhecimento dos códigos e leis
estaduais e municipais relativos à atividade. Liderança e gestão de
equipes. Conhecimentos e saberes relacionados à gestão de documentos.
Conhecimentos e saberes relacionados ao uso dos equipamentos de
prevenção e combate a incêndios, resgate e de primeiros socorros.

\textbf{CBOs correspondidos e seus perfis (presentes no CNCT, e presentes na predição da IA)}

\textbf{CBO: 510305  ( Supervisor de bombeiros )}

Atende clientes, informando-lhes sobre novas opções de prestação de
serviços de segurança e apresentando-lhes novas tecnologias utilizadas.
Dialoga com clientes sobre decisões que coloquem em risco a segurança.
Planeja atividades de segurança, discutindo proposta de trabalho e
elaborando cronograma. Prepara normas e procedimentos de trabalho.
Programa atividades de segurança, fazendo análise de risco e conduzindo
reuniões para apresentar e discutir os resultados dessa análise. Apura
fatos para prevenção de ocorrências. Apresenta projeto de segurança,
ressaltando o cumprimento da legislação pertinente. Mostra o
custo-benefício do projeto; sugere mudanças em equipamentos e no quadro
de funcionários efetivos; redimensiona os equipamentos de proteção a
incêndio; propõe medidas para cada tipo de sinistro; sugere providências
para redução de perdas; apresenta demarcação de áreas de risco de
incêndio em instalações; e propõe acompanhar vistorias e auditorias de
órgãos relacionados. Discute o projeto com os envolvidos. Dimensiona
pessoal necessário para realizar o projeto de segurança. Implanta postos
de trabalho e operacionaliza as atividades de bombeiros civis.
Administra pessoal, participando da seleção dos bombeiros civis,
verificando apresentação pessoal, elaborando escala de serviço e
conferindo a frequência. Desloca membro da equipe para suprir ausências.
Remaneja pessoal. Coordena o trabalho dos bombeiros civis, fiscalizando
o cumprimento dos procedimentos de trabalho; checando a disponibilidade
de equipamentos e recursos de trabalho; e conferindo a medição de
serviços. Realiza vistorias de rotina. Programa atividades de
treinamento e acompanha seu desenvolvimento, divulgando novas técnicas,
avaliando e validando os resultados da capacitação de pessoal. Programa
simulados de emergência. Supervisiona os serviços de segurança de
pessoas e de patrimônio, monitorando os locais e as atividades de risco
acentuado e propondo medidas preventivas e corretivas. Acompanha a
elaboração de boletim de ocorrência e investiga as causas do fato
ocorrido. Supervisiona brigadas de incêndio, monitora a inspeção de
equipamentos de proteção contra incêndio e libera equipamentos de
combate a incêndios para manutenção. Confere o estado de viaturas.
Relata ocorrências ao cliente contratante de serviços de segurança e
atende a solicitações extraordinárias. Coordena plano de emergência,
atendendo às ocorrências, dando apoio logístico, supervisionando plano
de abandono e providenciando o isolamento de vias de acesso e locais.
Solicita auxílio externo de equipe de emergência. Encaminha imprensa ao
setor competente. Presta esclarecimentos à gerência e/ou diretoria da
empresa, coletando e apresentando informações sobre o atendimento e/ou
serviço prestado. Interage com instituições, estabelecendo contatos com
órgãos públicos relacionados e participando de reuniões para troca de
experiência na área. Atende às autoridades e aos órgãos públicos.
Participa de reuniões para adoção de medidas preventivas. Atua junto à
defesa civil nas emergências. Representa a empresa em eventos externos.
Pesquisa novas tecnologias. Confere contratos de prestadoras de
serviços, para compará-los aos que utiliza. Contata empresas de
manutenção de equipamentos, para conhecer seus serviços. Pode prestar ou
providenciar primeiros socorros, quando necessário.

\begin{Form}
\textbf{Esta correspondência está adequada?} \\ 
\ChoiceMenu[radio,radiosymbol=\ding{52},bordercolor=0 0.5 0,name=cbovalid_ 510305 ]{}{CBO deve ficar=sim, \hspace{6em} CBO não fica aqui=nao}
\end{Form}

\textbf{CBO: 517105  ( Bombeiro de aeródromo )}

Prepara-se para o atendimento de ocorrências em aeródromos, conferindo o
funcionamento de equipamentos e viaturas. Requisita serviço de
manutenção de equipamentos e viaturas, quando necessário. Faz a
especificação de equipamentos, para aquisição. Pratica exercícios
físicos e outros condicionamentos, para manter a forma e a resistência
física e psicológica em momentos de atuação. Previne acidentes - como
incêndio, vazamento e explosão -, fazendo reconhecimento do local,
mapeando área de risco e vistoriando instalações e sistema de proteção
contra incêndio. Sinaliza locais de risco e estabelece rota de fuga.
Realiza análise de risco em diferentes tipos de aeronaves. Atende
ocorrência de incêndio no aeródromo, definindo plano de ação,
posicionando viatura e controlando tempo de resposta. Tria as
informações recebidas, para avaliar proporção e tipo de incêndio,
analisar situações de risco e classificar a ocorrência. Evacua o local e
isola a área do incêndio. Acopla mangueiras de água e confina o combate
à área atingida. Extingue o fogo. Revolve resíduos do incêndio,
eliminando possíveis focos de fogo e situações de risco. Preserva o
local para a perícia. Aplica táticas de resgate e realiza combate a
incêndio em aeronaves. Executa salvamento em altura e em operações de
busca para localizar, resgatar, desencarcerar e socorrer vítimas.
Executa salvamento terrestre, cavando o local de soterramento, cortando
ferragens e retirando vítimas. Executa salvamento aquático, retirando a
vítima da água e resgatando-a. Executa salvamento em altura, localizando
a vítima, aproximando-se dela, abordando-a e realizando o seu resgate.
Pode capturar aves e animais, balões, drones e outros objetos que
coloquem aeronaves em risco. Controla acidentes com produtos perigosos,
identificando os produtos e contendo o vazamento de recipiente. Demarca
a distância de segurança do produto vazado e afasta o público do local.
Monitora condições atmosféricas, para antecipar possível espalhamento do
produto vazado. Presta primeiros socorros, verificando parâmetros
básicos de frequência respiratória, frequência cardíaca, pressão
arterial, saturação de oxigênio, temperatura e nível de consciência da
vítima. Libera as vias aéreas da vítima. Em caso de parada cardíaca,
procede a ressuscitação cardiopulmonar (RCP) e a desfibrilação
automática externa no local. Constata a existência de lesões e
hemorragias e providencia a hemostasia. Em caso de fratura, faz a
imobilização cervical e vertebral no local, estabiliza, acalma e
transporta a vítima para o centro médico. Pode efetuar anamnese da
vítima, até a chegada ao socorro médico. Pode ministrar programa de
treinamento, capacitando brigadas de incêndio e corpo voluntário de
emergência. Simula ocorrências com funcionários de empresas e abandono
de local. Pode participar de campanhas educativas e proferir palestras.
Comunica-se orientando o público, conversando e ouvindo relatos de
vítimas e testemunhas. Comunica-se por sinais. Troca informações com o
centro de operações. Relata ocorrências em formulários, conforme normas
estabelecidas. Elabora relatórios, apresentando estatísticas. Atua
segundo as normas de segurança e biossegurança, submetendo-se a exames
periódicos, tomando vacinas e usando equipamentos de proteção individual
e de proteção coletiva. Seleciona vestimentas de acordo com a
ocorrência. Higieniza e descontamina equipamentos.

\begin{Form}
\textbf{Esta correspondência está adequada?} \\ 
\ChoiceMenu[radio,radiosymbol=\ding{52},bordercolor=0 0.5 0,name=cbovalid_ 517105 ]{}{CBO deve ficar=sim, \hspace{6em} CBO não fica aqui=nao}
\end{Form}

\textbf{CBO: 517110  ( Bombeiro civil )}

Prepara-se para o atendimento de ocorrências, examinando o funcionamento
de equipamentos e veículo. Requisita serviço de manutenção de
equipamentos e veículo, quando necessário. Pratica exercícios físicos e
outros condicionamentos, para manter a forma e a resistência física e
psicológica em momentos de atuação. Toma conhecimento de plano de
emergência do local, elaborado por profissional habilitado ou empresa
especializada. Previne acidentes - como incêndio e vazamentos -, fazendo
reconhecimento do local, mapeando áreas de risco e vistoriando
instalações. Realiza inspeção periódica dos equipamentos de combate a
incêndio - incluindo seus testes -, requisitando serviço de manutenção,
quando necessário. Inspeciona e sinaliza locais de risco e estabelece
rotas de fuga. Atende ocorrência de incêndio, definindo plano de ação.
Tria as informações recebidas, para avaliar proporção e tipo de incêndio
e analisar situações de risco. Evacua o local e isola a área do
incêndio. Desenergiza o local. Acopla mangueiras de água e confina o
combate à área atingida. Extingue o fogo. Revolve resíduos do incêndio,
eliminando possíveis focos de fogo e situações de risco. Preserva o
local para a perícia. Pode escorar paredes, pavimentos e telhados.
Executa salvamento em altura e em operações de busca. Localiza a vítima
em altura, aproximando-se dela, abordando-a e realizando o seu resgate.
Realiza operação de busca em terra, cavando o local de soterramento,
cortando ferragens e retirando vítimas. Executa salvamento aquático,
retirando a vítima da água e resgatando-a. Pode capturar animais em
geral, mantendo-os em segurança até que possam ser resgatados. Controla
acidentes com produtos perigosos, identificando os produtos. Contém o
vazamento de recipiente, criando diques de contenção. Demarca a
distância de segurança do produto vazado e afasta o público do local.
Monitora condições atmosféricas, para antecipar possível espalhamento do
produto vazado. Reembala os produtos perigosos. Pode montar corredor de
descontaminação. Presta primeiros socorros às vítimas de incidentes e
acidentes, verificando parâmetros básicos de frequência respiratória,
frequência cardíaca, pressão arterial, saturação de oxigênio,
temperatura e nível de consciência da vítima. Libera as vias aéreas da
vítima. Em caso de parada cardíaca, procede a ressuscitação
cardiopulmonar (RCP) e a desfibrilação automática externa no local.
Constata a existência de lesões e hemorragias e providencia a
hemostasia. Em caso de fratura, faz a imobilização cervical e vertebral
no local, estabiliza, acalma e transporta a vítima para o centro médico.
Pode efetuar anamnese da vítima, até a chegada ao socorro médico. Simula
ocorrências e abandono do local, para treinar pessoas que moram ou
trabalham em prédio e funcionários de empresas. Pode participar de
campanhas educativas e ministrar treinamentos para capacitar corpo
voluntário de emergência. Comunica-se orientando o público, conversando
e ouvindo relatos de vítimas. Troca informações com o centro de
operações. Elabora relatórios das ocorrências, conforme normas
estabelecidas. Pode apresentar eventuais sugestões para melhoria das
condições de segurança. Atua segundo as normas de segurança e
biossegurança, submetendo-se a exames periódicos, tomando vacinas e
usando equipamentos de proteção individual e de proteção coletiva.
Higieniza e descontamina equipamentos.

\begin{Form}
\textbf{Esta correspondência está adequada?} \\ 
\ChoiceMenu[radio,radiosymbol=\ding{52},bordercolor=0 0.5 0,name=cbovalid_ 517110 ]{}{CBO deve ficar=sim, \hspace{6em} CBO não fica aqui=nao}
\end{Form}

\newpage
\section{ 12005 -- Técnico em Segurança do Trabalho }

\textbf{Perfil do Curso (CNCT)}

O Técnico em Segurança do Trabalho será habilitado para:

Elaborar e implementar políticas de saúde no trabalho, identificando
variáveis de controle e ações educativas para prevenção e manutenção da
qualidade de vida do trabalhador. Desenvolver ações educativas na área
de saúde e segurança do trabalho. Investigar, analisar e recomendar
medidas de prevenção e controle de acidentes. Realizar estudo da relação
entre ocupações dos espaços físicos com as condições necessárias.
Promover a saúde e proteger a integridade do trabalhador em seu local de
atuação. Analisar os métodos e os processos laborais. Identificar
fatores de risco de acidentes do trabalho, de doenças profissionais e de
trabalho e de presença de agentes ambientais agressivos ao trabalhador.
Realizar procedimentos de orientação sobre medidas de eliminação e
neutralização de riscos. Elaborar procedimentos de acordo com a natureza
da empresa. Promover programas, eventos e capacitações de prevenção de
riscos ambientais. Divulgar normas e procedimentos de segurança e
higiene ocupacional. Indicar, solicitar e inspecionar equipamentos de
proteção coletiva e individual contra incêndio. Levantar e utilizar
dados estatísticos de doenças e acidentes de trabalho para ajustes das
ações prevencionistas. Produzir relatórios referentes à segurança e à
saúde do trabalhador.

Para atuação como Técnico em Segurança do Trabalho, são fundamentais:

Conhecimentos e saberes relacionados aos processos produtivos do ramo de
atividade de atuação. Conhecimento das normas técnicas e
regulamentadoras. Liderança e gestão de equipes. Conhecimentos e saberes
relacionados à gestão de documentos. Conhecimentos e saberes
relacionados ao uso de instrumentos de higiene ocupacional.

\textbf{CBOs correspondidos e seus perfis (presentes no CNCT, e presentes na predição da IA)}

\textbf{CBO: 351605  ( Técnico em segurança do trabalho )}

Elabora, implanta e implementa política de saúde e segurança do
trabalho, definindo metas, prioridades e responsabilidades. Divulga a
política entre os trabalhadores. Acompanha e avalia os resultados,
podendo promover alterações para aperfeiçoamento e atualização. Analisa
e avalia o ambiente de trabalho, as instalações e os processos laborais,
identificando fatores de risco de acidentes do trabalho, de doenças
profissionais e de presença de agentes ambientais agressivos ao
trabalhador. Participa da seleção de tecnologias e processos de
trabalho, avaliando implantação e impactos por meio de análise coletiva
do trabalho (ACT) e análise ergonômica do trabalho (AET). Emite parecer
sobre equipamentos, máquinas e processos. Coleta, organiza e analisa
comunicações de acidentes de trabalho (CAT). Investiga e analisa
acidentes para determinar causas e recomendar medidas de prevenção e
controle. Levanta e utiliza dados estatísticos de doenças e acidentes de
trabalho para ajustes das ações prevencionistas. Adota medidas de
controle de riscos ocupacionais por meio de ações e programas de saúde e
segurança do trabalho. Realiza procedimentos de orientação sobre medidas
de eliminação e neutralização de riscos. Seleciona, controla, orienta e
fiscaliza o uso de EPC (Equipamento de Proteção Coletiva) e EPI
(Equipamento de Proteção Individual). Participa do desenvolvimento de
ações educativas e capacitações, elaborando cronograma, preparando
recursos didáticos e formando multiplicadores. Produz relatórios
referentes à segurança e à saúde do trabalhador. Elabora manual do
sistema de gestão de saúde e segurança no trabalho. Controla atualização
de documentos, normas e legislação. Atua em equipes multidisciplinares
de saúde e segurança ocupacionais.

\begin{Form}
\textbf{Esta correspondência está adequada?} \\ 
\ChoiceMenu[radio,radiosymbol=\ding{52},bordercolor=0 0.5 0,name=cbovalid_ 351605 ]{}{CBO deve ficar=sim, \hspace{6em} CBO não fica aqui=nao}
\end{Form}

\textbf{CBO: 351610  ( Técnico em higiene ocupacional )}

Participa do planejamento da política de saúde e segurança no trabalho,
diagnosticando e analisando tecnicamente as condições ambientais,
colaborando na definição dos indicadores e mostrando os impactos -
econômicos e na saúde do trabalhador - da implantação da política.
Participa da implantação da política, divulgando na instituição ou
empresa e acompanhando resultados. Antecipa os riscos físicos, químicos
e biológicos, realizando avaliação dos riscos potenciais e estabelecendo
medidas preventivas. Reconhece, por meio de avaliação qualitativa, os
agentes físicos, químicos e biológicos presentes no ambiente de trabalho
que possam causar danos à saúde e à integridade dos trabalhadores. Para
tanto, analisa projetos, processos, instalações e equipamentos. Avalia,
quantitativamente, o grau de severidade dos processos e riscos em
relação aos limites de tolerância, definidos em normas. Controla a
exposição aos riscos, analisando os valores de concentração e
intensidade dos agentes em relação aos limites permitidos na legislação
nacional e internacional. Estabelece medidas para minimização ou
eliminação dos riscos, antecipados, reconhecidos e avaliados no ambiente
de trabalho. Monitora a eficácia das medidas aplicadas no controle de
risco, acompanhando a implantação de melhorias, fazendo a revisão
periódica das medidas de proteção aplicadas e avaliando os programas de
gerenciamento de riscos ambientais. Seleciona, controla, orienta e
fiscaliza o uso de EPC (Equipamento de Proteção Coletiva) e EPI
(Equipamento de Proteção Individual). Analisa fatores ergonômicos e a
ocorrência de acidentes nos locais de trabalho para entender seus
efeitos na saúde e bem-estar dos trabalhadores. Desenvolve ações
educativas, estruturadas a partir das necessidades levantadas. Elabora
cronograma, prepara os materiais didáticos e forma os multiplicadores.
Avalia os resultados das ações educativas. Participa de perícias e
fiscalizações. Elabora relatórios e pareceres técnicos. Documenta
procedimentos e normas de sistemas de segurança e controla sua
atualização. Atua em equipes multidisciplinares de saúde e segurança
ocupacionais.

\begin{Form}
\textbf{Esta correspondência está adequada?} \\ 
\ChoiceMenu[radio,radiosymbol=\ding{52},bordercolor=0 0.5 0,name=cbovalid_ 351610 ]{}{CBO deve ficar=sim, \hspace{6em} CBO não fica aqui=nao}
\end{Form}

\newpage
\section{ 13001 -- Técnico em Agenciamento de Viagem }

\textbf{Perfil do Curso (CNCT)}

O Técnico em Agenciamento de Viagens será habilitado para:

Planejar, organizar e customizar roteiros e serviços turísticos. Prestar
serviços de consultoria sobre viagens, fornecendo informações sobre
atrativos e destinos turísticos. Realizar reservas e intermediar a
contratação e a comercialização de serviços como: transportes,
hospedagem, locação de veículos, seguros viagens, guiamento, ingressos
de atrações turísticas e entretenimento e alimentação. Auxiliar no
receptivo turístico. Prestar assistência ao viajante, orientando sobre
documentação de viagens, cotações cambiais e procedimentos de embarque e
desembarque. Desenvolver atividades de pós-vendas. Executar demais
atividades administrativas relacionadas à função.

Para atuação como Técnico em Agenciamento de Viagens, são fundamentais:

Conhecimentos multidisciplinares relacionados aos aspectos
socioculturais, históricos, geográficos, legais, econômicos e
financeiros, bem como conhecimentos técnicos sobre produtos, serviços e
tecnologias pertinentes às operações turísticas. Comunicação clara e
cordial, respeito à diversidade, atenção à sustentabilidade, atitude
empreendedora, trabalho colaborativo, inovação, criatividade e
flexibilidade para a solução de problemas e gestão de conflitos.

\textbf{CBOs correspondidos e seus perfis (presentes no CNCT, e presentes na predição da IA)}

\textbf{CBO: 354810  ( Operador de turismo )}

Pesquisa destinos e oportunidades turísticas, conforme solicitação e
necessidades de diferentes perfis de clientes ou grupos de clientes,
atuando de forma autônoma ou vinculado a diferentes tipos de agências de
turismo. Opera sistemas de busca e reserva de destinos, roteiros e
itinerários, atrativos e oportunidades turísticas. Busca, também,
empresas para prestação de serviços, tais como hospedagem, alimentação e
espaços para visitação. Pode realizar visitas prévias de avaliação dos
serviços turísticos. Comercializa produtos e serviços turísticos
nacionais e internacionais. Pode vender pacotes de agências de turismo
ou contratar diretamente prestadores de serviços -- como fornecedores de
transporte, hospedagem e alimentação -, calculando custo e preços. Pode
preparar documentos para formalização de contratos. Pode realizar
cancelamento de produtos turísticos e traslados. Orienta turista ou
grupo sobre normas e legislações existentes nos países ou regiões dos
roteiros ou itinerários, bem como regras dos locais e dos serviços
utilizados, informando-os com clareza e objetividade. Esclarece dúvidas
sobre a necessidade de vistos e documentações, assim como fornece outras
informações referentes às especificidades das experiências turísticas.
Presta assistência no embarque. Atua na intermediação entre turistas,
empresas e prestadores de serviços turísticos, assistindo o cliente - ou
grupo de clientes - na solução de eventuais intercorrências durante a
viagem ou experiência turística. Verifica satisfação dos clientes, no
final do roteiro. Interage com empresas da cadeia produtiva do turismo
de diferentes áreas - como financeira, administrativa e vendas -, com
guias turísticos e com fornecedores. Elabora conteúdo promocional e
utiliza recursos tecnológicos para divulgação de atrativos e destinos
turísticos. Mantém e atualiza cadastro de público-alvo, em sistema
informatizado.

\begin{Form}
\textbf{Esta correspondência está adequada?} \\ 
\ChoiceMenu[radio,radiosymbol=\ding{52},bordercolor=0 0.5 0,name=cbovalid_ 354810 ]{}{CBO deve ficar=sim, \hspace{6em} CBO não fica aqui=nao}
\end{Form}

\textbf{CBO: 354815  ( Agente de viagem )}

Identifica necessidade do cliente -- inclusive autoridades e
profissionais que realizam viagens a negócios -, para atendimento com
montagem de pacotes de turismo ou com prestação de serviços de viagem.
Monta pacotes de turismo, realizando pesquisa para identificar atrativos
turísticos. Elabora roteiros personalizados e para grupos. Contata
prestadores de serviços turísticos e de apoio, podendo realizar visitas
prévias de avaliação do atendimento. Define cronogramas e atividades de
pacotes turísticos. Calcula custo e elabora orçamento para aprovação do
cliente. Presta serviços de viagem, vendendo e emitindo passagens
aéreas, hidroviárias ou terrestres; vendendo hospedagem e emitindo
vouchers; alugando carros; e vendendo seguro de viagem e assistência.
Vende excursões nacionais e internacionais, roteiros opcionais e pacotes
turísticos, oferecidos por agências. Para montagem de pacotes de turismo
ou serviços de viagem, contrata fornecedores de transporte, de
hospedagem e de apoio a eventos. Contrata guias de turismo e atrações
artísticas e culturais. Prepara documentos para formalização de
contratos. Encaminha relação de clientes aos prestadores de serviços
contratados. No final da prestação de serviços, libera ordem de
pagamento. Presta aos clientes atendimento personalizado, orientando-os
sobre documentação de viagem -- inclusive visto e comprovante de
vacinação -; conferindo emissão de bilhetes e ocorrências de alterações;
providenciando reserva de serviços; controlando prazos de reservas;
programando o pré-pagamento para datas especiais; solucionando problemas
surgidos durante a viagem; e prestando assistência no embarque ou no
desembarque da viagem. Verifica a satisfação dos clientes, analisando as
informações coletadas. Utiliza recursos tecnológicos para divulgação dos
destinos, roteiros e itinerários turísticos. Pode utilizar software de
Gerenciamento da Relação com o Cliente (Customer Relationships
Management-CRM) para comunicação com usuários. Empreende registro de
informações sobre emissão de documentos e reservas em sistemas internos
de agência de viagem, informatizados ou não. Elabora relatórios sobre
montagem de pacotes turísticos e venda de serviços.

\begin{Form}
\textbf{Esta correspondência está adequada?} \\ 
\ChoiceMenu[radio,radiosymbol=\ding{52},bordercolor=0 0.5 0,name=cbovalid_ 354815 ]{}{CBO deve ficar=sim, \hspace{6em} CBO não fica aqui=nao}
\end{Form}

\newpage
\section{ 13002 -- Técnico em Eventos }

\textbf{Perfil do Curso (CNCT)}

O Técnico em Eventos será habilitado para:

Prospectar e planejar eventos de acordo com o público-alvo, as
necessidades dos clientes e o mercado. Promover ações de comercialização
e divulgação relacionadas ao evento. Coordenar e realizar a execução do
evento: montagem, decoração, serviços técnicos, logísticos e
operacionais. Apoiar o planejamento e a operação de serviços de
alimentos e bebidas. Realizar procedimentos de cerimonial e protocolo.
Coordenar a recepção de eventos. Realizar o pós-evento.

Para atuação como Técnico em Eventos, são fundamentais:

Conhecimentos multidisciplinares sobre aspectos socioculturais e
econômicos dos locais onde serão realizados os eventos, bem como
conhecimentos técnicos sobre classificação e tipologias de eventos,
hospitalidade, sistemas de realização de eventos, além das legislações
que visam a garantir a integridade e a segurança dos participantes.
Comunicação clara e cordial, respeito às diversidades, atitude
empreendedora, trabalho colaborativo, atenção à sustentabilidade,
proatividade, criatividade, flexibilidade para solução de problemas e
gestão de conflitos.

\textbf{CBOs correspondidos e seus perfis (presentes no CNCT, e presentes na predição da IA)}

\textbf{CBO: 354820  ( Organizador de evento )}

Planeja o evento, analisando o público-alvo, as características
propostas pelo cliente e a finalidade, que pode ser social (como
casamentos e aniversários), empresarial ou de negócios (como feiras e
reuniões), comunitária (como jantares beneficentes), acadêmica (como
festas de formatura), cultural (como espetáculos teatrais e concertos
musicais), esportiva (como jogos) ou educacional (como palestras e
congressos). Realiza pesquisas sobre o tema do evento. Propõe
programação, local e data do evento. Levanta necessidade de recursos
humanos, materiais e financeiros. Elabora orçamento, para aprovação do
cliente. Organiza o evento, elaborando cronograma, dimensionando
leiaute, definindo programação visual e preparando cerimonial e/ou
roteiro. Seleciona empresas para prestar serviços, tais como apoio a
eventos, definição e preparação de alimentos e bebidas e apresentações
artísticas e culturais. Prepara documentos para formalização de
contratos. Solicita alvará e autorizações para realização do evento.
Garante o cumprimento de normas sanitárias estabelecidas nos âmbitos
local, estadual e nacional. Observa a legislação vigente em relação à
participação de crianças e adolescentes no evento. Controla o orçamento.
Pode buscar patrocínio ou fomento, junto a organizações e instituições
públicas ou privadas. Realiza a divulgação do evento. Capta cadastro do
público-alvo, expede convites e registra participantes. Contrata
fornecedores de serviços. Contrata atrações artísticas e culturais.
Coordena a montagem da infraestrutura do evento e a decoração do espaço.
Controla o armazenamento e o manuseio de alimentos a serem servidos no
evento. Durante o evento, atende participantes, convidados, imprensa,
autoridades e expositores. Monitora a recepção dos convidados e os
serviços prestados pelos contratados. Após o evento, coordena a
desmontagem e libera ordem de pagamento dos serviços prestados. Realiza
prestação de contas. Verifica satisfação do cliente. Pode aplicar
modelos de gestão e operar ferramentas e sistemas informatizados para
realizar estudo de custos; elaborar contratos; programar e efetuar o
cálculo do pagamento aos fornecedores e às equipes de trabalho; e
prestar contas ao cliente.

\begin{Form}
\textbf{Esta correspondência está adequada?} \\ 
\ChoiceMenu[radio,radiosymbol=\ding{52},bordercolor=0 0.5 0,name=cbovalid_ 354820 ]{}{CBO deve ficar=sim, \hspace{6em} CBO não fica aqui=nao}
\end{Form}

\textbf{CBO: 354825  ( Cerimonialista )}

Planeja o trabalho, identificando as preferências e a necessidade do
cliente ou autoridade e analisando as características do evento. Elabora
orçamento. Verifica local, data e horário do evento. Pode colaborar na
escolha de um local para a realização do evento. Pode utilizar
plataformas digitais, capazes de abrigar eventos e cerimônias on-line.
Contrata serviços, como decoração, fotografia, músicos, segurança,
limpeza, entre outros. Seleciona os prestadores de serviços de apoio a
eventos. Participa da preparação do cerimonial. Elabora o roteiro do
evento, respeitando as formalidades. Valida aspectos protocolares de
comunicação. Organiza solenidades e cerimônias, aplicando leis, normas e
princípios protocolares. Elabora lista de convidados e prepara os
convites. Elabora o plano de mesas. Orienta os participantes sobre as
vestimentas adequadas. Define o mestre de cerimônia. Seleciona os
protocolos adequados às diferentes cerimônias em que atua, estando à
frente do controle e da segurança dos convidados. Prospecta agenda de
autoridade para confirmar a possibilidade de participação em evento
formal. Reúne-se com a missão precursora, realizada após confirmação da
presença de autoridades, para verificar o local onde será realizada a
cerimônia e levantar todos os detalhes do acontecimento. Aplica normas
de cerimonial público e ordem geral de precedência, regras de composição
de mesas de autoridades, assim como as formas de tratamento e etiqueta.
Pode elaborar roteiro de viagens para autoridade e sua comitiva,
prestando assistência no embarque e no desembarque de aeronaves e
alugando carros para deslocamento até o local da cerimônia. Durante o
evento, supervisiona o cumprimento do roteiro. Inspeciona a conformidade
da prestação dos serviços contratados. Orienta garçons, recepcionistas e
outros profissionais envolvidos no evento, utilizando equipamentos de
comunicação interpessoal. Em solenidades e cerimônias, recepciona as
autoridades. Pode disponibilizar equipamentos multimídias e softwares
para apresentação de audiovisual. Após a cerimônia, providencia o
pagamento dos serviços contratados e realiza a prestação de contas para
o cliente ou autoridade.

\begin{Form}
\textbf{Esta correspondência está adequada?} \\ 
\ChoiceMenu[radio,radiosymbol=\ding{52},bordercolor=0 0.5 0,name=cbovalid_ 354825 ]{}{CBO deve ficar=sim, \hspace{6em} CBO não fica aqui=nao}
\end{Form}

\newpage
\section{ 13003 -- Técnico em Gastronomia }

\textbf{Perfil do Curso (CNCT)}

O Técnico em Gastronomia será habilitado para:

Coordenar a organização e a preparação do ambiente de trabalho da
cozinha. Monitorar o recebimento, a entrada, a saída e o armazenamento
de mercadorias em estoque. Supervisionar o pré-preparo e a aplicação de
técnicas de corte e cocção em alimentos. Preparar e finalizar produções
gastronômicas. Produzir e executar fichas técnicas operacionais para
produções gastronômicas. Colaborar com a elaboração e a revisão de
cardápios. Intermediar as relações entre as equipes de cozinha, salão e
bar, além de auxiliar na coordenação da equipe de cozinha.

Para atuação como Técnico em Gastronomia, são fundamentais:

Conhecimentos multidisciplinares sobre aspectos socioculturais e
econômicos relacionados à gastronomia, bem como conhecimentos técnicos
sobre utensílios, equipamentos e técnicas de preparação de alimentos.
Conhecimentos sobre práticas sustentáveis e manejo de resíduos e sobre o
cumprimento das legislações relacionadas à produção e à segurança dos
alimentos. Trabalho colaborativo, comunicação clara, comprometimento com
a sustentabilidade, atitude empreendedora, proatividade, criatividade e
flexibilidade para a solução de problemas e situações imprevistas.

\textbf{CBOs correspondidos e seus perfis (presentes no CNCT, e presentes na predição da IA)}

\textbf{CBO: 513205  ( Cozinheiro geral )}

Planeja a rotina de trabalho da cozinha, de acordo com os cardápios
elaborados e com as receitas definidas. Lista os ingredientes, de acordo
com o plano de produção e a capacidade de armazenamento. Seleciona
fornecedores e providencia a compra, considerando custo, qualidade e
disponibilidade no mercado. Recebe mercadorias, conferindo os itens e
suas características. Pré-higieniza as mercadorias recebidas e executa
seu armazenamento. Colabora na criação do cardápio, interpretando
pedidos de clientes. Participa do teste de novas receitas. Cozinha
legumes, verduras e cereais. Prepara entradas, saladas, fundos, molhos,
sopas e massas. Pré-prepara carnes, aves e peixes para cozimento,
cortando-os, limpando-os, pesando-os e separando-os de acordo com as
porções solicitadas. Segue métodos de cocção e padrões de qualidade para
cada tipo de alimento. Faz o empratamento das porções. Finaliza pratos,
respeitando princípios estéticos. Entrega produções culinárias, de
acordo com os pedidos e as normas de segurança de alimentos. Organiza o
ambiente de trabalho para elaboração de produções culinárias. Orienta a
higienização do ambiente, das bancadas, das louças, dos utensílios, das
máquinas e dos equipamentos da cozinha. Providencia o descarte de
resíduos e sobras. Mantém-se atualizado sobre as tendências e técnicas
do setor, identificando necessidade de novos utensílios, máquinas e
equipamentos. Organiza utensílios de trabalho. Verifica funcionamento
das máquinas e dos equipamentos de cozinha, providenciando a necessária
manutenção ou substituição. Zela pela conservação dos alimentos
estocados, providenciando as condições necessárias para evitar
deterioração e perdas. Pode distribuir tarefas para auxiliares,
definindo horários de execução e término de tarefas de acordo com
prioridades. Elabora fichas técnicas e auxilia na precificação de
cardápios. Comunica-se com a equipe, com maître e garçom e com clientes.

\begin{Form}
\textbf{Esta correspondência está adequada?} \\ 
\ChoiceMenu[radio,radiosymbol=\ding{52},bordercolor=0 0.5 0,name=cbovalid_ 513205 ]{}{CBO deve ficar=sim, \hspace{6em} CBO não fica aqui=nao}
\end{Form}

\newpage
\section{ 13004 -- Técnico em Guia de Turismo }

\textbf{Perfil do Curso (CNCT)}

O Técnico em Guia de Turismo será habilitado para:

Planejar e organizar a execução de roteiros e itinerários turísticos.
Conduzir e orientar visitantes na realização de traslados, passeios,
visitas e viagens. Prestar informações turísticas no contexto local,
regional e nacional. Intermediar as relações entre visitantes,
comunidade e prestadores de serviços turísticos. Prestar assistência aos
visitantes durante a realização dos roteiros e itinerários turísticos.

Para atuação como Técnico em Guia de Turismo, são fundamentais:

Conhecimentos multidisciplinares sobre aspectos socioculturais,
históricos, ambientais, geográficos, legais e econômicos, relacionados
aos roteiros e itinerários turísticos programados, bem como
conhecimentos técnicos relacionados à operação turística, marketing
pessoal e idiomas. Comunicação clara e empática, respeito à diversidade,
atenção à sustentabilidade dos produtos, atrativos e destinos
turísticos, atitude empreendedora, proatividade na tomada de decisões
táticas e operacionais relacionadas à atividade, criatividade e
flexibilidade para a solução de problemas e conflitos.

\textbf{CBOs correspondidos e seus perfis (presentes no CNCT, e presentes na predição da IA)}

\textbf{CBO: 511405  ( Guia de turismo )}

Estrutura roteiro e itinerário turístico, avaliando locais potenciais
para realização de turismo; selecionando itinerário; identificando
infraestrutura necessária para execução do novo roteiro; verificando
grau de dificuldade para acesso ao local ou realização de passeios;
coletando dados sobre atrativos -- tais como aspectos históricos,
culturais e naturais - do local ou da região; e avaliando potencial
turístico. Consolida informações do novo roteiro turístico, para
apresentá-lo às operadoras e às empresas especializadas em turismo.
Organiza plano de viagem e cronograma da programação do roteiro ou
itinerário, conforme demanda de operadora ou empresa de turismo. Planeja
e organiza a execução de roteiro turístico, verificando forma de
transporte -- avião, navio, transporte rodoviário ou transporte
ferroviário - ao local ou região do passeio; selecionando guia local
especializado; providenciando veículo para transporte local e
organizando seu trajeto e suas paradas técnicas; selecionando
prestadores de serviços -- hotéis, restaurantes, entre outros -- de
acordo com o perfil de visitantes ou grupos; e indicando passeios
opcionais. Elabora instruções sobre documentação exigida em visitas a
diversos locais do itinerário. Executa o roteiro turístico, conduzindo e
orientando grupo de pessoas na realização de traslados, visitas e
passeios previstos na programação. Confere documentos e passa a cuidar
da unidade do grupo pela lista de presença. Presta informações ao grupo
sobre os passeios a serem realizados, comentando programação geral;
explicando normas de conduta a serem observadas ao longo da excursão;
esclarecendo, em viagens internacionais, o câmbio de moeda no país; e
detalhando normas e regulamentos dos locais visitados no itinerário.
Salienta diferenças culturais existentes no local visitado. Orienta
sobre procedimentos de segurança. Informa sobre os equipamentos e
serviços de apoio em cada local do itinerário. Faz intermediação das
relações entre o grupo, a população local e os prestadores de serviços.
Pode traduzir textos e conversações durante a visita. Acompanha as
pessoas no período contratado, orientando o grupo nas atividades
previstas no roteiro de viagem. Alerta sobre apetrechos e indumentária
para as atividades do dia. Supervisiona refeições, inclusas ou não no
roteiro. Cumpre horário da programação, localizando pessoas do grupo
atrasadas para saída. Promove ações para integração e bom relacionamento
entre as pessoas, selecionando as estratégias de acordo com o perfil do
grupo. Apresenta guia local ao grupo, quando for o caso. Apresenta
pontos visitados, informando o grupo sobre infraestrutura e facilidades
do local, e descrevendo os atrativos culturais e naturais. Pode fazer
adaptações no roteiro de viagem, em decorrência de problemas na
localidade ou nos serviços turísticos contratados. Pode adequar o ritmo
do roteiro à capacidade física do grupo. Pode apresentar atividades
opcionais ao roteiro original para participantes que apresentem
restrições físicas para a atividade do dia. Assiste o grupo durante a
realização dos roteiros e itinerários turísticos, acompanhando pessoa
para solicitar, às autoridades locais, solução de problemas relacionados
à segurança e à documentação; providenciando assistência médica e
prestando primeiros-socorros em casos de emergência; e apoiando
processos de ``check-in'' e ``checkout'' em hotéis e aeroportos. Atua na
proteção do meio ambiente, orientando os visitantes para cumprimento das
normas ambientais durante os passeios. Realiza tarefas administrativas,
tais como confirmar voos; fazer reservas de atividades opcionais e
controlar ingressos, documentos e valores da excursão. Realiza prestação
de contas da excursão para agência de viagens. Encaminha, à agência, a
avaliação dos serviços prestados, feita pelo grupo que finalizou o
itinerário turístico. Elabora relatórios diário e final, incluindo
ocorrências durante o itinerário.

\begin{Form}
\textbf{Esta correspondência está adequada?} \\ 
\ChoiceMenu[radio,radiosymbol=\ding{52},bordercolor=0 0.5 0,name=cbovalid_ 511405 ]{}{CBO deve ficar=sim, \hspace{6em} CBO não fica aqui=nao}
\end{Form}

\newpage
\section{ 13005 -- Técnico em Hospedagem }

\textbf{Perfil do Curso (CNCT)}

O Técnico em Hospedagem será habilitado para:

Realizar atividades de recepção, reserva, governança, mensageria,
mordomia e conciergerie em meios de hospedagem. Prestar serviços de
atendimento e suporte aos hóspedes. Divulgar os serviços de hospedagem e
produtos turísticos. Supervisionar a manutenção de equipamentos e
estrutura física. Acompanhar e orientar procedimentos de higienização,
controle e arrumação das unidades habitacionais e dos espaços do
estabelecimento. Auxiliar na operacionalização de eventos, serviços,
alimentos e bebidas, articulando às necessidades dos hóspedes,
fornecedores e clientes.

Para atuação como Técnico em Hospedagem, são fundamentais:

Conhecimentos multidisciplinares sobre aspectos geográficos, históricos
e turísticos da região, bem como conhecimentos técnicos sobre
hospitalidade, classificação, estrutura, normas, procedimentos dos meios
de hospedagem e sistemas operacionais. Comunicação clara e cordial,
respeito à diversidade, atenção à sustentabilidade, trabalho
colaborativo, proatividade e flexibilidade para a solução de problemas e
conflitos.

\textbf{CBOs correspondidos e seus perfis (presentes no CNCT, e presentes na predição da IA)}

\textbf{CBO: 422120  ( Recepcionista de hotel )}

Planeja a rotina, organizando os materiais e os equipamentos de
trabalho, providenciando ambiente agradável e limpo na recepção e
verificando o funcionamento dos sistemas informatizados e de
comunicação. Imprime os relatórios de controle e as listagens de pessoas
e grupos com entrada ou saída prevista para o dia. Providencia o
atendimento às solicitações de reservas especiais. Verifica agenda
semanal de eventos. Participa de reuniões na troca de turnos. Recepciona
as visitas e as pessoas que vão se hospedar no hotel. Informa às visitas
sobre os serviços prestados pelo hotel, os preços e as condições para
hospedagem. Atende a pessoa ou o grupo que vai se hospedar no hotel,
realizando check-in. Efetua o cadastro, verifica a disponibilidade de
acomodações, confirma a reserva e fornece a chave do apartamento, quarto
ou outro local de hospedagem. Presta informações sobre recursos e normas
de funcionamento do hotel e sobre serviços de cofre, de alimentação e
bebidas e de lavanderia, entre outros. Pode liberar o uso de
estacionamento para o hóspede. Informa ao concierge sobre a chegada do
hóspede. Durante o período de hospedagem, presta informações ao hóspede
sobre acesso a Wi-Fi, pontos turísticos próximos, locais para câmbio de
moedas, realização de eventos e outras, quando solicitadas. Providencia
a prestação de atendimento médico ao hóspede. Auxilia na movimentação do
hóspede com dificuldade de locomoção. Pode receber e registrar chamadas
telefônicas e mensagens eletrônicas, transmitindo recados. Organiza a
distribuição de correspondências e de encomendas recebidas. No final do
período de hospedagem, realiza o checkout. Verifica as despesas do
hóspede em sistema informatizado e efetua acerto de contas, recebendo
pagamento e emitindo notas fiscais e recibos. Interage e comunica-se com
hóspedes e com os setores comercial, de gerência, de segurança, de
governança, de restaurantes, entre outros, do hotel. Pode comunicar-se e
prestar atendimento em língua estrangeira.

\begin{Form}
\textbf{Esta correspondência está adequada?} \\ 
\ChoiceMenu[radio,radiosymbol=\ding{52},bordercolor=0 0.5 0,name=cbovalid_ 422120 ]{}{CBO deve ficar=sim, \hspace{6em} CBO não fica aqui=nao}
\end{Form}

\textbf{CBO: 513110  ( Mordomo de hotelaria )}

Atende hóspedes, dando-lhes boas-vindas e apresentando-lhes apartamentos
e equipamentos disponíveis. Acompanha entradas e saídas dos hóspedes,
durante a estadia. Toma conhecimento das preferências dos hóspedes.
Providencia facilidades para utilização de equipamentos de informática.
Recepciona convidados e visitantes. Supervisiona a arrumação de
aposentos e demais áreas do hotel e verifica o estado de conservação de
móveis e equipamentos. Prepara aposentos personalizados. Verifica a
decoração das áreas. Controla a entrega e o recebimento de chaves.
Supervisiona a troca de apartamentos. Supervisiona a preparação das
refeições para os hóspedes, preenche comandas de alimentos e bebidas e
apronta coquetéis a serem servidos. Dispõe os hóspedes à mesa e serve os
alimentos e as bebidas. Pode atender pedidos especiais, como refeição
para hóspedes com restrições alimentares. Cuida do vestuário e objetos
dos hóspedes, podendo desfazer suas malas e arrumar vestuário e objetos
pessoais nos armários e em gavetas. Encaminha peças de vestuário para
reparos e ajustes e agiliza serviços de lavanderia, de passadoria e de
engraxate. Controla a guarda temporária de vestuário de hóspedes. Pode
prestar outros serviços - tais como reserva de sala de reuniões,
agendamento de consultas médicas, compra de remédios e aquisição de
entradas de teatro --, para propiciar bem-estar aos hóspedes. Executa
atividades administrativas, atualizando cadastro de perfis e históricos
de hóspedes; controlando livros de ocorrências; acompanhando registros
de entrada e saída de hóspedes; solicitando serviços de manutenção em
instalações, móveis e equipamentos; preparando e analisando relatórios.
Atende telefonemas, anota e transmite recados e mensagens, fazendo uso
de telefone fixo, de recursos de informática e de telefonia móvel. Pode
receber solicitações de hóspedes por e-mail. Coordena equipe, planejando
o trabalho dos empregados, dimensionando pessoal, estabelecendo rotinas,
atribuindo tarefas e orientando a execução dos serviços. Observa a
apresentação pessoal dos empregados e examina seus uniformes. Soluciona
situações de conflito na equipe. Administra a equipe de trabalho,
decidindo sobre admissão de novos empregados; fiscalizando o cumprimento
de horário de trabalho; programando escala de serviço e férias; e
decidindo sobre demissão de empregados subordinados. Indica membros da
equipe para promoções. Avalia a conduta e a qualidade dos serviços
realizados pelos membros da equipe. Pode promover treinamentos para
melhoria dos serviços prestados ou para o uso de novos equipamentos e
produtos. Zela pelo cumprimento das normas de segurança, prevenindo
acidentes e orientando a equipe para o uso de equipamentos de proteção
individual.

\begin{Form}
\textbf{Esta correspondência está adequada?} \\ 
\ChoiceMenu[radio,radiosymbol=\ding{52},bordercolor=0 0.5 0,name=cbovalid_ 513110 ]{}{CBO deve ficar=sim, \hspace{6em} CBO não fica aqui=nao}
\end{Form}

\textbf{CBO: 513115  ( Governanta de hotelaria )}

Recepciona os hóspedes, conhecendo suas preferências e providenciando os
serviços extras solicitados, como facilidades para utilização de
equipamento de informática nos aposentos. Planeja as atividades de
limpeza, de organização e de conservação dos aposentos e das demais
áreas, levando em consideração a necessidade de permanência e de
circulação dos hóspedes nas instalações. Supervisiona e orienta a
limpeza e a arrumação de aposentos e das demais áreas do hotel,
inspecionando a qualidade do trabalho realizado. Providencia aposentos
personalizados, para atender a exigência de hóspedes. Verifica o consumo
e a reposição de frigobar. Examina o estado de conservação de móveis e
equipamentos, providenciando serviços de reparação e manutenção. Pode
interditar apartamentos, para execução de reparos e manutenção nas
instalações. Supervisiona e orienta os serviços de lavanderia e
passadoria e as atividades do setor de costura, verificando a forma de
uso de produtos químicos, examinando os arremates de costura, observando
o processo de passar roupas e outros detalhes da execução dos serviços,
cuidando para que o hóspede receba suas roupas de acordo com suas
exigências. Controla entrada e saída de vestuário e enxovais lavados por
terceiros. Executa atividades administrativas, preparando relatórios,
recebendo auditores e atendendo fornecedores. Lança débito aos hóspedes
de serviços de lavanderia, passadoria e costura e de consumo do
frigobar. Controla objetos achados e perdidos. Acompanha registros de
entrada e saída de hóspedes. Controla o estoque de matérias-primas e
produtos para limpeza e conservação de instalações, para os serviços de
lavanderia e passadoria e para as atividades do setor de costura,
requisitando aquisição para reposição. Controla os custos das atividades
de governança, monitorando as despesas e considerando o orçamento da
área. Controla desperdícios, verificando o reaproveitamento de
materiais. Orienta o descarte de resíduos de acordo com as normas
ambientais. Coordena equipe, planejando o trabalho de camareiros, de
pessoal de lavanderia, passadoria e costura e de outros empregados que
realizam atividades relacionadas à governança. Dimensiona pessoal,
estabelece rotinas, atribui tarefas e orienta a execução dos serviços.
Observa a apresentação pessoal dos empregados e fornece seus uniformes.
Administra a equipe de trabalho, decidindo sobre admissão de novos
empregados; fiscalizando o cumprimento de horário de trabalho;
programando escala de serviço e férias; e decidindo sobre demissão de
empregados subordinados. Indica membros da equipe para promoções. Avalia
a conduta e a qualidade dos serviços realizados pelos membros da equipe.
Pode promover treinamentos para melhoria dos serviços prestados ou para
o uso de novos equipamentos e produtos. Utiliza recursos de novas
tecnologias para agilizar e aperfeiçoar seu trabalho, fazendo uso de
softwares para monitorar o andamento dos trabalhos de sua equipe,
controlar chegada de novas ordens de serviço em tempo real e informar à
equipe as preferências dos hóspedes para execução dos serviços. Zela
pelo cumprimento das normas de segurança, prevenindo acidentes e
orientando a equipe para o uso de equipamentos de proteção individual.

\begin{Form}
\textbf{Esta correspondência está adequada?} \\ 
\ChoiceMenu[radio,radiosymbol=\ding{52},bordercolor=0 0.5 0,name=cbovalid_ 513115 ]{}{CBO deve ficar=sim, \hspace{6em} CBO não fica aqui=nao}
\end{Form}

\newpage
\section{ 13006 -- Técnico em Lazer }

\textbf{Perfil do Curso (CNCT)}

O Técnico em Lazer será habilitado para:

Planejar atividades e programações de lazer para fins recreativos,
culturais e pedagógicos, de acordo com o público-alvo, recursos e
espaços disponíveis. Organizar e realizar atividades de lazer e
recreação de acordo com as necessidades do público. Aplicar técnicas de
recreação para a promoção da diversão, do lazer, da qualidade de vida e
do entretenimento.

Para atuação como Técnico em Lazer, são fundamentais:

Conhecimentos relacionados aos aspectos socioculturais, bem como
conhecimentos técnicos sobre hospitalidade e atividades de lazer,
recreação, programações culturais e esportivas para fins lúdicos, as
quais buscam garantir o bem-estar, a integridade e a segurança do
público. Comunicação clara e cordial, atuação de forma empática,
respeito às diversidades, trabalho colaborativo, proatividade,
criatividade e flexibilidade para solução de problemas e gestão de
conflitos.

\textbf{CBOs correspondidos e seus perfis (presentes no CNCT, e presentes na predição da IA)}

\textbf{CBO: 371410  ( Recreador )}

Atua geralmente em hotéis, clubes e complexos esportivos, associações
desportivas, navios, espaços de festas e eventos, parques temáticos,
colônias de férias, espaços culturais e em ambientes escolares. Elabora
programação de entretenimento e recreação, de acordo com o perfil do
público-alvo. Define objetivos, temáticas e sequência de atividades.
Programa atividades especiais para crianças, como contação de histórias,
dramatizações de situações e personagens, canto de músicas e cantigas,
danças, jogos e brincadeiras. Estabelece recursos necessários, definindo
espaço, instalações e equipamentos. Estima os recursos financeiros e
avalia a viabilidade de execução. Apresenta os resultados esperados da
programação. Pesquisa, requisita ou compra equipamentos e materiais.
Prepara, distribui e supervisiona o uso dos materiais e equipamentos,
garantindo as condições seguras e adequadas para a utilização. Recruta,
seleciona e prepara uma pessoa ou uma equipe para a execução das
atividades e para acolher diferentes públicos, inclusive pessoas com
deficiência. Supervisiona e coordena as atividades nas instalações
recreativas. Pode executar as atividades programadas, caracterizando-se
com maquiagens e fantasias, de acordo com a característica do público e
tipo de evento. Avalia as atividades realizadas, verificando a
satisfação dos participantes. Identifica, avalia e previne pessoas em
relação a situações de risco. Quando necessário, presta primeiros
socorros e encaminha participantes ao atendimento médico. Produz
cartazes e folhetos ou providencia outros meios de divulgação das
atividades. Registra as atividades em sistema informatizado. Elabora
relatório das atividades.

\begin{Form}
\textbf{Esta correspondência está adequada?} \\ 
\ChoiceMenu[radio,radiosymbol=\ding{52},bordercolor=0 0.5 0,name=cbovalid_ 371410 ]{}{CBO deve ficar=sim, \hspace{6em} CBO não fica aqui=nao}
\end{Form}

\newpage
\section{ 13008 -- Técnico em Serviços de Restaurante e Bar }

\textbf{Perfil do Curso (CNCT)}

O Técnico em Serviços de Restaurante e Bar será habilitado para:

Supervisionar o serviço de alimentos e bebidas no salão, no bar e em
eventos. Coordenar o atendimento ao cliente no estabelecimento e em
eventos. Coordenar equipes de serviço de salão e bar. Aplicar controles
operacionais em relação a vendas, equipamentos, utensílios e manutenção
da infraestrutura. Monitorar o recebimento, a entrada, a saída e o
armazenamento de mercadorias em estoque. Colaborar com a elaboração e a
revisão de cardápios.

Para atuação como Técnico em Serviços de Restaurante e Bar, são
fundamentais:

Conhecimentos técnicos relacionados à tipologia de serviços e eventos,
aos tipos de utensílios e equipamentos, à organização dos espaços, à
diferenciação de bebidas e legislações aplicadas aos estabelecimentos de
alimentação. Comunicação clara e cordial, trabalho colaborativo e
liderança de equipes, atenção à sustentabilidade, proatividade e
flexibilidade para a solução de problemas e gestão de conflitos.

\textbf{CBOs correspondidos e seus perfis (presentes no CNCT, e presentes na predição da IA)}

\textbf{CBO: 510130  ( Chefe de bar )}

Planeja as rotinas de trabalho em serviços de bar, no estabelecimento
especializado em servir bebidas ou no interior de estabelecimentos
maiores, tais como uma unidade do departamento de alimentos e bebidas de
hotéis ou uma área delimitada de restaurantes. Define as funções da
equipe de trabalho e descreve os procedimentos de execução dos serviços.
Programa o abastecimento do bar, pesquisando os preços de
matérias-primas e produtos. Atende clientes do bar, recepcionando-os e
orientando seus pedidos. Verifica o fundamento de reclamações dos
clientes. Controla a emissão de comandas. Realiza o controle do
suprimento de matérias-primas e produtos para o bar, monitorando o
consumo e realizando aquisição para reposição. Controla o recebimento de
mercadorias e supervisiona locais de armazenamento e acondicionamento
dos produtos adquiridos. Examina prazo de validade dos produtos.
Supervisiona a preparação de bebidas e petiscos. Pesquisa e cria novas
receitas. Avalia a execução dos serviços, analisando opiniões dos
clientes sobre produtos e serviços. Avalia o tempo de espera e a
qualidade na entrega dos produtos e serviços aos clientes. Avalia a
apresentação dos coquetéis, drinques, sucos e outras bebidas. Verifica o
estado de conservação das instalações do bar, realizando pequenos
reparos. Providencia serviços de manutenção, quando constatados
problemas maiores, acompanhando a execução e verificando o resultado dos
serviços. Monitora a realização dos serviços de dedetização e
desratização. Monitora o funcionamento dos equipamentos, realizando
pequenos reparos ou providenciando serviços de manutenção. Propõe a
substituição de equipamentos tecnologicamente ultrapassados por novos.
Realiza triagem de utensílios para descarte ou reparo. Coordena a equipe
de trabalho, dimensionando pessoal, atribuindo tarefas, delegando
responsabilidades e orientando a execução dos serviços. Controla o uso
de uniformes e a higiene pessoal da equipe. Soluciona situações de
conflito na equipe ou problemas entre a equipe e os clientes. Administra
recursos humanos, participando da seleção de pessoal, controlando
absenteísmo e estabelecendo escala de horários e folgas. Estabelece
rotinas de passagem de turnos. Supervisiona a equipe de trabalho,
determinando parâmetros para avaliação de execução de tarefas e
coletando dados sobre desempenho. Examina o tempo de execução e a
qualidade do trabalho da equipe na realização de tarefas. Realiza
reuniões de avaliação da equipe. Levanta as necessidades de treinamento
e providencia as atividades de capacitação. Instrui a equipe sobre
normas e regras do estabelecimento. Organiza treinamento para execução
de novas receitas, novos serviços e para uso de novos equipamentos e
produtos. Simula atividades para aprendizagem. Propõe programa de
treinamento externo para a equipe. Elabora relatórios de operação e de
avaliação, registrando ocorrências, expondo o consumo e a perda de
matérias-primas e produtos; elaborando documento com informações de
vendas; preparando informações sobre quebra e reposição de utensílios; e
apresentando o nível de satisfação dos clientes. Mantém a organização e
a limpeza do bar, supervisionando a aplicação de boas práticas de
higienização de instalações físicas e utensílios. Controla desperdícios,
com o reaproveitamento de material, sempre que possível. Orienta o
descarte de resíduos de acordo com as normas ambientais. Zela pela
segurança das pessoas, assegurando o cumprimento das normas;
identificando e eliminando as condições inseguras; orientando e
conscientizando a equipe para o uso de Equipamentos de Proteção
Individual (EPI).

\begin{Form}
\textbf{Esta correspondência está adequada?} \\ 
\ChoiceMenu[radio,radiosymbol=\ding{52},bordercolor=0 0.5 0,name=cbovalid_ 510130 ]{}{CBO deve ficar=sim, \hspace{6em} CBO não fica aqui=nao}
\end{Form}

\textbf{CBO: 510135  ( Maître )}

Planeja as rotinas de servir comidas preparadas e bebidas nos
restaurantes, nos eventos e em outros locais, definindo as funções da
equipe de trabalho e descrevendo os procedimentos para execução dos
serviços. Orienta os garçons em relação ao cardápio e às opções de
bebida, detalhando a composição dos pratos e expondo as características
de vinhos, licores e outras bebidas. Planeja operações para datas
especiais. Atende clientes, recepcionando-os e encaminhando-os à mesa
escolhida. Orienta os pedidos dos clientes, aplicando técnicas de
harmonização entre alimentos e bebidas. Atende reclamações dos clientes,
verificando seus fundamentos e procurando solucionar os problemas
apontados. Pode acompanhar visitações de clientes nas instalações do
restaurante. Coordena e supervisiona os serviços de mesa. Controla o
tempo para preparação de alimentos e bebidas e para entrega do pedido
aos clientes. Pode finalizar pratos. Avalia a execução dos serviços,
analisando opiniões dos clientes. Examina a preparação de ambientes onde
alimentos e bebidas são servidos. Analisa a apresentação dos pratos e
das bebidas. Verifica a aceitação de novos serviços e a utilização de
novos equipamentos e utensílios. Controla o suprimento de
matérias-primas e produtos no salão do restaurante ou em outro local
onde alimentos e bebidas são servidos. Verifica o consumo de
matérias-primas e produtos e requisita aquisição para reposição.
Controla recebimento de mercadorias e supervisiona locais de
armazenamento e acondicionamento de matérias-primas e produtos. Examina
o prazo de validade de produtos. Examina o estado de conservação das
instalações e monitora o funcionamento de equipamentos, executando
pequenos reparos e encaminhando requisição de serviços de manutenção.
Verifica os resultados de serviços de dedetização. Coordena a equipe de
trabalho, dimensionando pessoal, atribuindo tarefas, delegando
responsabilidades e orientando a execução dos serviços. Controla o uso
de uniformes e a higiene pessoal da equipe. Soluciona situações de
conflito na equipe e problemas entre a equipe e os clientes. Administra
recursos humanos, participando da seleção de pessoal, controlando
absenteísmo e estabelecendo escala de horários e folgas. Estabelece
rotinas de passagem de turnos. Supervisiona a equipe de trabalho,
determinando parâmetros para avaliação de execução de tarefas e
coletando dados sobre desempenho. Examina o tempo de execução e a
qualidade do trabalho da equipe na realização de tarefas. Realiza
reuniões de avaliação da equipe. Levanta as necessidades de treinamento
e providencia as atividades de capacitação. Instrui a equipe sobre
normas e regras do estabelecimento. Organiza treinamento para execução
de novos serviços e para uso de novos equipamentos e produtos. Simula
atividades para aprendizagem. Propõe programa de treinamento externo
para a equipe. Elabora relatórios de operação e de avaliação,
registrando ocorrências; expondo o consumo e a perda de matérias-primas
de alimentos e bebidas; relatando a avaliação de cardápio; preparando
informações sobre quebra e reposição de utensílios; e apresentando o
nível de satisfação dos clientes. Elabora relatório estatístico de
vendas do restaurante. Conserva o salão do restaurante limpo e
organizado, supervisionando a aplicação de boas práticas de higienização
de instalações físicas e utensílios. Controla desperdícios, com o
reaproveitamento de material, sempre que possível. Orienta o descarte de
resíduos de acordo com as normas ambientais. Zela pela segurança das
pessoas, assegurando o cumprimento das normas; identificando e
eliminando as condições inseguras; orientando e conscientizando a equipe
para o uso de Equipamentos de Proteção Individual (EPI).

\begin{Form}
\textbf{Esta correspondência está adequada?} \\ 
\ChoiceMenu[radio,radiosymbol=\ding{52},bordercolor=0 0.5 0,name=cbovalid_ 510135 ]{}{CBO deve ficar=sim, \hspace{6em} CBO não fica aqui=nao}
\end{Form}

\textbf{CBO: 513405  ( Garçom )}

Organiza o trabalho, examinando o estoque de bebidas e tomando
conhecimento do cardápio do dia. Recepciona o cliente, acompanhando-o à
mesa. Entrega o cardápio ao cliente, descrevendo pratos e prestando
outras informações solicitadas. Consulta o cliente sobre suas
preferências e apresenta sugestões, para auxiliá-lo na escolha dos
pratos e bebidas. Registra o pedido - em comanda eletrônica ou manual -,
encaminhando-o para cozinha e bar. Traz couvert e serve aperitivos.
Administra fluxo de pratos entre cozinha e mesas. Serve os pratos a cada
cliente, podendo fatiar carnes e destrinchar peixe. Pode fazer a
finalização dos pratos, quando for o caso. Pode preparar bebidas ou
encaminhar os pedidos ao barman. Serve vinho, café, licores e outras
bebidas. Atende solicitações específicas do cliente, buscando mantê-lo
satisfeito. Traz a conta quando for solicitada, podendo realizar o
recebimento do valor devido. Organiza a estrutura de apoio, limpando e
organizando mesas e cadeiras no início e ao término do serviço do
estabelecimento, selecionando copos ou taças, colocando pratos e
talheres, repondo aparadores e guéridons (mesas auxiliares), entre
outras providências. Higieniza utensílios e equipamentos. Limpa bandejas
e carrinhos. Mantém o salão abastecido, administrando e solicitando
reposição de produtos. Recebe as mercadorias, conferindo os itens
solicitados. Segue procedimentos para descarte de embalagens, sobras e
resíduos.

\begin{Form}
\textbf{Esta correspondência está adequada?} \\ 
\ChoiceMenu[radio,radiosymbol=\ding{52},bordercolor=0 0.5 0,name=cbovalid_ 513405 ]{}{CBO deve ficar=sim, \hspace{6em} CBO não fica aqui=nao}
\end{Form}

\textbf{CBO: 513415  ( Cumim )}

Auxilia na organização do salão, montando mesa com pratos, talheres,
copos ou taças, recipiente com sal e pimenta, guardanapos e outros
itens. Pode auxiliar na decoração das mesas. Auxilia na montagem de
balcão ou carrinho, com alimentos ou bebidas. Pode recepcionar e
acompanhar o cliente à mesa, entregando-lhe o cardápio. Pode anotar os
pedidos - manualmente ou por meio eletrônico -, encaminhando-os para a
cozinha e para o bar. Auxilia no atendimento ao cliente, servindo
aperitivos, entrada, pratos e bebidas, para agilizar o serviço. Pode
preparar sucos e cafés. Atende solicitações do cliente, buscando
mantê-lo satisfeito. Traz a conta quando for solicitada, podendo
realizar o recebimento do valor devido. Arruma a mesa para um novo
cliente, retirando a louça usada, limpando a mesa e trocando toalha e
acessórios. Recolhe os utensílios usados, levando-os para a copa. Ao
final do serviço de atendimento aos clientes, auxilia na arrumação do
espaço, limpando e arranjando mesas, tirando toalhas, guardando produtos
e desmontando carrinho e balcão. Reabastece o suprimento de toalhas,
guardanapos, talheres, copos e pratos limpos. Auxilia na higienização de
utensílios e equipamentos. Limpa bandejas, balcão, bancada e carrinho.
Separa e descarta lixo, seguindo os procedimentos adotados no
estabelecimento.

\begin{Form}
\textbf{Esta correspondência está adequada?} \\ 
\ChoiceMenu[radio,radiosymbol=\ding{52},bordercolor=0 0.5 0,name=cbovalid_ 513415 ]{}{CBO deve ficar=sim, \hspace{6em} CBO não fica aqui=nao}
\end{Form}

``` \footnotetext{Fonte: CNCT e QBQ, 17/07/2025}

\end{document}
